% Options for packages loaded elsewhere
% Options for packages loaded elsewhere
\PassOptionsToPackage{unicode}{hyperref}
\PassOptionsToPackage{hyphens}{url}
\PassOptionsToPackage{dvipsnames,svgnames,x11names}{xcolor}
%
\documentclass[
  letterpaper,
  DIV=11,
  numbers=noendperiod]{scrreprt}
\usepackage{xcolor}
\usepackage{amsmath,amssymb}
\setcounter{secnumdepth}{5}
\usepackage{iftex}
\ifPDFTeX
  \usepackage[T1]{fontenc}
  \usepackage[utf8]{inputenc}
  \usepackage{textcomp} % provide euro and other symbols
\else % if luatex or xetex
  \usepackage{unicode-math} % this also loads fontspec
  \defaultfontfeatures{Scale=MatchLowercase}
  \defaultfontfeatures[\rmfamily]{Ligatures=TeX,Scale=1}
\fi
\usepackage{lmodern}
\ifPDFTeX\else
  % xetex/luatex font selection
\fi
% Use upquote if available, for straight quotes in verbatim environments
\IfFileExists{upquote.sty}{\usepackage{upquote}}{}
\IfFileExists{microtype.sty}{% use microtype if available
  \usepackage[]{microtype}
  \UseMicrotypeSet[protrusion]{basicmath} % disable protrusion for tt fonts
}{}
\makeatletter
\@ifundefined{KOMAClassName}{% if non-KOMA class
  \IfFileExists{parskip.sty}{%
    \usepackage{parskip}
  }{% else
    \setlength{\parindent}{0pt}
    \setlength{\parskip}{6pt plus 2pt minus 1pt}}
}{% if KOMA class
  \KOMAoptions{parskip=half}}
\makeatother
% Make \paragraph and \subparagraph free-standing
\makeatletter
\ifx\paragraph\undefined\else
  \let\oldparagraph\paragraph
  \renewcommand{\paragraph}{
    \@ifstar
      \xxxParagraphStar
      \xxxParagraphNoStar
  }
  \newcommand{\xxxParagraphStar}[1]{\oldparagraph*{#1}\mbox{}}
  \newcommand{\xxxParagraphNoStar}[1]{\oldparagraph{#1}\mbox{}}
\fi
\ifx\subparagraph\undefined\else
  \let\oldsubparagraph\subparagraph
  \renewcommand{\subparagraph}{
    \@ifstar
      \xxxSubParagraphStar
      \xxxSubParagraphNoStar
  }
  \newcommand{\xxxSubParagraphStar}[1]{\oldsubparagraph*{#1}\mbox{}}
  \newcommand{\xxxSubParagraphNoStar}[1]{\oldsubparagraph{#1}\mbox{}}
\fi
\makeatother

\usepackage{color}
\usepackage{fancyvrb}
\newcommand{\VerbBar}{|}
\newcommand{\VERB}{\Verb[commandchars=\\\{\}]}
\DefineVerbatimEnvironment{Highlighting}{Verbatim}{commandchars=\\\{\}}
% Add ',fontsize=\small' for more characters per line
\usepackage{framed}
\definecolor{shadecolor}{RGB}{241,243,245}
\newenvironment{Shaded}{\begin{snugshade}}{\end{snugshade}}
\newcommand{\AlertTok}[1]{\textcolor[rgb]{0.68,0.00,0.00}{#1}}
\newcommand{\AnnotationTok}[1]{\textcolor[rgb]{0.37,0.37,0.37}{#1}}
\newcommand{\AttributeTok}[1]{\textcolor[rgb]{0.40,0.45,0.13}{#1}}
\newcommand{\BaseNTok}[1]{\textcolor[rgb]{0.68,0.00,0.00}{#1}}
\newcommand{\BuiltInTok}[1]{\textcolor[rgb]{0.00,0.23,0.31}{#1}}
\newcommand{\CharTok}[1]{\textcolor[rgb]{0.13,0.47,0.30}{#1}}
\newcommand{\CommentTok}[1]{\textcolor[rgb]{0.37,0.37,0.37}{#1}}
\newcommand{\CommentVarTok}[1]{\textcolor[rgb]{0.37,0.37,0.37}{\textit{#1}}}
\newcommand{\ConstantTok}[1]{\textcolor[rgb]{0.56,0.35,0.01}{#1}}
\newcommand{\ControlFlowTok}[1]{\textcolor[rgb]{0.00,0.23,0.31}{\textbf{#1}}}
\newcommand{\DataTypeTok}[1]{\textcolor[rgb]{0.68,0.00,0.00}{#1}}
\newcommand{\DecValTok}[1]{\textcolor[rgb]{0.68,0.00,0.00}{#1}}
\newcommand{\DocumentationTok}[1]{\textcolor[rgb]{0.37,0.37,0.37}{\textit{#1}}}
\newcommand{\ErrorTok}[1]{\textcolor[rgb]{0.68,0.00,0.00}{#1}}
\newcommand{\ExtensionTok}[1]{\textcolor[rgb]{0.00,0.23,0.31}{#1}}
\newcommand{\FloatTok}[1]{\textcolor[rgb]{0.68,0.00,0.00}{#1}}
\newcommand{\FunctionTok}[1]{\textcolor[rgb]{0.28,0.35,0.67}{#1}}
\newcommand{\ImportTok}[1]{\textcolor[rgb]{0.00,0.46,0.62}{#1}}
\newcommand{\InformationTok}[1]{\textcolor[rgb]{0.37,0.37,0.37}{#1}}
\newcommand{\KeywordTok}[1]{\textcolor[rgb]{0.00,0.23,0.31}{\textbf{#1}}}
\newcommand{\NormalTok}[1]{\textcolor[rgb]{0.00,0.23,0.31}{#1}}
\newcommand{\OperatorTok}[1]{\textcolor[rgb]{0.37,0.37,0.37}{#1}}
\newcommand{\OtherTok}[1]{\textcolor[rgb]{0.00,0.23,0.31}{#1}}
\newcommand{\PreprocessorTok}[1]{\textcolor[rgb]{0.68,0.00,0.00}{#1}}
\newcommand{\RegionMarkerTok}[1]{\textcolor[rgb]{0.00,0.23,0.31}{#1}}
\newcommand{\SpecialCharTok}[1]{\textcolor[rgb]{0.37,0.37,0.37}{#1}}
\newcommand{\SpecialStringTok}[1]{\textcolor[rgb]{0.13,0.47,0.30}{#1}}
\newcommand{\StringTok}[1]{\textcolor[rgb]{0.13,0.47,0.30}{#1}}
\newcommand{\VariableTok}[1]{\textcolor[rgb]{0.07,0.07,0.07}{#1}}
\newcommand{\VerbatimStringTok}[1]{\textcolor[rgb]{0.13,0.47,0.30}{#1}}
\newcommand{\WarningTok}[1]{\textcolor[rgb]{0.37,0.37,0.37}{\textit{#1}}}

\usepackage{longtable,booktabs,array}
\newcounter{none} % for unnumbered tables
\usepackage{calc} % for calculating minipage widths
% Correct order of tables after \paragraph or \subparagraph
\usepackage{etoolbox}
\makeatletter
\patchcmd\longtable{\par}{\if@noskipsec\mbox{}\fi\par}{}{}
\makeatother
% Allow footnotes in longtable head/foot
\IfFileExists{footnotehyper.sty}{\usepackage{footnotehyper}}{\usepackage{footnote}}
\makesavenoteenv{longtable}
\usepackage{graphicx}
\makeatletter
\newsavebox\pandoc@box
\newcommand*\pandocbounded[1]{% scales image to fit in text height/width
  \sbox\pandoc@box{#1}%
  \Gscale@div\@tempa{\textheight}{\dimexpr\ht\pandoc@box+\dp\pandoc@box\relax}%
  \Gscale@div\@tempb{\linewidth}{\wd\pandoc@box}%
  \ifdim\@tempb\p@<\@tempa\p@\let\@tempa\@tempb\fi% select the smaller of both
  \ifdim\@tempa\p@<\p@\scalebox{\@tempa}{\usebox\pandoc@box}%
  \else\usebox{\pandoc@box}%
  \fi%
}
% Set default figure placement to htbp
\def\fps@figure{htbp}
\makeatother





\setlength{\emergencystretch}{3em} % prevent overfull lines

\providecommand{\tightlist}{%
  \setlength{\itemsep}{0pt}\setlength{\parskip}{0pt}}



 


\KOMAoption{captions}{tableheading}
\makeatletter
\@ifpackageloaded{tcolorbox}{}{\usepackage[skins,breakable]{tcolorbox}}
\@ifpackageloaded{fontawesome5}{}{\usepackage{fontawesome5}}
\definecolor{quarto-callout-color}{HTML}{909090}
\definecolor{quarto-callout-note-color}{HTML}{0758E5}
\definecolor{quarto-callout-important-color}{HTML}{CC1914}
\definecolor{quarto-callout-warning-color}{HTML}{EB9113}
\definecolor{quarto-callout-tip-color}{HTML}{00A047}
\definecolor{quarto-callout-caution-color}{HTML}{FC5300}
\definecolor{quarto-callout-color-frame}{HTML}{acacac}
\definecolor{quarto-callout-note-color-frame}{HTML}{4582ec}
\definecolor{quarto-callout-important-color-frame}{HTML}{d9534f}
\definecolor{quarto-callout-warning-color-frame}{HTML}{f0ad4e}
\definecolor{quarto-callout-tip-color-frame}{HTML}{02b875}
\definecolor{quarto-callout-caution-color-frame}{HTML}{fd7e14}
\makeatother
\makeatletter
\@ifpackageloaded{bookmark}{}{\usepackage{bookmark}}
\makeatother
\makeatletter
\@ifpackageloaded{caption}{}{\usepackage{caption}}
\AtBeginDocument{%
\ifdefined\contentsname
  \renewcommand*\contentsname{Table of contents}
\else
  \newcommand\contentsname{Table of contents}
\fi
\ifdefined\listfigurename
  \renewcommand*\listfigurename{List of Figures}
\else
  \newcommand\listfigurename{List of Figures}
\fi
\ifdefined\listtablename
  \renewcommand*\listtablename{List of Tables}
\else
  \newcommand\listtablename{List of Tables}
\fi
\ifdefined\figurename
  \renewcommand*\figurename{Figure}
\else
  \newcommand\figurename{Figure}
\fi
\ifdefined\tablename
  \renewcommand*\tablename{Table}
\else
  \newcommand\tablename{Table}
\fi
}
\@ifpackageloaded{float}{}{\usepackage{float}}
\floatstyle{ruled}
\@ifundefined{c@chapter}{\newfloat{codelisting}{h}{lop}}{\newfloat{codelisting}{h}{lop}[chapter]}
\floatname{codelisting}{Listing}
\newcommand*\listoflistings{\listof{codelisting}{List of Listings}}
\usepackage{amsthm}
\theoremstyle{remark}
\AtBeginDocument{\renewcommand*{\proofname}{Proof}}
\newtheorem*{remark}{Remark}
\newtheorem*{solution}{Solution}
\newtheorem{refremark}{Remark}[chapter]
\newtheorem{refsolution}{Solution}[chapter]
\makeatother
\makeatletter
\makeatother
\makeatletter
\@ifpackageloaded{caption}{}{\usepackage{caption}}
\@ifpackageloaded{subcaption}{}{\usepackage{subcaption}}
\makeatother
\usepackage{bookmark}
\IfFileExists{xurl.sty}{\usepackage{xurl}}{} % add URL line breaks if available
\urlstyle{same}
\hypersetup{
  pdftitle={Foundations of Probability and Statistical Theory},
  pdfauthor={He Wang},
  colorlinks=true,
  linkcolor={blue},
  filecolor={Maroon},
  citecolor={Blue},
  urlcolor={Blue},
  pdfcreator={LaTeX via pandoc}}


\title{Foundations of Probability and Statistical Theory}
\usepackage{etoolbox}
\makeatletter
\providecommand{\subtitle}[1]{% add subtitle to \maketitle
  \apptocmd{\@title}{\par {\large #1 \par}}{}{}
}
\makeatother
\subtitle{A Modern Introduction with Computation, Simulation, and
Inference}
\author{He Wang}
\date{2025-12-31}
\begin{document}
\maketitle

\renewcommand*\contentsname{Table of contents}
{
\setcounter{tocdepth}{2}
\tableofcontents
}

\bookmarksetup{startatroot}

\chapter{Foundations of Probability and Statistical
Theory}\label{foundations-of-probability-and-statistical-theory}

A Modern Introduction with Computation, Simulation, and Inference

\hfill\break

\bookmarksetup{startatroot}

\chapter*{Welcome}\label{welcome}
\addcontentsline{toc}{chapter}{Welcome}

\markboth{Welcome}{Welcome}

Welcome to the online book for \textbf{MATH 5010: Foundations of
Statistical Theory \& Probability}.

This book is designed for graduate students who want a rigorous but
usable foundation in probability, mathematical statistics, and
statistical inference. The central goal is to connect three ways of
learning statistics:

\begin{enumerate}
\def\labelenumi{\arabic{enumi}.}
\tightlist
\item
  \textbf{Theory}: definitions, theorems, derivations, and proofs.
\item
  \textbf{Computation}: Python labs, simulation, Monte Carlo methods,
  and numerical verification.
\item
  \textbf{Interactive learning}: visual HTML modules that let students
  explore concepts dynamically.
\end{enumerate}

The book begins with probability theory and random variables, then
builds toward convergence theory, sampling distributions, sufficient
statistics, point estimation, hypothesis testing, interval estimation,
ANOVA, and Bayesian inference.

\section{How to use this book}\label{how-to-use-this-book}

Each main chapter is written as a lecture-ready Quarto page. Most
chapters are paired with an interactive HTML activity and a Python
notebook lab.

A good workflow for students is:

\begin{enumerate}
\def\labelenumi{\arabic{enumi}.}
\tightlist
\item
  read the chapter notes before class;
\item
  work through the examples during class;
\item
  use the interactive HTML page to visualize the main ideas;
\item
  open the notebook in Google Colab and run the simulations;
\item
  return to the practice problems and solutions for review.
\end{enumerate}

A good workflow for instructors is:

\begin{enumerate}
\def\labelenumi{\arabic{enumi}.}
\tightlist
\item
  use the chapter page as lecture notes;
\item
  use the examples as board work or in-class activities;
\item
  assign selected practice problems as homework or discussion problems;
\item
  use the notebook labs for computational reinforcement;
\item
  use the HTML modules for visual explanation and student exploration.
\end{enumerate}

\section{Book map}\label{book-map}

\subsection{Part I: Probability
foundations}\label{part-i-probability-foundations}

{\def\LTcaptype{none} % do not increment counter
\begin{longtable}[]{@{}
  >{\raggedleft\arraybackslash}p{(\linewidth - 4\tabcolsep) * \real{0.4000}}
  >{\raggedright\arraybackslash}p{(\linewidth - 4\tabcolsep) * \real{0.3000}}
  >{\raggedright\arraybackslash}p{(\linewidth - 4\tabcolsep) * \real{0.3000}}@{}}
\toprule\noalign{}
\begin{minipage}[b]{\linewidth}\raggedleft
Chapter
\end{minipage} & \begin{minipage}[b]{\linewidth}\raggedright
Topic
\end{minipage} & \begin{minipage}[b]{\linewidth}\raggedright
Main purpose
\end{minipage} \\
\midrule\noalign{}
\endhead
\bottomrule\noalign{}
\endlastfoot
1 & \href{chapters/chapter-01-basic-probability.qmd}{Basic Probability}
& Build the language of probability: sample spaces, events, probability
rules, conditioning, independence, Bayes theorem, and paradoxes. \\
2 & \href{chapters/chapter-02-random-variables-distributions.qmd}{Random
Variables and Distributions} & Introduce random variables, CDFs, PMFs,
PDFs, mixed distributions, and common distribution families. \\
3 & \href{chapters/chapter-03-joint-conditional-probability.qmd}{Joint
and Conditional Probability} & Study joint, marginal, and conditional
distributions; independence; covariance; correlation; and Bayes rules in
multivariate settings. \\
4 &
\href{chapters/chapter-04-expectations-moments-mgf.qmd}{Expectations,
Moments, and MGFs} & Develop expected value, variance, moments,
moment-generating functions, and multivariate moment tools. \\
5 & \href{chapters/chapter-05-transformations.qmd}{Transformations} &
Learn CDF methods, change of variables, Jacobians, convolution,
products, quotients, extrema, and polar transformations. \\
6 & \href{chapters/chapter-06-conditional-expectations.qmd}{Conditional
Expectations} & Use conditional distributions and conditional
expectation to compute means, variances, Bayesian updates, and
first-step decompositions. \\
6 Extra &
\href{chapters/chapter-06-extra-multinomial-multinormal.qmd}{Multinomial
and Multinormal Models} & Extend probability tools to multinomial,
Dirichlet, multivariate normal, conditional Gaussian, and linear
Gaussian models. \\
7 & \href{chapters/chapter-07-inequalities-identities.qmd}{Inequalities
and Identities} & Study Markov, Chebyshev, Chernoff, Hoeffding, Jensen,
Cauchy-Schwarz, and other core probability inequalities. \\
8 & \href{chapters/chapter-08-sampling-order-statistics.qmd}{Sampling
and Order Statistics} & Introduce random samples, sample mean and
variance, normal sampling theory, chi-square, \(t\), \(F\), and order
statistics. \\
9 & \href{chapters/chapter-09-convergence-theory.qmd}{Convergence
Theory} & Study convergence in distribution, probability, mean, almost
surely, LLN, CLT, continuous mapping, and Slutsky theorem. \\
10 & \href{chapters/chapter-10-monte-carlo-sampling.qmd}{Monte Carlo
Sampling} & Use simulation, Monte Carlo integration, inverse transform
sampling, rejection sampling, importance sampling, and SIR. \\
\end{longtable}
}

\subsection{Part II: Statistical
inference}\label{part-ii-statistical-inference}

{\def\LTcaptype{none} % do not increment counter
\begin{longtable}[]{@{}
  >{\raggedleft\arraybackslash}p{(\linewidth - 4\tabcolsep) * \real{0.4000}}
  >{\raggedright\arraybackslash}p{(\linewidth - 4\tabcolsep) * \real{0.3000}}
  >{\raggedright\arraybackslash}p{(\linewidth - 4\tabcolsep) * \real{0.3000}}@{}}
\toprule\noalign{}
\begin{minipage}[b]{\linewidth}\raggedleft
Chapter
\end{minipage} & \begin{minipage}[b]{\linewidth}\raggedright
Topic
\end{minipage} & \begin{minipage}[b]{\linewidth}\raggedright
Main purpose
\end{minipage} \\
\midrule\noalign{}
\endhead
\bottomrule\noalign{}
\endlastfoot
11 & \href{chapters/chapter-11-sufficient-statistics.qmd}{Sufficient
Statistics} & Understand statistics, estimators, sufficiency,
factorization, minimal sufficiency, complete sufficiency, and the
likelihood principle. \\
12 & \href{chapters/chapter-12-point-estimation1.qmd}{Point Estimation
I} & Derive estimators using method of moments, maximum likelihood,
Bayesian estimation, conjugacy, and MAP estimation. \\
13 & \href{chapters/chapter-13-point-estimation2.qmd}{Point Estimation
II} & Evaluate estimators using bias, variance, MSE, Fisher information,
Cramér-Rao bounds, Rao-Blackwell, and UMVUE theory. \\
14 & \href{chapters/chapter-14-hypothesis-tests1.qmd}{Hypothesis Tests
I} & Introduce null and alternative hypotheses, Type I/II errors,
\(p\)-values, likelihood ratio tests, Bayesian tests, and Neyman-Pearson
ideas. \\
15 & \href{chapters/chapter-15-hypothesis-tests2.qmd}{Hypothesis Tests
II} & Study power functions, size, unbiased tests, UMP tests, binomial
and normal examples, and risk-based testing. \\
16 & \href{chapters/chapter-16-interval-estimation1.qmd}{Interval
Estimation I} & Construct confidence intervals using pivots, test
inversion, normal theory, binomial, uniform, exponential, and Bayesian
credible intervals. \\
17 & \href{chapters/chapter-17-interval-estimation2.qmd}{Interval
Estimation II} & Study interval optimality, coverage, shortest
intervals, UMA confidence sets, HPD intervals, and risk of confidence
sets. \\
18 & \href{chapters/chapter-18-anova1.qmd}{ANOVA I} & Develop one-way
ANOVA, treatment effects, contrasts, sums of squares, ANOVA tables, and
the \(F\)-test. \\
19 & \href{chapters/chapter-19-anova2-applications.qmd}{ANOVA II and
Applications} & Apply ANOVA to real examples, multiple testing,
experimental design, randomized block designs, and Python workflows. \\
20 & \href{chapters/chapter-20-bayesian-inference.qmd}{Bayesian
Inference} & Unify Bayesian estimation, testing, credible intervals, HPD
intervals, conjugate models, Bayes risk, and Bayesian workflow. \\
\end{longtable}
}

\section{Labs and interactive pages}\label{labs-and-interactive-pages}

The computational and interactive materials are summarized on a separate
page:

\href{labs-summary.qmd}{Open the labs and interactive modules summary}

The labs are stored in the \texttt{labs/} folder and can be opened
directly in Google Colab from the GitHub repository. The interactive
pages are stored in the \texttt{htmls/} folder and are intended for
visual exploration before, during, or after lecture.

\bookmarksetup{startatroot}

\chapter*{Introduction}\label{introduction}
\addcontentsline{toc}{chapter}{Introduction}

\markboth{Introduction}{Introduction}

This book supports \textbf{MATH 5010: Foundations of Statistical Theory
\& Probability}. It is written as a complete course companion: lecture
notes, worked examples, practice problems, interactive visualizations,
and Python labs are organized around the same sequence of topics.

The course has two connected parts. The first part develops probability
theory: probability spaces, random variables, distributions,
expectation, transformations, conditional expectation, inequalities,
random samples, order statistics, convergence, and simulation. The
second part develops statistical inference: sufficient statistics, point
estimation, estimator evaluation, hypothesis testing, confidence
intervals, ANOVA, and Bayesian inference.

The guiding idea is that statistical theory should be learned in three
languages at once:

\begin{itemize}
\tightlist
\item
  the \textbf{mathematical language} of probability models, theorems,
  and derivations;
\item
  the \textbf{computational language} of simulation, numerical
  approximation, and Python implementation;
\item
  the \textbf{statistical language} of estimation, uncertainty, testing,
  modeling, and interpretation.
\end{itemize}

\section*{Course philosophy}\label{course-philosophy}
\addcontentsline{toc}{section}{Course philosophy}

\markright{Course philosophy}

Probability and statistics are not only collections of formulas. They
are ways of thinking about uncertainty.

A probability model starts with a sample space, events, random
variables, and distributions. From these objects we define expectations,
variances, conditional probabilities, and limiting behavior. Statistical
inference then reverses the direction: we observe data and try to learn
about the unknown model that may have generated it.

This book repeatedly moves through the same cycle:

\begin{enumerate}
\def\labelenumi{\arabic{enumi}.}
\tightlist
\item
  \textbf{Start with a model.} Specify random variables, parameters,
  assumptions, and the data-generating process.
\item
  \textbf{Derive the theory.} Compute distributions, expectations,
  likelihoods, test statistics, intervals, or posterior distributions.
\item
  \textbf{Simulate the model.} Use Python to check whether the theory
  matches repeated random experiments.
\item
  \textbf{Interpret the result.} Explain what the formula, estimate,
  test, or interval means in plain language.
\end{enumerate}

\section*{Learning outcomes}\label{learning-outcomes}
\addcontentsline{toc}{section}{Learning outcomes}

\markright{Learning outcomes}

After working through this book, students should be able to:

\begin{itemize}
\tightlist
\item
  define probability spaces, events, random variables, distributions,
  and conditional distributions;
\item
  compute probabilities using complements, partitions, total
  probability, Bayes theorem, and independence;
\item
  work with PMFs, PDFs, CDFs, joint distributions, marginal
  distributions, and conditional distributions;
\item
  compute expectations, variances, moments, covariance, correlation, and
  moment-generating functions;
\item
  transform random variables using CDF methods, Jacobians, convolution,
  and multivariate changes of variables;
\item
  use conditional expectation, conditional variance, and first-step
  analysis;
\item
  apply probability inequalities and convergence theorems, including LLN
  and CLT;
\item
  derive sampling distributions and order statistic distributions;
\item
  use Monte Carlo methods for integration, simulation, and
  approximation;
\item
  identify sufficient statistics and apply the Neyman-Fisher
  factorization theorem;
\item
  derive estimators using method of moments, maximum likelihood, MAP,
  and posterior means;
\item
  evaluate estimators using bias, variance, MSE, efficiency, Fisher
  information, and Cramér-Rao bounds;
\item
  construct and interpret hypothesis tests, power functions,
  \(p\)-values, and likelihood ratio tests;
\item
  construct and interpret confidence intervals and Bayesian credible
  intervals;
\item
  use ANOVA models, contrasts, sums of squares, and \(F\)-tests;
\item
  explain Bayesian inference through priors, likelihoods, posteriors,
  Bayes estimators, Bayesian tests, and credible sets;
\item
  use Python notebooks to simulate, visualize, and communicate
  statistical results.
\end{itemize}

\section*{Organization of the book}\label{organization-of-the-book}
\addcontentsline{toc}{section}{Organization of the book}

\markright{Organization of the book}

The chapters are arranged to follow the logical development of
statistical theory.

\subsection*{Probability theory}\label{probability-theory}
\addcontentsline{toc}{subsection}{Probability theory}

Chapters 1--10 build the probability foundation. These chapters begin
with basic probability and move through random variables, multivariate
distributions, moments, transformations, conditioning, inequalities,
sampling theory, convergence, and Monte Carlo methods.

The main transition is from finite probability calculations to random
variables and then to limiting approximations. By the end of Chapter 10,
students should be comfortable moving between exact probability
formulas, asymptotic approximations, and simulation-based
approximations.

\subsection*{Statistical inference}\label{statistical-inference}
\addcontentsline{toc}{subsection}{Statistical inference}

Chapters 11--20 build inference from the probability foundation. The
sequence begins with sufficient statistics and likelihood, then develops
point estimation, estimator evaluation, hypothesis testing, interval
estimation, ANOVA, and Bayesian inference.

The main transition is from knowing a probability model to learning
unknown quantities from data. Students learn classical and Bayesian
methods side by side so that they can compare estimators, tests,
confidence intervals, credible intervals, and decision rules.

\section*{How each chapter is
structured}\label{how-each-chapter-is-structured}
\addcontentsline{toc}{section}{How each chapter is structured}

\markright{How each chapter is structured}

Most chapters follow a common structure:

\begin{enumerate}
\def\labelenumi{\arabic{enumi}.}
\tightlist
\item
  \textbf{Motivation and learning objectives} introduce the central
  question of the chapter.
\item
  \textbf{Definitions and theorems} give the mathematical foundation.
\item
  \textbf{Worked examples} show the reasoning step by step.
\item
  \textbf{Practice problems} give students a chance to test their
  understanding.
\item
  \textbf{Solutions} provide detailed derivations for review.
\item
  \textbf{Interactive and lab links} connect theory with visualization
  and computation.
\end{enumerate}

This structure is intended to support both classroom teaching and
self-study.

\section*{Computational labs}\label{computational-labs}
\addcontentsline{toc}{section}{Computational labs}

\markright{Computational labs}

The labs are Jupyter notebooks in the \texttt{labs/} folder. They use
Python to simulate probability models, estimate parameters, visualize
sampling distributions, check asymptotic approximations, and implement
inference procedures.

Students may download the notebooks or open them in Google Colab from
the GitHub repository. The Colab workflow is useful because students can
run the notebooks without installing Python locally.

See the full list here:

\href{labs-summary.qmd}{Open the labs and interactive modules summary}

\section*{Interactive HTML modules}\label{interactive-html-modules}
\addcontentsline{toc}{section}{Interactive HTML modules}

\markright{Interactive HTML modules}

The interactive modules are standalone HTML files in the \texttt{htmls/}
folder. They are designed for visual and conceptual learning. Examples
include exploring probability rules, changing parameters of
distributions, visualizing transformations, simulating convergence, and
comparing classical and Bayesian procedures.

These pages can be used in lecture, recitation, homework review, or
independent study.

\section*{Mathematical notation}\label{mathematical-notation}
\addcontentsline{toc}{section}{Mathematical notation}

\markright{Mathematical notation}

This book uses standard notation from probability and mathematical
statistics. For example:

\begin{itemize}
\tightlist
\item
  \(X, Y, Z\) denote random variables;
\item
  \(F_X(x)\) denotes the CDF of \(X\);
\item
  \(f_X(x)\) denotes a PDF or PMF, depending on context;
\item
  \(E[X]\) denotes expectation;
\item
  \(\operatorname{Var}(X)\) denotes variance;
\item
  \(X_1, \ldots, X_n\) denotes a random sample;
\item
  \(ar X\) denotes the sample mean;
\item
  \(L(   heta)\) and \(\ell(   heta)\) denote likelihood and
  log-likelihood;
\item
  \(\hat heta\) denotes an estimator;
\item
  \(H_0\) and \(H_1\) denote null and alternative hypotheses.
\end{itemize}

\section*{Suggested study strategy}\label{suggested-study-strategy}
\addcontentsline{toc}{section}{Suggested study strategy}

\markright{Suggested study strategy}

For each chapter, students should first identify the main statistical
question. For example: Are we computing a probability? Finding a
distribution? Estimating a parameter? Testing a hypothesis? Building an
interval? Comparing groups?

Then follow this pattern:

\begin{enumerate}
\def\labelenumi{\arabic{enumi}.}
\tightlist
\item
  Write down the model assumptions.
\item
  Identify the relevant random variable or statistic.
\item
  Derive the formula carefully.
\item
  Check the result using simulation when possible.
\item
  Interpret the result in words.
\end{enumerate}

This process is more important than memorizing individual formulas. The
goal is to become fluent in the logic of probability and statistical
inference.

\section*{Repository and book site}\label{repository-and-book-site}
\addcontentsline{toc}{section}{Repository and book site}

\markright{Repository and book site}

The GitHub repository for this book is:

\url{https://github.com/wanghemath/Book-ProbabilityStatisticalTheory}

The published Quarto book site is:

\url{https://wanghemath.github.io/Book-ProbabilityStatisticalTheory/}

\bookmarksetup{startatroot}

\chapter*{Computer Labs and Interactive
Pages}\label{computer-labs-and-interactive-pages}
\addcontentsline{toc}{chapter}{Computer Labs and Interactive Pages}

\markboth{Computer Labs and Interactive Pages}{Computer Labs and
Interactive Pages}

This page collects all interactive HTML pages and Jupyter notebook labs
for \textbf{MATH 5010: Foundations of Statistical Theory and
Probability}.

\section*{Quick Start for Students}\label{quick-start-for-students}
\addcontentsline{toc}{section}{Quick Start for Students}

\markright{Quick Start for Students}

\begin{enumerate}
\def\labelenumi{\arabic{enumi}.}
\tightlist
\item
  Use the \textbf{Interactive HTML} pages for visual exploration before
  or after lecture.
\item
  Use the \textbf{Open in Google Colab} links to run the Python labs in
  the browser.
\item
  Use the \textbf{Local notebook} links if you want to download the
  \texttt{.ipynb} file and run it in JupyterLab, VS Code, or another
  local Python environment.
\end{enumerate}

\section*{Interactive HTML Pages}\label{interactive-html-pages}
\addcontentsline{toc}{section}{Interactive HTML Pages}

\markright{Interactive HTML Pages}

{\def\LTcaptype{none} % do not increment counter
\begin{longtable}[]{@{}
  >{\raggedright\arraybackslash}p{(\linewidth - 4\tabcolsep) * \real{0.3333}}
  >{\raggedright\arraybackslash}p{(\linewidth - 4\tabcolsep) * \real{0.3333}}
  >{\raggedright\arraybackslash}p{(\linewidth - 4\tabcolsep) * \real{0.3333}}@{}}
\toprule\noalign{}
\begin{minipage}[b]{\linewidth}\raggedright
Section
\end{minipage} & \begin{minipage}[b]{\linewidth}\raggedright
Topic
\end{minipage} & \begin{minipage}[b]{\linewidth}\raggedright
Interactive HTML
\end{minipage} \\
\midrule\noalign{}
\endhead
\bottomrule\noalign{}
\endlastfoot
Section 1 & Basic Probability &
\href{htmls/section1_probability_interactive.html}{Open interactive
page} \\
Section 2 & Random Variables and Distributions &
\href{htmls/section2_random_variables_intertributions_interactive.html}{Open
interactive page} \\
Section 3 & Joint and Conditional Distributions &
\href{htmls/section3_joint_conditional_interactive.html}{Open
interactive page} \\
Section 4 & Expectation, Moments, and MGF &
\href{htmls/section4_expectation_mgf_interactive.html}{Open interactive
page} \\
Section 5 & Transformations &
\href{htmls/section5_transformations_interactive.html}{Open interactive
page} \\
Section 6 & Conditional Distribution and Expectation &
\href{htmls/section6_conditional_distribution_interactive.html}{Open
interactive page} \\
Section 7 & Inequalities and Identities &
\href{htmls/section7_inequalities_interactive.html}{Open interactive
page} \\
Section 8 & Random Samples and Order Statistics &
\href{htmls/section8_random_samples_interactive.html}{Open interactive
page} \\
Section 9 & Convergence Theory &
\href{htmls/section9_convergence_interactive.html}{Open interactive
page} \\
Section 10 & Monte Carlo Sampling &
\href{htmls/section10_monte_carlo_interactive.html}{Open interactive
page} \\
Section 11 & Sufficient Statistics &
\href{htmls/section11_sufficient_statistics_interactive.html}{Open
interactive page} \\
Sections 12--13 & Point Estimation &
\href{htmls/sections12_13_point_estimation_interactive.html}{Open
interactive page} \\
Sections 14--15 & Hypothesis Tests &
\href{htmls/sections14_15_hypothesis_tests_interactive.html}{Open
interactive page} \\
Sections 16--17 & Interval Estimation &
\href{htmls/sections16_17_interval_estimations_interactive.html}{Open
interactive page} \\
Sections 18--19 & ANOVA &
\href{htmls/sections18_19_anova_interactive.html}{Open interactive
page} \\
Section 20 & Bayesian Inference &
\href{htmls/section20_bayesian_inference_interactive.html}{Open
interactive page} \\
\end{longtable}
}

\section*{Computer Labs in Google
Colab}\label{computer-labs-in-google-colab}
\addcontentsline{toc}{section}{Computer Labs in Google Colab}

\markright{Computer Labs in Google Colab}

{\def\LTcaptype{none} % do not increment counter
\begin{longtable}[]{@{}
  >{\raggedright\arraybackslash}p{(\linewidth - 6\tabcolsep) * \real{0.2500}}
  >{\raggedright\arraybackslash}p{(\linewidth - 6\tabcolsep) * \real{0.2500}}
  >{\raggedright\arraybackslash}p{(\linewidth - 6\tabcolsep) * \real{0.2500}}
  >{\raggedright\arraybackslash}p{(\linewidth - 6\tabcolsep) * \real{0.2500}}@{}}
\toprule\noalign{}
\begin{minipage}[b]{\linewidth}\raggedright
Section
\end{minipage} & \begin{minipage}[b]{\linewidth}\raggedright
Topic
\end{minipage} & \begin{minipage}[b]{\linewidth}\raggedright
Open in Colab
\end{minipage} & \begin{minipage}[b]{\linewidth}\raggedright
Local notebook
\end{minipage} \\
\midrule\noalign{}
\endhead
\bottomrule\noalign{}
\endlastfoot
Section 1 & Basic Probability &
\href{https://colab.research.google.com/github/wanghemath/Book-ProbabilityStatisticalTheory/blob/main/labs/lab_section1_basic_probability_full_solutions.ipynb}{Open
in Google Colab} &
\href{labs/lab_section1_basic_probability_full_solutions.ipynb}{Download
notebook} \\
Section 2 & Random Variables and Distributions &
\href{https://colab.research.google.com/github/wanghemath/Book-ProbabilityStatisticalTheory/blob/main/labs/lab_section2_random_variables_distributions_full_solutions.ipynb}{Open
in Google Colab} &
\href{labs/lab_section2_random_variables_distributions_full_solutions.ipynb}{Download
notebook} \\
Section 3 & Joint and Conditional Probability &
\href{https://colab.research.google.com/github/wanghemath/Book-ProbabilityStatisticalTheory/blob/main/labs/lab_section3_joint_conditional_probability_full_solutions.ipynb}{Open
in Google Colab} &
\href{labs/lab_section3_joint_conditional_probability_full_solutions.ipynb}{Download
notebook} \\
Section 4 & Expectations, Moments, and MGF &
\href{https://colab.research.google.com/github/wanghemath/Book-ProbabilityStatisticalTheory/blob/main/labs/lab_section4_expectations_moments_mgf_full_solutions_executed.ipynb}{Open
in Google Colab} &
\href{labs/lab_section4_expectations_moments_mgf_full_solutions_executed.ipynb}{Download
notebook} \\
Section 5 & Transformations &
\href{https://colab.research.google.com/github/wanghemath/Book-ProbabilityStatisticalTheory/blob/main/labs/lab_section5_transformations_full_solutions_executed.ipynb}{Open
in Google Colab} &
\href{labs/lab_section5_transformations_full_solutions_executed.ipynb}{Download
notebook} \\
Section 6 & Conditional Expectations &
\href{https://colab.research.google.com/github/wanghemath/Book-ProbabilityStatisticalTheory/blob/main/labs/lab_section6_conditional_expectations_full_solutions.ipynb}{Open
in Google Colab} &
\href{labs/lab_section6_conditional_expectations_full_solutions.ipynb}{Download
notebook} \\
Section 7 & Inequalities and Identities &
\href{https://colab.research.google.com/github/wanghemath/Book-ProbabilityStatisticalTheory/blob/main/labs/lab_section7_inequalities_identities_full_solutions.ipynb}{Open
in Google Colab} &
\href{labs/lab_section7_inequalities_identities_full_solutions.ipynb}{Download
notebook} \\
Section 8 & Sampling and Order Statistics &
\href{https://colab.research.google.com/github/wanghemath/Book-ProbabilityStatisticalTheory/blob/main/labs/lab_section8_sampling_order_statistics_full_solutions.ipynb}{Open
in Google Colab} &
\href{labs/lab_section8_sampling_order_statistics_full_solutions.ipynb}{Download
notebook} \\
Section 9 & Convergence Theory &
\href{https://colab.research.google.com/github/wanghemath/Book-ProbabilityStatisticalTheory/blob/main/labs/lab_section9_convergence_theory_full_solutions.ipynb}{Open
in Google Colab} &
\href{labs/lab_section9_convergence_theory_full_solutions.ipynb}{Download
notebook} \\
Section 10 & Monte Carlo Sampling &
\href{https://colab.research.google.com/github/wanghemath/Book-ProbabilityStatisticalTheory/blob/main/labs/lab_section10_monte_carlo_sampling_full_solutions.ipynb}{Open
in Google Colab} &
\href{labs/lab_section10_monte_carlo_sampling_full_solutions.ipynb}{Download
notebook} \\
Section 11 & Sufficient Statistics &
\href{https://colab.research.google.com/github/wanghemath/Book-ProbabilityStatisticalTheory/blob/main/labs/MATH5010_Lab11_Sufficient_Statistics_Full_Solutions.ipynb}{Open
in Google Colab} &
\href{labs/MATH5010_Lab11_Sufficient_Statistics_Full_Solutions.ipynb}{Download
notebook} \\
Sections 12--13 & Point Estimation &
\href{https://colab.research.google.com/github/wanghemath/Book-ProbabilityStatisticalTheory/blob/main/labs/MATH5010_Lab12_13_Point_Estimation_Full_Solutions.ipynb}{Open
in Google Colab} &
\href{labs/MATH5010_Lab12_13_Point_Estimation_Full_Solutions.ipynb}{Download
notebook} \\
Sections 14--15 & Hypothesis Tests and Interval Estimation &
\href{https://colab.research.google.com/github/wanghemath/Book-ProbabilityStatisticalTheory/blob/main/labs/MATH5010_Lab14_15_Hypothesis_Tests_and_Interval_Estimation_Full_Solutions.ipynb}{Open
in Google Colab} &
\href{labs/MATH5010_Lab14_15_Hypothesis_Tests_and_Interval_Estimation_Full_Solutions.ipynb}{Download
notebook} \\
Sections 16--17 & ANOVA and Bayesian Statistics &
\href{https://colab.research.google.com/github/wanghemath/Book-ProbabilityStatisticalTheory/blob/main/labs/MATH5010_Lab16_17_ANOVA_and_Bayesian_Statistics_Full_Solutions.ipynb}{Open
in Google Colab} &
\href{labs/MATH5010_Lab16_17_ANOVA_and_Bayesian_Statistics_Full_Solutions.ipynb}{Download
notebook} \\
Sections 18--19 & Bayesian Inference II and ANOVA II &
\href{https://colab.research.google.com/github/wanghemath/Book-ProbabilityStatisticalTheory/blob/main/labs/MATH5010_Lab18_19_Bayesian_Inference_II_and_ANOVA_II_Full_Solutions.ipynb}{Open
in Google Colab} &
\href{labs/MATH5010_Lab18_19_Bayesian_Inference_II_and_ANOVA_II_Full_Solutions.ipynb}{Download
notebook} \\
Section 20 & Bayesian Inference &
\href{https://colab.research.google.com/github/wanghemath/Book-ProbabilityStatisticalTheory/blob/main/labs/MATH5010_Lab20_Bayesian_Inference_Full_Solutions.ipynb}{Open
in Google Colab} &
\href{labs/MATH5010_Lab20_Bayesian_Inference_Full_Solutions.ipynb}{Download
notebook} \\
\end{longtable}
}

\section*{Suggested Weekly Workflow}\label{suggested-weekly-workflow}
\addcontentsline{toc}{section}{Suggested Weekly Workflow}

\markright{Suggested Weekly Workflow}

For each section, a good learning sequence is:

\begin{enumerate}
\def\labelenumi{\arabic{enumi}.}
\tightlist
\item
  Read the corresponding book chapter.
\item
  Open the interactive HTML page and experiment with the sliders, plots,
  or simulations.
\item
  Open the corresponding lab in Google Colab.
\item
  Run all cells once.
\item
  Modify parameters and rerun the code.
\item
  Write a short interpretation of the output in statistical language.
\end{enumerate}

\bookmarksetup{startatroot}

\chapter{Chapter 1: Basic
Probability}\label{chapter-1-basic-probability}

This chapter introduces the language of probability used throughout MATH
5010. Events are sets, probability is a function on events, and
conditional probability is the basic operation for updating uncertainty
after new information is observed.

\textbf{Topics.} Set theory; probability spaces; counting; conditional
probability; independence; total probability; Bayes' theorem; classical
examples and paradoxes.

\section{Why Probability Theory?}\label{why-probability-theory}

Probability theory gives a mathematical language for uncertainty. In
statistics, machine learning, scientific experiments, and data analysis,
we often observe only partial information about a system. Probability
provides a model for possible outcomes, while statistical inference uses
observed data to learn about the model.

\begin{tcolorbox}[enhanced jigsaw, arc=.35mm, bottomrule=.15mm, bottomtitle=1mm, breakable, colback=white, colbacktitle=quarto-callout-tip-color!10!white, colframe=quarto-callout-tip-color-frame, coltitle=black, left=2mm, leftrule=.75mm, opacityback=0, opacitybacktitle=0.6, rightrule=.15mm, title=\textcolor{quarto-callout-tip-color}{\faLightbulb}\hspace{0.5em}{Learning box}, titlerule=0mm, toprule=.15mm, toptitle=1mm]

Main distinction

\begin{itemize}
\item
  \textbf{Probability:} assume a probability model is known, and use it
  to predict the behavior of data.
\item
  \textbf{Statistical inference:} observe data, and use the data to
  infer an unknown probability model or its parameters.
\end{itemize}

\end{tcolorbox}

In this section, we build the foundational objects of probability
theory:

\begin{enumerate}
\def\labelenumi{\arabic{enumi}.}
\item
  set theory terminology;
\item
  sample spaces, events, and probability measures;
\item
  classical counting methods;
\item
  conditional probability and independence;
\item
  Bayes' theorem and its applications.
\end{enumerate}

\section{Set Theory Terminology
Review}\label{set-theory-terminology-review}

The goal of this section is to review the set language that will become
the language of events in probability.

Probability begins with sets. The set of all possible outcomes is called
the sample space, and events are subsets of the sample space. Therefore,
set operations become event operations.

\subsection{Basic set operations}\label{basic-set-operations}

This subsection introduces the basic operations that translate English
statements such as ``not'', ``or'', and ``and'' into set notation.

Let \(S\) be a universal set, and let \(A,B\subseteq S\).

\textbf{Definition 1} (Complement, union, intersection, and difference).
\[\begin{aligned}
A^c &= S\setminus A && \text{the complement of $A$; this means ``not $A$''},\\
A\cup B &= \{x\in S: x\in A \text{ or } x\in B\} && \text{the union; this means ``$A$ or $B$''},\\
A\cap B &= \{x\in S: x\in A \text{ and } x\in B\} && \text{the intersection; this means ``$A$ and $B$''},\\
A\setminus B &= \{x\in S: x\in A \text{ and } x\notin B\} && \text{the difference.}
\end{aligned}\]

\subsection{Algebra of sets}\label{algebra-of-sets}

This subsection records the algebraic rules that allow us to rewrite
complicated events in simpler equivalent forms.

The following identities are used constantly in probability.

\textbf{Proposition 2} (Properties of set operations). \emph{For sets
\(A,B,C\subseteq S\):}

\begin{enumerate}
\def\labelenumi{\arabic{enumi}.}
\item
  \emph{\textbf{Commutativity:}
  \[A\cup B=B\cup A,\qquad A\cap B=B\cap A.\]}
\item
  \emph{\textbf{Associativity:}
  \[(A\cup B)\cup C=A\cup(B\cup C),\qquad (A\cap B)\cap C=A\cap(B\cap C).\]}
\item
  \emph{\textbf{Distributive laws:}
  \[A\cup(B\cap C)=(A\cup B)\cap(A\cup C),\]
  \[A\cap(B\cup C)=(A\cap B)\cup(A\cap C).\]}
\item
  \emph{\textbf{DeMorgan's laws:}
  \[(A\cup B)^c=A^c\cap B^c,\qquad (A\cap B)^c=A^c\cup B^c.\]}
\end{enumerate}

\begin{remark}
\emph{Remark 3}. In probability language, DeMorgan's laws say:

\begin{itemize}
\item
  ``not \((A\) or \(B)\)'' is the same as ``not \(A\) and not \(B\)'';
\item
  ``not \((A\) and \(B)\)'' is the same as ``not \(A\) or not \(B\)''.
\end{itemize}

\end{remark}

\subsection{Disjoint events, partitions, cardinality, and
countability}\label{disjoint-events-partitions-cardinality-and-countability}

This subsection introduces disjointness, partitions, and sizes of sets,
which are especially important for additivity and counting.

\textbf{Definition 4} (Pairwise disjoint). A collection of sets
\(A_1,A_2,\ldots\) is called \textbf{pairwise disjoint} or
\textbf{mutually exclusive} if
\[A_i\cap A_j=\varnothing\qquad \text{for all } i\ne j.\]

\textbf{Definition 5} (Partition). A collection \(A_1,A_2,\ldots\) is
called a \textbf{partition} of \(S\) if

\begin{enumerate}
\def\labelenumi{\arabic{enumi}.}
\item
  \(A_i\cap A_j=\varnothing\) for all \(i\ne j\);
\item
  \(\displaystyle \bigcup_{i=1}^{\infty} A_i=S\).
\end{enumerate}

\textbf{Definition 6} (Cardinality and countability). The notation
\(|A|\) denotes the number of elements in \(A\), called the
\textbf{cardinality} of \(A\). A set \(S\) is \textbf{countable} if its
cardinality is less than or equal to the cardinality of the natural
numbers \(\mathbb{N}\).

\textbf{Example 7} (Partition of \([0,\infty)\)). Let
\[A_i=[i,i+1),\qquad i=0,1,2,\ldots.\] Then the sets \(A_i\) are
pairwise disjoint and \[[0,\infty)=\bigcup_{i=0}^{\infty} A_i.\] Thus
\(\{A_i\}_{i=0}^{\infty}\) is a countable partition of \([0,\infty)\).

\begin{tcolorbox}[enhanced jigsaw, arc=.35mm, bottomrule=.15mm, bottomtitle=1mm, breakable, colback=white, colbacktitle=quarto-callout-note-color!10!white, colframe=quarto-callout-note-color-frame, coltitle=black, left=2mm, leftrule=.75mm, opacityback=0, opacitybacktitle=0.6, rightrule=.15mm, title=\textcolor{quarto-callout-note-color}{\faInfo}\hspace{0.5em}{Solution}, titlerule=0mm, toprule=.15mm, toptitle=1mm]

The intervals \(A_i=[i,i+1)\) do not overlap because each real number
belongs to at most one such half-open interval. Every \(x\ge 0\) has a
unique integer part \(i=\lfloor x\rfloor\), so \(x\in [i,i+1)\).
Therefore the intervals are pairwise disjoint and their union is exactly
\([0,\infty)\).

\end{tcolorbox}

\textbf{Example 8} (Rational numbers). The set of rational numbers
\(\mathbb{Q}\) is countable. One can list rational numbers
systematically by arranging fractions \(p/q\) in a two-dimensional grid
and moving diagonally through the grid while skipping duplicates.

\begin{tcolorbox}[enhanced jigsaw, arc=.35mm, bottomrule=.15mm, bottomtitle=1mm, breakable, colback=white, colbacktitle=quarto-callout-note-color!10!white, colframe=quarto-callout-note-color-frame, coltitle=black, left=2mm, leftrule=.75mm, opacityback=0, opacitybacktitle=0.6, rightrule=.15mm, title=\textcolor{quarto-callout-note-color}{\faInfo}\hspace{0.5em}{Solution}, titlerule=0mm, toprule=.15mm, toptitle=1mm]

A rational number can be written as \(p/q\) with \(p\in\mathbb Z\) and
\(q\in\mathbb N\). The pairs \((p,q)\) form a countable grid, and a
countable union of countable rows is countable. Removing duplicate
fractions such as \(1/2=2/4\) leaves a subset of a countable set, so
\(\mathbb Q\) is countable.

\end{tcolorbox}

\section{Sample Spaces and Events}\label{sample-spaces-and-events}

The goal of this section is to identify the universe of possible
outcomes and the events whose probabilities we want to compute.

\textbf{Definition 9} (Experiment, sample space, outcome, event). An
\textbf{experiment} is a repeatable procedure with a set of possible
results. Each possible result is called an \textbf{outcome},
\textbf{realization}, or \textbf{element}. The \textbf{sample space} is
\[S=\{\text{all possible outcomes of the experiment}\}.\] An
\textbf{event} is a subset of \(S\).

\textbf{Example 10} (Flipping a coin once). If the experiment is
flipping a coin once, then \[S=\{H,T\},\] where \(H\) means heads and
\(T\) means tails. The event ``heads occurs'' is \(\{H\}\).

\begin{tcolorbox}[enhanced jigsaw, arc=.35mm, bottomrule=.15mm, bottomtitle=1mm, breakable, colback=white, colbacktitle=quarto-callout-note-color!10!white, colframe=quarto-callout-note-color-frame, coltitle=black, left=2mm, leftrule=.75mm, opacityback=0, opacitybacktitle=0.6, rightrule=.15mm, title=\textcolor{quarto-callout-note-color}{\faInfo}\hspace{0.5em}{Solution}, titlerule=0mm, toprule=.15mm, toptitle=1mm]

There are two possible outcomes, heads or tails, so the sample space is
\(S=\{H,T\}\). The event that heads occurs is the subset containing only
the outcome \(H\), namely \(\{H\}\).

\end{tcolorbox}

\textbf{Example 11} (Rolling a six-sided die once). If the experiment is
rolling a six-sided die once, then \[S=\{1,2,3,4,5,6\}.\] The event
``the result is even'' is \[E=\{2,4,6\}.\]

\begin{tcolorbox}[enhanced jigsaw, arc=.35mm, bottomrule=.15mm, bottomtitle=1mm, breakable, colback=white, colbacktitle=quarto-callout-note-color!10!white, colframe=quarto-callout-note-color-frame, coltitle=black, left=2mm, leftrule=.75mm, opacityback=0, opacitybacktitle=0.6, rightrule=.15mm, title=\textcolor{quarto-callout-note-color}{\faInfo}\hspace{0.5em}{Solution}, titlerule=0mm, toprule=.15mm, toptitle=1mm]

For a fair die, each singleton has probability \(1/6\). The event
\(E=\{2,4,6\}\) contains three equally likely outcomes, so
\(\mathbb{P}(E)=3(1/6)=1/2\).

\end{tcolorbox}

\textbf{Example 12} (Tossing a coin until heads appears). Toss a coin
repeatedly until heads appears, and record the number of tosses needed.
Then \[S=\{1,2,3,\ldots\}.\] This sample space is infinite, countable,
and discrete.

\begin{tcolorbox}[enhanced jigsaw, arc=.35mm, bottomrule=.15mm, bottomtitle=1mm, breakable, colback=white, colbacktitle=quarto-callout-note-color!10!white, colframe=quarto-callout-note-color-frame, coltitle=black, left=2mm, leftrule=.75mm, opacityback=0, opacitybacktitle=0.6, rightrule=.15mm, title=\textcolor{quarto-callout-note-color}{\faInfo}\hspace{0.5em}{Solution}, titlerule=0mm, toprule=.15mm, toptitle=1mm]

The first head could occur on toss \(1\), toss \(2\), toss \(3\), and so
on. Thus the possible recorded values are all positive integers, giving
\(S=\{1,2,3,\ldots\}\).

\end{tcolorbox}

\textbf{Example 13} (Dropping a point in a disk). Suppose a point is
randomly dropped inside the unit disk. The sample space is
\[S=\{(x_1,x_2)\in \mathbb{R}^2: x_1^2+x_2^2<1\}.\] This sample space is
uncountable and continuous.

\begin{tcolorbox}[enhanced jigsaw, arc=.35mm, bottomrule=.15mm, bottomtitle=1mm, breakable, colback=white, colbacktitle=quarto-callout-note-color!10!white, colframe=quarto-callout-note-color-frame, coltitle=black, left=2mm, leftrule=.75mm, opacityback=0, opacitybacktitle=0.6, rightrule=.15mm, title=\textcolor{quarto-callout-note-color}{\faInfo}\hspace{0.5em}{Solution}, titlerule=0mm, toprule=.15mm, toptitle=1mm]

A point in the plane is represented by coordinates \((x_1,x_2)\). Being
inside the unit disk means its squared distance from the origin is less
than \(1\), so \(x_1^2+x_2^2<1\). Hence the sample space is the open
unit disk.

\end{tcolorbox}

\textbf{Example 14} (Brownian motion sample space). A Brownian motion
experiment observes the path of a particle suspended in a liquid or
gaseous medium. The outcome is not a single number, but a continuous
path. A simplified way to write the sample space is
\[S=\{\text{all continuous paths } x:[0,\infty)\to \mathbb{R}^d\}.\]
This sample space is continuous but not finite-dimensional.

\begin{tcolorbox}[enhanced jigsaw, arc=.35mm, bottomrule=.15mm, bottomtitle=1mm, breakable, colback=white, colbacktitle=quarto-callout-note-color!10!white, colframe=quarto-callout-note-color-frame, coltitle=black, left=2mm, leftrule=.75mm, opacityback=0, opacitybacktitle=0.6, rightrule=.15mm, title=\textcolor{quarto-callout-note-color}{\faInfo}\hspace{0.5em}{Solution}, titlerule=0mm, toprule=.15mm, toptitle=1mm]

For Brownian motion, one observation is an entire function of time
rather than a finite vector of measurements. Therefore the sample space
is naturally a set of continuous paths, which is infinite-dimensional in
spirit.

\end{tcolorbox}

\textbf{Example 15} (Infinitely many coin tosses). If we toss a coin
infinitely many times, then an outcome is an infinite sequence such as
\[HHHTHTTHT\cdots.\] Thus \[S=\{H,T\}^{\mathbb{N}}.\] This sample space
is uncountable. It can be connected conceptually to binary expansions of
numbers in \([0,1]\).

\begin{tcolorbox}[enhanced jigsaw, arc=.35mm, bottomrule=.15mm, bottomtitle=1mm, breakable, colback=white, colbacktitle=quarto-callout-note-color!10!white, colframe=quarto-callout-note-color-frame, coltitle=black, left=2mm, leftrule=.75mm, opacityback=0, opacitybacktitle=0.6, rightrule=.15mm, title=\textcolor{quarto-callout-note-color}{\faInfo}\hspace{0.5em}{Solution}, titlerule=0mm, toprule=.15mm, toptitle=1mm]

Each outcome chooses either \(H\) or \(T\) at every positive integer
time. Therefore an outcome is an element of \(\{H,T\}^{\mathbb N}\).
Since infinite binary sequences can be matched with binary expansions of
numbers in \([0,1]\) up to the usual non-unique endpoints, this space is
uncountable.

\end{tcolorbox}

\begin{tcolorbox}[enhanced jigsaw, arc=.35mm, bottomrule=.15mm, bottomtitle=1mm, breakable, colback=white, colbacktitle=quarto-callout-warning-color!10!white, colframe=quarto-callout-warning-color-frame, coltitle=black, left=2mm, leftrule=.75mm, opacityback=0, opacitybacktitle=0.6, rightrule=.15mm, title=\textcolor{quarto-callout-warning-color}{\faExclamationTriangle}\hspace{0.5em}{Warning}, titlerule=0mm, toprule=.15mm, toptitle=1mm]

Important point A sample space can be finite, countably infinite,
uncountable, geometric, functional, or infinite-dimensional. The
probability framework is designed to handle all of these examples in one
language.

\end{tcolorbox}

\section{Probability Measures and Probability
Spaces}\label{probability-measures-and-probability-spaces}

The goal of this section is to define probability rigorously as a
function on a suitable collection of events.

The modern axiomatic definition of probability was formalized by Andrey
Kolmogorov in the 1930s. The core idea is that probability is a function
assigning numbers to events.

\subsection{Sigma-algebras}\label{sigma-algebras}

This subsection explains why we need a designated collection of
measurable events before assigning probabilities.

Before defining a probability measure, we need to specify which subsets
are allowed to be events. In finite sample spaces, every subset can be
treated as an event. In uncountable spaces such as intervals or disks,
not every subset behaves well enough to have a probability assigned to
it. This motivates the notion of a sigma-algebra.

\textbf{Definition 16} (Sigma-algebra). Let \(S\) be a sample space. A
\textbf{\(\sigma\)-algebra} \(\mathcal{B}\) on \(S\) is a collection of
subsets of \(S\) satisfying:

\begin{enumerate}
\def\labelenumi{\arabic{enumi}.}
\item
  \(S\in \mathcal{B}\) and \(\varnothing\in\mathcal{B}\);
\item
  if \(A\in\mathcal{B}\), then \(A^c\in\mathcal{B}\);
\item
  if \(A_1,A_2,\ldots\in\mathcal{B}\), then
  \(\displaystyle \bigcup_{i=1}^{\infty}A_i\in\mathcal{B}\).
\end{enumerate}

The sets in \(\mathcal{B}\) are called \textbf{measurable sets}. The
pair \((S,\mathcal{B})\) is called a \textbf{measurable space}.

\begin{remark}
\emph{Remark 17}. This course does not focus on measure theory. The main
practical idea is: \(\mathcal{B}\) is the collection of events to which
the probability function is allowed to assign probabilities.
\end{remark}

\subsection{Kolmogorov axioms}\label{kolmogorov-axioms}

This subsection states the three axioms that define probability in the
modern mathematical framework.

\textbf{Definition 18} (Probability measure). Let \(S\) be a sample
space and let \(\mathcal{B}\) be a \(\sigma\)-algebra of subsets of
\(S\). A \textbf{probability measure} is a function
\[\mathbb{P}:\mathcal{B}\to \mathbb{R}\] satisfying the following three
axioms:

\begin{enumerate}
\def\labelenumi{\arabic{enumi}.}
\item
  \textbf{Nonnegativity:} \(\mathbb{P}(A)\ge 0\) for all
  \(A\in\mathcal{B}\).
\item
  \textbf{Normalization:} \(\mathbb{P}(S)=1\).
\item
  \textbf{Countable additivity:} if \(A_i\cap A_j=\varnothing\) for all
  \(i\ne j\), then
  \[\mathbb{P}\left(\bigcup_{i=1}^{\infty}A_i\right)=\sum_{i=1}^{\infty}\mathbb{P}(A_i).\]
\end{enumerate}

The triple \((S,\mathcal{B},\mathbb{P})\) is called a
\textbf{probability space}.

\subsection{Consequences of the
axioms}\label{consequences-of-the-axioms}

This subsection derives the probability rules that we will use
repeatedly throughout the course.

\textbf{Proposition 19} (Basic properties of probability). \emph{Let
\((S,\mathcal{B},\mathbb{P})\) be a probability space and let
\(A,B\in\mathcal{B}\). Then:}

\begin{enumerate}
\def\labelenumi{\arabic{enumi}.}
\item
  \emph{\(\mathbb{P}(\varnothing)=0\);}
\item
  \emph{\(0\le \mathbb{P}(A)\le 1\);}
\item
  \emph{if \(A\subseteq B\), then \(\mathbb{P}(A)\le \mathbb{P}(B)\);}
\item
  \emph{\(\mathbb{P}(A^c)=1-\mathbb{P}(A)\);}
\item
  \emph{\(\mathbb{P}(A\cup B)=\mathbb{P}(A)+\mathbb{P}(B)-\mathbb{P}(A\cap B)\);}
\item
  \emph{\(\mathbb{P}(A\cup B)\le \mathbb{P}(A)+\mathbb{P}(B)\);}
\item
  \emph{\(\mathbb{P}(A\cap B)\ge \mathbb{P}(A)+\mathbb{P}(B)-1\).}
\end{enumerate}

\begin{proof}
\emph{Proof.} We prove the main identities.

First, \(S\) and \(\varnothing\) are disjoint and
\(S\cup\varnothing=S\). By additivity,
\[1=\mathbb{P}(S)=\mathbb{P}(S\cup\varnothing)=\mathbb{P}(S)+\mathbb{P}(\varnothing)=1+\mathbb{P}(\varnothing),\]
so \(\mathbb{P}(\varnothing)=0\).

Next, \(A\) and \(A^c\) are disjoint and \(A\cup A^c=S\). Therefore
\[1=\mathbb{P}(S)=\mathbb{P}(A)+\mathbb{P}(A^c),\] so
\(\mathbb{P}(A^c)=1-\mathbb{P}(A)\). Since probabilities are
nonnegative, \(\mathbb{P}(A^c)\ge0\), and hence \(\mathbb{P}(A)\le1\).

If \(A\subseteq B\), then \(B=A\cup(B\setminus A)\), and this is a
disjoint union. Thus
\[\mathbb{P}(B)=\mathbb{P}(A)+\mathbb{P}(B\setminus A)\ge \mathbb{P}(A).\]

Finally, write \[A\cup B=A\cup(B\setminus A),\] where the union is
disjoint. Also \[B=(B\setminus A)\cup(A\cap B),\] again disjoint. Hence
\(\mathbb{P}(B\setminus A)=\mathbb{P}(B)-\mathbb{P}(A\cap B)\), so
\[\mathbb{P}(A\cup B)=\mathbb{P}(A)+\mathbb{P}(B)-\mathbb{P}(A\cap B).\]
The inequalities follow from \(\mathbb{P}(A\cap B)\ge0\) and
\(\mathbb{P}(A\cup B)\le1\).~◻
\end{proof}

\begin{remark}
\emph{Remark 20} (Bonferroni's inequality). The inequality
\[\mathbb{P}(A\cap B)\ge \mathbb{P}(A)+\mathbb{P}(B)-1\] is a two-event
version of a Bonferroni inequality. It is useful when we know two
marginal probabilities but do not know the exact overlap.
\end{remark}

\subsection{Law of total probability for
partitions}\label{law-of-total-probability-for-partitions}

This subsection introduces the idea that a partition lets us break an
event into disjoint pieces.

If \(A_1,A_2,\ldots\) form a partition of \(S\), then any event \(B\)
can be decomposed into disjoint pieces:
\[B=(B\cap A_1)\cup(B\cap A_2)\cup(B\cap A_3)\cup\cdots.\] Therefore,
\[\mathbb{P}(B)=\mathbb{P}(B\cap A_1)+\mathbb{P}(B\cap A_2)+\mathbb{P}(B\cap A_3)+\cdots.\]

\section{Continuity of Probability}\label{continuity-of-probability}

The goal of this section is to explain how probabilities behave when
events grow or shrink through limits.

A probability measure behaves continuously with respect to increasing
and decreasing limits of events.

\textbf{Definition 21} (Increasing and decreasing sequences of events).
A sequence of events \(A_1,A_2,\ldots\) is called \textbf{increasing} if
\[A_n\subseteq A_{n+1}\qquad \text{for all } n\ge1.\] It is called
\textbf{decreasing} if
\[A_n\supseteq A_{n+1}\qquad \text{for all } n\ge1.\] For an increasing
sequence, define \[\lim_{n\to\infty} A_n=\bigcup_{n=1}^{\infty}A_n.\]
For a decreasing sequence, define
\[\lim_{n\to\infty} A_n=\bigcap_{n=1}^{\infty}A_n.\]

\textbf{Theorem 22} (Continuity of probability). \emph{Let
\((S,\mathcal{B},\mathbb{P})\) be a probability space.}

\begin{enumerate}
\def\labelenumi{\arabic{enumi}.}
\item
  \emph{If \(A_n\subseteq A_{n+1}\) for all \(n\), then
  \[\mathbb{P}\left(\bigcup_{n=1}^{\infty}A_n\right)=\lim_{n\to\infty}\mathbb{P}(A_n).\]}
\item
  \emph{If \(A_n\supseteq A_{n+1}\) for all \(n\), then
  \[\mathbb{P}\left(\bigcap_{n=1}^{\infty}A_n\right)=\lim_{n\to\infty}\mathbb{P}(A_n).\]}
\end{enumerate}

\emph{Equivalently, for monotone event sequences,
\[\mathbb{P}\left(\lim_{n\to\infty}A_n\right)=\lim_{n\to\infty}\mathbb{P}(A_n).\]}

\begin{proof}
\emph{Proof.} For the increasing case, define disjoint sets
\[B_1=A_1,\qquad B_n=A_n\setminus A_{n-1},\quad n\ge2.\] Then
\[\bigcup_{n=1}^{\infty}A_n=\bigcup_{n=1}^{\infty}B_n\] and the \(B_n\)
are pairwise disjoint. Thus
\[\mathbb{P}\left(\bigcup_{n=1}^{\infty}A_n\right)=\sum_{n=1}^{\infty}\mathbb{P}(B_n)=\lim_{N\to\infty}\sum_{n=1}^{N}\mathbb{P}(B_n)=\lim_{N\to\infty}\mathbb{P}(A_N).\]
For the decreasing case, apply the increasing result to the complements
\(A_n^c\), since \(A_n^c\subseteq A_{n+1}^c\), and use DeMorgan's law.~◻
\end{proof}

\section{Examples of Probability
Models}\label{examples-of-probability-models}

The goal of this section is to translate familiar random experiments
into explicit probability models.

\subsection{Unfair coin}\label{unfair-coin}

This subsection begins with the simplest non-equally-likely probability
model.

\textbf{Example 23} (Flipping an unfair coin once). Let \[S=\{H,T\}.\]
If the coin is unfair, we may assign
\[\mathbb{P}(H)=p,\qquad \mathbb{P}(T)=1-p,\] where \(0\le p\le1\). This
is a probability model because probabilities are nonnegative and sum to
one.

\begin{tcolorbox}[enhanced jigsaw, arc=.35mm, bottomrule=.15mm, bottomtitle=1mm, breakable, colback=white, colbacktitle=quarto-callout-note-color!10!white, colframe=quarto-callout-note-color-frame, coltitle=black, left=2mm, leftrule=.75mm, opacityback=0, opacitybacktitle=0.6, rightrule=.15mm, title=\textcolor{quarto-callout-note-color}{\faInfo}\hspace{0.5em}{Solution}, titlerule=0mm, toprule=.15mm, toptitle=1mm]

We must assign nonnegative probabilities that sum to one. If
\(\mathbb{P}(H)=p\) with \(0\le p\le 1\), then the only possible value
for \(\mathbb{P}(T)\) is \(1-p\). Thus
\(\mathbb{P}(H)+\mathbb{P}(T)=p+(1-p)=1\).

\end{tcolorbox}

\subsection{Rolling a die}\label{rolling-a-die}

This subsection shows how a finite equally likely sample space leads to
probability by counting.

\textbf{Example 24} (Rolling a six-sided die once). For a fair die,
\[S=\{1,2,3,4,5,6\},\qquad \mathbb{P}(\{i\})=\frac16,\quad i=1,\ldots,6.\]
For the event \(E=\{2,4,6\}\) of rolling an even number,
\[\mathbb{P}(E)=\frac{|E|}{|S|}=\frac{3}{6}=\frac12.\]

\begin{tcolorbox}[enhanced jigsaw, arc=.35mm, bottomrule=.15mm, bottomtitle=1mm, breakable, colback=white, colbacktitle=quarto-callout-note-color!10!white, colframe=quarto-callout-note-color-frame, coltitle=black, left=2mm, leftrule=.75mm, opacityback=0, opacitybacktitle=0.6, rightrule=.15mm, title=\textcolor{quarto-callout-note-color}{\faInfo}\hspace{0.5em}{Solution}, titlerule=0mm, toprule=.15mm, toptitle=1mm]

For a fair die, each singleton has probability \(1/6\). The event
\(E=\{2,4,6\}\) contains three equally likely outcomes, so
\(\mathbb{P}(E)=3(1/6)=1/2\).

\end{tcolorbox}

\subsection{Dart board model}\label{dart-board-model}

This subsection gives a geometric probability model where probability is
proportional to area.

\textbf{Example 25} (Dart board). Suppose a circular dart board has
radius \(r\), and the distance between adjacent rings is \(r/5\). Assume
the dart always hits the board and the hit location is uniformly
distributed over area.

The probability of scoring \(i\) points is
\[\mathbb{P}(\text{scoring } i \text{ points})=\frac{\text{area of region }i}{\text{area of dart board}}.\]
For \(i=1,2,3,4,5\), the ring for score \(i\) has outer radius
\((6-i)r/5\) and inner radius \((5-i)r/5\). Hence
\[\mathbb{P}(\text{scoring } i \text{ points})
=\frac{\pi\left((6-i)r/5\right)^2-\pi\left((5-i)r/5\right)^2}{\pi r^2}
=\frac{(6-i)^2-(5-i)^2}{5^2}.\]

\begin{tcolorbox}[enhanced jigsaw, arc=.35mm, bottomrule=.15mm, bottomtitle=1mm, breakable, colback=white, colbacktitle=quarto-callout-note-color!10!white, colframe=quarto-callout-note-color-frame, coltitle=black, left=2mm, leftrule=.75mm, opacityback=0, opacitybacktitle=0.6, rightrule=.15mm, title=\textcolor{quarto-callout-note-color}{\faInfo}\hspace{0.5em}{Solution}, titlerule=0mm, toprule=.15mm, toptitle=1mm]

The board is assumed uniformly hit over area, so probability equals area
ratio. The score-\(i\) ring has area
\(\pi((6-i)r/5)^2-\pi((5-i)r/5)^2\). Dividing by the total area
\(\pi r^2\) gives \(\big((6-i)^2-(5-i)^2\big)/25\).

\end{tcolorbox}

\section{Classical Probability and
Counting}\label{classical-probability-and-counting}

The goal of this section is to compute probabilities by counting equally
likely outcomes.

\subsection{Classical definition of
probability}\label{classical-definition-of-probability}

This subsection states the counting formula for finite equally likely
sample spaces.

If the outcomes of an experiment are all equally likely and the sample
space is finite, probability can be computed by counting.

\textbf{Definition 26} (Classical probability). Suppose \(S\) is finite
and all outcomes in \(S\) are equally likely. For an event
\(A\subseteq S\),
\[\mathbb{P}(A)=\frac{|A|}{|S|}=\frac{\text{number of ways the event can happen}}{\text{total number of possible outcomes}}.\]

\textbf{Example 27} (Flipping a fair coin two times). The sample space
is \[S=\{HH,HT,TH,TT\}.\] All four outcomes are equally likely. If \(A\)
is the event ``exactly one head occurs,'' then
\[A=\{HT,TH\},\qquad \mathbb{P}(A)=\frac{|A|}{|S|}=\frac24=\frac12.\]

\begin{tcolorbox}[enhanced jigsaw, arc=.35mm, bottomrule=.15mm, bottomtitle=1mm, breakable, colback=white, colbacktitle=quarto-callout-note-color!10!white, colframe=quarto-callout-note-color-frame, coltitle=black, left=2mm, leftrule=.75mm, opacityback=0, opacitybacktitle=0.6, rightrule=.15mm, title=\textcolor{quarto-callout-note-color}{\faInfo}\hspace{0.5em}{Solution}, titlerule=0mm, toprule=.15mm, toptitle=1mm]

The four outcomes \(HH,HT,TH,TT\) are equally likely. Exactly one head
occurs in \(HT\) and \(TH\), so there are two favorable outcomes out of
four total outcomes. Therefore the probability is \(2/4=1/2\).

\end{tcolorbox}

\subsection{Fundamental theorem of
counting}\label{fundamental-theorem-of-counting}

This subsection explains the product rule for multi-stage choices.

\textbf{Theorem 28} (Fundamental theorem of counting). \emph{If
operation A can be performed in \(m\) different ways and operation B can
be performed in \(n\) different ways, then the sequence ``operation A
followed by operation B'' can be performed in \[m\times n\] different
ways.}

\textbf{Example 29} (Two-stage experiment). If a password begins with
one of \(26\) letters and ends with one of \(10\) digits, then the
number of possible passwords is \[26\times 10=260.\]

\begin{tcolorbox}[enhanced jigsaw, arc=.35mm, bottomrule=.15mm, bottomtitle=1mm, breakable, colback=white, colbacktitle=quarto-callout-note-color!10!white, colframe=quarto-callout-note-color-frame, coltitle=black, left=2mm, leftrule=.75mm, opacityback=0, opacitybacktitle=0.6, rightrule=.15mm, title=\textcolor{quarto-callout-note-color}{\faInfo}\hspace{0.5em}{Solution}, titlerule=0mm, toprule=.15mm, toptitle=1mm]

Choose the first character in \(26\) ways and the second character in
\(10\) ways. By the product rule, the total number of passwords is
\(26\cdot 10=260\).

\end{tcolorbox}

\subsection{Permutations}\label{permutations}

This subsection studies ordered selections without replacement.

\textbf{Definition 30} (Permutation rule). The number of ways to arrange
\(k\) objects from a set of \(n\) distinct elements, where order matters
and no replacement is allowed, is
\[P_k^n=\frac{n!}{(n-k)!}=n(n-1)\cdots(n-k+1).\]

\textbf{Example 31} (Three-letter arrangements). From the five letters
\(A,B,C,D,E\), the number of ordered three-letter arrangements without
repetition is \[P_3^5=\frac{5!}{(5-3)!}=5\cdot4\cdot3=60.\]

\begin{tcolorbox}[enhanced jigsaw, arc=.35mm, bottomrule=.15mm, bottomtitle=1mm, breakable, colback=white, colbacktitle=quarto-callout-note-color!10!white, colframe=quarto-callout-note-color-frame, coltitle=black, left=2mm, leftrule=.75mm, opacityback=0, opacitybacktitle=0.6, rightrule=.15mm, title=\textcolor{quarto-callout-note-color}{\faInfo}\hspace{0.5em}{Solution}, titlerule=0mm, toprule=.15mm, toptitle=1mm]

There are \(5\) choices for the first letter, \(4\) choices for the
second, and \(3\) choices for the third. Therefore the total is
\(5\cdot4\cdot3=60\), which equals \(P_3^5\).

\end{tcolorbox}

\subsection{Combinations}\label{combinations}

This subsection studies unordered selections without replacement.

\textbf{Definition 32} (Combination rule). The number of ways to choose
a subset of \(k\) objects from \(n\) distinct objects, where order does
not matter and no replacement is allowed, is
\[C_k^n=\binom{n}{k}=\frac{n!}{(n-k)!k!}=\frac{P_k^n}{k!}.\]

\textbf{Example 33} (Committee selection). The number of ways to choose
\(3\) students from \(10\) students for a committee is
\[\binom{10}{3}=120.\]

\begin{tcolorbox}[enhanced jigsaw, arc=.35mm, bottomrule=.15mm, bottomtitle=1mm, breakable, colback=white, colbacktitle=quarto-callout-note-color!10!white, colframe=quarto-callout-note-color-frame, coltitle=black, left=2mm, leftrule=.75mm, opacityback=0, opacitybacktitle=0.6, rightrule=.15mm, title=\textcolor{quarto-callout-note-color}{\faInfo}\hspace{0.5em}{Solution}, titlerule=0mm, toprule=.15mm, toptitle=1mm]

Here order does not matter: choosing Alice, Bob, and Carol is the same
committee in any order. Thus the count is
\(\binom{10}{3}=10!/(7!3!)=120\).

\end{tcolorbox}

\subsection{Four basic counting
situations}\label{four-basic-counting-situations}

This subsection organizes counting problems by whether replacement is
allowed and whether order matters.

{\def\LTcaptype{none} % do not increment counter
\begin{longtable}[]{@{}llll@{}}
\toprule\noalign{}
\textbf{Replacement?} & \textbf{Order matters?} & \textbf{Number of
ways} & \textbf{Example} \\
\midrule\noalign{}
\endhead
\bottomrule\noalign{}
\endlastfoot
Yes & Yes & \(n^r\) & -digit codes using digits 0--9: \(10^3=1000\) \\
No & Yes & \(P_r^n=\dfrac{n!}{(n-r)!}\) & -letter arrangements from 5
letters \\
No & No & \(\binom{n}{r}\) & Choose 3 students from 10 \\
Yes & No & \(\binom{n+r-1}{r}\) & Choose 3 ice-cream scoops from 5
flavors, repetition allowed \\
\end{longtable}
}

\textbf{Example 34} (Stars and bars: ice cream scoops). Choosing \(3\)
scoops from \(5\) flavors with repetition allowed and order not
important is the same as distributing \(3\) identical scoops into \(5\)
flavor bins. The number of choices is
\[\binom{5+3-1}{3}=\binom{7}{3}=35.\]

\begin{tcolorbox}[enhanced jigsaw, arc=.35mm, bottomrule=.15mm, bottomtitle=1mm, breakable, colback=white, colbacktitle=quarto-callout-note-color!10!white, colframe=quarto-callout-note-color-frame, coltitle=black, left=2mm, leftrule=.75mm, opacityback=0, opacitybacktitle=0.6, rightrule=.15mm, title=\textcolor{quarto-callout-note-color}{\faInfo}\hspace{0.5em}{Solution}, titlerule=0mm, toprule=.15mm, toptitle=1mm]

Let \(x_j\) be the number of scoops of flavor \(j\). We need nonnegative
integer solutions to \(x_1+x_2+x_3+x_4+x_5=3\). By stars and bars, the
number of solutions is
\(\binom{3+5-1}{5-1}=\binom{7}{4}=\binom{7}{3}=35\).

\end{tcolorbox}

\section{Frequency Interpretation of
Probability}\label{frequency-interpretation-of-probability}

The goal of this section is to connect probability with long-run
relative frequencies from repeated experiments.

\textbf{Definition 35} (Frequency definition). The probability of an
event can be interpreted as the limiting relative frequency with which
that event occurs in a large number of repeated trials under identical
conditions:
\[\mathbb{P}(A)\approx \frac{\text{number of times event $A$ occurs}}{\text{total number of repeated trials}}.\]
In the limit of many trials, this relative frequency stabilizes to the
probability.

\textbf{Example 36} (Coin toss frequency). If a fair coin is flipped
many times, the relative frequency of heads should be close to \(0.5\)
in the long run. This does not mean that every short sequence must have
half heads and half tails. Rather, it means that the long-run relative
frequency tends to stabilize around \(1/2\).

\begin{tcolorbox}[enhanced jigsaw, arc=.35mm, bottomrule=.15mm, bottomtitle=1mm, breakable, colback=white, colbacktitle=quarto-callout-note-color!10!white, colframe=quarto-callout-note-color-frame, coltitle=black, left=2mm, leftrule=.75mm, opacityback=0, opacitybacktitle=0.6, rightrule=.15mm, title=\textcolor{quarto-callout-note-color}{\faInfo}\hspace{0.5em}{Solution}, titlerule=0mm, toprule=.15mm, toptitle=1mm]

In a short run, the proportion of heads can deviate substantially from
\(1/2\). In many repeated flips of a fair coin, however, the law of
large numbers says the relative frequency of heads tends to \(1/2\).

\end{tcolorbox}

\begin{tcolorbox}[enhanced jigsaw, arc=.35mm, bottomrule=.15mm, bottomtitle=1mm, breakable, colback=white, colbacktitle=quarto-callout-tip-color!10!white, colframe=quarto-callout-tip-color-frame, coltitle=black, left=2mm, leftrule=.75mm, opacityback=0, opacitybacktitle=0.6, rightrule=.15mm, title=\textcolor{quarto-callout-tip-color}{\faLightbulb}\hspace{0.5em}{Learning box}, titlerule=0mm, toprule=.15mm, toptitle=1mm]

Classical vs.~frequency interpretation

\begin{itemize}
\item
  The \textbf{classical} definition is based on equally likely finite
  outcomes.
\item
  The \textbf{frequency} interpretation is based on repeated experiments
  and long-run relative frequencies.
\end{itemize}

\end{tcolorbox}

\section{Conditional Probability}\label{conditional-probability}

The goal of this section is to update probabilities after new
information is observed.

Conditional probability is the probability of an event after we learn
that another event has occurred.

\textbf{Definition 37} (Conditional probability). Let \(A\) and \(B\) be
events with \(\mathbb{P}(B)>0\). The \textbf{conditional probability of
\(A\) given \(B\)} is
\[\mathbb{P}(A\mid B)=\frac{\mathbb{P}(A\cap B)}{\mathbb{P}(B)}.\]

\begin{remark}
\emph{Remark 38}. For fixed \(B\) with \(\mathbb{P}(B)>0\), the function
\[A\mapsto \mathbb{P}(A\mid B)\] is itself a probability measure on the
same sample space. In other words, conditional probability is ordinary
probability after restricting attention to the event \(B\).
\end{remark}

\textbf{Example 39} (Die roll given even). Roll a fair six-sided die
once. Let \[A=\{6\},\qquad B=\{2,4,6\}.\] The event \(B\) is ``the
result is even.'' We want the probability that the result is \(6\) given
that the result is even:
\[\mathbb{P}(A\mid B)=\frac{\mathbb{P}(A\cap B)}{\mathbb{P}(B)}=\frac{\mathbb{P}(\{6\})}{\mathbb{P}(\{2,4,6\})}=\frac{1/6}{3/6}=\frac13.\]

\begin{tcolorbox}[enhanced jigsaw, arc=.35mm, bottomrule=.15mm, bottomtitle=1mm, breakable, colback=white, colbacktitle=quarto-callout-note-color!10!white, colframe=quarto-callout-note-color-frame, coltitle=black, left=2mm, leftrule=.75mm, opacityback=0, opacitybacktitle=0.6, rightrule=.15mm, title=\textcolor{quarto-callout-note-color}{\faInfo}\hspace{0.5em}{Solution}, titlerule=0mm, toprule=.15mm, toptitle=1mm]

Conditioning on \(B=\{2,4,6\}\) restricts the possible outcomes to three
equally likely even results. Among these, only one result is \(6\). Thus
the conditional probability is \(1/3\).

\end{tcolorbox}

\textbf{Example 40} (Two dice). Roll a fair six-sided die twice. Let
\[A=\{\text{the first die is even}\},\qquad B=\{\text{the second die is }6\}.\]
There are \(36\) equally likely ordered outcomes. We have
\[\mathbb{P}(A)=\frac{18}{36}=\frac12,
\qquad
\mathbb{P}(B)=\frac{6}{36}=\frac16,\] and
\[\mathbb{P}(A\cap B)=\frac{3}{36}=\frac{1}{12}.\] Thus
\[\mathbb{P}(A\mid B)=\frac{\mathbb{P}(A\cap B)}{\mathbb{P}(B)}=\frac{1/12}{1/6}=\frac12=\mathbb{P}(A).\]
Knowing the second die is \(6\) does not change the probability that the
first die is even.

\begin{tcolorbox}[enhanced jigsaw, arc=.35mm, bottomrule=.15mm, bottomtitle=1mm, breakable, colback=white, colbacktitle=quarto-callout-note-color!10!white, colframe=quarto-callout-note-color-frame, coltitle=black, left=2mm, leftrule=.75mm, opacityback=0, opacitybacktitle=0.6, rightrule=.15mm, title=\textcolor{quarto-callout-note-color}{\faInfo}\hspace{0.5em}{Solution}, titlerule=0mm, toprule=.15mm, toptitle=1mm]

There are \(36\) equally likely ordered outcomes. Event \(A\) contains
\(3\cdot6=18\) outcomes, event \(B\) contains \(6\) outcomes, and
\(A\cap B\) contains \((2,6),(4,6),(6,6)\), so
\(\mathbb{P}(A\cap B)=3/36=1/12\). Hence
\(\mathbb{P}(A\mid B)=(1/12)/(1/6)=1/2\).

\end{tcolorbox}

\section{Independence}\label{independence}

The goal of this section is to formalize when information about one
event does not change the probability of another event.

\textbf{Definition 41} (Independence of two events). Two events \(A\)
and \(B\) are called \textbf{independent} if
\[\mathbb{P}(A\cap B)=\mathbb{P}(A)\mathbb{P}(B).\]

\textbf{Proposition 42} (Independence and conditional probability).
\emph{Suppose \(\mathbb{P}(A)>0\) and \(\mathbb{P}(B)>0\). Then the
following are equivalent:}

\begin{enumerate}
\def\labelenumi{\arabic{enumi}.}
\item
  \emph{\(A\) and \(B\) are independent;}
\item
  \emph{\(\mathbb{P}(A\mid B)=\mathbb{P}(A)\);}
\item
  \emph{\(\mathbb{P}(B\mid A)=\mathbb{P}(B)\).}
\end{enumerate}

\begin{proof}
\emph{Proof.} If \(A\) and \(B\) are independent, then
\[\mathbb{P}(A\mid B)=\frac{\mathbb{P}(A\cap B)}{\mathbb{P}(B)}=\frac{\mathbb{P}(A)\mathbb{P}(B)}{\mathbb{P}(B)}=\mathbb{P}(A).\]
Similarly, \(\mathbb{P}(B\mid A)=\mathbb{P}(B)\). Conversely, if
\(\mathbb{P}(A\mid B)=\mathbb{P}(A)\), then
\[\frac{\mathbb{P}(A\cap B)}{\mathbb{P}(B)}=\mathbb{P}(A),\] so
\(\mathbb{P}(A\cap B)=\mathbb{P}(A)\mathbb{P}(B)\).~◻
\end{proof}

\begin{tcolorbox}[enhanced jigsaw, arc=.35mm, bottomrule=.15mm, bottomtitle=1mm, breakable, colback=white, colbacktitle=quarto-callout-tip-color!10!white, colframe=quarto-callout-tip-color-frame, coltitle=black, left=2mm, leftrule=.75mm, opacityback=0, opacitybacktitle=0.6, rightrule=.15mm, title=\textcolor{quarto-callout-tip-color}{\faLightbulb}\hspace{0.5em}{Learning box}, titlerule=0mm, toprule=.15mm, toptitle=1mm]

Interpretation of independence Events \(A\) and \(B\) are independent if
knowing that one occurred does not change the probability of the other.

\end{tcolorbox}

\textbf{Example 43} (Two dice are independent). For the two-dice example
above,
\[\mathbb{P}(A\cap B)=\frac{1}{12}=\frac12\cdot\frac16=\mathbb{P}(A)\mathbb{P}(B).\]
Therefore \(A\) and \(B\) are independent.

\begin{tcolorbox}[enhanced jigsaw, arc=.35mm, bottomrule=.15mm, bottomtitle=1mm, breakable, colback=white, colbacktitle=quarto-callout-note-color!10!white, colframe=quarto-callout-note-color-frame, coltitle=black, left=2mm, leftrule=.75mm, opacityback=0, opacitybacktitle=0.6, rightrule=.15mm, title=\textcolor{quarto-callout-note-color}{\faInfo}\hspace{0.5em}{Solution}, titlerule=0mm, toprule=.15mm, toptitle=1mm]

Using the previous computation, \(\mathbb{P}(A\cap B)=1/12\) and
\(\mathbb{P}(A)\mathbb{P}(B)=(1/2)(1/6)=1/12\). Since the two quantities
are equal, \(A\) and \(B\) are independent.

\end{tcolorbox}

\begin{tcolorbox}[enhanced jigsaw, arc=.35mm, bottomrule=.15mm, bottomtitle=1mm, breakable, colback=white, colbacktitle=quarto-callout-warning-color!10!white, colframe=quarto-callout-warning-color-frame, coltitle=black, left=2mm, leftrule=.75mm, opacityback=0, opacitybacktitle=0.6, rightrule=.15mm, title=\textcolor{quarto-callout-warning-color}{\faExclamationTriangle}\hspace{0.5em}{Warning}, titlerule=0mm, toprule=.15mm, toptitle=1mm]

Disjoint is not the same as independent If \(A\) and \(B\) are disjoint
and both have positive probability, then \[\mathbb{P}(A\cap B)=0,\] but
\[\mathbb{P}(A)\mathbb{P}(B)>0.\] Thus disjoint events with positive
probabilities are not independent. Mutually exclusive means they cannot
happen together. Independent means one does not provide information
about the other.

\end{tcolorbox}

\section{Law of Total Probability and Bayes'
Theorem}\label{law-of-total-probability-and-bayes-theorem}

The goal of this section is to decompose probabilities through cases and
then reverse conditional probabilities using Bayes' theorem.

\subsection{Law of total probability}\label{law-of-total-probability}

This subsection gives the finite form of total probability using a
partition of the sample space.

\textbf{Theorem 44} (Law of total probability). \emph{Let
\(A_1,A_2,\ldots,A_n\) be a finite partition of \(S\); that is,
\[S=\bigcup_{i=1}^{n}A_i,\qquad A_i\cap A_j=\varnothing\quad \text{for }i\ne j.\]
Then, for any event \(B\),
\[\mathbb{P}(B)=\sum_{i=1}^{n}\mathbb{P}(B\cap A_i)=\sum_{i=1}^{n}\mathbb{P}(B\mid A_i)\mathbb{P}(A_i),\]
provided \(\mathbb{P}(A_i)>0\) for all \(i\) in the second expression.}

\begin{proof}
\emph{Proof.} Since \(A_1,\ldots,A_n\) partition \(S\), the sets
\[B\cap A_1,\ldots,B\cap A_n\] are pairwise disjoint and their union is
\(B\). Hence \[\mathbb{P}(B)=\sum_{i=1}^{n}\mathbb{P}(B\cap A_i).\]
Using the definition of conditional probability,
\[\mathbb{P}(B\cap A_i)=\mathbb{P}(B\mid A_i)\mathbb{P}(A_i).\]
Substituting gives the result.~◻
\end{proof}

\subsection{Bayes' theorem}\label{bayes-theorem}

This subsection shows how to reverse conditional probabilities using
total probability.

\textbf{Theorem 45} (Bayes' theorem). \emph{Let \(A_1,A_2,\ldots,A_n\)
be a finite partition of \(S\). Suppose \(\mathbb{P}(A_i)>0\) for all
\(i\) and \(\mathbb{P}(B)>0\). Then, for each \(j=1,\ldots,n\),
\[\mathbb{P}(A_j\mid B)
=\frac{\mathbb{P}(B\mid A_j)\mathbb{P}(A_j)}{\sum_{i=1}^{n}\mathbb{P}(B\mid A_i)\mathbb{P}(A_i)}.\]}

\begin{proof}
\emph{Proof.} By conditional probability,
\[\mathbb{P}(A_j\mid B)=\frac{\mathbb{P}(A_j\cap B)}{\mathbb{P}(B)}.\]
The numerator is
\[\mathbb{P}(A_j\cap B)=\mathbb{P}(B\mid A_j)\mathbb{P}(A_j),\] and by
the law of total probability,
\[\mathbb{P}(B)=\sum_{i=1}^{n}\mathbb{P}(B\mid A_i)\mathbb{P}(A_i).\]
Combining these identities gives Bayes' theorem.~◻
\end{proof}

\begin{tcolorbox}[enhanced jigsaw, arc=.35mm, bottomrule=.15mm, bottomtitle=1mm, breakable, colback=white, colbacktitle=quarto-callout-tip-color!10!white, colframe=quarto-callout-tip-color-frame, coltitle=black, left=2mm, leftrule=.75mm, opacityback=0, opacitybacktitle=0.6, rightrule=.15mm, title=\textcolor{quarto-callout-tip-color}{\faLightbulb}\hspace{0.5em}{Learning box}, titlerule=0mm, toprule=.15mm, toptitle=1mm]

Prior, likelihood, posterior Bayes' theorem is often read as
\[\text{posterior}\propto \text{likelihood}\times \text{prior}.\] Here
\(\mathbb{P}(A_j)\) is the prior probability of hypothesis \(A_j\),
\(\mathbb{P}(B\mid A_j)\) is the likelihood of observing evidence \(B\)
under hypothesis \(A_j\), and \(\mathbb{P}(A_j\mid B)\) is the posterior
probability after observing \(B\).

\end{tcolorbox}

\section{Applications of Bayes' Rule}\label{applications-of-bayes-rule}

The goal of this section is to practice Bayes' theorem in examples where
intuition can be misleading.

\subsection{Medical testing example}\label{medical-testing-example}

This subsection applies Bayes' theorem to a binary classification
problem with sensitivity, specificity, and prevalence.

\textbf{Example 46} (Testing for Covid-19). Let
\[D=1 \quad \text{mean infected},\qquad D=0 \quad \text{mean not infected},\]
and
\[Y=1 \quad \text{mean test positive},\qquad Y=0 \quad \text{mean test negative}.\]
The main quantities are: \[\begin{aligned}
\text{sensitivity} &= \mathbb{P}(Y=1\mid D=1),\\
\text{specificity} &= \mathbb{P}(Y=0\mid D=0),\\
\text{prevalence} &= \mathbb{P}(D=1).
\end{aligned}\] Suppose \[\mathbb{P}(D=1)=0.01,
\qquad
\mathbb{P}(Y=1\mid D=1)=0.875,
\qquad
\mathbb{P}(Y=0\mid D=0)=0.975.\] Then
\[\mathbb{P}(Y=1\mid D=0)=1-0.975=0.025,\] and \[\mathbb{P}(D=0)=0.99.\]
If a person tests positive, the probability that the person is truly
infected is \[\begin{aligned}
\mathbb{P}(D=1\mid Y=1)
&=\frac{\mathbb{P}(Y=1\mid D=1)\mathbb{P}(D=1)}{\mathbb{P}(Y=1\mid D=1)\mathbb{P}(D=1)+\mathbb{P}(Y=1\mid D=0)\mathbb{P}(D=0)}\\
&=\frac{0.875\cdot0.01}{0.875\cdot0.01+0.025\cdot0.99}\\
&\approx 0.2612.
\end{aligned}\] Thus the probability is about \(26\%\).

Similarly, the probability that a person is infected given a negative
test is \[\begin{aligned}
\mathbb{P}(D=1\mid Y=0)
&=\frac{\mathbb{P}(Y=0\mid D=1)\mathbb{P}(D=1)}{\mathbb{P}(Y=0\mid D=1)\mathbb{P}(D=1)+\mathbb{P}(Y=0\mid D=0)\mathbb{P}(D=0)}\\
&=\frac{0.125\cdot0.01}{0.125\cdot0.01+0.975\cdot0.99}\\
&\approx 0.00129.
\end{aligned}\] Thus the probability is about \(0.13\%\).

\begin{tcolorbox}[enhanced jigsaw, arc=.35mm, bottomrule=.15mm, bottomtitle=1mm, breakable, colback=white, colbacktitle=quarto-callout-note-color!10!white, colframe=quarto-callout-note-color-frame, coltitle=black, left=2mm, leftrule=.75mm, opacityback=0, opacitybacktitle=0.6, rightrule=.15mm, title=\textcolor{quarto-callout-note-color}{\faInfo}\hspace{0.5em}{Solution}, titlerule=0mm, toprule=.15mm, toptitle=1mm]

Apply Bayes' theorem with the disease status as the partition. The
positive predictive probability is approximately \(0.2612\), so a
positive test gives about a \(26\%\) chance of infection under these
assumptions. The negative-test infection probability is about
\(0.00129\), or \(0.13\%\).

\end{tcolorbox}

\begin{tcolorbox}[enhanced jigsaw, arc=.35mm, bottomrule=.15mm, bottomtitle=1mm, breakable, colback=white, colbacktitle=quarto-callout-warning-color!10!white, colframe=quarto-callout-warning-color-frame, coltitle=black, left=2mm, leftrule=.75mm, opacityback=0, opacitybacktitle=0.6, rightrule=.15mm, title=\textcolor{quarto-callout-warning-color}{\faExclamationTriangle}\hspace{0.5em}{Warning}, titlerule=0mm, toprule=.15mm, toptitle=1mm]

Base-rate effect Even with a good test, a positive result may not mean
the person is very likely infected when the disease prevalence is low.
Bayes' theorem correctly combines test accuracy with the base rate.

\end{tcolorbox}

\subsection{Monty Hall problem}\label{monty-hall-problem}

This subsection uses total probability to explain why switching doors is
advantageous.

\textbf{Example 47} (Monty Hall problem). There are three doors. Behind
one door is a prize, and behind the other two doors are goats. You first
choose one door. The host, who knows where the prize is, opens one of
the other doors and shows a goat. You are then offered the chance to
switch. Should you switch?

Let \[C=\{\text{the first choice is correct}\},\] so
\(\mathbb{P}(C)=1/3\) and \(\mathbb{P}(C^c)=2/3\). Let
\[S=\{\text{win after switching}\},\qquad T=\{\text{win without switching}\}.\]
If your first choice is correct, switching loses, so
\(\mathbb{P}(S\mid C)=0\). If your first choice is incorrect, switching
wins, so \(\mathbb{P}(S\mid C^c)=1\). Therefore
\[\mathbb{P}(S)=\mathbb{P}(S\mid C)\mathbb{P}(C)+\mathbb{P}(S\mid C^c)\mathbb{P}(C^c)
=0\cdot\frac13+1\cdot\frac23=\frac23.\] Similarly,
\[\mathbb{P}(T)=\mathbb{P}(T\mid C)\mathbb{P}(C)+\mathbb{P}(T\mid C^c)\mathbb{P}(C^c)
=1\cdot\frac13+0\cdot\frac23=\frac13.\] \textbf{Conclusion:} you should
switch.

\begin{tcolorbox}[enhanced jigsaw, arc=.35mm, bottomrule=.15mm, bottomtitle=1mm, breakable, colback=white, colbacktitle=quarto-callout-note-color!10!white, colframe=quarto-callout-note-color-frame, coltitle=black, left=2mm, leftrule=.75mm, opacityback=0, opacitybacktitle=0.6, rightrule=.15mm, title=\textcolor{quarto-callout-note-color}{\faInfo}\hspace{0.5em}{Solution}, titlerule=0mm, toprule=.15mm, toptitle=1mm]

Switching wins exactly when the original choice was wrong. The original
choice is wrong with probability \(2/3\), so switching wins with
probability \(2/3\). Staying wins exactly when the original choice was
correct, which has probability \(1/3\).

\end{tcolorbox}

\subsection{Bayesian inference
example}\label{bayesian-inference-example}

This subsection shows how data update prior beliefs into posterior
probabilities.

\textbf{Example 48} (Can someone always toss heads?). Suppose I tell you
that I can toss a coin so that it always comes up heads. You are
initially \(95\%\) certain that I am lying. Then I toss a coin \(5\)
times in front of you and get heads every time. How certain are you now
that I am lying?

Define two hypotheses: \[H_1=\{\text{I am lying; the coin is fair}\},
\qquad
H_2=\{\text{I can always toss heads}\}.\] The prior probabilities are
\[\mathbb{P}(H_1)=0.95,
\qquad
\mathbb{P}(H_2)=0.05.\] Let the data be
\[D=\{\text{5 tosses and 5 heads}\}.\] Under \(H_1\), the coin is fair,
so by independence, \[\mathbb{P}(D\mid H_1)=\left(\frac12\right)^5.\]
Under \(H_2\), heads always occurs, so \[\mathbb{P}(D\mid H_2)=1.\] By
Bayes' theorem, \[\begin{aligned}
\mathbb{P}(H_1\mid D)
&=\frac{\mathbb{P}(D\mid H_1)\mathbb{P}(H_1)}{\mathbb{P}(D\mid H_1)\mathbb{P}(H_1)+\mathbb{P}(D\mid H_2)\mathbb{P}(H_2)}\\
&=\frac{(0.5)^5(0.95)}{(0.5)^5(0.95)+1(0.05)}\\
&\approx 0.3725.
\end{aligned}\] After observing five heads, your probability that I am
lying decreases from \(95\%\) to about \(37\%\).

\begin{tcolorbox}[enhanced jigsaw, arc=.35mm, bottomrule=.15mm, bottomtitle=1mm, breakable, colback=white, colbacktitle=quarto-callout-note-color!10!white, colframe=quarto-callout-note-color-frame, coltitle=black, left=2mm, leftrule=.75mm, opacityback=0, opacitybacktitle=0.6, rightrule=.15mm, title=\textcolor{quarto-callout-note-color}{\faInfo}\hspace{0.5em}{Solution}, titlerule=0mm, toprule=.15mm, toptitle=1mm]

The data \(D\) are much more likely under \(H_2\) than under \(H_1\),
because \(\mathbb{P}(D\mid H_2)=1\) while
\(\mathbb{P}(D\mid H_1)=2^{-5}\). Bayes' theorem gives
\(\mathbb{P}(H_1\mid D)\approx0.3725\), so the belief that the person is
lying drops to about \(37\%\).

\end{tcolorbox}

\begin{remark}
\emph{Remark 49} (Further Bayes examples). Bayes' rule is used in many
applications, including spam filtering, medical diagnosis, Bayesian
inference, and classification. A common machine learning example is
Naive Bayes, where features are assumed conditionally independent given
the class label.
\end{remark}

\section{Conditional Independence}\label{conditional-independence}

The goal of this section is to distinguish ordinary independence from
independence after conditioning on additional information.

Independence and conditional independence are related concepts, but
neither one automatically implies the other.

\textbf{Definition 50} (Conditional independence of events). Two events
\(A\) and \(B\) are \textbf{conditionally independent given \(C\)} if
\(\mathbb{P}(C)>0\) and
\[\mathbb{P}(A\cap B\mid C)=\mathbb{P}(A\mid C)\mathbb{P}(B\mid C).\]

\textbf{Definition 51} (Conditional independence of random variables).
Random variables \(X\) and \(Y\) are \textbf{conditionally independent
given \(Z\)} if, for all relevant \(x,y,z\),
\[F_{X,Y\mid Z}(x,y\mid z)=F_{X\mid Z}(x\mid z)F_{Y\mid Z}(y\mid z).\]

\begin{remark}
\emph{Remark 52}. There is no direct implication between independence
and conditional independence:

\begin{itemize}
\item
  \(A\) and \(B\) can be independent but become dependent after
  conditioning on \(C\).
\item
  \(A\) and \(B\) can be dependent but become independent after
  conditioning on \(C\).
\end{itemize}

This point is fundamental in statistics, causal inference, and machine
learning.
\end{remark}

\section{Paradoxes and Cautions}\label{paradoxes-and-cautions}

The goal of this section is to see how conditioning, aggregation, and
selection can lead to surprising conclusions.

This section is not just about formulas. Probability is also about
learning how intuition can fail.

\subsection{Simpson's paradox}\label{simpsons-paradox}

This subsection illustrates how aggregation can reverse comparisons that
hold within subgroups.

\textbf{Example 53} (Kidney stone treatment). Consider two treatments A
and B for kidney stones. The success rates by stone size are:

{\def\LTcaptype{none} % do not increment counter
\begin{longtable}[]{@{}lcc@{}}
\toprule\noalign{}
\textbf{Stone size} & \textbf{Treatment A} & \textbf{Treatment B} \\
\midrule\noalign{}
\endhead
\bottomrule\noalign{}
\endlastfoot
Small stones & \(93\%\ (81/87)\) & \(87\%\ (234/270)\) \\
Large stones & \(73\%\ (192/263)\) & \(69\%\ (55/80)\) \\
Both & \(78\%\ (273/350)\) & \(83\%\ (289/350)\) \\
\end{longtable}
}

Treatment A has a higher success rate for both small stones and large
stones separately. However, treatment B has a higher overall success
rate. This reversal is called \textbf{Simpson's paradox}. It happens
because the two treatments were applied to different mixtures of small
and large stones.

\begin{tcolorbox}[enhanced jigsaw, arc=.35mm, bottomrule=.15mm, bottomtitle=1mm, breakable, colback=white, colbacktitle=quarto-callout-note-color!10!white, colframe=quarto-callout-note-color-frame, coltitle=black, left=2mm, leftrule=.75mm, opacityback=0, opacitybacktitle=0.6, rightrule=.15mm, title=\textcolor{quarto-callout-note-color}{\faInfo}\hspace{0.5em}{Solution}, titlerule=0mm, toprule=.15mm, toptitle=1mm]

Within each stone-size group, Treatment A has a higher success rate:
\(93\%>87\%\) for small stones and \(73\%>69\%\) for large stones.
Overall, Treatment B has \(289/350\approx83\%\) while Treatment A has
\(273/350\approx78\%\). The reversal occurs because Treatment A was used
more often on large stones, which are harder to treat, while Treatment B
was used more often on small stones.

\end{tcolorbox}

\begin{tcolorbox}[enhanced jigsaw, arc=.35mm, bottomrule=.15mm, bottomtitle=1mm, breakable, colback=white, colbacktitle=quarto-callout-tip-color!10!white, colframe=quarto-callout-tip-color-frame, coltitle=black, left=2mm, leftrule=.75mm, opacityback=0, opacitybacktitle=0.6, rightrule=.15mm, title=\textcolor{quarto-callout-tip-color}{\faLightbulb}\hspace{0.5em}{Learning box}, titlerule=0mm, toprule=.15mm, toptitle=1mm]

Lesson from Simpson's paradox Aggregated data can tell a different story
from stratified data. When comparing groups, always ask whether there
are important hidden variables or confounders.

\end{tcolorbox}

\textbf{Example 54} (Baseball batting averages). A similar reversal can
happen with batting averages. For example, one batter may have a higher
batting average in each of two years separately, while another batter
has a higher combined batting average because the numbers of at-bats are
distributed differently across the years.

\begin{tcolorbox}[enhanced jigsaw, arc=.35mm, bottomrule=.15mm, bottomtitle=1mm, breakable, colback=white, colbacktitle=quarto-callout-note-color!10!white, colframe=quarto-callout-note-color-frame, coltitle=black, left=2mm, leftrule=.75mm, opacityback=0, opacitybacktitle=0.6, rightrule=.15mm, title=\textcolor{quarto-callout-note-color}{\faInfo}\hspace{0.5em}{Solution}, titlerule=0mm, toprule=.15mm, toptitle=1mm]

The combined batting average is a weighted average, with weights
determined by the number of at-bats in each year. Different weights can
reverse the comparison even if one player is better in each separate
year.

\end{tcolorbox}

\subsection{Berkson's paradox}\label{berksons-paradox}

This subsection illustrates how conditioning on selection can create
spurious negative association.

\textbf{Example 55} (Selection bias). Suppose talent and attractiveness
are uncorrelated in the general population. Now suppose we only sample
celebrities, where becoming a celebrity may require some combination of
talent and attractiveness. In this restricted sample, talent and
attractiveness may appear negatively correlated: among celebrities,
someone with lower attractiveness may need higher talent to enter the
sample, and vice versa.

This phenomenon is an example of \textbf{Berkson's paradox}, a
selection-bias effect caused by conditioning on membership in a selected
group.

\begin{tcolorbox}[enhanced jigsaw, arc=.35mm, bottomrule=.15mm, bottomtitle=1mm, breakable, colback=white, colbacktitle=quarto-callout-note-color!10!white, colframe=quarto-callout-note-color-frame, coltitle=black, left=2mm, leftrule=.75mm, opacityback=0, opacitybacktitle=0.6, rightrule=.15mm, title=\textcolor{quarto-callout-note-color}{\faInfo}\hspace{0.5em}{Solution}, titlerule=0mm, toprule=.15mm, toptitle=1mm]

Conditioning on being a celebrity changes the population. If celebrity
status tends to require a high combined level of talent and
attractiveness, then within the selected celebrity group, very high
talent can compensate for lower attractiveness and vice versa. This
creates a negative association in the selected sample even if none
exists in the original population.

\end{tcolorbox}

\begin{tcolorbox}[enhanced jigsaw, arc=.35mm, bottomrule=.15mm, bottomtitle=1mm, breakable, colback=white, colbacktitle=quarto-callout-warning-color!10!white, colframe=quarto-callout-warning-color-frame, coltitle=black, left=2mm, leftrule=.75mm, opacityback=0, opacitybacktitle=0.6, rightrule=.15mm, title=\textcolor{quarto-callout-warning-color}{\faExclamationTriangle}\hspace{0.5em}{Warning}, titlerule=0mm, toprule=.15mm, toptitle=1mm]

Probability intuition warning Monty Hall, Simpson's paradox, and
Berkson's paradox all show the same message: conditioning changes the
probability space. Many mistakes come from forgetting what information
has been conditioned on.

\end{tcolorbox}

\section{Frequentist and Bayesian
Views}\label{frequentist-and-bayesian-views}

The goal of this section is to compare two common interpretations of
probability used throughout statistics.

There are different philosophical interpretations of probability. Both
are useful in statistics.

\begin{tcolorbox}[enhanced jigsaw, arc=.35mm, bottomrule=.15mm, bottomtitle=1mm, breakable, colback=white, colbacktitle=quarto-callout-tip-color!10!white, colframe=quarto-callout-tip-color-frame, coltitle=black, left=2mm, leftrule=.75mm, opacityback=0, opacitybacktitle=0.6, rightrule=.15mm, title=\textcolor{quarto-callout-tip-color}{\faLightbulb}\hspace{0.5em}{Learning box}, titlerule=0mm, toprule=.15mm, toptitle=1mm]

Two interpretations

\begin{itemize}
\item
  \textbf{Frequentist view:} probability is interpreted as a limiting
  frequency over repetitions in identical situations.
\item
  \textbf{Bayesian/subjective view:} probability quantifies degree of
  belief or uncertainty.
\end{itemize}

\end{tcolorbox}

One subjective interpretation is that the probability of an event \(E\)
is the price one is just willing to pay to enter a game in which one
wins one unit of money if \(E\) is true.

\textbf{Example 56} (Betting interpretation). If I believe a coin is
fair and I win \(1\) unit if heads appears, then I would pay \(1/2\)
unit to enter the bet. This corresponds to assigning probability \(1/2\)
to heads.

\begin{tcolorbox}[enhanced jigsaw, arc=.35mm, bottomrule=.15mm, bottomtitle=1mm, breakable, colback=white, colbacktitle=quarto-callout-note-color!10!white, colframe=quarto-callout-note-color-frame, coltitle=black, left=2mm, leftrule=.75mm, opacityback=0, opacitybacktitle=0.6, rightrule=.15mm, title=\textcolor{quarto-callout-note-color}{\faInfo}\hspace{0.5em}{Solution}, titlerule=0mm, toprule=.15mm, toptitle=1mm]

If the payoff is \(1\) unit when heads occurs and the subjective
probability of heads is \(1/2\), then the fair price is expected payoff
\((1/2)\cdot1=1/2\). Paying more than \(1/2\) would be unfavorable,
while paying less would be favorable.

\end{tcolorbox}

\section{Brief History of
Probability}\label{brief-history-of-probability}

The goal of this section is to place the mathematical ideas in
historical context.

\begin{itemize}
\item
  Ideas about chance can be traced back to ancient civilizations and
  Greek philosophy.
\item
  Modern probability began to develop in the study of games of chance.
  Gerolamo Cardano wrote \emph{Book on Games of Chance} in the sixteenth
  century.
\item
  Blaise Pascal and Pierre de Fermat laid groundwork for probability
  theory in the 1650s through problems about gambling and fair division
  of stakes.
\item
  Jacob Bernoulli's work, published in 1713, developed early versions of
  the law of large numbers.
\item
  Pierre-Simon Laplace further developed probability and Bayesian ideas
  in the early nineteenth century.
\item
  In the early twentieth century, R. A. Fisher developed many
  statistical methods for experimental design and data analysis.
\item
  Computational probability, including Monte Carlo simulation, became
  central with the development of modern computing.
\item
  Kolmogorov formalized probability theory using measure-theoretic
  axioms in the 1930s, leading to modern probability spaces,
  martingales, and stochastic processes.
\end{itemize}

\section{Practice Problems}\label{practice-problems}

The goal of this section is to give students a chance to check
understanding through short computations.

\textbf{Practice Problem 57} (Set operations). Let
\(S=\{1,2,3,4,5,6,7,8,9,10\}\), \(A=\{1,2,3,4,5\}\), and
\(B=\{2,4,6,8,10\}\). Find \(A^c\), \(A\cup B\), \(A\cap B\), and
\(A\setminus B\).

\begin{tcolorbox}[enhanced jigsaw, arc=.35mm, bottomrule=.15mm, bottomtitle=1mm, breakable, colback=white, colbacktitle=quarto-callout-note-color!10!white, colframe=quarto-callout-note-color-frame, coltitle=black, left=2mm, leftrule=.75mm, opacityback=0, opacitybacktitle=0.6, rightrule=.15mm, title=\textcolor{quarto-callout-note-color}{\faInfo}\hspace{0.5em}{Solution}, titlerule=0mm, toprule=.15mm, toptitle=1mm]

Since \(S=\{1,2,3,4,5,6,7,8,9,10\}\) and \(A=\{1,2,3,4,5\}\),
\[A^c=\{6,7,8,9,10\}.\] Also, \[A\cup B=\{1,2,3,4,5,6,8,10\},\]
\[A\cap B=\{2,4\},\] and \[A\setminus B=\{1,3,5\}.\]

\end{tcolorbox}

\textbf{Practice Problem 58} (Counting). How many passwords of length
\(4\) can be made using digits \(0\)--\(9\) if repetition is allowed?
How many if repetition is not allowed?

\begin{tcolorbox}[enhanced jigsaw, arc=.35mm, bottomrule=.15mm, bottomtitle=1mm, breakable, colback=white, colbacktitle=quarto-callout-note-color!10!white, colframe=quarto-callout-note-color-frame, coltitle=black, left=2mm, leftrule=.75mm, opacityback=0, opacitybacktitle=0.6, rightrule=.15mm, title=\textcolor{quarto-callout-note-color}{\faInfo}\hspace{0.5em}{Solution}, titlerule=0mm, toprule=.15mm, toptitle=1mm]

With repetition allowed, each of the \(4\) positions has \(10\) choices,
so the number is \[10^4=10000.\] Without repetition, the choices are
\(10,9,8,7\), so the number is \[10\cdot9\cdot8\cdot7=5040.\]

\end{tcolorbox}

\textbf{Practice Problem 59} (Conditional probability). Roll a fair die
once. Given that the result is at least \(4\), what is the probability
that the result is even?

\begin{tcolorbox}[enhanced jigsaw, arc=.35mm, bottomrule=.15mm, bottomtitle=1mm, breakable, colback=white, colbacktitle=quarto-callout-note-color!10!white, colframe=quarto-callout-note-color-frame, coltitle=black, left=2mm, leftrule=.75mm, opacityback=0, opacitybacktitle=0.6, rightrule=.15mm, title=\textcolor{quarto-callout-note-color}{\faInfo}\hspace{0.5em}{Solution}, titlerule=0mm, toprule=.15mm, toptitle=1mm]

Let \(B=\{4,5,6\}\) be the event that the result is at least \(4\), and
let \(A=\{2,4,6\}\) be the event that the result is even. Then
\[A\cap B=\{4,6\}.\] Thus
\[\mathbb{P}(A\mid B)=\frac{|A\cap B|}{|B|}=\frac{2}{3}.\]

\end{tcolorbox}

\textbf{Practice Problem 60} (Independence). Let \(A\) and \(B\) be
events with \(\mathbb{P}(A)=0.4\), \(\mathbb{P}(B)=0.5\), and
\(\mathbb{P}(A\cap B)=0.2\). Are \(A\) and \(B\) independent?

\begin{tcolorbox}[enhanced jigsaw, arc=.35mm, bottomrule=.15mm, bottomtitle=1mm, breakable, colback=white, colbacktitle=quarto-callout-note-color!10!white, colframe=quarto-callout-note-color-frame, coltitle=black, left=2mm, leftrule=.75mm, opacityback=0, opacitybacktitle=0.6, rightrule=.15mm, title=\textcolor{quarto-callout-note-color}{\faInfo}\hspace{0.5em}{Solution}, titlerule=0mm, toprule=.15mm, toptitle=1mm]

Check whether \(\mathbb{P}(A\cap B)=\mathbb{P}(A)\mathbb{P}(B)\). We
have \[\mathbb{P}(A)\mathbb{P}(B)=0.4\cdot0.5=0.2.\] This equals
\(\mathbb{P}(A\cap B)=0.2\), so \(A\) and \(B\) are independent.

\end{tcolorbox}

\textbf{Practice Problem 61} (Bayes' theorem). A disease has prevalence
\(0.02\). A test has sensitivity \(0.90\) and specificity \(0.95\). If a
person tests positive, what is the probability that the person has the
disease?

\begin{tcolorbox}[enhanced jigsaw, arc=.35mm, bottomrule=.15mm, bottomtitle=1mm, breakable, colback=white, colbacktitle=quarto-callout-note-color!10!white, colframe=quarto-callout-note-color-frame, coltitle=black, left=2mm, leftrule=.75mm, opacityback=0, opacitybacktitle=0.6, rightrule=.15mm, title=\textcolor{quarto-callout-note-color}{\faInfo}\hspace{0.5em}{Solution}, titlerule=0mm, toprule=.15mm, toptitle=1mm]

Let \(D\) be the event of disease and \(+\) be the event of a positive
test. Then
\[\mathbb{P}(D)=0.02,\quad \mathbb{P}(+\mid D)=0.90,\quad \mathbb{P}(-\mid D^c)=0.95.\]
So \(\mathbb{P}(+\mid D^c)=0.05\) and \(\mathbb{P}(D^c)=0.98\). By
Bayes' theorem,
\[\mathbb{P}(D\mid+)=\frac{0.90\cdot0.02}{0.90\cdot0.02+0.05\cdot0.98}
=\frac{0.018}{0.067}\approx0.2687.\] The probability is about
\(26.9\%\).

\end{tcolorbox}

\textbf{Practice Problem 62} (Monty Hall). Repeat the Monty Hall
calculation if there are \(4\) doors, one prize, and the host opens two
goat doors after your first choice. What is the probability of winning
if you switch?

\begin{tcolorbox}[enhanced jigsaw, arc=.35mm, bottomrule=.15mm, bottomtitle=1mm, breakable, colback=white, colbacktitle=quarto-callout-note-color!10!white, colframe=quarto-callout-note-color-frame, coltitle=black, left=2mm, leftrule=.75mm, opacityback=0, opacitybacktitle=0.6, rightrule=.15mm, title=\textcolor{quarto-callout-note-color}{\faInfo}\hspace{0.5em}{Solution}, titlerule=0mm, toprule=.15mm, toptitle=1mm]

With \(4\) doors, the initial choice is correct with probability \(1/4\)
and incorrect with probability \(3/4\). If the initial choice is
incorrect and the host opens two goat doors, switching to the only
remaining unopened door wins. Therefore
\[\mathbb{P}(\text{win by switching})=\frac34.\]

\end{tcolorbox}

\section{Summary}\label{summary}

The goal of this section is to collect the main takeaways before moving
to random variables.

\begin{itemize}
\item
  Events are sets; probability is a function on events.
\item
  A probability space is a triple \((S,\mathcal{B},\mathbb{P})\).
\item
  Kolmogorov's axioms are nonnegativity, normalization, and countable
  additivity.
\item
  Classical probability uses counting when the sample space is finite
  and equally likely.
\item
  Conditional probability is
  \(\mathbb{P}(A\mid B)=\mathbb{P}(A\cap B)/\mathbb{P}(B)\).
\item
  Independence means \(\mathbb{P}(A\cap B)=\mathbb{P}(A)\mathbb{P}(B)\).
\item
  The law of total probability decomposes an event using a partition.
\item
  Bayes' theorem updates probabilities after observing evidence.
\item
  Many probability mistakes come from ignoring conditioning, selection
  bias, or hidden variables.
\end{itemize}

\bookmarksetup{startatroot}

\chapter{Chapter 2: Random Variables and
Distributions}\label{chapter-2-random-variables-and-distributions}

This chapter turns probability spaces into numerical models. A random
variable is a function from outcomes to numbers, and its distribution
describes how probability is assigned to those numbers.

\textbf{Topics.} Random variables; cumulative distribution functions;
probability mass functions; probability density functions; discrete,
continuous, and mixed distributions; common distribution families;
exponential families; location-scale families.

\section{From Outcomes to Numbers}\label{from-outcomes-to-numbers}

This section begins the transition from probability spaces to random
variables: instead of studying raw outcomes directly, we often study
numerical summaries of the outcomes.

In Section 1, an experiment was modeled by a probability space
\((S,\mathcal B,\mathbb{P})\). The elements of \(S\) can be very
concrete, such as coin toss strings, or very complicated, such as
continuous paths in Brownian motion. For statistical analysis, we
usually need numbers. A random variable is the mathematical object that
turns outcomes into numbers.

\begin{tcolorbox}[enhanced jigsaw, arc=.35mm, bottomrule=.15mm, bottomtitle=1mm, breakable, colback=white, colbacktitle=quarto-callout-tip-color!10!white, colframe=quarto-callout-tip-color-frame, coltitle=black, left=2mm, leftrule=.75mm, opacityback=0, opacitybacktitle=0.6, rightrule=.15mm, title=\textcolor{quarto-callout-tip-color}{\faLightbulb}\hspace{0.5em}{Learning box}, titlerule=0mm, toprule=.15mm, toptitle=1mm]

\textbf{Main idea.} A random variable is not ``random'' by itself. It is
a function. Randomness enters because the input outcome \(s\in S\) is
random.

\end{tcolorbox}

\section{Random Variables}\label{random-variables}

The goal of this section is to define random variables and show how they
encode the numerical information we care about in an experiment.

\subsection{Informal motivation}\label{informal-motivation}

This subsection explains why random variables are useful before giving
the formal definition.

Suppose we flip a coin twice. The sample space is \[S=\{FF,FT,TF,TT\}.\]
If our question is ``how many heads occurred?'', then we do not really
care whether the outcome was \(FT\) or \(TF\); both have exactly one
head. We care about a function of the outcome:
\[X=\text{number of heads}.\] Thus,
\[X(FF)=0,\qquad X(FT)=1,\qquad X(TF)=1,\qquad X(TT)=2.\]

\textbf{Example 1} (Flipping an unfair coin twice). Suppose a coin has
probability \(\phi\) of landing heads and probability \(1-\phi\) of
landing tails. The coin is flipped twice independently. Let \(X\) be the
number of heads. Find the possible values of \(X\) and the probability
of each value.

\begin{tcolorbox}[enhanced jigsaw, arc=.35mm, bottomrule=.15mm, bottomtitle=1mm, breakable, colback=white, colbacktitle=quarto-callout-note-color!10!white, colframe=quarto-callout-note-color-frame, coltitle=black, left=2mm, leftrule=.75mm, opacityback=0, opacitybacktitle=0.6, rightrule=.15mm, title=\textcolor{quarto-callout-note-color}{\faInfo}\hspace{0.5em}{Solution}, titlerule=0mm, toprule=.15mm, toptitle=1mm]

The sample space is \[S=\{HH,HT,TH,TT\}.\] The random variable \(X\)
takes values \(0,1,2\): \[X(TT)=0,\qquad X(HT)=X(TH)=1,\qquad X(HH)=2.\]
Using independence, \[\mathbb{P}(X=0)=\mathbb{P}(TT)=(1-\phi)^2,\]
\[\mathbb{P}(X=1)=\mathbb{P}(HT)+\mathbb{P}(TH)=\phi(1-\phi)+(1-\phi)\phi=2\phi(1-\phi),\]
\[\mathbb{P}(X=2)=\mathbb{P}(HH)=\phi^2.\] These probabilities add to
\(1\): \[(1-\phi)^2+2\phi(1-\phi)+\phi^2=1.\]

\end{tcolorbox}

\subsection{Formal definition}\label{formal-definition}

This subsection gives the formal definition: a random variable is a
function from the sample space to the real line.

\textbf{Definition 2} (Random variable). Let
\((S,\mathcal B,\mathbb{P})\) be a probability space. A real-valued
random variable is a function \[X:S\to \mathbb{R}.\] More precisely, in
measure-theoretic probability, \(X\) must be measurable, meaning that
events of the form \(\{s\in S:X(s)\le x\}\) are measurable. In this
course, we will use this condition implicitly.

\textbf{Example 3} (Sum of two dice). Toss two fair six-sided dice. Let
\[X=\text{sum of the numbers}.\] Find the sample space, the range of
\(X\), and \(\mathbb{P}(X=7)\).

\begin{tcolorbox}[enhanced jigsaw, arc=.35mm, bottomrule=.15mm, bottomtitle=1mm, breakable, colback=white, colbacktitle=quarto-callout-note-color!10!white, colframe=quarto-callout-note-color-frame, coltitle=black, left=2mm, leftrule=.75mm, opacityback=0, opacitybacktitle=0.6, rightrule=.15mm, title=\textcolor{quarto-callout-note-color}{\faInfo}\hspace{0.5em}{Solution}, titlerule=0mm, toprule=.15mm, toptitle=1mm]

The sample space is \[S=\{(i,j):i,j\in\{1,2,3,4,5,6\}\},\] with \(36\)
equally likely outcomes. The random variable is \[X(i,j)=i+j.\] The
range is \[\operatorname{Range}(X)=\{2,3,4,5,6,7,8,9,10,11,12\}.\] The
event \(\{X=7\}\) consists of \[(1,6),(2,5),(3,4),(4,3),(5,2),(6,1),\]
so \[\mathbb{P}(X=7)=\frac{6}{36}=\frac16.\]

\end{tcolorbox}

\textbf{Example 4} (Fertilizer and crop yield). Suppose an agricultural
experiment applies different amounts of fertilizer to corn plants. Let
\[X=\text{yield per acre}.\] Explain why \(X\) is a random variable.

\begin{tcolorbox}[enhanced jigsaw, arc=.35mm, bottomrule=.15mm, bottomtitle=1mm, breakable, colback=white, colbacktitle=quarto-callout-note-color!10!white, colframe=quarto-callout-note-color-frame, coltitle=black, left=2mm, leftrule=.75mm, opacityback=0, opacitybacktitle=0.6, rightrule=.15mm, title=\textcolor{quarto-callout-note-color}{\faInfo}\hspace{0.5em}{Solution}, titlerule=0mm, toprule=.15mm, toptitle=1mm]

The sample space contains possible outcomes of the agricultural
experiment: weather conditions, soil response, plant growth, measurement
noise, and other uncontrolled factors. The quantity \(X\) assigns to
each outcome a real number, the yield per acre. Therefore
\(X:S\to\mathbb{R}\) is a random variable. In applications, \(X\) would
usually be modeled as continuous because yield can vary over an interval
of values.

\end{tcolorbox}

\subsection{Indicator random
variables}\label{indicator-random-variables}

This subsection introduces one of the most useful random variables in
probability: the indicator of an event.

\textbf{Definition 5} (Indicator random variable). For an event
\(E\subseteq S\), the indicator random variable of \(E\) is
\[\mathbf{1}_E(s)=
\begin{cases}
1, & s\in E,\\
0, & s\notin E.
\end{cases}\] It records whether the event \(E\) occurred.

\textbf{Example 6} (Battery lifetime). Suppose an experiment records how
long a battery operates before wearing down. We are only interested in
whether the battery lasts at least two years. Define
\[E=\{\text{battery lasts at least two years}\}.\] Construct an
indicator random variable for \(E\).

\begin{tcolorbox}[enhanced jigsaw, arc=.35mm, bottomrule=.15mm, bottomtitle=1mm, breakable, colback=white, colbacktitle=quarto-callout-note-color!10!white, colframe=quarto-callout-note-color-frame, coltitle=black, left=2mm, leftrule=.75mm, opacityback=0, opacitybacktitle=0.6, rightrule=.15mm, title=\textcolor{quarto-callout-note-color}{\faInfo}\hspace{0.5em}{Solution}, titlerule=0mm, toprule=.15mm, toptitle=1mm]

Define \[I=\mathbf{1}_E=
\begin{cases}
1, & \text{if the lifetime of the battery is two or more years},\\
0, & \text{otherwise}.
\end{cases}\] The variable \(I\) has range \(\{0,1\}\). It reduces a
potentially continuous lifetime measurement to a binary success/failure
variable.

\end{tcolorbox}

\textbf{Practice Problem 7} (Indicator for a die event). Roll one fair
die and let \(E=\{\text{the result is even}\}\). Define \(\mathbf{1}_E\)
and find \(\mathbb{P}(\mathbf{1}_E=1)\).

\begin{tcolorbox}[enhanced jigsaw, arc=.35mm, bottomrule=.15mm, bottomtitle=1mm, breakable, colback=white, colbacktitle=quarto-callout-note-color!10!white, colframe=quarto-callout-note-color-frame, coltitle=black, left=2mm, leftrule=.75mm, opacityback=0, opacitybacktitle=0.6, rightrule=.15mm, title=\textcolor{quarto-callout-note-color}{\faInfo}\hspace{0.5em}{Solution}, titlerule=0mm, toprule=.15mm, toptitle=1mm]

The sample space is \(S=\{1,2,3,4,5,6\}\). The event is \(E=\{2,4,6\}\).
The indicator is \[\mathbf{1}_E(s)=
\begin{cases}
1, & s\in\{2,4,6\},\\
0, & s\in\{1,3,5\}.
\end{cases}\] Therefore
\[\mathbb{P}(\mathbf{1}_E=1)=\mathbb{P}(E)=\frac{3}{6}=\frac12.\]

\end{tcolorbox}

\section{Probability Distributions}\label{probability-distributions}

This section explains how the probability measure on the sample space
induces probabilities for the values of a random variable.

A random variable \(X\) has two essential ingredients:

\begin{enumerate}
\def\labelenumi{\arabic{enumi}.}
\item
  the set of possible values, \(\operatorname{Range}(X)\);
\item
  the probabilities associated with those values or ranges of values.
\end{enumerate}

The distribution of \(X\) describes how probability is assigned to
values of \(X\).

\textbf{Definition 8} (Probability distribution). The probability
distribution of a random variable \(X\) is the rule that assigns
probabilities to events involving \(X\), such as \(\{X\le x\}\),
\(\{X=x\}\), or \(\{a\le X\le b\}\).

\begin{tcolorbox}[enhanced jigsaw, arc=.35mm, bottomrule=.15mm, bottomtitle=1mm, breakable, colback=white, colbacktitle=quarto-callout-warning-color!10!white, colframe=quarto-callout-warning-color-frame, coltitle=black, left=2mm, leftrule=.75mm, opacityback=0, opacitybacktitle=0.6, rightrule=.15mm, title=\textcolor{quarto-callout-warning-color}{\faExclamationTriangle}\hspace{0.5em}{Warning}, titlerule=0mm, toprule=.15mm, toptitle=1mm]

\textbf{Notation warning.} Strictly speaking, for continuous random
variables, \(\mathbb{P}(X=x)=0\) for every single point \(x\). Therefore
a probability density \(f_X(x)\) is not equal to \(\mathbb{P}(X=x)\); it
is a density whose area gives probability.

\end{tcolorbox}

\section{Cumulative Distribution
Functions}\label{cumulative-distribution-functions}

This section introduces the cumulative distribution function, which is
the most general way to describe the distribution of a real-valued
random variable.

\subsection{Definition of the CDF}\label{definition-of-the-cdf}

This subsection defines the CDF and explains its meaning as an
accumulated probability.

\textbf{Definition 9} (Cumulative distribution function). The cumulative
distribution function (CDF) of a random variable \(X\) is
\[F_X(x)=\mathbb{P}(X\le x)=\mathbb{P}(\{s\in S:X(s)\le x\}).\] When
there is no confusion, we write \(F(x)\) instead of \(F_X(x)\).

The CDF answers the question: ``What is the probability that \(X\) is no
larger than \(x\)?''

\textbf{Example 10} (CDF for number of heads in two fair tosses). Flip a
fair coin twice and let \(X\) be the number of heads. Find the CDF
\(F_X(x)\).

\begin{tcolorbox}[enhanced jigsaw, arc=.35mm, bottomrule=.15mm, bottomtitle=1mm, breakable, colback=white, colbacktitle=quarto-callout-note-color!10!white, colframe=quarto-callout-note-color-frame, coltitle=black, left=2mm, leftrule=.75mm, opacityback=0, opacitybacktitle=0.6, rightrule=.15mm, title=\textcolor{quarto-callout-note-color}{\faInfo}\hspace{0.5em}{Solution}, titlerule=0mm, toprule=.15mm, toptitle=1mm]

From the sample space \(\{HH,HT,TH,TT\}\),
\[\mathbb{P}(X=0)=\frac14,\qquad \mathbb{P}(X=1)=\frac24=\frac12,\qquad \mathbb{P}(X=2)=\frac14.\]
Therefore \[F_X(x)=
\begin{cases}
0, & x<0,\\[2mm]
\frac14, & 0\le x<1,\\[2mm]
\frac34, & 1\le x<2,\\[2mm]
1, & x\ge 2.
\end{cases}\] The CDF is a step function because \(X\) is discrete.

\end{tcolorbox}

\subsection{Properties of CDFs}\label{properties-of-cdfs}

This subsection lists the structural properties that every CDF must
satisfy.

\textbf{Theorem 11} (Basic properties of a CDF). \emph{For any random
variable \(X\), its CDF \(F_X\) satisfies:}

\begin{enumerate}
\def\labelenumi{\arabic{enumi}.}
\item
  \emph{\(\displaystyle \lim_{x\to -\infty}F_X(x)=0\);}
\item
  \emph{\(\displaystyle \lim_{x\to +\infty}F_X(x)=1\);}
\item
  \emph{\(F_X\) is nondecreasing;}
\item
  \emph{\(F_X\) is right-continuous:
  \(\displaystyle \lim_{x\downarrow x_0}F_X(x)=F_X(x_0)\).}
\end{enumerate}

\begin{tcolorbox}[enhanced jigsaw, arc=.35mm, bottomrule=.15mm, bottomtitle=1mm, breakable, colback=white, colbacktitle=quarto-callout-note-color!10!white, colframe=quarto-callout-note-color-frame, coltitle=black, left=2mm, leftrule=.75mm, opacityback=0, opacitybacktitle=0.6, rightrule=.15mm, title=\textcolor{quarto-callout-note-color}{\faInfo}\hspace{0.5em}{Solution}, titlerule=0mm, toprule=.15mm, toptitle=1mm]

We give the intuition for each property. If \(x\) is extremely small,
the event \(\{X\le x\}\) becomes impossible, so the probability tends to
\(0\). If \(x\) is extremely large, the event \(\{X\le x\}\) becomes
almost sure, so the probability tends to \(1\). If \(x_1<x_2\), then
\(\{X\le x_1\}\subseteq\{X\le x_2\}\), so monotonicity of probability
gives \(F_X(x_1)\le F_X(x_2)\). Right-continuity is a consequence of
continuity of probability for decreasing sequences of events.

\end{tcolorbox}

\subsection{Discrete and continuous random
variables}\label{discrete-and-continuous-random-variables}

This subsection distinguishes the two main types of random variables
encountered in the course.

\textbf{Definition 12} (Discrete random variable). A random variable
\(X\) is discrete if its range is finite or countably infinite.
Equivalently, its CDF is a step function.

\textbf{Definition 13} (Continuous random variable). In this course, a
random variable \(X\) is called continuous if its CDF is continuous. In
most examples we study, continuous random variables also have
probability density functions.

\begin{tcolorbox}[enhanced jigsaw, arc=.35mm, bottomrule=.15mm, bottomtitle=1mm, breakable, colback=white, colbacktitle=quarto-callout-tip-color!10!white, colframe=quarto-callout-tip-color-frame, coltitle=black, left=2mm, leftrule=.75mm, opacityback=0, opacitybacktitle=0.6, rightrule=.15mm, title=\textcolor{quarto-callout-tip-color}{\faLightbulb}\hspace{0.5em}{Learning box}, titlerule=0mm, toprule=.15mm, toptitle=1mm]

\textbf{Discrete versus continuous.}

\begin{itemize}
\item
  Discrete random variables assign positive probability to individual
  points.
\item
  Continuous random variables assign probability to intervals through
  area under a density curve.
\end{itemize}

\end{tcolorbox}

\textbf{Practice Problem 14} (Checking a CDF). Is the function \[F(x)=
\begin{cases}
0, & x<0,\\
0.3, & 0\le x<2,\\
0.8, & 2\le x<5,\\
1, & x\ge 5
\end{cases}\] a valid CDF? If yes, find the corresponding point
probabilities.

\begin{tcolorbox}[enhanced jigsaw, arc=.35mm, bottomrule=.15mm, bottomtitle=1mm, breakable, colback=white, colbacktitle=quarto-callout-note-color!10!white, colframe=quarto-callout-note-color-frame, coltitle=black, left=2mm, leftrule=.75mm, opacityback=0, opacitybacktitle=0.6, rightrule=.15mm, title=\textcolor{quarto-callout-note-color}{\faInfo}\hspace{0.5em}{Solution}, titlerule=0mm, toprule=.15mm, toptitle=1mm]

The function is nondecreasing, right-continuous, tends to \(0\) as
\(x\to -\infty\), and tends to \(1\) as \(x\to\infty\). Hence it is a
valid CDF. The jumps give the point probabilities:
\[\mathbb{P}(X=0)=0.3-0=0.3,\] \[\mathbb{P}(X=2)=0.8-0.3=0.5,\]
\[\mathbb{P}(X=5)=1-0.8=0.2.\]

\end{tcolorbox}

\section{Identical Distributions}\label{identical-distributions}

This section explains what it means for two random variables to have the
same probability law, even if they are not the same function on the
sample space.

\textbf{Definition 15} (Identically distributed). Random variables \(X\)
and \(Y\) are identically distributed if
\[\mathbb{P}(X\in A)=\mathbb{P}(Y\in A)\] for every appropriate set
\(A\subseteq\mathbb{R}\). Equivalently,
\[F_X(x)=F_Y(x)\qquad \text{for all }x\in\mathbb{R}.\] We write
\(X\sim Y\) to indicate that \(X\) and \(Y\) have the same distribution.

\textbf{Example 16} (Heads and tails in two fair flips). Flip a fair
coin twice. Let \(X\) be the number of heads and let \(Y\) be the number
of tails. Show that \(X\) and \(Y\) are identically distributed.

\begin{tcolorbox}[enhanced jigsaw, arc=.35mm, bottomrule=.15mm, bottomtitle=1mm, breakable, colback=white, colbacktitle=quarto-callout-note-color!10!white, colframe=quarto-callout-note-color-frame, coltitle=black, left=2mm, leftrule=.75mm, opacityback=0, opacitybacktitle=0.6, rightrule=.15mm, title=\textcolor{quarto-callout-note-color}{\faInfo}\hspace{0.5em}{Solution}, titlerule=0mm, toprule=.15mm, toptitle=1mm]

The sample space is \(\{HH,HT,TH,TT\}\). For \(X=\) number of heads,
\[\mathbb{P}(X=0)=\frac14,
\qquad \mathbb{P}(X=1)=\frac12,
\qquad \mathbb{P}(X=2)=\frac14.\] For \(Y=\) number of tails,
\[\mathbb{P}(Y=0)=\frac14,
\qquad \mathbb{P}(Y=1)=\frac12,
\qquad \mathbb{P}(Y=2)=\frac14.\] Thus \(X\) and \(Y\) have the same
probability mass function and the same CDF. Therefore \(X\sim Y\).
Notice that \(X\) and \(Y\) are not the same random variable, because
for example on outcome \(HH\), \(X=2\) and \(Y=0\).

\end{tcolorbox}

\section{Probability Mass Functions}\label{probability-mass-functions}

This section studies the probability mass function, the standard way to
describe a discrete random variable.

\subsection{Definition of the pmf}\label{definition-of-the-pmf}

This subsection defines the pmf and connects it with probabilities of
individual values.

\textbf{Definition 17} (Probability mass function). For a discrete
random variable \(X\), the probability mass function (pmf) is
\[p_X(x)=\mathbb{P}(X=x)=\mathbb{P}(\{s\in S:X(s)=x\}),\] where \(x\)
ranges over the possible values of \(X\).

\textbf{Theorem 18} (Properties of a pmf). \emph{A function \(p_X\) is a
pmf if it satisfies \[p_X(x)\ge 0\] for all \(x\), and
\[\sum_{x\in\operatorname{Range}(X)}p_X(x)=1.\]}

\begin{tcolorbox}[enhanced jigsaw, arc=.35mm, bottomrule=.15mm, bottomtitle=1mm, breakable, colback=white, colbacktitle=quarto-callout-note-color!10!white, colframe=quarto-callout-note-color-frame, coltitle=black, left=2mm, leftrule=.75mm, opacityback=0, opacitybacktitle=0.6, rightrule=.15mm, title=\textcolor{quarto-callout-note-color}{\faInfo}\hspace{0.5em}{Solution}, titlerule=0mm, toprule=.15mm, toptitle=1mm]

Nonnegativity follows from the nonnegativity axiom of probability. The
values \(\{X=x\}\) for \(x\in\operatorname{Range}(X)\) are disjoint and
their union is the whole sample space, so countable additivity gives
\[1=\mathbb{P}(S)=\sum_{x\in\operatorname{Range}(X)}\mathbb{P}(X=x)=\sum_{x\in\operatorname{Range}(X)}p_X(x).\]

\end{tcolorbox}

\subsection{Relation between pmf and
CDF}\label{relation-between-pmf-and-cdf}

This subsection explains how the jumps of the CDF recover the pmf.

For a discrete random variable, \[F_X(x)=\sum_{t\le x}p_X(t).\] If \(X\)
has a possible value \(k\), then \[p_X(k)=F_X(k)-F_X(k^-),\] where
\[F_X(k^-)=\lim_{\epsilon\downarrow 0}F_X(k-\epsilon).\]

\textbf{Example 19} (Recovering a pmf from a CDF). Suppose \[F(x)=
\begin{cases}
0, & x<1,\\
0.2, & 1\le x<3,\\
0.7, & 3\le x<4,\\
1, & x\ge 4.
\end{cases}\] Find the pmf.

\begin{tcolorbox}[enhanced jigsaw, arc=.35mm, bottomrule=.15mm, bottomtitle=1mm, breakable, colback=white, colbacktitle=quarto-callout-note-color!10!white, colframe=quarto-callout-note-color-frame, coltitle=black, left=2mm, leftrule=.75mm, opacityback=0, opacitybacktitle=0.6, rightrule=.15mm, title=\textcolor{quarto-callout-note-color}{\faInfo}\hspace{0.5em}{Solution}, titlerule=0mm, toprule=.15mm, toptitle=1mm]

The jumps occur at \(1,3,4\). Therefore \[p(1)=F(1)-F(1^-)=0.2-0=0.2,\]
\[p(3)=F(3)-F(3^-)=0.7-0.2=0.5,\] \[p(4)=F(4)-F(4^-)=1-0.7=0.3.\] So
\(X\) takes values \(1,3,4\) with probabilities \(0.2,0.5,0.3\).

\end{tcolorbox}

\subsection{Bernoulli and categorical
variables}\label{bernoulli-and-categorical-variables}

This subsection introduces two basic discrete models: Bernoulli for
binary outcomes and categorical for finite multi-class outcomes.

\textbf{Definition 20} (Bernoulli distribution). A random variable \(X\)
has a Bernoulli distribution with parameter \(\phi\), written
\[X\sim \operatorname{Bernoulli}(\phi),\] if \(X\in\{0,1\}\) and
\[p_X(k)=\phi^k(1-\phi)^{1-k}
=
\begin{cases}
1-\phi, & k=0,\\
\phi, & k=1.
\end{cases}\]

\textbf{Example 21} (Unfair coin). Flip an unfair coin once, where
\(\mathbb{P}(H)=\phi\). Define \(X=1\) if the result is heads and
\(X=0\) otherwise. Find the distribution of \(X\).

\begin{tcolorbox}[enhanced jigsaw, arc=.35mm, bottomrule=.15mm, bottomtitle=1mm, breakable, colback=white, colbacktitle=quarto-callout-note-color!10!white, colframe=quarto-callout-note-color-frame, coltitle=black, left=2mm, leftrule=.75mm, opacityback=0, opacitybacktitle=0.6, rightrule=.15mm, title=\textcolor{quarto-callout-note-color}{\faInfo}\hspace{0.5em}{Solution}, titlerule=0mm, toprule=.15mm, toptitle=1mm]

By definition, \[\mathbb{P}(X=1)=\mathbb{P}(H)=\phi,
\qquad
\mathbb{P}(X=0)=\mathbb{P}(T)=1-\phi.\] Thus
\(X\sim \operatorname{Bernoulli}(\phi)\).

\end{tcolorbox}

\textbf{Definition 22} (Categorical distribution). A random variable
\(Y\) has a categorical distribution with parameters
\((\phi_1,\ldots,\phi_K)\) if
\[Y\in\{1,2,\ldots,K\},\qquad \phi_i\ge 0,\qquad \sum_{i=1}^K\phi_i=1,\]
and \[\mathbb{P}(Y=i)=\phi_i.\] Equivalently,
\[p_Y(y)=\prod_{i=1}^K \phi_i^{\mathbf{1}(y=i)}.\]

\textbf{Example 23} (Unfair six-sided die). Roll an unfair six-sided die
with probabilities \(\phi_1,\ldots,\phi_6\). Write the pmf.

\begin{tcolorbox}[enhanced jigsaw, arc=.35mm, bottomrule=.15mm, bottomtitle=1mm, breakable, colback=white, colbacktitle=quarto-callout-note-color!10!white, colframe=quarto-callout-note-color-frame, coltitle=black, left=2mm, leftrule=.75mm, opacityback=0, opacitybacktitle=0.6, rightrule=.15mm, title=\textcolor{quarto-callout-note-color}{\faInfo}\hspace{0.5em}{Solution}, titlerule=0mm, toprule=.15mm, toptitle=1mm]

The possible values are \(1,2,3,4,5,6\). The pmf is
\[p_Y(y)=\phi_y,\qquad y=1,2,3,4,5,6,\] where \(\phi_i\ge 0\) and
\(\sum_{i=1}^6\phi_i=1\). In indicator form,
\[p_Y(y)=\phi_1^{\mathbf{1}(y=1)}\phi_2^{\mathbf{1}(y=2)}\cdots \phi_6^{\mathbf{1}(y=6)}.\]

\end{tcolorbox}

\section{Probability Density
Functions}\label{probability-density-functions}

This section studies the probability density function, the standard way
to describe most continuous random variables in this course.

\subsection{Definition of a pdf}\label{definition-of-a-pdf}

This subsection defines density through the derivative of the CDF and
through area under a curve.

\textbf{Definition 24} (Probability density function). For an absolutely
continuous random variable \(X\), the probability density function (pdf)
is a function \(f_X\) such that
\[F_X(x)=\mathbb{P}(X\le x)=\int_{-\infty}^x f_X(u)\,du.\] When \(F_X\)
is differentiable, \[f_X(x)=F_X'(x)=\frac{d}{dx}F_X(x).\]

\textbf{Theorem 25} (Properties of a pdf). \emph{A pdf \(f_X\) satisfies
\[f_X(x)\ge 0\] for all \(x\), and
\[\int_{-\infty}^{\infty} f_X(x)\,dx=1.\] For an interval \([a,b]\),
\[\mathbb{P}(a\le X\le b)=\int_a^b f_X(x)\,dx.\]}

\begin{tcolorbox}[enhanced jigsaw, arc=.35mm, bottomrule=.15mm, bottomtitle=1mm, breakable, colback=white, colbacktitle=quarto-callout-note-color!10!white, colframe=quarto-callout-note-color-frame, coltitle=black, left=2mm, leftrule=.75mm, opacityback=0, opacitybacktitle=0.6, rightrule=.15mm, title=\textcolor{quarto-callout-note-color}{\faInfo}\hspace{0.5em}{Solution}, titlerule=0mm, toprule=.15mm, toptitle=1mm]

The density must be nonnegative because probabilities of intervals are
nonnegative. The total area is \(1\) because
\[1=\mathbb{P}(-\infty<X<\infty)=\int_{-\infty}^{\infty}f_X(x)\,dx.\]
The interval formula follows from
\[\mathbb{P}(a\le X\le b)=F_X(b)-F_X(a)=\int_{-\infty}^b f_X(x)\,dx-\int_{-\infty}^a f_X(x)\,dx=\int_a^b f_X(x)\,dx.\]
For continuous random variables, endpoint choices do not affect the
probability because single points have probability zero.

\end{tcolorbox}

\subsection{Logistic CDF and density}\label{logistic-cdf-and-density}

This subsection uses the logistic function as an example of deriving a
pdf from a CDF.

\textbf{Example 26} (Logistic distribution). Suppose a continuous random
variable \(X\) has CDF \[F_X(x)=\frac{1}{1+e^{-x}}.\] Find its pdf.

\begin{tcolorbox}[enhanced jigsaw, arc=.35mm, bottomrule=.15mm, bottomtitle=1mm, breakable, colback=white, colbacktitle=quarto-callout-note-color!10!white, colframe=quarto-callout-note-color-frame, coltitle=black, left=2mm, leftrule=.75mm, opacityback=0, opacitybacktitle=0.6, rightrule=.15mm, title=\textcolor{quarto-callout-note-color}{\faInfo}\hspace{0.5em}{Solution}, titlerule=0mm, toprule=.15mm, toptitle=1mm]

Differentiate the CDF: \[f_X(x)=F_X'(x)=\frac{d}{dx}(1+e^{-x})^{-1}.\]
By the chain rule,
\[f_X(x)=-(1+e^{-x})^{-2}(-e^{-x})=\frac{e^{-x}}{(1+e^{-x})^2}.\] This
density is nonnegative and integrates to
\[\int_{-\infty}^{\infty}f_X(x)\,dx=F_X(\infty)-F_X(-\infty)=1-0=1.\]

\end{tcolorbox}

\subsection{Uniform distribution}\label{uniform-distribution}

This subsection introduces the simplest continuous density: constant
density on an interval.

\textbf{Definition 27} (Continuous uniform distribution). A random
variable \(X\) has a uniform distribution on \([a,b]\), written
\[X\sim \operatorname{Uniform}(a,b),\] if \[f_X(x)=
\begin{cases}
\dfrac{1}{b-a}, & a\le x\le b,\\[2mm]
0, & x<a\text{ or }x>b.
\end{cases}\]

\textbf{Example 28} (Uniform probability). Let
\(X\sim \operatorname{Uniform}(2,8)\). Compute
\(\mathbb{P}(3\le X\le 5)\).

\begin{tcolorbox}[enhanced jigsaw, arc=.35mm, bottomrule=.15mm, bottomtitle=1mm, breakable, colback=white, colbacktitle=quarto-callout-note-color!10!white, colframe=quarto-callout-note-color-frame, coltitle=black, left=2mm, leftrule=.75mm, opacityback=0, opacitybacktitle=0.6, rightrule=.15mm, title=\textcolor{quarto-callout-note-color}{\faInfo}\hspace{0.5em}{Solution}, titlerule=0mm, toprule=.15mm, toptitle=1mm]

The density is \(f_X(x)=1/(8-2)=1/6\) for \(2\le x\le 8\). Therefore
\[\mathbb{P}(3\le X\le 5)=\int_3^5\frac16\,dx=\frac{5-3}{6}=\frac13.\]

\end{tcolorbox}

\textbf{Practice Problem 29} (Find the uniform CDF). Let
\(X\sim \operatorname{Uniform}(a,b)\). Derive the CDF \(F_X(x)\).

\begin{tcolorbox}[enhanced jigsaw, arc=.35mm, bottomrule=.15mm, bottomtitle=1mm, breakable, colback=white, colbacktitle=quarto-callout-note-color!10!white, colframe=quarto-callout-note-color-frame, coltitle=black, left=2mm, leftrule=.75mm, opacityback=0, opacitybacktitle=0.6, rightrule=.15mm, title=\textcolor{quarto-callout-note-color}{\faInfo}\hspace{0.5em}{Solution}, titlerule=0mm, toprule=.15mm, toptitle=1mm]

For \(x<a\), no mass has accumulated, so \(F_X(x)=0\). For
\(a\le x\le b\), \[F_X(x)=\int_a^x \frac{1}{b-a}\,du=\frac{x-a}{b-a}.\]
For \(x>b\), all mass has accumulated, so \(F_X(x)=1\). Hence \[F_X(x)=
\begin{cases}
0, & x<a,\\[1mm]
\dfrac{x-a}{b-a}, & a\le x\le b,\\[2mm]
1, & x>b.
\end{cases}\]

\end{tcolorbox}

\subsection{Mixed distributions}\label{mixed-distributions}

This subsection explains why the CDF is more general than either a pmf
alone or a pdf alone.

\textbf{Example 30} (A random variable with no single pmf/pdf
description). Let \(X\) be defined as follows: with probability \(0.5\),
\(X\) takes the fixed value \(2\); with probability \(0.5\), \(X\) is
drawn uniformly from \([0,1]\). Find the CDF.

\begin{tcolorbox}[enhanced jigsaw, arc=.35mm, bottomrule=.15mm, bottomtitle=1mm, breakable, colback=white, colbacktitle=quarto-callout-note-color!10!white, colframe=quarto-callout-note-color-frame, coltitle=black, left=2mm, leftrule=.75mm, opacityback=0, opacitybacktitle=0.6, rightrule=.15mm, title=\textcolor{quarto-callout-note-color}{\faInfo}\hspace{0.5em}{Solution}, titlerule=0mm, toprule=.15mm, toptitle=1mm]

For \(x\le 0\), \(F_X(x)=0\). For \(0<x\le 1\), only the uniform part
contributes: \[F_X(x)=0.5\cdot x=\frac{x}{2}.\] For \(1<x<2\), the
uniform part is fully accumulated, but the point mass at \(2\) has not
yet occurred: \[F_X(x)=\frac12.\] For \(x\ge 2\), both parts are
included: \[F_X(x)=1.\] Thus \[F_X(x)=
\begin{cases}
0, & x\le 0,\\[1mm]
\dfrac{x}{2}, & 0<x\le 1,\\[2mm]
\dfrac12, & 1<x<2,\\[2mm]
1, & x\ge 2.
\end{cases}\] This distribution has a continuous component on \([0,1]\)
and a point mass at \(2\), so a single ordinary pdf or pmf does not
describe the whole distribution.

\end{tcolorbox}

\subsection{Continuous versus absolutely
continuous}\label{continuous-versus-absolutely-continuous}

This subsection clarifies a technical distinction that will usually be
suppressed in this course.

\textbf{Definition 31} (Continuous random variable). A random variable
\(X\) is continuous if \[\mathbb{P}(X=x)=0\] for all \(x\in\mathbb{R}\).

\textbf{Definition 32} (Absolutely continuous random variable). A random
variable \(X\) is absolutely continuous if there exists a function \(f\)
such that \[\mathbb{P}(X\in A)=\int_A f(x)\,dx\] for every Borel set
\(A\).

Absolutely continuous random variables are continuous, but the converse
is not always true. The Cantor distribution is a standard
counterexample. In this course, almost all continuous random variables
we use are absolutely continuous, so we will usually not distinguish the
two.

\section{Common Discrete
Distributions}\label{common-discrete-distributions}

This section collects the main discrete distribution families used
throughout probability and statistical inference.

\subsection{Binomial distribution}\label{binomial-distribution}

This subsection introduces the binomial distribution as the number of
successes in independent Bernoulli trials.

\textbf{Definition 33} (Binomial distribution). Suppose we perform \(n\)
independent trials, each with two outcomes: success with probability
\(p\) and failure with probability \(1-p\). Let \(X\) be the number of
successes. Then \[X\sim \operatorname{Binomial}(n,p),\] and
\[\mathbb{P}(X=k)=\binom{n}{k}p^k(1-p)^{n-k},\qquad k=0,1,\ldots,n.\]

\textbf{Example 34} (Coin flips). Flip a coin \(n\) times, where the
probability of heads is \(p\). Let \(X\) be the number of heads. Find
\(\mathbb{P}(X=k)\).

\begin{tcolorbox}[enhanced jigsaw, arc=.35mm, bottomrule=.15mm, bottomtitle=1mm, breakable, colback=white, colbacktitle=quarto-callout-note-color!10!white, colframe=quarto-callout-note-color-frame, coltitle=black, left=2mm, leftrule=.75mm, opacityback=0, opacitybacktitle=0.6, rightrule=.15mm, title=\textcolor{quarto-callout-note-color}{\faInfo}\hspace{0.5em}{Solution}, titlerule=0mm, toprule=.15mm, toptitle=1mm]

To get exactly \(k\) heads, choose which \(k\) of the \(n\) trials are
heads. There are \(\binom{n}{k}\) such choices. Each particular sequence
with \(k\) heads and \(n-k\) tails has probability \(p^k(1-p)^{n-k}\).
Therefore \[\mathbb{P}(X=k)=\binom{n}{k}p^k(1-p)^{n-k}.\]

\end{tcolorbox}

\textbf{Example 35} (Airline overbooking). An airline knows that \(5\%\)
of people do not show up for a flight. It sells \(52\) tickets for a
plane with \(50\) seats. Assuming passengers show up independently, what
is the probability that nobody is bumped?

\begin{tcolorbox}[enhanced jigsaw, arc=.35mm, bottomrule=.15mm, bottomtitle=1mm, breakable, colback=white, colbacktitle=quarto-callout-note-color!10!white, colframe=quarto-callout-note-color-frame, coltitle=black, left=2mm, leftrule=.75mm, opacityback=0, opacitybacktitle=0.6, rightrule=.15mm, title=\textcolor{quarto-callout-note-color}{\faInfo}\hspace{0.5em}{Solution}, titlerule=0mm, toprule=.15mm, toptitle=1mm]

Let \(X\) be the number of passengers who show up. Then
\[X\sim \operatorname{Binomial}(52,0.95).\] Nobody is bumped if
\(X\le 50\). Thus
\[\mathbb{P}(\text{nobody bumped})=\mathbb{P}(X\le 50)=1-\mathbb{P}(X=51)-\mathbb{P}(X=52).\]
Compute \[\mathbb{P}(X=51)=\binom{52}{51}(0.95)^{51}(0.05),\]
\[\mathbb{P}(X=52)=(0.95)^{52}.\] Therefore
\[\mathbb{P}(X\le 50)=1-52(0.95)^{51}(0.05)-(0.95)^{52}\approx 0.7405.\]
So the probability that nobody is bumped is about \(74.05\%\).

\end{tcolorbox}

\subsection{Multinomial distribution}\label{multinomial-distribution}

This subsection generalizes the binomial distribution from two
categories to several categories.

\textbf{Definition 36} (Multinomial distribution). Suppose each of \(n\)
independent trials results in one of \(m\) categories \(O_1,\ldots,O_m\)
with probabilities \(\phi_1,\ldots,\phi_m\), where
\(\sum_{i=1}^m\phi_i=1\). Let \(X_i\) be the number of times category
\(O_i\) appears. Then
\[(X_1,\ldots,X_m)\sim \operatorname{Multinomial}(n;\phi_1,\ldots,\phi_m),\]
and for \(n_1+\cdots+n_m=n\),
\[\mathbb{P}(X_1=n_1,\ldots,X_m=n_m)=\frac{n!}{n_1!\cdots n_m!}\phi_1^{n_1}\cdots\phi_m^{n_m}.\]

\textbf{Example 37} (Rolling an \(m\)-sided die). A possibly unfair
\(m\)-sided die has probabilities \(\phi_1,\ldots,\phi_m\). Roll it
\(n\) times. Find the probability of seeing side \(i\) exactly \(n_i\)
times for all \(i\).

\begin{tcolorbox}[enhanced jigsaw, arc=.35mm, bottomrule=.15mm, bottomtitle=1mm, breakable, colback=white, colbacktitle=quarto-callout-note-color!10!white, colframe=quarto-callout-note-color-frame, coltitle=black, left=2mm, leftrule=.75mm, opacityback=0, opacitybacktitle=0.6, rightrule=.15mm, title=\textcolor{quarto-callout-note-color}{\faInfo}\hspace{0.5em}{Solution}, titlerule=0mm, toprule=.15mm, toptitle=1mm]

The vector of counts has a multinomial distribution. Thus, if
\(n_1+\cdots+n_m=n\),
\[\mathbb{P}(X_1=n_1,\ldots,X_m=n_m)=\frac{n!}{n_1!\cdots n_m!}\phi_1^{n_1}\cdots\phi_m^{n_m}.\]
The coefficient counts how many sequences have those category counts.

\end{tcolorbox}

\subsection{Geometric distribution}\label{geometric-distribution}

This subsection introduces the geometric distribution as the waiting
time until the first success.

\textbf{Definition 38} (Geometric distribution). If independent
Bernoulli trials have success probability \(p\), and \(X\) is the trial
number on which the first success occurs, then
\[X\sim \operatorname{Geometric}(p),\] and
\[\mathbb{P}(X=n)=(1-p)^{n-1}p,
\qquad n=1,2,3,\ldots.\]

\textbf{Example 39} (First heads). A coin has probability \(p\) of
heads. Flip until the first heads appears. Find the probability that the
first heads appears on the fifth flip.

\begin{tcolorbox}[enhanced jigsaw, arc=.35mm, bottomrule=.15mm, bottomtitle=1mm, breakable, colback=white, colbacktitle=quarto-callout-note-color!10!white, colframe=quarto-callout-note-color-frame, coltitle=black, left=2mm, leftrule=.75mm, opacityback=0, opacitybacktitle=0.6, rightrule=.15mm, title=\textcolor{quarto-callout-note-color}{\faInfo}\hspace{0.5em}{Solution}, titlerule=0mm, toprule=.15mm, toptitle=1mm]

For the first heads to appear on flip \(5\), the first four flips must
be tails and the fifth must be heads. Therefore
\[\mathbb{P}(X=5)=(1-p)^4p.\]

\end{tcolorbox}

\textbf{Theorem 40} (CDF of a geometric random variable). \emph{If
\(X\sim\operatorname{Geometric}(p)\), then for integer \(x\ge 1\),
\[\mathbb{P}(X\le x)=1-(1-p)^x.\]}

\begin{tcolorbox}[enhanced jigsaw, arc=.35mm, bottomrule=.15mm, bottomtitle=1mm, breakable, colback=white, colbacktitle=quarto-callout-note-color!10!white, colframe=quarto-callout-note-color-frame, coltitle=black, left=2mm, leftrule=.75mm, opacityback=0, opacitybacktitle=0.6, rightrule=.15mm, title=\textcolor{quarto-callout-note-color}{\faInfo}\hspace{0.5em}{Solution}, titlerule=0mm, toprule=.15mm, toptitle=1mm]

Using the geometric sum,
\[\mathbb{P}(X\le x)=\sum_{i=1}^x\mathbb{P}(X=i)=\sum_{i=1}^x(1-p)^{i-1}p.\]
Thus \[\mathbb{P}(X\le x)=p\frac{1-(1-p)^x}{1-(1-p)}=1-(1-p)^x.\]

\end{tcolorbox}

\subsection{Poisson distribution}\label{poisson-distribution}

This subsection introduces the Poisson distribution for counts of rare
or randomly occurring events in time or space.

\textbf{Definition 41} (Poisson distribution). A random variable \(X\)
has a Poisson distribution with rate \(\lambda>0\), written
\[X\sim \operatorname{Poisson}(\lambda),\] if
\[\mathbb{P}(X=k)=\frac{\lambda^k e^{-\lambda}}{k!},
\qquad k=0,1,2,\ldots.\]

\textbf{Example 42} (Telephone calls). A telephone operator handles on
average \(5\) calls every \(3\) minutes. Model calls by a Poisson
process. What is the probability of no calls in the next minute? What is
the probability of at least two calls in the next minute?

\begin{tcolorbox}[enhanced jigsaw, arc=.35mm, bottomrule=.15mm, bottomtitle=1mm, breakable, colback=white, colbacktitle=quarto-callout-note-color!10!white, colframe=quarto-callout-note-color-frame, coltitle=black, left=2mm, leftrule=.75mm, opacityback=0, opacitybacktitle=0.6, rightrule=.15mm, title=\textcolor{quarto-callout-note-color}{\faInfo}\hspace{0.5em}{Solution}, titlerule=0mm, toprule=.15mm, toptitle=1mm]

The average rate per minute is \[\lambda=\frac{5}{3}.\] Let \(X\) be the
number of calls in the next minute. Then
\(X\sim\operatorname{Poisson}(5/3)\). The probability of no calls is
\[\mathbb{P}(X=0)=e^{-5/3}\approx 0.1889.\] The probability of at least
two calls is \[\mathbb{P}(X\ge 2)=1-\mathbb{P}(X=0)-\mathbb{P}(X=1)
=1-e^{-5/3}-\frac{5}{3}e^{-5/3}
\approx 0.4963.\]

\end{tcolorbox}

\textbf{Theorem 43} (Poisson approximation to the binomial). \emph{If
\(n\to\infty\), \(p\to 0\), and \(np=\lambda\) stays constant, then
\[\binom{n}{k}p^k(1-p)^{n-k}\to \frac{\lambda^k e^{-\lambda}}{k!}.\]}

\begin{tcolorbox}[enhanced jigsaw, arc=.35mm, bottomrule=.15mm, bottomtitle=1mm, breakable, colback=white, colbacktitle=quarto-callout-note-color!10!white, colframe=quarto-callout-note-color-frame, coltitle=black, left=2mm, leftrule=.75mm, opacityback=0, opacitybacktitle=0.6, rightrule=.15mm, title=\textcolor{quarto-callout-note-color}{\faInfo}\hspace{0.5em}{Solution}, titlerule=0mm, toprule=.15mm, toptitle=1mm]

Set \(p=\lambda/n\). Then \[\binom{n}{k}p^k(1-p)^{n-k}
=\frac{n(n-1)\cdots(n-k+1)}{k!}\left(\frac{\lambda}{n}\right)^k\left(1-\frac{\lambda}{n}\right)^{n-k}.\]
The first product satisfies \[\frac{n(n-1)\cdots(n-k+1)}{n^k}\to 1,\]
and \[\left(1-\frac{\lambda}{n}\right)^{n-k}\to e^{-\lambda}.\]
Therefore the limit is \[\frac{\lambda^k}{k!}e^{-\lambda}.\]

\end{tcolorbox}

\subsection{Discrete uniform
distribution}\label{discrete-uniform-distribution}

This subsection introduces equal probabilities over a finite set of
integer-valued outcomes.

\textbf{Definition 44} (Discrete uniform distribution). A random
variable \(X\) has a discrete uniform distribution on
\(\{a,a+1,\ldots,b\}\) if \[\mathbb{P}(X=x)=\frac1n,
\qquad x=a,a+1,\ldots,b,\] where \[n=b-a+1.\]

\textbf{Example 45} (Fair six-sided die). Let \(X\) be the result of
rolling a fair six-sided die. Identify the distribution and compute
\(\mathbb{P}(X\ge 5)\).

\begin{tcolorbox}[enhanced jigsaw, arc=.35mm, bottomrule=.15mm, bottomtitle=1mm, breakable, colback=white, colbacktitle=quarto-callout-note-color!10!white, colframe=quarto-callout-note-color-frame, coltitle=black, left=2mm, leftrule=.75mm, opacityback=0, opacitybacktitle=0.6, rightrule=.15mm, title=\textcolor{quarto-callout-note-color}{\faInfo}\hspace{0.5em}{Solution}, titlerule=0mm, toprule=.15mm, toptitle=1mm]

The die roll is discrete uniform on \(\{1,2,3,4,5,6\}\). Therefore
\[\mathbb{P}(X=x)=\frac16,\qquad x=1,2,3,4,5,6.\] The event \(X\ge 5\)
is \(\{5,6\}\), so \[\mathbb{P}(X\ge 5)=\frac26=\frac13.\]

\end{tcolorbox}

\textbf{Example 46} (German tank problem setup). Suppose enemy tanks are
numbered \(1,2,\ldots,N\), and a captured tank has serial number \(X\).
If the captured tank is equally likely to be any of the \(N\) tanks,
write the distribution of \(X\).

\begin{tcolorbox}[enhanced jigsaw, arc=.35mm, bottomrule=.15mm, bottomtitle=1mm, breakable, colback=white, colbacktitle=quarto-callout-note-color!10!white, colframe=quarto-callout-note-color-frame, coltitle=black, left=2mm, leftrule=.75mm, opacityback=0, opacitybacktitle=0.6, rightrule=.15mm, title=\textcolor{quarto-callout-note-color}{\faInfo}\hspace{0.5em}{Solution}, titlerule=0mm, toprule=.15mm, toptitle=1mm]

The random variable \(X\) is discrete uniform on \(\{1,\ldots,N\}\).
Hence \[\mathbb{P}(X=x)=\frac{1}{N},\qquad x=1,2,\ldots,N.\] In the
German tank problem, observations from this distribution are used to
estimate the unknown maximum \(N\).

\end{tcolorbox}

\subsection{Hypergeometric
distribution}\label{hypergeometric-distribution}

This subsection introduces sampling without replacement from a finite
population.

\textbf{Definition 47} (Hypergeometric distribution). Suppose an urn
contains \(N\) balls, \(M\) of which are red and \(N-M\) of which are
green. Draw \(K\) balls without replacement. Let \(X\) be the number of
red balls drawn. Then
\[\mathbb{P}(X=x)=\frac{\binom{M}{x}\binom{N-M}{K-x}}{\binom{N}{K}},\]
for feasible values of \(x\). The mean and variance are
\[\mathbb{E}[X]=\frac{KM}{N},\]
\[\operatorname{Var}(X)=\frac{KM}{N}\left(\frac{(N-M)(N-K)}{N(N-1)}\right).\]

\textbf{Example 48} (Urn without replacement). An urn contains \(20\)
red balls and \(30\) green balls. Draw \(10\) balls without replacement.
What is the probability of drawing exactly \(3\) red balls?

\begin{tcolorbox}[enhanced jigsaw, arc=.35mm, bottomrule=.15mm, bottomtitle=1mm, breakable, colback=white, colbacktitle=quarto-callout-note-color!10!white, colframe=quarto-callout-note-color-frame, coltitle=black, left=2mm, leftrule=.75mm, opacityback=0, opacitybacktitle=0.6, rightrule=.15mm, title=\textcolor{quarto-callout-note-color}{\faInfo}\hspace{0.5em}{Solution}, titlerule=0mm, toprule=.15mm, toptitle=1mm]

Here \(N=50\), \(M=20\), \(K=10\), and \(x=3\). Therefore
\[\mathbb{P}(X=3)=\frac{\binom{20}{3}\binom{30}{7}}{\binom{50}{10}}.\]
This counts favorable samples with exactly \(3\) red and \(7\) green
divided by all samples of size \(10\).

\end{tcolorbox}

\textbf{Example 49} (Acceptance sampling). A retailer receives a lot of
\(N=25\) machine parts, of which \(M=6\) are defective. A sample of
\(K=10\) parts is selected without replacement. What is the probability
that all sampled parts are acceptable?

\begin{tcolorbox}[enhanced jigsaw, arc=.35mm, bottomrule=.15mm, bottomtitle=1mm, breakable, colback=white, colbacktitle=quarto-callout-note-color!10!white, colframe=quarto-callout-note-color-frame, coltitle=black, left=2mm, leftrule=.75mm, opacityback=0, opacitybacktitle=0.6, rightrule=.15mm, title=\textcolor{quarto-callout-note-color}{\faInfo}\hspace{0.5em}{Solution}, titlerule=0mm, toprule=.15mm, toptitle=1mm]

Let \(X\) be the number of defective parts in the sample. Then \(X\) is
hypergeometric with \(N=25\), \(M=6\), \(K=10\). All sampled parts are
acceptable means \(X=0\). Therefore
\[\mathbb{P}(X=0)=\frac{\binom{6}{0}\binom{19}{10}}{\binom{25}{10}}
\approx 0.0342.\] This event is quite unlikely if there are six
defectives in the lot.

\end{tcolorbox}

\textbf{Practice Problem 50} (With replacement versus without
replacement). In the urn example, what distribution would be used if the
balls were drawn with replacement?

\begin{tcolorbox}[enhanced jigsaw, arc=.35mm, bottomrule=.15mm, bottomtitle=1mm, breakable, colback=white, colbacktitle=quarto-callout-note-color!10!white, colframe=quarto-callout-note-color-frame, coltitle=black, left=2mm, leftrule=.75mm, opacityback=0, opacitybacktitle=0.6, rightrule=.15mm, title=\textcolor{quarto-callout-note-color}{\faInfo}\hspace{0.5em}{Solution}, titlerule=0mm, toprule=.15mm, toptitle=1mm]

With replacement, each draw has the same probability of red:
\[p=\frac{M}{N}.\] The draws are independent, so the number of red balls
in \(K\) draws follows
\[X\sim \operatorname{Binomial}\left(K,\frac{M}{N}\right).\] Without
replacement gives the hypergeometric distribution; with replacement
gives the binomial distribution.

\end{tcolorbox}

\subsection{Negative binomial
distribution}\label{negative-binomial-distribution}

This subsection introduces the negative binomial distribution as a
waiting-count model for a fixed number of successes.

\textbf{Definition 51} (Negative binomial distribution: failures before
\(r\) successes). Consider independent Bernoulli trials with success
probability \(p\). Let \(Y\) be the number of failures observed before
the \(r\)-th success. Then
\[\mathbb{P}(Y=y)=\binom{y+r-1}{y}(1-p)^y p^r,
\qquad y=0,1,2,\ldots.\] The mean and variance are
\[\mathbb{E}[Y]=\frac{r(1-p)}{p},
\qquad
\operatorname{Var}(Y)=\frac{r(1-p)}{p^2}.\]

\textbf{Definition 52} (Negative binomial distribution: total trials).
If \(X\) is the total number of trials needed to obtain \(r\) successes,
then \(X=r+Y\) and \[\mathbb{P}(X=n)=\binom{n-1}{r-1}p^r(1-p)^{n-r},
\qquad n=r,r+1,r+2,\ldots.\]

\textbf{Example 53} (Survey completion). Jim goes door to door asking
people to fill out a survey. At each house, there is a \(0.6\)
probability that the survey is completed. Jim must collect \(30\)
completed surveys. What is the probability that the last survey is
completed at the \(n\)-th house?

\begin{tcolorbox}[enhanced jigsaw, arc=.35mm, bottomrule=.15mm, bottomtitle=1mm, breakable, colback=white, colbacktitle=quarto-callout-note-color!10!white, colframe=quarto-callout-note-color-frame, coltitle=black, left=2mm, leftrule=.75mm, opacityback=0, opacitybacktitle=0.6, rightrule=.15mm, title=\textcolor{quarto-callout-note-color}{\faInfo}\hspace{0.5em}{Solution}, titlerule=0mm, toprule=.15mm, toptitle=1mm]

Let \(X\) be the total number of houses visited to obtain \(r=30\)
completed surveys. The last survey is completed at the \(n\)-th house if
among the first \(n-1\) houses there were exactly \(29\) completions,
and the \(n\)-th house is a completion. Therefore
\[\mathbb{P}(X=n)=\binom{n-1}{29}(0.6)^{30}(0.4)^{n-30},
\qquad n=30,31,32,\ldots.\]

\end{tcolorbox}

\textbf{Theorem 54} (Poisson limit of negative binomial). \emph{Under an
appropriate limiting regime, the negative binomial distribution can
converge to a Poisson distribution. In particular, when \(r\to\infty\),
\(p\to 1\), and \[r(1-p)\to \lambda,\] the number of failures before
\(r\) successes converges in distribution to
\(\operatorname{Poisson}(\lambda)\).}

\begin{tcolorbox}[enhanced jigsaw, arc=.35mm, bottomrule=.15mm, bottomtitle=1mm, breakable, colback=white, colbacktitle=quarto-callout-note-color!10!white, colframe=quarto-callout-note-color-frame, coltitle=black, left=2mm, leftrule=.75mm, opacityback=0, opacitybacktitle=0.6, rightrule=.15mm, title=\textcolor{quarto-callout-note-color}{\faInfo}\hspace{0.5em}{Solution}, titlerule=0mm, toprule=.15mm, toptitle=1mm]

The negative binomial model counts rare failures before many successes.
If \(p\to 1\), failures become rare; if \(r\to\infty\) while \(r(1-p)\)
stays near \(\lambda\), the expected number of failures remains finite.
This is the same rare-event principle behind the Poisson approximation.

\end{tcolorbox}

\section{Common Continuous
Distributions}\label{common-continuous-distributions}

This section collects the main continuous distribution families used
throughout probability and statistical inference.

\subsection{Normal distribution}\label{normal-distribution}

This subsection introduces the normal distribution, the central
continuous model in statistics.

\textbf{Definition 55} (Normal distribution). A random variable \(X\)
has a normal distribution with mean \(\mu\) and variance \(\sigma^2\),
written \[X\sim \operatorname{Normal}(\mu,\sigma^2),\] if
\[f_X(x)=\frac{1}{\sqrt{2\pi}\sigma}\exp\left(-\frac{(x-\mu)^2}{2\sigma^2}\right),
\qquad x\in\mathbb{R}.\]

The normal distribution is one of the most important distributions
because of the Central Limit Theorem: sums and averages of many small
independent effects are often approximately normal.

\textbf{Example 56} (Standard normal probability notation). Let
\(Z\sim\operatorname{Normal}(0,1)\). Express \(\mathbb{P}(a\le Z\le b)\)
using the standard normal CDF \(\Phi\).

\begin{tcolorbox}[enhanced jigsaw, arc=.35mm, bottomrule=.15mm, bottomtitle=1mm, breakable, colback=white, colbacktitle=quarto-callout-note-color!10!white, colframe=quarto-callout-note-color-frame, coltitle=black, left=2mm, leftrule=.75mm, opacityback=0, opacitybacktitle=0.6, rightrule=.15mm, title=\textcolor{quarto-callout-note-color}{\faInfo}\hspace{0.5em}{Solution}, titlerule=0mm, toprule=.15mm, toptitle=1mm]

The standard normal CDF is \[\Phi(x)=\mathbb{P}(Z\le x).\] Therefore
\[\mathbb{P}(a\le Z\le b)=\mathbb{P}(Z\le b)-\mathbb{P}(Z<a)=\Phi(b)-\Phi(a).\]
For a continuous random variable, using \(Z<a\) or \(Z\le a\) gives the
same probability.

\end{tcolorbox}

\subsection{Cauchy distribution}\label{cauchy-distribution}

This subsection introduces a heavy-tailed distribution whose mean does
not exist.

\textbf{Definition 57} (Cauchy distribution). A Cauchy random variable
with location \(\mu\) and scale \(\sigma>0\) has density
\[f_X(x)=\frac{1}{\pi\sigma\left[1+\left(\frac{x-\mu}{\sigma}\right)^2\right]}
=\frac{\sigma}{\pi\left[\sigma^2+(x-\mu)^2\right]}.\] The Cauchy
distribution has no mean; the parameter \(\mu\) is the median and
location parameter.

\textbf{Example 58} (Kicking a ball at a random angle). A person stands
distance \(d\) from a straight line and kicks a ball toward the line.
The angle \(\theta\) is uniformly distributed on \((-\pi/2,\pi/2)\). Let
\(X\) be the hitting location on the line, measured from the closest
point. Show why \(X\) has a Cauchy-type density.

\begin{tcolorbox}[enhanced jigsaw, arc=.35mm, bottomrule=.15mm, bottomtitle=1mm, breakable, colback=white, colbacktitle=quarto-callout-note-color!10!white, colframe=quarto-callout-note-color-frame, coltitle=black, left=2mm, leftrule=.75mm, opacityback=0, opacitybacktitle=0.6, rightrule=.15mm, title=\textcolor{quarto-callout-note-color}{\faInfo}\hspace{0.5em}{Solution}, titlerule=0mm, toprule=.15mm, toptitle=1mm]

Geometry gives \[\tan\theta=\frac{X}{d},
\qquad X=d\tan\theta.\] Since
\(\theta\sim\operatorname{Uniform}(-\pi/2,\pi/2)\),
\[f_\theta(\theta)=\frac1\pi.\] Use the transformation
\(x=d\tan\theta\), so \[\theta=\arctan(x/d),
\qquad \frac{d\theta}{dx}=\frac{d}{d^2+x^2}.\] Thus
\[f_X(x)=f_\theta(\arctan(x/d))\left|\frac{d\theta}{dx}\right|
=\frac1\pi\frac{d}{d^2+x^2}.\] This is a Cauchy density with location
\(0\) and scale \(d\).

\end{tcolorbox}

\subsection{Exponential distribution}\label{exponential-distribution}

This subsection introduces the exponential distribution as a
waiting-time model.

\textbf{Definition 59} (Exponential distribution). A random variable
\(X\) has an exponential distribution with rate \(\lambda>0\), written
\[X\sim\operatorname{Exponential}(\lambda),\] if \[f_X(x)=
\begin{cases}
\lambda e^{-\lambda x}, & x\ge 0,\\
0, & x<0.
\end{cases}\] The exponential distribution is the continuous analogue of
the geometric distribution.

\textbf{Example 60} (Time until failure). Suppose the lifetime of a
device follows \(X\sim\operatorname{Exponential}(\lambda)\). Find
\(\mathbb{P}(X>t)\) and \(\mathbb{P}(X\le t)\) for \(t\ge 0\).

\begin{tcolorbox}[enhanced jigsaw, arc=.35mm, bottomrule=.15mm, bottomtitle=1mm, breakable, colback=white, colbacktitle=quarto-callout-note-color!10!white, colframe=quarto-callout-note-color-frame, coltitle=black, left=2mm, leftrule=.75mm, opacityback=0, opacitybacktitle=0.6, rightrule=.15mm, title=\textcolor{quarto-callout-note-color}{\faInfo}\hspace{0.5em}{Solution}, titlerule=0mm, toprule=.15mm, toptitle=1mm]

For \(t\ge 0\),
\[\mathbb{P}(X>t)=\int_t^\infty \lambda e^{-\lambda x}\,dx=e^{-\lambda t}.\]
Therefore \[\mathbb{P}(X\le t)=1-\mathbb{P}(X>t)=1-e^{-\lambda t}.\] The
survival probability decays exponentially.

\end{tcolorbox}

\begin{tcolorbox}[enhanced jigsaw, arc=.35mm, bottomrule=.15mm, bottomtitle=1mm, breakable, colback=white, colbacktitle=quarto-callout-tip-color!10!white, colframe=quarto-callout-tip-color-frame, coltitle=black, left=2mm, leftrule=.75mm, opacityback=0, opacitybacktitle=0.6, rightrule=.15mm, title=\textcolor{quarto-callout-tip-color}{\faLightbulb}\hspace{0.5em}{Learning box}, titlerule=0mm, toprule=.15mm, toptitle=1mm]

Exponential versus Poisson

\begin{itemize}
\item
  Poisson distribution: number of events in a time interval.
\item
  Exponential distribution: waiting time between events or time until
  the first event.
\end{itemize}

\end{tcolorbox}

\subsection{Laplace distribution}\label{laplace-distribution}

This subsection introduces the double exponential distribution, also
called the Laplace distribution.

\textbf{Definition 61} (Laplace distribution). A random variable \(X\)
has a Laplace distribution with location \(\mu\) and scale \(b>0\),
written \(X\sim\operatorname{Laplace}(\mu,b)\), if
\[f_X(x)=\frac{1}{2b}\exp\left(-\frac{|x-\mu|}{b}\right),
\qquad -\infty<x<\infty.\] Its mean and variance are
\[\mathbb{E}[X]=\mu,
\qquad
\operatorname{Var}(X)=2b^2.\]

The Laplace distribution can be viewed as reflecting an exponential
distribution around its center. The maximum likelihood estimator of the
location \(\mu\) is the sample median.

\textbf{Example 62} (Difference of exponential variables). Let \(U\) and
\(V\) be independent \(\operatorname{Exponential}(\lambda)\) random
variables. The difference \(X=U-V\) has a Laplace distribution centered
at \(0\). State its scale parameter.

\begin{tcolorbox}[enhanced jigsaw, arc=.35mm, bottomrule=.15mm, bottomtitle=1mm, breakable, colback=white, colbacktitle=quarto-callout-note-color!10!white, colframe=quarto-callout-note-color-frame, coltitle=black, left=2mm, leftrule=.75mm, opacityback=0, opacitybacktitle=0.6, rightrule=.15mm, title=\textcolor{quarto-callout-note-color}{\faInfo}\hspace{0.5em}{Solution}, titlerule=0mm, toprule=.15mm, toptitle=1mm]

The difference of two independent exponential random variables with
common rate \(\lambda\) follows a Laplace distribution with location
\(0\) and scale \[b=\frac{1}{\lambda}.\] Thus
\[f_X(x)=\frac{\lambda}{2}e^{-\lambda |x|}.\] This explains why the
Laplace distribution is sometimes called the double exponential
distribution.

\end{tcolorbox}

\subsection{Gamma distribution}\label{gamma-distribution}

This subsection introduces the gamma distribution and its connection to
waiting times in Poisson processes.

\textbf{Definition 63} (Gamma distribution). A random variable \(X\) has
a gamma distribution with shape \(\alpha>0\) and scale \(\theta>0\),
written \[X\sim\operatorname{Gamma}(\alpha,\theta),\] if
\[f_X(x;\alpha,\theta)=\frac{x^{\alpha-1}e^{-x/\theta}}{\theta^\alpha\Gamma(\alpha)},
\qquad x\ge 0,\] where the gamma function is
\[\Gamma(\alpha)=\int_0^\infty t^{\alpha-1}e^{-t}\,dt.\]

\begin{tcolorbox}[enhanced jigsaw, arc=.35mm, bottomrule=.15mm, bottomtitle=1mm, breakable, colback=white, colbacktitle=quarto-callout-tip-color!10!white, colframe=quarto-callout-tip-color-frame, coltitle=black, left=2mm, leftrule=.75mm, opacityback=0, opacitybacktitle=0.6, rightrule=.15mm, title=\textcolor{quarto-callout-tip-color}{\faLightbulb}\hspace{0.5em}{Learning box}, titlerule=0mm, toprule=.15mm, toptitle=1mm]

Parameterization warning Some books use the rate parameter
\(\beta=1/\theta\) instead of the scale parameter \(\theta\). Under the
rate parameterization,
\[f_X(x;\alpha,\beta)=\frac{\beta^\alpha}{\Gamma(\alpha)}x^{\alpha-1}e^{-\beta x}.\]

\end{tcolorbox}

\textbf{Example 64} (Exponential as a special gamma). Show that the
exponential distribution is a special case of the gamma distribution.

\begin{tcolorbox}[enhanced jigsaw, arc=.35mm, bottomrule=.15mm, bottomtitle=1mm, breakable, colback=white, colbacktitle=quarto-callout-note-color!10!white, colframe=quarto-callout-note-color-frame, coltitle=black, left=2mm, leftrule=.75mm, opacityback=0, opacitybacktitle=0.6, rightrule=.15mm, title=\textcolor{quarto-callout-note-color}{\faInfo}\hspace{0.5em}{Solution}, titlerule=0mm, toprule=.15mm, toptitle=1mm]

Take \(\alpha=1\) and \(\theta=1/\lambda\). Since \(\Gamma(1)=1\),
\[f_X(x)=\frac{x^{0}e^{-x/(1/\lambda)}}{(1/\lambda)^1\Gamma(1)}=\lambda e^{-\lambda x},
\qquad x\ge 0.\] This is the exponential density with rate \(\lambda\).

\end{tcolorbox}

\textbf{Example 65} (Erlang distribution). Explain why
\(\operatorname{Gamma}(k,\theta)\) with integer \(k\) is called an
Erlang distribution.

\begin{tcolorbox}[enhanced jigsaw, arc=.35mm, bottomrule=.15mm, bottomtitle=1mm, breakable, colback=white, colbacktitle=quarto-callout-note-color!10!white, colframe=quarto-callout-note-color-frame, coltitle=black, left=2mm, leftrule=.75mm, opacityback=0, opacitybacktitle=0.6, rightrule=.15mm, title=\textcolor{quarto-callout-note-color}{\faInfo}\hspace{0.5em}{Solution}, titlerule=0mm, toprule=.15mm, toptitle=1mm]

If \(k\) is a positive integer, then \[\Gamma(k)=(k-1)!.\] In a Poisson
process, the waiting time until the \(k\)-th event is the sum of \(k\)
independent exponential waiting times. This sum has a gamma distribution
with integer shape \(k\), called the Erlang distribution.

\end{tcolorbox}

\textbf{Theorem 66} (Gamma-Poisson relationship). \emph{If
\(X\sim \operatorname{Gamma}(\alpha,\beta)\) under the rate
parameterization, with \(\alpha\) a positive integer, and if
\(Y\sim \operatorname{Poisson}(\beta x)\), then
\[\mathbb{P}(X\le x)=\mathbb{P}(Y\ge \alpha).\]}

\begin{tcolorbox}[enhanced jigsaw, arc=.35mm, bottomrule=.15mm, bottomtitle=1mm, breakable, colback=white, colbacktitle=quarto-callout-note-color!10!white, colframe=quarto-callout-note-color-frame, coltitle=black, left=2mm, leftrule=.75mm, opacityback=0, opacitybacktitle=0.6, rightrule=.15mm, title=\textcolor{quarto-callout-note-color}{\faInfo}\hspace{0.5em}{Solution}, titlerule=0mm, toprule=.15mm, toptitle=1mm]

For integer \(\alpha\), \(X\) can be interpreted as the waiting time
until the \(\alpha\)-th event in a Poisson process of rate \(\beta\).
The event \(\{X\le x\}\) means that by time \(x\), at least \(\alpha\)
events have occurred. If \(Y\) is the number of events by time \(x\),
then \(Y\sim\operatorname{Poisson}(\beta x)\) and
\[\{X\le x\}=\{Y\ge \alpha\}.\] Thus the probabilities are equal.

\end{tcolorbox}

\subsection{Weibull distribution}\label{weibull-distribution}

This subsection introduces the Weibull distribution as a transformation
of an exponential random variable.

\textbf{Definition 67} (Weibull distribution). If
\(X\sim \operatorname{Exponential}(\beta)\) and \[Y=X^{1/\gamma},\] then
\(Y\) has a Weibull distribution with parameters \((\gamma,\beta)\) and
density
\[f_Y(y\mid\gamma,\beta)=\frac{\gamma}{\beta}y^{\gamma-1}e^{-y^\gamma/\beta},
\qquad y>0,\] where \(\gamma>0\) and \(\beta>0\).

\textbf{Example 68} (Deriving the Weibull density). Let
\(X\sim\operatorname{Exponential}(1/\beta)\) in the scale-style notation
with density \(f_X(x)=\frac{1}{\beta}e^{-x/\beta}\) for \(x>0\), and let
\(Y=X^{1/\gamma}\). Derive the density of \(Y\).

\begin{tcolorbox}[enhanced jigsaw, arc=.35mm, bottomrule=.15mm, bottomtitle=1mm, breakable, colback=white, colbacktitle=quarto-callout-note-color!10!white, colframe=quarto-callout-note-color-frame, coltitle=black, left=2mm, leftrule=.75mm, opacityback=0, opacitybacktitle=0.6, rightrule=.15mm, title=\textcolor{quarto-callout-note-color}{\faInfo}\hspace{0.5em}{Solution}, titlerule=0mm, toprule=.15mm, toptitle=1mm]

The transformation is \(X=Y^\gamma\). Then
\[\frac{dx}{dy}=\gamma y^{\gamma-1}.\] Thus
\[f_Y(y)=f_X(y^\gamma)\left|\frac{dx}{dy}\right|
=\frac1\beta e^{-y^\gamma/\beta}\gamma y^{\gamma-1}
=\frac{\gamma}{\beta}y^{\gamma-1}e^{-y^\gamma/\beta},\] for \(y>0\).

\end{tcolorbox}

\subsection{Chi-squared distribution}\label{chi-squared-distribution}

This subsection introduces the chi-squared distribution as a gamma
special case and as a sum of squared standard normals.

\textbf{Definition 69} (Chi-squared distribution). A chi-squared random
variable with \(k\) degrees of freedom is a gamma special case:
\[\chi_k^2\sim \operatorname{Gamma}\left(\alpha=\frac{k}{2},\theta=2\right).\]
Equivalently, if \(Z_1,\ldots,Z_k\) are independent standard normal
random variables, then \[Z_1^2+\cdots+Z_k^2\sim \chi_k^2.\]

\textbf{Example 70} (Sum of squared standard normals). If
\(Z_1,Z_2,Z_3\) are independent \(\operatorname{Normal}(0,1)\) random
variables, identify the distribution of \[Q=Z_1^2+Z_2^2+Z_3^2.\]

\begin{tcolorbox}[enhanced jigsaw, arc=.35mm, bottomrule=.15mm, bottomtitle=1mm, breakable, colback=white, colbacktitle=quarto-callout-note-color!10!white, colframe=quarto-callout-note-color-frame, coltitle=black, left=2mm, leftrule=.75mm, opacityback=0, opacitybacktitle=0.6, rightrule=.15mm, title=\textcolor{quarto-callout-note-color}{\faInfo}\hspace{0.5em}{Solution}, titlerule=0mm, toprule=.15mm, toptitle=1mm]

By the definition of a chi-squared distribution, \[Q\sim\chi_3^2.\]
Equivalently, \[Q\sim\operatorname{Gamma}\left(\frac32,2\right).\]

\end{tcolorbox}

\subsection{Beta distribution}\label{beta-distribution}

This subsection introduces the beta distribution, a flexible
distribution on the interval \((0,1)\).

\textbf{Definition 71} (Beta distribution). A random variable \(X\) has
a beta distribution with parameters \(\alpha>0\) and \(\beta>0\),
written \[X\sim\operatorname{Beta}(\alpha,\beta),\] if
\[f_X(x;\alpha,\beta)=\frac{\Gamma(\alpha+\beta)}{\Gamma(\alpha)\Gamma(\beta)}x^{\alpha-1}(1-x)^{\beta-1},
\qquad 0<x<1.\] Equivalently,
\[f_X(x;\alpha,\beta)=\frac{1}{B(\alpha,\beta)}x^{\alpha-1}(1-x)^{\beta-1},\]
where
\[B(\alpha,\beta)=\frac{\Gamma(\alpha)\Gamma(\beta)}{\Gamma(\alpha+\beta)}.\]

The beta distribution is often used as a model for proportions and
percentages. It is also a conjugate prior for the binomial model.

\textbf{Example 72} (Uniform as a special beta). Show that
\(\operatorname{Beta}(1,1)\) is the uniform distribution on \((0,1)\).

\begin{tcolorbox}[enhanced jigsaw, arc=.35mm, bottomrule=.15mm, bottomtitle=1mm, breakable, colback=white, colbacktitle=quarto-callout-note-color!10!white, colframe=quarto-callout-note-color-frame, coltitle=black, left=2mm, leftrule=.75mm, opacityback=0, opacitybacktitle=0.6, rightrule=.15mm, title=\textcolor{quarto-callout-note-color}{\faInfo}\hspace{0.5em}{Solution}, titlerule=0mm, toprule=.15mm, toptitle=1mm]

If \(\alpha=\beta=1\), then
\[f_X(x)=\frac{\Gamma(2)}{\Gamma(1)\Gamma(1)}x^0(1-x)^0.\] Since
\(\Gamma(2)=1! =1\) and \(\Gamma(1)=1\), \[f_X(x)=1,
\qquad 0<x<1.\] This is the \(\operatorname{Uniform}(0,1)\) density.

\end{tcolorbox}

\subsection{Logistic distribution}\label{logistic-distribution}

This subsection revisits the logistic distribution as a flexible
continuous distribution and as a model related to logistic regression.

\textbf{Definition 73} (Logistic distribution). A logistic distribution
can be described by the CDF \[F(x)=\frac{1}{1+e^{-\alpha-\beta x}},\]
with corresponding density
\[f(x)=\frac{\beta e^{-\alpha-\beta x}}{(1+e^{-\alpha-\beta x})^2}
=\frac{\beta e^{\alpha+\beta x}}{(1+e^{\alpha+\beta x})^2}.\] Another
location-scale form is
\[F(x)=\frac12+\frac12\tanh\left(\frac{x-\mu}{2s}\right).\]

The logistic distribution is central in logistic regression for
categorical dependent variables. It has also been used in rating
systems, such as chess rating models.

\textbf{Example 74} (Logit link). Suppose \[\mu=\frac{1}{1+e^{-\eta}}.\]
Solve for \(\eta\) in terms of \(\mu\).

\begin{tcolorbox}[enhanced jigsaw, arc=.35mm, bottomrule=.15mm, bottomtitle=1mm, breakable, colback=white, colbacktitle=quarto-callout-note-color!10!white, colframe=quarto-callout-note-color-frame, coltitle=black, left=2mm, leftrule=.75mm, opacityback=0, opacitybacktitle=0.6, rightrule=.15mm, title=\textcolor{quarto-callout-note-color}{\faInfo}\hspace{0.5em}{Solution}, titlerule=0mm, toprule=.15mm, toptitle=1mm]

Start with \[\mu=\frac{1}{1+e^{-\eta}}.\] Then
\[\frac{1}{\mu}=1+e^{-\eta},
\qquad
\frac{1-\mu}{\mu}=e^{-\eta}.\] Taking logs gives
\[-\eta=\log\left(\frac{1-\mu}{\mu}\right),\] so
\[\eta=\log\left(\frac{\mu}{1-\mu}\right).\] This is the logit
transformation.

\end{tcolorbox}

\subsection{Lognormal distribution}\label{lognormal-distribution}

This subsection introduces a distribution for positive, right-skewed
quantities.

\textbf{Definition 75} (Lognormal distribution). A positive random
variable \(X\) has a lognormal distribution if
\[\log X\sim \operatorname{Normal}(\mu,\sigma^2).\] Its mean and
variance are \[\mathbb{E}[X]=e^{\mu+\sigma^2/2},\]
\[\operatorname{Var}(X)=e^{2\mu+\sigma^2}\left(e^{\sigma^2}-1\right).\]

The lognormal distribution is popular for modeling right-skewed
variables such as incomes, file sizes, internet comment lengths, and
biological measurements.

\textbf{Example 76} (Median of a lognormal). If \(X\) is lognormal and
\(\log X\sim\operatorname{Normal}(\mu,\sigma^2)\), find the median of
\(X\).

\begin{tcolorbox}[enhanced jigsaw, arc=.35mm, bottomrule=.15mm, bottomtitle=1mm, breakable, colback=white, colbacktitle=quarto-callout-note-color!10!white, colframe=quarto-callout-note-color-frame, coltitle=black, left=2mm, leftrule=.75mm, opacityback=0, opacitybacktitle=0.6, rightrule=.15mm, title=\textcolor{quarto-callout-note-color}{\faInfo}\hspace{0.5em}{Solution}, titlerule=0mm, toprule=.15mm, toptitle=1mm]

The median \(m\) satisfies \[\mathbb{P}(X\le m)=\frac12.\] Since
\(\log\) is increasing,
\[\mathbb{P}(X\le m)=\mathbb{P}(\log X\le \log m).\] For a normal random
variable \(Y\sim\operatorname{Normal}(\mu,\sigma^2)\), the median is
\(\mu\), so we need \[\log m=\mu.\] Therefore \[m=e^\mu.\]

\end{tcolorbox}

\section{Summary Tables for Common
Distributions}\label{summary-tables-for-common-distributions}

This section organizes the common distribution families by their
stories, formulas, means, and variances.

\subsection{Distribution stories}\label{distribution-stories}

This subsection summarizes how the most common distributions arise in
applications.

{\def\LTcaptype{none} % do not increment counter
\begin{longtable}[]{@{}ll@{}}
\toprule\noalign{}
\textbf{Name} & \textbf{Useful story} \\
\midrule\noalign{}
\endhead
\bottomrule\noalign{}
\endlastfoot
\(\operatorname{Bernoulli}(p)\) & Toss a coin with probability \(p\) of
heads. \(X=\) number of heads in one toss. \\
\(\operatorname{Binomial}(n,p)\) & Toss a coin with probability \(p\) of
heads \(n\) times. \(X=\) number of heads in \(n\) tosses. It is the sum
of \(n\) independent Bernoulli variables. \\
\(\operatorname{Geometric}(p)\) & Toss a coin until the first heads.
\(X=\) number of tosses until the first heads. \\
\(\operatorname{Poisson}(\lambda)\) & Random calls arrive with rate
\(\lambda\). \(X=\) number of calls in one time unit. \\
\(\operatorname{Exponential}(\lambda)\) & Random calls arrive with rate
\(\lambda\). \(X=\) time until the first arrival. \\
\(\operatorname{Gamma}(n,\lambda)\) & Random calls arrive with rate
\(\lambda\). \(X=\) time until the \(n\)-th arrival. \\
\(\operatorname{Uniform}(a,b)\) & Pick a random number between \(a\) and
\(b\). \\
\(\operatorname{Normal}(\mu,\sigma^2)\) & Pick an individual from a
large population. \(X=\) height or another aggregate measurement. \\
\(\operatorname{Beta}(\alpha,\beta)\) & Model a random probability or
proportion, often after observing successes and failures. \\
\end{longtable}
}

\subsection{Pmf/pdf table}\label{pmfpdf-table}

This subsection lists common pmfs and pdfs in one place for quick
reference.

{\def\LTcaptype{none} % do not increment counter
\begin{longtable}[]{@{}ll@{}}
\toprule\noalign{}
\textbf{Name} & \textbf{pmf/pdf} \\
\midrule\noalign{}
\endhead
\bottomrule\noalign{}
\endlastfoot
\(\operatorname{Bernoulli}(p)\) & \(p_X(k)=p^k(1-p)^{1-k}\),
\(k\in\{0,1\}\). \\
\(\operatorname{Binomial}(n,p)\) & \(p_X(k)=\binom nk p^k(1-p)^{n-k}\),
\(k=0,\ldots,n\). \\
\(\operatorname{Geometric}(p)\) & \(p_X(k)=(1-p)^{k-1}p\),
\(k=1,2,\ldots\). \\
\(\operatorname{Poisson}(\lambda)\) &
\(p_X(k)=\dfrac{\lambda^k e^{-\lambda}}{k!}\), \(k=0,1,2,\ldots\). \\
\(\operatorname{Exponential}(\lambda)\) &
\(f_X(x)=\lambda e^{-\lambda x}\), \(x\ge 0\). \\
\(\operatorname{Gamma}(\alpha,\theta)\) &
\(f_X(x)=\dfrac{x^{\alpha-1}e^{-x/\theta}}{\theta^\alpha\Gamma(\alpha)}\),
\(x\ge0\). \\
\(\operatorname{Uniform}(a,b)\) & \(f_X(x)=\dfrac{1}{b-a}\),
\(a\le x\le b\). \\
\(\operatorname{Normal}(\mu,\sigma^2)\) &
\(f_X(x)=\dfrac{1}{\sqrt{2\pi}\sigma}\exp\left[-\dfrac{(x-\mu)^2}{2\sigma^2}\right]\). \\
\(\operatorname{Beta}(\alpha,\beta)\) &
\(f_X(x)=\dfrac{\Gamma(\alpha+\beta)}{\Gamma(\alpha)\Gamma(\beta)}x^{\alpha-1}(1-x)^{\beta-1}\),
\(0<x<1\). \\
\end{longtable}
}

\subsection{Mean and variance table}\label{mean-and-variance-table}

This subsection lists common means and variances, which will be used
later in expectation, sampling, and inference.

{\def\LTcaptype{none} % do not increment counter
\begin{longtable}[]{@{}lll@{}}
\toprule\noalign{}
\textbf{Name} & \textbf{Mean} & \textbf{Variance} \\
\midrule\noalign{}
\endhead
\bottomrule\noalign{}
\endlastfoot
\(\operatorname{Bernoulli}(p)\) & \(p\) & \(p(1-p)\) \\
\(\operatorname{Binomial}(n,p)\) & \(np\) & \(np(1-p)\) \\
\(\operatorname{Geometric}(p)\) & \(1/p\) & \((1-p)/p^2\) \\
\(\operatorname{Poisson}(\lambda)\) & \(\lambda\) & \(\lambda\) \\
\(\operatorname{Exponential}(\lambda)\) & \(1/\lambda\) &
\(1/\lambda^2\) \\
\(\operatorname{Gamma}(n,\lambda)\), rate form & \(n/\lambda\) &
\(n/\lambda^2\) \\
\(\operatorname{Uniform}(a,b)\) & \((a+b)/2\) & \((b-a)^2/12\) \\
\(\operatorname{Normal}(\mu,\sigma^2)\) & \(\mu\) & \(\sigma^2\) \\
\(\operatorname{Beta}(\alpha,\beta)\) & \(\dfrac{\alpha}{\alpha+\beta}\)
& \(\dfrac{\alpha\beta}{(\alpha+\beta)^2(\alpha+\beta+1)}\) \\
\end{longtable}
}

\section{Exponential Families}\label{exponential-families}

This section introduces the exponential family, a unifying framework
that includes many common distributions and prepares students for
generalized linear models and sufficient statistics.

\subsection{Motivation and examples}\label{motivation-and-examples}

This subsection explains why exponential families are important and
lists common members.

The exponential family includes many distributions used in statistics:
Gaussian, Bernoulli, binomial, multinomial, Poisson, exponential, gamma,
beta, von Mises, Dirichlet, Weibull, and Wishart distributions, often
under suitable fixed-parameter assumptions. Some families are
exponential families only when certain parameters are fixed, such as
binomial with fixed number of trials or multinomial with fixed total
count. Some common distributions are not exponential families in full
generality, such as Student's \(t\), most mixture distributions, and
uniform distributions when the bounds are both unknown.

\subsection{Definition}\label{definition}

This subsection gives the mathematical form of an exponential family
distribution.

\textbf{Definition 77} (Exponential family). A family of densities or
probability mass functions is an exponential family if it can be written
in the form \[p(y;\eta)=h(y)\exp\left(\eta^T T(y)-A(\eta)\right).\]
Here:

\begin{itemize}
\item
  \(\eta\in\mathbb{R}^d\) is the natural parameter;
\item
  \(T(y)\in\mathbb{R}^d\) is the sufficient statistic;
\item
  \(h(y)\) is the base measure;
\item
  \(A(\eta)\) is the log partition function or log normalizer.
\end{itemize}

The log partition function is defined by
\[A(\eta)=\log\int h(y)\exp\left(\eta^T T(y)\right)\,dy,\] with sums
replacing integrals in the discrete case. It makes the total probability
equal to \(1\).

Another common form is
\[p(y;\eta)=\exp\left(\eta^T T(y)-A(\eta)+C(y)\right),\] where
\(C(y)=\log h(y)\).

For generalized linear models, a dispersion parameter \(\phi\) is often
included:
\[p(y;\eta,\phi)=\exp\left(\frac{\eta^T T(y)-A(\eta)}{\phi}+C(y,\phi)\right).\]

\subsection{Link and response
functions}\label{link-and-response-functions}

This subsection connects the natural parameter to the mean parameter,
which is essential in generalized linear models.

The natural parameter \(\eta\) is usually not the mean of the
distribution. We often view \(\eta\) as a function of the mean
\(\mu=\mathbb{E}(Y)\): \[\eta=g(\mu),\] where \(g\) is the link
function. The inverse function \[\mu=g^{-1}(\eta)\] is the response
function.

\textbf{Example 78} (Bernoulli as an exponential family). Let
\(Y\sim\operatorname{Bernoulli}(\mu)\). Show that it belongs to the
exponential family and find its canonical link.

\begin{tcolorbox}[enhanced jigsaw, arc=.35mm, bottomrule=.15mm, bottomtitle=1mm, breakable, colback=white, colbacktitle=quarto-callout-note-color!10!white, colframe=quarto-callout-note-color-frame, coltitle=black, left=2mm, leftrule=.75mm, opacityback=0, opacitybacktitle=0.6, rightrule=.15mm, title=\textcolor{quarto-callout-note-color}{\faInfo}\hspace{0.5em}{Solution}, titlerule=0mm, toprule=.15mm, toptitle=1mm]

The pmf is \[p(y;\mu)=\mu^y(1-\mu)^{1-y},\qquad y\in\{0,1\}.\] Rewrite
it as \[p(y;\mu)=\exp\left(y\log\mu+(1-y)\log(1-\mu)\right).\] Then
\[p(y;\mu)=\exp\left(y\log\frac{\mu}{1-\mu}+\log(1-\mu)\right).\]
Compare with \[p(y;\eta)=h(y)\exp(\eta T(y)-A(\eta)).\] We identify
\[T(y)=y,
\qquad h(y)=1,
\qquad \eta=\log\frac{\mu}{1-\mu}.\] Solving for \(\mu\) gives
\[\mu=\frac{1}{1+e^{-\eta}}.\] Also, \[A(\eta)=\log(1+e^\eta).\] The
canonical link is the logit link: \[g(\mu)=\log\frac{\mu}{1-\mu}.\]

\end{tcolorbox}

\textbf{Example 79} (Categorical distribution as an exponential family).
Let \(Y\) take values \(1,\ldots,K\) with probabilities
\(\phi_1,\ldots,\phi_K\). Express the categorical distribution in
exponential family form using category \(K\) as the baseline.

\begin{tcolorbox}[enhanced jigsaw, arc=.35mm, bottomrule=.15mm, bottomtitle=1mm, breakable, colback=white, colbacktitle=quarto-callout-note-color!10!white, colframe=quarto-callout-note-color-frame, coltitle=black, left=2mm, leftrule=.75mm, opacityback=0, opacitybacktitle=0.6, rightrule=.15mm, title=\textcolor{quarto-callout-note-color}{\faInfo}\hspace{0.5em}{Solution}, titlerule=0mm, toprule=.15mm, toptitle=1mm]

The pmf is \[p(y;\phi)=\prod_{i=1}^K\phi_i^{\mathbf{1}(y=i)}.\] Use
category \(K\) as the baseline and define, for \(i=1,\ldots,K-1\),
\[\eta_i=\log\frac{\phi_i}{\phi_K},
\qquad T_i(y)=\mathbf{1}(y=i).\] Then
\[p(y;\phi)=\exp\left(\sum_{i=1}^{K-1}\mathbf{1}(y=i)\log\frac{\phi_i}{\phi_K}+\log\phi_K\right).\]
Since \[\phi_K=\frac{1}{1+\sum_{i=1}^{K-1}e^{\eta_i}},\] we have
\[A(\eta)=\log\left(1+\sum_{i=1}^{K-1}e^{\eta_i}\right).\] Thus the
categorical distribution is an exponential family.

\end{tcolorbox}

\textbf{Example 80} (Binomial as an exponential family). For fixed
\(n\), show that \(X\sim\operatorname{Binomial}(n,\phi)\) is an
exponential family.

\begin{tcolorbox}[enhanced jigsaw, arc=.35mm, bottomrule=.15mm, bottomtitle=1mm, breakable, colback=white, colbacktitle=quarto-callout-note-color!10!white, colframe=quarto-callout-note-color-frame, coltitle=black, left=2mm, leftrule=.75mm, opacityback=0, opacitybacktitle=0.6, rightrule=.15mm, title=\textcolor{quarto-callout-note-color}{\faInfo}\hspace{0.5em}{Solution}, titlerule=0mm, toprule=.15mm, toptitle=1mm]

The pmf is \[p(x;n,\phi)=\binom nx\phi^x(1-\phi)^{n-x}.\] Rewrite:
\[p(x;n,\phi)=\exp\left(x\log\frac{\phi}{1-\phi}+n\log(1-\phi)+\log\binom nx\right).\]
Set \[\eta=\log\frac{\phi}{1-\phi},
\qquad T(x)=x.\] Since \(\phi=e^\eta/(1+e^\eta)\),
\[-n\log(1-\phi)=n\log(1+e^\eta).\] Thus \[A(\eta)=n\log(1+e^\eta),
\qquad C(x)=\log\binom nx.\] Therefore the binomial distribution with
fixed \(n\) is an exponential family.

\end{tcolorbox}

\textbf{Example 81} (Normal distribution as an exponential family). Show
how \(Y\sim\operatorname{Normal}(\mu,\sigma^2)\) can be written as an
exponential family when both \(\mu\) and \(\sigma^2\) are treated as
parameters.

\begin{tcolorbox}[enhanced jigsaw, arc=.35mm, bottomrule=.15mm, bottomtitle=1mm, breakable, colback=white, colbacktitle=quarto-callout-note-color!10!white, colframe=quarto-callout-note-color-frame, coltitle=black, left=2mm, leftrule=.75mm, opacityback=0, opacitybacktitle=0.6, rightrule=.15mm, title=\textcolor{quarto-callout-note-color}{\faInfo}\hspace{0.5em}{Solution}, titlerule=0mm, toprule=.15mm, toptitle=1mm]

The density is
\[p(y;\mu,\sigma^2)=\frac{1}{\sqrt{2\pi}\sigma}\exp\left(-\frac{(y-\mu)^2}{2\sigma^2}\right).\]
Expand the square: \[-\frac{(y-\mu)^2}{2\sigma^2}
=-\frac{y^2}{2\sigma^2}+\frac{\mu}{\sigma^2}y-\frac{\mu^2}{2\sigma^2}.\]
Thus
\[p(y;\mu,\sigma^2)=\frac{1}{\sqrt{2\pi}\sigma}\exp\left(\frac{\mu}{\sigma^2}y-\frac{1}{2\sigma^2}y^2-\frac{\mu^2}{2\sigma^2}\right).\]
This is exponential family form with sufficient statistic
\[T(y)=\begin{pmatrix}y\\y^2\end{pmatrix},\] and natural parameters
\[\eta_1=\frac{\mu}{\sigma^2},
\qquad
\eta_2=-\frac{1}{2\sigma^2}.\]

\end{tcolorbox}

\textbf{Example 82} (Normal with known variance). Assume
\(Y\sim\operatorname{Normal}(\mu,\sigma^2)\) with known \(\sigma^2\).
Write it as a one-parameter exponential family.

\begin{tcolorbox}[enhanced jigsaw, arc=.35mm, bottomrule=.15mm, bottomtitle=1mm, breakable, colback=white, colbacktitle=quarto-callout-note-color!10!white, colframe=quarto-callout-note-color-frame, coltitle=black, left=2mm, leftrule=.75mm, opacityback=0, opacitybacktitle=0.6, rightrule=.15mm, title=\textcolor{quarto-callout-note-color}{\faInfo}\hspace{0.5em}{Solution}, titlerule=0mm, toprule=.15mm, toptitle=1mm]

When \(\sigma^2\) is fixed, only \(\mu\) varies. The density can be
written as
\[p(y;\mu)=\frac{1}{\sqrt{2\pi}\sigma}\exp\left(-\frac{y^2}{2\sigma^2}+\frac{\mu}{\sigma^2}y-\frac{\mu^2}{2\sigma^2}\right).\]
Set \[\eta=\frac{\mu}{\sigma^2},
\qquad T(y)=y,
\qquad A(\eta)=\frac{\sigma^2\eta^2}{2}.\] Then
\[p(y;\eta)=h(y)\exp(\eta y-A(\eta)),\] where
\[h(y)=\frac{1}{\sqrt{2\pi}\sigma}\exp\left(-\frac{y^2}{2\sigma^2}\right).\]
In GLM dispersion notation, the canonical link for a normal mean model
is the identity link.

\end{tcolorbox}

\textbf{Example 83} (Poisson as an exponential family). Show that
\(Y\sim\operatorname{Poisson}(\lambda)\) is an exponential family.

\begin{tcolorbox}[enhanced jigsaw, arc=.35mm, bottomrule=.15mm, bottomtitle=1mm, breakable, colback=white, colbacktitle=quarto-callout-note-color!10!white, colframe=quarto-callout-note-color-frame, coltitle=black, left=2mm, leftrule=.75mm, opacityback=0, opacitybacktitle=0.6, rightrule=.15mm, title=\textcolor{quarto-callout-note-color}{\faInfo}\hspace{0.5em}{Solution}, titlerule=0mm, toprule=.15mm, toptitle=1mm]

The pmf is \[p(y\mid\lambda)=\frac{\lambda^y e^{-\lambda}}{y!}
=\frac1{y!}\exp(y\log\lambda-\lambda).\] Set \[\eta=\log\lambda,
\qquad T(y)=y,
\qquad h(y)=\frac1{y!}.\] Because \(\lambda=e^\eta\),
\[A(\eta)=e^\eta.\] Thus \[p(y;\eta)=h(y)\exp(\eta y-A(\eta)).\] Also,
\[\mathbb{E}[Y]=\lambda,
\qquad \operatorname{Var}(Y)=\lambda.\]

\end{tcolorbox}

\textbf{Example 84} (Exponential distribution as an exponential family).
Show that the exponential density
\(p(y;\lambda)=\lambda e^{-\lambda y}\) for \(y\ge0\) is an exponential
family.

\begin{tcolorbox}[enhanced jigsaw, arc=.35mm, bottomrule=.15mm, bottomtitle=1mm, breakable, colback=white, colbacktitle=quarto-callout-note-color!10!white, colframe=quarto-callout-note-color-frame, coltitle=black, left=2mm, leftrule=.75mm, opacityback=0, opacitybacktitle=0.6, rightrule=.15mm, title=\textcolor{quarto-callout-note-color}{\faInfo}\hspace{0.5em}{Solution}, titlerule=0mm, toprule=.15mm, toptitle=1mm]

Write \[p(y;\lambda)=\lambda e^{-\lambda y}\mathbf{1}(y\ge0)
=\mathbf{1}(y\ge0)\exp\left(-\lambda y+\log\lambda\right).\] One
possible exponential family representation is \[\eta=\lambda,
\qquad T(y)=-y,
\qquad h(y)=\mathbf{1}(y\ge0),\] with \[A(\eta)=-\log\eta.\] Indeed,
\[h(y)\exp(\eta T(y)-A(\eta))=\mathbf{1}(y\ge0)\exp(-\lambda y+\log\lambda)=\lambda e^{-\lambda y}\mathbf{1}(y\ge0).\]

\end{tcolorbox}

\textbf{Example 85} (Laplace distribution and exponential family form).
For the Laplace density
\[f(x\mid\mu,b)=\frac{1}{2b}\exp\left(-\frac{|x-\mu|}{b}\right),\]
explain why it is an exponential family when \(\mu\) is known and \(b\)
varies.

\begin{tcolorbox}[enhanced jigsaw, arc=.35mm, bottomrule=.15mm, bottomtitle=1mm, breakable, colback=white, colbacktitle=quarto-callout-note-color!10!white, colframe=quarto-callout-note-color-frame, coltitle=black, left=2mm, leftrule=.75mm, opacityback=0, opacitybacktitle=0.6, rightrule=.15mm, title=\textcolor{quarto-callout-note-color}{\faInfo}\hspace{0.5em}{Solution}, titlerule=0mm, toprule=.15mm, toptitle=1mm]

If \(\mu\) is known, write
\[f(x\mid\mu,b)=\frac1{2b}\exp\left(-\frac{1}{b}|x-\mu|\right).\] Set
\[\eta=-\frac1b,
\qquad T(x)=|x-\mu|.\] Then the density has the exponential form with a
normalizing term depending on \(b\) (or equivalently \(\eta\)). If both
\(\mu\) and \(b\) vary freely, the absolute value centered at an unknown
\(\mu\) prevents a regular finite-dimensional exponential family
representation in the same simple canonical form.

\end{tcolorbox}

\textbf{Example 86} (Beta distribution as an exponential family). Show
the exponential-family structure of the beta distribution.

\begin{tcolorbox}[enhanced jigsaw, arc=.35mm, bottomrule=.15mm, bottomtitle=1mm, breakable, colback=white, colbacktitle=quarto-callout-note-color!10!white, colframe=quarto-callout-note-color-frame, coltitle=black, left=2mm, leftrule=.75mm, opacityback=0, opacitybacktitle=0.6, rightrule=.15mm, title=\textcolor{quarto-callout-note-color}{\faInfo}\hspace{0.5em}{Solution}, titlerule=0mm, toprule=.15mm, toptitle=1mm]

The beta density is
\[f(x;\alpha,\beta)=\frac{\Gamma(\alpha+\beta)}{\Gamma(\alpha)\Gamma(\beta)}x^{\alpha-1}(1-x)^{\beta-1},
\qquad 0<x<1.\] Taking logs,
\[\log f(x;\alpha,\beta)=(\alpha-1)\log x+(\beta-1)\log(1-x)-\log B(\alpha,\beta).\]
Thus \[T(x)=\begin{pmatrix}\log x\\ \log(1-x)\end{pmatrix},
\qquad
\eta=\begin{pmatrix}\alpha-1\\ \beta-1\end{pmatrix},\] and
\[A(\eta)=\log B(\alpha,\beta).\] This gives an exponential family
representation on \((0,1)\).

\end{tcolorbox}

\textbf{Example 87} (Gamma distribution as an exponential family). Show
the exponential-family structure of the gamma distribution with shape
\(k\) and scale \(\theta\).

\begin{tcolorbox}[enhanced jigsaw, arc=.35mm, bottomrule=.15mm, bottomtitle=1mm, breakable, colback=white, colbacktitle=quarto-callout-note-color!10!white, colframe=quarto-callout-note-color-frame, coltitle=black, left=2mm, leftrule=.75mm, opacityback=0, opacitybacktitle=0.6, rightrule=.15mm, title=\textcolor{quarto-callout-note-color}{\faInfo}\hspace{0.5em}{Solution}, titlerule=0mm, toprule=.15mm, toptitle=1mm]

The gamma density is
\[f(x;k,\theta)=\frac{x^{k-1}e^{-x/\theta}}{\theta^k\Gamma(k)},
\qquad x\ge0.\] Taking logs,
\[\log f(x;k,\theta)=(k-1)\log x-\frac{x}{\theta}-k\log\theta-\log\Gamma(k).\]
This can be written as an exponential family with sufficient statistics
\(\log x\) and \(x\) when both \(k\) and \(\theta\) vary, or as a
one-parameter exponential family when one parameter is fixed.

\end{tcolorbox}

\subsection{Moments of exponential
families}\label{moments-of-exponential-families}

This subsection explains why the log partition function is so useful:
its derivatives give moments of the sufficient statistic.

\textbf{Theorem 88} (Moments from the log partition function). \emph{For
an exponential family \[p(y;\eta)=h(y)\exp(\eta^T T(y)-A(\eta)),\] we
have \[\nabla_\eta A(\eta)=\mathbb{E}_\eta[T(Y)]\] and
\[\nabla_\eta^2 A(\eta)=\operatorname{Cov}_\eta(T(Y)).\] In the
one-dimensional case, \[A'(\eta)=\mathbb{E}[T(Y)],
\qquad
A''(\eta)=\operatorname{Var}(T(Y)).\]}

\begin{tcolorbox}[enhanced jigsaw, arc=.35mm, bottomrule=.15mm, bottomtitle=1mm, breakable, colback=white, colbacktitle=quarto-callout-note-color!10!white, colframe=quarto-callout-note-color-frame, coltitle=black, left=2mm, leftrule=.75mm, opacityback=0, opacitybacktitle=0.6, rightrule=.15mm, title=\textcolor{quarto-callout-note-color}{\faInfo}\hspace{0.5em}{Solution}, titlerule=0mm, toprule=.15mm, toptitle=1mm]

In one dimension, \[A(\eta)=\log\int h(y)e^{\eta T(y)}\,dy.\]
Differentiate:
\[A'(\eta)=\frac{\int T(y)h(y)e^{\eta T(y)}\,dy}{\int h(y)e^{\eta T(y)}\,dy}.\]
Since \[p(y;\eta)=h(y)e^{\eta T(y)-A(\eta)},\] the derivative becomes
\[A'(\eta)=\int T(y)p(y;\eta)\,dy=\mathbb{E}[T(Y)].\] Differentiating
again gives
\[A''(\eta)=\mathbb{E}[T(Y)^2]-(\mathbb{E}[T(Y)])^2=\operatorname{Var}(T(Y)).\]
The multidimensional version replaces derivatives by gradients and
Hessians.

\end{tcolorbox}

\textbf{Example 89} (Bernoulli moments from \(A(\eta)\)). For Bernoulli
in canonical form, \[A(\eta)=\log(1+e^\eta).\] Use derivatives of \(A\)
to find \(\mathbb{E}[Y]\) and \(\operatorname{Var}(Y)\).

\begin{tcolorbox}[enhanced jigsaw, arc=.35mm, bottomrule=.15mm, bottomtitle=1mm, breakable, colback=white, colbacktitle=quarto-callout-note-color!10!white, colframe=quarto-callout-note-color-frame, coltitle=black, left=2mm, leftrule=.75mm, opacityback=0, opacitybacktitle=0.6, rightrule=.15mm, title=\textcolor{quarto-callout-note-color}{\faInfo}\hspace{0.5em}{Solution}, titlerule=0mm, toprule=.15mm, toptitle=1mm]

First, \[A'(\eta)=\frac{e^\eta}{1+e^\eta}=\frac{1}{1+e^{-\eta}}=\mu.\]
Therefore \[\mathbb{E}[Y]=\mu.\] Second,
\[A''(\eta)=\frac{e^\eta}{(1+e^\eta)^2}=\mu(1-\mu).\] Thus
\[\operatorname{Var}(Y)=\mu(1-\mu).\]

\end{tcolorbox}

\section{Location and Scale Families}\label{location-and-scale-families}

This section introduces another unifying idea: many distributions are
obtained by shifting and rescaling a standard distribution.

\subsection{Definition}\label{definition-1}

This subsection defines the location-scale family through a
transformation of a standard random variable.

\textbf{Definition 90} (Location-scale family). Let \(X\) be a random
variable with CDF \(F\) and pdf \(f\). For constants
\(\mu\in\mathbb{R}\) and \(\sigma>0\), define \[Y=\mu+\sigma X.\] Then
\(Y\) belongs to the location-scale family generated by \(X\). Its pdf
is \[f_Y(y)=\frac1\sigma f\left(\frac{y-\mu}{\sigma}\right).\]

\begin{tcolorbox}[enhanced jigsaw, arc=.35mm, bottomrule=.15mm, bottomtitle=1mm, breakable, colback=white, colbacktitle=quarto-callout-note-color!10!white, colframe=quarto-callout-note-color-frame, coltitle=black, left=2mm, leftrule=.75mm, opacityback=0, opacitybacktitle=0.6, rightrule=.15mm, title=\textcolor{quarto-callout-note-color}{\faInfo}\hspace{0.5em}{Solution}, titlerule=0mm, toprule=.15mm, toptitle=1mm]

The formula follows from the change of variables
\[x=\frac{y-\mu}{\sigma},
\qquad
\frac{dx}{dy}=\frac1\sigma.\] Therefore
\[f_Y(y)=f_X\left(\frac{y-\mu}{\sigma}\right)\left|\frac1\sigma\right|.\]

\end{tcolorbox}

\subsection{Examples of location-scale
families}\label{examples-of-location-scale-families}

This subsection lists familiar distributions that fit the location-scale
framework.

Examples include:

\begin{itemize}
\item
  normal distribution;
\item
  Cauchy distribution;
\item
  continuous uniform distribution;
\item
  discrete uniform distribution in certain shifted/scaled forms;
\item
  logistic distribution;
\item
  Laplace distribution;
\item
  Student's \(t\) distribution.
\end{itemize}

\textbf{Example 91} (Normal location-scale transformation). If
\(Z\sim\operatorname{Normal}(0,1)\) and \(Y=\mu+\sigma Z\), find the
distribution of \(Y\).

\begin{tcolorbox}[enhanced jigsaw, arc=.35mm, bottomrule=.15mm, bottomtitle=1mm, breakable, colback=white, colbacktitle=quarto-callout-note-color!10!white, colframe=quarto-callout-note-color-frame, coltitle=black, left=2mm, leftrule=.75mm, opacityback=0, opacitybacktitle=0.6, rightrule=.15mm, title=\textcolor{quarto-callout-note-color}{\faInfo}\hspace{0.5em}{Solution}, titlerule=0mm, toprule=.15mm, toptitle=1mm]

The standard normal density is \[f_Z(z)=\frac1{\sqrt{2\pi}}e^{-z^2/2}.\]
Using the location-scale formula,
\[f_Y(y)=\frac1\sigma \frac1{\sqrt{2\pi}}\exp\left[-\frac12\left(\frac{y-\mu}{\sigma}\right)^2\right]
=\frac{1}{\sqrt{2\pi}\sigma}\exp\left[-\frac{(y-\mu)^2}{2\sigma^2}\right].\]
Thus \[Y\sim\operatorname{Normal}(\mu,\sigma^2).\]

\end{tcolorbox}

\textbf{Practice Problem 92} (Uniform location-scale transformation).
Let \(X\sim\operatorname{Uniform}(0,1)\) and \(Y=a+(b-a)X\) where
\(a<b\). Find the distribution of \(Y\).

\begin{tcolorbox}[enhanced jigsaw, arc=.35mm, bottomrule=.15mm, bottomtitle=1mm, breakable, colback=white, colbacktitle=quarto-callout-note-color!10!white, colframe=quarto-callout-note-color-frame, coltitle=black, left=2mm, leftrule=.75mm, opacityback=0, opacitybacktitle=0.6, rightrule=.15mm, title=\textcolor{quarto-callout-note-color}{\faInfo}\hspace{0.5em}{Solution}, titlerule=0mm, toprule=.15mm, toptitle=1mm]

Here \(\mu=a\) and \(\sigma=b-a\). Since \(X\) is uniform on \([0,1]\),
the transformation maps \([0,1]\) to \([a,b]\). The density is
\[f_Y(y)=\frac1{b-a},
\qquad a\le y\le b.\] Therefore \[Y\sim\operatorname{Uniform}(a,b).\]

\end{tcolorbox}

\section{Additional Practice
Problems}\label{additional-practice-problems}

This section provides extra problems with solutions to reinforce the
core skills of Section 2.

\textbf{Practice Problem 93} (Classify random variables). Classify each
random variable as discrete, continuous, or mixed.

\begin{enumerate}
\def\labelenumi{\arabic{enumi}.}
\item
  Number of emails received in one hour.
\item
  Time until the next email arrives.
\item
  A variable that equals \(0\) with probability \(0.3\) and otherwise is
  uniform on \([1,2]\).
\end{enumerate}

\begin{tcolorbox}[enhanced jigsaw, arc=.35mm, bottomrule=.15mm, bottomtitle=1mm, breakable, colback=white, colbacktitle=quarto-callout-note-color!10!white, colframe=quarto-callout-note-color-frame, coltitle=black, left=2mm, leftrule=.75mm, opacityback=0, opacitybacktitle=0.6, rightrule=.15mm, title=\textcolor{quarto-callout-note-color}{\faInfo}\hspace{0.5em}{Solution}, titlerule=0mm, toprule=.15mm, toptitle=1mm]

(a) The number of emails is a count, so it is discrete. A Poisson model
may be appropriate.

(b) Waiting time is measured on a continuum, so it is continuous. An
exponential model may be appropriate.

(c) This variable has a point mass at \(0\) and a continuous component
on \([1,2]\), so it is mixed.

\end{tcolorbox}

\textbf{Practice Problem 94} (From density to probability). Let \[f(x)=
\begin{cases}
cx, & 0\le x\le 2,\\
0, & \text{otherwise}.
\end{cases}\] Find \(c\) and compute \(\mathbb{P}(1\le X\le 2)\).

\begin{tcolorbox}[enhanced jigsaw, arc=.35mm, bottomrule=.15mm, bottomtitle=1mm, breakable, colback=white, colbacktitle=quarto-callout-note-color!10!white, colframe=quarto-callout-note-color-frame, coltitle=black, left=2mm, leftrule=.75mm, opacityback=0, opacitybacktitle=0.6, rightrule=.15mm, title=\textcolor{quarto-callout-note-color}{\faInfo}\hspace{0.5em}{Solution}, titlerule=0mm, toprule=.15mm, toptitle=1mm]

The total area must be \(1\):
\[1=\int_0^2 cx\,dx=c\left[\frac{x^2}{2}\right]_0^2=2c.\] Thus
\(c=1/2\). Then
\[\mathbb{P}(1\le X\le 2)=\int_1^2 \frac{x}{2}\,dx=\left[\frac{x^2}{4}\right]_1^2=1-\frac14=\frac34.\]

\end{tcolorbox}

\textbf{Practice Problem 95} (Choosing the right distribution). For each
story, choose a reasonable distribution.

\begin{enumerate}
\def\labelenumi{\arabic{enumi}.}
\item
  Count the number of defective items in \(20\) independent items, each
  defective with probability \(0.02\).
\item
  Count the number of arrivals in one minute when arrivals occur at
  average rate \(3\) per minute.
\item
  Time until the first arrival when arrivals occur at average rate \(3\)
  per minute.
\item
  Draw \(5\) cards from a deck without replacement and count the number
  of aces.
\end{enumerate}

\begin{tcolorbox}[enhanced jigsaw, arc=.35mm, bottomrule=.15mm, bottomtitle=1mm, breakable, colback=white, colbacktitle=quarto-callout-note-color!10!white, colframe=quarto-callout-note-color-frame, coltitle=black, left=2mm, leftrule=.75mm, opacityback=0, opacitybacktitle=0.6, rightrule=.15mm, title=\textcolor{quarto-callout-note-color}{\faInfo}\hspace{0.5em}{Solution}, titlerule=0mm, toprule=.15mm, toptitle=1mm]

(a) Binomial: \(X\sim\operatorname{Binomial}(20,0.02)\).

(b) Poisson: \(X\sim\operatorname{Poisson}(3)\).

(c) Exponential: \(X\sim\operatorname{Exponential}(3)\).

(d) Hypergeometric: \(N=52\), \(M=4\) aces, \(K=5\) draws, so
\[\mathbb{P}(X=x)=\frac{\binom4x\binom{48}{5-x}}{\binom{52}{5}}.\]

\end{tcolorbox}

\textbf{Practice Problem 96} (Exponential family calculation). Let
\(Y\sim\operatorname{Poisson}(\lambda)\) and use the natural parameter
\(\eta=\log\lambda\). Verify that \(A'(\eta)=\lambda\) and
\(A''(\eta)=\lambda\).

\begin{tcolorbox}[enhanced jigsaw, arc=.35mm, bottomrule=.15mm, bottomtitle=1mm, breakable, colback=white, colbacktitle=quarto-callout-note-color!10!white, colframe=quarto-callout-note-color-frame, coltitle=black, left=2mm, leftrule=.75mm, opacityback=0, opacitybacktitle=0.6, rightrule=.15mm, title=\textcolor{quarto-callout-note-color}{\faInfo}\hspace{0.5em}{Solution}, titlerule=0mm, toprule=.15mm, toptitle=1mm]

For the Poisson exponential family, \[A(\eta)=e^\eta.\] Therefore
\[A'(\eta)=e^\eta=\lambda,\] and \[A''(\eta)=e^\eta=\lambda.\] Since
\(T(Y)=Y\), this agrees with \[\mathbb{E}[Y]=\lambda,
\qquad \operatorname{Var}(Y)=\lambda.\]

\end{tcolorbox}

\bookmarksetup{startatroot}

\chapter{Chapter 3: Joint and Conditional
Probability}\label{chapter-3-joint-and-conditional-probability}

This chapter moves from one random variable to several random variables
studied together. The main objects are joint distributions, marginal
distributions, conditional distributions, independence, covariance,
total probability, Bayes theorem, and conditional independence.

\textbf{Topics.} Expected value review; variance and covariance; joint,
marginal, and conditional distributions; independence; total
probability; Bayes theorem; conditional independence.

\section{Overview}\label{overview}

This section moves from one random variable to several random variables
studied together.

In Section 2, we learned how a single random variable is described by a
CDF, a pmf, or a pdf. In applications, however, data usually come in
groups: height and weight, treatment and outcome, first die and second
die, test result and disease status, or multiple measurements from the
same experiment. The correct mathematical language for this is the
\emph{joint distribution}. Once we have a joint distribution, we can
recover marginal distributions, define conditional distributions, test
independence, and update probabilities using Bayes theorem.

\begin{tcolorbox}[enhanced jigsaw, arc=.35mm, bottomrule=.15mm, bottomtitle=1mm, breakable, colback=white, colbacktitle=quarto-callout-tip-color!10!white, colframe=quarto-callout-tip-color-frame, coltitle=black, left=2mm, leftrule=.75mm, opacityback=0, opacitybacktitle=0.6, rightrule=.15mm, title=\textcolor{quarto-callout-tip-color}{\faLightbulb}\hspace{0.5em}{Main message}, titlerule=0mm, toprule=.15mm, toptitle=1mm]

A joint distribution describes the complete probabilistic relationship
among variables. Marginal distributions describe individual variables.
Conditional distributions describe what remains after partial
information is known.

\end{tcolorbox}

\section{Expected Values: Review}\label{expected-values-review}

This section reviews expected value and variance because they are the
basic numerical summaries used throughout joint and conditional
probability.

\subsection{Expected value}\label{expected-value}

This subsection recalls the definition and interpretation of expected
value for discrete and continuous random variables.

\textbf{Definition (Expected value).}

Let \(X\) be a random variable. If \(X\) is discrete with pmf
\(p_X(k)\), then {[} \mathbb{E}{[}X{]}=\sum\emph{\{\text{all }k\} k
p\_X(k). {]} If \(X\) is continuous with pdf \(f_X(x)\), then {[}
\mathbb{E}{[}X{]}=\int}\{-\infty\}\^{}\{\infty\} x f\_X(x),d x. {]}

Expected value is a generalization of the idea of an average.
Operationally, if we measure \(X\) in many independent trials and obtain
\(X_1,\ldots,X_n\), then the sample average {[}
\overline X\_n=\frac{1}{n}(X\_1+\cdots+X\_n) {]} should approach
\(\mathbb{E}[X]\) as \(n\) becomes large. This is formalized later by
the Law of Large Numbers.

\begin{tcolorbox}[enhanced jigsaw, arc=.35mm, bottomrule=.15mm, bottomtitle=1mm, breakable, colback=white, colbacktitle=quarto-callout-important-color!10!white, colframe=quarto-callout-important-color-frame, coltitle=black, left=2mm, leftrule=.75mm, opacityback=0, opacitybacktitle=0.6, rightrule=.15mm, title=\textcolor{quarto-callout-important-color}{\faExclamation}\hspace{0.5em}{Proposition: Linearity for one random variable}, titlerule=0mm, toprule=.15mm, toptitle=1mm]

For constants \(a,b\in\mathbb{R}\), {[}
\mathbb{E}{[}aX+b{]}=a\mathbb{E}{[}X{]}+b. {]}

\end{tcolorbox}

\textbf{Example (Bernoulli expected value).}

Let \(X\sim\operatorname{Bernoulli}(p)\), so {[}
\mathbb{P}(X=0)=1-p,\qquad \mathbb{P}(X=1)=p. {]} Find
\(\mathbb{E}[X]\).

\begin{tcolorbox}[enhanced jigsaw, arc=.35mm, bottomrule=.15mm, bottomtitle=1mm, breakable, colback=white, colbacktitle=quarto-callout-note-color!10!white, colframe=quarto-callout-note-color-frame, coltitle=black, left=2mm, leftrule=.75mm, opacityback=0, opacitybacktitle=0.6, rightrule=.15mm, title=\textcolor{quarto-callout-note-color}{\faInfo}\hspace{0.5em}{Solution}, titlerule=0mm, toprule=.15mm, toptitle=1mm]

Using the discrete expectation formula, {[} \mathbb{E}{[}X{]}=\sum\_k k
p\_X(k)=0(1-p)+1(p)=p. {]} Thus the expected value of a Bernoulli random
variable is its success probability: {[} \boxed{\mathbb{E}[X]=p.} {]}

\end{tcolorbox}

\textbf{Example (Expected outcome of a fair die).}

Let \(X\) be the outcome of rolling a fair six-sided die. Find
\(\mathbb{E}[X]\).

\begin{tcolorbox}[enhanced jigsaw, arc=.35mm, bottomrule=.15mm, bottomtitle=1mm, breakable, colback=white, colbacktitle=quarto-callout-note-color!10!white, colframe=quarto-callout-note-color-frame, coltitle=black, left=2mm, leftrule=.75mm, opacityback=0, opacitybacktitle=0.6, rightrule=.15mm, title=\textcolor{quarto-callout-note-color}{\faInfo}\hspace{0.5em}{Solution}, titlerule=0mm, toprule=.15mm, toptitle=1mm]

The pmf is \(p_X(k)=1/6\) for \(k=1,2,3,4,5,6\). Therefore {[}
\mathbb{E}{[}X{]}=\sum\_\{k=1\}\^{}\{6\} k\frac{1}{6}
=\frac{1+2+3+4+5+6}{6} =\frac{21}{6}=\frac{7}{2}. {]} So the long-run
average die value is {[} \boxed{\mathbb{E}[X]=3.5.} {]}

\end{tcolorbox}

\subsection{Variance and standard
deviation}\label{variance-and-standard-deviation}

This subsection reviews how variance and standard deviation measure the
spread of a random variable around its mean.

\textbf{Definition (Variance and standard deviation).}

The variance of a random variable \(X\) is {[}
\operatorname{Var}(X)=\mathbb{E}\left[(X-\mathbb{E}[X])\^{}2\right{]}.
{]} The standard deviation is {[}
\operatorname{STD}(X)=\sqrt{\operatorname{Var}(X)}. {]}

A useful computational formula is {[}
\operatorname{Var}(X)=\mathbb{E}{[}X\^{}2{]}-\left(\mathbb{E}{[}X{]}\right)\^{}2.
{]} For constants \(a,b\in\mathbb{R}\), {[}
\operatorname{Var}(aX+b)=a\^{}2\operatorname{Var}(X). {]} Adding a
constant changes location but not spread; multiplying by \(a\) scales
spread by \(a^2\) in variance and by \(|a|\) in standard deviation.

\textbf{Example (Variance of a Bernoulli random variable).}

Let \(X\sim\operatorname{Bernoulli}(p)\). Find
\(\operatorname{Var}(X)\).

\begin{tcolorbox}[enhanced jigsaw, arc=.35mm, bottomrule=.15mm, bottomtitle=1mm, breakable, colback=white, colbacktitle=quarto-callout-note-color!10!white, colframe=quarto-callout-note-color-frame, coltitle=black, left=2mm, leftrule=.75mm, opacityback=0, opacitybacktitle=0.6, rightrule=.15mm, title=\textcolor{quarto-callout-note-color}{\faInfo}\hspace{0.5em}{Solution}, titlerule=0mm, toprule=.15mm, toptitle=1mm]

Since \(X\) only takes values \(0\) and \(1\), we have \(X^2=X\). Thus
{[} \mathbb{E}{[}X\^{}2{]}=\mathbb{E}{[}X{]}=p. {]} Using the
computational formula, {[}
\operatorname{Var}(X)=\mathbb{E}{[}X\^{}2{]}-(\mathbb{E}{[}X{]})\textsuperscript{2=p-p}2=p(1-p).
{]} Therefore {[} \boxed{\operatorname{Var}(X)=p(1-p).} {]}

\end{tcolorbox}

\textbf{Practice Problem (Review: expectation and variance).}

Suppose \(X\) takes values \(-1,0,2\) with probabilities
\(1/4,1/2,1/4\), respectively. Find \(\mathbb{E}[X]\) and
\(\operatorname{Var}(X)\).

\begin{tcolorbox}[enhanced jigsaw, arc=.35mm, bottomrule=.15mm, bottomtitle=1mm, breakable, colback=white, colbacktitle=quarto-callout-note-color!10!white, colframe=quarto-callout-note-color-frame, coltitle=black, left=2mm, leftrule=.75mm, opacityback=0, opacitybacktitle=0.6, rightrule=.15mm, title=\textcolor{quarto-callout-note-color}{\faInfo}\hspace{0.5em}{Solution}, titlerule=0mm, toprule=.15mm, toptitle=1mm]

First compute the mean: {[}
\mathbb{E}{[}X{]}=(-1)\frac14+0\frac12+2\frac14=-\frac14+\frac12=\frac14.
{]} Next compute the second moment: {[}
\mathbb{E}{[}X\^{}2{]}=(-1)\textsuperscript{2\frac14+0}2\frac12+2\^{}2\frac14=\frac14+1=\frac54.
{]} Therefore {[}
\operatorname{Var}(X)=\mathbb{E}{[}X\^{}2{]}-(\mathbb{E}{[}X{]})\textsuperscript{2=\frac54-\left(\frac14\right)}2
=\frac{20}{16}-\frac{1}{16}=\frac{19}{16}. {]} So {[}
\boxed{\mathbb{E}[X]=\frac14,\qquad \operatorname{Var}(X)=\frac{19}{16}.}
{]}

\end{tcolorbox}

\section{Joint Distributions}\label{joint-distributions}

This section introduces the joint distribution, which is the main object
for studying two or more random variables together.

\subsection{Joint CDF, joint pmf, and joint
pdf}\label{joint-cdf-joint-pmf-and-joint-pdf}

This subsection defines the joint distribution for two random variables
in both discrete and continuous cases.

\textbf{Definition (Joint CDF).}

For two random variables \(X\) and \(Y\), the joint cumulative
distribution function is {[} F\_\{X,Y\}(x,y)=\mathbb{P}(X\le x,,Y\le y).
{]}

If \(X\) and \(Y\) are both discrete, then their joint pmf is {[}
p\_\{X,Y\}(x,y)=\mathbb{P}(X=x,Y=y), {]} for all
\(x\in\mathbb{R}ange(X)\) and \(y\in\mathbb{R}ange(Y)\). It must satisfy
{[} p\_\{X,Y\}(x,y)\ge 0, \qquad
\sum\emph{\{x\}\sum}\{y\}p\_\{X,Y\}(x,y)=1. {]}

If \(X\) and \(Y\) are absolutely continuous, then their joint pdf
\(f_{X,Y}(x,y)\) satisfies {[} f\_\{X,Y\}(x,y)\ge 0, \qquad
\int\emph{\{-\infty\}\textsuperscript{\{\infty\}\int\_\{-\infty\}}\{\infty\}f}\{X,Y\}(x,y),d
x,d y=1. {]} When the joint CDF is differentiable, {[}
f\_\{X,Y\}(x,y)=\frac{\partial^2}{\partial x\partial y}F\_\{X,Y\}(x,y).
{]} For a region \(R\) in the \(xy\)-plane, {[}
\mathbb{P}((X,Y)\in R)=\iint\emph{R f}\{X,Y\}(x,y),d x,d y. {]}

\subsection{Marginal distributions}\label{marginal-distributions}

This subsection explains how to recover the distribution of one variable
from the joint distribution.

The marginal distribution of \(X\) is obtained by summing or integrating
out \(Y\). Similarly, the marginal distribution of \(Y\) is obtained by
summing or integrating out \(X\).

For discrete random variables, {[} p\_X(x)=\sum\emph{y p}\{X,Y\}(x,y),
\qquad p\_Y(y)=\sum\emph{x p}\{X,Y\}(x,y). {]} For continuous random
variables, {[} f\_X(x)=\int\emph{\{-\infty\}\^{}\{\infty\}
f}\{X,Y\}(x,y),d y, \qquad f\_Y(y)=\int\emph{\{-\infty\}\^{}\{\infty\}
f}\{X,Y\}(x,y),d x. {]}

\begin{tcolorbox}[enhanced jigsaw, arc=.35mm, bottomrule=.15mm, bottomtitle=1mm, breakable, colback=white, colbacktitle=quarto-callout-warning-color!10!white, colframe=quarto-callout-warning-color-frame, coltitle=black, left=2mm, leftrule=.75mm, opacityback=0, opacitybacktitle=0.6, rightrule=.15mm, title=\textcolor{quarto-callout-warning-color}{\faExclamationTriangle}\hspace{0.5em}{Common mistake}, titlerule=0mm, toprule=.15mm, toptitle=1mm]

When finding \(f_X(x)\), integrate with respect to \(y\). When finding
\(f_Y(y)\), integrate with respect to \(x\).

\end{tcolorbox}

\subsection{Discrete joint distribution example: two
dice}\label{discrete-joint-distribution-example-two-dice}

This subsection works out a full discrete joint distribution using two
dice.

\textbf{Example (Difference and maximum of two dice).}

Roll two fair dice. Let {[}
X=\text{absolute difference of the two values}, \qquad
Y=\text{maximum of the two values}. {]} Find the ranges, marginal pmfs,
and joint pmf table of \((X,Y)\).

\begin{tcolorbox}[enhanced jigsaw, arc=.35mm, bottomrule=.15mm, bottomtitle=1mm, breakable, colback=white, colbacktitle=quarto-callout-note-color!10!white, colframe=quarto-callout-note-color-frame, coltitle=black, left=2mm, leftrule=.75mm, opacityback=0, opacitybacktitle=0.6, rightrule=.15mm, title=\textcolor{quarto-callout-note-color}{\faInfo}\hspace{0.5em}{Solution}, titlerule=0mm, toprule=.15mm, toptitle=1mm]

The sample space has \(36\) equally likely outcomes \((i,j)\) with
\(i,j\in\{1,2,3,4,5,6\}\). The range of \(X\) is {[}
\mathbb{R}ange(X)=\{0,1,2,3,4,5\}, {]} and the range of \(Y\) is {[}
\mathbb{R}ange(Y)=\{1,2,3,4,5,6\}. {]}

For \(X=0\), the two dice match. There are \(6\) outcomes. For
\(X=d\ge 1\), the two dice differ by \(d\); there are \(2(6-d)\) ordered
outcomes. Hence {[}

\begin{array}{c|cccccc}
x&0&1&2&3&4&5\\ \hline
p_X(x)&6/36&10/36&8/36&6/36&4/36&2/36
\end{array}

{]}

For \(Y=y\), at least one die equals \(y\) and both dice are at most
\(y\). There are {[} y\textsuperscript{2-(y-1)}2=2y-1 {]} outcomes. Thus
{[}

\begin{array}{c|cccccc}
y&1&2&3&4&5&6\\ \hline
p_Y(y)&1/36&3/36&5/36&7/36&9/36&11/36
\end{array}

{]}

The joint pmf table is {[}

\begin{array}{c|cccccc|c}
& X=0&X=1&X=2&X=3&X=4&X=5&p_Y(y)\\ \hline
Y=1&1/36&0&0&0&0&0&1/36\\
Y=2&1/36&2/36&0&0&0&0&3/36\\
Y=3&1/36&2/36&2/36&0&0&0&5/36\\
Y=4&1/36&2/36&2/36&2/36&0&0&7/36\\
Y=5&1/36&2/36&2/36&2/36&2/36&0&9/36\\
Y=6&1/36&2/36&2/36&2/36&2/36&2/36&11/36\\ \hline
p_X(x)&6/36&10/36&8/36&6/36&4/36&2/36&1
\end{array}

{]} The entries sum to \(1\), and the row and column sums agree with the
marginals.

\end{tcolorbox}

\textbf{Example (Same marginals, different joint pmf).}

Show that marginal distributions alone do not determine the joint
distribution.

\begin{tcolorbox}[enhanced jigsaw, arc=.35mm, bottomrule=.15mm, bottomtitle=1mm, breakable, colback=white, colbacktitle=quarto-callout-note-color!10!white, colframe=quarto-callout-note-color-frame, coltitle=black, left=2mm, leftrule=.75mm, opacityback=0, opacitybacktitle=0.6, rightrule=.15mm, title=\textcolor{quarto-callout-note-color}{\faInfo}\hspace{0.5em}{Solution}, titlerule=0mm, toprule=.15mm, toptitle=1mm]

Let \(X,Y\in\{0,1\}\), and suppose both marginals satisfy {[}
\mathbb{P}(X=0)=\mathbb{P}(X=1)=\frac12, \qquad
\mathbb{P}(Y=0)=\mathbb{P}(Y=1)=\frac12. {]} Consider the two joint pmf
tables {[}

\begin{array}{c|cc}
 &Y=0&Y=1\\ \hline
X=0&1/2&0\\
X=1&0&1/2
\end{array}
\qquad\text{and}\qquad
\begin{array}{c|cc}
 &Y=0&Y=1\\ \hline
X=0&1/4&1/4\\
X=1&1/4&1/4
\end{array}

. {]} Both tables have the same marginal distributions for \(X\) and
\(Y\). However, the dependence structure is different. In the first
table, \(Y=X\) always. In the second table, \(X\) and \(Y\) are
independent. Therefore marginals alone do not determine the joint
distribution.

\end{tcolorbox}

\subsection{Continuous joint distribution
example}\label{continuous-joint-distribution-example}

This subsection gives a full continuous example with normalization, a
region probability, a marginal density, and a one-dimensional
probability.

\textbf{Example (A joint pdf on the unit square).}

A joint pdf is defined by {[} f(x,y)=6xy\^{}2,
\qquad 0\textless x\textless1,\quad 0\textless y\textless1, {]} and
\(f(x,y)=0\) otherwise.

\begin{itemize}
\tightlist
\item
  Check that it is a valid joint pdf.
\item
  Calculate \(\mathbb{P}(X+Y\ge 1)\).
\item
  Calculate the marginal pdf of \(X\), \(f_X(x)\).
\item
  Calculate \(\mathbb{P}\left(\frac12<X<\frac34\right)\).
\end{itemize}

\begin{tcolorbox}[enhanced jigsaw, arc=.35mm, bottomrule=.15mm, bottomtitle=1mm, breakable, colback=white, colbacktitle=quarto-callout-note-color!10!white, colframe=quarto-callout-note-color-frame, coltitle=black, left=2mm, leftrule=.75mm, opacityback=0, opacitybacktitle=0.6, rightrule=.15mm, title=\textcolor{quarto-callout-note-color}{\faInfo}\hspace{0.5em}{Solution}, titlerule=0mm, toprule=.15mm, toptitle=1mm]

\textbf{(1) Validity.} Clearly \(f(x,y)\ge 0\) on its support. Also, {[}
\int\_0\textsuperscript{1\int\_0}1 6xy\^{}2,d x,d y =\int\_0\^{}1
3y\^{}2,d y =\left.y\textsuperscript{3\right\textbar\_0}1=1. {]} Thus
\(f\) is a valid joint pdf.

\textbf{(2) Region probability.} The event \(X+Y\ge 1\) corresponds to
\(0<x<1\) and \(1-x\le y\le 1\). Hence {[} \mathbb{P}(X+Y\ge 1)
=\int\emph{0\textsuperscript{1\int\_\{1-x\}}\{1\}6xy\^{}2,d y,d x. {]}
Compute the inner integral: {[}
\int}\{1-x\}\textsuperscript{\{1\}6xy}2,d y
=2x\left(1-(1-x)\^{}3\right). {]} Therefore {[}
\mathbb{P}(X+Y\ge 1)=\int\_0\^{}1 2x\left(1-(1-x)\^{}3\right),d
x=\frac{9}{10}. {]}

\textbf{(3) Marginal of \(X\).} {[} f\_X(x)=\int\_0\^{}1 6xy\^{}2,d
y=6x\cdot \frac13=2x, \qquad 0\textless x\textless1. {]}

\textbf{(4) Probability for \(X\).} {[}
\mathbb{P}\left(\frac12\textless X\textless{}\frac34\right)=\int\_\{1/2\}\^{}\{3/4\}2x,d
x =\left.x\textsuperscript{2\right\textbar\_\{1/2\}}\{3/4\}
=\frac{9}{16}-\frac{4}{16}=\frac{5}{16}. {]} Thus {[}
\boxed{\mathbb{P}(X+Y\ge1)=\frac{9}{10},\qquad f_X(x)=2x,
\qquad \mathbb{P}\left(\frac12<X<\frac34\right)=\frac{5}{16}.} {]}

\end{tcolorbox}

\textbf{Practice Problem (Marginal density practice).}

Suppose {[}
f(x,y)=c(x+y),\qquad 0\textless x\textless1,\quad 0\textless y\textless1,
{]} and \(0\) otherwise. Find \(c\), \(f_X(x)\), and
\(\mathbb{P}(X\le 1/2)\).

\begin{tcolorbox}[enhanced jigsaw, arc=.35mm, bottomrule=.15mm, bottomtitle=1mm, breakable, colback=white, colbacktitle=quarto-callout-note-color!10!white, colframe=quarto-callout-note-color-frame, coltitle=black, left=2mm, leftrule=.75mm, opacityback=0, opacitybacktitle=0.6, rightrule=.15mm, title=\textcolor{quarto-callout-note-color}{\faInfo}\hspace{0.5em}{Solution}, titlerule=0mm, toprule=.15mm, toptitle=1mm]

Normalize: {[} 1=\int\_0\textsuperscript{1\int\_0}1 c(x+y),d y,d x
=c\int\_0\^{}1\left(x+\frac12\right),d x
=c\left(\frac12+\frac12\right)=c. {]} So \(c=1\). Then {[}
f\_X(x)=\int\_0\^{}1(x+y),d y=x+\frac12, \qquad 0\textless x\textless1.
{]} Finally, {[}
\mathbb{P}(X\le 1/2)=\int\_0\^{}\{1/2\}\left(x+\frac12\right),d x
=\left.\left(\frac{x^2}{2}+\frac{x}{2}\right)\right\textbar\_0\^{}\{1/2\}
=\frac18+\frac14=\frac38. {]}

\end{tcolorbox}

\section{Multivariate Normal
Distribution}\label{multivariate-normal-distribution}

This section introduces the most important multivariate continuous
distribution: the multivariate normal distribution.

\subsection{Definition and geometry}\label{definition-and-geometry}

This subsection defines the multivariate normal density and explains the
roles of the mean vector and covariance matrix.

\textbf{Definition (Multivariate normal distribution).}

Let {[} \vec X=

\begin{pmatrix}X_1\\ \vdots\\ X_d\end{pmatrix}

. {]} We say {[} \vec X\sim\mathbb{N}ormal(\vec\mu,\Sigma) {]} if
\(\vec\mu\in\mathbb{R}^d\) and \(\Sigma\) is a \(d\times d\) symmetric
positive definite matrix, and the joint pdf is {[} f\_\{\vec X\}(\vec x)
=\frac{1}{\sqrt{(2\pi)^d|\Sigma|}}
\exp\left[-\frac12(\vec x-\vec\mu)^T\Sigma^{-1}(\vec x-\vec\mu)\right].
{]}

The mean vector is {[} \mathbb{E}{[}\vec X{]}=\vec\mu, {]} and the
covariance matrix is {[} \operatorname{Cov}(\vec X)=\Sigma. {]} The
matrix \(\Sigma\) controls the spread and orientation of the density. A
nearly diagonal covariance matrix produces contours aligned with the
coordinate axes; a covariance matrix with nonzero off-diagonal entries
produces tilted elliptical contours.

\begin{tcolorbox}[enhanced jigsaw, arc=.35mm, bottomrule=.15mm, bottomtitle=1mm, breakable, colback=white, colbacktitle=quarto-callout-note-color!10!white, colframe=quarto-callout-note-color-frame, coltitle=black, left=2mm, leftrule=.75mm, opacityback=0, opacitybacktitle=0.6, rightrule=.15mm, title=\textcolor{quarto-callout-note-color}{\faInfo}\hspace{0.5em}{Visual idea}, titlerule=0mm, toprule=.15mm, toptitle=1mm]

Multivariate normal contours are ellipses. Diagonal covariance gives
axis-aligned contours, while nonzero covariance tilts the contours.

\end{tcolorbox}

\subsection{Covariance and
correlation}\label{covariance-and-correlation}

This subsection defines covariance and correlation as second-order
summaries of the relationship between two random variables.

\textbf{Definition (Covariance and correlation).}

For two random variables \(X\) and \(Y\), the covariance is {[}
\operatorname{Cov}(X,Y)=\mathbb{E}{[}(X-\mathbb{E}{[}X{]})(Y-\mathbb{E}{[}Y{]}){]}
=\mathbb{E}{[}XY{]}-\mathbb{E}{[}X{]}\mathbb{E}{[}Y{]}. {]} The Pearson
correlation coefficient is {[}
\operatorname{Corr}(X,Y)=\frac{\operatorname{Cov}(X,Y)}{\operatorname{STD}(X)\operatorname{STD}(Y)}.
{]} It satisfies {[} -1\le \operatorname{Corr}(X,Y)\le 1. {]}

Covariance has units depending on \(X\) and \(Y\), while correlation is
dimensionless. Positive correlation means that large values of \(X\)
tend to occur with large values of \(Y\); negative correlation means
that large values of \(X\) tend to occur with small values of \(Y\).

\textbf{Example (Covariance from a small joint table).}

Let \(X,Y\in\{0,1\}\) with joint pmf {[}

\begin{array}{c|cc}
&Y=0&Y=1\\ \hline
X=0&0.40&0.10\\
X=1&0.10&0.40
\end{array}

. {]} Find \(\operatorname{Cov}(X,Y)\).

\begin{tcolorbox}[enhanced jigsaw, arc=.35mm, bottomrule=.15mm, bottomtitle=1mm, breakable, colback=white, colbacktitle=quarto-callout-note-color!10!white, colframe=quarto-callout-note-color-frame, coltitle=black, left=2mm, leftrule=.75mm, opacityback=0, opacitybacktitle=0.6, rightrule=.15mm, title=\textcolor{quarto-callout-note-color}{\faInfo}\hspace{0.5em}{Solution}, titlerule=0mm, toprule=.15mm, toptitle=1mm]

The marginals are {[} \mathbb{P}(X=1)=0.10+0.40=0.50, \qquad
\mathbb{P}(Y=1)=0.10+0.40=0.50. {]} Thus
\(\mathbb{E}[X]=\mathbb{E}[Y]=0.5\). Also, since \(XY=1\) only when
\(X=1,Y=1\), {[} \mathbb{E}{[}XY{]}=\mathbb{P}(X=1,Y=1)=0.40. {]}
Therefore {[}
\operatorname{Cov}(X,Y)=\mathbb{E}{[}XY{]}-\mathbb{E}{[}X{]}\mathbb{E}{[}Y{]}=0.40-(0.50)(0.50)=0.15.
{]} The covariance is positive.

\end{tcolorbox}

\section{Conditional Probability and Conditional
Distributions}\label{conditional-probability-and-conditional-distributions}

This section formalizes the idea of updating probabilities after some
information has been observed.

\subsection{Conditional probability for
events}\label{conditional-probability-for-events}

This subsection recalls the event-level definition from Section 1.

\textbf{Definition (Conditional probability).}

If \(\mathbb{P}(B)>0\), then the probability of \(A\) given \(B\) is {[}
\mathbb{P}(A\mid B)=\frac{\mathbb{P}(A\cap B)}{\mathbb{P}(B)}. {]} When
\(B\) is fixed, \(\mathbb{P}(\cdot\mid B)\) is another probability
measure.

\subsection{Conditional distributions}\label{conditional-distributions}

This subsection extends conditional probability from events to random
variables.

For two random variables \(X\) and \(Y\), the conditional pmf/pdf of
\(Y\) given \(X=x\) is {[}
p\_\{Y\mid X\}(y\mid x)=\frac{p_{X,Y}(x,y)}{p_X(x)}, {]} provided
\(p_X(x)>0\). In the continuous case the same formula is written using
densities: {[} f\_\{Y\mid X\}(y\mid x)=\frac{f_{X,Y}(x,y)}{f_X(x)}, {]}
provided \(f_X(x)>0\).

\textbf{Example (Triangle uniform distribution).}

Let \((X,Y)\) be uniformly distributed on the triangular region {[}
D=\{(x,y):x\ge0,~y\ge0,~x+y\le1\}. {]} The joint pdf is {[} f(x,y)=

\begin{cases}
2, & (x,y)\in D,\\
0, & \text{otherwise}.
\end{cases}

{]} Find \(f_X(x)\) and \(f_{Y\mid X}(y\mid x)\).

\begin{tcolorbox}[enhanced jigsaw, arc=.35mm, bottomrule=.15mm, bottomtitle=1mm, breakable, colback=white, colbacktitle=quarto-callout-note-color!10!white, colframe=quarto-callout-note-color-frame, coltitle=black, left=2mm, leftrule=.75mm, opacityback=0, opacitybacktitle=0.6, rightrule=.15mm, title=\textcolor{quarto-callout-note-color}{\faInfo}\hspace{0.5em}{Solution}, titlerule=0mm, toprule=.15mm, toptitle=1mm]

For a fixed \(x\in[0,1]\), the variable \(y\) ranges from \(0\) to
\(1-x\). Hence {[} f\_X(x)=\int\emph{0\^{}\{1-x\}2,d y=2(1-x),
\qquad 0\le x\le 1. {]} Therefore {[} f}\{Y\mid X\}(y\mid x)
=\frac{f_{X,Y}(x,y)}{f_X(x)} =\frac{2}{2(1-x)}=\frac{1}{1-x}, {]} for
\(0\le y\le 1-x\). Thus, conditional on \(X=x\), \(Y\) is uniform on
\([0,1-x]\): {[}
\boxed{f_{Y\mid X}(y\mid x)=\frac{1}{1-x},\qquad 0\le y\le 1-x.} {]}

\end{tcolorbox}

\textbf{Example (Beta-Bernoulli joint distribution).}

The Beta-Bernoulli model is a basic Bayesian model. Suppose {[}
Y\sim \operatorname{Beta}(\alpha,\beta) {]} and, conditional on \(Y=y\),
{[} X\mid Y=y\sim\operatorname{Bernoulli}(y). {]} Find the joint
distribution \(p(x,y)\) and identify the posterior distribution of \(Y\)
given \(X=x\).

\begin{tcolorbox}[enhanced jigsaw, arc=.35mm, bottomrule=.15mm, bottomtitle=1mm, breakable, colback=white, colbacktitle=quarto-callout-note-color!10!white, colframe=quarto-callout-note-color-frame, coltitle=black, left=2mm, leftrule=.75mm, opacityback=0, opacitybacktitle=0.6, rightrule=.15mm, title=\textcolor{quarto-callout-note-color}{\faInfo}\hspace{0.5em}{Solution}, titlerule=0mm, toprule=.15mm, toptitle=1mm]

The conditional pmf of \(X\) given \(Y=y\) is {[}
p(x\mid y)=y\textsuperscript{x(1-y)}\{1-x\}, \qquad x\in\{0,1\}. {]} The
density of \(Y\) is {[}
f\_Y(y)=\frac{1}{B(\alpha,\beta)}y\textsuperscript{\{\alpha-1\}(1-y)}\{\beta-1\},
\qquad 0\textless y\textless1. {]} Thus the joint distribution is {[}
p(x,y)=p(x\mid y)f\_Y(y)
=\frac{1}{B(\alpha,\beta)}y\textsuperscript{\{\alpha+x-1\}(1-y)}\{\beta+1-x-1\}.
{]} As a function of \(y\) with \(x\) fixed, this has the kernel of a
beta density: {[}
y\textsuperscript{\{(\alpha+x)-1\}(1-y)}\{(\beta+1-x)-1\}. {]} Therefore
{[} \boxed{Y\mid X=x\sim \operatorname{Beta}(\alpha+x,\,\beta+1-x).} {]}
This is why the beta distribution is called a conjugate prior for the
Bernoulli model.

\end{tcolorbox}

\textbf{Practice Problem (Conditional density practice).}

Let {[} f(x,y)=8xy, \qquad 0\textless y\textless x\textless1, {]} and
\(0\) otherwise. Find \(f_X(x)\) and \(f_{Y\mid X}(y\mid x)\).

\begin{tcolorbox}[enhanced jigsaw, arc=.35mm, bottomrule=.15mm, bottomtitle=1mm, breakable, colback=white, colbacktitle=quarto-callout-note-color!10!white, colframe=quarto-callout-note-color-frame, coltitle=black, left=2mm, leftrule=.75mm, opacityback=0, opacitybacktitle=0.6, rightrule=.15mm, title=\textcolor{quarto-callout-note-color}{\faInfo}\hspace{0.5em}{Solution}, titlerule=0mm, toprule=.15mm, toptitle=1mm]

For a fixed \(x\in(0,1)\), \(y\) ranges from \(0\) to \(x\). Thus {[}
f\_X(x)=\int\emph{0\^{}x 8xy,d y=8x\cdot\frac{x^2}{2}=4x\^{}3,
\qquad 0\textless x\textless1. {]} Then {[}
f}\{Y\mid X\}(y\mid x)=\frac{8xy}{4x^3}=\frac{2y}{x^2},
\qquad 0\textless y\textless x. {]} The conditional density integrates
to one: {[} \int\_0\^{}x \frac{2y}{x^2},d y=1. {]}

\end{tcolorbox}

\section{Independence}\label{independence-1}

This section explains independence for events and for random variables,
and shows how independence simplifies joint distributions and
expectations.

\subsection{Independence for events}\label{independence-for-events}

This subsection recalls the event-level definition of independence.

\textbf{Definition (Independent events).}

Events \(A\) and \(B\) are independent if {[}
\mathbb{P}(A\cap B)=\mathbb{P}(A)\mathbb{P}(B). {]} If
\(\mathbb{P}(A)>0\) and \(\mathbb{P}(B)>0\), this is equivalent to {[}
\mathbb{P}(A\mid B)=\mathbb{P}(A) \qquad\text{and}\qquad
\mathbb{P}(B\mid A)=\mathbb{P}(B). {]}

The interpretation is that knowing one event occurred does not change
the probability of the other event.

\subsection{Independent random
variables}\label{independent-random-variables}

This subsection defines independence using joint and marginal
distributions.

\textbf{Definition (Independent random variables).}

Discrete random variables \(X\) and \(Y\) are independent if {[}
p\_\{X,Y\}(x,y)=p\_X(x)p\_Y(y) {]} for all possible \(x\) and \(y\).
Continuous random variables \(X\) and \(Y\) are independent if {[}
f\_\{X,Y\}(x,y)=f\_X(x)f\_Y(y) {]} for all \(x,y\).

Similarly, \(X_1,\ldots,X_n\) are mutually independent if their joint
pmf/pdf factors into the product of their marginal pmfs/pdfs.

\textbf{Example (Independent coin tosses).}

Toss a biased coin twice. Let {[} X=

\begin{cases}1,&\text{first toss is Heads},\\0,&\text{first toss is Tails},\end{cases}
\qquad

Y=

\begin{cases}1,&\text{second toss is Heads},\\0,&\text{second toss is Tails}.
\end{cases}

{]} Suppose \(\mathbb{P}(\text{Heads})=p\). Show that \(X\) and \(Y\)
are independent.

\begin{tcolorbox}[enhanced jigsaw, arc=.35mm, bottomrule=.15mm, bottomtitle=1mm, breakable, colback=white, colbacktitle=quarto-callout-note-color!10!white, colframe=quarto-callout-note-color-frame, coltitle=black, left=2mm, leftrule=.75mm, opacityback=0, opacitybacktitle=0.6, rightrule=.15mm, title=\textcolor{quarto-callout-note-color}{\faInfo}\hspace{0.5em}{Solution}, titlerule=0mm, toprule=.15mm, toptitle=1mm]

Because the two tosses are independent, the probability of any pair of
outcomes factors. For example, {[}
\mathbb{P}(X=1,Y=0)=\mathbb{P}(\text{Heads then Tails})=p(1-p). {]}
Also, {[} \mathbb{P}(X=1)=p, \qquad \mathbb{P}(Y=0)=1-p. {]} Hence {[}
\mathbb{P}(X=1,Y=0)=\mathbb{P}(X=1)\mathbb{P}(Y=0). {]} The same check
holds for the other three pairs \((0,0),(0,1),(1,1)\). Therefore \(X\)
and \(Y\) are independent.

\end{tcolorbox}

\textbf{Example (Checking dependence in the two-dice example).}

In the two-dice example, \(X\) is the absolute difference and \(Y\) is
the maximum. Are \(X\) and \(Y\) independent?

\begin{tcolorbox}[enhanced jigsaw, arc=.35mm, bottomrule=.15mm, bottomtitle=1mm, breakable, colback=white, colbacktitle=quarto-callout-note-color!10!white, colframe=quarto-callout-note-color-frame, coltitle=black, left=2mm, leftrule=.75mm, opacityback=0, opacitybacktitle=0.6, rightrule=.15mm, title=\textcolor{quarto-callout-note-color}{\faInfo}\hspace{0.5em}{Solution}, titlerule=0mm, toprule=.15mm, toptitle=1mm]

Use one entry of the joint table. We have {[}
p\_\{X,Y\}(0,1)=\frac{1}{36}. {]} The marginals are {[}
p\_X(0)=\frac{6}{36}=\frac16, \qquad p\_Y(1)=\frac{1}{36}. {]} If \(X\)
and \(Y\) were independent, then {[}
p\_\{X,Y\}(0,1)=p\_X(0)p\_Y(1)=\frac16\cdot\frac1{36}=\frac1{216}. {]}
But \(1/36\ne 1/216\). Therefore {[}
\boxed{X\text{ and }Y\text{ are not independent}.} {]}

\end{tcolorbox}

\subsection{Checking independence by
factorization}\label{checking-independence-by-factorization}

This subsection gives a useful theorem for checking independence without
explicitly computing both marginals first.

\begin{tcolorbox}[enhanced jigsaw, arc=.35mm, bottomrule=.15mm, bottomtitle=1mm, breakable, colback=white, colbacktitle=quarto-callout-important-color!10!white, colframe=quarto-callout-important-color-frame, coltitle=black, left=2mm, leftrule=.75mm, opacityback=0, opacitybacktitle=0.6, rightrule=.15mm, title=\textcolor{quarto-callout-important-color}{\faExclamation}\hspace{0.5em}{Theorem: Factorization criterion}, titlerule=0mm, toprule=.15mm, toptitle=1mm]

Let \((X,Y)\) have joint pmf or pdf \(f(x,y)\). Then \(X\) and \(Y\) are
independent if and only if there exist nonnegative functions \(g(x)\)
and \(h(y)\) such that {[} f(x,y)=g(x)h(y) {]} for every \((x,y)\) in
the support, and the support itself factors as a product of a set in
\(x\) and a set in \(y\).

\end{tcolorbox}

\textbf{Example (Factorized joint density).}

Suppose a joint density has the form {[}
f(x,y)=\frac{1}{384}x\textsuperscript{2y}4
e\textsuperscript{\{-x/2\}e}\{-y\},
\qquad x\textgreater0,\textbackslash{} y\textgreater0. {]} Determine
whether \(X\) and \(Y\) are independent.

\begin{tcolorbox}[enhanced jigsaw, arc=.35mm, bottomrule=.15mm, bottomtitle=1mm, breakable, colback=white, colbacktitle=quarto-callout-note-color!10!white, colframe=quarto-callout-note-color-frame, coltitle=black, left=2mm, leftrule=.75mm, opacityback=0, opacitybacktitle=0.6, rightrule=.15mm, title=\textcolor{quarto-callout-note-color}{\faInfo}\hspace{0.5em}{Solution}, titlerule=0mm, toprule=.15mm, toptitle=1mm]

We can write {[}
f(x,y)=\left(x\textsuperscript{2e}\{-x/2\}\right)\left(\frac{1}{384}y\textsuperscript{4e}\{-y\}\right)
=g(x)h(y), {]} where {[}
g(x)=x\textsuperscript{2e}\{-x/2\},\qquad h(y)=\frac{1}{384}y\textsuperscript{4e}\{-y\}.
{]} The support is also a product region: \(x>0\) and \(y>0\). Hence the
joint density factors into a function of \(x\) times a function of
\(y\), so \(X\) and \(Y\) are independent.

\end{tcolorbox}

\begin{tcolorbox}[enhanced jigsaw, arc=.35mm, bottomrule=.15mm, bottomtitle=1mm, breakable, colback=white, colbacktitle=quarto-callout-warning-color!10!white, colframe=quarto-callout-warning-color-frame, coltitle=black, left=2mm, leftrule=.75mm, opacityback=0, opacitybacktitle=0.6, rightrule=.15mm, title=\textcolor{quarto-callout-warning-color}{\faExclamationTriangle}\hspace{0.5em}{Support matters}, titlerule=0mm, toprule=.15mm, toptitle=1mm]

A formula may look factorized, but if the support couples \(x\) and
\(y\), the variables may not be independent. For example, \(f(x,y)=2\)
on \(0<x<1\), \(0<y<1-x\) does not have a product support.

\end{tcolorbox}

\subsection{Independence and
information}\label{independence-and-information}

This subsection explains why independence is essential in
likelihood-based statistical inference.

Suppose we observe \(X_1,\ldots,X_n\in\{0,1\}\), where the observations
are independent and identically distributed as
\(\operatorname{Bernoulli}(\theta)\). Then the joint pmf factors as {[}
p(x\_1,\ldots,x\_n;\theta)=p(x\_1;\theta)p(x\_2;\theta)\cdots p(x\_n;\theta).
{]} Taking logarithms gives {[}
\log p(x\_1,\ldots,x\_n;\theta)=\log p(x\_1;\theta)+\cdots+\log p(x\_n;\theta).
{]} Thus the log-likelihood is a sum of individual log-likelihood
contributions: {[}
\ell(\theta;x\_1,\ldots,x\_n)=\ell(\theta;x\_1)+\cdots+\ell(\theta;x\_n).
{]} This is the statistical meaning of the phrase: total information
equals the sum of individual information.

\section{Consequences of
Independence}\label{consequences-of-independence}

This section collects important algebraic consequences of independence
for expectation, variance, and moment generating functions.

\subsection{Expectations of products}\label{expectations-of-products}

This subsection gives the key expectation rule for independent
variables.

\begin{tcolorbox}[enhanced jigsaw, arc=.35mm, bottomrule=.15mm, bottomtitle=1mm, breakable, colback=white, colbacktitle=quarto-callout-important-color!10!white, colframe=quarto-callout-important-color-frame, coltitle=black, left=2mm, leftrule=.75mm, opacityback=0, opacitybacktitle=0.6, rightrule=.15mm, title=\textcolor{quarto-callout-important-color}{\faExclamation}\hspace{0.5em}{Theorem: Products of functions of independent random variables}, titlerule=0mm, toprule=.15mm, toptitle=1mm]

If \(X\) and \(Y\) are independent, then for suitable functions \(g\)
and \(h\), {[}
\mathbb{E}{[}g(X)h(Y){]}=\mathbb{E}{[}g(X){]}\mathbb{E}{[}h(Y){]}. {]}
In particular, {[}
\mathbb{E}{[}XY{]}=\mathbb{E}{[}X{]}\mathbb{E}{[}Y{]}. {]}

\end{tcolorbox}

\begin{tcolorbox}[enhanced jigsaw, arc=.35mm, bottomrule=.15mm, bottomtitle=1mm, breakable, colback=white, colbacktitle=quarto-callout-note-color!10!white, colframe=quarto-callout-note-color-frame, coltitle=black, left=2mm, leftrule=.75mm, opacityback=0, opacitybacktitle=0.6, rightrule=.15mm, title=\textcolor{quarto-callout-note-color}{\faInfo}\hspace{0.5em}{Solution}, titlerule=0mm, toprule=.15mm, toptitle=1mm]

For the continuous case, {[} \mathbb{E}{[}g(X)h(Y){]}
=\int\int g(x)h(y)f\_\{X,Y\}(x,y),d x,d y. {]} Independence gives
\(f_{X,Y}(x,y)=f_X(x)f_Y(y)\), so {[} \mathbb{E}{[}g(X)h(Y){]}
=\int\int g(x)h(y)f\_X(x)f\_Y(y),d x,d y =\left(\int g(x)f\_X(x),d
x\right) \left(\int h(y)f\_Y(y),d y\right). {]} Thus {[}
\mathbb{E}{[}g(X)h(Y){]}=\mathbb{E}{[}g(X){]}\mathbb{E}{[}h(Y){]}. {]}
The discrete proof is the same with sums replacing integrals.

\end{tcolorbox}

\subsection{Moment generating functions and
sums}\label{moment-generating-functions-and-sums}

This subsection explains why moment generating functions are especially
useful for sums of independent random variables.

\begin{tcolorbox}[enhanced jigsaw, arc=.35mm, bottomrule=.15mm, bottomtitle=1mm, breakable, colback=white, colbacktitle=quarto-callout-important-color!10!white, colframe=quarto-callout-important-color-frame, coltitle=black, left=2mm, leftrule=.75mm, opacityback=0, opacitybacktitle=0.6, rightrule=.15mm, title=\textcolor{quarto-callout-important-color}{\faExclamation}\hspace{0.5em}{Theorem: MGF of a sum of independent random variables}, titlerule=0mm, toprule=.15mm, toptitle=1mm]

If \(X\) and \(Y\) are independent and \(Z=X+Y\), then {[}
M\_Z(t)=M\_X(t)M\_Y(t), {]} where \(M_X(t)=\mathbb{E}[e^{tX}]\).

\end{tcolorbox}

\begin{tcolorbox}[enhanced jigsaw, arc=.35mm, bottomrule=.15mm, bottomtitle=1mm, breakable, colback=white, colbacktitle=quarto-callout-note-color!10!white, colframe=quarto-callout-note-color-frame, coltitle=black, left=2mm, leftrule=.75mm, opacityback=0, opacitybacktitle=0.6, rightrule=.15mm, title=\textcolor{quarto-callout-note-color}{\faInfo}\hspace{0.5em}{Solution}, titlerule=0mm, toprule=.15mm, toptitle=1mm]

By definition, {[}
M\_Z(t)=\mathbb{E}{[}e\^{}\{tZ\}{]}=\mathbb{E}{[}e\^{}\{t(X+Y)\}{]}=\mathbb{E}{[}e\textsuperscript{\{tX\}e}\{tY\}{]}.
{]} Since \(X\) and \(Y\) are independent, \(e^{tX}\) and \(e^{tY}\) are
independent functions of them, so {[}
M\_Z(t)=\mathbb{E}{[}e\^{}\{tX\}{]}\mathbb{E}{[}e\^{}\{tY\}{]}=M\_X(t)M\_Y(t).
{]}

\end{tcolorbox}

\begin{tcolorbox}[enhanced jigsaw, arc=.35mm, bottomrule=.15mm, bottomtitle=1mm, breakable, colback=white, colbacktitle=quarto-callout-important-color!10!white, colframe=quarto-callout-important-color-frame, coltitle=black, left=2mm, leftrule=.75mm, opacityback=0, opacitybacktitle=0.6, rightrule=.15mm, title=\textcolor{quarto-callout-important-color}{\faExclamation}\hspace{0.5em}{Theorem: Sum of independent normal random variables}, titlerule=0mm, toprule=.15mm, toptitle=1mm]

Suppose {[} X\sim\mathbb{N}ormal(\mu\_1,\sigma\_1\^{}2), \qquad
Y\sim\mathbb{N}ormal(\mu\_2,\sigma\_2\^{}2), {]} and \(X\) and \(Y\) are
independent. Then {[}
X+Y\sim\mathbb{N}ormal(\mu\_1+\mu\_2,\sigma\_1\textsuperscript{2+\sigma\_2}2).
{]}

\end{tcolorbox}

\begin{tcolorbox}[enhanced jigsaw, arc=.35mm, bottomrule=.15mm, bottomtitle=1mm, breakable, colback=white, colbacktitle=quarto-callout-note-color!10!white, colframe=quarto-callout-note-color-frame, coltitle=black, left=2mm, leftrule=.75mm, opacityback=0, opacitybacktitle=0.6, rightrule=.15mm, title=\textcolor{quarto-callout-note-color}{\faInfo}\hspace{0.5em}{Solution}, titlerule=0mm, toprule=.15mm, toptitle=1mm]

The MGF of a normal random variable
\(X\sim\mathbb{N}ormal(\mu,\sigma^2)\) is {[}
M\_X(t)=\exp\left(\mu t+\frac{\sigma^2t^2}{2}\right). {]} Let \(Z=X+Y\).
Since \(X\) and \(Y\) are independent, {[} M\_Z(t)=M\_X(t)M\_Y(t). {]}
Therefore {[} M\_Z(t) =\exp\left(\mu\_1t+\frac{\sigma_1^2t^2}{2}\right)
\exp\left(\mu\_2t+\frac{\sigma_2^2t^2}{2}\right). {]} Combining
exponents, {[}
M\_Z(t)=\exp\left((\mu\_1+\mu\_2)t+\frac{(\sigma_1^2+\sigma_2^2)t^2}{2}\right),
{]} which is the MGF of
\(\mathbb{N}ormal(\mu_1+\mu_2,\sigma_1^2+\sigma_2^2)\).

\end{tcolorbox}

\subsection{Linearity, covariance, and variance of
sums}\label{linearity-covariance-and-variance-of-sums}

This subsection separates what is always true from what requires
independence.

For any random variables \(X\) and \(Y\) on the same sample space, {[}
\mathbb{E}{[}aX+bY{]}=a\mathbb{E}{[}X{]}+b\mathbb{E}{[}Y{]}. {]} No
independence is required for linearity of expectation.

The variance of a linear combination is {[}
\operatorname{Var}(aX+bY)=a\textsuperscript{2\operatorname{Var}(X)+b}2\operatorname{Var}(Y)+2ab\operatorname{Cov}(X,Y).
{]} If \(X\) and \(Y\) are independent, then {[}
\operatorname{Cov}(X,Y)=0, {]} and therefore {[}
\operatorname{Var}(aX+bY)=a\textsuperscript{2\operatorname{Var}(X)+b}2\operatorname{Var}(Y).
{]}

\begin{tcolorbox}[enhanced jigsaw, arc=.35mm, bottomrule=.15mm, bottomtitle=1mm, breakable, colback=white, colbacktitle=quarto-callout-warning-color!10!white, colframe=quarto-callout-warning-color-frame, coltitle=black, left=2mm, leftrule=.75mm, opacityback=0, opacitybacktitle=0.6, rightrule=.15mm, title=\textcolor{quarto-callout-warning-color}{\faExclamationTriangle}\hspace{0.5em}{Zero covariance is not the same as independence}, titlerule=0mm, toprule=.15mm, toptitle=1mm]

Independence implies zero covariance, provided the variances exist. The
converse is not true: zero covariance does not necessarily imply
independence.

\end{tcolorbox}

\textbf{Practice Problem (A zero covariance but not independent
example).}

Let \(X\) take values \(-1,0,1\) with probabilities \(1/3,1/3,1/3\), and
define \(Y=X^2\). Show that \(\operatorname{Cov}(X,Y)=0\) but \(X\) and
\(Y\) are not independent.

\begin{tcolorbox}[enhanced jigsaw, arc=.35mm, bottomrule=.15mm, bottomtitle=1mm, breakable, colback=white, colbacktitle=quarto-callout-note-color!10!white, colframe=quarto-callout-note-color-frame, coltitle=black, left=2mm, leftrule=.75mm, opacityback=0, opacitybacktitle=0.6, rightrule=.15mm, title=\textcolor{quarto-callout-note-color}{\faInfo}\hspace{0.5em}{Solution}, titlerule=0mm, toprule=.15mm, toptitle=1mm]

We compute {[} \mathbb{E}{[}X{]}=\frac{-1+0+1}{3}=0, \qquad Y=X\^{}2,
\qquad XY=X\^{}3. {]} Thus {[}
\mathbb{E}{[}XY{]}=\mathbb{E}{[}X\^{}3{]}=\frac{(-1)^3+0^3+1^3}{3}=0.
{]} Hence {[}
\operatorname{Cov}(X,Y)=\mathbb{E}{[}XY{]}-\mathbb{E}{[}X{]}\mathbb{E}{[}Y{]}=0-0\cdot \mathbb{E}{[}Y{]}=0.
{]} But \(X\) and \(Y\) are not independent. For example, {[}
\mathbb{P}(Y=0\mid X=0)=1, {]} while {[}
\mathbb{P}(Y=0)=\mathbb{P}(X=0)=\frac13. {]} Knowing \(X=0\) changes the
probability of \(Y=0\), so \(X\) and \(Y\) are dependent.

\end{tcolorbox}

\section{Law of Total Probability and Bayes Theorem for Random
Variables}\label{law-of-total-probability-and-bayes-theorem-for-random-variables}

This section extends total probability and Bayes theorem from events to
random variables.

\subsection{Law of total probability for random
variables}\label{law-of-total-probability-for-random-variables}

This subsection shows how a marginal distribution can be obtained by
averaging conditional distributions.

\begin{tcolorbox}[enhanced jigsaw, arc=.35mm, bottomrule=.15mm, bottomtitle=1mm, breakable, colback=white, colbacktitle=quarto-callout-important-color!10!white, colframe=quarto-callout-important-color-frame, coltitle=black, left=2mm, leftrule=.75mm, opacityback=0, opacitybacktitle=0.6, rightrule=.15mm, title=\textcolor{quarto-callout-important-color}{\faExclamation}\hspace{0.5em}{Theorem: Law of total probability for random variables}, titlerule=0mm, toprule=.15mm, toptitle=1mm]

If \(X\) is discrete, then {[} p\_Y(y)=\sum\emph{\{x'\}
p}\{Y\mid X\}(y\mid x')p\_X(x'). {]} If \(X\) is continuous, then {[}
p\_Y(y)=\int p\_\{Y\mid X\}(y\mid x')p\_X(x'),d x'. {]}

\end{tcolorbox}

\textbf{Example (Poisson-Binomial thinning).}

Let {[} X\sim\operatorname{Poisson}(\lambda), \qquad
Y\mid X=x\sim\operatorname{Binomial}(x,p). {]} Find the marginal
distribution of \(Y\).

\begin{tcolorbox}[enhanced jigsaw, arc=.35mm, bottomrule=.15mm, bottomtitle=1mm, breakable, colback=white, colbacktitle=quarto-callout-note-color!10!white, colframe=quarto-callout-note-color-frame, coltitle=black, left=2mm, leftrule=.75mm, opacityback=0, opacitybacktitle=0.6, rightrule=.15mm, title=\textcolor{quarto-callout-note-color}{\faInfo}\hspace{0.5em}{Solution}, titlerule=0mm, toprule=.15mm, toptitle=1mm]

By the law of total probability, {[} \mathbb{P}(Y=y)=\sum\emph{\{x\}
\mathbb{P}(Y=y\mid X=x)\mathbb{P}(X=x). {]} Since
\(Y\mid X=x\sim\operatorname{Binomial}(x,p)\), the sum starts at
\(x=y\): {[}
\mathbb{P}(Y=y)=\sum}\{x\ge y\}\binom{x}{y}p\textsuperscript{y(1-p)}\{x-y\}\frac{\lambda^x e^{-\lambda}}{x!}.
{]} Using \(\binom{x}{y}=x!/[y!(x-y)!]\), we get {[}
\mathbb{P}(Y=y)=\sum\emph{\{x\ge y\}\frac{p^y(1-p)^{x-y}\lambda^x e^{-\lambda}}{y!(x-y)!}.
{]} Let \(k=x-y\). Then {[}
\mathbb{P}(Y=y)=\frac{(\lambda p)^y e^{-\lambda}}{y!}
\sum}\{k=0\}\textsuperscript{\{\infty\}\textbackslash frac\{\left(\lambda(1-p)\right)}k\}\{k!\}.
{]} The sum is \(e^{\lambda(1-p)}\), so {[}
\mathbb{P}(Y=y)=\frac{(\lambda p)^y e^{-\lambda}e^{\lambda(1-p)}}{y!}
=\frac{(\lambda p)^y e^{-\lambda p}}{y!}. {]} Therefore {[}
\boxed{Y\sim\operatorname{Poisson}(\lambda p).} {]}

\end{tcolorbox}

\subsection{Bayes theorem for random
variables}\label{bayes-theorem-for-random-variables}

This subsection gives Bayes theorem in pmf/pdf notation.

\begin{tcolorbox}[enhanced jigsaw, arc=.35mm, bottomrule=.15mm, bottomtitle=1mm, breakable, colback=white, colbacktitle=quarto-callout-important-color!10!white, colframe=quarto-callout-important-color-frame, coltitle=black, left=2mm, leftrule=.75mm, opacityback=0, opacitybacktitle=0.6, rightrule=.15mm, title=\textcolor{quarto-callout-important-color}{\faExclamation}\hspace{0.5em}{Theorem: Bayes theorem for random variables}, titlerule=0mm, toprule=.15mm, toptitle=1mm]

For random variables \(X\) and \(Y\), {[} p\_\{X\mid Y\}(x\mid y)
=\frac{p_{X,Y}(x,y)}{p_Y(y)} =\frac{p_{Y\mid X}(y\mid x)p_X(x)}{p_Y(y)}.
{]} If \(X\) is discrete, then {[} p\_\{X\mid Y\}(x\mid y)
=\frac{p_{Y\mid X}(y\mid x)p_X(x)}{\sum_{x'}p_{Y\mid X}(y\mid x')p_X(x')}.
{]} If \(X\) is continuous, then {[} p\_\{X\mid Y\}(x\mid y)
=\frac{p_{Y\mid X}(y\mid x)p_X(x)}{\int p_{Y\mid X}(y\mid x')p_X(x')\,d x'}.
{]}

\end{tcolorbox}

\textbf{Example (Bayesian update with a beta prior).}

Suppose \(Y\sim\operatorname{Beta}(\alpha,\beta)\) and
\(X\mid Y=y\sim\operatorname{Bernoulli}(y)\). Use Bayes theorem to find
\(Y\mid X=1\).

\begin{tcolorbox}[enhanced jigsaw, arc=.35mm, bottomrule=.15mm, bottomtitle=1mm, breakable, colback=white, colbacktitle=quarto-callout-note-color!10!white, colframe=quarto-callout-note-color-frame, coltitle=black, left=2mm, leftrule=.75mm, opacityback=0, opacitybacktitle=0.6, rightrule=.15mm, title=\textcolor{quarto-callout-note-color}{\faInfo}\hspace{0.5em}{Solution}, titlerule=0mm, toprule=.15mm, toptitle=1mm]

By Bayes theorem, {[}
f\_\{Y\mid X\}(y\mid 1)\propto \mathbb{P}(X=1\mid Y=y)f\_Y(y). {]} Now
{[} \mathbb{P}(X=1\mid Y=y)=y, \qquad
f\_Y(y)\propto y\textsuperscript{\{\alpha-1\}(1-y)}\{\beta-1\}. {]}
Therefore {[}
f\_\{Y\mid X\}(y\mid 1)\propto y\textsuperscript{\{\alpha\}(1-y)}\{\beta-1\}.
{]} This is the kernel of \(\operatorname{Beta}(\alpha+1,\beta)\), so
{[} \boxed{Y\mid X=1\sim\operatorname{Beta}(\alpha+1,\beta).} {]}
Similarly, if \(X=0\), then
\(Y\mid X=0\sim\operatorname{Beta}(\alpha,\beta+1)\).

\end{tcolorbox}

\textbf{Practice Problem (Total probability practice).}

Suppose \(X\sim\operatorname{Bernoulli}(q)\) and {[}
Y\mid X=0\sim\operatorname{Bernoulli}(p\_0), \qquad
Y\mid X=1\sim\operatorname{Bernoulli}(p\_1). {]} Find
\(\mathbb{P}(Y=1)\).

\begin{tcolorbox}[enhanced jigsaw, arc=.35mm, bottomrule=.15mm, bottomtitle=1mm, breakable, colback=white, colbacktitle=quarto-callout-note-color!10!white, colframe=quarto-callout-note-color-frame, coltitle=black, left=2mm, leftrule=.75mm, opacityback=0, opacitybacktitle=0.6, rightrule=.15mm, title=\textcolor{quarto-callout-note-color}{\faInfo}\hspace{0.5em}{Solution}, titlerule=0mm, toprule=.15mm, toptitle=1mm]

Use total probability over the two possible values of \(X\): {[}
\mathbb{P}(Y=1)=\mathbb{P}(Y=1\mid X=0)\mathbb{P}(X=0)
+\mathbb{P}(Y=1\mid X=1)\mathbb{P}(X=1). {]} Since \(\mathbb{P}(X=1)=q\)
and \(\mathbb{P}(X=0)=1-q\), {[} \boxed{\mathbb{P}(Y=1)=p_0(1-q)+p_1q.}
{]}

\end{tcolorbox}

\section{Conditional Independence}\label{conditional-independence-1}

This section introduces conditional independence, a central concept in
Bayesian statistics, causal inference, and graphical models.

\subsection{Definition and equivalent
forms}\label{definition-and-equivalent-forms}

This subsection defines conditional independence and lists several
equivalent factorization properties.

\textbf{Definition (Conditional independence).}

Random variables \(X\) and \(Y\) are conditionally independent given
\(Z\) if {[}
p\_\{X,Y\mid Z\}(x,y\mid z)=p\_\{X\mid Z\}(x\mid z)p\_\{Y\mid Z\}(y\mid z)
{]} for all relevant \(x,y,z\).

\begin{tcolorbox}[enhanced jigsaw, arc=.35mm, bottomrule=.15mm, bottomtitle=1mm, breakable, colback=white, colbacktitle=quarto-callout-important-color!10!white, colframe=quarto-callout-important-color-frame, coltitle=black, left=2mm, leftrule=.75mm, opacityback=0, opacitybacktitle=0.6, rightrule=.15mm, title=\textcolor{quarto-callout-important-color}{\faExclamation}\hspace{0.5em}{Theorem: Equivalent forms}, titlerule=0mm, toprule=.15mm, toptitle=1mm]

If \(X\) and \(Y\) are conditionally independent given \(Z\), then the
following equivalent statements hold whenever the conditional densities
are defined:

\begin{itemize}
\tightlist
\item
  {[}
  p\_\{X,Y\mid Z\}(x,y\mid z)=p\_\{X\mid Z\}(x\mid z)p\_\{Y\mid Z\}(y\mid z).
  {]}
\item
  {[} p\_\{X\mid Y,Z\}(x\mid y,z)=p\_\{X\mid Z\}(x\mid z). {]}
\item
  {[} p\_\{X,Y,Z\}(x,y,z)=\frac{p_{X,Z}(x,z)p_{Y,Z}(y,z)}{p_Z(z)}. {]}
\item
  There exist functions \(g\) and \(h\) such that {[}
  p\_\{X,Y,Z\}(x,y,z)=g(x,z)h(y,z). {]}
\end{itemize}

\end{tcolorbox}

The key interpretation is: after \(Z\) is known, learning \(Y\) gives no
additional information about \(X\).

\begin{tcolorbox}[enhanced jigsaw, arc=.35mm, bottomrule=.15mm, bottomtitle=1mm, breakable, colback=white, colbacktitle=quarto-callout-warning-color!10!white, colframe=quarto-callout-warning-color-frame, coltitle=black, left=2mm, leftrule=.75mm, opacityback=0, opacitybacktitle=0.6, rightrule=.15mm, title=\textcolor{quarto-callout-warning-color}{\faExclamationTriangle}\hspace{0.5em}{Independence versus conditional independence}, titlerule=0mm, toprule=.15mm, toptitle=1mm]

Independence and conditional independence are different concepts.
Independence does not always imply conditional independence, and
conditional independence does not always imply independence.

\end{tcolorbox}

\subsection{Unit disk example}\label{unit-disk-example}

This subsection gives a full example involving a joint density on a
non-rectangular support.

\textbf{Example (Uniform distribution on the unit disk).}

Let \((X,Y)\) be uniformly distributed over the unit disk {[}
D=\{(x,y):x\textsuperscript{2+y}2\le1\}. {]} The joint pdf is {[}
f\_\{X,Y\}(x,y)=

\begin{cases}
c, & (x,y)\in D,\\
0, & \text{otherwise}.
\end{cases}

{]}

\begin{itemize}
\tightlist
\item
  Find \(c\).
\item
  Find the marginal pdfs \(f_X(x)\) and \(f_Y(y)\).
\item
  Find the conditional pdf \(f_{X\mid Y}(x\mid y)\) for \(-1\le y\le1\).
\item
  Are \(X\) and \(Y\) independent?
\end{itemize}

\begin{tcolorbox}[enhanced jigsaw, arc=.35mm, bottomrule=.15mm, bottomtitle=1mm, breakable, colback=white, colbacktitle=quarto-callout-note-color!10!white, colframe=quarto-callout-note-color-frame, coltitle=black, left=2mm, leftrule=.75mm, opacityback=0, opacitybacktitle=0.6, rightrule=.15mm, title=\textcolor{quarto-callout-note-color}{\faInfo}\hspace{0.5em}{Solution}, titlerule=0mm, toprule=.15mm, toptitle=1mm]

\textbf{(a) Find \(c\).} Since the area of the unit disk is \(\pi\), {[}
1=\iint\_D c,d x,d y=c\cdot \pi. {]} Thus {[} \boxed{c=\frac1\pi.} {]}

\textbf{(b) Find the marginals.} For fixed \(x\in[-1,1]\), \(y\) ranges
from \(-\sqrt{1-x^2}\) to \(\sqrt{1-x^2}\). Hence {[}
f\_X(x)=\int\_\{-\sqrt{1-x^2}\}\textsuperscript{\{\textbackslash sqrt\{1-x}2\}\}\frac1\pi,d
y =\frac{2\sqrt{1-x^2}}{\pi}, \qquad -1\le x\le1. {]} By symmetry, {[}
f\_Y(y)=\frac{2\sqrt{1-y^2}}{\pi}, \qquad -1\le y\le1. {]}

\textbf{(c) Find the conditional pdf.} For fixed \(y\in[-1,1]\), \(x\)
ranges from \(-\sqrt{1-y^2}\) to \(\sqrt{1-y^2}\). Therefore {[}
f\_\{X\mid Y\}(x\mid y)=\frac{f_{X,Y}(x,y)}{f_Y(y)}
=\frac{1/\pi}{2\sqrt{1-y^2}/\pi} =\frac{1}{2\sqrt{1-y^2}}, {]} for {[}
-\sqrt{1-y^2}\le x\le \sqrt{1-y^2}. {]} Thus {[} \boxed{
f_{X\mid Y}(x\mid y)=
\begin{cases}
\dfrac{1}{2\sqrt{1-y^2}}, & -\sqrt{1-y^2}\le x\le \sqrt{1-y^2},\\[6pt]
0, & \text{otherwise}.
\end{cases}} {]}

\textbf{(d) Independence.} If \(X\) and \(Y\) were independent, we would
have {[} f\_\{X,Y\}(x,y)=f\_X(x)f\_Y(y). {]} But inside the disk, {[}
f\_\{X,Y\}(x,y)=\frac1\pi, {]} whereas {[}
f\_X(x)f\_Y(y)=\frac{4\sqrt{(1-x^2)(1-y^2)}}{\pi^2}, {]} which is not
constant on the disk. Also, \(f_{X\mid Y}(x\mid y)\) depends on \(y\).
Therefore {[} \boxed{X\text{ and }Y\text{ are not independent}.} {]}

\end{tcolorbox}

\section{Summary}\label{summary-1}

This section summarizes the main formulas and conceptual points from
joint and conditional probability.

\begin{tcolorbox}[enhanced jigsaw, arc=.35mm, bottomrule=.15mm, bottomtitle=1mm, breakable, colback=white, colbacktitle=quarto-callout-tip-color!10!white, colframe=quarto-callout-tip-color-frame, coltitle=black, left=2mm, leftrule=.75mm, opacityback=0, opacitybacktitle=0.6, rightrule=.15mm, title=\textcolor{quarto-callout-tip-color}{\faLightbulb}\hspace{0.5em}{Core formulas}, titlerule=0mm, toprule=.15mm, toptitle=1mm]

\begin{align*}
F_{X,Y}(x,y)&=\mathbb{P}(X\le x,Y\le y),\\
p_X(x)&=\sum_y p_{X,Y}(x,y),\\
f_X(x)&=\int f_{X,Y}(x,y)\,d y,\\
p_{Y\mid X}(y\mid x)&=\frac{p_{X,Y}(x,y)}{p_X(x)},\\
p_Y(y)&=\sum_x p_{Y\mid X}(y\mid x)p_X(x),\\
p_{X\mid Y}(x\mid y)&=\frac{p_{Y\mid X}(y\mid x)p_X(x)}{p_Y(y)},\\
\operatorname{Cov}(X,Y)&=\mathbb{E}[XY]-\mathbb{E}[X]\mathbb{E}[Y].
\end{align*}

\end{tcolorbox}

\begin{tcolorbox}[enhanced jigsaw, arc=.35mm, bottomrule=.15mm, bottomtitle=1mm, breakable, colback=white, colbacktitle=quarto-callout-tip-color!10!white, colframe=quarto-callout-tip-color-frame, coltitle=black, left=2mm, leftrule=.75mm, opacityback=0, opacitybacktitle=0.6, rightrule=.15mm, title=\textcolor{quarto-callout-tip-color}{\faLightbulb}\hspace{0.5em}{Conceptual checklist}, titlerule=0mm, toprule=.15mm, toptitle=1mm]

\begin{itemize}
\tightlist
\item
  A joint distribution contains the full relationship between variables.
\item
  A marginal distribution describes one variable after summing or
  integrating out the others.
\item
  A conditional distribution describes one variable after another
  variable has been fixed or observed.
\item
  Independence means the joint distribution factors into marginals.
\item
  Conditional independence means the conditional joint distribution
  factors after conditioning on another variable.
\item
  Bayes theorem reverses conditioning: it turns \(p_{Y\mid X}\) and
  \(p_X\) into \(p_{X\mid Y}\).
\end{itemize}

\end{tcolorbox}

\bookmarksetup{startatroot}

\chapter{Chapter 4: Expectations, Moments, and Moment Generating
Functions}\label{chapter-4-expectations-moments-and-moment-generating-functions}

This chapter develops numerical summaries of random variables and
introduces moment generating functions as a compact way to encode
moments and distributions.

\textbf{Topics.} Expected value; variance and standard deviation;
expectations of functions; moments; skewness; kurtosis; moment
generating functions; multivariate MGFs.

\section{Overview}\label{overview-1}

This section develops numerical summaries of random variables and
introduces moment generating functions as a compact way to encode
distributions.

In earlier sections, we described random variables by distributions:
CDFs, pmfs, and pdfs. In this section, we ask what numbers can summarize
a distribution. The most important summaries are expected value,
variance, and higher moments. Moment generating functions then package
all raw moments into one function, and they become especially powerful
for sums of independent random variables and normal distributions.

\begin{tcolorbox}[enhanced jigsaw, arc=.35mm, bottomrule=.15mm, bottomtitle=1mm, breakable, colback=white, colbacktitle=quarto-callout-tip-color!10!white, colframe=quarto-callout-tip-color-frame, coltitle=black, left=2mm, leftrule=.75mm, opacityback=0, opacitybacktitle=0.6, rightrule=.15mm, title=\textcolor{quarto-callout-tip-color}{\faLightbulb}\hspace{0.5em}{Main message}, titlerule=0mm, toprule=.15mm, toptitle=1mm]

Expectation is the mathematical version of long-run average. Variance
measures spread. Higher moments describe shape. The moment generating
function, when it exists near \(0\), generates all moments and can
identify a distribution.

\end{tcolorbox}

\section{Expected Value}\label{expected-value-1}

This section reviews the central idea of expectation and explains how it
connects probability theory with averages observed in data.

\subsection{Definition for discrete and continuous random
variables}\label{definition-for-discrete-and-continuous-random-variables}

This subsection gives the two most common formulas for expected value:
one for sums and one for integrals.

\textbf{Definition 1} (Expected value). Let \(X\) be a random variable.

\begin{itemize}
\item
  If \(X\) is discrete with pmf \(p_X(k)\), then
  \[\mathbb{E}[X]=\sum_{\text{all }k} k p_X(k).\]
\item
  If \(X\) is continuous with pdf \(f_X(x)\), then
  \[\mathbb{E}[X]=\int_{-\infty}^{\infty} x f_X(x)\,dx.\]
\end{itemize}

Expected value is a generalization of the concept ``average.'' It is not
necessarily the most likely value, and it does not need to be a value
that \(X\) can actually take. It is the weighted average of possible
values, where the weights are probabilities.

\textbf{Example 2} (Bernoulli random variable). Let
\(X\sim \operatorname{Bernoulli}(p)\), so
\[\mathbb{P}(X=0)=1-p,\qquad \mathbb{P}(X=1)=p.\] Find
\(\mathbb{E}[X]\).

\begin{tcolorbox}[enhanced jigsaw, arc=.35mm, bottomrule=.15mm, bottomtitle=1mm, breakable, colback=white, colbacktitle=quarto-callout-note-color!10!white, colframe=quarto-callout-note-color-frame, coltitle=black, left=2mm, leftrule=.75mm, opacityback=0, opacitybacktitle=0.6, rightrule=.15mm, title=\textcolor{quarto-callout-note-color}{\faInfo}\hspace{0.5em}{Solution}, titlerule=0mm, toprule=.15mm, toptitle=1mm]

Using the definition of expectation for a discrete random variable,
\[\mathbb{E}[X]=\sum_k k p_X(k)=0\cdot (1-p)+1\cdot p=p.\] Thus the mean
of a Bernoulli random variable is its success probability.

\end{tcolorbox}

\textbf{Example 3} (Outcome of a fair die). Let \(X\) be the outcome of
rolling a fair six-sided die. Find \(\mathbb{E}[X]\).

\begin{tcolorbox}[enhanced jigsaw, arc=.35mm, bottomrule=.15mm, bottomtitle=1mm, breakable, colback=white, colbacktitle=quarto-callout-note-color!10!white, colframe=quarto-callout-note-color-frame, coltitle=black, left=2mm, leftrule=.75mm, opacityback=0, opacitybacktitle=0.6, rightrule=.15mm, title=\textcolor{quarto-callout-note-color}{\faInfo}\hspace{0.5em}{Solution}, titlerule=0mm, toprule=.15mm, toptitle=1mm]

Here \(X\in\{1,2,3,4,5,6\}\) and each outcome has probability \(1/6\).
Therefore \[\mathbb{E}[X]=\sum_{k=1}^6 k\cdot \frac16
=\frac{1+2+3+4+5+6}{6}=\frac{21}{6}=3.5.\] Notice that \(3.5\) is not a
possible die outcome. The expected value is a long-run average, not
necessarily a possible observation.

\end{tcolorbox}

\subsection{Operational meaning: long-run
average}\label{operational-meaning-long-run-average}

This subsection explains why the expected value is the number we expect
empirical averages to approach after many repeated trials.

Suppose we measure a random variable \(X\) in \(n\) independent trials
and record \[X_1,X_2,\ldots,X_n.\] The sample average is
\[\overline X_n=\frac{1}{n}(X_1+\cdots+X_n).\] The operational meaning
of expected value is that \(\mathbb{E}[X]\) is the long-run average
value of repeated measurements of \(X\). Later, the Law of Large Numbers
will make this statement precise:
\[\lim_{n\to\infty}\overline X_n=\mathbb{E}[X]\] in an appropriate
probabilistic sense.

\begin{tcolorbox}[enhanced jigsaw, arc=.35mm, bottomrule=.15mm, bottomtitle=1mm, breakable, colback=white, colbacktitle=quarto-callout-tip-color!10!white, colframe=quarto-callout-tip-color-frame, coltitle=black, left=2mm, leftrule=.75mm, opacityback=0, opacitybacktitle=0.6, rightrule=.15mm, title=\textcolor{quarto-callout-tip-color}{\faLightbulb}\hspace{0.5em}{Interpretation}, titlerule=0mm, toprule=.15mm, toptitle=1mm]

If a casino game has expected gain \(-0.05\) dollars per play, then one
play can be positive or negative, but over many plays the average gain
per play tends to be close to \(-0.05\).

\end{tcolorbox}

\subsection{Linearity of expectation}\label{linearity-of-expectation}

This subsection presents one of the most useful properties in
probability: expectation is linear.

\textbf{Theorem 4} (Linearity for one random variable). \emph{For
constants \(a,b\in\mathbb{R}\), \[\mathbb{E}[aX+b]=a\mathbb{E}[X]+b.\]}

\begin{proof}
\emph{Proof.} For the discrete case,
\[\mathbb{E}[aX+b]=\sum_x (ax+b)p_X(x)
=a\sum_x xp_X(x)+b\sum_x p_X(x)=a\mathbb{E}[X]+b.\] The continuous case
is the same argument with integrals:
\[\mathbb{E}[aX+b]=\int_{-\infty}^{\infty} (ax+b)f_X(x)\,dx
=a\mathbb{E}[X]+b.\]~◻
\end{proof}

\textbf{Theorem 5} (Linearity for multiple random variables). \emph{For
any random variables \(X\) and \(Y\) defined on the same probability
space, \[\mathbb{E}[aX+bY]=a\mathbb{E}[X]+b\mathbb{E}[Y].\] This is true
whether or not \(X\) and \(Y\) are independent.}

\begin{tcolorbox}[enhanced jigsaw, arc=.35mm, bottomrule=.15mm, bottomtitle=1mm, breakable, colback=white, colbacktitle=quarto-callout-warning-color!10!white, colframe=quarto-callout-warning-color-frame, coltitle=black, left=2mm, leftrule=.75mm, opacityback=0, opacitybacktitle=0.6, rightrule=.15mm, title=\textcolor{quarto-callout-warning-color}{\faExclamationTriangle}\hspace{0.5em}{Common mistake}, titlerule=0mm, toprule=.15mm, toptitle=1mm]

Linearity says \(\mathbb{E}[aX+bY]=a\mathbb{E}[X]+b\mathbb{E}[Y]\). It
does \emph{not} say \(\mathbb{E}[g(X)]=g(\mathbb{E}[X])\) for a
nonlinear function \(g\).

\end{tcolorbox}

\textbf{Example 6} (Linearity without independence). Let \(X\) be any
random variable with \(\mathbb{E}[X]=3\), and let \(Y=2X+1\). Find
\(\mathbb{E}[4X-5Y]\).

\begin{tcolorbox}[enhanced jigsaw, arc=.35mm, bottomrule=.15mm, bottomtitle=1mm, breakable, colback=white, colbacktitle=quarto-callout-note-color!10!white, colframe=quarto-callout-note-color-frame, coltitle=black, left=2mm, leftrule=.75mm, opacityback=0, opacitybacktitle=0.6, rightrule=.15mm, title=\textcolor{quarto-callout-note-color}{\faInfo}\hspace{0.5em}{Solution}, titlerule=0mm, toprule=.15mm, toptitle=1mm]

Even though \(X\) and \(Y\) are clearly dependent, linearity still
applies. First, \[\mathbb{E}[Y]=\mathbb{E}[2X+1]=2\mathbb{E}[X]+1=7.\]
Therefore,
\[\mathbb{E}[4X-5Y]=4\mathbb{E}[X]-5\mathbb{E}[Y]=4(3)-5(7)=12-35=-23.\]

\end{tcolorbox}

\section{Variance and Standard
Deviation}\label{variance-and-standard-deviation-1}

This section introduces variance as a measurement of spread around the
expected value.

\subsection{Definition and calculation
formula}\label{definition-and-calculation-formula}

This subsection gives both the conceptual definition and the
computational formula for variance.

\textbf{Definition 7} (Variance and standard deviation). The variance of
a random variable \(X\) is
\[\operatorname{Var}(X)=\mathbb{E}\big[(X-\mathbb{E}[X])^2\big].\] The
standard deviation is
\[\operatorname{SD}(X)=\sqrt{\operatorname{Var}(X)}.\]

Variance is the expected squared distance from the mean. It measures how
spread out the distribution is. Standard deviation puts the measurement
back into the original units of \(X\).

\textbf{Proposition 8} (Computational formula).
\emph{\[\operatorname{Var}(X)=\mathbb{E}[X^2]-\big(\mathbb{E}[X]\big)^2.\]}

\begin{proof}
\emph{Proof.} Let \(\mu=\mathbb{E}[X]\). Then
\[\operatorname{Var}(X)=\mathbb{E}[(X-\mu)^2]
=\mathbb{E}[X^2-2\mu X+\mu^2]
=\mathbb{E}[X^2]-2\mu\mathbb{E}[X]+\mu^2.\] Since \(\mu=\mathbb{E}[X]\),
this becomes
\[\operatorname{Var}(X)=\mathbb{E}[X^2]-2\mu^2+\mu^2=\mathbb{E}[X^2]-\mu^2.\]~◻
\end{proof}

\textbf{Proposition 9} (Scaling and shifting). \emph{For constants
\(a,b\in\mathbb{R}\),
\[\operatorname{Var}(aX+b)=a^2\operatorname{Var}(X).\]}

\begin{proof}
\emph{Proof.} Since \(\mathbb{E}[aX+b]=a\mathbb{E}[X]+b\),
\[(aX+b)-\mathbb{E}[aX+b]=aX+b-(a\mathbb{E}[X]+b)=a(X-\mathbb{E}[X]).\]
Therefore,
\[\operatorname{Var}(aX+b)=\mathbb{E}\left[a^2(X-\mathbb{E}[X])^2\right]=a^2\operatorname{Var}(X).\]~◻
\end{proof}

\textbf{Example 10} (Variance of a Bernoulli random variable). Let
\(X\sim\operatorname{Bernoulli}(p)\). Find \(\operatorname{Var}(X)\).

\begin{tcolorbox}[enhanced jigsaw, arc=.35mm, bottomrule=.15mm, bottomtitle=1mm, breakable, colback=white, colbacktitle=quarto-callout-note-color!10!white, colframe=quarto-callout-note-color-frame, coltitle=black, left=2mm, leftrule=.75mm, opacityback=0, opacitybacktitle=0.6, rightrule=.15mm, title=\textcolor{quarto-callout-note-color}{\faInfo}\hspace{0.5em}{Solution}, titlerule=0mm, toprule=.15mm, toptitle=1mm]

Since \(X\) takes values \(0\) and \(1\), we have \(X^2=X\). Therefore
\[\mathbb{E}[X^2]=\mathbb{E}[X]=p.\] Using the computational formula,
\[\operatorname{Var}(X)=\mathbb{E}[X^2]-(\mathbb{E}[X])^2=p-p^2=p(1-p).\]

\end{tcolorbox}

\textbf{Example 11} (Variance of a fair die). Let \(X\) be the outcome
of rolling a fair six-sided die. Find \(\operatorname{Var}(X)\).

\begin{tcolorbox}[enhanced jigsaw, arc=.35mm, bottomrule=.15mm, bottomtitle=1mm, breakable, colback=white, colbacktitle=quarto-callout-note-color!10!white, colframe=quarto-callout-note-color-frame, coltitle=black, left=2mm, leftrule=.75mm, opacityback=0, opacitybacktitle=0.6, rightrule=.15mm, title=\textcolor{quarto-callout-note-color}{\faInfo}\hspace{0.5em}{Solution}, titlerule=0mm, toprule=.15mm, toptitle=1mm]

We already know \(\mathbb{E}[X]=3.5=7/2\). Next,
\[\mathbb{E}[X^2]=\frac{1^2+2^2+3^2+4^2+5^2+6^2}{6}
=\frac{91}{6}.\] Thus
\[\operatorname{Var}(X)=\frac{91}{6}-\left(\frac72\right)^2
=\frac{91}{6}-\frac{49}{4}
=\frac{182-147}{12}=\frac{35}{12}.\] So the standard deviation is
\[\operatorname{SD}(X)=\sqrt{\frac{35}{12}}.\]

\end{tcolorbox}

\section{Expectation of a Function}\label{expectation-of-a-function}

This section explains how to compute the expectation of a transformed
random variable without first finding the full distribution of the
transformed variable.

\subsection{The law of the unconscious
statistician}\label{the-law-of-the-unconscious-statistician}

This subsection introduces a practical formula: to compute
\(\mathbb{E}[g(X)]\), average \(g(x)\) with respect to the distribution
of \(X\).

Suppose \(Y=g(X)\). One way to compute \(\mathbb{E}[Y]\) is to first
find the distribution of \(Y\) and then sum or integrate over \(Y\).
Often this is unnecessary.

\textbf{Theorem 12} (Expectation of a function). \emph{Let \(Y=g(X)\).}

\begin{itemize}
\item
  \emph{If \(X\) is discrete, then
  \[\mathbb{E}[g(X)]=\sum_x g(x)p_X(x).\]}
\item
  \emph{If \(X\) is continuous, then
  \[\mathbb{E}[g(X)]=\int_{-\infty}^{\infty} g(x)f_X(x)\,dx.\]}
\end{itemize}

\begin{proof}
\emph{Proof idea for the discrete case.} Let the possible values of
\(X\) be \(x_k\). Then \[\mathbb{E}[Y]=\sum_y y\mathbb{P}(Y=y)
=\sum_y y\sum_{k:g(x_k)=y}\mathbb{P}(X=x_k).\] Since \(y=g(x_k)\) on the
event \(\{g(X)=y\}\),
\[\mathbb{E}[Y]=\sum_k g(x_k)\mathbb{P}(X=x_k).\]~◻
\end{proof}

\textbf{Example 13} (Computing a second moment directly). Let \(X\) be
the outcome of a fair six-sided die. Use the function formula to compute
\(\mathbb{E}[X^2]\).

\begin{tcolorbox}[enhanced jigsaw, arc=.35mm, bottomrule=.15mm, bottomtitle=1mm, breakable, colback=white, colbacktitle=quarto-callout-note-color!10!white, colframe=quarto-callout-note-color-frame, coltitle=black, left=2mm, leftrule=.75mm, opacityback=0, opacitybacktitle=0.6, rightrule=.15mm, title=\textcolor{quarto-callout-note-color}{\faInfo}\hspace{0.5em}{Solution}, titlerule=0mm, toprule=.15mm, toptitle=1mm]

Here \(g(x)=x^2\) and \(p_X(x)=1/6\) for \(x=1,2,3,4,5,6\). Therefore
\[\mathbb{E}[X^2]=\sum_{x=1}^6 x^2\frac16
=\frac{1+4+9+16+25+36}{6}=\frac{91}{6}.\]

\end{tcolorbox}

\textbf{Example 14} (Uniform distribution). Let
\(X\sim \operatorname{Uniform}(0,1)\). Compute \(\mathbb{E}[X^m]\) for
an integer \(m\ge 1\).

\begin{tcolorbox}[enhanced jigsaw, arc=.35mm, bottomrule=.15mm, bottomtitle=1mm, breakable, colback=white, colbacktitle=quarto-callout-note-color!10!white, colframe=quarto-callout-note-color-frame, coltitle=black, left=2mm, leftrule=.75mm, opacityback=0, opacitybacktitle=0.6, rightrule=.15mm, title=\textcolor{quarto-callout-note-color}{\faInfo}\hspace{0.5em}{Solution}, titlerule=0mm, toprule=.15mm, toptitle=1mm]

The pdf is \(f_X(x)=1\) for \(0\le x\le 1\). Thus
\[\mathbb{E}[X^m]=\int_0^1 x^m\,dx=\frac{1}{m+1}.\] In particular,
\(\mathbb{E}[X]=1/2\) and \(\mathbb{E}[X^2]=1/3\).

\end{tcolorbox}

\begin{tcolorbox}[enhanced jigsaw, arc=.35mm, bottomrule=.15mm, bottomtitle=1mm, breakable, colback=white, colbacktitle=quarto-callout-warning-color!10!white, colframe=quarto-callout-warning-color-frame, coltitle=black, left=2mm, leftrule=.75mm, opacityback=0, opacitybacktitle=0.6, rightrule=.15mm, title=\textcolor{quarto-callout-warning-color}{\faExclamationTriangle}\hspace{0.5em}{Important distinction}, titlerule=0mm, toprule=.15mm, toptitle=1mm]

In general, \[\mathbb{E}[g(X)]\ne g(\mathbb{E}[X]).\] For example, if
\(X\) is not constant, then \(\mathbb{E}[X^2]\ne (\mathbb{E}[X])^2\) in
general. Their difference is the variance.

\end{tcolorbox}

\subsection{Products of independent random
variables}\label{products-of-independent-random-variables}

This subsection records the special multiplication rule for independent
variables.

\textbf{Theorem 15} (Expectation of products under independence).
\emph{If \(X\) and \(Y\) are independent, then
\[\mathbb{E}[XY]=\mathbb{E}[X]\mathbb{E}[Y].\] More generally, for
suitable functions \(f\) and \(g\),
\[\mathbb{E}[f(X)g(Y)]=\mathbb{E}[f(X)]\mathbb{E}[g(Y)].\]}

\begin{tcolorbox}[enhanced jigsaw, arc=.35mm, bottomrule=.15mm, bottomtitle=1mm, breakable, colback=white, colbacktitle=quarto-callout-note-color!10!white, colframe=quarto-callout-note-color-frame, coltitle=black, left=2mm, leftrule=.75mm, opacityback=0, opacitybacktitle=0.6, rightrule=.15mm, title=\textcolor{quarto-callout-note-color}{\faInfo}\hspace{0.5em}{Remark}, titlerule=0mm, toprule=.15mm, toptitle=1mm]

\emph{Remark 16}. The converse is not true in general. The equality
\(\mathbb{E}[XY]=\mathbb{E}[X]\mathbb{E}[Y]\) only says \(X\) and \(Y\)
are uncorrelated; it does not necessarily imply independence.

\end{tcolorbox}

\textbf{Example 17} (Independent Bernoulli variables). Let
\(X\sim\operatorname{Bernoulli}(p)\) and
\(Y\sim\operatorname{Bernoulli}(q)\) be independent. Find
\(\mathbb{E}[XY]\).

\begin{tcolorbox}[enhanced jigsaw, arc=.35mm, bottomrule=.15mm, bottomtitle=1mm, breakable, colback=white, colbacktitle=quarto-callout-note-color!10!white, colframe=quarto-callout-note-color-frame, coltitle=black, left=2mm, leftrule=.75mm, opacityback=0, opacitybacktitle=0.6, rightrule=.15mm, title=\textcolor{quarto-callout-note-color}{\faInfo}\hspace{0.5em}{Solution}, titlerule=0mm, toprule=.15mm, toptitle=1mm]

By independence, \[\mathbb{E}[XY]=\mathbb{E}[X]\mathbb{E}[Y]=pq.\] This
also has a direct probability interpretation: \(XY=1\) exactly when
\(X=1\) and \(Y=1\), so \[\mathbb{E}[XY]=\mathbb{P}(X=1,Y=1)=pq.\]

\end{tcolorbox}

\section{Moments}\label{moments}

This section introduces moments as numerical summaries that describe
center, spread, asymmetry, and tail behavior.

\subsection{Raw and centered moments}\label{raw-and-centered-moments}

This subsection defines the main moment quantities used in probability
and statistics.

\textbf{Definition 18} (Raw moment). The \(m\)-th moment, or \(m\)-th
raw moment, of a random variable \(X\) is \[\mathbb{E}[X^m].\]

\textbf{Definition 19} (Centered moment). The \(m\)-th centered moment
of a random variable \(X\) is
\[\mathbb{E}\left[(X-\mathbb{E}[X])^m\right].\]

The first raw moment is the mean. The second centered moment is the
variance:
\[\operatorname{Var}(X)=\mathbb{E}\left[(X-\mathbb{E}[X])^2\right].\]
The square root of variance is the standard deviation:
\[\sigma=\sqrt{\operatorname{Var}(X)}.\]

\textbf{Example 20} (First two moments of
\(\operatorname{Uniform}(0,1)\)). Let
\(X\sim \operatorname{Uniform}(0,1)\). Compute \(\mathbb{E}[X]\),
\(\mathbb{E}[X^2]\), and \(\operatorname{Var}(X)\).

\begin{tcolorbox}[enhanced jigsaw, arc=.35mm, bottomrule=.15mm, bottomtitle=1mm, breakable, colback=white, colbacktitle=quarto-callout-note-color!10!white, colframe=quarto-callout-note-color-frame, coltitle=black, left=2mm, leftrule=.75mm, opacityback=0, opacitybacktitle=0.6, rightrule=.15mm, title=\textcolor{quarto-callout-note-color}{\faInfo}\hspace{0.5em}{Solution}, titlerule=0mm, toprule=.15mm, toptitle=1mm]

Using the previous formula \(\mathbb{E}[X^m]=1/(m+1)\),
\[\mathbb{E}[X]=\frac12,\qquad \mathbb{E}[X^2]=\frac13.\] Therefore
\[\operatorname{Var}(X)=\mathbb{E}[X^2]-(\mathbb{E}[X])^2=\frac13-\frac14=\frac{1}{12}.\]

\end{tcolorbox}

\subsection{Mean and variance table}\label{mean-and-variance-table-1}

This subsection summarizes common means and variances that will be used
repeatedly in the course.

{\def\LTcaptype{none} % do not increment counter
\begin{longtable}[]{@{}
  >{\raggedright\arraybackslash}p{(\linewidth - 4\tabcolsep) * \real{0.4143}}
  >{\raggedright\arraybackslash}p{(\linewidth - 4\tabcolsep) * \real{0.1786}}
  >{\raggedright\arraybackslash}p{(\linewidth - 4\tabcolsep) * \real{0.4071}}@{}}
\toprule\noalign{}
\begin{minipage}[b]{\linewidth}\raggedright
\textbf{Name}
\end{minipage} & \begin{minipage}[b]{\linewidth}\raggedright
\textbf{Mean}
\end{minipage} & \begin{minipage}[b]{\linewidth}\raggedright
\textbf{Variance}
\end{minipage} \\
\midrule\noalign{}
\endhead
\bottomrule\noalign{}
\endlastfoot
\(\operatorname{Bernoulli}(p)\) & \(p\) & \(p(1-p)\) \\
\(\operatorname{Binomial}(n,p)\) & \(np\) & \(np(1-p)\) \\
\(\operatorname{Geometric}(p)\) & \(1/p\) & \((1-p)/p^2\) \\
\(\operatorname{Poisson}(\lambda)\) & \(\lambda\) & \(\lambda\) \\
\(\operatorname{Exponential}(\lambda)\) & \(1/\lambda\) &
\(1/\lambda^2\) \\
\(\operatorname{Gamma}(n,\lambda)\), rate parameterization &
\(n/\lambda\) & \(n/\lambda^2\) \\
\(\operatorname{Uniform}(a,b)\) & \((a+b)/2\) & \((b-a)^2/12\) \\
\(\operatorname{Normal}(\mu,\sigma^2)\) & \(\mu\) & \(\sigma^2\) \\
\(\operatorname{Beta}(\alpha,\beta)\) & \(\alpha/(\alpha+\beta)\) &
\(\dfrac{\alpha\beta}{(\alpha+\beta)^2(\alpha+\beta+1)}\) \\
\end{longtable}
}

\begin{tcolorbox}[enhanced jigsaw, arc=.35mm, bottomrule=.15mm, bottomtitle=1mm, breakable, colback=white, colbacktitle=quarto-callout-note-color!10!white, colframe=quarto-callout-note-color-frame, coltitle=black, left=2mm, leftrule=.75mm, opacityback=0, opacitybacktitle=0.6, rightrule=.15mm, title=\textcolor{quarto-callout-note-color}{\faInfo}\hspace{0.5em}{Remark}, titlerule=0mm, toprule=.15mm, toptitle=1mm]

\emph{Remark 21}. Some texts parameterize the Gamma distribution by
scale \(\theta\) instead of rate \(\lambda\). If
\(X\sim\operatorname{Gamma}(n,\theta)\) in the scale parameterization,
then \(\mathbb{E}[X]=n\theta\) and \(\operatorname{Var}(X)=n\theta^2\).
If \(\lambda=1/\theta\), then this becomes \(n/\lambda\) and
\(n/\lambda^2\).

\end{tcolorbox}

\subsection{Skewness}\label{skewness}

This subsection describes how the third standardized moment measures
asymmetry.

\textbf{Definition 22} (Skewness). Let \(\mu=\mathbb{E}[X]\) and
\(\sigma=\sqrt{\operatorname{Var}(X)}\). The skewness of \(X\) is
\[\mathbb{E}\left[\left(\frac{X-\mu}{\sigma}\right)^3\right]
=\frac{\mathbb{E}[(X-\mu)^3]}{\sigma^3}.\]

Skewness measures asymmetry of a probability distribution. A
distribution with a long right tail usually has positive skewness. A
distribution with a long left tail usually has negative skewness.

\textbf{Example 23} (Symmetric distributions have zero skewness).
Suppose \(X\) has a distribution symmetric around its mean \(\mu\). What
is the skewness?

\begin{tcolorbox}[enhanced jigsaw, arc=.35mm, bottomrule=.15mm, bottomtitle=1mm, breakable, colback=white, colbacktitle=quarto-callout-note-color!10!white, colframe=quarto-callout-note-color-frame, coltitle=black, left=2mm, leftrule=.75mm, opacityback=0, opacitybacktitle=0.6, rightrule=.15mm, title=\textcolor{quarto-callout-note-color}{\faInfo}\hspace{0.5em}{Solution}, titlerule=0mm, toprule=.15mm, toptitle=1mm]

Let \(Z=X-\mu\). Symmetry around \(\mu\) means \(Z\) and \(-Z\) have the
same distribution. Therefore
\(\mathbb{E}[Z^3]=\mathbb{E}[(-Z)^3]=-\mathbb{E}[Z^3]\), so
\(\mathbb{E}[Z^3]=0\). Hence
\[\frac{\mathbb{E}[(X-\mu)^3]}{\sigma^3}=0.\] Thus the skewness is
\(0\).

\end{tcolorbox}

\subsection{Kurtosis}\label{kurtosis}

This subsection describes how the fourth standardized moment measures
tail behavior and peakedness.

\textbf{Definition 24} (Kurtosis). Let \(\mu=\mathbb{E}[X]\) and
\(\sigma=\sqrt{\operatorname{Var}(X)}\). The kurtosis of \(X\) is
\[\mathbb{E}\left[\left(\frac{X-\mu}{\sigma}\right)^4\right]
=\frac{\mathbb{E}[(X-\mu)^4]}{\sigma^4}.\]

Kurtosis characterizes the ``tailedness'' of a distribution.
Distributions with heavier tails often have larger kurtosis. Examples
often compared by tail behavior include the Laplace (double
exponential), hyperbolic secant, logistic, normal, raised cosine, Wigner
semicircle, and uniform distributions.

\textbf{Example 25} (Kurtosis of the standard normal). Let
\(Z\sim \operatorname{Normal}(0,1)\). What is the kurtosis of \(Z\)?

\begin{tcolorbox}[enhanced jigsaw, arc=.35mm, bottomrule=.15mm, bottomtitle=1mm, breakable, colback=white, colbacktitle=quarto-callout-note-color!10!white, colframe=quarto-callout-note-color-frame, coltitle=black, left=2mm, leftrule=.75mm, opacityback=0, opacitybacktitle=0.6, rightrule=.15mm, title=\textcolor{quarto-callout-note-color}{\faInfo}\hspace{0.5em}{Solution}, titlerule=0mm, toprule=.15mm, toptitle=1mm]

For the standard normal distribution,
\[\mathbb{E}[Z]=0,\qquad \operatorname{Var}(Z)=1,\qquad \mathbb{E}[Z^4]=3.\]
Thus
\[\text{kurtosis}=\mathbb{E}\left[\left(\frac{Z-0}{1}\right)^4\right]=\mathbb{E}[Z^4]=3.\]
The excess kurtosis is often defined as kurtosis minus \(3\), so the
standard normal has excess kurtosis \(0\).

\end{tcolorbox}

\section{Moment Generating Functions}\label{moment-generating-functions}

This section introduces moment generating functions and explains why
they are useful for moments, distribution identification, and sums.

\subsection{Definition and moment
generation}\label{definition-and-moment-generation}

This subsection defines the MGF and shows how derivatives of the MGF
produce moments.

\textbf{Definition 26} (Moment generating function). The moment
generating function (MGF) of a random variable \(X\) is
\[M_X(t)=\mathbb{E}[e^{tX}],\] for values of \(t\) where the expectation
exists.

When \(M_X(t)\) exists in an open neighborhood of \(0\), Taylor
expansion gives
\[e^{tX}=1+tX+\frac{(tX)^2}{2!}+\frac{(tX)^3}{3!}+\cdots.\] Taking
expectations,
\[M_X(t)=1+t\mathbb{E}[X]+\frac{t^2\mathbb{E}[X^2]}{2!}+\frac{t^3\mathbb{E}[X^3]}{3!}+\cdots.\]
Therefore,
\[\mathbb{E}[X^m]=M_X^{(m)}(0)=\left.\frac{d^m}{dt^m}M_X(t)\right|_{t=0}.\]

\begin{tcolorbox}[enhanced jigsaw, arc=.35mm, bottomrule=.15mm, bottomtitle=1mm, breakable, colback=white, colbacktitle=quarto-callout-tip-color!10!white, colframe=quarto-callout-tip-color-frame, coltitle=black, left=2mm, leftrule=.75mm, opacityback=0, opacitybacktitle=0.6, rightrule=.15mm, title=\textcolor{quarto-callout-tip-color}{\faLightbulb}\hspace{0.5em}{Why the name ``moment generating''?}, titlerule=0mm, toprule=.15mm, toptitle=1mm]

The derivatives of \(M_X(t)\) at \(t=0\) generate the raw moments
\(\mathbb{E}[X]\), \(\mathbb{E}[X^2]\), \(\mathbb{E}[X^3]\), and so on.

\end{tcolorbox}

\textbf{Example 27} (Using an MGF to find moments). Suppose a random
variable has MGF \[M_X(t)=\frac{1}{1-t},\qquad t<1.\] Find
\(\mathbb{E}[X]\) and \(\mathbb{E}[X^2]\).

\begin{tcolorbox}[enhanced jigsaw, arc=.35mm, bottomrule=.15mm, bottomtitle=1mm, breakable, colback=white, colbacktitle=quarto-callout-note-color!10!white, colframe=quarto-callout-note-color-frame, coltitle=black, left=2mm, leftrule=.75mm, opacityback=0, opacitybacktitle=0.6, rightrule=.15mm, title=\textcolor{quarto-callout-note-color}{\faInfo}\hspace{0.5em}{Solution}, titlerule=0mm, toprule=.15mm, toptitle=1mm]

Differentiate: \[M_X'(t)=\frac{1}{(1-t)^2},
\qquad
M_X''(t)=\frac{2}{(1-t)^3}.\] Thus \[\mathbb{E}[X]=M_X'(0)=1,
\qquad
\mathbb{E}[X^2]=M_X''(0)=2.\] Therefore
\[\operatorname{Var}(X)=\mathbb{E}[X^2]-(\mathbb{E}[X])^2=2-1^2=1.\]

\end{tcolorbox}

\subsection{Uniqueness and transformation
rules}\label{uniqueness-and-transformation-rules}

This subsection states two key reasons MGFs are useful: they can
identify distributions and they behave nicely under shifts, scales, and
independent sums.

\textbf{Theorem 28} (MGF determines the distribution). \emph{Suppose
\(M_X(t)\) and \(M_Y(t)\) exist for all \(t\) in an open neighborhood of
\(0\). If \[M_X(t)=M_Y(t)\] for all such \(t\), then \(X\) and \(Y\)
have the same distribution.}

\textbf{Theorem 29} (Bounded support and moments). \emph{Suppose all
moments exist for random variables \(X\) and \(Y\). If \(X\) and \(Y\)
have bounded support, then the CDFs of \(X\) and \(Y\) are equal if and
only if all moments are equal.}

\textbf{Theorem 30} (MGF transformation rules). \emph{Let
\(a,b\in\mathbb{R}\). \[M_{aX+b}(t)=e^{bt}M_X(at).\] If \(X\) and \(Y\)
are independent, then \[M_{X+Y}(t)=M_X(t)M_Y(t).\]}

\begin{proof}
\emph{Proof.} For the affine transformation,
\[M_{aX+b}(t)=\mathbb{E}[e^{t(aX+b)}]=e^{bt}\mathbb{E}[e^{(at)X}]=e^{bt}M_X(at).\]
If \(X\) and \(Y\) are independent,
\[M_{X+Y}(t)=\mathbb{E}[e^{t(X+Y)}]=\mathbb{E}[e^{tX}e^{tY}]=\mathbb{E}[e^{tX}]\mathbb{E}[e^{tY}]=M_X(t)M_Y(t).\]~◻
\end{proof}

\begin{tcolorbox}[enhanced jigsaw, arc=.35mm, bottomrule=.15mm, bottomtitle=1mm, breakable, colback=white, colbacktitle=quarto-callout-important-color!10!white, colframe=quarto-callout-important-color-frame, coltitle=black, left=2mm, leftrule=.75mm, opacityback=0, opacitybacktitle=0.6, rightrule=.15mm, title=\textcolor{quarto-callout-important-color}{\faExclamation}\hspace{0.5em}{Practice Problem}, titlerule=0mm, toprule=.15mm, toptitle=1mm]

\textbf{Practice Problem 31} (Affine transformation). If \(M_X(t)\) is
known, find the MGF of \(Y=3X-2\).

\end{tcolorbox}

\begin{tcolorbox}[enhanced jigsaw, arc=.35mm, bottomrule=.15mm, bottomtitle=1mm, breakable, colback=white, colbacktitle=quarto-callout-note-color!10!white, colframe=quarto-callout-note-color-frame, coltitle=black, left=2mm, leftrule=.75mm, opacityback=0, opacitybacktitle=0.6, rightrule=.15mm, title=\textcolor{quarto-callout-note-color}{\faInfo}\hspace{0.5em}{Solution}, titlerule=0mm, toprule=.15mm, toptitle=1mm]

Use \(a=3\) and \(b=-2\): \[M_Y(t)=M_{3X-2}(t)=e^{-2t}M_X(3t).\]

\end{tcolorbox}

\section{MGFs of Common
Distributions}\label{mgfs-of-common-distributions}

This section computes MGFs for several common distributions and uses
them to recover moments.

\subsection{Bernoulli and Binomial
distributions}\label{bernoulli-and-binomial-distributions}

This subsection shows how the binomial MGF follows from the Bernoulli
MGF and independence.

\textbf{Example 32} (Bernoulli MGF). Let
\(X\sim\operatorname{Bernoulli}(p)\). Find \(M_X(t)\).

\begin{tcolorbox}[enhanced jigsaw, arc=.35mm, bottomrule=.15mm, bottomtitle=1mm, breakable, colback=white, colbacktitle=quarto-callout-note-color!10!white, colframe=quarto-callout-note-color-frame, coltitle=black, left=2mm, leftrule=.75mm, opacityback=0, opacitybacktitle=0.6, rightrule=.15mm, title=\textcolor{quarto-callout-note-color}{\faInfo}\hspace{0.5em}{Solution}, titlerule=0mm, toprule=.15mm, toptitle=1mm]

Since \(X=1\) with probability \(p\) and \(X=0\) with probability
\(1-p\),
\[M_X(t)=\mathbb{E}[e^{tX}]=e^{t\cdot 1}p+e^{t\cdot 0}(1-p)=pe^t+(1-p).\]
Therefore, \[M_X(t)=1-p+pe^t.\] As a check,
\[M_X'(t)=pe^t,\qquad M_X'(0)=p=\mathbb{E}[X].\]

\end{tcolorbox}

\textbf{Example 33} (Binomial MGF). Let
\(Y\sim\operatorname{Binomial}(n,p)\). Find \(M_Y(t)\).

\begin{tcolorbox}[enhanced jigsaw, arc=.35mm, bottomrule=.15mm, bottomtitle=1mm, breakable, colback=white, colbacktitle=quarto-callout-note-color!10!white, colframe=quarto-callout-note-color-frame, coltitle=black, left=2mm, leftrule=.75mm, opacityback=0, opacitybacktitle=0.6, rightrule=.15mm, title=\textcolor{quarto-callout-note-color}{\faInfo}\hspace{0.5em}{Solution}, titlerule=0mm, toprule=.15mm, toptitle=1mm]

Write \[Y=X_1+\cdots+X_n,\] where \(X_i\sim\operatorname{Bernoulli}(p)\)
are independent. Then
\[M_Y(t)=\prod_{i=1}^n M_{X_i}(t)=\left(1-p+pe^t\right)^n.\] Thus
\[M_Y(t)=\left(pe^t+1-p\right)^n.\]

\end{tcolorbox}

\begin{tcolorbox}[enhanced jigsaw, arc=.35mm, bottomrule=.15mm, bottomtitle=1mm, breakable, colback=white, colbacktitle=quarto-callout-important-color!10!white, colframe=quarto-callout-important-color-frame, coltitle=black, left=2mm, leftrule=.75mm, opacityback=0, opacitybacktitle=0.6, rightrule=.15mm, title=\textcolor{quarto-callout-important-color}{\faExclamation}\hspace{0.5em}{Practice Problem}, titlerule=0mm, toprule=.15mm, toptitle=1mm]

\textbf{Practice Problem 34} (Mean and variance from the binomial MGF).
Use the binomial MGF \(M_Y(t)=(1-p+pe^t)^n\) to find \(\mathbb{E}[Y]\)
and \(\operatorname{Var}(Y)\).

\end{tcolorbox}

\begin{tcolorbox}[enhanced jigsaw, arc=.35mm, bottomrule=.15mm, bottomtitle=1mm, breakable, colback=white, colbacktitle=quarto-callout-note-color!10!white, colframe=quarto-callout-note-color-frame, coltitle=black, left=2mm, leftrule=.75mm, opacityback=0, opacitybacktitle=0.6, rightrule=.15mm, title=\textcolor{quarto-callout-note-color}{\faInfo}\hspace{0.5em}{Solution}, titlerule=0mm, toprule=.15mm, toptitle=1mm]

Let \(q=1-p+pe^t\). Then \[M_Y'(t)=nq^{n-1}pe^t,\] so
\[\mathbb{E}[Y]=M_Y'(0)=n(1)^{n-1}p=np.\] For the second derivative,
\[M_Y''(t)=n(n-1)q^{n-2}(pe^t)^2+nq^{n-1}pe^t.\] Thus
\[\mathbb{E}[Y^2]=M_Y''(0)=n(n-1)p^2+np.\] Therefore
\[\operatorname{Var}(Y)=\mathbb{E}[Y^2]-(\mathbb{E}[Y])^2
=n(n-1)p^2+np-n^2p^2=np(1-p).\]

\end{tcolorbox}

\subsection{Poisson distribution}\label{poisson-distribution-1}

This subsection derives the MGF of the Poisson distribution using the
exponential series.

\textbf{Example 35} (Poisson MGF). Let
\(X\sim\operatorname{Poisson}(\lambda)\). Find \(M_X(t)\).

\begin{tcolorbox}[enhanced jigsaw, arc=.35mm, bottomrule=.15mm, bottomtitle=1mm, breakable, colback=white, colbacktitle=quarto-callout-note-color!10!white, colframe=quarto-callout-note-color-frame, coltitle=black, left=2mm, leftrule=.75mm, opacityback=0, opacitybacktitle=0.6, rightrule=.15mm, title=\textcolor{quarto-callout-note-color}{\faInfo}\hspace{0.5em}{Solution}, titlerule=0mm, toprule=.15mm, toptitle=1mm]

The pmf is
\[\mathbb{P}(X=k)=\frac{\lambda^k e^{-\lambda}}{k!},\qquad k=0,1,2,\ldots.\]
Therefore \[\begin{aligned}
M_X(t)
&=\mathbb{E}[e^{tX}]
=\sum_{k=0}^{\infty} e^{tk}\frac{\lambda^k e^{-\lambda}}{k!}\\
&=e^{-\lambda}\sum_{k=0}^{\infty}\frac{(\lambda e^t)^k}{k!}
=e^{-\lambda}\exp(\lambda e^t)\\
&=\exp\{\lambda(e^t-1)\}.
\end{aligned}\] Thus \[M_X(t)=e^{\lambda(e^t-1)}.\]

\end{tcolorbox}

\begin{tcolorbox}[enhanced jigsaw, arc=.35mm, bottomrule=.15mm, bottomtitle=1mm, breakable, colback=white, colbacktitle=quarto-callout-important-color!10!white, colframe=quarto-callout-important-color-frame, coltitle=black, left=2mm, leftrule=.75mm, opacityback=0, opacitybacktitle=0.6, rightrule=.15mm, title=\textcolor{quarto-callout-important-color}{\faExclamation}\hspace{0.5em}{Practice Problem}, titlerule=0mm, toprule=.15mm, toptitle=1mm]

\textbf{Practice Problem 36} (Sum of independent Poisson variables).
Suppose \(X\sim\operatorname{Poisson}(\lambda_1)\) and
\(Y\sim\operatorname{Poisson}(\lambda_2)\) are independent. Use MGFs to
find the distribution of \(X+Y\).

\end{tcolorbox}

\begin{tcolorbox}[enhanced jigsaw, arc=.35mm, bottomrule=.15mm, bottomtitle=1mm, breakable, colback=white, colbacktitle=quarto-callout-note-color!10!white, colframe=quarto-callout-note-color-frame, coltitle=black, left=2mm, leftrule=.75mm, opacityback=0, opacitybacktitle=0.6, rightrule=.15mm, title=\textcolor{quarto-callout-note-color}{\faInfo}\hspace{0.5em}{Solution}, titlerule=0mm, toprule=.15mm, toptitle=1mm]

By independence, \[M_{X+Y}(t)=M_X(t)M_Y(t)
=e^{\lambda_1(e^t-1)}e^{\lambda_2(e^t-1)}
=e^{(\lambda_1+\lambda_2)(e^t-1)}.\] This is the MGF of a
\(\operatorname{Poisson}(\lambda_1+\lambda_2)\) random variable.
Therefore \[X+Y\sim\operatorname{Poisson}(\lambda_1+\lambda_2).\]

\end{tcolorbox}

\subsection{Exponential distribution}\label{exponential-distribution-1}

This subsection derives the MGF of the exponential distribution by
evaluating an integral.

\textbf{Example 37} (Exponential MGF). Let
\(X\sim\operatorname{Exponential}(\lambda)\) with pdf
\[f_X(x)=\lambda e^{-\lambda x},\qquad x\ge 0.\] Find \(M_X(t)\).

\begin{tcolorbox}[enhanced jigsaw, arc=.35mm, bottomrule=.15mm, bottomtitle=1mm, breakable, colback=white, colbacktitle=quarto-callout-note-color!10!white, colframe=quarto-callout-note-color-frame, coltitle=black, left=2mm, leftrule=.75mm, opacityback=0, opacitybacktitle=0.6, rightrule=.15mm, title=\textcolor{quarto-callout-note-color}{\faInfo}\hspace{0.5em}{Solution}, titlerule=0mm, toprule=.15mm, toptitle=1mm]

For \(t<\lambda\), \[\begin{aligned}
M_X(t)&=\mathbb{E}[e^{tX}]
=\int_0^{\infty} e^{tx}\lambda e^{-\lambda x}\,dx\\
&=\lambda\int_0^{\infty}e^{-(\lambda-t)x}\,dx
=\lambda\cdot \frac{1}{\lambda-t}.
\end{aligned}\] Thus
\[M_X(t)=\frac{\lambda}{\lambda-t},\qquad t<\lambda.\]

\end{tcolorbox}

\begin{tcolorbox}[enhanced jigsaw, arc=.35mm, bottomrule=.15mm, bottomtitle=1mm, breakable, colback=white, colbacktitle=quarto-callout-important-color!10!white, colframe=quarto-callout-important-color-frame, coltitle=black, left=2mm, leftrule=.75mm, opacityback=0, opacitybacktitle=0.6, rightrule=.15mm, title=\textcolor{quarto-callout-important-color}{\faExclamation}\hspace{0.5em}{Practice Problem}, titlerule=0mm, toprule=.15mm, toptitle=1mm]

\textbf{Practice Problem 38} (Mean and variance of an exponential random
variable). Use \(M_X(t)=\lambda/(\lambda-t)\) to find \(\mathbb{E}[X]\)
and \(\operatorname{Var}(X)\).

\end{tcolorbox}

\begin{tcolorbox}[enhanced jigsaw, arc=.35mm, bottomrule=.15mm, bottomtitle=1mm, breakable, colback=white, colbacktitle=quarto-callout-note-color!10!white, colframe=quarto-callout-note-color-frame, coltitle=black, left=2mm, leftrule=.75mm, opacityback=0, opacitybacktitle=0.6, rightrule=.15mm, title=\textcolor{quarto-callout-note-color}{\faInfo}\hspace{0.5em}{Solution}, titlerule=0mm, toprule=.15mm, toptitle=1mm]

Differentiate: \[M_X'(t)=\frac{\lambda}{(\lambda-t)^2},
\qquad
M_X''(t)=\frac{2\lambda}{(\lambda-t)^3}.\] Hence
\[\mathbb{E}[X]=M_X'(0)=\frac{1}{\lambda},
\qquad
\mathbb{E}[X^2]=M_X''(0)=\frac{2}{\lambda^2}.\] Therefore
\[\operatorname{Var}(X)=\mathbb{E}[X^2]-(\mathbb{E}[X])^2
=\frac{2}{\lambda^2}-\frac{1}{\lambda^2}
=\frac{1}{\lambda^2}.\]

\end{tcolorbox}

\subsection{Normal distribution}\label{normal-distribution-1}

This subsection derives the normal MGF and uses it to show closure under
independent sums.

\textbf{Example 39} (Normal MGF). Let
\(X\sim\operatorname{Normal}(\mu,\sigma^2)\). Find \(M_X(t)\).

\begin{tcolorbox}[enhanced jigsaw, arc=.35mm, bottomrule=.15mm, bottomtitle=1mm, breakable, colback=white, colbacktitle=quarto-callout-note-color!10!white, colframe=quarto-callout-note-color-frame, coltitle=black, left=2mm, leftrule=.75mm, opacityback=0, opacitybacktitle=0.6, rightrule=.15mm, title=\textcolor{quarto-callout-note-color}{\faInfo}\hspace{0.5em}{Solution}, titlerule=0mm, toprule=.15mm, toptitle=1mm]

Write \(X=\mu+\sigma Z\), where \(Z\sim\operatorname{Normal}(0,1)\).
First compute the standard normal MGF: \[\begin{aligned}
M_Z(t)&=\mathbb{E}[e^{tZ}]
=\frac{1}{\sqrt{2\pi}}\int_{-\infty}^{\infty} e^{tz}e^{-z^2/2}\,dz\\
&=\frac{1}{\sqrt{2\pi}}\int_{-\infty}^{\infty}\exp\left(-\frac12(z^2-2tz)\right)\,dz\\
&=\frac{1}{\sqrt{2\pi}}\int_{-\infty}^{\infty}\exp\left(-\frac12(z-t)^2+\frac{t^2}{2}\right)\,dz\\
&=e^{t^2/2}.
\end{aligned}\] Now apply the affine transformation rule:
\[M_X(t)=M_{\mu+\sigma Z}(t)=e^{\mu t}M_Z(\sigma t)
=e^{\mu t}e^{\sigma^2t^2/2}.\] Therefore
\[M_X(t)=\exp\left(\mu t+\frac{\sigma^2t^2}{2}\right).\]

\end{tcolorbox}

\textbf{Example 40} (Sum of independent normal random variables).
Suppose \(X\sim\operatorname{Normal}(\mu_1,\sigma_1^2)\) and
\(Y\sim\operatorname{Normal}(\mu_2,\sigma_2^2)\) are independent. Find
the distribution of \(Z=X+Y\).

\begin{tcolorbox}[enhanced jigsaw, arc=.35mm, bottomrule=.15mm, bottomtitle=1mm, breakable, colback=white, colbacktitle=quarto-callout-note-color!10!white, colframe=quarto-callout-note-color-frame, coltitle=black, left=2mm, leftrule=.75mm, opacityback=0, opacitybacktitle=0.6, rightrule=.15mm, title=\textcolor{quarto-callout-note-color}{\faInfo}\hspace{0.5em}{Solution}, titlerule=0mm, toprule=.15mm, toptitle=1mm]

By the product rule for MGFs, \[\begin{aligned}
M_Z(t)&=M_X(t)M_Y(t)\\
&=\exp\left(\mu_1t+\frac{\sigma_1^2t^2}{2}\right)
\exp\left(\mu_2t+\frac{\sigma_2^2t^2}{2}\right)\\
&=\exp\left((\mu_1+\mu_2)t+\frac{(\sigma_1^2+\sigma_2^2)t^2}{2}\right).
\end{aligned}\] This is the MGF of a normal random variable with mean
\(\mu_1+\mu_2\) and variance \(\sigma_1^2+\sigma_2^2\). Therefore
\[X+Y\sim\operatorname{Normal}(\mu_1+\mu_2,\sigma_1^2+\sigma_2^2).\]

\end{tcolorbox}

\section{Multivariate Moment Generating
Functions}\label{multivariate-moment-generating-functions}

This section extends the MGF idea from one random variable to random
vectors.

\subsection{Definition for random
vectors}\label{definition-for-random-vectors}

This subsection defines the multivariate MGF and explains its role in
describing joint distributions.

\textbf{Definition 41} (Multivariate MGF). Let \(X\in\mathbb{R}^d\) be a
random vector and let \(t\in\mathbb{R}^d\). The multivariate MGF of
\(X\) is \[M_X(t)=\mathbb{E}\left[e^{t^T X}\right],\] for values of
\(t\) where the expectation exists.

The multivariate MGF contains joint moment information. For example, in
two dimensions, \[M_{X,Y}(s,t)=\mathbb{E}[e^{sX+tY}],\] and mixed
derivatives generate mixed moments such as \(\mathbb{E}[X^aY^b]\).

\subsection{Multivariate normal MGF}\label{multivariate-normal-mgf}

This subsection records the fundamental MGF formula for the multivariate
normal distribution.

\textbf{Theorem 42} (Multivariate normal MGF). \emph{If
\(X\sim\operatorname{Normal}(\mu,\Sigma)\) in \(\mathbb{R}^d\), then
\[M_X(t)=\exp\left(t^T\mu+\frac12 t^T\Sigma t\right).\]}

\textbf{Example 43} (Linear transformation of a multivariate normal).
Let \(X\sim\operatorname{Normal}(\mu,\Sigma)\) and define \[Z=AX+b,\]
where \(A\) is a matrix and \(b\) is a vector. Find the distribution of
\(Z\).

\begin{tcolorbox}[enhanced jigsaw, arc=.35mm, bottomrule=.15mm, bottomtitle=1mm, breakable, colback=white, colbacktitle=quarto-callout-note-color!10!white, colframe=quarto-callout-note-color-frame, coltitle=black, left=2mm, leftrule=.75mm, opacityback=0, opacitybacktitle=0.6, rightrule=.15mm, title=\textcolor{quarto-callout-note-color}{\faInfo}\hspace{0.5em}{Solution}, titlerule=0mm, toprule=.15mm, toptitle=1mm]

The MGF of \(Z\) is \[\begin{aligned}
M_Z(t)&=\mathbb{E}[e^{t^T(AX+b)}]
=e^{t^Tb}\mathbb{E}[e^{(A^Tt)^TX}]\\
&=e^{t^Tb}M_X(A^Tt)\\
&=\exp\left(t^Tb+(A^Tt)^T\mu+\frac12(A^Tt)^T\Sigma(A^Tt)\right)\\
&=\exp\left(t^T(b+A\mu)+\frac12 t^T(A\Sigma A^T)t\right).
\end{aligned}\] This is the MGF of a multivariate normal distribution
with mean \(b+A\mu\) and covariance matrix \(A\Sigma A^T\). Therefore
\[Z\sim\operatorname{Normal}(b+A\mu,A\Sigma A^T).\]

\end{tcolorbox}

\subsection{Sum of multivariate normal random
vectors}\label{sum-of-multivariate-normal-random-vectors}

This subsection presents the general normal-sum formula, including the
covariance between the two vectors.

\textbf{Example 44} (Sum of multivariate normal distributions). Let
\[X\sim\operatorname{Normal}(\mu_1,\Sigma_1),\qquad
Y\sim\operatorname{Normal}(\mu_2,\Sigma_2),\] and suppose the joint
covariance matrix of \((X,Y)\) is \[\Sigma=\begin{pmatrix}
\Sigma_1 & \Sigma_{12}\\
\Sigma_{21} & \Sigma_2
\end{pmatrix},
\qquad \Sigma_{21}=\operatorname{Cov}(Y,X).\] Find the distribution of
\(Z=X+Y\) when the joint distribution is multivariate normal.

\begin{tcolorbox}[enhanced jigsaw, arc=.35mm, bottomrule=.15mm, bottomtitle=1mm, breakable, colback=white, colbacktitle=quarto-callout-note-color!10!white, colframe=quarto-callout-note-color-frame, coltitle=black, left=2mm, leftrule=.75mm, opacityback=0, opacitybacktitle=0.6, rightrule=.15mm, title=\textcolor{quarto-callout-note-color}{\faInfo}\hspace{0.5em}{Solution}, titlerule=0mm, toprule=.15mm, toptitle=1mm]

Since \(Z=X+Y\) is a linear transformation of the jointly normal vector
\((X,Y)\), it is normal. Its mean is
\[\mathbb{E}[Z]=\mathbb{E}[X]+\mathbb{E}[Y]=\mu_1+
\mu_2.\] Its covariance matrix is \[\begin{aligned}
\operatorname{Cov}(Z)&=\operatorname{Cov}(X+Y)\\
&=\operatorname{Cov}(X)+\operatorname{Cov}(Y)+\operatorname{Cov}(X,Y)+\operatorname{Cov}(Y,X)\\
&=\Sigma_1+\Sigma_2+\Sigma_{12}+\Sigma_{21}.
\end{aligned}\] If \(\Sigma_{12}=\Sigma_{21}^T\), this is often written
as \[\Sigma_1+\Sigma_2+\Sigma_{12}+\Sigma_{12}^T.\] In the scalar case,
this becomes
\[\operatorname{Var}(X+Y)=\operatorname{Var}(X)+\operatorname{Var}(Y)+2\operatorname{Cov}(X,Y).\]
If \(X\) and \(Y\) are independent, then \(\Sigma_{12}=\Sigma_{21}=0\),
so \[X+Y\sim\operatorname{Normal}(\mu_1+\mu_2,\Sigma_1+
\Sigma_2).\]

\end{tcolorbox}

\section{Practice Problems}\label{practice-problems-1}

This section gives additional practice problems that reinforce the main
computational skills from the lecture.

\begin{tcolorbox}[enhanced jigsaw, arc=.35mm, bottomrule=.15mm, bottomtitle=1mm, breakable, colback=white, colbacktitle=quarto-callout-important-color!10!white, colframe=quarto-callout-important-color-frame, coltitle=black, left=2mm, leftrule=.75mm, opacityback=0, opacitybacktitle=0.6, rightrule=.15mm, title=\textcolor{quarto-callout-important-color}{\faExclamation}\hspace{0.5em}{Practice Problem}, titlerule=0mm, toprule=.15mm, toptitle=1mm]

\textbf{Practice Problem 45} (Expectation and variance). Let \(X\) have
pmf \[\mathbb{P}(X=-1)=\frac14,
\qquad
\mathbb{P}(X=0)=\frac12,
\qquad
\mathbb{P}(X=2)=\frac14.\] Find \(\mathbb{E}[X]\), \(\mathbb{E}[X^2]\),
and \(\operatorname{Var}(X)\).

\end{tcolorbox}

\begin{tcolorbox}[enhanced jigsaw, arc=.35mm, bottomrule=.15mm, bottomtitle=1mm, breakable, colback=white, colbacktitle=quarto-callout-note-color!10!white, colframe=quarto-callout-note-color-frame, coltitle=black, left=2mm, leftrule=.75mm, opacityback=0, opacitybacktitle=0.6, rightrule=.15mm, title=\textcolor{quarto-callout-note-color}{\faInfo}\hspace{0.5em}{Solution}, titlerule=0mm, toprule=.15mm, toptitle=1mm]

\[\mathbb{E}[X]=(-1)\frac14+0\frac12+2\frac14=-\frac14+\frac12=\frac14.\]
Also,
\[\mathbb{E}[X^2]=(-1)^2\frac14+0^2\frac12+2^2\frac14=\frac14+1=\frac54.\]
Therefore
\[\operatorname{Var}(X)=\mathbb{E}[X^2]-(\mathbb{E}[X])^2=\frac54-\left(\frac14\right)^2=\frac{20}{16}-\frac{1}{16}=\frac{19}{16}.\]

\end{tcolorbox}

\begin{tcolorbox}[enhanced jigsaw, arc=.35mm, bottomrule=.15mm, bottomtitle=1mm, breakable, colback=white, colbacktitle=quarto-callout-important-color!10!white, colframe=quarto-callout-important-color-frame, coltitle=black, left=2mm, leftrule=.75mm, opacityback=0, opacitybacktitle=0.6, rightrule=.15mm, title=\textcolor{quarto-callout-important-color}{\faExclamation}\hspace{0.5em}{Practice Problem}, titlerule=0mm, toprule=.15mm, toptitle=1mm]

\textbf{Practice Problem 46} (Expectation of a function). Let
\(X\sim \operatorname{Exponential}(\lambda)\). Compute
\(\mathbb{E}[e^{-sX}]\) for \(s>-\lambda\).

\end{tcolorbox}

\begin{tcolorbox}[enhanced jigsaw, arc=.35mm, bottomrule=.15mm, bottomtitle=1mm, breakable, colback=white, colbacktitle=quarto-callout-note-color!10!white, colframe=quarto-callout-note-color-frame, coltitle=black, left=2mm, leftrule=.75mm, opacityback=0, opacitybacktitle=0.6, rightrule=.15mm, title=\textcolor{quarto-callout-note-color}{\faInfo}\hspace{0.5em}{Solution}, titlerule=0mm, toprule=.15mm, toptitle=1mm]

Using the expectation of a function formula,
\[\mathbb{E}[e^{-sX}]=\int_0^\infty e^{-sx}\lambda e^{-\lambda x}\,dx
=\lambda\int_0^\infty e^{-(\lambda+s)x}\,dx
=\frac{\lambda}{\lambda+s}.\] This is the Laplace transform of the
exponential distribution.

\end{tcolorbox}

\begin{tcolorbox}[enhanced jigsaw, arc=.35mm, bottomrule=.15mm, bottomtitle=1mm, breakable, colback=white, colbacktitle=quarto-callout-important-color!10!white, colframe=quarto-callout-important-color-frame, coltitle=black, left=2mm, leftrule=.75mm, opacityback=0, opacitybacktitle=0.6, rightrule=.15mm, title=\textcolor{quarto-callout-important-color}{\faExclamation}\hspace{0.5em}{Practice Problem}, titlerule=0mm, toprule=.15mm, toptitle=1mm]

\textbf{Practice Problem 47} (MGF and distribution identification).
Suppose \(X\) has MGF
\[M_X(t)=\left(\frac{1}{3}+\frac{2}{3}e^t\right)^5.\] Identify the
distribution of \(X\).

\end{tcolorbox}

\begin{tcolorbox}[enhanced jigsaw, arc=.35mm, bottomrule=.15mm, bottomtitle=1mm, breakable, colback=white, colbacktitle=quarto-callout-note-color!10!white, colframe=quarto-callout-note-color-frame, coltitle=black, left=2mm, leftrule=.75mm, opacityback=0, opacitybacktitle=0.6, rightrule=.15mm, title=\textcolor{quarto-callout-note-color}{\faInfo}\hspace{0.5em}{Solution}, titlerule=0mm, toprule=.15mm, toptitle=1mm]

The binomial MGF is \[M(t)=(1-p+pe^t)^n.\] Here \(n=5\) and \(p=2/3\).
Therefore \[X\sim\operatorname{Binomial}\left(5,\frac23\right).\]

\end{tcolorbox}

\begin{tcolorbox}[enhanced jigsaw, arc=.35mm, bottomrule=.15mm, bottomtitle=1mm, breakable, colback=white, colbacktitle=quarto-callout-important-color!10!white, colframe=quarto-callout-important-color-frame, coltitle=black, left=2mm, leftrule=.75mm, opacityback=0, opacitybacktitle=0.6, rightrule=.15mm, title=\textcolor{quarto-callout-important-color}{\faExclamation}\hspace{0.5em}{Practice Problem}, titlerule=0mm, toprule=.15mm, toptitle=1mm]

\textbf{Practice Problem 48} (Independent sum). Let \(X_1,\ldots,X_n\)
be independent exponential random variables with rate \(\lambda\). Write
the MGF of \(S_n=X_1+\cdots+X_n\).

\end{tcolorbox}

\begin{tcolorbox}[enhanced jigsaw, arc=.35mm, bottomrule=.15mm, bottomtitle=1mm, breakable, colback=white, colbacktitle=quarto-callout-note-color!10!white, colframe=quarto-callout-note-color-frame, coltitle=black, left=2mm, leftrule=.75mm, opacityback=0, opacitybacktitle=0.6, rightrule=.15mm, title=\textcolor{quarto-callout-note-color}{\faInfo}\hspace{0.5em}{Solution}, titlerule=0mm, toprule=.15mm, toptitle=1mm]

The MGF of each \(X_i\) is
\[M_{X_i}(t)=\frac{\lambda}{\lambda-t},\qquad t<\lambda.\] By
independence,
\[M_{S_n}(t)=\prod_{i=1}^nM_{X_i}(t)=\left(\frac{\lambda}{\lambda-t}\right)^n.\]
This is the MGF of a Gamma/Erlang distribution with shape \(n\) and rate
\(\lambda\).

\end{tcolorbox}

\section{Summary}\label{summary-2}

This section summarizes the most important ideas and formulas from the
lecture.

\begin{tcolorbox}[enhanced jigsaw, arc=.35mm, bottomrule=.15mm, bottomtitle=1mm, breakable, colback=white, colbacktitle=quarto-callout-tip-color!10!white, colframe=quarto-callout-tip-color-frame, coltitle=black, left=2mm, leftrule=.75mm, opacityback=0, opacitybacktitle=0.6, rightrule=.15mm, title=\textcolor{quarto-callout-tip-color}{\faLightbulb}\hspace{0.5em}{Core formulas}, titlerule=0mm, toprule=.15mm, toptitle=1mm]

\[\begin{aligned}
\mathbb{E}[X]&=\sum_x xp_X(x) \quad \text{or}\quad \mathbb{E}[X]=\int xf_X(x)\,dx,\\
\operatorname{Var}(X)&=\mathbb{E}[(X-\mathbb{E}[X])^2]=\mathbb{E}[X^2]-(\mathbb{E}[X])^2,\\
\mathbb{E}[g(X)]&=\sum_x g(x)p_X(x) \quad \text{or}\quad \mathbb{E}[g(X)]=\int g(x)f_X(x)\,dx,\\
M_X(t)&=\mathbb{E}[e^{tX}],\\
\mathbb{E}[X^m]&=M_X^{(m)}(0),\\
M_{aX+b}(t)&=e^{bt}M_X(at),\\
M_{X+Y}(t)&=M_X(t)M_Y(t)\quad \text{if }X,Y\text{ are independent.}
\end{aligned}\]

\end{tcolorbox}

\begin{tcolorbox}[enhanced jigsaw, arc=.35mm, bottomrule=.15mm, bottomtitle=1mm, breakable, colback=white, colbacktitle=quarto-callout-tip-color!10!white, colframe=quarto-callout-tip-color-frame, coltitle=black, left=2mm, leftrule=.75mm, opacityback=0, opacitybacktitle=0.6, rightrule=.15mm, title=\textcolor{quarto-callout-tip-color}{\faLightbulb}\hspace{0.5em}{Common MGFs}, titlerule=0mm, toprule=.15mm, toptitle=1mm]

\[\begin{aligned}
X\sim\operatorname{Bernoulli}(p):\quad &M_X(t)=1-p+pe^t,\\
X\sim\operatorname{Binomial}(n,p):\quad &M_X(t)=(1-p+pe^t)^n,\\
X\sim\operatorname{Poisson}(\lambda):\quad &M_X(t)=\exp\{\lambda(e^t-1)\},\\
X\sim\operatorname{Exponential}(\lambda):\quad &M_X(t)=\frac{\lambda}{\lambda-t},\quad t<\lambda,\\
X\sim\operatorname{Normal}(\mu,\sigma^2):\quad &M_X(t)=\exp\left(\mu t+\frac{\sigma^2t^2}{2}\right).
\end{aligned}\]

\end{tcolorbox}

\bookmarksetup{startatroot}

\chapter{Chapter 5: Transformations}\label{chapter-5-transformations}

This chapter explains how probability distributions change when random
variables are transformed by functions, sums, products, quotients,
maxima, minima, and coordinate changes.

\textbf{Topics.} Functions of random variables; CDF method;
one-dimensional change of variables; transformations of several random
variables; Jacobians; convolution; quotient, product, maximum, and
minimum distributions; Box--Muller method.

\section{Overview}\label{overview-2}

This section explains how to find the distribution of a new random
variable created from old random variables.

In applications, we rarely observe only the original random variable. We
often transform it: we square it, take logarithms, add several
variables, take a maximum or minimum, form a ratio, or change
coordinates. The central problem is:
\[\text{given the distribution of }X,\text{ find the distribution of }Y=g(X),\]
or more generally,
\[\text{given the joint distribution of }(X_1,\ldots,X_n),\text{ find the distribution of }Y=g(X_1,\ldots,X_n).\]

\begin{tcolorbox}[enhanced jigsaw, arc=.35mm, bottomrule=.15mm, bottomtitle=1mm, breakable, colback=white, colbacktitle=quarto-callout-tip-color!10!white, colframe=quarto-callout-tip-color-frame, coltitle=black, left=2mm, leftrule=.75mm, opacityback=0, opacitybacktitle=0.6, rightrule=.15mm, title=\textcolor{quarto-callout-tip-color}{\faLightbulb}\hspace{0.5em}{Main message}, titlerule=0mm, toprule=.15mm, toptitle=1mm]

A transformation changes probability by changing the set of outcomes
being counted or integrated. For one variable, the main tools are the
CDF method and the change-of-variables formula. For several variables,
the main tool is the Jacobian determinant.

\end{tcolorbox}

\section{Function of One Random
Variable}\label{function-of-one-random-variable}

This section begins with the basic transformation problem for a single
random variable.

Let \(X\) be a random variable and let \(g\) be a function. Define
\[Y=g(X).\] Our goal is to find the distribution of \(Y\). The key idea
is to translate events about \(Y\) into events about \(X\).

\subsection{Inverse images of events}\label{inverse-images-of-events}

This subsection introduces the set-level viewpoint that works for both
discrete and continuous random variables.

For any subset \(A\) of possible values of \(Y\),
\[\mathbb{P}(Y\in A)=\mathbb{P}(g(X)\in A)=\mathbb{P}(X\in g^{-1}(A)),\]
where \[g^{-1}(A)=\{x:g(x)\in A\}.\] Here \(g^{-1}\) does not
necessarily mean an ordinary inverse function. It means the inverse
image of a set.

\textbf{Example 1} (Set inverse image). Let \(Y=X^2\). If \(A=[0,4]\),
then \[g^{-1}(A)=\{x:x^2\in[0,4]\}=[-2,2].\] Thus
\[\mathbb{P}(Y\in[0,4])=\mathbb{P}(-2\leq X\leq 2).\]

\begin{tcolorbox}[enhanced jigsaw, arc=.35mm, bottomrule=.15mm, bottomtitle=1mm, breakable, colback=white, colbacktitle=quarto-callout-note-color!10!white, colframe=quarto-callout-note-color-frame, coltitle=black, left=2mm, leftrule=.75mm, opacityback=0, opacitybacktitle=0.6, rightrule=.15mm, title=\textcolor{quarto-callout-note-color}{\faInfo}\hspace{0.5em}{Solution}, titlerule=0mm, toprule=.15mm, toptitle=1mm]

The event \(Y\in[0,4]\) means \(X^2\in[0,4]\). This is equivalent to
\(|X|\leq 2\), or \(-2\leq X\leq 2\). Therefore the probability of the
event for \(Y\) is computed from the corresponding interval for \(X\).

\end{tcolorbox}

\subsection{Discrete transformations}\label{discrete-transformations}

This subsection gives the formula for transforming a discrete random
variable.

If \(X\) is discrete and \(Y=g(X)\), then the pmf of \(Y\) is
\[p_Y(y)=\mathbb{P}(Y=y)=\mathbb{P}(X\in g^{-1}(\{y\}))
      =\sum_{x:g(x)=y}\mathbb{P}(X=x)
      =\sum_{x:g(x)=y}p_X(x).\]

\textbf{Example 2} (A nonlinear discrete transformation). Let \(X\) be a
discrete random variable with
\[\mathbb{P}(X=x)=\frac15,\qquad x\in\{-2,-1,0,1,2\}.\] Let
\[Y=g(X)=2|X|.\] Find the range and pmf of \(Y\).

\begin{tcolorbox}[enhanced jigsaw, arc=.35mm, bottomrule=.15mm, bottomtitle=1mm, breakable, colback=white, colbacktitle=quarto-callout-note-color!10!white, colframe=quarto-callout-note-color-frame, coltitle=black, left=2mm, leftrule=.75mm, opacityback=0, opacitybacktitle=0.6, rightrule=.15mm, title=\textcolor{quarto-callout-note-color}{\faInfo}\hspace{0.5em}{Solution}, titlerule=0mm, toprule=.15mm, toptitle=1mm]

Compute \(Y=2|X|\) for each possible value of \(X\):
\[\begin{array}{c|ccccc}
 x & -2 & -1 & 0 & 1 & 2\\ \hline
 y=2|x| & 4 & 2 & 0 & 2 & 4
\end{array}\] Hence \[\operatorname{Range}(Y)=\{0,2,4\}.\] Now sum the
probabilities over all preimages: \[p_Y(0)=\mathbb{P}(X=0)=\frac15,\]
\[p_Y(2)=\mathbb{P}(X=-1)+\mathbb{P}(X=1)=\frac15+\frac15=\frac25,\]
\[p_Y(4)=\mathbb{P}(X=-2)+\mathbb{P}(X=2)=\frac15+\frac15=\frac25.\]
Therefore \[\begin{array}{c|ccc}
 y & 0 & 2 & 4\\ \hline
 p_Y(y) & 1/5 & 2/5 & 2/5
\end{array}\] The probabilities add to \(1\), so this is a valid pmf.

\end{tcolorbox}

\textbf{Practice Problem 3} (Practice: discrete transformation). Let
\(X\) be uniformly distributed on \(\{-2,-1,0,1,2\}\) and let
\(Y=X^2+1\). Find the pmf of \(Y\).

\begin{tcolorbox}[enhanced jigsaw, arc=.35mm, bottomrule=.15mm, bottomtitle=1mm, breakable, colback=white, colbacktitle=quarto-callout-note-color!10!white, colframe=quarto-callout-note-color-frame, coltitle=black, left=2mm, leftrule=.75mm, opacityback=0, opacitybacktitle=0.6, rightrule=.15mm, title=\textcolor{quarto-callout-note-color}{\faInfo}\hspace{0.5em}{Solution}, titlerule=0mm, toprule=.15mm, toptitle=1mm]

The possible values are \[Y=X^2+1\in\{1,2,5\}.\] The preimages are
\[Y=1 \Longleftrightarrow X=0,\] \[Y=2 \Longleftrightarrow X=\pm 1,\]
\[Y=5 \Longleftrightarrow X=\pm 2.\] Since each value of \(X\) has
probability \(1/5\),
\[p_Y(1)=\frac15,\qquad p_Y(2)=\frac25,\qquad p_Y(5)=\frac25.\]

\end{tcolorbox}

\subsection{Continuous transformations by the CDF
method}\label{continuous-transformations-by-the-cdf-method}

This subsection gives the most reliable method for continuous
transformations: first find the CDF, then differentiate.

If \(X\) is continuous with pdf \(f_X\), then the CDF of \(Y=g(X)\) is
\[F_Y(y)=\mathbb{P}(Y\leq y)=\mathbb{P}(g(X)\leq y)
      =\int_{\{x:g(x)\leq y\}} f_X(x)\,dx.\] If \(F_Y\) is
differentiable, then \[f_Y(y)=F_Y'(y).\]

\begin{tcolorbox}[enhanced jigsaw, arc=.35mm, bottomrule=.15mm, bottomtitle=1mm, breakable, colback=white, colbacktitle=quarto-callout-tip-color!10!white, colframe=quarto-callout-tip-color-frame, coltitle=black, left=2mm, leftrule=.75mm, opacityback=0, opacitybacktitle=0.6, rightrule=.15mm, title=\textcolor{quarto-callout-tip-color}{\faLightbulb}\hspace{0.5em}{CDF method}, titlerule=0mm, toprule=.15mm, toptitle=1mm]

To find the pdf of \(Y=g(X)\):

\begin{enumerate}
\def\labelenumi{\arabic{enumi}.}
\item
  Write \(F_Y(y)=\mathbb{P}(g(X)\leq y)\).
\item
  Convert the event into an event involving \(X\).
\item
  Compute the probability using \(F_X\) or \(f_X\).
\item
  Differentiate with respect to \(y\).
\end{enumerate}

\end{tcolorbox}

\subsection{Monotone change of
variables}\label{monotone-change-of-variables}

This subsection gives the standard formula when the transformation is
one-to-one.

\textbf{Theorem 4} (Increasing transformation). \emph{Suppose \(X\) has
pdf \(f_X\) on an interval \((a,b)\). Let \(Y=g(X)\), where \(g\) is
differentiable and strictly increasing on \((a,b)\) with \(g'(x)>0\).
Then \[f_Y(y)=
\begin{cases}
\displaystyle \frac{f_X(g^{-1}(y))}{g'(g^{-1}(y))}, & g(a)<y<g(b),\\[1.2em]
0, & \text{otherwise.}
\end{cases}\]}

\begin{tcolorbox}[enhanced jigsaw, arc=.35mm, bottomrule=.15mm, bottomtitle=1mm, breakable, colback=white, colbacktitle=quarto-callout-note-color!10!white, colframe=quarto-callout-note-color-frame, coltitle=black, left=2mm, leftrule=.75mm, opacityback=0, opacitybacktitle=0.6, rightrule=.15mm, title=\textcolor{quarto-callout-note-color}{\faInfo}\hspace{0.5em}{Proof}, titlerule=0mm, toprule=.15mm, toptitle=1mm]

Because \(g\) is increasing,
\[F_Y(y)=\mathbb{P}(g(X)\leq y)=\mathbb{P}(X\leq g^{-1}(y))=F_X(g^{-1}(y)).\]
Differentiate with respect to \(y\):
\[f_Y(y)=f_X(g^{-1}(y))\frac{\,d}{\,dy}g^{-1}(y)
      =\frac{f_X(g^{-1}(y))}{g'(g^{-1}(y))}.\]

\end{tcolorbox}

\textbf{Theorem 5} (General one-to-one transformation). \emph{If \(g\)
is differentiable and strictly monotone, then
\[f_Y(y)=f_X(g^{-1}(y))\left|\frac{\,d}{\,dy}g^{-1}(y)\right|
      =\frac{f_X(g^{-1}(y))}{|g'(g^{-1}(y))|}.\]}

\begin{tcolorbox}[enhanced jigsaw, arc=.35mm, bottomrule=.15mm, bottomtitle=1mm, breakable, colback=white, colbacktitle=quarto-callout-warning-color!10!white, colframe=quarto-callout-warning-color-frame, coltitle=black, left=2mm, leftrule=.75mm, opacityback=0, opacitybacktitle=0.6, rightrule=.15mm, title=\textcolor{quarto-callout-warning-color}{\faExclamationTriangle}\hspace{0.5em}{Common mistake}, titlerule=0mm, toprule=.15mm, toptitle=1mm]

For decreasing transformations, use an absolute value. Densities must be
nonnegative.

\end{tcolorbox}

\subsection{Linear transformations}\label{linear-transformations}

This subsection records the simplest and most frequently used
transformation.

\textbf{Example 6} (Linear transformation of a uniform random variable).
Let \(X\sim \operatorname{Uniform}(0,1)\) and let \[Y=aX+b,\] where
\(a>0\). Find the distribution of \(Y\).

\begin{tcolorbox}[enhanced jigsaw, arc=.35mm, bottomrule=.15mm, bottomtitle=1mm, breakable, colback=white, colbacktitle=quarto-callout-note-color!10!white, colframe=quarto-callout-note-color-frame, coltitle=black, left=2mm, leftrule=.75mm, opacityback=0, opacitybacktitle=0.6, rightrule=.15mm, title=\textcolor{quarto-callout-note-color}{\faInfo}\hspace{0.5em}{Solution}, titlerule=0mm, toprule=.15mm, toptitle=1mm]

Since \(0\leq X\leq 1\), the range of \(Y\) is \[b\leq Y\leq a+b.\] The
inverse transformation is \[x=\frac{y-b}{a},\] and
\[\left|\frac{\,dx}{\,dy}\right|=\frac1a.\] Since \(f_X(x)=1\) on
\([0,1]\), \[f_Y(y)=f_X\left(\frac{y-b}{a}\right)\frac1a=\frac1a,
\qquad b\leq y\leq a+b.\] Thus \[Y\sim \operatorname{Uniform}(b,a+b).\]
If \(a<0\), the interval endpoints should be reordered, and
\(Y\sim \operatorname{Uniform}(a+b,b)\).

\end{tcolorbox}

\textbf{Example 7} (Linear transformation of a normal random variable).
Let \(X\sim \operatorname{Normal}(\mu,\sigma^2)\) and let \(Y=aX+b\).
Find the mean, variance, and distribution of \(Y\).

\begin{tcolorbox}[enhanced jigsaw, arc=.35mm, bottomrule=.15mm, bottomtitle=1mm, breakable, colback=white, colbacktitle=quarto-callout-note-color!10!white, colframe=quarto-callout-note-color-frame, coltitle=black, left=2mm, leftrule=.75mm, opacityback=0, opacitybacktitle=0.6, rightrule=.15mm, title=\textcolor{quarto-callout-note-color}{\faInfo}\hspace{0.5em}{Solution}, titlerule=0mm, toprule=.15mm, toptitle=1mm]

The expectation and variance are
\[\mathbb{E}[Y]=\mathbb{E}[aX+b]=a\mu+b,\]
\[\operatorname{Var}(Y)=\operatorname{Var}(aX+b)=a^2\sigma^2.\] A linear
transformation of a normal random variable is still normal, so
\[Y\sim \operatorname{Normal}(a\mu+b,a^2\sigma^2).\] If \(a=0\), then
\(Y=b\) is degenerate at \(b\).

\end{tcolorbox}

\textbf{Example 8} (General linear density formula). Let \(X\) have pdf
\(f_X(x)\) and let \(Y=aX+b\), where \(a\ne0\). Find \(f_Y\).

\begin{tcolorbox}[enhanced jigsaw, arc=.35mm, bottomrule=.15mm, bottomtitle=1mm, breakable, colback=white, colbacktitle=quarto-callout-note-color!10!white, colframe=quarto-callout-note-color-frame, coltitle=black, left=2mm, leftrule=.75mm, opacityback=0, opacitybacktitle=0.6, rightrule=.15mm, title=\textcolor{quarto-callout-note-color}{\faInfo}\hspace{0.5em}{Solution}, titlerule=0mm, toprule=.15mm, toptitle=1mm]

The inverse is \[x=\frac{y-b}{a},\] and
\[\left|\frac{\,dx}{\,dy}\right|=\frac1{|a|}.\] Therefore
\[f_Y(y)=\frac1{|a|}f_X\left(\frac{y-b}{a}\right).\] Also,
\[\mathbb{E}[Y]=a\mathbb{E}[X]+b,
\qquad
\operatorname{Var}(Y)=a^2\operatorname{Var}(X).\]

\end{tcolorbox}

\subsection{Nonlinear one-dimensional
examples}\label{nonlinear-one-dimensional-examples}

This subsection develops several important nonlinear transformations:
square, logarithm, chi-square, and inverse transform sampling.

\textbf{Example 9} (Quadratic transformation). Let
\(X\sim \operatorname{Uniform}(0,1)\) and let \[Y=X^2.\] Find the pdf of
\(Y\).

\begin{tcolorbox}[enhanced jigsaw, arc=.35mm, bottomrule=.15mm, bottomtitle=1mm, breakable, colback=white, colbacktitle=quarto-callout-note-color!10!white, colframe=quarto-callout-note-color-frame, coltitle=black, left=2mm, leftrule=.75mm, opacityback=0, opacitybacktitle=0.6, rightrule=.15mm, title=\textcolor{quarto-callout-note-color}{\faInfo}\hspace{0.5em}{Solution}, titlerule=0mm, toprule=.15mm, toptitle=1mm]

The function \(g(x)=x^2\) is increasing on \([0,1]\). Its inverse is
\[g^{-1}(y)=\sqrt{y},\] and \[g'(x)=2x.\] Thus, for \(0<y<1\),
\[f_Y(y)=\frac{f_X(\sqrt{y})}{2\sqrt{y}}=\frac1{2\sqrt{y}}.\] Therefore
\[f_Y(y)=\frac1{2\sqrt{y}}\mathbf{1}_{(0,1)}(y).\] Check normalization:
\[\int_0^1 \frac1{2\sqrt{y}}\,dy=\left[\sqrt{y}\right]_0^1=1.\]

\end{tcolorbox}

\textbf{Example 10} (Logarithmic transformation). Let
\(X\sim \operatorname{Uniform}(0,1)\) and define \[Y=-2\log X.\] Find
the pdf of \(Y\) and identify its distribution.

\begin{tcolorbox}[enhanced jigsaw, arc=.35mm, bottomrule=.15mm, bottomtitle=1mm, breakable, colback=white, colbacktitle=quarto-callout-note-color!10!white, colframe=quarto-callout-note-color-frame, coltitle=black, left=2mm, leftrule=.75mm, opacityback=0, opacitybacktitle=0.6, rightrule=.15mm, title=\textcolor{quarto-callout-note-color}{\faInfo}\hspace{0.5em}{Solution}, titlerule=0mm, toprule=.15mm, toptitle=1mm]

The transformation \(g(x)=-2\log x\) is decreasing on \((0,1)\). Its
inverse is \[x=e^{-y/2},\qquad y>0.\] Then
\[\left|\frac{\,dx}{\,dy}\right|=\frac12 e^{-y/2}.\] Since \(f_X(x)=1\)
on \((0,1)\), \[f_Y(y)=\frac12 e^{-y/2},\qquad y>0.\] Hence
\[Y\sim \operatorname{Exp}\left(\lambda=\frac12\right),\] where the
exponential pdf is \(\lambda e^{-\lambda y}\) for \(y\geq0\).

\end{tcolorbox}

\textbf{Example 11} (Chi-square distribution from a standard normal).
Let \(X\sim \operatorname{Normal}(0,1)\) and define \[Y=X^2.\] Find the
pdf of \(Y\).

\begin{tcolorbox}[enhanced jigsaw, arc=.35mm, bottomrule=.15mm, bottomtitle=1mm, breakable, colback=white, colbacktitle=quarto-callout-note-color!10!white, colframe=quarto-callout-note-color-frame, coltitle=black, left=2mm, leftrule=.75mm, opacityback=0, opacitybacktitle=0.6, rightrule=.15mm, title=\textcolor{quarto-callout-note-color}{\faInfo}\hspace{0.5em}{Solution}, titlerule=0mm, toprule=.15mm, toptitle=1mm]

The transformation \(x\mapsto x^2\) is not one-to-one on \(\mathbb{R}\),
so use the CDF method. For \(y\geq0\),
\[F_Y(y)=\mathbb{P}(Y\leq y)=\mathbb{P}(X^2\leq y)
       =\mathbb{P}(-\sqrt{y}\leq X\leq \sqrt{y})
       =F_X(\sqrt{y})-F_X(-\sqrt{y}).\] Differentiate:
\[f_Y(y)=f_X(\sqrt{y})\frac1{2\sqrt{y}}+f_X(-\sqrt{y})\frac1{2\sqrt{y}}.\]
Because the standard normal pdf is symmetric,
\[f_X(\sqrt{y})=f_X(-\sqrt{y})=\frac1{\sqrt{2\pi}}e^{-y/2}.\] Therefore
\[f_Y(y)=\frac1{\sqrt{2\pi y}}e^{-y/2},\qquad y>0.\] This is the
chi-square distribution with \(1\) degree of freedom:
\[Y\sim \chi_1^2.\] Equivalently,
\[Y\sim \operatorname{Gamma}\left(\alpha=\frac12,\theta=2\right),\]
using the shape-scale parameterization.

\end{tcolorbox}

\textbf{Example 12} (Inverse-transform sampling). Suppose
\(U\sim \operatorname{Uniform}(0,1)\) and suppose \(Y\) has CDF
\(H(y)\). If \(H\) is continuous and strictly increasing, show that
\[Y=H^{-1}(U)\] has CDF \(H\).

\begin{tcolorbox}[enhanced jigsaw, arc=.35mm, bottomrule=.15mm, bottomtitle=1mm, breakable, colback=white, colbacktitle=quarto-callout-note-color!10!white, colframe=quarto-callout-note-color-frame, coltitle=black, left=2mm, leftrule=.75mm, opacityback=0, opacitybacktitle=0.6, rightrule=.15mm, title=\textcolor{quarto-callout-note-color}{\faInfo}\hspace{0.5em}{Solution}, titlerule=0mm, toprule=.15mm, toptitle=1mm]

For any \(y\), \[\mathbb{P}(Y\leq y)=\mathbb{P}(H^{-1}(U)\leq y).\]
Since \(H\) is increasing,
\[H^{-1}(U)\leq y \Longleftrightarrow U\leq H(y).\] Thus
\[\mathbb{P}(Y\leq y)=\mathbb{P}(U\leq H(y))=H(y),\] because
\(U\sim\operatorname{Uniform}(0,1)\). Therefore \(Y\) has CDF \(H\).

\end{tcolorbox}

\textbf{Practice Problem 13} (Practice: inverse transform for
exponential). Use inverse-transform sampling to generate an exponential
random variable with rate \(\lambda\) from
\(U\sim \operatorname{Uniform}(0,1)\).

\begin{tcolorbox}[enhanced jigsaw, arc=.35mm, bottomrule=.15mm, bottomtitle=1mm, breakable, colback=white, colbacktitle=quarto-callout-note-color!10!white, colframe=quarto-callout-note-color-frame, coltitle=black, left=2mm, leftrule=.75mm, opacityback=0, opacitybacktitle=0.6, rightrule=.15mm, title=\textcolor{quarto-callout-note-color}{\faInfo}\hspace{0.5em}{Solution}, titlerule=0mm, toprule=.15mm, toptitle=1mm]

The exponential CDF is \[F(y)=1-e^{-\lambda y},\qquad y\geq0.\] Set
\(U=F(y)\) and solve for \(y\): \[U=1-e^{-\lambda y}
\quad\Longrightarrow\quad
1-U=e^{-\lambda y}
\quad\Longrightarrow\quad
y=-\frac1\lambda\log(1-U).\] Since
\(1-U\sim \operatorname{Uniform}(0,1)\), we often write
\[Y=-\frac1\lambda\log U.\] Then \(Y\sim \operatorname{Exp}(\lambda)\).

\end{tcolorbox}

\section{Functions of Two or More Random
Variables}\label{functions-of-two-or-more-random-variables}

This section studies transformations involving several random variables,
such as sums, maxima, minima, and sample means.

Let \(X_1,\ldots,X_n\) be random variables and let
\[Y=g(X_1,\ldots,X_n).\] As before, the basic method is to compute the
CDF \[F_Y(y)=\mathbb{P}(g(X_1,\ldots,X_n)\leq y),\] and then
differentiate when a pdf exists.

\subsection{A motivating example: binomial as a sum of Bernoulli
variables}\label{a-motivating-example-binomial-as-a-sum-of-bernoulli-variables}

This subsection shows that many familiar distributions arise as
transformations of several random variables.

\textbf{Example 14} (Binomial distribution as a sum). Let
\[X_1,\ldots,X_n\stackrel{\text{iid}}{\sim}\operatorname{Bernoulli}(p),\]
and define \[X=X_1+\cdots+X_n.\] Find the distribution of \(X\).

\begin{tcolorbox}[enhanced jigsaw, arc=.35mm, bottomrule=.15mm, bottomtitle=1mm, breakable, colback=white, colbacktitle=quarto-callout-note-color!10!white, colframe=quarto-callout-note-color-frame, coltitle=black, left=2mm, leftrule=.75mm, opacityback=0, opacitybacktitle=0.6, rightrule=.15mm, title=\textcolor{quarto-callout-note-color}{\faInfo}\hspace{0.5em}{Solution}, titlerule=0mm, toprule=.15mm, toptitle=1mm]

The random variable \(X\) counts the number of successes in \(n\)
independent Bernoulli trials. Therefore
\[X\sim \operatorname{Binomial}(n,p),\] and
\[\mathbb{P}(X=k)=\binom{n}{k}p^k(1-p)^{n-k},
\qquad k=0,1,\ldots,n.\] The transformation is a sum:
\[(X_1,\ldots,X_n)\mapsto X_1+\cdots+X_n.\]

\end{tcolorbox}

\subsection{Two independent uniforms on the unit
square}\label{two-independent-uniforms-on-the-unit-square}

This subsection uses geometry to find distributions of sums, maxima, and
minima.

Let \((X,Y)\) be uniformly distributed on the unit square
\([0,1]\times[0,1]\). Equivalently, \(X\) and \(Y\) are independent
\(\operatorname{Uniform}(0,1)\) random variables.

\textbf{Example 15} (Sum of two independent uniforms). Let \[U=X+Y.\]
Find the CDF and pdf of \(U\).

\begin{tcolorbox}[enhanced jigsaw, arc=.35mm, bottomrule=.15mm, bottomtitle=1mm, breakable, colback=white, colbacktitle=quarto-callout-note-color!10!white, colframe=quarto-callout-note-color-frame, coltitle=black, left=2mm, leftrule=.75mm, opacityback=0, opacitybacktitle=0.6, rightrule=.15mm, title=\textcolor{quarto-callout-note-color}{\faInfo}\hspace{0.5em}{Solution}, titlerule=0mm, toprule=.15mm, toptitle=1mm]

Because \((X,Y)\) is uniform on the unit square, probabilities are
areas.

For \(u\leq0\), \[F_U(u)=0.\] For \(0\leq u\leq1\), the event
\(X+Y\leq u\) is a right triangle with side length \(u\), so
\[F_U(u)=\frac{u^2}{2}.\] For \(1\leq u\leq2\), it is easier to subtract
the upper-right triangle where \(X+Y>u\). That triangle has side length
\(2-u\), so \[F_U(u)=1-\frac{(2-u)^2}{2}.\] For \(u\geq2\),
\[F_U(u)=1.\] Thus \[F_U(u)=
\begin{cases}
0, & u\leq0,\\
\dfrac{u^2}{2}, & 0\leq u\leq1,\\[0.8em]
1-\dfrac{(2-u)^2}{2}, & 1\leq u\leq2,\\[0.8em]
1, & u\geq2.
\end{cases}\] Differentiating gives \[f_U(u)=
\begin{cases}
u, & 0\leq u\leq1,\\
2-u, & 1\leq u\leq2,\\
0, & \text{otherwise.}
\end{cases}\] This is the triangular distribution on \([0,2]\).

\begin{figure}[H]

\centering{

\pandocbounded{\includegraphics[keepaspectratio]{chapters/chapter-05-transformations_files/figure-pdf/fig-ch5-sum-uniforms-output-1.pdf}}

}

\caption{\label{fig-ch5-sum-uniforms}Triangular PDF of U = X + Y for
independent Uniform(0, 1) variables.}

\end{figure}%

\end{tcolorbox}

\textbf{Example 16} (Maximum of two independent uniforms). Let
\[V=\max\{X,Y\}.\] Find the CDF and pdf of \(V\).

\begin{tcolorbox}[enhanced jigsaw, arc=.35mm, bottomrule=.15mm, bottomtitle=1mm, breakable, colback=white, colbacktitle=quarto-callout-note-color!10!white, colframe=quarto-callout-note-color-frame, coltitle=black, left=2mm, leftrule=.75mm, opacityback=0, opacitybacktitle=0.6, rightrule=.15mm, title=\textcolor{quarto-callout-note-color}{\faInfo}\hspace{0.5em}{Solution}, titlerule=0mm, toprule=.15mm, toptitle=1mm]

For \(0\leq v\leq1\),
\[F_V(v)=\mathbb{P}(V\leq v)=\mathbb{P}(X\leq v,Y\leq v).\] By
independence, \[F_V(v)=\mathbb{P}(X\leq v)\mathbb{P}(Y\leq v)=v^2.\]
Therefore \[f_V(v)=2v,
\qquad 0\leq v\leq1.\] So \[f_V(v)=2v\mathbf{1}_{[0,1]}(v).\]

\end{tcolorbox}

\textbf{Example 17} (Minimum of two independent uniforms). Let
\[W=\min\{X,Y\}.\] Find the CDF and pdf of \(W\).

\begin{tcolorbox}[enhanced jigsaw, arc=.35mm, bottomrule=.15mm, bottomtitle=1mm, breakable, colback=white, colbacktitle=quarto-callout-note-color!10!white, colframe=quarto-callout-note-color-frame, coltitle=black, left=2mm, leftrule=.75mm, opacityback=0, opacitybacktitle=0.6, rightrule=.15mm, title=\textcolor{quarto-callout-note-color}{\faInfo}\hspace{0.5em}{Solution}, titlerule=0mm, toprule=.15mm, toptitle=1mm]

It is convenient to use the complement. For \(0\leq w\leq1\),
\[\mathbb{P}(W>w)=\mathbb{P}(X>w,Y>w).\] By independence,
\[\mathbb{P}(W>w)=(1-w)^2.\] Hence \[F_W(w)=1-(1-w)^2,
\qquad 0\leq w\leq1.\] Differentiate: \[f_W(w)=2(1-w),
\qquad 0\leq w\leq1.\] Thus \[f_W(w)=2(1-w)\mathbf{1}_{[0,1]}(w).\]

\end{tcolorbox}

\subsection{Minimum of many uniforms and an exponential
limit}\label{minimum-of-many-uniforms-and-an-exponential-limit}

This subsection shows how transformations can reveal limiting
distributions.

\textbf{Example 18} (Minimum of several uniform distributions). Let
\[X_1,\ldots,X_n\stackrel{\text{iid}}{\sim}\operatorname{Uniform}(0,1),\]
and define \[W_n=n\min\{X_1,\ldots,X_n\}.\] Show that \(W_n\) converges
in distribution to an \(\operatorname{Exp}(1)\) random variable.

\begin{tcolorbox}[enhanced jigsaw, arc=.35mm, bottomrule=.15mm, bottomtitle=1mm, breakable, colback=white, colbacktitle=quarto-callout-note-color!10!white, colframe=quarto-callout-note-color-frame, coltitle=black, left=2mm, leftrule=.75mm, opacityback=0, opacitybacktitle=0.6, rightrule=.15mm, title=\textcolor{quarto-callout-note-color}{\faInfo}\hspace{0.5em}{Solution}, titlerule=0mm, toprule=.15mm, toptitle=1mm]

For \(w\geq0\),
\[\mathbb{P}(W_n>w)=\mathbb{P}\left(\min\{X_1,\ldots,X_n\}>\frac{w}{n}\right).\]
For \(0\leq w\leq n\),
\[\mathbb{P}(W_n>w)=\mathbb{P}\left(X_1>\frac{w}{n},\ldots,X_n>\frac{w}{n}\right).\]
Using independence, \[\mathbb{P}(W_n>w)=\left(1-\frac{w}{n}\right)^n.\]
As \(n\to\infty\), \[\left(1-\frac{w}{n}\right)^n\to e^{-w}.\] Thus
\[\mathbb{P}(W_n\leq w)\to 1-e^{-w},\qquad w\geq0,\] which is the CDF of
\(\operatorname{Exp}(1)\). Therefore
\[W_n\Rightarrow \operatorname{Exp}(1).\]

\end{tcolorbox}

\subsection{Normal transformations}\label{normal-transformations}

This subsection records three important transformations involving normal
random variables.

\textbf{Example 19} (Sum of independent normal random variables).
Suppose \[X\sim \operatorname{Normal}(\mu_1,\sigma_1^2),
\qquad
Y\sim \operatorname{Normal}(\mu_2,\sigma_2^2),\] and \(X\) and \(Y\) are
independent. Find the distribution of \(X+Y\).

\begin{tcolorbox}[enhanced jigsaw, arc=.35mm, bottomrule=.15mm, bottomtitle=1mm, breakable, colback=white, colbacktitle=quarto-callout-note-color!10!white, colframe=quarto-callout-note-color-frame, coltitle=black, left=2mm, leftrule=.75mm, opacityback=0, opacitybacktitle=0.6, rightrule=.15mm, title=\textcolor{quarto-callout-note-color}{\faInfo}\hspace{0.5em}{Solution}, titlerule=0mm, toprule=.15mm, toptitle=1mm]

The sum of independent normal random variables is normal. The mean and
variance add: \[\mathbb{E}[X+Y]=\mu_1+\mu_2,\]
\[\operatorname{Var}(X+Y)=\sigma_1^2+\sigma_2^2.\] Therefore
\[X+Y\sim \operatorname{Normal}(\mu_1+\mu_2,\sigma_1^2+\sigma_2^2).\]

\end{tcolorbox}

\textbf{Example 20} (Sample mean of normal observations). Suppose
\[X_1,\ldots,X_n\stackrel{\text{iid}}{\sim}\operatorname{Normal}(\mu,\sigma^2),\]
and define \[\bar X_n=\frac1n\sum_{i=1}^n X_i.\] Find the distribution
of \(\bar X_n\).

\begin{tcolorbox}[enhanced jigsaw, arc=.35mm, bottomrule=.15mm, bottomtitle=1mm, breakable, colback=white, colbacktitle=quarto-callout-note-color!10!white, colframe=quarto-callout-note-color-frame, coltitle=black, left=2mm, leftrule=.75mm, opacityback=0, opacitybacktitle=0.6, rightrule=.15mm, title=\textcolor{quarto-callout-note-color}{\faInfo}\hspace{0.5em}{Solution}, titlerule=0mm, toprule=.15mm, toptitle=1mm]

A linear combination of independent normal random variables is normal.
The mean is
\[\mathbb{E}[\bar X_n]=\frac1n\sum_{i=1}^n\mathbb{E}[X_i]=\mu.\] The
variance is
\[\operatorname{Var}(\bar X_n)=\frac1{n^2}\sum_{i=1}^n\operatorname{Var}(X_i)=\frac{n\sigma^2}{n^2}=\frac{\sigma^2}{n}.\]
Therefore
\[\bar X_n\sim \operatorname{Normal}\left(\mu,\frac{\sigma^2}{n}\right).\]

\end{tcolorbox}

\textbf{Example 21} (Chi-square distribution). Suppose
\[X_1,\ldots,X_n\stackrel{\text{iid}}{\sim}\operatorname{Normal}(0,1).\]
Define \[Z_n=X_1^2+\cdots+X_n^2.\] Identify the distribution of \(Z_n\).

\begin{tcolorbox}[enhanced jigsaw, arc=.35mm, bottomrule=.15mm, bottomtitle=1mm, breakable, colback=white, colbacktitle=quarto-callout-note-color!10!white, colframe=quarto-callout-note-color-frame, coltitle=black, left=2mm, leftrule=.75mm, opacityback=0, opacitybacktitle=0.6, rightrule=.15mm, title=\textcolor{quarto-callout-note-color}{\faInfo}\hspace{0.5em}{Solution}, titlerule=0mm, toprule=.15mm, toptitle=1mm]

The sum of squares of \(n\) independent standard normal random variables
is chi-square with \(n\) degrees of freedom. Thus \[Z_n\sim \chi_n^2.\]
Equivalently, \[Z_n\sim \operatorname{Gamma}\left(\frac n2,2\right)\]
under the shape-scale parameterization.

\end{tcolorbox}

\subsection{Transformations of exponential random
variables}\label{transformations-of-exponential-random-variables}

This subsection summarizes common transformations involving independent
exponential random variables.

\textbf{Example 22} (Sum of two independent exponentials). Let \(X\) and
\(Y\) be iid \(\operatorname{Exp}(1)\) random variables, so
\[f_X(x)=e^{-x},\qquad x\geq0,\] and similarly for \(Y\). Let \[U=X+Y.\]
Find the pdf of \(U\).

\begin{tcolorbox}[enhanced jigsaw, arc=.35mm, bottomrule=.15mm, bottomtitle=1mm, breakable, colback=white, colbacktitle=quarto-callout-note-color!10!white, colframe=quarto-callout-note-color-frame, coltitle=black, left=2mm, leftrule=.75mm, opacityback=0, opacitybacktitle=0.6, rightrule=.15mm, title=\textcolor{quarto-callout-note-color}{\faInfo}\hspace{0.5em}{Solution}, titlerule=0mm, toprule=.15mm, toptitle=1mm]

Using convolution,
\[f_U(u)=\int_{-\infty}^{\infty} f_X(x)f_Y(u-x)\,dx.\] The support
requires \(x\geq0\) and \(u-x\geq0\), so \(0\leq x\leq u\). Hence, for
\(u\geq0\), \[f_U(u)=\int_0^u e^{-x}e^{-(u-x)}\,dx
      =\int_0^u e^{-u}\,dx
      =ue^{-u}.\] For \(u<0\), \(f_U(u)=0\). Thus
\[f_U(u)=ue^{-u}\mathbf{1}_{[0,\infty)}(u).\] This is a
\(\operatorname{Gamma}(2,1)\) distribution using the rate
parameterization, or \(\operatorname{Gamma}(2,1)\) using shape-scale
because the rate and scale are both \(1\) in this special case.

\end{tcolorbox}

\textbf{Example 23} (Minimum of two independent exponentials). Let
\(X,Y\stackrel{\text{iid}}{\sim}\operatorname{Exp}(1)\) and let
\[V=\min\{X,Y\}.\] Find the distribution of \(V\).

\begin{tcolorbox}[enhanced jigsaw, arc=.35mm, bottomrule=.15mm, bottomtitle=1mm, breakable, colback=white, colbacktitle=quarto-callout-note-color!10!white, colframe=quarto-callout-note-color-frame, coltitle=black, left=2mm, leftrule=.75mm, opacityback=0, opacitybacktitle=0.6, rightrule=.15mm, title=\textcolor{quarto-callout-note-color}{\faInfo}\hspace{0.5em}{Solution}, titlerule=0mm, toprule=.15mm, toptitle=1mm]

For \(v\geq0\),
\[\mathbb{P}(V>v)=\mathbb{P}(X>v,Y>v)=\mathbb{P}(X>v)\mathbb{P}(Y>v)=e^{-v}e^{-v}=e^{-2v}.\]
Thus \[F_V(v)=1-e^{-2v},\qquad v\geq0.\] Therefore
\[V\sim \operatorname{Exp}(2).\]

\end{tcolorbox}

\textbf{Example 24} (Ratio of exponentials). Let
\(X,Y\stackrel{\text{iid}}{\sim}\operatorname{Exp}(1)\) and define
\[U=\frac{X}{X+Y}.\] Show that \(U\sim\operatorname{Uniform}(0,1)\).

\begin{tcolorbox}[enhanced jigsaw, arc=.35mm, bottomrule=.15mm, bottomtitle=1mm, breakable, colback=white, colbacktitle=quarto-callout-note-color!10!white, colframe=quarto-callout-note-color-frame, coltitle=black, left=2mm, leftrule=.75mm, opacityback=0, opacitybacktitle=0.6, rightrule=.15mm, title=\textcolor{quarto-callout-note-color}{\faInfo}\hspace{0.5em}{Solution}, titlerule=0mm, toprule=.15mm, toptitle=1mm]

Use the transformation \[U=\frac{X}{X+Y},\qquad T=X+Y.\] Then \[X=UT,
\qquad
Y=(1-U)T.\] The support is \(0<u<1\) and \(t>0\). The Jacobian is
\[\left|\frac{\partial(x,y)}{\partial(u,t)}\right|
=\left|
\begin{matrix}
t & u\\
-t & 1-u
\end{matrix}
\right|=t.\] The joint density of \((X,Y)\) is
\[f_{X,Y}(x,y)=e^{-x-y},\qquad x,y>0.\] Therefore
\[f_{U,T}(u,t)=e^{-t}t,
\qquad 0<u<1,\; t>0.\] The marginal density of \(U\) is
\[f_U(u)=\int_0^\infty t e^{-t}\,dt=1,
\qquad 0<u<1.\] Thus \[U\sim\operatorname{Uniform}(0,1).\]

\end{tcolorbox}

\textbf{Example 25} (Difference of two independent exponentials). Let
\(X,Y\stackrel{\text{iid}}{\sim}\operatorname{Exp}(1)\) and define
\[W=X-Y.\] Find the distribution of \(W\).

\begin{tcolorbox}[enhanced jigsaw, arc=.35mm, bottomrule=.15mm, bottomtitle=1mm, breakable, colback=white, colbacktitle=quarto-callout-note-color!10!white, colframe=quarto-callout-note-color-frame, coltitle=black, left=2mm, leftrule=.75mm, opacityback=0, opacitybacktitle=0.6, rightrule=.15mm, title=\textcolor{quarto-callout-note-color}{\faInfo}\hspace{0.5em}{Solution}, titlerule=0mm, toprule=.15mm, toptitle=1mm]

Use the density of a difference. Since \(W=X-Y\), for
\(w\in\mathbb{R}\),
\[f_W(w)=\int_{-\infty}^{\infty} f_X(w+y)f_Y(y)\,dy,\] where the support
requires \(y\geq0\) and \(w+y\geq0\).

If \(w\geq0\), then \(y\geq0\), so
\[f_W(w)=\int_0^\infty e^{-(w+y)}e^{-y}\,dy
=e^{-w}\int_0^\infty e^{-2y}\,dy
=\frac12 e^{-w}.\] If \(w<0\), then \(y\geq -w\), so
\[f_W(w)=\int_{-w}^\infty e^{-(w+y)}e^{-y}\,dy
=e^{-w}\int_{-w}^\infty e^{-2y}\,dy
=\frac12 e^{w}.\] Thus
\[f_W(w)=\frac12 e^{-|w|},\qquad -\infty<w<\infty.\] This is the
Laplace, or double exponential, distribution with location \(0\) and
scale \(1\).

\end{tcolorbox}

\section{Transformations of Several Random Variables and
Jacobians}\label{transformations-of-several-random-variables-and-jacobians}

This section gives the general multivariable change-of-variables
theorem.

The one-dimensional formula contains a derivative. In several
dimensions, that derivative is replaced by the absolute value of a
determinant, called the Jacobian determinant.

\subsection{The Jacobian theorem}\label{the-jacobian-theorem}

This subsection states the main theorem for one-to-one transformations
in several dimensions.

\textbf{Theorem 26} (Multivariable change of variables). \emph{Let
\(X=(X_1,\ldots,X_n)\) have joint pdf \(f_X(x_1,\ldots,x_n)\). Suppose
\[Y_i=g_i(X_1,\ldots,X_n),\qquad i=1,\ldots,n,\] and suppose the
transformation is one-to-one on the relevant support with inverse
\[X_i=h_i(Y_1,\ldots,Y_n),\qquad i=1,\ldots,n.\] Then the joint pdf of
\(Y=(Y_1,\ldots,Y_n)\) is \[f_Y(y_1,\ldots,y_n)
=f_X(h_1(y),\ldots,h_n(y))
\left|\frac{\partial(x_1,\ldots,x_n)}{\partial(y_1,\ldots,y_n)}\right|,\]
where
\[\frac{\partial(x_1,\ldots,x_n)}{\partial(y_1,\ldots,y_n)}=\begin{vmatrix}
\dfrac{\partial x_1}{\partial y_1} & \dfrac{\partial x_1}{\partial y_2} & \cdots & \dfrac{\partial x_1}{\partial y_n}\\[0.7em]
\dfrac{\partial x_2}{\partial y_1} & \dfrac{\partial x_2}{\partial y_2} & \cdots & \dfrac{\partial x_2}{\partial y_n}\\
\vdots & \vdots & \ddots & \vdots\\
\dfrac{\partial x_n}{\partial y_1} & \dfrac{\partial x_n}{\partial y_2} & \cdots & \dfrac{\partial x_n}{\partial y_n}
\end{vmatrix}.\]}

\begin{tcolorbox}[enhanced jigsaw, arc=.35mm, bottomrule=.15mm, bottomtitle=1mm, breakable, colback=white, colbacktitle=quarto-callout-warning-color!10!white, colframe=quarto-callout-warning-color-frame, coltitle=black, left=2mm, leftrule=.75mm, opacityback=0, opacitybacktitle=0.6, rightrule=.15mm, title=\textcolor{quarto-callout-warning-color}{\faExclamationTriangle}\hspace{0.5em}{Support matters}, titlerule=0mm, toprule=.15mm, toptitle=1mm]

When using the Jacobian formula, always transform the support as well as
the formula. A correct density with the wrong support is not a correct
answer.

\end{tcolorbox}

\subsection{Spacing transformation
example}\label{spacing-transformation-example}

This subsection works through a triangular transformation from ordered
variables to spacings.

\textbf{Example 27} (Ordered exponential-type joint density). Let
\(X_1,X_2,X_3,X_4\) have joint pdf
\[f_X(x_1,x_2,x_3,x_4)=24e^{-x_1-x_2-x_3-x_4},\] for
\[0<x_1<x_2<x_3<x_4<\infty.\] Define \[Y_1=X_1,
\qquad
Y_2=X_2-X_1,
\qquad
Y_3=X_3-X_2,
\qquad
Y_4=X_4-X_3.\] Find the joint pdf of \((Y_1,Y_2,Y_3,Y_4)\).

\begin{tcolorbox}[enhanced jigsaw, arc=.35mm, bottomrule=.15mm, bottomtitle=1mm, breakable, colback=white, colbacktitle=quarto-callout-note-color!10!white, colframe=quarto-callout-note-color-frame, coltitle=black, left=2mm, leftrule=.75mm, opacityback=0, opacitybacktitle=0.6, rightrule=.15mm, title=\textcolor{quarto-callout-note-color}{\faInfo}\hspace{0.5em}{Solution}, titlerule=0mm, toprule=.15mm, toptitle=1mm]

The inverse transformation is \[X_1=Y_1,\] \[X_2=Y_1+Y_2,\]
\[X_3=Y_1+Y_2+Y_3,\] \[X_4=Y_1+Y_2+Y_3+Y_4.\] The support becomes
\[y_1>0,
\qquad y_2>0,
\qquad y_3>0,
\qquad y_4>0.\] The Jacobian matrix is lower triangular:
\[\frac{\partial(x_1,x_2,x_3,x_4)}{\partial(y_1,y_2,y_3,y_4)}
=
\begin{pmatrix}
1&0&0&0\\
1&1&0&0\\
1&1&1&0\\
1&1&1&1
\end{pmatrix},\] so its determinant is \(1\).

Substitute the inverse into the original density: \[x_1+x_2+x_3+x_4
=y_1+(y_1+y_2)+(y_1+y_2+y_3)+(y_1+y_2+y_3+y_4).\] Thus
\[x_1+x_2+x_3+x_4=4y_1+3y_2+2y_3+y_4.\] Therefore \[f_Y(y_1,y_2,y_3,y_4)
=24e^{-4y_1-3y_2-2y_3-y_4},\] for \(y_1,y_2,y_3,y_4>0\), and \(0\)
otherwise.

\end{tcolorbox}

\subsection{Polar coordinates and the standard normal
distribution}\label{polar-coordinates-and-the-standard-normal-distribution}

This subsection uses the Jacobian to explain the geometry behind the
Box--Muller method.

\textbf{Example 28} (Polar transformation of two independent standard
normals). Let \(Z_1\) and \(Z_2\) be independent standard normal random
variables. Their joint pdf is
\[f_{Z_1,Z_2}(z_1,z_2)=\frac1{2\pi}e^{-(z_1^2+z_2^2)/2}.\] Use the polar
transformation \[Z_1=R\cos\Theta,
\qquad
Z_2=R\sin\Theta.\] Find the joint density of \((R,\Theta)\).

\begin{tcolorbox}[enhanced jigsaw, arc=.35mm, bottomrule=.15mm, bottomtitle=1mm, breakable, colback=white, colbacktitle=quarto-callout-note-color!10!white, colframe=quarto-callout-note-color-frame, coltitle=black, left=2mm, leftrule=.75mm, opacityback=0, opacitybacktitle=0.6, rightrule=.15mm, title=\textcolor{quarto-callout-note-color}{\faInfo}\hspace{0.5em}{Solution}, titlerule=0mm, toprule=.15mm, toptitle=1mm]

The support is \[r\geq0,
\qquad 0\leq \theta<2\pi.\] The Jacobian determinant for polar
coordinates is
\[\left|\frac{\partial(z_1,z_2)}{\partial(r,\theta)}\right|=r.\] Also,
\[z_1^2+z_2^2=r^2.\] Therefore
\[f_{R,\Theta}(r,\theta)=\frac1{2\pi}e^{-r^2/2}r,
\qquad r\geq0,
\quad 0\leq\theta<2\pi.\] This factors as
\[f_{R,\Theta}(r,\theta)=\left(re^{-r^2/2}\right)\left(\frac1{2\pi}\right).\]
Hence \(R\) and \(\Theta\) are independent,
\[\Theta\sim \operatorname{Uniform}(0,2\pi),\] and
\[f_R(r)=re^{-r^2/2},\qquad r\geq0.\] The distribution of \(R\) is
Rayleigh.

\end{tcolorbox}

\textbf{Example 29} (Box--Muller method). Let
\(U_1,U_2\stackrel{\text{iid}}{\sim}\operatorname{Uniform}(0,1)\).
Define \[R=\sqrt{-2\log U_1},
\qquad
\Theta=2\pi U_2.\] Then define \[Z_1=R\cos\Theta,
\qquad
Z_2=R\sin\Theta.\] Show that \(Z_1\) and \(Z_2\) are independent
standard normal random variables.

\begin{tcolorbox}[enhanced jigsaw, arc=.35mm, bottomrule=.15mm, bottomtitle=1mm, breakable, colback=white, colbacktitle=quarto-callout-note-color!10!white, colframe=quarto-callout-note-color-frame, coltitle=black, left=2mm, leftrule=.75mm, opacityback=0, opacitybacktitle=0.6, rightrule=.15mm, title=\textcolor{quarto-callout-note-color}{\faInfo}\hspace{0.5em}{Solution}, titlerule=0mm, toprule=.15mm, toptitle=1mm]

From the previous example, two independent standard normals have polar
representation where \[\Theta\sim\operatorname{Uniform}(0,2\pi)\] and
\[f_R(r)=re^{-r^2/2},\qquad r\geq0.\] The CDF of \(R\) is
\[F_R(r)=\int_0^r se^{-s^2/2}\,ds=1-e^{-r^2/2}.\] Using inverse
transform sampling, if \(U_1\sim\operatorname{Uniform}(0,1)\), then
\[R=\sqrt{-2\log U_1}\] has this Rayleigh distribution. Also,
\(\Theta=2\pi U_2\) is uniform on \([0,2\pi]\). Since \(U_1\) and
\(U_2\) are independent, \(R\) and \(\Theta\) are independent.
Transforming back by \[Z_1=R\cos\Theta,
\qquad
Z_2=R\sin\Theta\] produces the joint density
\[\frac1{2\pi}e^{-(z_1^2+z_2^2)/2},\] which factors into the product of
two standard normal densities. Hence \(Z_1\) and \(Z_2\) are independent
standard normal random variables.

\end{tcolorbox}

\section{Sums, Quotients, Products, Maxima, and
Minima}\label{sums-quotients-products-maxima-and-minima}

This section records general formulas for common transformations of two
or more random variables.

These formulas are important because many statistical quantities are
built from sums, ratios, products, maxima, or minima.

\subsection{Sum of independent random
variables}\label{sum-of-independent-random-variables}

This subsection introduces convolution, one of the most important
operations in probability.

\textbf{Theorem 30} (Sum: discrete case). \emph{Suppose \(X\) and \(Y\)
are independent discrete random variables with pmfs \(p_X\) and \(p_Y\).
Let \[W=X+Y.\] Then \[p_W(w)=\sum_x p_X(x)p_Y(w-x).\]}

\begin{tcolorbox}[enhanced jigsaw, arc=.35mm, bottomrule=.15mm, bottomtitle=1mm, breakable, colback=white, colbacktitle=quarto-callout-note-color!10!white, colframe=quarto-callout-note-color-frame, coltitle=black, left=2mm, leftrule=.75mm, opacityback=0, opacitybacktitle=0.6, rightrule=.15mm, title=\textcolor{quarto-callout-note-color}{\faInfo}\hspace{0.5em}{Proof}, titlerule=0mm, toprule=.15mm, toptitle=1mm]

We decompose the event \(W=w\) according to all possible values of
\(X\): \[p_W(w)=\mathbb{P}(W=w)=\mathbb{P}(X+Y=w).\] Thus
\[\mathbb{P}(X+Y=w)=\sum_x \mathbb{P}(X=x,Y=w-x).\] By independence,
\[\mathbb{P}(X=x,Y=w-x)=\mathbb{P}(X=x)\mathbb{P}(Y=w-x).\] Therefore
\[p_W(w)=\sum_x p_X(x)p_Y(w-x).\]

\end{tcolorbox}

\textbf{Theorem 31} (Sum: continuous case). \emph{Suppose \(X\) and
\(Y\) are independent continuous random variables with pdfs \(f_X\) and
\(f_Y\). Let \[W=X+Y.\] Then
\[f_W(w)=\int_{-\infty}^{\infty} f_X(x)f_Y(w-x)\,dx.\]}

\begin{tcolorbox}[enhanced jigsaw, arc=.35mm, bottomrule=.15mm, bottomtitle=1mm, breakable, colback=white, colbacktitle=quarto-callout-note-color!10!white, colframe=quarto-callout-note-color-frame, coltitle=black, left=2mm, leftrule=.75mm, opacityback=0, opacitybacktitle=0.6, rightrule=.15mm, title=\textcolor{quarto-callout-note-color}{\faInfo}\hspace{0.5em}{Proof}, titlerule=0mm, toprule=.15mm, toptitle=1mm]

Using the CDF method, \[F_W(w)=\mathbb{P}(X+Y\leq w)
=\int_{-\infty}^{\infty}\int_{-\infty}^{w-x} f_X(x)f_Y(y)\,dy\,dx.\] The
inner integral is \(F_Y(w-x)\), so
\[F_W(w)=\int_{-\infty}^{\infty} f_X(x)F_Y(w-x)\,dx.\] Differentiate
with respect to \(w\):
\[f_W(w)=\int_{-\infty}^{\infty} f_X(x)f_Y(w-x)\,dx.\]

\end{tcolorbox}

\textbf{Practice Problem 32} (Practice: sum of two dice). Let \(X\) and
\(Y\) be independent rolls of a fair six-sided die. Find
\(\mathbb{P}(X+Y=7)\).

\begin{tcolorbox}[enhanced jigsaw, arc=.35mm, bottomrule=.15mm, bottomtitle=1mm, breakable, colback=white, colbacktitle=quarto-callout-note-color!10!white, colframe=quarto-callout-note-color-frame, coltitle=black, left=2mm, leftrule=.75mm, opacityback=0, opacitybacktitle=0.6, rightrule=.15mm, title=\textcolor{quarto-callout-note-color}{\faInfo}\hspace{0.5em}{Solution}, titlerule=0mm, toprule=.15mm, toptitle=1mm]

Using convolution,
\[\mathbb{P}(X+Y=7)=\sum_x \mathbb{P}(X=x)\mathbb{P}(Y=7-x).\] The valid
pairs are \[(1,6),(2,5),(3,4),(4,3),(5,2),(6,1),\] so there are \(6\)
favorable outcomes out of \(36\). Therefore
\[\mathbb{P}(X+Y=7)=\frac6{36}=\frac16.\]

\end{tcolorbox}

\subsection{Quotient of two independent continuous random
variables}\label{quotient-of-two-independent-continuous-random-variables}

This subsection gives the density formula for a ratio.

\textbf{Theorem 33} (Quotient). \emph{Suppose \(X\) and \(Y\) are
independent continuous random variables with pdfs \(f_X\) and \(f_Y\).
Let \[W=\frac{Y}{X}.\] Then
\[f_W(w)=\int_{-\infty}^{\infty} |x|f_X(x)f_Y(wx)\,dx.\]}

\begin{tcolorbox}[enhanced jigsaw, arc=.35mm, bottomrule=.15mm, bottomtitle=1mm, breakable, colback=white, colbacktitle=quarto-callout-note-color!10!white, colframe=quarto-callout-note-color-frame, coltitle=black, left=2mm, leftrule=.75mm, opacityback=0, opacitybacktitle=0.6, rightrule=.15mm, title=\textcolor{quarto-callout-note-color}{\faInfo}\hspace{0.5em}{Proof}, titlerule=0mm, toprule=.15mm, toptitle=1mm]

Use the transformation \[W=\frac{Y}{X},
\qquad
T=X.\] Then \[X=T,
\qquad
Y=WT.\] The Jacobian determinant is
\[\left|\frac{\partial(x,y)}{\partial(w,t)}\right|
=\left|
\begin{matrix}
0&1\\
t&w
\end{matrix}
\right|=|t|.\] By independence, \[f_{W,T}(w,t)=f_X(t)f_Y(wt)|t|.\]
Integrating out \(t\) gives
\[f_W(w)=\int_{-\infty}^{\infty}|t|f_X(t)f_Y(wt)\,dt.\] Renaming \(t\)
as \(x\) gives the formula.

\end{tcolorbox}

\subsection{Product of two independent continuous random
variables}\label{product-of-two-independent-continuous-random-variables}

This subsection gives the density formula for a product.

\textbf{Theorem 34} (Product). \emph{Suppose \(X\) and \(Y\) are
independent continuous random variables with pdfs \(f_X\) and \(f_Y\).
Let \[W=XY.\] Then
\[f_W(w)=\int_{-\infty}^{\infty}\frac1{|x|}f_X(x)f_Y\left(\frac{w}{x}\right)\,dx,\]
where the integral is over values \(x\ne0\) for which the densities are
defined.}

\begin{tcolorbox}[enhanced jigsaw, arc=.35mm, bottomrule=.15mm, bottomtitle=1mm, breakable, colback=white, colbacktitle=quarto-callout-note-color!10!white, colframe=quarto-callout-note-color-frame, coltitle=black, left=2mm, leftrule=.75mm, opacityback=0, opacitybacktitle=0.6, rightrule=.15mm, title=\textcolor{quarto-callout-note-color}{\faInfo}\hspace{0.5em}{Proof}, titlerule=0mm, toprule=.15mm, toptitle=1mm]

Use the transformation \[W=XY,
\qquad
T=X.\] Then \[X=T,
\qquad
Y=\frac WT.\] The Jacobian determinant is
\[\left|\frac{\partial(x,y)}{\partial(w,t)}\right|
=\left|
\begin{matrix}
0&1\\[0.3em]
1/t&-w/t^2
\end{matrix}
\right|=\frac1{|t|}.\] Thus
\[f_{W,T}(w,t)=f_X(t)f_Y\left(\frac wt\right)\frac1{|t|}.\] Integrating
out \(t\) gives
\[f_W(w)=\int_{-\infty}^{\infty}\frac1{|t|}f_X(t)f_Y\left(\frac wt\right)\,dt.\]

\end{tcolorbox}

\begin{remark}
\emph{Remark 35}. If \(X\) and \(Y\) are not independent, replace the
product \(f_X(x)f_Y(y)\) by the joint density \(f_{X,Y}(x,y)\) and use
the appropriate Jacobian transformation.
\end{remark}

\subsection{Maximum and minimum of iid continuous random
variables}\label{maximum-and-minimum-of-iid-continuous-random-variables}

This subsection derives the distributions of order-statistic extremes.

\textbf{Theorem 36} (Maximum and minimum). \emph{Suppose
\(Y_1,\ldots,Y_n\) are iid continuous random variables with CDF \(F_Y\)
and pdf \(f_Y\).}

\begin{enumerate}
\def\labelenumi{\arabic{enumi}.}
\item
  \emph{If \[Y_{\max}=\max\{Y_1,\ldots,Y_n\},\] then
  \[f_{Y_{\max}}(y)=n[F_Y(y)]^{n-1}f_Y(y).\]}
\item
  \emph{If \[Y_{\min}=\min\{Y_1,\ldots,Y_n\},\] then
  \[f_{Y_{\min}}(y)=n[1-F_Y(y)]^{n-1}f_Y(y).\]}
\end{enumerate}

\begin{tcolorbox}[enhanced jigsaw, arc=.35mm, bottomrule=.15mm, bottomtitle=1mm, breakable, colback=white, colbacktitle=quarto-callout-note-color!10!white, colframe=quarto-callout-note-color-frame, coltitle=black, left=2mm, leftrule=.75mm, opacityback=0, opacitybacktitle=0.6, rightrule=.15mm, title=\textcolor{quarto-callout-note-color}{\faInfo}\hspace{0.5em}{Proof}, titlerule=0mm, toprule=.15mm, toptitle=1mm]

For the maximum, \[F_{Y_{\max}}(y)=\mathbb{P}(Y_{\max}\leq y)
=\mathbb{P}(Y_1\leq y,\ldots,Y_n\leq y).\] By independence,
\[F_{Y_{\max}}(y)=[F_Y(y)]^n.\] Differentiating gives
\[f_{Y_{\max}}(y)=n[F_Y(y)]^{n-1}f_Y(y).\]

For the minimum, use the complement:
\[F_{Y_{\min}}(y)=\mathbb{P}(Y_{\min}\leq y)
=1-\mathbb{P}(Y_{\min}>y).\] Now
\[\mathbb{P}(Y_{\min}>y)=\mathbb{P}(Y_1>y,\ldots,Y_n>y)=[1-F_Y(y)]^n.\]
Thus \[F_{Y_{\min}}(y)=1-[1-F_Y(y)]^n.\] Differentiating gives
\[f_{Y_{\min}}(y)=n[1-F_Y(y)]^{n-1}f_Y(y).\]

\end{tcolorbox}

\textbf{Practice Problem 37} (Practice: maximum of iid uniforms). Let
\(Y_1,\ldots,Y_n\stackrel{\text{iid}}{\sim}\operatorname{Uniform}(0,1)\)
and let \(M=\max\{Y_1,\ldots,Y_n\}\). Find the pdf of \(M\).

\begin{tcolorbox}[enhanced jigsaw, arc=.35mm, bottomrule=.15mm, bottomtitle=1mm, breakable, colback=white, colbacktitle=quarto-callout-note-color!10!white, colframe=quarto-callout-note-color-frame, coltitle=black, left=2mm, leftrule=.75mm, opacityback=0, opacitybacktitle=0.6, rightrule=.15mm, title=\textcolor{quarto-callout-note-color}{\faInfo}\hspace{0.5em}{Solution}, titlerule=0mm, toprule=.15mm, toptitle=1mm]

For \(0\leq y\leq1\), \[F_Y(y)=y,
\qquad
f_Y(y)=1.\] By the maximum formula,
\[f_M(y)=n[F_Y(y)]^{n-1}f_Y(y)=ny^{n-1},
\qquad 0\leq y\leq1.\] Thus \[f_M(y)=ny^{n-1}\mathbf{1}_{[0,1]}(y).\]

\end{tcolorbox}

\section{Summary and Study Guide}\label{summary-and-study-guide}

This section summarizes the main transformation tools and when to use
each one.

\begin{tcolorbox}[enhanced jigsaw, arc=.35mm, bottomrule=.15mm, bottomtitle=1mm, breakable, colback=white, colbacktitle=quarto-callout-tip-color!10!white, colframe=quarto-callout-tip-color-frame, coltitle=black, left=2mm, leftrule=.75mm, opacityback=0, opacitybacktitle=0.6, rightrule=.15mm, title=\textcolor{quarto-callout-tip-color}{\faLightbulb}\hspace{0.5em}{Which method should I use?}, titlerule=0mm, toprule=.15mm, toptitle=1mm]

\begin{itemize}
\item
  \textbf{Discrete transformation:} group the preimages and add
  probabilities.
\item
  \textbf{One continuous variable:} use the CDF method, or use the
  one-to-one change-of-variables formula.
\item
  \textbf{Several continuous variables:} use a one-to-one transformation
  and the Jacobian.
\item
  \textbf{Sum of independent variables:} use convolution.
\item
  \textbf{Maximum/minimum of iid variables:} use CDF complements and
  independence.
\item
  \textbf{Simulation from a target distribution:} use inverse-transform
  sampling when the inverse CDF is available.
\end{itemize}

\end{tcolorbox}

{\def\LTcaptype{none} % do not increment counter
\begin{longtable}[]{@{}ll@{}}
\toprule\noalign{}
\textbf{Transformation type} & \textbf{Formula or method} \\
\midrule\noalign{}
\endhead
\bottomrule\noalign{}
\endlastfoot
Discrete \(Y=g(X)\) & \(p_Y(y)=\sum_{x:g(x)=y}p_X(x)\) \\
Continuous CDF method & \(F_Y(y)=\mathbb{P}(g(X)\leq y)\), then
\(f_Y(y)=F_Y'(y)\) \\
One-to-one \(Y=g(X)\) &
\(f_Y(y)=f_X(g^{-1}(y))\left|\dfrac{\,d}{\,dy}g^{-1}(y)\right|\) \\
Multivariable transform &
\(f_Y(y)=f_X(h(y))\left|\dfrac{\partial x}{\partial y}\right|\) \\
Sum \(W=X+Y\) independent & \(f_W(w)=\int f_X(x)f_Y(w-x)\,dx\) \\
Quotient \(W=Y/X\) independent & \(f_W(w)=\int |x|f_X(x)f_Y(wx)\,dx\) \\
Product \(W=XY\) independent &
\(f_W(w)=\int \dfrac1{|x|}f_X(x)f_Y(w/x)\,dx\) \\
Maximum of iid variables & \(f_{\max}(y)=n[F_Y(y)]^{n-1}f_Y(y)\) \\
Minimum of iid variables & \(f_{\min}(y)=n[1-F_Y(y)]^{n-1}f_Y(y)\) \\
\end{longtable}
}

\textbf{Practice Problem 38} (Comprehensive practice). Let
\(X\sim\operatorname{Uniform}(0,1)\) and define \(Y=-\log X\). Then let
\(Y_1,Y_2\) be independent copies of \(Y\) and set \(S=Y_1+Y_2\).

\begin{enumerate}
\def\labelenumi{\arabic{enumi}.}
\item
  Find the distribution of \(Y\).
\item
  Find the pdf of \(S\).
\end{enumerate}

\begin{tcolorbox}[enhanced jigsaw, arc=.35mm, bottomrule=.15mm, bottomtitle=1mm, breakable, colback=white, colbacktitle=quarto-callout-note-color!10!white, colframe=quarto-callout-note-color-frame, coltitle=black, left=2mm, leftrule=.75mm, opacityback=0, opacitybacktitle=0.6, rightrule=.15mm, title=\textcolor{quarto-callout-note-color}{\faInfo}\hspace{0.5em}{Solution}, titlerule=0mm, toprule=.15mm, toptitle=1mm]

\begin{enumerate}
\def\labelenumi{\arabic{enumi}.}
\item
  Since \(Y=-\log X\), the inverse is \(x=e^{-y}\) for \(y\geq0\), and
  \[\left|\frac{\,dx}{\,dy}\right|=e^{-y}.\] Thus
  \[f_Y(y)=e^{-y},\qquad y\geq0.\] So \[Y\sim \operatorname{Exp}(1).\]
\item
  Since \(Y_1,Y_2\) are independent \(\operatorname{Exp}(1)\) random
  variables, use convolution: \[f_S(s)=\int_0^s e^{-y}e^{-(s-y)}\,dy
  =s e^{-s},
  \qquad s\geq0.\] Thus \[S\sim \operatorname{Gamma}(2,1)\] under the
  rate-\(1\) parameterization.
\end{enumerate}

\end{tcolorbox}

\bookmarksetup{startatroot}

\chapter{Chapter 6: Conditional
Expectations}\label{chapter-6-conditional-expectations}

This chapter explains how conditional distributions and conditional
expectations turn difficult probability, expectation, and variance
calculations into easier two-step calculations: condition first, then
average.

\textbf{Topics.} Conditional distributions; conditioning as a
problem-solving method; law of total probability; memoryless property;
conditional expectation; law of total expectation; law of total
variance; random sums; Bayesian updating; inverse probability weighting;
importance weighting.

\section{Overview}\label{overview-3}

This section explains how conditioning turns a difficult probability or
expectation into an average of easier conditional probabilities or
conditional expectations.

Conditional probability was introduced earlier as
\[\mathbb{P}(A\mid B)=\frac{\mathbb{P}(A\cap B)}{\mathbb{P}(B)}.\] In
this section, we extend the idea from events to random variables. We
will use conditional distributions such as \(p_{X\mid Y}(x\mid y)\) and
conditional expectations such as \(\mathbb{E}[X\mid Y]\).

\begin{tcolorbox}[enhanced jigsaw, arc=.35mm, bottomrule=.15mm, bottomtitle=1mm, breakable, colback=white, colbacktitle=quarto-callout-tip-color!10!white, colframe=quarto-callout-tip-color-frame, coltitle=black, left=2mm, leftrule=.75mm, opacityback=0, opacitybacktitle=0.6, rightrule=.15mm, title=\textcolor{quarto-callout-tip-color}{\faLightbulb}\hspace{0.5em}{Main message}, titlerule=0mm, toprule=.15mm, toptitle=1mm]

Conditioning is a way to break a problem into two steps: first solve the
problem after some information is fixed, then average over the
information that was fixed.
\[\text{unconditional quantity}=\text{average of conditional quantities}.\]

\end{tcolorbox}

\section{Conditional Distributions}\label{conditional-distributions-1}

This section reviews conditional distributions and shows how they are
used to compute probabilities by conditioning on another random
variable.

\subsection{Discrete conditional
distributions}\label{discrete-conditional-distributions}

This subsection begins with the discrete case, where conditional
distributions are ratios of joint probabilities and marginals.

Let \(X\) and \(Y\) be discrete random variables. If \(p_Y(y)>0\), then
\[p_{X\mid Y}(x\mid y)
=\mathbb{P}(X=x\mid Y=y)
=\frac{p_{X,Y}(x,y)}{p_Y(y)}.\] Equivalently,
\[p_{X,Y}(x,y)=p_{X\mid Y}(x\mid y)p_Y(y).\] Therefore, summing over all
possible values of \(Y\) gives the total probability formula
\[\mathbb{P}(X=x)=\sum_y \mathbb{P}(X=x\mid Y=y)\mathbb{P}(Y=y).\]

\begin{tcolorbox}[enhanced jigsaw, arc=.35mm, bottomrule=.15mm, bottomtitle=1mm, breakable, colback=white, colbacktitle=quarto-callout-tip-color!10!white, colframe=quarto-callout-tip-color-frame, coltitle=black, left=2mm, leftrule=.75mm, opacityback=0, opacitybacktitle=0.6, rightrule=.15mm, title=\textcolor{quarto-callout-tip-color}{\faLightbulb}\hspace{0.5em}{Why conditioning helps}, titlerule=0mm, toprule=.15mm, toptitle=1mm]

It is often easier to compute \(\mathbb{P}(X=x\mid Y=y)\) after \(Y\) is
fixed. Then the unconditional probability is obtained by summing over
all possible values of \(Y\).

\end{tcolorbox}

\textbf{Example 1} (Best prize / secretary problem). There are \(n\)
distinct prizes arriving in a random order. Exactly one prize is the
best. You must accept or reject each prize when it arrives, and you
cannot return to earlier prizes. Consider the strategy:

\begin{quote}
Reject the first \(k\) prizes, then choose the first later prize that is
better than all previous prizes.
\end{quote}

Let \(X\) be the position of the best prize. Find the approximate
probability of selecting the best prize under this strategy and find the
approximate optimal value of \(k\).

\begin{tcolorbox}[enhanced jigsaw, arc=.35mm, bottomrule=.15mm, bottomtitle=1mm, breakable, colback=white, colbacktitle=quarto-callout-note-color!10!white, colframe=quarto-callout-note-color-frame, coltitle=black, left=2mm, leftrule=.75mm, opacityback=0, opacitybacktitle=0.6, rightrule=.15mm, title=\textcolor{quarto-callout-note-color}{\faInfo}\hspace{0.5em}{Solution}, titlerule=0mm, toprule=.15mm, toptitle=1mm]

The best prize is equally likely to occur in any position, so
\[\mathbb{P}(X=i)=\frac1n,\qquad i=1,\ldots,n.\] If \(i\le k\), the best
prize occurs during the rejection period, so the strategy cannot win.

If \(i>k\), then the strategy wins exactly when the best prize among the
first \(i-1\) positions occurs among the first \(k\) positions. Since
the best among the first \(i-1\) positions is equally likely to be in
any of those positions,
\[\mathbb{P}(\text{win}\mid X=i)=\frac{k}{i-1},\qquad i=k+1,
\ldots,n.\] Hence \[\mathbb{P}_k(\text{win})
=\sum_{i=k+1}^n \mathbb{P}(\text{win}\mid X=i)\mathbb{P}(X=i)
=\sum_{i=k+1}^n \frac{k}{i-1}\frac1n.\] Thus \[\mathbb{P}_k(\text{win})
=\frac{k}{n}\sum_{j=k}^{n-1}\frac1j
\approx \frac{k}{n}\log\left(\frac{n}{k}\right).\] Let \(x=k/n\). Then
approximately \[f(x)=x\log\left(\frac1x\right)=-x\log x.\]
Differentiate: \[f'(x)=-\log x-1.\] Set \(f'(x)=0\):
\[-\log x-1=0 \quad\Longrightarrow\quad x=e^{-1}.\] So the approximate
optimal choice is \[\boxed{k\approx \frac{n}{e}},\] and the
corresponding maximum probability is approximately
\[\boxed{\frac1e\approx 0.368.}\]

\end{tcolorbox}

\subsection{Conditioning on the first
step}\label{conditioning-on-the-first-step}

This subsection shows a common recursive method: condition on the first
move of a stochastic process.

\textbf{Example 2} (Gambler's ruin). A gambler starts with \(k\)
dollars. In each round, the gambler wins \(1\) dollar with probability
\(p\) and loses \(1\) dollar with probability \(q=1-p\). The game stops
when the gambler reaches \(N\) dollars or \(0\) dollars. Let
\[P_k=\mathbb{P}(\text{reach }N\text{ before }0\mid\text{start at }k).\]
Use conditioning on the first bet to derive the recursion for \(P_k\).

\begin{tcolorbox}[enhanced jigsaw, arc=.35mm, bottomrule=.15mm, bottomtitle=1mm, breakable, colback=white, colbacktitle=quarto-callout-note-color!10!white, colframe=quarto-callout-note-color-frame, coltitle=black, left=2mm, leftrule=.75mm, opacityback=0, opacitybacktitle=0.6, rightrule=.15mm, title=\textcolor{quarto-callout-note-color}{\faInfo}\hspace{0.5em}{Solution}, titlerule=0mm, toprule=.15mm, toptitle=1mm]

Let \(A\) be the event that the gambler eventually reaches \(N\) before
going broke. Let \(X\) be the result of the first bet, where \(X=+1\)
with probability \(p\) and \(X=-1\) with probability \(q\).

Conditioning on the first bet gives \[P_k=\mathbb{P}(A)
=\mathbb{P}(A\mid X=+1)\mathbb{P}(X=+1)+\mathbb{P}(A\mid X=-1)\mathbb{P}(X=-1).\]
If the first bet is a win, the gambler moves to \(k+1\) dollars. If it
is a loss, the gambler moves to \(k-1\) dollars. Therefore
\[\boxed{P_k=pP_{k+1}+qP_{k-1}},\qquad k=1,\ldots,N-1,\] with boundary
conditions \[P_0=0,\qquad P_N=1.\] If desired, solving the difference
equation gives \[P_k=\begin{cases}
\displaystyle \frac{1-(q/p)^k}{1-(q/p)^N}, & p\ne q,\\[1.2em]
\displaystyle \frac{k}{N}, & p=q=\frac12.
\end{cases}\]

\end{tcolorbox}

\textbf{Practice Problem 3} (First-step recursion). In the gambler's
ruin problem, suppose \(p=q=1/2\), \(N=10\), and the gambler starts with
\(k=4\) dollars. What is the probability that the gambler reaches \(10\)
dollars before going broke?

\begin{tcolorbox}[enhanced jigsaw, arc=.35mm, bottomrule=.15mm, bottomtitle=1mm, breakable, colback=white, colbacktitle=quarto-callout-note-color!10!white, colframe=quarto-callout-note-color-frame, coltitle=black, left=2mm, leftrule=.75mm, opacityback=0, opacitybacktitle=0.6, rightrule=.15mm, title=\textcolor{quarto-callout-note-color}{\faInfo}\hspace{0.5em}{Solution}, titlerule=0mm, toprule=.15mm, toptitle=1mm]

For the fair game \(p=q=1/2\), the solution is \[P_k=\frac{k}{N}.\] Thus
\[P_4=\frac{4}{10}=0.4.\] So the probability is \[\boxed{0.4}.\]

\end{tcolorbox}

\subsection{Continuous conditioning and total
probability}\label{continuous-conditioning-and-total-probability}

This subsection extends the law of total probability to conditioning on
a continuous random variable.

Let \(X\) be a continuous random variable with density \(f_X(x)\). For
any event \(A\),
\[\mathbb{P}(A)=\int_{-\infty}^{\infty}\mathbb{P}(A\mid X=x)f_X(x)\,dx.\]
A shorthand notation is
\[\mathbb{P}(A)=\mathbb{E}[\mathbb{P}(A\mid X)],\] where
\(\mathbb{P}(A\mid X)\) is a random variable that is a function of
\(X\).

More generally, if \(X\) and \(Y\) have joint density or joint pmf, then
\[p_{X\mid Y}(x\mid y)=\frac{p_{X,Y}(x,y)}{p_Y(y)}\] whenever
\(p_Y(y)>0\), and the marginal density or pmf can be recovered by
\[p_X(x)=\int p_{X\mid Y}(x\mid y)p_Y(y)\,dy\] in the continuous case,
or \[p_X(x)=\sum_y p_{X\mid Y}(x\mid y)p_Y(y)\] in the discrete case.

\textbf{Example 4} (Tail probability for a sum of exponentials). Suppose
\(X\) and \(Y\) are independent exponential random variables with mean
\(1\), so \(f_X(x)=e^{-x}\) for \(x\ge0\). For \(z\ge0\), compute
\[\mathbb{P}(X+Y\ge z).\]

\begin{tcolorbox}[enhanced jigsaw, arc=.35mm, bottomrule=.15mm, bottomtitle=1mm, breakable, colback=white, colbacktitle=quarto-callout-note-color!10!white, colframe=quarto-callout-note-color-frame, coltitle=black, left=2mm, leftrule=.75mm, opacityback=0, opacitybacktitle=0.6, rightrule=.15mm, title=\textcolor{quarto-callout-note-color}{\faInfo}\hspace{0.5em}{Solution}, titlerule=0mm, toprule=.15mm, toptitle=1mm]

Condition on \(X=x\). Since \(X\) and \(Y\) are independent,
\[\mathbb{P}(X+Y\ge z\mid X=x)=\mathbb{P}(Y\ge z-x).\] For an
exponential random variable with mean \(1\),
\[\mathbb{P}(Y\ge a)=e^{-a}\quad\text{for }a\ge0.\] Thus
\[\mathbb{P}(Y\ge z-x)=
\begin{cases}
e^{-(z-x)}, & 0\le x\le z,\\
1, & x>z.
\end{cases}\] Therefore \[\begin{aligned}
\mathbb{P}(X+Y\ge z)
&=\int_0^\infty \mathbb{P}(X+Y\ge z\mid X=x)e^{-x}\,dx\\
&=\int_0^z e^{-(z-x)}e^{-x}\,dx+\int_z^\infty e^{-x}\,dx\\
&=\int_0^z e^{-z}\,dx+e^{-z}\\
&=ze^{-z}+e^{-z}.
\end{aligned}\] Hence
\[\boxed{\mathbb{P}(X+Y\ge z)=(z+1)e^{-z}},\qquad z\ge0.\] This agrees
with the fact that \(X+Y\sim\operatorname{Gamma}(2,1)\) in rate
parameterization.

\end{tcolorbox}

\textbf{Example 5} (A dependent example). Suppose
\(X\sim\operatorname{Uniform}(0,1)\) and, conditional on \(X=x\), the
random variable \(Y\) is uniform on \([0,x]\). Calculate
\(\mathbb{E}[Y]\).

\begin{tcolorbox}[enhanced jigsaw, arc=.35mm, bottomrule=.15mm, bottomtitle=1mm, breakable, colback=white, colbacktitle=quarto-callout-note-color!10!white, colframe=quarto-callout-note-color-frame, coltitle=black, left=2mm, leftrule=.75mm, opacityback=0, opacitybacktitle=0.6, rightrule=.15mm, title=\textcolor{quarto-callout-note-color}{\faInfo}\hspace{0.5em}{Solution}, titlerule=0mm, toprule=.15mm, toptitle=1mm]

Conditioning on \(X=x\), we have
\[Y\mid X=x\sim\operatorname{Uniform}(0,x),\] so
\[\mathbb{E}[Y\mid X=x]=\frac{x}{2}.\] Therefore
\[\mathbb{E}[Y\mid X]=\frac{X}{2}.\] Using the law of total expectation,
\[\mathbb{E}[Y]=\mathbb{E}[\mathbb{E}(Y\mid X)]=\mathbb{E}\left[\frac{X}{2}\right]
=\frac12\mathbb{E}[X]=\frac12\cdot\frac12=\frac14.\] Thus
\[\boxed{\mathbb{E}[Y]=\frac14.}\]

\end{tcolorbox}

\subsection{Memoryless property of the exponential
distribution}\label{memoryless-property-of-the-exponential-distribution}

This subsection studies an important example where conditioning does not
change the remaining lifetime distribution.

Let \(X\sim\operatorname{Exp}(\lambda)\), with density
\[f_X(t)=\lambda e^{-\lambda t},\qquad t\ge0.\] Then
\[\mathbb{P}(X\ge t)=e^{-\lambda t},\qquad t\ge0.\] For \(s,t\ge0\),
\[\begin{aligned}
\mathbb{P}(X\ge t+s\mid X>s)
&=\frac{\mathbb{P}(X\ge t+s, X>s)}{\mathbb{P}(X>s)}\\
&=\frac{\mathbb{P}(X\ge t+s)}{\mathbb{P}(X>s)}\\
&=\frac{e^{-\lambda(t+s)}}{e^{-\lambda s}}\\
&=e^{-\lambda t}.
\end{aligned}\] Thus
\[\boxed{\mathbb{P}(X\ge t+s\mid X>s)=\mathbb{P}(X\ge t).}\]

\begin{tcolorbox}[enhanced jigsaw, arc=.35mm, bottomrule=.15mm, bottomtitle=1mm, breakable, colback=white, colbacktitle=quarto-callout-tip-color!10!white, colframe=quarto-callout-tip-color-frame, coltitle=black, left=2mm, leftrule=.75mm, opacityback=0, opacitybacktitle=0.6, rightrule=.15mm, title=\textcolor{quarto-callout-tip-color}{\faLightbulb}\hspace{0.5em}{Interpretation}, titlerule=0mm, toprule=.15mm, toptitle=1mm]

If \(X\) is the lifetime of a device and the device has survived until
time \(s\), then the additional waiting time has the same distribution
as a brand-new lifetime. The exponential distribution has no memory.

\end{tcolorbox}

\textbf{Example 6} (Waiting for the next car). Cars pass a point on a
highway. The times between successive cars are independent exponential
random variables with mean \(m\). Suppose at a random time you stand at
the point on the highway. What is the mean time until the next car
passes?

\begin{tcolorbox}[enhanced jigsaw, arc=.35mm, bottomrule=.15mm, bottomtitle=1mm, breakable, colback=white, colbacktitle=quarto-callout-note-color!10!white, colframe=quarto-callout-note-color-frame, coltitle=black, left=2mm, leftrule=.75mm, opacityback=0, opacitybacktitle=0.6, rightrule=.15mm, title=\textcolor{quarto-callout-note-color}{\faInfo}\hspace{0.5em}{Solution}, titlerule=0mm, toprule=.15mm, toptitle=1mm]

An exponential interarrival time with mean \(m\) has rate
\[\lambda=\frac1m.\] Because of the memoryless property, the remaining
waiting time until the next car is again exponential with rate
\(\lambda\). Therefore its mean is \[\frac1\lambda=m.\] Thus the mean
time until the next car is \[\boxed{m}.\]

\end{tcolorbox}

\subsection{Mixed conditional
distributions}\label{mixed-conditional-distributions}

This subsection considers examples where one variable is discrete and
the other is continuous.

\textbf{Example 7} (Poisson--Exponential--Gamma example). Suppose
\(X\in\{0,1,2,\ldots\}\) is discrete and \(Y\ge0\) is continuous with
joint density/mass function
\[p_{X,Y}(x,y)=\frac{\lambda y^x e^{-(\lambda+1)y}}{x!},
\qquad x=0,1,2,\ldots,\quad y\ge0.\] Find the marginal distribution of
\(Y\), the conditional distribution of \(X\mid Y=y\), and the
conditional distribution of \(Y\mid X=x\).

\begin{tcolorbox}[enhanced jigsaw, arc=.35mm, bottomrule=.15mm, bottomtitle=1mm, breakable, colback=white, colbacktitle=quarto-callout-note-color!10!white, colframe=quarto-callout-note-color-frame, coltitle=black, left=2mm, leftrule=.75mm, opacityback=0, opacitybacktitle=0.6, rightrule=.15mm, title=\textcolor{quarto-callout-note-color}{\faInfo}\hspace{0.5em}{Solution}, titlerule=0mm, toprule=.15mm, toptitle=1mm]

First compute the marginal density of \(Y\): \[\begin{aligned}
p_Y(y)
&=\sum_{x=0}^{\infty}p_{X,Y}(x,y)\\
&=\sum_{x=0}^{\infty}\frac{\lambda y^x e^{-(\lambda+1)y}}{x!}\\
&=\lambda e^{-(\lambda+1)y}\sum_{x=0}^{\infty}\frac{y^x}{x!}\\
&=\lambda e^{-(\lambda+1)y}e^y\\
&=\lambda e^{-\lambda y},\qquad y\ge0.
\end{aligned}\] Thus \[\boxed{Y\sim\operatorname{Exp}(\lambda).}\]

Next, \[\begin{aligned}
p_{X\mid Y}(x\mid y)
&=\frac{p_{X,Y}(x,y)}{p_Y(y)}\\
&=\frac{\lambda y^x e^{-(\lambda+1)y}/x!}{\lambda e^{-\lambda y}}\\
&=\frac{y^x e^{-y}}{x!}.
\end{aligned}\] Therefore
\[\boxed{X\mid Y=y\sim\operatorname{Poisson}(y).}\]

Finally, to identify \(Y\mid X=x\), treat \(x\) as fixed and keep only
factors depending on \(y\): \[p_{Y\mid X}(y\mid x)\propto p_{X,Y}(x,y)
\propto y^x e^{-(\lambda+1)y},\qquad y\ge0.\] This is a Gamma density
with shape \(x+1\) and rate \(\lambda+1\). Hence
\[\boxed{Y\mid X=x\sim\operatorname{Gamma}(x+1,\lambda+1)}\] where the
second parameter is the rate.

\end{tcolorbox}

\section{Conditional Expectations}\label{conditional-expectations}

This section introduces conditional expectation as the expected value
computed under a conditional distribution.

\subsection{Definition}\label{definition-2}

This subsection defines conditional expectation for both discrete and
continuous random variables.

\textbf{Definition 8} (Conditional expectation). Let \(X\) and \(Y\) be
random variables.

If \(X\) is discrete, then \[\mathbb{E}[X\mid Y=y]
=\sum_x x\,p_{X\mid Y}(x\mid y).\] If \(X\) is continuous, then
\[\mathbb{E}[X\mid Y=y]
=\int_{-\infty}^{\infty}x\,p_{X\mid Y}(x\mid y)\,dx.\]

For each possible value \(y\) of \(Y\), the expression
\(\mathbb{E}[X\mid Y=y]\) is a number. Therefore \(\mathbb{E}[X\mid Y]\)
is a random variable that is a function of \(Y\).

\begin{tcolorbox}[enhanced jigsaw, arc=.35mm, bottomrule=.15mm, bottomtitle=1mm, breakable, colback=white, colbacktitle=quarto-callout-tip-color!10!white, colframe=quarto-callout-tip-color-frame, coltitle=black, left=2mm, leftrule=.75mm, opacityback=0, opacitybacktitle=0.6, rightrule=.15mm, title=\textcolor{quarto-callout-tip-color}{\faLightbulb}\hspace{0.5em}{How to read \(\mathbb{E}[X\mid Y]\)}, titlerule=0mm, toprule=.15mm, toptitle=1mm]

The conditional expectation \(\mathbb{E}[X\mid Y]\) is the best
prediction of \(X\) after observing \(Y\), when prediction is measured
by squared error. In this course, the most important point is that it is
a random variable determined by \(Y\).

\end{tcolorbox}

If \(X\) and \(Y\) are independent, then conditioning on \(Y\) does not
change the distribution of \(X\). Hence
\[\mathbb{E}[X\mid Y=y]=\mathbb{E}[X],\] and therefore
\[\mathbb{E}[X\mid Y]=\mathbb{E}[X].\] Also, if \(X\) and \(Y\) are
independent and the expectations exist, then
\[\mathbb{E}[XY]=\mathbb{E}[X]\mathbb{E}[Y].\]

\subsection{Law of total expectation}\label{law-of-total-expectation}

This subsection states the most important computational identity for
conditional expectation.

\textbf{Theorem 9} (Law of total expectation). \emph{If the relevant
expectations exist, then
\[\boxed{\mathbb{E}[\mathbb{E}[X\mid Y]]=\mathbb{E}[X].}\] More
generally,
\[\boxed{\mathbb{E}[g(X,Y)]=\mathbb{E}\big[\mathbb{E}[g(X,Y)\mid X]\big]}\]
for any measurable function \(g\) for which the expectations exist.}

\begin{proof}
\emph{Proof in the discrete case.} Assume \(X\) and \(Y\) are discrete.
Then \[\begin{aligned}
\mathbb{E}_Y[\mathbb{E}_X(X\mid Y)]
&=\sum_y \mathbb{E}[X\mid Y=y]\mathbb{P}(Y=y)\\
&=\sum_y\left(\sum_x x\mathbb{P}(X=x\mid Y=y)\right)\mathbb{P}(Y=y)\\
&=\sum_x\sum_y x\mathbb{P}(X=x\mid Y=y)\mathbb{P}(Y=y)\\
&=\sum_x\sum_y x\mathbb{P}(X=x,Y=y)\\
&=\sum_x x\mathbb{P}(X=x)\\
&=\mathbb{E}[X].
\end{aligned}\]~◻
\end{proof}

\begin{proof}
\emph{Proof in the continuous case.} Assume \(X\) and \(Y\) are
continuous. Then \[\begin{aligned}
\mathbb{E}_Y[\mathbb{E}_X(X\mid Y)]
&=\int_y \mathbb{E}[X\mid Y=y]p_Y(y)\,dy\\
&=\int_y\int_x x p_{X\mid Y}(x\mid y)p_Y(y)\,dx\,dy\\
&=\int_y\int_x x p_{X,Y}(x,y)\,dx\,dy\\
&=\int_x x\left(\int_y p_{X,Y}(x,y)\,dy\right)\,dx\\
&=\int_x x p_X(x)\,dx\\
&=\mathbb{E}[X].
\end{aligned}\]~◻
\end{proof}

\textbf{Example 10} (Unit disk: uncorrelated but not independent). Let
\((X,Y)\) be uniformly distributed over the unit disk
\[D=\{(x,y):x^2+y^2\le1\}.\] The joint density is
\[f_{X,Y}(x,y)=\frac1\pi,\qquad (x,y)\in D.\] Are \(X\) and \(Y\)
uncorrelated?

\begin{tcolorbox}[enhanced jigsaw, arc=.35mm, bottomrule=.15mm, bottomtitle=1mm, breakable, colback=white, colbacktitle=quarto-callout-note-color!10!white, colframe=quarto-callout-note-color-frame, coltitle=black, left=2mm, leftrule=.75mm, opacityback=0, opacitybacktitle=0.6, rightrule=.15mm, title=\textcolor{quarto-callout-note-color}{\faInfo}\hspace{0.5em}{Solution}, titlerule=0mm, toprule=.15mm, toptitle=1mm]

From symmetry, \[\mathbb{E}[X]=0,\qquad \mathbb{E}[Y]=0.\] We need to
compute
\[\operatorname{Cov}(X,Y)=\mathbb{E}[XY]-\mathbb{E}[X]\mathbb{E}[Y]=\mathbb{E}[XY].\]
Condition on \(Y\). For a fixed \(Y=y\), the possible values of \(X\)
are \[-\sqrt{1-y^2}\le X\le \sqrt{1-y^2}.\] Moreover,
\[X\mid Y=y\sim \operatorname{Uniform}\left(-\sqrt{1-y^2},\sqrt{1-y^2}\right).\]
Thus \[\mathbb{E}[X\mid Y=y]=0.\] Therefore
\[\mathbb{E}[XY]=\mathbb{E}\big[\mathbb{E}[XY\mid Y]\big]
=\mathbb{E}\big[Y\mathbb{E}[X\mid Y]\big]
=\mathbb{E}[Y\cdot0]=0.\] Hence \[\boxed{\operatorname{Cov}(X,Y)=0.}\]
So \(X\) and \(Y\) are uncorrelated. However, they are not independent
because the conditional support of \(X\mid Y=y\) depends on \(y\).

\end{tcolorbox}

\textbf{Example 11} (Computing covariance by conditioning). Suppose
\[X\sim\operatorname{Uniform}(1,2),\] and conditional on \(X=x\),
\[Y\mid X=x\sim\operatorname{Exp}(x),\] where \(x\) is the rate
parameter. Find \(\operatorname{Cov}(X,Y)\).

\begin{tcolorbox}[enhanced jigsaw, arc=.35mm, bottomrule=.15mm, bottomtitle=1mm, breakable, colback=white, colbacktitle=quarto-callout-note-color!10!white, colframe=quarto-callout-note-color-frame, coltitle=black, left=2mm, leftrule=.75mm, opacityback=0, opacitybacktitle=0.6, rightrule=.15mm, title=\textcolor{quarto-callout-note-color}{\faInfo}\hspace{0.5em}{Solution}, titlerule=0mm, toprule=.15mm, toptitle=1mm]

We use
\[\operatorname{Cov}(X,Y)=\mathbb{E}[XY]-\mathbb{E}[X]\mathbb{E}[Y].\]
First, \[\mathbb{E}[X]=\frac{1+2}{2}=\frac32.\] Because
\(Y\mid X=x\sim\operatorname{Exp}(x)\) with rate \(x\),
\[\mathbb{E}[Y\mid X=x]=\frac1x.\] Thus \[\mathbb{E}[Y\mid X]=\frac1X.\]
Using total expectation,
\[\mathbb{E}[Y]=\mathbb{E}\left[\frac1X\right]=\int_1^2\frac1x\,dx=\log2.\]
Next, \[\mathbb{E}[XY]=\mathbb{E}[\mathbb{E}[XY\mid X]]
=\mathbb{E}[X\mathbb{E}(Y\mid X)]
=\mathbb{E}\left[X\cdot\frac1X\right]=1.\] Therefore
\[\boxed{\operatorname{Cov}(X,Y)=1-\frac32\log2.}\] This value is
negative because large \(X\) implies a larger exponential rate and
therefore a smaller conditional mean for \(Y\).

\end{tcolorbox}

\subsection{Random sums of random
variables}\label{random-sums-of-random-variables}

This subsection uses conditioning to evaluate expectations when the
number of terms is itself random.

\textbf{Example 12} (Random sum). Let \(N,X_1,X_2,\\ldots\) be
independent, where \(X_i\) are IID with \[\mathbb{E}[X_i]=\mu.\] Define
\[Y=\sum_{i=1}^N X_i.\] For example, \(N\) could be the number of
insurance claims in a month and \(X_i\) the size of the \(i\)-th claim.
Find \(\mathbb{E}[Y]\).

\begin{tcolorbox}[enhanced jigsaw, arc=.35mm, bottomrule=.15mm, bottomtitle=1mm, breakable, colback=white, colbacktitle=quarto-callout-note-color!10!white, colframe=quarto-callout-note-color-frame, coltitle=black, left=2mm, leftrule=.75mm, opacityback=0, opacitybacktitle=0.6, rightrule=.15mm, title=\textcolor{quarto-callout-note-color}{\faInfo}\hspace{0.5em}{Solution}, titlerule=0mm, toprule=.15mm, toptitle=1mm]

Condition on \(N=n\). Then the number of terms is fixed:
\[Y\mid N=n=\sum_{i=1}^n X_i.\] Since \(N\) is independent of the
\(X_i\)'s, \[\mathbb{E}[Y\mid N=n]
=\mathbb{E}\left[\sum_{i=1}^n X_i\right]
=\sum_{i=1}^n\mathbb{E}[X_i]
=n\mu.\] Therefore \[\mathbb{E}[Y\mid N]=N\mu.\] Using the law of total
expectation,
\[\mathbb{E}[Y]=\mathbb{E}[\mathbb{E}(Y\mid N)]=\mathbb{E}[N\mu]=\mu\mathbb{E}[N].\]
Thus
\[\boxed{\mathbb{E}\left[\sum_{i=1}^N X_i\right]=\mathbb{E}[N]\mathbb{E}[X_1].}\]
This identity is a basic form of Wald's equation.

\end{tcolorbox}

\textbf{Practice Problem 13} (Variance of a random sum). In the random
sum example, assume also that \[\operatorname{Var}(X_i)=\sigma^2,\] and
\(N\) is independent of the \(X_i\)'s. Use the law of total variance to
show that
\[\operatorname{Var}(Y)=\sigma^2\mathbb{E}[N]+\mu^2\operatorname{Var}(N).\]

\begin{tcolorbox}[enhanced jigsaw, arc=.35mm, bottomrule=.15mm, bottomtitle=1mm, breakable, colback=white, colbacktitle=quarto-callout-note-color!10!white, colframe=quarto-callout-note-color-frame, coltitle=black, left=2mm, leftrule=.75mm, opacityback=0, opacitybacktitle=0.6, rightrule=.15mm, title=\textcolor{quarto-callout-note-color}{\faInfo}\hspace{0.5em}{Solution}, titlerule=0mm, toprule=.15mm, toptitle=1mm]

Condition on \(N=n\). Then
\[\operatorname{Var}(Y\mid N=n)=\operatorname{Var}\left(\sum_{i=1}^n X_i\right)=n\sigma^2,\]
so \[\operatorname{Var}(Y\mid N)=N\sigma^2.\] Also,
\[\mathbb{E}[Y\mid N]=N\mu.\] By the law of total variance,
\[\begin{aligned}
\operatorname{Var}(Y)
&=\mathbb{E}[\operatorname{Var}(Y\mid N)]+\operatorname{Var}(\mathbb{E}[Y\mid N])\\
&=\mathbb{E}[N\sigma^2]+\operatorname{Var}(N\mu)\\
&=\sigma^2\mathbb{E}[N]+\mu^2\operatorname{Var}(N).
\end{aligned}\] Thus
\[\boxed{\operatorname{Var}(Y)=\sigma^2\mathbb{E}[N]+\mu^2\operatorname{Var}(N).}\]

\end{tcolorbox}

\subsection{Law of total variance}\label{law-of-total-variance}

This subsection decomposes total variation into average within-group
variation and between-group variation.

\textbf{Theorem 14} (Law of total variance). \emph{If the relevant
second moments exist, then
\[\boxed{\operatorname{Var}(Y)=\mathbb{E}[\operatorname{Var}(Y\mid X)]+\operatorname{Var}(\mathbb{E}[Y\mid X]).}\]}

\begin{proof}
\emph{Proof.} Start from
\[\operatorname{Var}(Y)=\mathbb{E}[Y^2]-(\mathbb{E}[Y])^2.\] Using total
expectation, \[\mathbb{E}[Y^2]=\mathbb{E}[\mathbb{E}(Y^2\mid X)]\] and
\[\mathbb{E}[Y]=\mathbb{E}[\mathbb{E}(Y\mid X)].\] For fixed \(X\),
\[\operatorname{Var}(Y\mid X)=\mathbb{E}(Y^2\mid X)-\big(\mathbb{E}(Y\mid X)\big)^2.\]
Thus
\[\mathbb{E}(Y^2\mid X)=\operatorname{Var}(Y\mid X)+\big(\mathbb{E}(Y\mid X)\big)^2.\]
Taking expectations,
\[\mathbb{E}[Y^2]=\mathbb{E}[\operatorname{Var}(Y\mid X)]+\mathbb{E}\left[\big(\mathbb{E}(Y\mid X)\big)^2\right].\]
Therefore \[\begin{aligned}
\operatorname{Var}(Y)
&=\mathbb{E}[\operatorname{Var}(Y\mid X)]+\mathbb{E}\left[\big(\mathbb{E}(Y\mid X)\big)^2\right]-\left(\mathbb{E}[\mathbb{E}(Y\mid X)]\right)^2\\
&=\mathbb{E}[\operatorname{Var}(Y\mid X)]+\operatorname{Var}(\mathbb{E}[Y\mid X]).
\end{aligned}\]~◻
\end{proof}

\begin{tcolorbox}[enhanced jigsaw, arc=.35mm, bottomrule=.15mm, bottomtitle=1mm, breakable, colback=white, colbacktitle=quarto-callout-tip-color!10!white, colframe=quarto-callout-tip-color-frame, coltitle=black, left=2mm, leftrule=.75mm, opacityback=0, opacitybacktitle=0.6, rightrule=.15mm, title=\textcolor{quarto-callout-tip-color}{\faLightbulb}\hspace{0.5em}{Interpretation}, titlerule=0mm, toprule=.15mm, toptitle=1mm]

The identity
\[\operatorname{Var}(Y)=\mathbb{E}[\operatorname{Var}(Y\mid X)]+\operatorname{Var}(\mathbb{E}[Y\mid X])\]
says:
\[\text{total variation}=\text{average conditional variation}+\text{variation of conditional means}.\]

\end{tcolorbox}

\subsection{A covariance identity}\label{a-covariance-identity}

This subsection records a useful projection identity for covariance.

Let \(g(X)\) be a function of \(X\) and \(h(Y)\) be a function of \(Y\).
Then
\[\boxed{\operatorname{Cov}(g(X),h(Y))=\operatorname{Cov}(g(X),\mathbb{E}[h(Y)\mid X]).}\]
Indeed, \[\begin{aligned}
\operatorname{Cov}(g(X),h(Y))
&=\mathbb{E}[g(X)h(Y)]-\mathbb{E}[g(X)]\mathbb{E}[h(Y)]\\
&=\mathbb{E}\left[\mathbb{E}[g(X)h(Y)\mid X]\right]-\mathbb{E}[g(X)]\mathbb{E}\left[\mathbb{E}[h(Y)\mid X]\right]\\
&=\mathbb{E}\left[g(X)\mathbb{E}[h(Y)\mid X]\right]-\mathbb{E}[g(X)]\mathbb{E}\left[\mathbb{E}[h(Y)\mid X]\right]\\
&=\operatorname{Cov}(g(X),\mathbb{E}[h(Y)\mid X]).
\end{aligned}\]

\begin{tcolorbox}[enhanced jigsaw, arc=.35mm, bottomrule=.15mm, bottomtitle=1mm, breakable, colback=white, colbacktitle=quarto-callout-tip-color!10!white, colframe=quarto-callout-tip-color-frame, coltitle=black, left=2mm, leftrule=.75mm, opacityback=0, opacitybacktitle=0.6, rightrule=.15mm, title=\textcolor{quarto-callout-tip-color}{\faLightbulb}\hspace{0.5em}{Projection viewpoint}, titlerule=0mm, toprule=.15mm, toptitle=1mm]

The random variable \(\mathbb{E}[h(Y)\mid X]\) may be viewed as the part
of \(h(Y)\) that can be predicted from \(X\). Therefore, when computing
covariance with a function of \(X\), we can replace \(h(Y)\) by its
conditional expectation given \(X\).

\end{tcolorbox}

\subsection{Binomial--uniform example and Bayesian
updating}\label{binomialuniform-example-and-bayesian-updating}

This subsection connects conditional expectation, total variance, and
Bayesian inference.

\textbf{Example 15} (Binomial--uniform). Suppose
\[X\mid Y\sim\operatorname{Binomial}(n,Y),
\qquad
Y\sim\operatorname{Uniform}(0,1).\] Find \(\mathbb{E}[X]\) and
\(\operatorname{Var}(X)\). Then determine the conditional distribution
of \(Y\mid X=x\).

\begin{tcolorbox}[enhanced jigsaw, arc=.35mm, bottomrule=.15mm, bottomtitle=1mm, breakable, colback=white, colbacktitle=quarto-callout-note-color!10!white, colframe=quarto-callout-note-color-frame, coltitle=black, left=2mm, leftrule=.75mm, opacityback=0, opacitybacktitle=0.6, rightrule=.15mm, title=\textcolor{quarto-callout-note-color}{\faInfo}\hspace{0.5em}{Solution}, titlerule=0mm, toprule=.15mm, toptitle=1mm]

Given \(Y\), the conditional mean and variance of \(X\) are
\[\mathbb{E}[X\mid Y]=nY,
\qquad
\operatorname{Var}(X\mid Y)=nY(1-Y).\] By the law of total expectation,
\[\mathbb{E}[X]=\mathbb{E}[\mathbb{E}(X\mid Y)]=\mathbb{E}[nY]=n\mathbb{E}[Y]=\frac n2.\]

By the law of total variance, \[\begin{aligned}
\operatorname{Var}(X)
&=\mathbb{E}[\operatorname{Var}(X\mid Y)]+\operatorname{Var}(\mathbb{E}[X\mid Y])\\
&=\mathbb{E}[nY(1-Y)]+\operatorname{Var}(nY)\\
&=n\left(\mathbb{E}[Y]-\mathbb{E}[Y^2]\right)+n^2\operatorname{Var}(Y).
\end{aligned}\] For \(Y\sim\operatorname{Uniform}(0,1)\),
\[\mathbb{E}[Y]=\frac12,
\qquad
\mathbb{E}[Y^2]=\frac13,
\qquad
\operatorname{Var}(Y)=\frac1{12}.\] Thus
\[\operatorname{Var}(X)=n\left(\frac12-\frac13\right)+n^2\cdot\frac1{12}
=\frac n6+\frac{n^2}{12}.\] Therefore \[\boxed{\mathbb{E}[X]=\frac n2,
\qquad
\operatorname{Var}(X)=\frac n6+\frac{n^2}{12}.}\]

Now find \(Y\mid X=x\). Since \(Y\sim\operatorname{Uniform}(0,1)\), its
density is constant on \([0,1]\). Therefore \[\begin{aligned}
p_{Y\mid X}(y\mid x)
&\propto p_{X\mid Y}(x\mid y)p_Y(y)\\
&\propto \binom{n}{x}y^x(1-y)^{n-x}\cdot 1\\
&\propto y^x(1-y)^{n-x},\qquad 0<y<1.
\end{aligned}\] This is the kernel of a beta distribution with
parameters \[\alpha=x+1,
\qquad
\beta=n-x+1.\] Thus
\[\boxed{Y\mid X=x\sim\operatorname{Beta}(x+1,n-x+1).}\] This is a
Bayesian update: the prior \(Y\sim\operatorname{Beta}(1,1)\) becomes the
posterior \(Y\mid X=x\sim\operatorname{Beta}(x+1,n-x+1)\) after
observing \(x\) successes in \(n\) trials.

\end{tcolorbox}

\textbf{Practice Problem 16} (Posterior mean). In the binomial--uniform
example, compute \(\mathbb{E}[Y\mid X=x]\).

\begin{tcolorbox}[enhanced jigsaw, arc=.35mm, bottomrule=.15mm, bottomtitle=1mm, breakable, colback=white, colbacktitle=quarto-callout-note-color!10!white, colframe=quarto-callout-note-color-frame, coltitle=black, left=2mm, leftrule=.75mm, opacityback=0, opacitybacktitle=0.6, rightrule=.15mm, title=\textcolor{quarto-callout-note-color}{\faInfo}\hspace{0.5em}{Solution}, titlerule=0mm, toprule=.15mm, toptitle=1mm]

If \[Y\mid X=x\sim\operatorname{Beta}(x+1,n-x+1),\] then the mean of a
\(\operatorname{Beta}(\alpha,\beta)\) random variable is
\[\frac{\alpha}{\alpha+\beta}.\] Therefore \[\mathbb{E}[Y\mid X=x]
=\frac{x+1}{(x+1)+(n-x+1)}
=\frac{x+1}{n+2}.\] Thus
\[\boxed{\mathbb{E}[Y\mid X=x]=\frac{x+1}{n+2}.}\]

\end{tcolorbox}

\section{Applications: Weighting and Missing
Data}\label{applications-weighting-and-missing-data}

This section shows how conditional expectation explains important
weighting methods in statistics and data analysis.

\subsection{Inverse probability weighting for missing
data}\label{inverse-probability-weighting-for-missing-data}

This subsection studies a missing-data problem where income is sometimes
unobserved.

\textbf{Example 17} (Missing data and IPW). Consider a survey with two
variables: \[X=\text{age of a participant},
\qquad
Y=\text{income of a participant}.\] We want \[\mu=\mathbb{E}[Y].\]
However, \(Y\) may be missing because some people refuse to provide
income information. Let \(R\) be the response indicator:
\[R=1 \quad\text{if }Y\text{ is observed},
\qquad
R=0 \quad\text{if }Y\text{ is missing}.\] Assume \[R\perp Y\mid X,\]
which means the response probability depends only on \(X\):
\[\mathbb{P}(R=1\mid X,Y)=\mathbb{P}(R=1\mid X)=\pi(X).\] Assume
\(\pi(X)\) is known. Show that \[W=\frac{RY}{\pi(X)}\] satisfies
\[\mathbb{E}[W]=\mathbb{E}[Y].\]

\begin{tcolorbox}[enhanced jigsaw, arc=.35mm, bottomrule=.15mm, bottomtitle=1mm, breakable, colback=white, colbacktitle=quarto-callout-note-color!10!white, colframe=quarto-callout-note-color-frame, coltitle=black, left=2mm, leftrule=.75mm, opacityback=0, opacitybacktitle=0.6, rightrule=.15mm, title=\textcolor{quarto-callout-note-color}{\faInfo}\hspace{0.5em}{Solution}, titlerule=0mm, toprule=.15mm, toptitle=1mm]

We compute using conditional expectation. First,
\[\mathbb{E}[W]=\mathbb{E}\left[\frac{RY}{\pi(X)}\right]
=\mathbb{E}\left[\mathbb{E}\left(\frac{RY}{\pi(X)}\mid X\right)\right].\]
Because \(1/\pi(X)\) is determined by \(X\),
\[\mathbb{E}\left(\frac{RY}{\pi(X)}\mid X\right)
=\frac1{\pi(X)}\mathbb{E}[RY\mid X].\] Using the conditional
independence assumption \(R\perp Y\mid X\),
\[\mathbb{E}[RY\mid X]=\mathbb{E}[R\mid X]\mathbb{E}[Y\mid X].\] But
\[\mathbb{E}[R\mid X]=\mathbb{P}(R=1\mid X)=\pi(X).\] Hence
\[\mathbb{E}\left(\frac{RY}{\pi(X)}\mid X\right)
=\frac1{\pi(X)}\pi(X)\mathbb{E}[Y\mid X]
=\mathbb{E}[Y\mid X].\] Therefore
\[\mathbb{E}[W]=\mathbb{E}[\mathbb{E}(Y\mid X)]=\mathbb{E}[Y].\] Thus
\[\boxed{\mathbb{E}\left[\frac{RY}{\pi(X)}\right]=\mathbb{E}[Y].}\]

\end{tcolorbox}

In practice, if we observe IID data, we estimate \(\mu=\mathbb{E}[Y]\)
by the inverse probability weighting estimator
\[\widehat\mu_{\operatorname{IPW}}
=\frac1n\sum_{i=1}^n \frac{R_iY_i}{\pi(X_i)}.\] This estimator upweights
observed responses that had a smaller probability of being observed.

\subsection{Survey sampling and importance
weighting}\label{survey-sampling-and-importance-weighting}

This subsection explains importance weighting when the sample
distribution differs from the population distribution.

\textbf{Example 18} (Survey sampling). A city has three districts \(A\),
\(B\), and \(C\). The population proportions are
\[\mathbb{P}_{\text{pop}}(X=A)=0.6,
\qquad
\mathbb{P}_{\text{pop}}(X=B)=0.3,
\qquad
\mathbb{P}_{\text{pop}}(X=C)=0.1.\] Let \(Y\) be income. The target
average income is
\[\mu=0.6\mathbb{E}[Y\mid X=A]+0.3\mathbb{E}[Y\mid X=B]+0.1\mathbb{E}[Y\mid X=C].\]
However, the survey samples the same number of individuals from each
district, so in the sample
\[\mathbb{P}_{\text{sample}}(X=A)=\mathbb{P}_{\text{sample}}(X=B)=\mathbb{P}_{\text{sample}}(X=C)=\frac13.\]
Construct a quantity \(Z=g(X,Y)\) such that
\[\mathbb{E}_{\text{sample}}[Z]=\mu.\]

\begin{tcolorbox}[enhanced jigsaw, arc=.35mm, bottomrule=.15mm, bottomtitle=1mm, breakable, colback=white, colbacktitle=quarto-callout-note-color!10!white, colframe=quarto-callout-note-color-frame, coltitle=black, left=2mm, leftrule=.75mm, opacityback=0, opacitybacktitle=0.6, rightrule=.15mm, title=\textcolor{quarto-callout-note-color}{\faInfo}\hspace{0.5em}{Solution}, titlerule=0mm, toprule=.15mm, toptitle=1mm]

We weight each observation by
\[\frac{\text{population probability}}{\text{sample probability}}.\]
Thus define \[\begin{aligned}
Z
&=\frac{0.6}{1/3}\mathbbm{1}(X=A)Y
+\frac{0.3}{1/3}\mathbbm{1}(X=B)Y
+\frac{0.1}{1/3}\mathbbm{1}(X=C)Y\\
&=1.8\mathbbm{1}(X=A)Y+0.9\mathbbm{1}(X=B)Y+0.3\mathbbm{1}(X=C)Y.
\end{aligned}\] Then \[\begin{aligned}
\mathbb{E}_{\text{sample}}[Z]
&=1.8\mathbb{E}[\mathbbm{1}(X=A)Y]+0.9\mathbb{E}[\mathbbm{1}(X=B)Y]+0.3\mathbb{E}[\mathbbm{1}(X=C)Y]\\
&=1.8\mathbb{P}_{\text{sample}}(X=A)\mathbb{E}[Y\mid X=A]\\
&\quad+0.9\mathbb{P}_{\text{sample}}(X=B)\mathbb{E}[Y\mid X=B]\\
&\quad+0.3\mathbb{P}_{\text{sample}}(X=C)\mathbb{E}[Y\mid X=C]\\
&=1.8\cdot\frac13\mathbb{E}[Y\mid X=A]+0.9\cdot\frac13\mathbb{E}[Y\mid X=B]+0.3\cdot\frac13\mathbb{E}[Y\mid X=C]\\
&=0.6\mathbb{E}[Y\mid X=A]+0.3\mathbb{E}[Y\mid X=B]+0.1\mathbb{E}[Y\mid X=C]\\
&=\mu.
\end{aligned}\] Thus
\[\boxed{Z=1.8\mathbbm{1}(X=A)Y+0.9\mathbbm{1}(X=B)Y+0.3\mathbbm{1}(X=C)Y}\]
has the desired property.

\end{tcolorbox}

\begin{tcolorbox}[enhanced jigsaw, arc=.35mm, bottomrule=.15mm, bottomtitle=1mm, breakable, colback=white, colbacktitle=quarto-callout-tip-color!10!white, colframe=quarto-callout-tip-color-frame, coltitle=black, left=2mm, leftrule=.75mm, opacityback=0, opacitybacktitle=0.6, rightrule=.15mm, title=\textcolor{quarto-callout-tip-color}{\faLightbulb}\hspace{0.5em}{Statistical meaning}, titlerule=0mm, toprule=.15mm, toptitle=1mm]

When the sampling design does not match the target population,
conditional expectation tells us how to reweight the sample so that the
weighted average targets the correct population quantity.

\end{tcolorbox}

\section{Summary}\label{summary-3}

This section summarizes the main formulas that should be remembered from
conditional distributions and conditional expectations.

{\def\LTcaptype{none} % do not increment counter
\begin{longtable}[]{@{}ll@{}}
\toprule\noalign{}
\textbf{Concept} & \textbf{Formula / message} \\
\midrule\noalign{}
\endhead
\bottomrule\noalign{}
\endlastfoot
Conditional pmf/pdf &
\(\displaystyle p_{X\mid Y}(x\mid y)=\frac{p_{X,Y}(x,y)}{p_Y(y)}\) \\
Law of total probability &
\(\displaystyle p_X(x)=\sum_y p_{X\mid Y}(x\mid y)p_Y(y)\) or
\(\displaystyle p_X(x)=\int p_{X\mid Y}(x\mid y)p_Y(y)\,dy\) \\
Conditional expectation &
\(\displaystyle \mathbb{E}[X\mid Y=y]=\sum_x x p_{X\mid Y}(x\mid y)\) or
\(\displaystyle \int x p_{X\mid Y}(x\mid y)\,dx\) \\
Law of total expectation &
\(\displaystyle \mathbb{E}[\mathbb{E}[X\mid Y]]=\mathbb{E}[X]\) \\
Law of total variance &
\(\displaystyle \operatorname{Var}(Y)=\mathbb{E}[\operatorname{Var}(Y\mid X)]+\operatorname{Var}(\mathbb{E}[Y\mid X])\) \\
Random sum & If \(Y=\sum_{i=1}^N X_i\), with \(N\) independent of IID
\(X_i\) and \(\mathbb{E}[X_i]=\mu\), then
\(\displaystyle \mathbb{E}[Y]=\mu\mathbb{E}[N]\). \\
Memoryless exponential & If \(X\sim\operatorname{Exp}(\lambda)\), then
\(\displaystyle \mathbb{P}(X\ge t+s\mid X>s)=\mathbb{P}(X\ge t)\). \\
IPW & If \(R\perp Y\mid X\) and \(\mathbb{P}(R=1\mid X)=\pi(X)\), then
\(\displaystyle \mathbb{E}\left[\frac{RY}{\pi(X)}\right]=\mathbb{E}[Y]\). \\
\end{longtable}
}

\section{Additional Practice
Problems}\label{additional-practice-problems-1}

This final section gives extra practice problems with full solutions.

\textbf{Practice Problem 19} (Conditional expectation from a table).
Suppose \(X,Y\in\{0,1\}\) have joint pmf \[\begin{array}{c|cc}
 & Y=0 & Y=1\\
\hline
X=0 & 0.2 & 0.3\\
X=1 & 0.1 & 0.4
\end{array}\] Find \(\mathbb{E}[X\mid Y=1]\) and \(\mathbb{E}[X]\).

\begin{tcolorbox}[enhanced jigsaw, arc=.35mm, bottomrule=.15mm, bottomtitle=1mm, breakable, colback=white, colbacktitle=quarto-callout-note-color!10!white, colframe=quarto-callout-note-color-frame, coltitle=black, left=2mm, leftrule=.75mm, opacityback=0, opacitybacktitle=0.6, rightrule=.15mm, title=\textcolor{quarto-callout-note-color}{\faInfo}\hspace{0.5em}{Solution}, titlerule=0mm, toprule=.15mm, toptitle=1mm]

First, \[\mathbb{P}(Y=1)=0.3+0.4=0.7.\] Thus
\[\mathbb{P}(X=1\mid Y=1)=\frac{\mathbb{P}(X=1,Y=1)}{\mathbb{P}(Y=1)}=\frac{0.4}{0.7}=\frac47.\]
Since \(X\) is Bernoulli conditional on \(Y=1\),
\[\boxed{\mathbb{E}[X\mid Y=1]=\frac47.}\] Also
\[\mathbb{P}(X=1)=0.1+0.4=0.5,\] so \[\boxed{\mathbb{E}[X]=0.5.}\]

\end{tcolorbox}

\textbf{Practice Problem 20} (Total expectation with a Poisson mixture).
Suppose \[X\mid\Lambda=\lambda\sim\operatorname{Poisson}(\lambda)\] and
\(\mathbb{E}[\Lambda]\) exists. Find \(\mathbb{E}[X]\).

\begin{tcolorbox}[enhanced jigsaw, arc=.35mm, bottomrule=.15mm, bottomtitle=1mm, breakable, colback=white, colbacktitle=quarto-callout-note-color!10!white, colframe=quarto-callout-note-color-frame, coltitle=black, left=2mm, leftrule=.75mm, opacityback=0, opacitybacktitle=0.6, rightrule=.15mm, title=\textcolor{quarto-callout-note-color}{\faInfo}\hspace{0.5em}{Solution}, titlerule=0mm, toprule=.15mm, toptitle=1mm]

For a Poisson random variable with rate \(\lambda\),
\[\mathbb{E}[X\mid\Lambda=\lambda]=\lambda.\] Thus
\[\mathbb{E}[X\mid\Lambda]=\Lambda.\] By total expectation,
\[\mathbb{E}[X]=\mathbb{E}[\mathbb{E}(X\mid\Lambda)]=\mathbb{E}[\Lambda].\]
So \[\boxed{\mathbb{E}[X]=\mathbb{E}[\Lambda].}\]

\end{tcolorbox}

\textbf{Practice Problem 21} (Total variance with a Poisson mixture). In
the previous problem, find \(\operatorname{Var}(X)\) in terms of
\(\mathbb{E}[\Lambda]\) and \(\operatorname{Var}(\Lambda)\).

\begin{tcolorbox}[enhanced jigsaw, arc=.35mm, bottomrule=.15mm, bottomtitle=1mm, breakable, colback=white, colbacktitle=quarto-callout-note-color!10!white, colframe=quarto-callout-note-color-frame, coltitle=black, left=2mm, leftrule=.75mm, opacityback=0, opacitybacktitle=0.6, rightrule=.15mm, title=\textcolor{quarto-callout-note-color}{\faInfo}\hspace{0.5em}{Solution}, titlerule=0mm, toprule=.15mm, toptitle=1mm]

For \(X\mid\Lambda=\lambda\sim\operatorname{Poisson}(\lambda)\),
\[\mathbb{E}[X\mid\Lambda]=\Lambda,
\qquad
\operatorname{Var}(X\mid\Lambda)=\Lambda.\] Therefore, by total
variance, \[\begin{aligned}
\operatorname{Var}(X)
&=\mathbb{E}[\operatorname{Var}(X\mid\Lambda)]+\operatorname{Var}(\mathbb{E}[X\mid\Lambda])\\
&=\mathbb{E}[\Lambda]+\operatorname{Var}(\Lambda).
\end{aligned}\] Thus
\[\boxed{\operatorname{Var}(X)=\mathbb{E}[\Lambda]+\operatorname{Var}(\Lambda).}\]

\end{tcolorbox}

\textbf{Practice Problem 22} (Memoryless calculation). Let
\(X\sim\operatorname{Exp}(2)\). Compute \[\mathbb{P}(X>7\mid X>3).\]

\begin{tcolorbox}[enhanced jigsaw, arc=.35mm, bottomrule=.15mm, bottomtitle=1mm, breakable, colback=white, colbacktitle=quarto-callout-note-color!10!white, colframe=quarto-callout-note-color-frame, coltitle=black, left=2mm, leftrule=.75mm, opacityback=0, opacitybacktitle=0.6, rightrule=.15mm, title=\textcolor{quarto-callout-note-color}{\faInfo}\hspace{0.5em}{Solution}, titlerule=0mm, toprule=.15mm, toptitle=1mm]

By the memoryless property, \[\mathbb{P}(X>7\mid X>3)=\mathbb{P}(X>4).\]
Since \(X\sim\operatorname{Exp}(2)\),
\[\mathbb{P}(X>4)=e^{-2\cdot 4}=e^{-8}.\] Thus
\[\boxed{\mathbb{P}(X>7\mid X>3)=e^{-8}.}\]

\end{tcolorbox}

\bookmarksetup{startatroot}

\chapter*{Chapter 6 Extra: Multinomial and Multinormal
Distributions}\label{chapter-6-extra-multinomial-and-multinormal-distributions}
\addcontentsline{toc}{chapter}{Chapter 6 Extra: Multinomial and
Multinormal Distributions}

\markboth{Chapter 6 Extra: Multinomial and Multinormal
Distributions}{Chapter 6 Extra: Multinomial and Multinormal
Distributions}

This extra chapter extends Section 6 in two directions. First, it
develops the multinomial distribution and the Dirichlet prior for
categorical data. Second, it introduces the multivariate Gaussian
distribution and the linear-algebra formulas behind marginalization,
conditioning, affine transformations, and linear Gaussian models.

\textbf{Topics.} Multinomial distribution; Dirichlet distribution;
Dirichlet-multinomial Bayesian model; multivariate Gaussian
distribution; Mahalanobis distance; spectral geometry; marginal and
conditional Gaussian distributions; affine transformations; products and
convolutions of Gaussians; linear Gaussian models.

\section*{Multinomial Distribution}\label{multinomial-distribution-1}
\addcontentsline{toc}{section}{Multinomial Distribution}

\markright{Multinomial Distribution}

This section extends the categorical and binomial distributions to
experiments with more than two possible outcomes.

\subsection*{From categorical trials to multinomial
counts}\label{from-categorical-trials-to-multinomial-counts}
\addcontentsline{toc}{subsection}{From categorical trials to multinomial
counts}

A categorical random variable records one outcome from several
categories, while a multinomial random vector records the counts of each
outcome after many independent categorical trials.

\textbf{Definition 1} (Categorical trial). A categorical trial has \(m\)
possible outcomes \(O_1,\ldots,O_m\) with probabilities
\[\phi_1,\ldots,\phi_m,\qquad \phi_i\ge 0,\qquad \sum_{i=1}^m \phi_i=1.\]
If one trial is performed, then the outcome has a categorical
distribution.

\textbf{Definition 2} (Multinomial distribution). Suppose we perform
\(n\) independent trials, each with outcomes \(O_1,\ldots,O_m\) and
constant category probabilities \((\phi_1,\ldots,\phi_m)\). Let
\[X_i=\text{the number of times outcome }O_i\text{ appears in the }n\text{ trials}.\]
Then \[X=(X_1,\ldots,X_m)\sim \operatorname{Multinomial}(n,\phi_1,
\ldots,\phi_m),\] and for nonnegative integers \(n_1,\ldots,n_m\)
satisfying \(n_1+\cdots+n_m=n\), \[\mathbb{P}(X_1=n_1,\ldots,X_m=n_m)
=
\frac{n!}{n_1!\cdots n_m!}\phi_1^{n_1}\cdots \phi_m^{n_m}.\]

\begin{tcolorbox}[enhanced jigsaw, arc=.35mm, bottomrule=.15mm, bottomtitle=1mm, breakable, colback=white, colbacktitle=quarto-callout-tip-color!10!white, colframe=quarto-callout-tip-color-frame, coltitle=black, left=2mm, leftrule=.75mm, opacityback=0, opacitybacktitle=0.6, rightrule=.15mm, title=\textcolor{quarto-callout-tip-color}{\faLightbulb}\hspace{0.5em}{Course connection}, titlerule=0mm, toprule=.15mm, toptitle=1mm]

The multinomial distribution is the direct generalization of the
categorical distribution. When \(n=1\), it reduces to a categorical
trial. When \(m=2\), it reduces to the binomial distribution.

\end{tcolorbox}

\textbf{Example 3} (Tossing an \(m\)-sided die). An \(m\)-sided die has
probabilities \(\phi_1,\ldots,\phi_m\) for sides \(1,\ldots,m\). Toss
the die \(n\) times and let \(X_i\) be the number of times side \(i\)
appears. Find the joint pmf of \((X_1,\ldots,X_m)\).

\begin{tcolorbox}[enhanced jigsaw, arc=.35mm, bottomrule=.15mm, bottomtitle=1mm, breakable, colback=white, colbacktitle=quarto-callout-note-color!10!white, colframe=quarto-callout-note-color-frame, coltitle=black, left=2mm, leftrule=.75mm, opacityback=0, opacitybacktitle=0.6, rightrule=.15mm, title=\textcolor{quarto-callout-note-color}{\faInfo}\hspace{0.5em}{Solution}, titlerule=0mm, toprule=.15mm, toptitle=1mm]

The vector of counts follows
\[(X_1,\ldots,X_m)\sim \operatorname{Multinomial}(n,\phi_1,\ldots,\phi_m).\]
Therefore, for \(n_1+\cdots+n_m=n\),
\[\mathbb{P}(X_1=n_1,\ldots,X_m=n_m)
=
\frac{n!}{n_1!\cdots n_m!}\prod_{i=1}^m\phi_i^{n_i}.\] The coefficient
\(n!/(n_1!\cdots n_m!)\) counts how many orderings of the \(n\) trials
produce the same vector of counts.

\end{tcolorbox}

\textbf{Example 4} (Three-category experiment). Suppose a website visit
leads to one of three outcomes: purchase, sign-up only, or no action.
The probabilities are \((0.1,0.2,0.7)\). Among \(n=10\) visitors, find
the probability of exactly \(2\) purchases, \(3\) sign-ups only, and
\(5\) no-action visits.

\begin{tcolorbox}[enhanced jigsaw, arc=.35mm, bottomrule=.15mm, bottomtitle=1mm, breakable, colback=white, colbacktitle=quarto-callout-note-color!10!white, colframe=quarto-callout-note-color-frame, coltitle=black, left=2mm, leftrule=.75mm, opacityback=0, opacitybacktitle=0.6, rightrule=.15mm, title=\textcolor{quarto-callout-note-color}{\faInfo}\hspace{0.5em}{Solution}, titlerule=0mm, toprule=.15mm, toptitle=1mm]

Let \(X=(X_1,X_2,X_3)\sim \operatorname{Multinomial}(10,0.1,0.2,0.7)\).
Then \[\mathbb{P}(X_1=2,X_2=3,X_3=5)
=
\frac{10!}{2!3!5!}(0.1)^2(0.2)^3(0.7)^5.\] Numerically,
\[\frac{10!}{2!3!5!}(0.1)^2(0.2)^3(0.7)^5
=2520(0.01)(0.008)(0.16807)\approx 0.0339.\]

\end{tcolorbox}

\subsection*{Mean, variance, and
covariance}\label{mean-variance-and-covariance}
\addcontentsline{toc}{subsection}{Mean, variance, and covariance}

The components of a multinomial vector are individually binomial, but
they are dependent because the total count is fixed at \(n\).

\textbf{Theorem 5} (Moments of the multinomial distribution). \emph{If
\(X=(X_1,\ldots,X_m)\sim \operatorname{Multinomial}(n,\phi_1,\ldots,\phi_m)\),
then \[\mathbb{E}[X_i]=n\phi_i,
\qquad
\operatorname{Var}(X_i)=n\phi_i(1-\phi_i),\] and for \(i\ne j\),
\[\operatorname{Cov}(X_i,X_j)=-n\phi_i\phi_j.\] Equivalently,
\[\operatorname{Cov}(X)=n\left(\operatorname{diag}(\phi)-\phi\phi^T\right),\]
where \(\phi=(\phi_1,\ldots,\phi_m)^T\).}

\begin{tcolorbox}[enhanced jigsaw, arc=.35mm, bottomrule=.15mm, bottomtitle=1mm, breakable, colback=white, colbacktitle=quarto-callout-note-color!10!white, colframe=quarto-callout-note-color-frame, coltitle=black, left=2mm, leftrule=.75mm, opacityback=0, opacitybacktitle=0.6, rightrule=.15mm, title=\textcolor{quarto-callout-note-color}{\faInfo}\hspace{0.5em}{Proof}, titlerule=0mm, toprule=.15mm, toptitle=1mm]

Write \(X_i=\sum_{r=1}^n I_{ri}\), where \(I_{ri}=1\) if trial \(r\)
produces outcome \(O_i\), and \(0\) otherwise. Then
\(I_{ri}\sim \operatorname{Bernoulli}(\phi_i)\), so
\[\mathbb{E}[X_i]=\sum_{r=1}^n \mathbb{E}[I_{ri}]=n\phi_i,
\qquad
\operatorname{Var}(X_i)=\sum_{r=1}^n \operatorname{Var}(I_{ri})=n\phi_i(1-\phi_i).\]
For \(i\ne j\), in one trial \(I_{ri}I_{rj}=0\), so
\[\operatorname{Cov}(I_{ri},I_{rj})=\mathbb{E}[I_{ri}I_{rj}]-\mathbb{E}[I_{ri}]\mathbb{E}[I_{rj}]=0-\phi_i\phi_j=-\phi_i\phi_j.\]
Different trials are independent, so
\[\operatorname{Cov}(X_i,X_j)=\sum_{r=1}^n \operatorname{Cov}(I_{ri},I_{rj})=-n\phi_i\phi_j.\]

\end{tcolorbox}

\begin{remark}
\emph{Remark 6}. The negative covariance is not a paradox. If the total
number of trials is fixed, seeing more observations in one category
leaves fewer observations available for the other categories.
\end{remark}

\textbf{Practice Problem 7} (Checking the binomial special case). Show
that if \(m=2\), then the multinomial distribution reduces to the
binomial distribution.

\begin{tcolorbox}[enhanced jigsaw, arc=.35mm, bottomrule=.15mm, bottomtitle=1mm, breakable, colback=white, colbacktitle=quarto-callout-note-color!10!white, colframe=quarto-callout-note-color-frame, coltitle=black, left=2mm, leftrule=.75mm, opacityback=0, opacitybacktitle=0.6, rightrule=.15mm, title=\textcolor{quarto-callout-note-color}{\faInfo}\hspace{0.5em}{Solution}, titlerule=0mm, toprule=.15mm, toptitle=1mm]

Let \(X=(X_1,X_2)\sim \operatorname{Multinomial}(n,p,1-p)\). Because
\(X_1+X_2=n\), the entire vector is determined by \(X_1\). For
\(X_1=k\), we have \(X_2=n-k\). Thus
\[\mathbb{P}(X_1=k)=\mathbb{P}(X_1=k,X_2=n-k)
=\frac{n!}{k!(n-k)!}p^k(1-p)^{n-k},\] which is exactly the
\(\operatorname{Binomial}(n,p)\) pmf.

\end{tcolorbox}

\section*{Dirichlet Distribution}\label{dirichlet-distribution}
\addcontentsline{toc}{section}{Dirichlet Distribution}

\markright{Dirichlet Distribution}

This section introduces the Dirichlet distribution, the continuous
distribution on probability vectors that generalizes the Beta
distribution.

\subsection*{The simplex and the Dirichlet
density}\label{the-simplex-and-the-dirichlet-density}
\addcontentsline{toc}{subsection}{The simplex and the Dirichlet density}

The Dirichlet distribution is used when the unknown quantity is itself a
vector of probabilities.

\textbf{Definition 8} (Probability simplex). The \((K-1)\)-dimensional
probability simplex is
\[\Delta^{K-1}=\left\{\mu=(\mu_1,\ldots,\mu_K):\; \mu_k\ge 0,\; \sum_{k=1}^K \mu_k=1\right\}.\]
A point in the simplex represents a probability vector over \(K\)
categories.

\textbf{Definition 9} (Dirichlet distribution). Let
\(\alpha=(\alpha_1,\ldots,\alpha_K)\) with \(\alpha_k>0\) and let
\[\alpha_0=\sum_{k=1}^K \alpha_k.\] A random probability vector
\(\mu=(\mu_1,\ldots,\mu_K)\) has a Dirichlet distribution with parameter
\(\alpha\), written
\[\mu\sim \operatorname{Dirichlet}(\alpha_1,\ldots,\alpha_K),\] if its
density on the simplex is \[p(\mu\mid\alpha)=
 \frac{\Gamma(\alpha_0)}{\Gamma(\alpha_1)\cdots \Gamma(\alpha_K)}
 \prod_{k=1}^K \mu_k^{\alpha_k-1},\] where \(\Gamma(\cdot)\) is the
gamma function.

\begin{tcolorbox}[enhanced jigsaw, arc=.35mm, bottomrule=.15mm, bottomtitle=1mm, breakable, colback=white, colbacktitle=quarto-callout-tip-color!10!white, colframe=quarto-callout-tip-color-frame, coltitle=black, left=2mm, leftrule=.75mm, opacityback=0, opacitybacktitle=0.6, rightrule=.15mm, title=\textcolor{quarto-callout-tip-color}{\faLightbulb}\hspace{0.5em}{Interpretation of the parameters}, titlerule=0mm, toprule=.15mm, toptitle=1mm]

The parameter \(\alpha_k\) acts like a prior pseudo-count for category
\(k\). Large values of \(\alpha_0\) mean the distribution is more
concentrated around its mean; small values can place more mass near the
corners of the simplex.

\end{tcolorbox}

\textbf{Theorem 10} (Dirichlet moments). \emph{If
\(\mu\sim \operatorname{Dirichlet}(\alpha_1,\ldots,\alpha_K)\) and
\(\alpha_0=\sum_k\alpha_k\), then
\[\mathbb{E}[\mu_k]=\frac{\alpha_k}{\alpha_0},\]
\[\operatorname{Var}(\mu_k)=\frac{\alpha_k(\alpha_0-\alpha_k)}{\alpha_0^2(\alpha_0+1)},
\qquad
\operatorname{Cov}(\mu_i,\mu_j)= -\frac{\alpha_i\alpha_j}{\alpha_0^2(\alpha_0+1)}\quad (i\ne j).\]}

\textbf{Example 11} (The Beta distribution as a special case). Show that
when \(K=2\), the Dirichlet distribution becomes the Beta distribution.

\begin{tcolorbox}[enhanced jigsaw, arc=.35mm, bottomrule=.15mm, bottomtitle=1mm, breakable, colback=white, colbacktitle=quarto-callout-note-color!10!white, colframe=quarto-callout-note-color-frame, coltitle=black, left=2mm, leftrule=.75mm, opacityback=0, opacitybacktitle=0.6, rightrule=.15mm, title=\textcolor{quarto-callout-note-color}{\faInfo}\hspace{0.5em}{Solution}, titlerule=0mm, toprule=.15mm, toptitle=1mm]

When \(K=2\), the simplex condition says \(\mu_1+\mu_2=1\), so
\(\mu_2=1-\mu_1\). The Dirichlet density is proportional to
\[\mu_1^{\alpha_1-1}\mu_2^{\alpha_2-1}
=\mu_1^{\alpha_1-1}(1-\mu_1)^{\alpha_2-1}.\] The normalized density is
\[\frac{\Gamma(\alpha_1+\alpha_2)}{\Gamma(\alpha_1)\Gamma(\alpha_2)}
\mu_1^{\alpha_1-1}(1-\mu_1)^{\alpha_2-1},\] which is the
\(\operatorname{Beta}(\alpha_1,\alpha_2)\) density. Therefore,
\[\operatorname{Dirichlet}(\alpha_1,\alpha_2)=\operatorname{Beta}(\alpha_1,\alpha_2).\]

\end{tcolorbox}

\subsection*{Dirichlet-multinomial Bayesian
model}\label{dirichlet-multinomial-bayesian-model}
\addcontentsline{toc}{subsection}{Dirichlet-multinomial Bayesian model}

The Dirichlet distribution is the natural conjugate prior for
categorical and multinomial probability vectors.

\textbf{Definition 12} (Dirichlet-multinomial model). A Bayesian model
for categorical or multinomial data is
\[\mu\sim \operatorname{Dirichlet}(\alpha_1,
\ldots,\alpha_K),
\qquad
X\mid\mu\sim \operatorname{Multinomial}(n,\mu_1,\ldots,\mu_K).\] Here
the likelihood is multinomial, the prior is Dirichlet, and the posterior
is again Dirichlet.

\textbf{Theorem 13} (Dirichlet-multinomial conjugacy). \emph{Suppose
\[\mu\sim \operatorname{Dirichlet}(\alpha_1,
\ldots,\alpha_K),
\qquad
X=(X_1,\ldots,X_K)\mid\mu\sim \operatorname{Multinomial}(n,\mu_1,\ldots,\mu_K).\]
If the observed counts are \(X_1=n_1,\ldots,X_K=n_K\), then
\[\mu\mid X=(n_1,\ldots,n_K)
\sim
\operatorname{Dirichlet}(\alpha_1+n_1,
\ldots,
\alpha_K+n_K).\]}

\begin{tcolorbox}[enhanced jigsaw, arc=.35mm, bottomrule=.15mm, bottomtitle=1mm, breakable, colback=white, colbacktitle=quarto-callout-note-color!10!white, colframe=quarto-callout-note-color-frame, coltitle=black, left=2mm, leftrule=.75mm, opacityback=0, opacitybacktitle=0.6, rightrule=.15mm, title=\textcolor{quarto-callout-note-color}{\faInfo}\hspace{0.5em}{Proof}, titlerule=0mm, toprule=.15mm, toptitle=1mm]

The prior density is \[p(\mu)\propto \prod_{k=1}^K \mu_k^{\alpha_k-1}.\]
The multinomial likelihood, as a function of \(\mu\), is
\[p(x\mid\mu)\propto \prod_{k=1}^K \mu_k^{n_k}.\] Therefore, by Bayes'
theorem, \[p(\mu\mid x)\propto p(x\mid\mu)p(\mu)
\propto
\prod_{k=1}^K \mu_k^{n_k}\prod_{k=1}^K\mu_k^{\alpha_k-1}
=
\prod_{k=1}^K \mu_k^{\alpha_k+n_k-1}.\] This is the kernel of a
Dirichlet distribution with parameters \(\alpha_k+n_k\).

\end{tcolorbox}

\textbf{Example 14} (Updating category probabilities). Suppose there are
three categories and the prior is
\[\mu\sim \operatorname{Dirichlet}(2,2,2).\] We observe \(n=10\) trials
with counts \((5,3,2)\). Find the posterior distribution and posterior
mean.

\begin{tcolorbox}[enhanced jigsaw, arc=.35mm, bottomrule=.15mm, bottomtitle=1mm, breakable, colback=white, colbacktitle=quarto-callout-note-color!10!white, colframe=quarto-callout-note-color-frame, coltitle=black, left=2mm, leftrule=.75mm, opacityback=0, opacitybacktitle=0.6, rightrule=.15mm, title=\textcolor{quarto-callout-note-color}{\faInfo}\hspace{0.5em}{Solution}, titlerule=0mm, toprule=.15mm, toptitle=1mm]

By conjugacy,
\[\mu\mid X=(5,3,2)\sim \operatorname{Dirichlet}(2+5,2+3,2+2)=\operatorname{Dirichlet}(7,5,4).\]
The posterior total is \(7+5+4=16\). Therefore,
\[\mathbb{E}[\mu_1\mid X]=\frac{7}{16},
\qquad
\mathbb{E}[\mu_2\mid X]=\frac{5}{16},
\qquad
\mathbb{E}[\mu_3\mid X]=\frac{4}{16}.\] The posterior mean is a weighted
compromise between the prior pseudo-counts and the observed counts.

\end{tcolorbox}

\textbf{Practice Problem 15} (Uniform prior on the simplex). Let
\(\mu\sim \operatorname{Dirichlet}(1,1,1)\) and suppose counts
\((n_1,n_2,n_3)=(4,1,5)\) are observed. Find the posterior distribution.

\begin{tcolorbox}[enhanced jigsaw, arc=.35mm, bottomrule=.15mm, bottomtitle=1mm, breakable, colback=white, colbacktitle=quarto-callout-note-color!10!white, colframe=quarto-callout-note-color-frame, coltitle=black, left=2mm, leftrule=.75mm, opacityback=0, opacitybacktitle=0.6, rightrule=.15mm, title=\textcolor{quarto-callout-note-color}{\faInfo}\hspace{0.5em}{Solution}, titlerule=0mm, toprule=.15mm, toptitle=1mm]

The posterior is
\[\mu\mid X\sim \operatorname{Dirichlet}(1+4,1+1,1+5)=\operatorname{Dirichlet}(5,2,6).\]
Because \(\operatorname{Dirichlet}(1,1,1)\) is uniform over the
three-category simplex, the posterior parameters are simply one plus the
observed counts.

\end{tcolorbox}

\section*{Single Variable Gaussian
Distribution}\label{single-variable-gaussian-distribution}
\addcontentsline{toc}{section}{Single Variable Gaussian Distribution}

\markright{Single Variable Gaussian Distribution}

This section reviews the ordinary one-dimensional Gaussian distribution
before moving to the multivariate Gaussian distribution.

\subsection*{Density, mean, variance, and
probability}\label{density-mean-variance-and-probability}
\addcontentsline{toc}{subsection}{Density, mean, variance, and
probability}

The one-dimensional Gaussian distribution is the prototype for the
multivariate Gaussian distribution.

\textbf{Definition 16} (Univariate Gaussian). A random variable \(X\)
has a Gaussian or normal distribution with mean \(\mu\) and variance
\(\sigma^2\), written \[X\sim \operatorname{Normal}(\mu,\sigma^2),\] if
its density is
\[f_X(x)=\frac{1}{\sqrt{2\pi}\sigma}\exp\left\{-\frac{1}{2}\left(\frac{x-\mu}{\sigma}\right)^2\right\},
\qquad -\infty<x<\infty.\] Then \[\mathbb{E}[X]=\mu,
\qquad
\operatorname{Var}(X)=\sigma^2.\] For any \(a<b\),
\[\mathbb{P}(a<X<b)=\int_a^b f_X(x)\,dx.\]

\textbf{Example 17} (Standardizing a normal random variable). Let
\(X\sim \operatorname{Normal}(\mu,\sigma^2)\). Show that
\[Z=\frac{X-\mu}{\sigma}\] has the standard normal distribution.

\begin{tcolorbox}[enhanced jigsaw, arc=.35mm, bottomrule=.15mm, bottomtitle=1mm, breakable, colback=white, colbacktitle=quarto-callout-note-color!10!white, colframe=quarto-callout-note-color-frame, coltitle=black, left=2mm, leftrule=.75mm, opacityback=0, opacitybacktitle=0.6, rightrule=.15mm, title=\textcolor{quarto-callout-note-color}{\faInfo}\hspace{0.5em}{Solution}, titlerule=0mm, toprule=.15mm, toptitle=1mm]

The transformation is \(Z=(X-\mu)/\sigma\), so \(X=\mu+\sigma Z\). The
change-of-variables formula gives
\[f_Z(z)=f_X(\mu+\sigma z)\cdot \sigma.\] Substituting the density of
\(X\),
\[f_Z(z)=\frac{1}{\sqrt{2\pi}\sigma}\exp\left(-\frac{z^2}{2}\right)\sigma
=\frac{1}{\sqrt{2\pi}}e^{-z^2/2}.\] Thus
\(Z\sim \operatorname{Normal}(0,1)\).

\end{tcolorbox}

\section*{Multivariate Gaussian
Distribution}\label{multivariate-gaussian-distribution}
\addcontentsline{toc}{section}{Multivariate Gaussian Distribution}

\markright{Multivariate Gaussian Distribution}

This section introduces the Gaussian distribution for random vectors and
explains how the covariance matrix controls its geometry.

\subsection*{Definition and covariance
matrix}\label{definition-and-covariance-matrix}
\addcontentsline{toc}{subsection}{Definition and covariance matrix}

A multivariate Gaussian distribution is completely determined by its
mean vector and covariance matrix.

\textbf{Definition 18} (Multivariate Gaussian). Let \(X\in\mathbb{R}^d\)
be a random vector. We say \[X\sim \mathcal{N}_d(\mu,\Sigma),\] where
\(\mu\in\mathbb{R}^d\) and \(\Sigma\) is a \(d\times d\) symmetric
positive definite matrix, if \[p(x\mid\mu,\Sigma)
=
\frac{1}{\sqrt{(2\pi)^d|\Sigma|}}
\exp\left\{-\frac12 (x-\mu)^T\Sigma^{-1}(x-\mu)\right\}.\] The
covariance matrix is
\[\Sigma=\operatorname{Cov}(X,X)=\mathbb{E}\left[(X-\mu)(X-\mu)^T\right].\]

\begin{tcolorbox}[enhanced jigsaw, arc=.35mm, bottomrule=.15mm, bottomtitle=1mm, breakable, colback=white, colbacktitle=quarto-callout-tip-color!10!white, colframe=quarto-callout-tip-color-frame, coltitle=black, left=2mm, leftrule=.75mm, opacityback=0, opacitybacktitle=0.6, rightrule=.15mm, title=\textcolor{quarto-callout-tip-color}{\faLightbulb}\hspace{0.5em}{Normalization constant}, titlerule=0mm, toprule=.15mm, toptitle=1mm]

The normalization constant is
\[Z=\sqrt{\det(2\pi\Sigma)}=\sqrt{(2\pi)^d|\Sigma|}.\] This is the
multivariate analogue of \(\sqrt{2\pi}\sigma\) in the one-dimensional
normal density.

\end{tcolorbox}

\subsection*{Mahalanobis distance and
geometry}\label{mahalanobis-distance-and-geometry}
\addcontentsline{toc}{subsection}{Mahalanobis distance and geometry}

The shape of a multivariate Gaussian is controlled by a quadratic form
that generalizes the one-dimensional \(z\)-score.

\textbf{Definition 19} (Mahalanobis distance). For
\(X\sim \mathcal{N}_d(\mu,\Sigma)\), define
\[\Delta^2=(x-\mu)^T\Sigma^{-1}(x-\mu).\] The value \(\Delta\) is called
the Mahalanobis distance from \(x\) to \(\mu\) with respect to
\(\Sigma\).

\begin{remark}
\emph{Remark 20}. In one dimension,
\[\Delta^2=\frac{(x-\mu)^2}{\sigma^2},
\qquad
\Delta=\left|\frac{x-\mu}{\sigma}\right|.\] Thus Mahalanobis distance
generalizes the absolute value of the usual \(z\)-score. If
\(\Sigma=I\), then \[\Delta^2=(x-\mu)^T(x-\mu)=\|x-\mu\|^2,\] so
Mahalanobis distance reduces to Euclidean distance.
\end{remark}

\textbf{Example 21} (Two-dimensional covariance matrices). Compare the
covariance matrices \[\Sigma_1=\begin{pmatrix}1&0\\0&1\end{pmatrix},
\qquad
\Sigma_2=\begin{pmatrix}1&0.5\\0.5&1\end{pmatrix},
\qquad
\Sigma_3=\begin{pmatrix}1&0.8\\0.8&1\end{pmatrix}.\] Describe the
qualitative effect on the Gaussian density.

\begin{tcolorbox}[enhanced jigsaw, arc=.35mm, bottomrule=.15mm, bottomtitle=1mm, breakable, colback=white, colbacktitle=quarto-callout-note-color!10!white, colframe=quarto-callout-note-color-frame, coltitle=black, left=2mm, leftrule=.75mm, opacityback=0, opacitybacktitle=0.6, rightrule=.15mm, title=\textcolor{quarto-callout-note-color}{\faInfo}\hspace{0.5em}{Solution}, titlerule=0mm, toprule=.15mm, toptitle=1mm]

For \(\Sigma_1\), the coordinates are uncorrelated with equal variance,
so the contours are circles centered at \(\mu\).

For \(\Sigma_2\), the positive off-diagonal covariance creates
positively tilted elliptical contours. Larger \(x_1\) values tend to be
associated with larger \(x_2\) values.

For \(\Sigma_3\), the positive dependence is stronger, so the ellipse is
more stretched along the positive diagonal direction. The density is
more concentrated around a narrow tilted ridge than in the case
\(\Sigma_2\).

\end{tcolorbox}

\subsection*{Spectral decomposition and precision
matrix}\label{spectral-decomposition-and-precision-matrix}
\addcontentsline{toc}{subsection}{Spectral decomposition and precision
matrix}

The eigenvalues and eigenvectors of the covariance matrix reveal the
principal axes of a Gaussian distribution.

\textbf{Theorem 22} (Spectral form of the quadratic term). \emph{Let
\(\Sigma\) be symmetric positive definite. Then
\[\Sigma=UDU^T=\sum_{i=1}^d \lambda_i u_i u_i^T,\] where
\(U=[u_1\ \cdots\ u_d]\) is orthogonal and
\(D=\operatorname{diag}(\lambda_1,\ldots,\lambda_d)\) with
\(\lambda_i>0\). The precision matrix is
\[\Sigma^{-1}=UD^{-1}U^T=\sum_{i=1}^d\frac{1}{\lambda_i}u_i u_i^T.\] If
\[y_i=u_i^T(x-\mu),\] then
\[\Delta^2=(x-\mu)^T\Sigma^{-1}(x-\mu)=\sum_{i=1}^d \frac{y_i^2}{\lambda_i}.\]}

\begin{tcolorbox}[enhanced jigsaw, arc=.35mm, bottomrule=.15mm, bottomtitle=1mm, breakable, colback=white, colbacktitle=quarto-callout-note-color!10!white, colframe=quarto-callout-note-color-frame, coltitle=black, left=2mm, leftrule=.75mm, opacityback=0, opacitybacktitle=0.6, rightrule=.15mm, title=\textcolor{quarto-callout-note-color}{\faInfo}\hspace{0.5em}{Proof}, titlerule=0mm, toprule=.15mm, toptitle=1mm]

Since \(\Sigma=UDU^T\) and \(U^TU=I\), the inverse is
\[\Sigma^{-1}=UD^{-1}U^T.\] Thus \[(x-\mu)^T\Sigma^{-1}(x-\mu)
=(x-\mu)^TUD^{-1}U^T(x-\mu).\] Let \(y=U^T(x-\mu)\). Then the quadratic
form becomes \[y^TD^{-1}y=\sum_{i=1}^d\frac{y_i^2}{\lambda_i}.\]

\end{tcolorbox}

\begin{tcolorbox}[enhanced jigsaw, arc=.35mm, bottomrule=.15mm, bottomtitle=1mm, breakable, colback=white, colbacktitle=quarto-callout-warning-color!10!white, colframe=quarto-callout-warning-color-frame, coltitle=black, left=2mm, leftrule=.75mm, opacityback=0, opacitybacktitle=0.6, rightrule=.15mm, title=\textcolor{quarto-callout-warning-color}{\faExclamationTriangle}\hspace{0.5em}{Terminology}, titlerule=0mm, toprule=.15mm, toptitle=1mm]

The matrix \(\Sigma^{-1}\) is called the \emph{precision matrix}. The
covariance matrix describes spread; the precision matrix describes
inverse spread and appears naturally in the exponent of the Gaussian
density.

\end{tcolorbox}

\subsection*{Why Gaussian distributions are
central}\label{why-gaussian-distributions-are-central}
\addcontentsline{toc}{subsection}{Why Gaussian distributions are
central}

Gaussian distributions appear frequently because of the central limit
theorem and because Gaussianity is preserved by many important
operations.

\begin{tcolorbox}[enhanced jigsaw, arc=.35mm, bottomrule=.15mm, bottomtitle=1mm, breakable, colback=white, colbacktitle=quarto-callout-tip-color!10!white, colframe=quarto-callout-tip-color-frame, coltitle=black, left=2mm, leftrule=.75mm, opacityback=0, opacitybacktitle=0.6, rightrule=.15mm, title=\textcolor{quarto-callout-tip-color}{\faLightbulb}\hspace{0.5em}{Once Gaussian, always Gaussian}, titlerule=0mm, toprule=.15mm, toptitle=1mm]

Gaussian random vectors are completely determined by mean vector \(\mu\)
and covariance matrix \(\Sigma\). Linear transformations,
marginalization, conditioning, and sums of independent Gaussian vectors
all produce Gaussian distributions again.

\end{tcolorbox}

\begin{remark}
\emph{Remark 23} (Central limit theorem motivation). If
\(X_1,\ldots,X_m\) are independent and identically distributed with
finite mean and variance, then the arithmetic mean
\[\bar X=\frac{1}{m}(X_1+\cdots+X_m)\] is approximately normal for large
\(m\), after appropriate centering and scaling. This explains why
Gaussian distributions occur often in real data.
\end{remark}

\section*{Marginal and Conditional Gaussian
Distributions}\label{marginal-and-conditional-gaussian-distributions}
\addcontentsline{toc}{section}{Marginal and Conditional Gaussian
Distributions}

\markright{Marginal and Conditional Gaussian Distributions}

This section gives the most important computational formulas for
partitioned multivariate Gaussian distributions.

\subsection*{Partitioned Gaussian
vectors}\label{partitioned-gaussian-vectors}
\addcontentsline{toc}{subsection}{Partitioned Gaussian vectors}

Partitioning a Gaussian vector allows us to study subsets of variables
and conditional distributions.

Suppose \[X=\begin{pmatrix}X_1\\X_2\end{pmatrix}
\sim
\mathcal{N}\left(
\begin{pmatrix}\mu_1\\\mu_2\end{pmatrix},
\begin{pmatrix}
\Sigma_{11}&\Sigma_{12}\\
\Sigma_{21}&\Sigma_{22}
\end{pmatrix}
\right).\] Here \(X_1\) and \(X_2\) may themselves be vectors, and the
blocks of \(\Sigma\) have compatible dimensions.

\subsection*{Marginalization}\label{marginalization}
\addcontentsline{toc}{subsection}{Marginalization}

The marginal distribution of a subvector of a Gaussian vector is
Gaussian with the corresponding mean block and covariance block.

\textbf{Theorem 24} (Gaussian marginal distribution). \emph{If
\[X=\begin{pmatrix}X_1\\X_2\end{pmatrix}
\sim
\mathcal{N}\left(
\begin{pmatrix}\mu_1\\\mu_2\end{pmatrix},
\begin{pmatrix}
\Sigma_{11}&\Sigma_{12}\\
\Sigma_{21}&\Sigma_{22}
\end{pmatrix}
\right),\] then \[X_1\sim \mathcal{N}(\mu_1,\Sigma_{11}),
\qquad
X_2\sim \mathcal{N}(\mu_2,\Sigma_{22}).\]}

\textbf{Example 25} (Marginal of a bivariate normal). Let
\[\begin{pmatrix}X_1\\X_2\end{pmatrix}
\sim
\mathcal{N}\left(
\begin{pmatrix}1\\2\end{pmatrix},
\begin{pmatrix}4&1\\1&9\end{pmatrix}
\right).\] Find the marginal distributions of \(X_1\) and \(X_2\).

\begin{tcolorbox}[enhanced jigsaw, arc=.35mm, bottomrule=.15mm, bottomtitle=1mm, breakable, colback=white, colbacktitle=quarto-callout-note-color!10!white, colframe=quarto-callout-note-color-frame, coltitle=black, left=2mm, leftrule=.75mm, opacityback=0, opacitybacktitle=0.6, rightrule=.15mm, title=\textcolor{quarto-callout-note-color}{\faInfo}\hspace{0.5em}{Solution}, titlerule=0mm, toprule=.15mm, toptitle=1mm]

By the marginalization theorem, \[X_1\sim \operatorname{Normal}(1,4),
\qquad
X_2\sim \operatorname{Normal}(2,9).\] The off-diagonal covariance \(1\)
affects the dependence between \(X_1\) and \(X_2\), but it does not
change the marginal variances.

\end{tcolorbox}

\subsection*{Conditional Gaussian
distributions}\label{conditional-gaussian-distributions}
\addcontentsline{toc}{subsection}{Conditional Gaussian distributions}

Conditioning on one part of a Gaussian vector gives another Gaussian
distribution.

\textbf{Theorem 26} (Conditional Gaussian distribution). \emph{Suppose
\[X=\begin{pmatrix}X_1\\X_2\end{pmatrix}
\sim
\mathcal{N}\left(
\begin{pmatrix}\mu_1\\\mu_2\end{pmatrix},
\begin{pmatrix}
\Sigma_{11}&\Sigma_{12}\\
\Sigma_{21}&\Sigma_{22}
\end{pmatrix}
\right),\] where \(\Sigma_{22}\) is invertible. Then
\[X_1\mid X_2=x_2\sim \mathcal{N}(\mu_{1\mid 2},\Sigma_{11\mid 2}),\]
where \[\mu_{1\mid 2}=\mu_1+\Sigma_{12}\Sigma_{22}^{-1}(x_2-\mu_2),\]
and
\[\Sigma_{11\mid 2}=\Sigma_{11}-\Sigma_{12}\Sigma_{22}^{-1}\Sigma_{21}.\]}

\begin{tcolorbox}[enhanced jigsaw, arc=.35mm, bottomrule=.15mm, bottomtitle=1mm, breakable, colback=white, colbacktitle=quarto-callout-note-color!10!white, colframe=quarto-callout-note-color-frame, coltitle=black, left=2mm, leftrule=.75mm, opacityback=0, opacitybacktitle=0.6, rightrule=.15mm, title=\textcolor{quarto-callout-note-color}{\faInfo}\hspace{0.5em}{Proof}, titlerule=0mm, toprule=.15mm, toptitle=1mm]

\emph{Proof idea by completing the square.} Let the precision matrix be
\[\Lambda=\Sigma^{-1}=\begin{pmatrix}\Lambda_{11}&\Lambda_{12}\\\Lambda_{21}&\Lambda_{22}\end{pmatrix}.\]
The exponent of the joint Gaussian density is
\[-\frac12(x-\mu)^T\Sigma^{-1}(x-\mu).\] If \(x_2\) is fixed, then as a
function of \(x_1\) this expression is a quadratic form in \(x_1\).
Therefore the conditional distribution must be Gaussian. Completing the
square gives \[\Sigma_{11\mid 2}=\Lambda_{11}^{-1},
\qquad
\mu_{1\mid 2}=\mu_1-\Lambda_{11}^{-1}\Lambda_{12}(x_2-\mu_2).\] Using
the block inverse formula and Schur complement identities,
\[\Lambda_{11}^{-1}=\Sigma_{11}-\Sigma_{12}\Sigma_{22}^{-1}\Sigma_{21},\]
and \[-\Lambda_{11}^{-1}\Lambda_{12}=\Sigma_{12}\Sigma_{22}^{-1}.\] This
gives the stated formula.

\end{tcolorbox}

\begin{tcolorbox}[enhanced jigsaw, arc=.35mm, bottomrule=.15mm, bottomtitle=1mm, breakable, colback=white, colbacktitle=quarto-callout-tip-color!10!white, colframe=quarto-callout-tip-color-frame, coltitle=black, left=2mm, leftrule=.75mm, opacityback=0, opacitybacktitle=0.6, rightrule=.15mm, title=\textcolor{quarto-callout-tip-color}{\faLightbulb}\hspace{0.5em}{Interpretation}, titlerule=0mm, toprule=.15mm, toptitle=1mm]

The conditional mean is a linear function of the observed value \(x_2\).
The conditional covariance does not depend on the observed value
\(x_2\); it only depends on the covariance blocks.

\end{tcolorbox}

\textbf{Example 27} (Conditional distribution in two dimensions). Let
\[\begin{pmatrix}X_1\\X_2\end{pmatrix}
\sim
\mathcal{N}\left(
\begin{pmatrix}0\\0\end{pmatrix},
\begin{pmatrix}1&\rho\\\rho&1\end{pmatrix}
\right),
\qquad -1<\rho<1.\] Find \(X_1\mid X_2=x_2\).

\begin{tcolorbox}[enhanced jigsaw, arc=.35mm, bottomrule=.15mm, bottomtitle=1mm, breakable, colback=white, colbacktitle=quarto-callout-note-color!10!white, colframe=quarto-callout-note-color-frame, coltitle=black, left=2mm, leftrule=.75mm, opacityback=0, opacitybacktitle=0.6, rightrule=.15mm, title=\textcolor{quarto-callout-note-color}{\faInfo}\hspace{0.5em}{Solution}, titlerule=0mm, toprule=.15mm, toptitle=1mm]

Here \[\mu_1=0,
\quad \mu_2=0,
\quad \Sigma_{11}=1,
\quad \Sigma_{22}=1,
\quad \Sigma_{12}=\rho.\] Therefore,
\[\mu_{1\mid 2}=0+\rho(1)^{-1}(x_2-0)=\rho x_2,\] and
\[\Sigma_{11\mid 2}=1-\rho(1)^{-1}\rho=1-\rho^2.\] Thus
\[X_1\mid X_2=x_2\sim \operatorname{Normal}(\rho x_2,1-\rho^2).\] If
\(\rho=0\), the conditional distribution is the same as the marginal
distribution, reflecting independence.

\end{tcolorbox}

\subsection*{Linear algebra review: Schur complements and block
inverses}\label{linear-algebra-review-schur-complements-and-block-inverses}
\addcontentsline{toc}{subsection}{Linear algebra review: Schur
complements and block inverses}

The Schur complement is the matrix identity behind the conditional
covariance formula.

\textbf{Definition 28} (Schur complement). Let
\[M=\begin{pmatrix}A&B\\C&D\end{pmatrix},\] where \(D\) is invertible.
The Schur complement of \(D\) in \(M\) is \[M/D=A-BD^{-1}C.\]

\textbf{Theorem 29} (Block inverse formula). \emph{If \(D\) and \(M/D\)
are invertible, then \[M^{-1}=\begin{pmatrix}
(M/D)^{-1} & -(M/D)^{-1}BD^{-1}\\
-D^{-1}C(M/D)^{-1} & D^{-1}+D^{-1}C(M/D)^{-1}BD^{-1}
\end{pmatrix}.\]}

\begin{tcolorbox}[enhanced jigsaw, arc=.35mm, bottomrule=.15mm, bottomtitle=1mm, breakable, colback=white, colbacktitle=quarto-callout-note-color!10!white, colframe=quarto-callout-note-color-frame, coltitle=black, left=2mm, leftrule=.75mm, opacityback=0, opacitybacktitle=0.6, rightrule=.15mm, title=\textcolor{quarto-callout-note-color}{\faInfo}\hspace{0.5em}{Proof}, titlerule=0mm, toprule=.15mm, toptitle=1mm]

Block Gaussian elimination gives \[\begin{pmatrix}A&B\\C&D\end{pmatrix}
=
\begin{pmatrix}I&BD^{-1}\\0&I\end{pmatrix}
\begin{pmatrix}A-BD^{-1}C&0\\0&D\end{pmatrix}
\begin{pmatrix}I&0\\D^{-1}C&I\end{pmatrix}.\] Taking inverses of the
three factors and multiplying them gives the stated block inverse
formula.

\end{tcolorbox}

\textbf{Practice Problem 30} (Schur complement calculation). Let
\[M=\begin{pmatrix}4&2\\2&3\end{pmatrix}.\] Compute the Schur complement
of \(D=3\).

\begin{tcolorbox}[enhanced jigsaw, arc=.35mm, bottomrule=.15mm, bottomtitle=1mm, breakable, colback=white, colbacktitle=quarto-callout-note-color!10!white, colframe=quarto-callout-note-color-frame, coltitle=black, left=2mm, leftrule=.75mm, opacityback=0, opacitybacktitle=0.6, rightrule=.15mm, title=\textcolor{quarto-callout-note-color}{\faInfo}\hspace{0.5em}{Solution}, titlerule=0mm, toprule=.15mm, toptitle=1mm]

Here \(A=4\), \(B=2\), \(C=2\), and \(D=3\). Therefore
\[M/D=A-BD^{-1}C=4-2\cdot\frac13\cdot 2=4-\frac43=\frac83.\]

\end{tcolorbox}

\section*{Operations on Gaussian Random
Variables}\label{operations-on-gaussian-random-variables}
\addcontentsline{toc}{section}{Operations on Gaussian Random Variables}

\markright{Operations on Gaussian Random Variables}

This section summarizes several operations that preserve Gaussian
structure.

\subsection*{Affine transformations}\label{affine-transformations}
\addcontentsline{toc}{subsection}{Affine transformations}

A linear transformation plus a shift maps a Gaussian vector to another
Gaussian vector.

\textbf{Theorem 31} (Affine transformation of a Gaussian). \emph{If
\[X\sim \mathcal{N}(\mu,\Sigma),
\qquad
Y=AX+b,\] then \[Y\sim \mathcal{N}(A\mu+b,A\Sigma A^T).\]}

\begin{tcolorbox}[enhanced jigsaw, arc=.35mm, bottomrule=.15mm, bottomtitle=1mm, breakable, colback=white, colbacktitle=quarto-callout-note-color!10!white, colframe=quarto-callout-note-color-frame, coltitle=black, left=2mm, leftrule=.75mm, opacityback=0, opacitybacktitle=0.6, rightrule=.15mm, title=\textcolor{quarto-callout-note-color}{\faInfo}\hspace{0.5em}{Proof}, titlerule=0mm, toprule=.15mm, toptitle=1mm]

The mean is \[\mathbb{E}[Y]=\mathbb{E}[AX+b]=A\mathbb{E}[X]+b=A\mu+b.\]
The covariance is
\[\operatorname{Var}(Y)=\operatorname{Var}(AX+b)=A\operatorname{Var}(X)A^T=A\Sigma A^T.\]
Because any affine transformation of a Gaussian vector is Gaussian,
these two quantities determine the distribution.

\end{tcolorbox}

\textbf{Example 32} (Linear combination of a Gaussian vector). Let
\(X\sim \mathcal{N}(\mu,\Sigma)\) and let \(a\) be a fixed vector. Find
the distribution of \(Y=a^TX\).

\begin{tcolorbox}[enhanced jigsaw, arc=.35mm, bottomrule=.15mm, bottomtitle=1mm, breakable, colback=white, colbacktitle=quarto-callout-note-color!10!white, colframe=quarto-callout-note-color-frame, coltitle=black, left=2mm, leftrule=.75mm, opacityback=0, opacitybacktitle=0.6, rightrule=.15mm, title=\textcolor{quarto-callout-note-color}{\faInfo}\hspace{0.5em}{Solution}, titlerule=0mm, toprule=.15mm, toptitle=1mm]

This is the affine transformation with \(A=a^T\) and \(b=0\). Hence
\[Y=a^TX\sim \operatorname{Normal}(a^T\mu,a^T\Sigma a).\]

\end{tcolorbox}

\subsection*{Products and convolutions of Gaussian
densities}\label{products-and-convolutions-of-gaussian-densities}
\addcontentsline{toc}{subsection}{Products and convolutions of Gaussian
densities}

Products and convolutions of Gaussian densities are central in Bayesian
updating, filtering, and linear models.

\textbf{Theorem 33} (Pointwise product of two Gaussian densities).
\emph{Consider two Gaussian densities over the same variable \(x\):
\[p_1(x)=\mathcal{N}(x\mid\mu_1,\Sigma_1),
\qquad
p_2(x)=\mathcal{N}(x\mid\mu_2,\Sigma_2).\] Their pointwise product is
proportional to another Gaussian density:
\[p_1(x)p_2(x)\propto \mathcal{N}(x\mid\mu,\Sigma),\] where
\[\Sigma^{-1}=\Sigma_1^{-1}+\Sigma_2^{-1},
\qquad
\mu=\Sigma\left(\Sigma_1^{-1}\mu_1+\Sigma_2^{-1}\mu_2\right).\]}

\begin{tcolorbox}[enhanced jigsaw, arc=.35mm, bottomrule=.15mm, bottomtitle=1mm, breakable, colback=white, colbacktitle=quarto-callout-note-color!10!white, colframe=quarto-callout-note-color-frame, coltitle=black, left=2mm, leftrule=.75mm, opacityback=0, opacitybacktitle=0.6, rightrule=.15mm, title=\textcolor{quarto-callout-note-color}{\faInfo}\hspace{0.5em}{Proof}, titlerule=0mm, toprule=.15mm, toptitle=1mm]

The product has exponent \[-\frac12(x-\mu_1)^T\Sigma_1^{-1}(x-\mu_1)
-\frac12(x-\mu_2)^T\Sigma_2^{-1}(x-\mu_2).\] Collecting quadratic and
linear terms in \(x\) gives \[-\frac12 x^T(\Sigma_1^{-1}+\Sigma_2^{-1})x
+x^T(\Sigma_1^{-1}\mu_1+\Sigma_2^{-1}\mu_2)+\text{constant}.\]
Completing the square gives the stated covariance and mean.

\end{tcolorbox}

\textbf{Theorem 34} (Convolution and sums). \emph{If
\(X\sim \mathcal{N}(\mu_1,\Sigma_1)\) and
\(Y\sim \mathcal{N}(\mu_2,\Sigma_2)\) are independent, then
\[Z=X+Y\sim \mathcal{N}(\mu_1+\mu_2,\Sigma_1+\Sigma_2).\] Equivalently,
the convolution \[p_Z(z)=\int p_X(x)p_Y(z-x)\,dx\] is Gaussian.}

\begin{tcolorbox}[enhanced jigsaw, arc=.35mm, bottomrule=.15mm, bottomtitle=1mm, breakable, colback=white, colbacktitle=quarto-callout-note-color!10!white, colframe=quarto-callout-note-color-frame, coltitle=black, left=2mm, leftrule=.75mm, opacityback=0, opacitybacktitle=0.6, rightrule=.15mm, title=\textcolor{quarto-callout-note-color}{\faInfo}\hspace{0.5em}{Proof}, titlerule=0mm, toprule=.15mm, toptitle=1mm]

Since \(X\) and \(Y\) are independent,
\[\mathbb{E}[Z]=\mathbb{E}[X]+\mathbb{E}[Y]=\mu_1+\mu_2,\] and
\[\operatorname{Var}(Z)=\operatorname{Var}(X)+\operatorname{Var}(Y)=\Sigma_1+\Sigma_2.\]
The sum of independent Gaussian random vectors is Gaussian, so the
distribution is determined by these two quantities.

\end{tcolorbox}

\textbf{Example 35} (Bayesian normal update as product of Gaussians).
Suppose a prior density for a scalar parameter is
\(\theta\sim \operatorname{Normal}(\mu_0,\sigma_0^2)\), and a Gaussian
likelihood kernel is proportional to
\(\operatorname{Normal}(\theta\mid y,\sigma^2)\). Find the posterior
variance and mean.

\begin{tcolorbox}[enhanced jigsaw, arc=.35mm, bottomrule=.15mm, bottomtitle=1mm, breakable, colback=white, colbacktitle=quarto-callout-note-color!10!white, colframe=quarto-callout-note-color-frame, coltitle=black, left=2mm, leftrule=.75mm, opacityback=0, opacitybacktitle=0.6, rightrule=.15mm, title=\textcolor{quarto-callout-note-color}{\faInfo}\hspace{0.5em}{Solution}, titlerule=0mm, toprule=.15mm, toptitle=1mm]

Using the product formula in one dimension,
\[\frac{1}{\sigma_{post}^2}=\frac{1}{\sigma_0^2}+\frac{1}{\sigma^2}.\]
Thus
\[\sigma_{post}^2=\left(\frac{1}{\sigma_0^2}+\frac{1}{\sigma^2}\right)^{-1}.\]
The posterior mean is
\[\mu_{post}=\sigma_{post}^2\left(\frac{\mu_0}{\sigma_0^2}+\frac{y}{\sigma^2}\right).\]
This is a precision-weighted average of the prior mean and the data
value.

\end{tcolorbox}

\section*{Linear Gaussian Models}\label{linear-gaussian-models}
\addcontentsline{toc}{section}{Linear Gaussian Models}

\markright{Linear Gaussian Models}

This section combines marginal and conditional Gaussian formulas in a
common model used in Bayesian statistics and machine learning.

\subsection*{Marginal and conditional distributions in a linear Gaussian
model}\label{marginal-and-conditional-distributions-in-a-linear-gaussian-model}
\addcontentsline{toc}{subsection}{Marginal and conditional distributions
in a linear Gaussian model}

A linear Gaussian model assumes a Gaussian prior on an input vector and
a Gaussian conditional distribution for an output vector.

\textbf{Theorem 36} (Marginal and conditional Gaussians in a linear
model). \emph{Suppose \[p(x)=\mathcal{N}(x\mid\mu,\Lambda^{-1}),
\qquad
p(y\mid x)=\mathcal{N}(y\mid Ax+b,L^{-1}),\] where \(\Lambda\) and \(L\)
are precision matrices. Then \[p(y)=\mathcal{N}\left(y\mid A\mu+b,
L^{-1}+A\Lambda^{-1}A^T\right).\] Moreover,
\[p(x\mid y)=\mathcal{N}(x\mid m,S),\] where \[S=(\Lambda+A^TLA)^{-1},\]
and \[m=S\left(A^TL(y-b)+\Lambda\mu\right).\]}

\begin{tcolorbox}[enhanced jigsaw, arc=.35mm, bottomrule=.15mm, bottomtitle=1mm, breakable, colback=white, colbacktitle=quarto-callout-note-color!10!white, colframe=quarto-callout-note-color-frame, coltitle=black, left=2mm, leftrule=.75mm, opacityback=0, opacitybacktitle=0.6, rightrule=.15mm, title=\textcolor{quarto-callout-note-color}{\faInfo}\hspace{0.5em}{Proof}, titlerule=0mm, toprule=.15mm, toptitle=1mm]

Write the model as \[y=Ax+b+\varepsilon,
\qquad
\varepsilon\sim \mathcal{N}(0,L^{-1}),
\qquad
x\sim \mathcal{N}(\mu,\Lambda^{-1}),\] with \(x\) independent of
\(\varepsilon\). Therefore \[\mathbb{E}[y]=A\mathbb{E}[x]+b=A\mu+b,\]
and
\[\operatorname{Var}(y)=A\operatorname{Var}(x)A^T+\operatorname{Var}(\varepsilon)=A\Lambda^{-1}A^T+L^{-1}.\]
This gives the marginal distribution of \(y\).

For the conditional distribution, multiply the prior and likelihood as
functions of \(x\): \[p(x\mid y)\propto
\exp\left\{-\frac12(x-\mu)^T\Lambda(x-\mu)
-\frac12(y-Ax-b)^TL(y-Ax-b)\right\}.\] Collecting quadratic terms in
\(x\) gives precision \[S^{-1}=\Lambda+A^TLA.\] Collecting linear terms
gives \[S^{-1}m=\Lambda\mu+A^TL(y-b),\] so
\[m=S\left(A^TL(y-b)+\Lambda\mu\right).\]

\end{tcolorbox}

\textbf{Example 37} (Scalar linear Gaussian model). Let
\[X\sim \operatorname{Normal}(\mu,\tau^2),
\qquad
Y\mid X=x\sim \operatorname{Normal}(ax+b,\sigma^2).\] Find the marginal
distribution of \(Y\).

\begin{tcolorbox}[enhanced jigsaw, arc=.35mm, bottomrule=.15mm, bottomtitle=1mm, breakable, colback=white, colbacktitle=quarto-callout-note-color!10!white, colframe=quarto-callout-note-color-frame, coltitle=black, left=2mm, leftrule=.75mm, opacityback=0, opacitybacktitle=0.6, rightrule=.15mm, title=\textcolor{quarto-callout-note-color}{\faInfo}\hspace{0.5em}{Solution}, titlerule=0mm, toprule=.15mm, toptitle=1mm]

We can write \[Y=aX+b+\varepsilon,
\qquad \varepsilon\sim \operatorname{Normal}(0,\sigma^2),\] with
\(\varepsilon\) independent of \(X\). Hence \[\mathbb{E}[Y]=a\mu+b,\]
and \[\operatorname{Var}(Y)=a^2\tau^2+\sigma^2.\] Therefore
\[Y\sim \operatorname{Normal}(a\mu+b,a^2\tau^2+\sigma^2).\]

\end{tcolorbox}

\textbf{Practice Problem 38} (Posterior in the scalar linear Gaussian
model). For the same model, \[X\sim \operatorname{Normal}(\mu,\tau^2),
\qquad
Y\mid X=x\sim \operatorname{Normal}(ax+b,\sigma^2),\] find the
conditional distribution of \(X\mid Y=y\).

\begin{tcolorbox}[enhanced jigsaw, arc=.35mm, bottomrule=.15mm, bottomtitle=1mm, breakable, colback=white, colbacktitle=quarto-callout-note-color!10!white, colframe=quarto-callout-note-color-frame, coltitle=black, left=2mm, leftrule=.75mm, opacityback=0, opacitybacktitle=0.6, rightrule=.15mm, title=\textcolor{quarto-callout-note-color}{\faInfo}\hspace{0.5em}{Solution}, titlerule=0mm, toprule=.15mm, toptitle=1mm]

Here the prior precision is \(1/\tau^2\) and the observation precision
is \(1/\sigma^2\). The posterior precision is
\[\frac{1}{s^2}=\frac{1}{\tau^2}+\frac{a^2}{\sigma^2}.\] Thus
\[s^2=\left(\frac{1}{\tau^2}+\frac{a^2}{\sigma^2}\right)^{-1}.\] The
posterior mean is
\[m=s^2\left(\frac{\mu}{\tau^2}+\frac{a(y-b)}{\sigma^2}\right).\]
Therefore \[X\mid Y=y\sim \operatorname{Normal}(m,s^2).\]

\end{tcolorbox}

\section*{Summary and Practice}\label{summary-and-practice}
\addcontentsline{toc}{section}{Summary and Practice}

\markright{Summary and Practice}

This final section summarizes the main formulas and provides practice
problems that reinforce the section.

\subsection*{Core formulas}\label{core-formulas}
\addcontentsline{toc}{subsection}{Core formulas}

The main ideas of this section are that multinomial counts live on a
constrained simplex of counts, Dirichlet distributions live on a simplex
of probabilities, and Gaussian random vectors are stable under many
operations.

\begin{tcolorbox}[enhanced jigsaw, arc=.35mm, bottomrule=.15mm, bottomtitle=1mm, breakable, colback=white, colbacktitle=quarto-callout-tip-color!10!white, colframe=quarto-callout-tip-color-frame, coltitle=black, left=2mm, leftrule=.75mm, opacityback=0, opacitybacktitle=0.6, rightrule=.15mm, title=\textcolor{quarto-callout-tip-color}{\faLightbulb}\hspace{0.5em}{Multinomial and Dirichlet formulas}, titlerule=0mm, toprule=.15mm, toptitle=1mm]

\[\mathbb{P}(X_1=n_1,\ldots,X_K=n_K)
=\frac{n!}{n_1!\cdots n_K!}\prod_{k=1}^K\phi_k^{n_k}.\]
\[p(\mu\mid\alpha)=
 \frac{\Gamma(\alpha_0)}{\prod_{k=1}^K\Gamma(\alpha_k)}
 \prod_{k=1}^K \mu_k^{\alpha_k-1},
 \qquad \alpha_0=\sum_{k=1}^K\alpha_k.\] \[\mu\mid X=(n_1,
\ldots,n_K)
\sim \operatorname{Dirichlet}(\alpha_1+n_1,
\ldots,
\alpha_K+n_K).\]

\end{tcolorbox}

\begin{tcolorbox}[enhanced jigsaw, arc=.35mm, bottomrule=.15mm, bottomtitle=1mm, breakable, colback=white, colbacktitle=quarto-callout-tip-color!10!white, colframe=quarto-callout-tip-color-frame, coltitle=black, left=2mm, leftrule=.75mm, opacityback=0, opacitybacktitle=0.6, rightrule=.15mm, title=\textcolor{quarto-callout-tip-color}{\faLightbulb}\hspace{0.5em}{Multivariate Gaussian formulas}, titlerule=0mm, toprule=.15mm, toptitle=1mm]

\[p(x\mid\mu,\Sigma)=\frac{1}{\sqrt{(2\pi)^d|\Sigma|}}
\exp\left\{-\frac12(x-\mu)^T\Sigma^{-1}(x-\mu)\right\}.\]
\[X_1\mid X_2=x_2\sim \mathcal{N}\left(\mu_1+\Sigma_{12}\Sigma_{22}^{-1}(x_2-\mu_2),
\Sigma_{11}-\Sigma_{12}\Sigma_{22}^{-1}\Sigma_{21}\right).\]
\[AX+b\sim \mathcal{N}(A\mu+b,A\Sigma A^T).\]

\end{tcolorbox}

\textbf{Practice Problem 39} (Multinomial probability). A four-sided die
has probabilities \((0.1,0.2,0.3,0.4)\). It is tossed \(8\) times. Find
the probability that the counts are \((1,2,2,3)\).

\begin{tcolorbox}[enhanced jigsaw, arc=.35mm, bottomrule=.15mm, bottomtitle=1mm, breakable, colback=white, colbacktitle=quarto-callout-note-color!10!white, colframe=quarto-callout-note-color-frame, coltitle=black, left=2mm, leftrule=.75mm, opacityback=0, opacitybacktitle=0.6, rightrule=.15mm, title=\textcolor{quarto-callout-note-color}{\faInfo}\hspace{0.5em}{Solution}, titlerule=0mm, toprule=.15mm, toptitle=1mm]

Use the multinomial pmf: \[\mathbb{P}(X=(1,2,2,3))=
\frac{8!}{1!2!2!3!}(0.1)^1(0.2)^2(0.3)^2(0.4)^3.\] The coefficient is
\[\frac{8!}{1!2!2!3!}=1680.\] Thus
\[\mathbb{P}(X=(1,2,2,3))=1680(0.1)(0.04)(0.09)(0.064)
\approx 0.0387.\]

\end{tcolorbox}

\textbf{Practice Problem 40} (Dirichlet posterior mean). Suppose
\(\mu\sim \operatorname{Dirichlet}(3,1,2)\) and the observed counts are
\((2,4,1)\). Find the posterior distribution and posterior mean.

\begin{tcolorbox}[enhanced jigsaw, arc=.35mm, bottomrule=.15mm, bottomtitle=1mm, breakable, colback=white, colbacktitle=quarto-callout-note-color!10!white, colframe=quarto-callout-note-color-frame, coltitle=black, left=2mm, leftrule=.75mm, opacityback=0, opacitybacktitle=0.6, rightrule=.15mm, title=\textcolor{quarto-callout-note-color}{\faInfo}\hspace{0.5em}{Solution}, titlerule=0mm, toprule=.15mm, toptitle=1mm]

The posterior is
\[\mu\mid X\sim \operatorname{Dirichlet}(3+2,1+4,2+1)=\operatorname{Dirichlet}(5,5,3).\]
The posterior total is \(13\), so
\[\mathbb{E}[\mu\mid X]=\left(\frac{5}{13},\frac{5}{13},\frac{3}{13}\right).\]

\end{tcolorbox}

\textbf{Practice Problem 41} (Conditional normal calculation). Let
\[\begin{pmatrix}X_1\\X_2\end{pmatrix}
\sim
\mathcal{N}\left(
\begin{pmatrix}2\\1\end{pmatrix},
\begin{pmatrix}9&3\\3&4\end{pmatrix}
\right).\] Find \(X_1\mid X_2=5\).

\begin{tcolorbox}[enhanced jigsaw, arc=.35mm, bottomrule=.15mm, bottomtitle=1mm, breakable, colback=white, colbacktitle=quarto-callout-note-color!10!white, colframe=quarto-callout-note-color-frame, coltitle=black, left=2mm, leftrule=.75mm, opacityback=0, opacitybacktitle=0.6, rightrule=.15mm, title=\textcolor{quarto-callout-note-color}{\faInfo}\hspace{0.5em}{Solution}, titlerule=0mm, toprule=.15mm, toptitle=1mm]

We have \[\mu_1=2,
\quad \mu_2=1,
\quad \Sigma_{11}=9,
\quad \Sigma_{12}=3,
\quad \Sigma_{22}=4,
\quad \Sigma_{21}=3.\] The conditional mean is
\[\mu_{1\mid 2}=2+3\cdot 4^{-1}(5-1)=2+3=5.\] The conditional variance
is \[\Sigma_{11\mid 2}=9-3\cdot 4^{-1}\cdot 3=9-\frac94=\frac{27}{4}.\]
Therefore
\[X_1\mid X_2=5\sim \operatorname{Normal}\left(5,\frac{27}{4}\right).\]

\end{tcolorbox}

\textbf{Practice Problem 42} (Affine transformation). Let
\(X\sim \mathcal{N}_2\left(\begin{pmatrix}1\\2\end{pmatrix},\begin{pmatrix}1&0.5\\0.5&2\end{pmatrix}\right)\)
and let \(Y=2X_1-X_2+3\). Find the distribution of \(Y\).

\begin{tcolorbox}[enhanced jigsaw, arc=.35mm, bottomrule=.15mm, bottomtitle=1mm, breakable, colback=white, colbacktitle=quarto-callout-note-color!10!white, colframe=quarto-callout-note-color-frame, coltitle=black, left=2mm, leftrule=.75mm, opacityback=0, opacitybacktitle=0.6, rightrule=.15mm, title=\textcolor{quarto-callout-note-color}{\faInfo}\hspace{0.5em}{Solution}, titlerule=0mm, toprule=.15mm, toptitle=1mm]

Write \(Y=a^TX+3\), where \[a=\begin{pmatrix}2\\-1\end{pmatrix}.\] The
mean is \[\mathbb{E}[Y]=a^T\mu+3=2(1)-1(2)+3=3.\] The variance is
\[\operatorname{Var}(Y)=a^T\Sigma a.\] Compute \[\Sigma a=
\begin{pmatrix}1&0.5\\0.5&2\end{pmatrix}
\begin{pmatrix}2\\-1\end{pmatrix}
=\begin{pmatrix}1.5\\-1\end{pmatrix}.\] Thus
\[a^T\Sigma a=(2,-1)\begin{pmatrix}1.5\\-1\end{pmatrix}=3+1=4.\]
Therefore \[Y\sim \operatorname{Normal}(3,4).\]

\end{tcolorbox}

9 George Casella and Roger L. Berger, \emph{Statistical Inference}, 2nd
edition. Larry Wasserman, \emph{All of Statistics}. C. M. Grinstead and
J. L. Snell, \emph{Introduction to Probability}, American Mathematical
Society, 2012. Sheldon Ross, \emph{Introduction to Probability Models},
12th edition.

\bookmarksetup{startatroot}

\chapter{Chapter 7: Inequalities and
Identities}\label{chapter-7-inequalities-and-identities}

This chapter introduces the main probability inequalities used to
control tail probabilities, prove convergence results, and analyze
estimators. The focus is not only on remembering formulas, but also on
learning when each inequality is the right tool.

\textbf{Topics.} Boole's inequality; inclusion-exclusion; Bonferroni
inequalities; Markov's inequality; Chebyshev's inequality; Chernoff
bounds; Hoeffding's inequality; Cauchy-Schwarz, Holder, Minkowski, and
Jensen inequalities; equality identities for common distributions.

\subsection{Overview}\label{overview-4}

This section introduces probability inequalities and distribution
identities that are used to control tail probabilities, prove limit
theorems, and analyze statistical estimators.

Probability inequalities provide upper and lower bounds for
probabilities of events, especially events involving deviations of
random variables from their expectations. The important point is that a
useful bound may require only partial information, such as the mean,
variance, support, or moment generating function, rather than a full
distributional calculation.

\begin{tcolorbox}[enhanced jigsaw, arc=.35mm, bottomrule=.15mm, bottomtitle=1mm, breakable, colback=white, colbacktitle=quarto-callout-tip-color!10!white, colframe=quarto-callout-tip-color-frame, coltitle=black, left=2mm, leftrule=.75mm, opacityback=0, opacitybacktitle=0.6, rightrule=.15mm, title=\textcolor{quarto-callout-tip-color}{\faLightbulb}\hspace{0.5em}{Main message}, titlerule=0mm, toprule=.15mm, toptitle=1mm]

Main message Probability inequalities answer questions of the form
\[\mathbb{P}\{Z \geq \mathbb{E}[Z]+t\},\qquad \mathbb{P}\{Z\leq \mathbb{E}[Z]-t\},\qquad \mathbb{P}\{|Z-\mathbb{E}[Z]|\geq t\}.\]
They are useful when exact probabilities are hard to compute or when we
want distribution-free guarantees.

\end{tcolorbox}

The inequalities in this section include:

\begin{itemize}
\item
  union bounds: Boole's inequality, inclusion-exclusion, and Bonferroni
  inequalities;
\item
  tail bounds from moments: Markov's inequality and Chebyshev's
  inequality;
\item
  exponential tail bounds: Chernoff bounds and Hoeffding's inequality;
\item
  functional inequalities: Cauchy-Schwarz, Holder, Minkowski, and
  Jensen;
\item
  equality identities for special distributions.
\end{itemize}

\section{Union Bounds and
Inclusion-Exclusion}\label{union-bounds-and-inclusion-exclusion}

This section begins with inequalities for the probability that at least
one event in a collection occurs.

\subsection{Boole's inequality}\label{booles-inequality}

This subsection states the most basic and most useful union bound.

\begin{tcolorbox}[enhanced jigsaw, arc=.35mm, bottomrule=.15mm, bottomtitle=1mm, breakable, colback=white, colbacktitle=quarto-callout-important-color!10!white, colframe=quarto-callout-important-color-frame, coltitle=black, left=2mm, leftrule=.75mm, opacityback=0, opacitybacktitle=0.6, rightrule=.15mm, title=\textcolor{quarto-callout-important-color}{\faExclamation}\hspace{0.5em}{Theorem}, titlerule=0mm, toprule=.15mm, toptitle=1mm]

\textbf{Theorem 1} (Boole's inequality). \emph{Suppose
\((S,\mathcal B,\mathbb{P})\) is a probability space and
\(E_1,E_2,\ldots\in \mathcal B\) are events. Then
\[\mathbb{P}\left(\bigcup_{i=1}^{\infty}E_i\right)
\leq
\sum_{i=1}^{\infty}\mathbb{P}(E_i).\] For finitely many events,
\[\mathbb{P}\left(\bigcup_{i=1}^{n}E_i\right)
\leq
\sum_{i=1}^{n}\mathbb{P}(E_i).\]}

\end{tcolorbox}

Boole's inequality says that the probability that at least one of the
events happens is no greater than the sum of the individual
probabilities.

\begin{tcolorbox}[enhanced jigsaw, arc=.35mm, bottomrule=.15mm, bottomtitle=1mm, breakable, colback=white, colbacktitle=quarto-callout-note-color!10!white, colframe=quarto-callout-note-color-frame, coltitle=black, left=2mm, leftrule=.75mm, opacityback=0, opacitybacktitle=0.6, rightrule=.15mm, title=\textcolor{quarto-callout-note-color}{\faInfo}\hspace{0.5em}{Proof}, titlerule=0mm, toprule=.15mm, toptitle=1mm]

\emph{Proof.} For two events,
\[\mathbb{P}(A\cup B)=\mathbb{P}(A)+\mathbb{P}(B)-\mathbb{P}(A\cap B)\leq \mathbb{P}(A)+\mathbb{P}(B).\]
The finite result follows by induction. The countable result follows by
continuity of probability from below:
\[\mathbb{P}\left(\bigcup_{i=1}^{\infty}E_i\right)
=
\lim_{n\to\infty}\mathbb{P}\left(\bigcup_{i=1}^{n}E_i\right)
\leq
\lim_{n\to\infty}\sum_{i=1}^{n}\mathbb{P}(E_i).\]~◻

\end{tcolorbox}

\begin{tcolorbox}[enhanced jigsaw, arc=.35mm, bottomrule=.15mm, bottomtitle=1mm, breakable, colback=white, colbacktitle=quarto-callout-tip-color!10!white, colframe=quarto-callout-tip-color-frame, coltitle=black, left=2mm, leftrule=.75mm, opacityback=0, opacitybacktitle=0.6, rightrule=.15mm, title=\textcolor{quarto-callout-tip-color}{\faLightbulb}\hspace{0.5em}{Example}, titlerule=0mm, toprule=.15mm, toptitle=1mm]

\textbf{Example 2} (A quick union-bound estimate). Suppose \(100\)
students each independently submit an assignment late with probability
at most \(0.01\). Without assuming independence, bound the probability
that at least one student submits late.

\end{tcolorbox}

\begin{tcolorbox}[enhanced jigsaw, arc=.35mm, bottomrule=.15mm, bottomtitle=1mm, breakable, colback=white, colbacktitle=quarto-callout-note-color!10!white, colframe=quarto-callout-note-color-frame, coltitle=black, left=2mm, leftrule=.75mm, opacityback=0, opacitybacktitle=0.6, rightrule=.15mm, title=\textcolor{quarto-callout-note-color}{\faInfo}\hspace{0.5em}{Solution}, titlerule=0mm, toprule=.15mm, toptitle=1mm]

Let \(E_i\) be the event that student \(i\) submits late. Then
\[\mathbb{P}\left(\bigcup_{i=1}^{100}E_i\right)
\leq \sum_{i=1}^{100}\mathbb{P}(E_i)
\leq 100(0.01)=1.\] This bound is not informative because it only gives
an upper bound of \(1\). If the late probability were \(0.001\), then
Boole's inequality would give
\[\mathbb{P}(\text{at least one late})\leq 100(0.001)=0.1.\] The example
shows that the union bound is most useful when individual probabilities
are small.

\end{tcolorbox}

\subsection{Inclusion-exclusion}\label{inclusion-exclusion}

This subsection gives the exact formula for a finite union by adding and
subtracting intersections.

Suppose \(E_1,\ldots,E_n\) are events. Define
\[S_1=\sum_{i=1}^{n}\mathbb{P}(E_i),\]
\[S_2=\sum_{1\leq i<j\leq n}\mathbb{P}(E_i\cap E_j),\] and, more
generally, \[S_k=\sum_{1\leq i_1<\cdots<i_k\leq n}
\mathbb{P}(E_{i_1}\cap\cdots\cap E_{i_k}).\]

\begin{tcolorbox}[enhanced jigsaw, arc=.35mm, bottomrule=.15mm, bottomtitle=1mm, breakable, colback=white, colbacktitle=quarto-callout-important-color!10!white, colframe=quarto-callout-important-color-frame, coltitle=black, left=2mm, leftrule=.75mm, opacityback=0, opacitybacktitle=0.6, rightrule=.15mm, title=\textcolor{quarto-callout-important-color}{\faExclamation}\hspace{0.5em}{Theorem}, titlerule=0mm, toprule=.15mm, toptitle=1mm]

\textbf{Theorem 3} (Inclusion-exclusion). \emph{For events
\(E_1,\ldots,E_n\), \[\mathbb{P}\left(\bigcup_{i=1}^{n}E_i\right)
=S_1-S_2+S_3-\cdots+(-1)^{n-1}S_n.\]}

\end{tcolorbox}

\begin{tcolorbox}[enhanced jigsaw, arc=.35mm, bottomrule=.15mm, bottomtitle=1mm, breakable, colback=white, colbacktitle=quarto-callout-tip-color!10!white, colframe=quarto-callout-tip-color-frame, coltitle=black, left=2mm, leftrule=.75mm, opacityback=0, opacitybacktitle=0.6, rightrule=.15mm, title=\textcolor{quarto-callout-tip-color}{\faLightbulb}\hspace{0.5em}{Example}, titlerule=0mm, toprule=.15mm, toptitle=1mm]

\textbf{Example 4} (Three events). Write the inclusion-exclusion formula
for three events \(A,B,C\).

\end{tcolorbox}

\begin{tcolorbox}[enhanced jigsaw, arc=.35mm, bottomrule=.15mm, bottomtitle=1mm, breakable, colback=white, colbacktitle=quarto-callout-note-color!10!white, colframe=quarto-callout-note-color-frame, coltitle=black, left=2mm, leftrule=.75mm, opacityback=0, opacitybacktitle=0.6, rightrule=.15mm, title=\textcolor{quarto-callout-note-color}{\faInfo}\hspace{0.5em}{Solution}, titlerule=0mm, toprule=.15mm, toptitle=1mm]

For three events, \[\mathbb{P}(A\cup B\cup C)
=
\mathbb{P}(A)+\mathbb{P}(B)+\mathbb{P}(C)
-\mathbb{P}(A\cap B)-\mathbb{P}(A\cap C)-\mathbb{P}(B\cap C)
+\mathbb{P}(A\cap B\cap C).\] The first line adds the probability of
being in any one event. The second line removes double-counting. The
last term adds back the triple intersection, which was removed too many
times.

\end{tcolorbox}

\subsection{Bonferroni inequalities}\label{bonferroni-inequalities}

This subsection turns the inclusion-exclusion expansion into useful
upper and lower bounds.

\begin{tcolorbox}[enhanced jigsaw, arc=.35mm, bottomrule=.15mm, bottomtitle=1mm, breakable, colback=white, colbacktitle=quarto-callout-important-color!10!white, colframe=quarto-callout-important-color-frame, coltitle=black, left=2mm, leftrule=.75mm, opacityback=0, opacitybacktitle=0.6, rightrule=.15mm, title=\textcolor{quarto-callout-important-color}{\faExclamation}\hspace{0.5em}{Theorem}, titlerule=0mm, toprule=.15mm, toptitle=1mm]

\textbf{Theorem 5} (Bonferroni inequalities). \emph{Let
\(E_1,\ldots,E_n\) be events and let \(S_k\) be defined as above. For
\(1\leq k\leq n\), \[B_k=\sum_{j=1}^{k}(-1)^{j-1}S_j.\] If \(k\) is odd,
then \[\mathbb{P}\left(\bigcup_{i=1}^{n}E_i\right)\leq B_k.\] If \(k\)
is even, then \[\mathbb{P}\left(\bigcup_{i=1}^{n}E_i\right)\geq B_k.\]}

\end{tcolorbox}

The case \(k=1\) is exactly Boole's inequality. The case \(k=2\) gives
\[\mathbb{P}\left(\bigcup_{i=1}^{n}E_i\right)
\geq
\sum_i\mathbb{P}(E_i)-\sum_{i<j}\mathbb{P}(E_i\cap E_j).\]

\begin{tcolorbox}[enhanced jigsaw, arc=.35mm, bottomrule=.15mm, bottomtitle=1mm, breakable, colback=white, colbacktitle=quarto-callout-tip-color!10!white, colframe=quarto-callout-tip-color-frame, coltitle=black, left=2mm, leftrule=.75mm, opacityback=0, opacitybacktitle=0.6, rightrule=.15mm, title=\textcolor{quarto-callout-tip-color}{\faLightbulb}\hspace{0.5em}{Example}, titlerule=0mm, toprule=.15mm, toptitle=1mm]

\textbf{Example 6} (Empty boxes). Place \(n\) distinguishable balls into
\(m\) distinguishable boxes uniformly at random, where \(n>m\). Let
\(E\) be the event that at least one box is empty. Use
inclusion-exclusion and Bonferroni inequalities to estimate
\(\mathbb{P}(E)\).

\end{tcolorbox}

\begin{tcolorbox}[enhanced jigsaw, arc=.35mm, bottomrule=.15mm, bottomtitle=1mm, breakable, colback=white, colbacktitle=quarto-callout-note-color!10!white, colframe=quarto-callout-note-color-frame, coltitle=black, left=2mm, leftrule=.75mm, opacityback=0, opacitybacktitle=0.6, rightrule=.15mm, title=\textcolor{quarto-callout-note-color}{\faInfo}\hspace{0.5em}{Solution}, titlerule=0mm, toprule=.15mm, toptitle=1mm]

The sample space is
\[S=\{\omega=(\omega_1,\ldots,\omega_n):1\leq \omega_i\leq m\},
\qquad |S|=m^n.\] Let
\[E_\ell=\{\omega:\omega_i\neq \ell \text{ for all } i=1,\ldots,n\}\] be
the event that box \(\ell\) is empty. Then \[E=E_1\cup\cdots\cup E_m.\]
For any fixed \(k\) distinct boxes \(i_1,\ldots,i_k\),
\[\mathbb{P}(E_{i_1}\cap\cdots\cap E_{i_k})
=\left(\frac{m-k}{m}\right)^n
=\left(1-\frac{k}{m}\right)^n.\] Therefore inclusion-exclusion gives
\[\mathbb{P}(E)
=
\binom{m}{1}\left(1-\frac1m\right)^n
-
\binom{m}{2}\left(1-\frac2m\right)^n
+\cdots
+(-1)^{m-1}\binom{m}{m}\left(1-\frac{m}{m}\right)^n.\] The first two
Bonferroni bounds give \[m\left(1-\frac1m\right)^n
-
\binom{m}{2}\left(1-\frac2m\right)^n
\leq \mathbb{P}(E)
\leq
m\left(1-\frac1m\right)^n.\] For \(m=10\),
\[10(0.9)^n-45(0.8)^n\leq \mathbb{P}(E)\leq 10(0.9)^n.\] For \(n=50\),
the gap between the upper and lower bounds is
\[45(0.8)^{50}\approx 0.00064226146,\] so the two-sided estimate is
quite accurate.

\end{tcolorbox}

\begin{tcolorbox}[enhanced jigsaw, arc=.35mm, bottomrule=.15mm, bottomtitle=1mm, breakable, colback=white, colbacktitle=quarto-callout-caution-color!10!white, colframe=quarto-callout-caution-color-frame, coltitle=black, left=2mm, leftrule=.75mm, opacityback=0, opacitybacktitle=0.6, rightrule=.15mm, title=\textcolor{quarto-callout-caution-color}{\faFire}\hspace{0.5em}{Practice Problem}, titlerule=0mm, toprule=.15mm, toptitle=1mm]

\textbf{Practice Problem 7} (At least one duplicated birthday). Assume
each birthday is equally likely among \(365\) days. Let \(E_{ij}\) be
the event that people \(i\) and \(j\) share a birthday in a class of
\(r\) people. Use Boole's inequality to bound the probability that at
least two people share a birthday.

\end{tcolorbox}

\begin{tcolorbox}[enhanced jigsaw, arc=.35mm, bottomrule=.15mm, bottomtitle=1mm, breakable, colback=white, colbacktitle=quarto-callout-note-color!10!white, colframe=quarto-callout-note-color-frame, coltitle=black, left=2mm, leftrule=.75mm, opacityback=0, opacitybacktitle=0.6, rightrule=.15mm, title=\textcolor{quarto-callout-note-color}{\faInfo}\hspace{0.5em}{Solution}, titlerule=0mm, toprule=.15mm, toptitle=1mm]

There are \(\binom r2\) pairs. For each pair,
\[\mathbb{P}(E_{ij})=\frac1{365}.\] If \(E\) is the event that at least
one pair shares a birthday, then \[E=\bigcup_{1\leq i<j\leq r}E_{ij}.\]
By Boole's inequality,
\[\mathbb{P}(E)\leq \sum_{i<j}\mathbb{P}(E_{ij})=\binom r2\frac1{365}.\]
This is a simple upper bound; it can exceed \(1\) for large \(r\), in
which case the best probability upper bound is \(1\).

\end{tcolorbox}

\section{Markov's Inequality}\label{markovs-inequality}

This section introduces the fundamental inequality behind many tail
bounds.

\subsection{Basic Markov inequality}\label{basic-markov-inequality}

This subsection bounds a nonnegative random variable using only its
expectation.

\begin{tcolorbox}[enhanced jigsaw, arc=.35mm, bottomrule=.15mm, bottomtitle=1mm, breakable, colback=white, colbacktitle=quarto-callout-important-color!10!white, colframe=quarto-callout-important-color-frame, coltitle=black, left=2mm, leftrule=.75mm, opacityback=0, opacitybacktitle=0.6, rightrule=.15mm, title=\textcolor{quarto-callout-important-color}{\faExclamation}\hspace{0.5em}{Theorem}, titlerule=0mm, toprule=.15mm, toptitle=1mm]

\textbf{Theorem 8} (Markov's inequality). \emph{Let \(X\) be a
nonnegative random variable. Then, for any \(a>0\),
\[\mathbb{P}(X\geq a)\leq \frac{\mathbb{E}[X]}{a}.\]}

\end{tcolorbox}

\begin{tcolorbox}[enhanced jigsaw, arc=.35mm, bottomrule=.15mm, bottomtitle=1mm, breakable, colback=white, colbacktitle=quarto-callout-note-color!10!white, colframe=quarto-callout-note-color-frame, coltitle=black, left=2mm, leftrule=.75mm, opacityback=0, opacitybacktitle=0.6, rightrule=.15mm, title=\textcolor{quarto-callout-note-color}{\faInfo}\hspace{0.5em}{Proof}, titlerule=0mm, toprule=.15mm, toptitle=1mm]

\emph{Proof for continuous \(X\).} If \(X\) has density \(f\), then
\[\mathbb{E}[X]=\int_0^\infty x f(x)\,dx
=\int_0^a x f(x)\,dx+\int_a^\infty x f(x)\,dx
\geq \int_a^\infty x f(x)\,dx
\geq \int_a^\infty a f(x)\,dx
=a\mathbb{P}(X\geq a).\] Dividing by \(a\) gives the result. The same
proof works in the discrete case by replacing integrals with sums.~◻

\end{tcolorbox}

\begin{tcolorbox}[enhanced jigsaw, arc=.35mm, bottomrule=.15mm, bottomtitle=1mm, breakable, colback=white, colbacktitle=quarto-callout-tip-color!10!white, colframe=quarto-callout-tip-color-frame, coltitle=black, left=2mm, leftrule=.75mm, opacityback=0, opacitybacktitle=0.6, rightrule=.15mm, title=\textcolor{quarto-callout-tip-color}{\faLightbulb}\hspace{0.5em}{Example}, titlerule=0mm, toprule=.15mm, toptitle=1mm]

\textbf{Example 9} (A binomial upper-tail bound). Suppose
\(X\sim \operatorname{Binomial}(n,p)\). For \(p<q<1\), use Markov's
inequality to bound \(\mathbb{P}(X\geq qn)\).

\end{tcolorbox}

\begin{tcolorbox}[enhanced jigsaw, arc=.35mm, bottomrule=.15mm, bottomtitle=1mm, breakable, colback=white, colbacktitle=quarto-callout-note-color!10!white, colframe=quarto-callout-note-color-frame, coltitle=black, left=2mm, leftrule=.75mm, opacityback=0, opacitybacktitle=0.6, rightrule=.15mm, title=\textcolor{quarto-callout-note-color}{\faInfo}\hspace{0.5em}{Solution}, titlerule=0mm, toprule=.15mm, toptitle=1mm]

Since \(X\geq 0\) and \(\mathbb{E}[X]=np\), Markov's inequality gives
\[\mathbb{P}(X\geq qn)
\leq
\frac{\mathbb{E}[X]}{qn}
=
\frac{np}{qn}
=\frac pq.\] This bound does not improve as \(n\) increases, so it is
often weak for sums of many independent variables.

\end{tcolorbox}

\subsection{General Markov inequality}\label{general-markov-inequality}

This subsection applies Markov's idea to a nonnegative function of a
random variable.

\begin{tcolorbox}[enhanced jigsaw, arc=.35mm, bottomrule=.15mm, bottomtitle=1mm, breakable, colback=white, colbacktitle=quarto-callout-important-color!10!white, colframe=quarto-callout-important-color-frame, coltitle=black, left=2mm, leftrule=.75mm, opacityback=0, opacitybacktitle=0.6, rightrule=.15mm, title=\textcolor{quarto-callout-important-color}{\faExclamation}\hspace{0.5em}{Theorem}, titlerule=0mm, toprule=.15mm, toptitle=1mm]

\textbf{Theorem 10} (General Markov inequality). \emph{Let \(g(X)\) be a
nonnegative random variable. For any \(r>0\), \[\mathbb{P}(g(X)\geq r)
\leq
\frac{\mathbb{E}[g(X)]}{r}.\]}

\end{tcolorbox}

\begin{tcolorbox}[enhanced jigsaw, arc=.35mm, bottomrule=.15mm, bottomtitle=1mm, breakable, colback=white, colbacktitle=quarto-callout-note-color!10!white, colframe=quarto-callout-note-color-frame, coltitle=black, left=2mm, leftrule=.75mm, opacityback=0, opacitybacktitle=0.6, rightrule=.15mm, title=\textcolor{quarto-callout-note-color}{\faInfo}\hspace{0.5em}{Proof}, titlerule=0mm, toprule=.15mm, toptitle=1mm]

\emph{Proof.} Apply Markov's inequality to the nonnegative random
variable \(g(X)\).~◻

\end{tcolorbox}

\begin{tcolorbox}[enhanced jigsaw, arc=.35mm, bottomrule=.15mm, bottomtitle=1mm, breakable, colback=white, colbacktitle=quarto-callout-tip-color!10!white, colframe=quarto-callout-tip-color-frame, coltitle=black, left=2mm, leftrule=.75mm, opacityback=0, opacitybacktitle=0.6, rightrule=.15mm, title=\textcolor{quarto-callout-tip-color}{\faLightbulb}\hspace{0.5em}{Example}, titlerule=0mm, toprule=.15mm, toptitle=1mm]

\textbf{Example 11} (Universal deviation bound). Let \(X\) be a random
variable with mean \(\mu\) and variance \(\sigma^2\). Show that
\[\mathbb{P}(|X-\mu|\geq t\sigma)\leq \frac1{t^2}\] for \(t>0\).

\end{tcolorbox}

\begin{tcolorbox}[enhanced jigsaw, arc=.35mm, bottomrule=.15mm, bottomtitle=1mm, breakable, colback=white, colbacktitle=quarto-callout-note-color!10!white, colframe=quarto-callout-note-color-frame, coltitle=black, left=2mm, leftrule=.75mm, opacityback=0, opacitybacktitle=0.6, rightrule=.15mm, title=\textcolor{quarto-callout-note-color}{\faInfo}\hspace{0.5em}{Solution}, titlerule=0mm, toprule=.15mm, toptitle=1mm]

Let \[g(X)=\frac{(X-\mu)^2}{\sigma^2}.\] Then \(g(X)\geq 0\) and
\[\mathbb{E}[g(X)]=\frac{\mathbb{E}[(X-\mu)^2]}{\sigma^2}=1.\] The event
\(|X-\mu|\geq t\sigma\) is the same as
\[\frac{(X-\mu)^2}{\sigma^2}\geq t^2.\] By the general Markov
inequality, \[\mathbb{P}(|X-\mu|\geq t\sigma)
=
\mathbb{P}(g(X)\geq t^2)
\leq \frac{\mathbb{E}[g(X)]}{t^2}
=\frac1{t^2}.\] For example, if \(t=2\), then
\[\mathbb{P}(|X-\mu|\geq 2\sigma)\leq \frac14,\] so there is at least a
\(75\%\) chance that \(X\) lies within two standard deviations of its
mean, no matter what the distribution is.

\end{tcolorbox}

\section{Chebyshev's Inequality}\label{chebyshevs-inequality}

This section improves Markov's inequality by using variance to control
deviations from the mean.

\subsection{Statement and proof}\label{statement-and-proof}

This subsection proves Chebyshev's inequality as a direct consequence of
Markov's inequality.

\begin{tcolorbox}[enhanced jigsaw, arc=.35mm, bottomrule=.15mm, bottomtitle=1mm, breakable, colback=white, colbacktitle=quarto-callout-important-color!10!white, colframe=quarto-callout-important-color-frame, coltitle=black, left=2mm, leftrule=.75mm, opacityback=0, opacitybacktitle=0.6, rightrule=.15mm, title=\textcolor{quarto-callout-important-color}{\faExclamation}\hspace{0.5em}{Theorem}, titlerule=0mm, toprule=.15mm, toptitle=1mm]

\textbf{Theorem 12} (Chebyshev's inequality). \emph{Let \(X\) be a
random variable with finite variance. Then, for any \(k>0\),
\[\mathbb{P}(|X-\mathbb{E}[X]|\geq k)
\leq
\frac{\operatorname{Var}(X)}{k^2}.\] Equivalently, for \(t>0\),
\[\mathbb{P}(|X-\mathbb{E}[X]|\geq t\sqrt{\operatorname{Var}(X)})\leq \frac1{t^2}.\]}

\end{tcolorbox}

\begin{tcolorbox}[enhanced jigsaw, arc=.35mm, bottomrule=.15mm, bottomtitle=1mm, breakable, colback=white, colbacktitle=quarto-callout-note-color!10!white, colframe=quarto-callout-note-color-frame, coltitle=black, left=2mm, leftrule=.75mm, opacityback=0, opacitybacktitle=0.6, rightrule=.15mm, title=\textcolor{quarto-callout-note-color}{\faInfo}\hspace{0.5em}{Proof}, titlerule=0mm, toprule=.15mm, toptitle=1mm]

\emph{Proof.} Let \[Y=(X-\mathbb{E}[X])^2.\] Then \(Y\geq 0\). By
Markov's inequality, \[\mathbb{P}(|X-\mathbb{E}[X]|\geq k)
=\mathbb{P}((X-\mathbb{E}[X])^2\geq k^2)
\leq
\frac{\mathbb{E}[(X-\mathbb{E}[X])^2]}{k^2}
=\frac{\operatorname{Var}(X)}{k^2}.\]~◻

\end{tcolorbox}

\subsection{Comparison with Markov for a binomial
tail}\label{comparison-with-markov-for-a-binomial-tail}

This subsection shows that Chebyshev's inequality improves Markov's
inequality for a binomial upper tail because it uses the variance.

\begin{tcolorbox}[enhanced jigsaw, arc=.35mm, bottomrule=.15mm, bottomtitle=1mm, breakable, colback=white, colbacktitle=quarto-callout-tip-color!10!white, colframe=quarto-callout-tip-color-frame, coltitle=black, left=2mm, leftrule=.75mm, opacityback=0, opacitybacktitle=0.6, rightrule=.15mm, title=\textcolor{quarto-callout-tip-color}{\faLightbulb}\hspace{0.5em}{Example}, titlerule=0mm, toprule=.15mm, toptitle=1mm]

\textbf{Example 13} (Binomial tail: Markov versus Chebyshev). Suppose
\(X\sim \operatorname{Binomial}(n,p)\) and \(p<q<1\). Compare Markov's
and Chebyshev's bounds for \(\mathbb{P}(X\geq qn)\).

\end{tcolorbox}

\begin{tcolorbox}[enhanced jigsaw, arc=.35mm, bottomrule=.15mm, bottomtitle=1mm, breakable, colback=white, colbacktitle=quarto-callout-note-color!10!white, colframe=quarto-callout-note-color-frame, coltitle=black, left=2mm, leftrule=.75mm, opacityback=0, opacitybacktitle=0.6, rightrule=.15mm, title=\textcolor{quarto-callout-note-color}{\faInfo}\hspace{0.5em}{Solution}, titlerule=0mm, toprule=.15mm, toptitle=1mm]

Markov's inequality gives \[\mathbb{P}(X\geq qn)\leq \frac pq.\] For
Chebyshev's inequality, note that \[\mathbb{E}[X]=np,
\qquad
\operatorname{Var}(X)=np(1-p).\] Since \(q>p\),
\[\{X\geq qn\}\subseteq \{|X-np|\geq (q-p)n\}.\] Therefore
\[\mathbb{P}(X\geq qn)
\leq
\mathbb{P}(|X-np|\geq (q-p)n)
\leq
\frac{np(1-p)}{(q-p)^2n^2}
=
\frac{p(1-p)}{(q-p)^2n}.\] Unlike the Markov bound, this bound goes to
\(0\) as \(n\to\infty\).

\end{tcolorbox}

\begin{tcolorbox}[enhanced jigsaw, arc=.35mm, bottomrule=.15mm, bottomtitle=1mm, breakable, colback=white, colbacktitle=quarto-callout-caution-color!10!white, colframe=quarto-callout-caution-color-frame, coltitle=black, left=2mm, leftrule=.75mm, opacityback=0, opacitybacktitle=0.6, rightrule=.15mm, title=\textcolor{quarto-callout-caution-color}{\faFire}\hspace{0.5em}{Practice Problem}, titlerule=0mm, toprule=.15mm, toptitle=1mm]

\textbf{Practice Problem 14} (Sample mean bound). Let \(X_1,\ldots,X_n\)
be iid with mean \(\mu\) and variance \(\sigma^2\). Let
\(\bar X_n=n^{-1}\sum_{i=1}^n X_i\). Use Chebyshev's inequality to bound
\(\mathbb{P}(|\bar X_n-\mu|\geq \varepsilon)\).

\end{tcolorbox}

\begin{tcolorbox}[enhanced jigsaw, arc=.35mm, bottomrule=.15mm, bottomtitle=1mm, breakable, colback=white, colbacktitle=quarto-callout-note-color!10!white, colframe=quarto-callout-note-color-frame, coltitle=black, left=2mm, leftrule=.75mm, opacityback=0, opacitybacktitle=0.6, rightrule=.15mm, title=\textcolor{quarto-callout-note-color}{\faInfo}\hspace{0.5em}{Solution}, titlerule=0mm, toprule=.15mm, toptitle=1mm]

We have \[\mathbb{E}[\bar X_n]=\mu,
\qquad
\operatorname{Var}(\bar X_n)=\frac{\sigma^2}{n}.\] Chebyshev's
inequality gives \[\mathbb{P}(|\bar X_n-\mu|\geq \varepsilon)
\leq
\frac{\operatorname{Var}(\bar X_n)}{\varepsilon^2}
=
\frac{\sigma^2}{n\varepsilon^2}.\] This is the basic second-moment
concentration bound for the sample mean.

\end{tcolorbox}

\section{Chernoff Bounds}\label{chernoff-bounds}

This section introduces exponential tail bounds based on moment
generating functions.

\subsection{The Chernoff method}\label{the-chernoff-method}

This subsection derives the Chernoff bound by applying Markov's
inequality to \(e^{tX}\).

\begin{tcolorbox}[enhanced jigsaw, arc=.35mm, bottomrule=.15mm, bottomtitle=1mm, breakable, colback=white, colbacktitle=quarto-callout-important-color!10!white, colframe=quarto-callout-important-color-frame, coltitle=black, left=2mm, leftrule=.75mm, opacityback=0, opacitybacktitle=0.6, rightrule=.15mm, title=\textcolor{quarto-callout-important-color}{\faExclamation}\hspace{0.5em}{Theorem}, titlerule=0mm, toprule=.15mm, toptitle=1mm]

\textbf{Theorem 15} (Chernoff bounds). \emph{Let \(X\) be a random
variable with moment generating function \[M_X(t)=\mathbb{E}[e^{tX}].\]
For any \(a\in\mathbb{R}\), \[\mathbb{P}(X\geq a)
\leq
\inf_{t\geq 0} e^{-ta}M_X(t),\] and \[\mathbb{P}(X\leq a)
\leq
\inf_{t\geq 0} e^{ta}M_X(-t).\]}

\end{tcolorbox}

\begin{tcolorbox}[enhanced jigsaw, arc=.35mm, bottomrule=.15mm, bottomtitle=1mm, breakable, colback=white, colbacktitle=quarto-callout-note-color!10!white, colframe=quarto-callout-note-color-frame, coltitle=black, left=2mm, leftrule=.75mm, opacityback=0, opacitybacktitle=0.6, rightrule=.15mm, title=\textcolor{quarto-callout-note-color}{\faInfo}\hspace{0.5em}{Proof}, titlerule=0mm, toprule=.15mm, toptitle=1mm]

\emph{Proof.} For \(t\geq 0\), the function \(e^{tX}\) is nonnegative
and \[\{X\geq a\}=\{e^{tX}\geq e^{ta}\}.\] By Markov's inequality,
\[\mathbb{P}(X\geq a)
\leq
\frac{\mathbb{E}[e^{tX}]}{e^{ta}}
=e^{-ta}M_X(t).\] Taking the infimum over \(t\geq 0\) gives the
upper-tail bound. For the lower tail, apply the same argument to \(-X\):
\[\mathbb{P}(X\leq a)=\mathbb{P}(-X\geq -a)
\leq e^{ta}M_X(-t).\]~◻

\end{tcolorbox}

\subsection{Binomial Chernoff bound}\label{binomial-chernoff-bound}

This subsection applies the Chernoff method to a binomial random
variable.

\begin{tcolorbox}[enhanced jigsaw, arc=.35mm, bottomrule=.15mm, bottomtitle=1mm, breakable, colback=white, colbacktitle=quarto-callout-tip-color!10!white, colframe=quarto-callout-tip-color-frame, coltitle=black, left=2mm, leftrule=.75mm, opacityback=0, opacitybacktitle=0.6, rightrule=.15mm, title=\textcolor{quarto-callout-tip-color}{\faLightbulb}\hspace{0.5em}{Example}, titlerule=0mm, toprule=.15mm, toptitle=1mm]

\textbf{Example 16} (Chernoff bound for a binomial tail). Let
\(X\sim \operatorname{Binomial}(n,p)\). For \(p<q<1\), derive an
exponential upper bound for \(\mathbb{P}(X\geq qn)\).

\end{tcolorbox}

\begin{tcolorbox}[enhanced jigsaw, arc=.35mm, bottomrule=.15mm, bottomtitle=1mm, breakable, colback=white, colbacktitle=quarto-callout-note-color!10!white, colframe=quarto-callout-note-color-frame, coltitle=black, left=2mm, leftrule=.75mm, opacityback=0, opacitybacktitle=0.6, rightrule=.15mm, title=\textcolor{quarto-callout-note-color}{\faInfo}\hspace{0.5em}{Solution}, titlerule=0mm, toprule=.15mm, toptitle=1mm]

The moment generating function of \(X\sim \operatorname{Binomial}(n,p)\)
is \[M_X(t)=(pe^t+1-p)^n.\] The Chernoff bound gives
\[\mathbb{P}(X\geq qn)
\leq
\inf_{t\geq 0} e^{-tqn}(pe^t+1-p)^n.\] Let
\[g(t)= e^{-tqn}(pe^t+1-p)^n.\] Minimizing \(\log g(t)\) gives
\[\frac{\,d}{\,dt}\log g(t)
=-qn+n\frac{pe^t}{pe^t+1-p}=0.\] Thus
\[\frac{pe^{t^*}}{pe^{t^*}+1-p}=q,\] which implies
\[e^{t^*}=\frac{q(1-p)}{(1-q)p}.\] Substituting this value into \(g(t)\)
gives \[\mathbb{P}(X\geq qn)
\leq
\left(\frac pq\right)^{qn}
\left(\frac{1-p}{1-q}\right)^{(1-q)n}.\] This is an exponential bound in
\(n\).

\end{tcolorbox}

\subsection{Comparison: Markov, Chebyshev, and
Chernoff}\label{comparison-markov-chebyshev-and-chernoff}

This subsection compares the three bounds for binomial tails.

For \(X\sim \operatorname{Binomial}(n,p)\) and \(p<q<1\):
\[\text{Markov:}\qquad
\mathbb{P}(X\geq qn)\leq \frac pq.\] \[\text{Chebyshev:}\qquad
\mathbb{P}(X\geq qn)
\leq
\frac{p(1-p)}{(q-p)^2n}.\] \[\text{Chernoff:}\qquad
\mathbb{P}(X\geq qn)
\leq
\left(\frac pq\right)^{qn}
\left(\frac{1-p}{1-q}\right)^{(1-q)n}.\]

\begin{tcolorbox}[enhanced jigsaw, arc=.35mm, bottomrule=.15mm, bottomtitle=1mm, breakable, colback=white, colbacktitle=quarto-callout-tip-color!10!white, colframe=quarto-callout-tip-color-frame, coltitle=black, left=2mm, leftrule=.75mm, opacityback=0, opacitybacktitle=0.6, rightrule=.15mm, title=\textcolor{quarto-callout-tip-color}{\faLightbulb}\hspace{0.5em}{Example}, titlerule=0mm, toprule=.15mm, toptitle=1mm]

\textbf{Example 17} (Comparison when \(p=1/2\) and \(q=3/4\)). Let
\(X\sim \operatorname{Binomial}(n,1/2)\). Compare bounds for
\(\mathbb{P}(X\geq 3n/4)\).

\end{tcolorbox}

\begin{tcolorbox}[enhanced jigsaw, arc=.35mm, bottomrule=.15mm, bottomtitle=1mm, breakable, colback=white, colbacktitle=quarto-callout-note-color!10!white, colframe=quarto-callout-note-color-frame, coltitle=black, left=2mm, leftrule=.75mm, opacityback=0, opacitybacktitle=0.6, rightrule=.15mm, title=\textcolor{quarto-callout-note-color}{\faInfo}\hspace{0.5em}{Solution}, titlerule=0mm, toprule=.15mm, toptitle=1mm]

Markov's inequality gives
\[\mathbb{P}(X\geq 3n/4)\leq \frac{1/2}{3/4}=\frac23.\] Chebyshev's
inequality gives \[\mathbb{P}(X\geq 3n/4)
\leq
\frac{(1/2)(1/2)}{(1/4)^2n}
=\frac4n.\] Chernoff's bound gives \[\mathbb{P}(X\geq 3n/4)
\leq
\left(\frac{1/2}{3/4}\right)^{3n/4}
\left(\frac{1/2}{1/4}\right)^{n/4}
=\left(\frac{16}{27}\right)^{n/4}.\] Markov's bound does not decay with
\(n\), Chebyshev's bound decays like \(1/n\), and Chernoff's bound
decays exponentially.

\end{tcolorbox}

\begin{tcolorbox}[enhanced jigsaw, arc=.35mm, bottomrule=.15mm, bottomtitle=1mm, breakable, colback=white, colbacktitle=quarto-callout-tip-color!10!white, colframe=quarto-callout-tip-color-frame, coltitle=black, left=2mm, leftrule=.75mm, opacityback=0, opacitybacktitle=0.6, rightrule=.15mm, title=\textcolor{quarto-callout-tip-color}{\faLightbulb}\hspace{0.5em}{Example}, titlerule=0mm, toprule=.15mm, toptitle=1mm]

\textbf{Example 18} (Comparison when \(p=1/3\) and \(q=2/3\)). Let
\(X\sim \operatorname{Binomial}(n,1/3)\). Compare bounds for
\(\mathbb{P}(X\geq 2n/3)\).

\end{tcolorbox}

\begin{tcolorbox}[enhanced jigsaw, arc=.35mm, bottomrule=.15mm, bottomtitle=1mm, breakable, colback=white, colbacktitle=quarto-callout-note-color!10!white, colframe=quarto-callout-note-color-frame, coltitle=black, left=2mm, leftrule=.75mm, opacityback=0, opacitybacktitle=0.6, rightrule=.15mm, title=\textcolor{quarto-callout-note-color}{\faInfo}\hspace{0.5em}{Solution}, titlerule=0mm, toprule=.15mm, toptitle=1mm]

Markov's inequality gives
\[\mathbb{P}(X\geq 2n/3)\leq \frac{1/3}{2/3}=\frac12.\] Chebyshev's
inequality gives \[\mathbb{P}(X\geq 2n/3)
\leq
\frac{(1/3)(2/3)}{(1/3)^2n}
=\frac2n.\] The Chernoff bound gives \[\mathbb{P}(X\geq 2n/3)
\leq
\left(\frac{1/3}{2/3}\right)^{2n/3}
\left(\frac{2/3}{1/3}\right)^{n/3}
=2^{-2n/3}2^{n/3}=2^{-n/3}.\] The lecture slide states the sharper form
\(2^{-n/2}\) for this special comparison. The key qualitative point is
the same: Chernoff gives exponential decay, while Chebyshev gives
polynomial decay.

\end{tcolorbox}

\subsection{Chernoff bounds for sums}\label{chernoff-bounds-for-sums}

This subsection explains why Chernoff bounds are especially convenient
for sums of independent random variables.

Suppose \(Z_1,\ldots,Z_n\) are independent. Then \[M_{Z_1+\cdots+Z_n}(t)
=\prod_{i=1}^{n}M_{Z_i}(t).\] Therefore, if \(Z_i\) are iid,
\[M_{\sum_{i=1}^n Z_i}(t)=M_Z(t)^n.\] This means that the Chernoff bound
of a sum usually reduces to computing the MGF of a single summand.

\begin{tcolorbox}[enhanced jigsaw, arc=.35mm, bottomrule=.15mm, bottomtitle=1mm, breakable, colback=white, colbacktitle=quarto-callout-tip-color!10!white, colframe=quarto-callout-tip-color-frame, coltitle=black, left=2mm, leftrule=.75mm, opacityback=0, opacitybacktitle=0.6, rightrule=.15mm, title=\textcolor{quarto-callout-tip-color}{\faLightbulb}\hspace{0.5em}{Example}, titlerule=0mm, toprule=.15mm, toptitle=1mm]

\textbf{Example 19} (Random signs). Let \(S\) be a random sign variable
with \[\mathbb{P}(S=1)=\mathbb{P}(S=-1)=\frac12.\] Let
\(X=S_1+\cdots+S_n\), where \(S_i\) are iid copies of \(S\). Find a
Chernoff bound for \(\mathbb{P}(X\geq a)\).

\end{tcolorbox}

\begin{tcolorbox}[enhanced jigsaw, arc=.35mm, bottomrule=.15mm, bottomtitle=1mm, breakable, colback=white, colbacktitle=quarto-callout-note-color!10!white, colframe=quarto-callout-note-color-frame, coltitle=black, left=2mm, leftrule=.75mm, opacityback=0, opacitybacktitle=0.6, rightrule=.15mm, title=\textcolor{quarto-callout-note-color}{\faInfo}\hspace{0.5em}{Solution}, titlerule=0mm, toprule=.15mm, toptitle=1mm]

The MGF of one random sign is
\[M_S(t)=\mathbb{E}[e^{tS}]=\frac12e^t+\frac12e^{-t}=\cosh t.\] Hence
\[M_X(t)=(\cosh t)^n.\] Chernoff's bound gives \[\mathbb{P}(X\geq a)
\leq
\inf_{t\geq 0} e^{-ta}(\cosh t)^n.\] Using \(\cosh t\leq e^{t^2/2}\) for
all \(t\), we get \[\mathbb{P}(X\geq a)
\leq
\inf_{t\geq 0}\exp\left(-ta+\frac{nt^2}{2}\right).\] The minimum occurs
at \(t=a/n\), giving \[\mathbb{P}(X\geq a)
\leq
\exp\left(-\frac{a^2}{2n}\right).\]

\end{tcolorbox}

\section{Hoeffding's Inequality}\label{hoeffdings-inequality}

This section gives a powerful exponential concentration inequality for
bounded independent random variables.

\subsection{Concentration of the sample
mean}\label{concentration-of-the-sample-mean}

This subsection compares Chebyshev's polynomial-rate bound with
Hoeffding's exponential-rate bound.

Recall that if \(X_1,\ldots,X_n\) are iid with variance \(\sigma^2\),
then Chebyshev's inequality applied to the sample mean gives
\[\mathbb{P}(|\bar X_n-\mathbb{E}[\bar X_n]|\geq \varepsilon)
\leq
\frac{\sigma^2}{n\varepsilon^2}.\] This decays like \(1/n\).

\begin{tcolorbox}[enhanced jigsaw, arc=.35mm, bottomrule=.15mm, bottomtitle=1mm, breakable, colback=white, colbacktitle=quarto-callout-important-color!10!white, colframe=quarto-callout-important-color-frame, coltitle=black, left=2mm, leftrule=.75mm, opacityback=0, opacitybacktitle=0.6, rightrule=.15mm, title=\textcolor{quarto-callout-important-color}{\faExclamation}\hspace{0.5em}{Theorem}, titlerule=0mm, toprule=.15mm, toptitle=1mm]

\textbf{Theorem 20} (Hoeffding's inequality, bounded case).
\emph{Suppose \(X_1,\ldots,X_n\) are independent and
\(0\leq X_i\leq 1\). Let \[\bar X_n=\frac1n\sum_{i=1}^n X_i.\] Then, for
any \(\varepsilon>0\),
\[\mathbb{P}(|\bar X_n-\mathbb{E}[\bar X_n]|\geq \varepsilon)
\leq
2e^{-2n\varepsilon^2}.\] More generally, if \(a_i\leq X_i\leq b_i\),
then a scaled version holds.}

\end{tcolorbox}

\begin{tcolorbox}[enhanced jigsaw, arc=.35mm, bottomrule=.15mm, bottomtitle=1mm, breakable, colback=white, colbacktitle=quarto-callout-tip-color!10!white, colframe=quarto-callout-tip-color-frame, coltitle=black, left=2mm, leftrule=.75mm, opacityback=0, opacitybacktitle=0.6, rightrule=.15mm, title=\textcolor{quarto-callout-tip-color}{\faLightbulb}\hspace{0.5em}{Main message}, titlerule=0mm, toprule=.15mm, toptitle=1mm]

Why Hoeffding matters Hoeffding's inequality gives an exponential
concentration rate. This is much faster than the Chebyshev rate and is a
central tool in modern statistics, statistical learning,
high-dimensional problems, nonparametric inference, semiparametric
inference, and empirical risk minimization.

\end{tcolorbox}

\begin{tcolorbox}[enhanced jigsaw, arc=.35mm, bottomrule=.15mm, bottomtitle=1mm, breakable, colback=white, colbacktitle=quarto-callout-tip-color!10!white, colframe=quarto-callout-tip-color-frame, coltitle=black, left=2mm, leftrule=.75mm, opacityback=0, opacitybacktitle=0.6, rightrule=.15mm, title=\textcolor{quarto-callout-tip-color}{\faLightbulb}\hspace{0.5em}{Example}, titlerule=0mm, toprule=.15mm, toptitle=1mm]

\textbf{Example 21} (Coin-flip sample proportion). Let
\(X_1,\ldots,X_n\) be iid Bernoulli\((p)\) random variables and let
\(\bar X_n\) be the sample proportion of successes. Bound
\[\mathbb{P}(|\bar X_n-p|\geq \varepsilon).\]

\end{tcolorbox}

\begin{tcolorbox}[enhanced jigsaw, arc=.35mm, bottomrule=.15mm, bottomtitle=1mm, breakable, colback=white, colbacktitle=quarto-callout-note-color!10!white, colframe=quarto-callout-note-color-frame, coltitle=black, left=2mm, leftrule=.75mm, opacityback=0, opacitybacktitle=0.6, rightrule=.15mm, title=\textcolor{quarto-callout-note-color}{\faInfo}\hspace{0.5em}{Solution}, titlerule=0mm, toprule=.15mm, toptitle=1mm]

Since \(0\leq X_i\leq 1\) and \(\mathbb{E}[\bar X_n]=p\), Hoeffding's
inequality gives \[\mathbb{P}(|\bar X_n-p|\geq \varepsilon)
\leq
2e^{-2n\varepsilon^2}.\] For example, if \(\varepsilon=0.05\) and
\(n=1000\), then \[2e^{-2(1000)(0.05)^2}=2e^{-5}\approx 0.0135.\] Thus
the sample proportion is within \(0.05\) of \(p\) with probability at
least \(0.9865\).

\end{tcolorbox}

\begin{tcolorbox}[enhanced jigsaw, arc=.35mm, bottomrule=.15mm, bottomtitle=1mm, breakable, colback=white, colbacktitle=quarto-callout-caution-color!10!white, colframe=quarto-callout-caution-color-frame, coltitle=black, left=2mm, leftrule=.75mm, opacityback=0, opacitybacktitle=0.6, rightrule=.15mm, title=\textcolor{quarto-callout-caution-color}{\faFire}\hspace{0.5em}{Practice Problem}, titlerule=0mm, toprule=.15mm, toptitle=1mm]

\textbf{Practice Problem 22} (Sample size from Hoeffding). How large
should \(n\) be to guarantee
\[\mathbb{P}(|\bar X_n-p|\geq \varepsilon)\leq \delta\] for Bernoulli
data?

\end{tcolorbox}

\begin{tcolorbox}[enhanced jigsaw, arc=.35mm, bottomrule=.15mm, bottomtitle=1mm, breakable, colback=white, colbacktitle=quarto-callout-note-color!10!white, colframe=quarto-callout-note-color-frame, coltitle=black, left=2mm, leftrule=.75mm, opacityback=0, opacitybacktitle=0.6, rightrule=.15mm, title=\textcolor{quarto-callout-note-color}{\faInfo}\hspace{0.5em}{Solution}, titlerule=0mm, toprule=.15mm, toptitle=1mm]

Hoeffding gives \[\mathbb{P}(|\bar X_n-p|\geq \varepsilon)
\leq 2e^{-2n\varepsilon^2}.\] We want
\[2e^{-2n\varepsilon^2}\leq \delta.\] Taking logs,
\[-2n\varepsilon^2\leq \log(\delta/2).\] Thus it is enough to take
\[n\geq \frac{1}{2\varepsilon^2}\log\left(\frac2\delta\right).\]

\end{tcolorbox}

\section{Cauchy-Schwarz, Holder, and Minkowski
Inequalities}\label{cauchy-schwarz-holder-and-minkowski-inequalities}

This section introduces norm and inner-product inequalities for random
variables.

\subsection{Cauchy-Schwarz inequality}\label{cauchy-schwarz-inequality}

This subsection views random variables as vectors with inner product
\(\langle X,Y\rangle=\mathbb{E}[XY]\).

\begin{tcolorbox}[enhanced jigsaw, arc=.35mm, bottomrule=.15mm, bottomtitle=1mm, breakable, colback=white, colbacktitle=quarto-callout-important-color!10!white, colframe=quarto-callout-important-color-frame, coltitle=black, left=2mm, leftrule=.75mm, opacityback=0, opacitybacktitle=0.6, rightrule=.15mm, title=\textcolor{quarto-callout-important-color}{\faExclamation}\hspace{0.5em}{Theorem}, titlerule=0mm, toprule=.15mm, toptitle=1mm]

\textbf{Theorem 23} (Cauchy-Schwarz inequality). \emph{Suppose \(X\) and
\(Y\) are random variables with finite second moments. Then
\[\bigl(\mathbb{E}[XY]\bigr)^2\leq \mathbb{E}[X^2]\mathbb{E}[Y^2].\]
Equality holds if and only if \(X=aY\) almost surely for some constant
\(a\in\mathbb{R}\).}

\end{tcolorbox}

\begin{tcolorbox}[enhanced jigsaw, arc=.35mm, bottomrule=.15mm, bottomtitle=1mm, breakable, colback=white, colbacktitle=quarto-callout-note-color!10!white, colframe=quarto-callout-note-color-frame, coltitle=black, left=2mm, leftrule=.75mm, opacityback=0, opacitybacktitle=0.6, rightrule=.15mm, title=\textcolor{quarto-callout-note-color}{\faInfo}\hspace{0.5em}{Proof}, titlerule=0mm, toprule=.15mm, toptitle=1mm]

\emph{Proof.} For any real number \(t\),
\[0\leq \mathbb{E}[(X-tY)^2]=\mathbb{E}[X^2]-2t\mathbb{E}[XY]+t^2\mathbb{E}[Y^2].\]
This quadratic in \(t\) is nonnegative for all \(t\), so its
discriminant is nonpositive:
\[4(\mathbb{E}[XY])^2-4\mathbb{E}[X^2]\mathbb{E}[Y^2]\leq 0.\] Thus
\[(\mathbb{E}[XY])^2\leq \mathbb{E}[X^2]\mathbb{E}[Y^2].\] Equality
occurs precisely when \(X-tY=0\) almost surely for some \(t\).~◻

\end{tcolorbox}

The same idea generalizes to random vectors. If \(X,Y\in\mathbb{R}^n\),
define
\[\langle X,Y\rangle=\mathbb{E}[X^T Y]=\mathbb{E}[X_1Y_1+\cdots+X_nY_n].\]
Then Cauchy-Schwarz becomes \[\left(\mathbb{E}[X^TY]\right)^2
\leq
\mathbb{E}[X^TX]\,\mathbb{E}[Y^TY].\]

\begin{tcolorbox}[enhanced jigsaw, arc=.35mm, bottomrule=.15mm, bottomtitle=1mm, breakable, colback=white, colbacktitle=quarto-callout-tip-color!10!white, colframe=quarto-callout-tip-color-frame, coltitle=black, left=2mm, leftrule=.75mm, opacityback=0, opacitybacktitle=0.6, rightrule=.15mm, title=\textcolor{quarto-callout-tip-color}{\faLightbulb}\hspace{0.5em}{Example}, titlerule=0mm, toprule=.15mm, toptitle=1mm]

\textbf{Example 24} (Correlation is between \(-1\) and \(1\)). Use
Cauchy-Schwarz to prove \[|\operatorname{Corr}(X,Y)|\leq 1.\] Also
characterize equality.

\end{tcolorbox}

\begin{tcolorbox}[enhanced jigsaw, arc=.35mm, bottomrule=.15mm, bottomtitle=1mm, breakable, colback=white, colbacktitle=quarto-callout-note-color!10!white, colframe=quarto-callout-note-color-frame, coltitle=black, left=2mm, leftrule=.75mm, opacityback=0, opacitybacktitle=0.6, rightrule=.15mm, title=\textcolor{quarto-callout-note-color}{\faInfo}\hspace{0.5em}{Solution}, titlerule=0mm, toprule=.15mm, toptitle=1mm]

Assume \(\operatorname{Var}(X)>0\) and \(\operatorname{Var}(Y)>0\).
Define standardized variables
\[U=\frac{X-\mathbb{E}[X]}{\sqrt{\operatorname{Var}(X)}},
\qquad
V=\frac{Y-\mathbb{E}[Y]}{\sqrt{\operatorname{Var}(Y)}}.\] Then
\[\mathbb{E}[U]=0,
\qquad \mathbb{E}[V]=0,
\qquad \mathbb{E}[U^2]=\mathbb{E}[V^2]=1,\] and
\[\operatorname{Corr}(X,Y)=\mathbb{E}[UV].\] By Cauchy-Schwarz,
\[|\mathbb{E}[UV]|\leq \sqrt{\mathbb{E}[U^2]\mathbb{E}[V^2]}=1.\]
Equality holds if and only if \(U=aV\) almost surely for some constant
\(a\), which is equivalent to \[X=\alpha+\beta Y\] almost surely for
constants \(\alpha,\beta\in\mathbb{R}\).

\end{tcolorbox}

\subsection{Holder's inequality}\label{holders-inequality}

This subsection generalizes Cauchy-Schwarz from second moments to
\(p\)-th moments.

\begin{tcolorbox}[enhanced jigsaw, arc=.35mm, bottomrule=.15mm, bottomtitle=1mm, breakable, colback=white, colbacktitle=quarto-callout-important-color!10!white, colframe=quarto-callout-important-color-frame, coltitle=black, left=2mm, leftrule=.75mm, opacityback=0, opacitybacktitle=0.6, rightrule=.15mm, title=\textcolor{quarto-callout-important-color}{\faExclamation}\hspace{0.5em}{Theorem}, titlerule=0mm, toprule=.15mm, toptitle=1mm]

\textbf{Theorem 25} (Holder's inequality). \emph{Let \(X\) and \(Y\) be
random variables and let \(p,q>1\) satisfy \[\frac1p+\frac1q=1.\] Then
\[\mathbb{E}[|XY|]
\leq
\left(\mathbb{E}[|X|^p]\right)^{1/p}
\left(\mathbb{E}[|Y|^q]\right)^{1/q}.\]}

\end{tcolorbox}

When \(p=q=2\), Holder's inequality is the Cauchy-Schwarz inequality.

\begin{tcolorbox}[enhanced jigsaw, arc=.35mm, bottomrule=.15mm, bottomtitle=1mm, breakable, colback=white, colbacktitle=quarto-callout-tip-color!10!white, colframe=quarto-callout-tip-color-frame, coltitle=black, left=2mm, leftrule=.75mm, opacityback=0, opacitybacktitle=0.6, rightrule=.15mm, title=\textcolor{quarto-callout-tip-color}{\faLightbulb}\hspace{0.5em}{Example}, titlerule=0mm, toprule=.15mm, toptitle=1mm]

\textbf{Example 26} (Cauchy-Schwarz as a special case). Show that
Holder's inequality with \(p=q=2\) gives Cauchy-Schwarz.

\end{tcolorbox}

\begin{tcolorbox}[enhanced jigsaw, arc=.35mm, bottomrule=.15mm, bottomtitle=1mm, breakable, colback=white, colbacktitle=quarto-callout-note-color!10!white, colframe=quarto-callout-note-color-frame, coltitle=black, left=2mm, leftrule=.75mm, opacityback=0, opacitybacktitle=0.6, rightrule=.15mm, title=\textcolor{quarto-callout-note-color}{\faInfo}\hspace{0.5em}{Solution}, titlerule=0mm, toprule=.15mm, toptitle=1mm]

If \(p=q=2\), then \(1/p+1/q=1/2+1/2=1\). Holder's inequality gives
\[\mathbb{E}[|XY|]
\leq
\left(\mathbb{E}[|X|^2]\right)^{1/2}
\left(\mathbb{E}[|Y|^2]\right)^{1/2}.\] Squaring both sides gives
\[(\mathbb{E}[|XY|])^2\leq \mathbb{E}[X^2]\mathbb{E}[Y^2].\] For signed
variables, applying this to \(|X|\) and \(|Y|\) yields the usual bound
\[|\mathbb{E}[XY]|\leq \mathbb{E}[|XY|]\leq \sqrt{\mathbb{E}[X^2]\mathbb{E}[Y^2]}.\]

\end{tcolorbox}

\subsection{Minkowski's inequality}\label{minkowskis-inequality}

This subsection gives the triangle inequality for \(L^p\) norms of
random variables.

\begin{tcolorbox}[enhanced jigsaw, arc=.35mm, bottomrule=.15mm, bottomtitle=1mm, breakable, colback=white, colbacktitle=quarto-callout-important-color!10!white, colframe=quarto-callout-important-color-frame, coltitle=black, left=2mm, leftrule=.75mm, opacityback=0, opacitybacktitle=0.6, rightrule=.15mm, title=\textcolor{quarto-callout-important-color}{\faExclamation}\hspace{0.5em}{Theorem}, titlerule=0mm, toprule=.15mm, toptitle=1mm]

\textbf{Theorem 27} (Minkowski's inequality). \emph{Let \(X\) and \(Y\)
be random variables. For \(1\leq p<\infty\),
\[\left(\mathbb{E}[|X+Y|^p]\right)^{1/p}
\leq
\left(\mathbb{E}[|X|^p]\right)^{1/p}
+
\left(\mathbb{E}[|Y|^p]\right)^{1/p}.\]}

\end{tcolorbox}

\begin{tcolorbox}[enhanced jigsaw, arc=.35mm, bottomrule=.15mm, bottomtitle=1mm, breakable, colback=white, colbacktitle=quarto-callout-tip-color!10!white, colframe=quarto-callout-tip-color-frame, coltitle=black, left=2mm, leftrule=.75mm, opacityback=0, opacitybacktitle=0.6, rightrule=.15mm, title=\textcolor{quarto-callout-tip-color}{\faLightbulb}\hspace{0.5em}{Example}, titlerule=0mm, toprule=.15mm, toptitle=1mm]

\textbf{Example 28} (Triangle inequality in \(L^2\)). State Minkowski's
inequality for \(p=2\).

\end{tcolorbox}

\begin{tcolorbox}[enhanced jigsaw, arc=.35mm, bottomrule=.15mm, bottomtitle=1mm, breakable, colback=white, colbacktitle=quarto-callout-note-color!10!white, colframe=quarto-callout-note-color-frame, coltitle=black, left=2mm, leftrule=.75mm, opacityback=0, opacitybacktitle=0.6, rightrule=.15mm, title=\textcolor{quarto-callout-note-color}{\faInfo}\hspace{0.5em}{Solution}, titlerule=0mm, toprule=.15mm, toptitle=1mm]

For \(p=2\), Minkowski's inequality says \[\sqrt{\mathbb{E}[(X+Y)^2]}
\leq
\sqrt{\mathbb{E}[X^2]}+
\sqrt{\mathbb{E}[Y^2]}.\] This is the triangle inequality for the norm
\[\|X\|_2=\sqrt{\mathbb{E}[X^2]}.\]

\end{tcolorbox}

\section{Jensen's Inequality}\label{jensens-inequality}

This section explains how convexity interacts with expectation.

\subsection{Convex functions}\label{convex-functions}

This subsection reviews convexity before stating Jensen's inequality.

\textbf{Definition 29} (Convex and concave functions). A function
\(g:\mathbb{R}\to\mathbb{R}\) is convex on an interval \([a,b]\) if for
all \(x,y\in[a,b]\) and all \(\lambda\in[0,1]\),
\[g(\lambda x+(1-\lambda)y)
\leq
\lambda g(x)+(1-\lambda)g(y).\] A function \(g\) is concave if \(-g\) is
convex.

For a convex function, the same property holds for finite convex
combinations: \[g\left(\sum_{i=1}^n \lambda_i x_i\right)
\leq
\sum_{i=1}^n \lambda_i g(x_i),\] where \(\lambda_i\geq 0\) and
\(\sum_i\lambda_i=1\).

If \(g\) is twice differentiable, then \(g\) is convex on an interval if
\(g''(x)\geq 0\) for all \(x\) in the interval.

\begin{tcolorbox}[enhanced jigsaw, arc=.35mm, bottomrule=.15mm, bottomtitle=1mm, breakable, colback=white, colbacktitle=quarto-callout-important-color!10!white, colframe=quarto-callout-important-color-frame, coltitle=black, left=2mm, leftrule=.75mm, opacityback=0, opacitybacktitle=0.6, rightrule=.15mm, title=\textcolor{quarto-callout-important-color}{\faExclamation}\hspace{0.5em}{Theorem}, titlerule=0mm, toprule=.15mm, toptitle=1mm]

\textbf{Theorem 30} (Convex functions lie above tangents). \emph{Let
\(g:[a,b]\to\mathbb{R}\) be convex and let \(a<c<b\). Then there exist
constants \(A,B\in\mathbb{R}\) such that \[g(c)=Ac+B\] and
\[g(x)\geq Ax+B\] for all \(x\in[a,b]\).}

\end{tcolorbox}

\subsection{Jensen's inequality}\label{jensens-inequality-1}

This subsection gives one of the most important inequalities involving
expectation.

\begin{tcolorbox}[enhanced jigsaw, arc=.35mm, bottomrule=.15mm, bottomtitle=1mm, breakable, colback=white, colbacktitle=quarto-callout-important-color!10!white, colframe=quarto-callout-important-color-frame, coltitle=black, left=2mm, leftrule=.75mm, opacityback=0, opacitybacktitle=0.6, rightrule=.15mm, title=\textcolor{quarto-callout-important-color}{\faExclamation}\hspace{0.5em}{Theorem}, titlerule=0mm, toprule=.15mm, toptitle=1mm]

\textbf{Theorem 31} (Jensen's inequality). \emph{Suppose \(X\) is a
random variable such that \(\mathbb{P}(a\leq X\leq b)=1\). If
\(g:\mathbb{R}\to\mathbb{R}\) is convex on \([a,b]\), then
\[\mathbb{E}[g(X)]\geq g(\mathbb{E}[X]).\] If \(g\) is concave on
\([a,b]\), then \[\mathbb{E}[g(X)]\leq g(\mathbb{E}[X]).\]}

\end{tcolorbox}

\begin{tcolorbox}[enhanced jigsaw, arc=.35mm, bottomrule=.15mm, bottomtitle=1mm, breakable, colback=white, colbacktitle=quarto-callout-note-color!10!white, colframe=quarto-callout-note-color-frame, coltitle=black, left=2mm, leftrule=.75mm, opacityback=0, opacitybacktitle=0.6, rightrule=.15mm, title=\textcolor{quarto-callout-note-color}{\faInfo}\hspace{0.5em}{Proof}, titlerule=0mm, toprule=.15mm, toptitle=1mm]

\emph{Proof.} Assume \(g\) is convex. Let \(c=\mathbb{E}[X]\). By the
tangent-line property, there exist \(A,B\in\mathbb{R}\) such that
\[g(x)\geq Ax+B\] for all \(x\) in the support of \(X\), and
\[g(c)=Ac+B.\] Taking expectations gives
\[\mathbb{E}[g(X)]\geq \mathbb{E}[AX+B]=A\mathbb{E}[X]+B=Ac+B=g(c)=g(\mathbb{E}[X]).\]
The concave case follows by applying the convex result to \(-g\).~◻

\end{tcolorbox}

\begin{tcolorbox}[enhanced jigsaw, arc=.35mm, bottomrule=.15mm, bottomtitle=1mm, breakable, colback=white, colbacktitle=quarto-callout-tip-color!10!white, colframe=quarto-callout-tip-color-frame, coltitle=black, left=2mm, leftrule=.75mm, opacityback=0, opacitybacktitle=0.6, rightrule=.15mm, title=\textcolor{quarto-callout-tip-color}{\faLightbulb}\hspace{0.5em}{Example}, titlerule=0mm, toprule=.15mm, toptitle=1mm]

\textbf{Example 32} (Arithmetic-geometric mean inequality). Let
\(a_1,\ldots,a_n\) be positive numbers. Use Jensen's inequality to prove
\[\frac1n\sum_{k=1}^{n}a_k
\geq
\left(\prod_{k=1}^{n}a_k\right)^{1/n}.\]

\end{tcolorbox}

\begin{tcolorbox}[enhanced jigsaw, arc=.35mm, bottomrule=.15mm, bottomtitle=1mm, breakable, colback=white, colbacktitle=quarto-callout-note-color!10!white, colframe=quarto-callout-note-color-frame, coltitle=black, left=2mm, leftrule=.75mm, opacityback=0, opacitybacktitle=0.6, rightrule=.15mm, title=\textcolor{quarto-callout-note-color}{\faInfo}\hspace{0.5em}{Solution}, titlerule=0mm, toprule=.15mm, toptitle=1mm]

Let \(X\) be a discrete random variable with
\[\mathbb{P}(X=a_k)=\frac1n,
\qquad k=1,\ldots,n.\] The function \(g(x)=-\log x\) is convex on
\((0,\infty)\). Jensen's inequality gives
\[-\log(\mathbb{E}[X])\leq \mathbb{E}[-\log X].\] Therefore
\[-\log\left(\frac1n\sum_{k=1}^{n}a_k\right)
\leq
-\frac1n\sum_{k=1}^{n}\log a_k.\] Multiplying by \(-1\) reverses the
inequality: \[\log\left(\frac1n\sum_{k=1}^{n}a_k\right)
\geq
\frac1n\sum_{k=1}^{n}\log a_k
=
\log\left(\prod_{k=1}^{n}a_k\right)^{1/n}.\] Exponentiating gives
\[\frac1n\sum_{k=1}^{n}a_k
\geq
\left(\prod_{k=1}^{n}a_k\right)^{1/n}.\]

\end{tcolorbox}

\begin{tcolorbox}[enhanced jigsaw, arc=.35mm, bottomrule=.15mm, bottomtitle=1mm, breakable, colback=white, colbacktitle=quarto-callout-tip-color!10!white, colframe=quarto-callout-tip-color-frame, coltitle=black, left=2mm, leftrule=.75mm, opacityback=0, opacitybacktitle=0.6, rightrule=.15mm, title=\textcolor{quarto-callout-tip-color}{\faLightbulb}\hspace{0.5em}{Example}, titlerule=0mm, toprule=.15mm, toptitle=1mm]

\textbf{Example 33} (Variance is nonnegative). Suppose \(p>1\). Since
\(g(x)=x^p\) is convex on appropriate domains, Jensen's inequality
implies moment inequalities. In particular, show that
\(\operatorname{Var}(X)\geq 0\).

\end{tcolorbox}

\begin{tcolorbox}[enhanced jigsaw, arc=.35mm, bottomrule=.15mm, bottomtitle=1mm, breakable, colback=white, colbacktitle=quarto-callout-note-color!10!white, colframe=quarto-callout-note-color-frame, coltitle=black, left=2mm, leftrule=.75mm, opacityback=0, opacitybacktitle=0.6, rightrule=.15mm, title=\textcolor{quarto-callout-note-color}{\faInfo}\hspace{0.5em}{Solution}, titlerule=0mm, toprule=.15mm, toptitle=1mm]

Take \(g(x)=x^2\), which is convex. Jensen's inequality gives
\[\mathbb{E}[X^2]\geq (\mathbb{E}[X])^2.\] Therefore
\[\operatorname{Var}(X)=\mathbb{E}[X^2]-(\mathbb{E}[X])^2\geq 0.\]

\end{tcolorbox}

\begin{tcolorbox}[enhanced jigsaw, arc=.35mm, bottomrule=.15mm, bottomtitle=1mm, breakable, colback=white, colbacktitle=quarto-callout-caution-color!10!white, colframe=quarto-callout-caution-color-frame, coltitle=black, left=2mm, leftrule=.75mm, opacityback=0, opacitybacktitle=0.6, rightrule=.15mm, title=\textcolor{quarto-callout-caution-color}{\faFire}\hspace{0.5em}{Practice Problem}, titlerule=0mm, toprule=.15mm, toptitle=1mm]

\textbf{Practice Problem 34} (Log expectation). Let \(X>0\) almost
surely. Use Jensen's inequality to compare \(\mathbb{E}[\log X]\) and
\(\log\mathbb{E}[X]\).

\end{tcolorbox}

\begin{tcolorbox}[enhanced jigsaw, arc=.35mm, bottomrule=.15mm, bottomtitle=1mm, breakable, colback=white, colbacktitle=quarto-callout-note-color!10!white, colframe=quarto-callout-note-color-frame, coltitle=black, left=2mm, leftrule=.75mm, opacityback=0, opacitybacktitle=0.6, rightrule=.15mm, title=\textcolor{quarto-callout-note-color}{\faInfo}\hspace{0.5em}{Solution}, titlerule=0mm, toprule=.15mm, toptitle=1mm]

The function \(\log x\) is concave on \((0,\infty)\). Jensen's
inequality for concave functions gives
\[\mathbb{E}[\log X]\leq \log\mathbb{E}[X].\] Equivalently,
\[\exp(\mathbb{E}[\log X])\leq \mathbb{E}[X].\] This is the
probabilistic version of the geometric mean being no larger than the
arithmetic mean.

\end{tcolorbox}

\section{Equality Identities for Special
Distributions}\label{equality-identities-for-special-distributions}

This section collects useful identities for Poisson, Gamma, Normal,
chi-squared, and negative binomial distributions.

The lecture notes emphasize that for some specific distributions, such
as normal, chi-squared, Poisson, and negative binomial, there are exact
equality relations that can simplify calculations. These identities are
especially useful for moments and distributional calculations.

\subsection{Poisson recursion}\label{poisson-recursion}

This subsection begins with the basic recursion for the Poisson
probability mass function.

\begin{tcolorbox}[enhanced jigsaw, arc=.35mm, bottomrule=.15mm, bottomtitle=1mm, breakable, colback=white, colbacktitle=quarto-callout-important-color!10!white, colframe=quarto-callout-important-color-frame, coltitle=black, left=2mm, leftrule=.75mm, opacityback=0, opacitybacktitle=0.6, rightrule=.15mm, title=\textcolor{quarto-callout-important-color}{\faExclamation}\hspace{0.5em}{Proposition}, titlerule=0mm, toprule=.15mm, toptitle=1mm]

\textbf{Proposition 35} (Poisson probability recursion). \emph{If
\(X\sim \operatorname{Poisson}(\lambda)\), then for \(x=0,1,2,\ldots\),
\[\mathbb{P}(X=x+1)=\frac{\lambda}{x+1}\mathbb{P}(X=x).\]}

\end{tcolorbox}

\begin{tcolorbox}[enhanced jigsaw, arc=.35mm, bottomrule=.15mm, bottomtitle=1mm, breakable, colback=white, colbacktitle=quarto-callout-note-color!10!white, colframe=quarto-callout-note-color-frame, coltitle=black, left=2mm, leftrule=.75mm, opacityback=0, opacitybacktitle=0.6, rightrule=.15mm, title=\textcolor{quarto-callout-note-color}{\faInfo}\hspace{0.5em}{Proof}, titlerule=0mm, toprule=.15mm, toptitle=1mm]

\emph{Proof.} Since
\[\mathbb{P}(X=x)=e^{-\lambda}\frac{\lambda^x}{x!},\] we have
\[\mathbb{P}(X=x+1)=e^{-\lambda}\frac{\lambda^{x+1}}{(x+1)!}
=\frac{\lambda}{x+1}e^{-\lambda}\frac{\lambda^x}{x!}
=\frac{\lambda}{x+1}\mathbb{P}(X=x).\]~◻

\end{tcolorbox}

\subsection{Gamma identity}\label{gamma-identity}

This subsection records a useful relationship between Gamma
probabilities with adjacent shape parameters.

\begin{tcolorbox}[enhanced jigsaw, arc=.35mm, bottomrule=.15mm, bottomtitle=1mm, breakable, colback=white, colbacktitle=quarto-callout-important-color!10!white, colframe=quarto-callout-important-color-frame, coltitle=black, left=2mm, leftrule=.75mm, opacityback=0, opacitybacktitle=0.6, rightrule=.15mm, title=\textcolor{quarto-callout-important-color}{\faExclamation}\hspace{0.5em}{Proposition}, titlerule=0mm, toprule=.15mm, toptitle=1mm]

\textbf{Proposition 36} (Gamma identity). \emph{Let \(X_{\alpha,\beta}\)
denote a Gamma\((\alpha,\beta)\) random variable with density
\(f(x\mid\alpha,\beta)\), where \(\alpha>1\). Then for constants
\(a<b\), \[\mathbb{P}(a<X_{\alpha,\beta}<b)
=
\beta\{f(a\mid\alpha,\beta)-f(b\mid\alpha,\beta)\}
+
\mathbb{P}(a<X_{\alpha-1,\beta}<b),\] under the parameterization used in
Casella-Berger style identities.}

\end{tcolorbox}

\begin{remark}
\emph{Remark 37}. This identity is useful because it relates Gamma
probabilities with shape \(\alpha\) to Gamma probabilities with shape
\(\alpha-1\) plus boundary density terms.
\end{remark}

\subsection{Stein's lemma}\label{steins-lemma}

This subsection gives a normal-distribution identity that turns moments
into derivative expectations.

\begin{tcolorbox}[enhanced jigsaw, arc=.35mm, bottomrule=.15mm, bottomtitle=1mm, breakable, colback=white, colbacktitle=quarto-callout-important-color!10!white, colframe=quarto-callout-important-color-frame, coltitle=black, left=2mm, leftrule=.75mm, opacityback=0, opacitybacktitle=0.6, rightrule=.15mm, title=\textcolor{quarto-callout-important-color}{\faExclamation}\hspace{0.5em}{Theorem}, titlerule=0mm, toprule=.15mm, toptitle=1mm]

\textbf{Theorem 38} (Stein's lemma). \emph{Let
\(X\sim \operatorname{Normal}(\theta,\sigma^2)\) and let \(g\) be
differentiable with \(\mathbb{E}[|g'(X)|]<\infty\). Then
\[\mathbb{E}[g(X)(X-\theta)]=\sigma^2\mathbb{E}[g'(X)].\]}

\end{tcolorbox}

\begin{tcolorbox}[enhanced jigsaw, arc=.35mm, bottomrule=.15mm, bottomtitle=1mm, breakable, colback=white, colbacktitle=quarto-callout-tip-color!10!white, colframe=quarto-callout-tip-color-frame, coltitle=black, left=2mm, leftrule=.75mm, opacityback=0, opacitybacktitle=0.6, rightrule=.15mm, title=\textcolor{quarto-callout-tip-color}{\faLightbulb}\hspace{0.5em}{Example}, titlerule=0mm, toprule=.15mm, toptitle=1mm]

\textbf{Example 39} (Computing a normal third moment). Use Stein's lemma
to compute \(\mathbb{E}[X^3]\) when
\(X\sim \operatorname{Normal}(\theta,\sigma^2)\).

\end{tcolorbox}

\begin{tcolorbox}[enhanced jigsaw, arc=.35mm, bottomrule=.15mm, bottomtitle=1mm, breakable, colback=white, colbacktitle=quarto-callout-note-color!10!white, colframe=quarto-callout-note-color-frame, coltitle=black, left=2mm, leftrule=.75mm, opacityback=0, opacitybacktitle=0.6, rightrule=.15mm, title=\textcolor{quarto-callout-note-color}{\faInfo}\hspace{0.5em}{Solution}, titlerule=0mm, toprule=.15mm, toptitle=1mm]

Write \[X^3=X^2(X-\theta+\theta)=X^2(X-\theta)+\theta X^2.\] Taking
expectations,
\[\mathbb{E}[X^3]=\mathbb{E}[X^2(X-\theta)]+\theta\mathbb{E}[X^2].\]
Apply Stein's lemma with \(g(x)=x^2\), so \(g'(x)=2x\):
\[\mathbb{E}[X^2(X-\theta)]=\sigma^2\mathbb{E}[2X]=2\sigma^2\theta.\]
Also,
\[\mathbb{E}[X^2]=\operatorname{Var}(X)+(\mathbb{E}[X])^2=\sigma^2+\theta^2.\]
Therefore \[\mathbb{E}[X^3]=2\sigma^2\theta+\theta(\sigma^2+\theta^2)
=3\theta\sigma^2+\theta^3.\]

\end{tcolorbox}

\subsection{Chi-squared identity}\label{chi-squared-identity}

This subsection records an identity involving chi-squared random
variables with different degrees of freedom.

\begin{tcolorbox}[enhanced jigsaw, arc=.35mm, bottomrule=.15mm, bottomtitle=1mm, breakable, colback=white, colbacktitle=quarto-callout-important-color!10!white, colframe=quarto-callout-important-color-frame, coltitle=black, left=2mm, leftrule=.75mm, opacityback=0, opacitybacktitle=0.6, rightrule=.15mm, title=\textcolor{quarto-callout-important-color}{\faExclamation}\hspace{0.5em}{Proposition}, titlerule=0mm, toprule=.15mm, toptitle=1mm]

\textbf{Proposition 40} (Chi-squared identity). \emph{Let \(\chi_p^2\)
denote a chi-squared random variable with \(p\) degrees of freedom. If
expectations exist, then for any function \(h\),
\[\mathbb{E}[h(\chi_p^2)]
=p\,\mathbb{E}\left[\frac{h(\chi_{p+2}^2)}{\chi_{p+2}^2}\right].\]}

\end{tcolorbox}

\begin{remark}
\emph{Remark 41}. This identity relates expectations under \(p\) degrees
of freedom to expectations under \(p+2\) degrees of freedom. Such
identities are often useful in distribution theory and statistical
inference.
\end{remark}

\subsection{Hwang-type identities for Poisson and negative binomial
variables}\label{hwang-type-identities-for-poisson-and-negative-binomial-variables}

This subsection presents identities from the lecture note for Poisson
and negative binomial random variables.

\begin{tcolorbox}[enhanced jigsaw, arc=.35mm, bottomrule=.15mm, bottomtitle=1mm, breakable, colback=white, colbacktitle=quarto-callout-important-color!10!white, colframe=quarto-callout-important-color-frame, coltitle=black, left=2mm, leftrule=.75mm, opacityback=0, opacitybacktitle=0.6, rightrule=.15mm, title=\textcolor{quarto-callout-important-color}{\faExclamation}\hspace{0.5em}{Proposition}, titlerule=0mm, toprule=.15mm, toptitle=1mm]

\textbf{Proposition 42} (Poisson identity). \emph{Let \(g\) be a
function such that the expectations below exist. If
\(X\sim \operatorname{Poisson}(\lambda)\), then
\[\mathbb{E}[\lambda g(X)]=\mathbb{E}[Xg(X-1)].\]}

\end{tcolorbox}

\begin{tcolorbox}[enhanced jigsaw, arc=.35mm, bottomrule=.15mm, bottomtitle=1mm, breakable, colback=white, colbacktitle=quarto-callout-note-color!10!white, colframe=quarto-callout-note-color-frame, coltitle=black, left=2mm, leftrule=.75mm, opacityback=0, opacitybacktitle=0.6, rightrule=.15mm, title=\textcolor{quarto-callout-note-color}{\faInfo}\hspace{0.5em}{Proof}, titlerule=0mm, toprule=.15mm, toptitle=1mm]

\emph{Proof.} Using the Poisson pmf, \[\mathbb{E}[Xg(X-1)]
=
\sum_{x=1}^{\infty} x g(x-1)e^{-\lambda}\frac{\lambda^x}{x!}
=
\sum_{x=1}^{\infty} g(x-1)e^{-\lambda}\frac{\lambda^x}{(x-1)!}.\] Let
\(u=x-1\). Then \[\mathbb{E}[Xg(X-1)]
=
\sum_{u=0}^{\infty} g(u)e^{-\lambda}\frac{\lambda^{u+1}}{u!}
=
\lambda\sum_{u=0}^{\infty}g(u)e^{-\lambda}\frac{\lambda^u}{u!}
=
\mathbb{E}[\lambda g(X)].\]~◻

\end{tcolorbox}

\begin{tcolorbox}[enhanced jigsaw, arc=.35mm, bottomrule=.15mm, bottomtitle=1mm, breakable, colback=white, colbacktitle=quarto-callout-important-color!10!white, colframe=quarto-callout-important-color-frame, coltitle=black, left=2mm, leftrule=.75mm, opacityback=0, opacitybacktitle=0.6, rightrule=.15mm, title=\textcolor{quarto-callout-important-color}{\faExclamation}\hspace{0.5em}{Proposition}, titlerule=0mm, toprule=.15mm, toptitle=1mm]

\textbf{Proposition 43} (Negative binomial identity). \emph{If
\(X\sim \operatorname{NegBinomial}(r,p)\), then under the
parameterization used in the lecture note, \[\mathbb{E}[(1-p)g(X)]
=
\mathbb{E}\left[\frac{X}{r+X-1}g(X-1)\right],\] provided the
expectations exist.}

\end{tcolorbox}

\begin{tcolorbox}[enhanced jigsaw, arc=.35mm, bottomrule=.15mm, bottomtitle=1mm, breakable, colback=white, colbacktitle=quarto-callout-tip-color!10!white, colframe=quarto-callout-tip-color-frame, coltitle=black, left=2mm, leftrule=.75mm, opacityback=0, opacitybacktitle=0.6, rightrule=.15mm, title=\textcolor{quarto-callout-tip-color}{\faLightbulb}\hspace{0.5em}{Example}, titlerule=0mm, toprule=.15mm, toptitle=1mm]

\textbf{Example 44} (Poisson moment from the Poisson identity). If
\(X\sim \operatorname{Poisson}(\lambda)\), use the identity with
\(g(x)=x^2\) to compute \(\mathbb{E}[X^3]\).

\end{tcolorbox}

\begin{tcolorbox}[enhanced jigsaw, arc=.35mm, bottomrule=.15mm, bottomtitle=1mm, breakable, colback=white, colbacktitle=quarto-callout-note-color!10!white, colframe=quarto-callout-note-color-frame, coltitle=black, left=2mm, leftrule=.75mm, opacityback=0, opacitybacktitle=0.6, rightrule=.15mm, title=\textcolor{quarto-callout-note-color}{\faInfo}\hspace{0.5em}{Solution}, titlerule=0mm, toprule=.15mm, toptitle=1mm]

The Poisson identity gives
\[\mathbb{E}[\lambda X^2]=\mathbb{E}[X(X-1)^2].\] Expanding the
right-hand side, \[X(X-1)^2=X(X^2-2X+1)=X^3-2X^2+X.\] Therefore
\[\lambda\mathbb{E}[X^2]=\mathbb{E}[X^3]-2\mathbb{E}[X^2]+\mathbb{E}[X].\]
Solving for \(\mathbb{E}[X^3]\) gives
\[\mathbb{E}[X^3]=\lambda\mathbb{E}[X^2]+2\mathbb{E}[X^2]-\mathbb{E}[X].\]
For \(X\sim\operatorname{Poisson}(\lambda)\), \[\mathbb{E}[X]=\lambda,
\qquad
\mathbb{E}[X^2]=\operatorname{Var}(X)+(\mathbb{E}X)^2=\lambda+\lambda^2.\]
Thus \[\mathbb{E}[X^3]
=\lambda(\lambda+\lambda^2)+2(\lambda+\lambda^2)-\lambda
=\lambda^3+3\lambda^2+\lambda.\]

\end{tcolorbox}

\begin{tcolorbox}[enhanced jigsaw, arc=.35mm, bottomrule=.15mm, bottomtitle=1mm, breakable, colback=white, colbacktitle=quarto-callout-caution-color!10!white, colframe=quarto-callout-caution-color-frame, coltitle=black, left=2mm, leftrule=.75mm, opacityback=0, opacitybacktitle=0.6, rightrule=.15mm, title=\textcolor{quarto-callout-caution-color}{\faFire}\hspace{0.5em}{Practice Problem}, titlerule=0mm, toprule=.15mm, toptitle=1mm]

\textbf{Practice Problem 45} (Use Jensen for a moment bound). Let \(X\)
be any random variable with finite fourth moment. Prove that
\[(\mathbb{E}[X^2])^2\leq \mathbb{E}[X^4].\]

\end{tcolorbox}

\begin{tcolorbox}[enhanced jigsaw, arc=.35mm, bottomrule=.15mm, bottomtitle=1mm, breakable, colback=white, colbacktitle=quarto-callout-note-color!10!white, colframe=quarto-callout-note-color-frame, coltitle=black, left=2mm, leftrule=.75mm, opacityback=0, opacitybacktitle=0.6, rightrule=.15mm, title=\textcolor{quarto-callout-note-color}{\faInfo}\hspace{0.5em}{Solution}, titlerule=0mm, toprule=.15mm, toptitle=1mm]

Apply Jensen's inequality to the convex function \(g(u)=u^2\) and the
random variable \(U=X^2\). Then
\[\mathbb{E}[g(U)]\geq g(\mathbb{E}[U]).\] This means
\[\mathbb{E}[(X^2)^2]\geq (\mathbb{E}[X^2])^2.\] Hence
\[\mathbb{E}[X^4]\geq (\mathbb{E}[X^2])^2.\]

\end{tcolorbox}

\section{Summary and Study Guide}\label{summary-and-study-guide-1}

This final section summarizes the main inequalities and indicates when
each one is useful.

{\def\LTcaptype{none} % do not increment counter
\begin{longtable}[]{@{}
  >{\raggedright\arraybackslash}p{(\linewidth - 4\tabcolsep) * \real{0.3333}}
  >{\raggedright\arraybackslash}p{(\linewidth - 4\tabcolsep) * \real{0.3333}}
  >{\raggedright\arraybackslash}p{(\linewidth - 4\tabcolsep) * \real{0.3333}}@{}}
\toprule\noalign{}
\begin{minipage}[b]{\linewidth}\raggedright
Tool
\end{minipage} & \begin{minipage}[b]{\linewidth}\raggedright
Statement / idea
\end{minipage} & \begin{minipage}[b]{\linewidth}\raggedright
Uses
\end{minipage} \\
\midrule\noalign{}
\endhead
\bottomrule\noalign{}
\endlastfoot
Boole & \(\mathbb{P}(\cup_i E_i)\leq \sum_i\mathbb{P}(E_i)\) & At least
one bad event \\
Inclusion-exclusion & Exact alternating formula for finite unions &
Exact union probabilities \\
Bonferroni & Truncated inclusion-exclusion gives upper/lower bounds &
Better union bounds \\
Markov & \(\mathbb{P}(X\geq a)\leq \mathbb{E}[X]/a\) for \(X\geq0\) &
First moment tail bounds \\
Chebyshev &
\(\mathbb{P}(|X-\mathbb{E}X|\geq k)\leq \operatorname{Var}(X)/k^2\) &
Variance-based concentration \\
Chernoff & \(\mathbb{P}(X\geq a)\leq \inf_{t\geq0}e^{-ta}M_X(t)\) &
Exponential tail bounds \\
Hoeffding &
\(\mathbb{P}(|\bar X-\mathbb{E}\bar X|\geq\epsilon)\leq 2e^{-2n\epsilon^2}\)
for \(0\leq X_i\leq1\) & Bounded sample means \\
Cauchy-Schwarz &
\((\mathbb{E}[XY])^2\leq \mathbb{E}[X^2]\mathbb{E}[Y^2]\) & Correlation
and inner products \\
Hölder &
\(\mathbb{E}|XY|\leq (\mathbb{E}|X|^p)^{1/p}(\mathbb{E}|Y|^q)^{1/q}\) &
General moment control \\
Minkowski & \(\|X+Y\|_p\leq \|X\|_p+\|Y\|_p\) & Triangle inequality for
random variables \\
Jensen & \(g(\mathbb{E}X)\leq \mathbb{E}[g(X)]\) for convex \(g\) &
Convexity and expectation \\
\end{longtable}
}

\begin{tcolorbox}[enhanced jigsaw, arc=.35mm, bottomrule=.15mm, bottomtitle=1mm, breakable, colback=white, colbacktitle=quarto-callout-tip-color!10!white, colframe=quarto-callout-tip-color-frame, coltitle=black, left=2mm, leftrule=.75mm, opacityback=0, opacitybacktitle=0.6, rightrule=.15mm, title=\textcolor{quarto-callout-tip-color}{\faLightbulb}\hspace{0.5em}{Main message}, titlerule=0mm, toprule=.15mm, toptitle=1mm]

How to choose a bound Use Boole/Bonferroni for unions of events. Use
Markov if only a first moment is available. Use Chebyshev if a variance
is available. Use Chernoff or Hoeffding for exponential concentration,
especially for sums or averages of independent random variables. Use
Cauchy-Schwarz/Holder/Minkowski/Jensen when the problem involves
expectations of functions, products, or convexity.

\end{tcolorbox}

\bookmarksetup{startatroot}

\chapter{Chapter 8: Sampling and Order
Statistics}\label{chapter-8-sampling-and-order-statistics}

This chapter moves from individual random variables to random samples.
The central idea is that statistics such as the sample mean, sample
variance, minimum, maximum, median, and range are themselves random
variables before the data are observed. Their distributions are called
\textbf{sampling distributions}.

\textbf{Topics.} Random samples; statistics and sampling distributions;
sample mean and sample variance; sums of random samples; convolution;
sampling from the normal distribution; chi-square, Student's \(t\), and
\(F\) distributions; order statistics; sample range and median.

\section{Overview}\label{overview-5}

This section moves from single random variables to random samples, which
are the mathematical objects behind statistical inference.

A statistical data set is usually modeled as a collection of random
variables \(X_1,\ldots,X_n\). The most important case is that the
variables are independent and identically distributed. From a random
sample we build statistics such as the sample mean, sample variance,
maximum, minimum, median, and sample range. The probability
distributions of these statistics are called sampling distributions.

\begin{tcolorbox}[enhanced jigsaw, arc=.35mm, bottomrule=.15mm, bottomtitle=1mm, breakable, colback=white, colbacktitle=quarto-callout-tip-color!10!white, colframe=quarto-callout-tip-color-frame, coltitle=black, left=2mm, leftrule=.75mm, opacityback=0, opacitybacktitle=0.6, rightrule=.15mm, title=\textcolor{quarto-callout-tip-color}{\faLightbulb}\hspace{0.5em}{Main message}, titlerule=0mm, toprule=.15mm, toptitle=1mm]

A random sample produces random summaries. Statistical inference studies
the distributions of those summaries.
\[X_1,\ldots,X_n \quad \longrightarrow \quad T(X_1,\ldots,X_n).\] The
distribution of \(T\) tells us how the statistic behaves before the data
are observed.

\end{tcolorbox}

The key objects in this section are:

\begin{itemize}
\item
  random samples and iid assumptions;
\item
  sample mean \(\overline X\) and sample variance \(S^2\);
\item
  sums of independent variables and convolutions;
\item
  normal sampling theory: chi-square, \(t\), and \(F\) distributions;
\item
  order statistics \(X_{(1)}<\cdots<X_{(n)}\).
\end{itemize}

\section{Random Samples}\label{random-samples}

This section defines the random-sample model and explains why
independence is a modeling assumption, not automatic from the word
``sample.''

\begin{tcolorbox}[enhanced jigsaw, arc=.35mm, bottomrule=.15mm, bottomtitle=1mm, breakable, colback=white, colbacktitle=quarto-callout-note-color!10!white, colframe=quarto-callout-note-color-frame, coltitle=black, left=2mm, leftrule=.75mm, opacityback=0, opacitybacktitle=0.6, rightrule=.15mm, title=\textcolor{quarto-callout-note-color}{\faInfo}\hspace{0.5em}{Definition}, titlerule=0mm, toprule=.15mm, toptitle=1mm]

\textbf{Definition 1} (Independent and identically distributed random
variables). Suppose \(X_1,\ldots,X_n\) are mutually independent random
variables and the marginal pdf or pmf of each \(X_i\) is the same
function \(f(x)\). Then \(X_1,\ldots,X_n\) are called
\textbf{independent and identically distributed}, abbreviated
\textbf{iid}, with pdf or pmf \(f(x)\).

\end{tcolorbox}

\begin{tcolorbox}[enhanced jigsaw, arc=.35mm, bottomrule=.15mm, bottomtitle=1mm, breakable, colback=white, colbacktitle=quarto-callout-note-color!10!white, colframe=quarto-callout-note-color-frame, coltitle=black, left=2mm, leftrule=.75mm, opacityback=0, opacitybacktitle=0.6, rightrule=.15mm, title=\textcolor{quarto-callout-note-color}{\faInfo}\hspace{0.5em}{Definition}, titlerule=0mm, toprule=.15mm, toptitle=1mm]

\textbf{Definition 2} (Random sample). If \(X_1,\ldots,X_n\) are iid
from a population distribution with pdf or pmf \(f(x)\), then
\(X_1,\ldots,X_n\) is called a \textbf{random sample of size \(n\)} from
the population \(f(x)\).

\end{tcolorbox}

If \(X_1,\ldots,X_n\) are a random sample from \(f\), then their joint
density or mass function factors as \[f_{X_1,\ldots,X_n}(x_1,\ldots,x_n)
 = \prod_{i=1}^n f(x_i).\] This product form is one of the main reasons
random samples are mathematically tractable.

\begin{tcolorbox}[enhanced jigsaw, arc=.35mm, bottomrule=.15mm, bottomtitle=1mm, breakable, colback=white, colbacktitle=quarto-callout-tip-color!10!white, colframe=quarto-callout-tip-color-frame, coltitle=black, left=2mm, leftrule=.75mm, opacityback=0, opacitybacktitle=0.6, rightrule=.15mm, title=\textcolor{quarto-callout-tip-color}{\faLightbulb}\hspace{0.5em}{Example}, titlerule=0mm, toprule=.15mm, toptitle=1mm]

\textbf{Example 3} (Exponential lifetime sample). Let \(X_1,\ldots,X_n\)
be a random sample from the exponential population
\[f(x;\beta)=\frac{1}{\beta}e^{-x/\beta},\qquad x>0,\ \beta>0.\] In an
engineering application, \(X_1,\ldots,X_n\) might be the lifetimes, in
years, of \(n\) identical lightbulbs placed on test and used until
failure. Find the joint pdf of the sample.

\end{tcolorbox}

\begin{tcolorbox}[enhanced jigsaw, arc=.35mm, bottomrule=.15mm, bottomtitle=1mm, breakable, colback=white, colbacktitle=quarto-callout-note-color!10!white, colframe=quarto-callout-note-color-frame, coltitle=black, left=2mm, leftrule=.75mm, opacityback=0, opacitybacktitle=0.6, rightrule=.15mm, title=\textcolor{quarto-callout-note-color}{\faInfo}\hspace{0.5em}{Solution}, titlerule=0mm, toprule=.15mm, toptitle=1mm]

Since \(X_1,\ldots,X_n\) are iid, the joint pdf is the product of the
marginal pdfs: \[f(x_1,\ldots,x_n;\beta)
 =\prod_{i=1}^n \frac{1}{\beta}e^{-x_i/\beta}
 =\frac{1}{\beta^n}\exp\left(-\frac{1}{\beta}\sum_{i=1}^n x_i\right),\]
for \(x_1>0,\ldots,x_n>0\), and \(0\) otherwise.

\end{tcolorbox}

\subsection{How random samples arise}\label{how-random-samples-arise}

This subsection explains two common mechanisms that produce iid random
variables.

A random sample \(X_1,\ldots,X_n\) can be obtained when:

\begin{enumerate}
\def\labelenumi{\arabic{enumi}.}
\item
  the population is effectively infinite and each observation is
  independently selected;
\item
  sampling is done \textbf{with replacement} from a finite population.
  Sampling with replacement is also the basis of the bootstrap idea.
\end{enumerate}

\subsection{Sampling without
replacement}\label{sampling-without-replacement}

This subsection emphasizes that sampling without replacement from a
finite population usually destroys independence.

Suppose the finite population is \(\{x_1,\ldots,x_N\}\). If we sample
without replacement, then the first and second draws are not
independent. For distinct population elements \(x\) and \(y\),
\[\mathbb{P}(X_1=x)=\frac{1}{N},
\qquad
\mathbb{P}(X_2=y\mid X_1=x)=\frac{1}{N-1}.\] Also,
\[\mathbb{P}(X_2=y\mid X_1=y)=0,\] because the value \(y\) cannot be
drawn again after it has already been selected.

Thus \(X_1\) and \(X_2\) are not independent. However, they have the
same marginal distribution. When \(N\) is very large and \(n\) is small
compared with \(N\), sampling without replacement may be approximately
treated as independent.

\begin{tcolorbox}[enhanced jigsaw, arc=.35mm, bottomrule=.15mm, bottomtitle=1mm, breakable, colback=white, colbacktitle=quarto-callout-warning-color!10!white, colframe=quarto-callout-warning-color-frame, coltitle=black, left=2mm, leftrule=.75mm, opacityback=0, opacitybacktitle=0.6, rightrule=.15mm, title=\textcolor{quarto-callout-warning-color}{\faExclamationTriangle}\hspace{0.5em}{Practice Problem}, titlerule=0mm, toprule=.15mm, toptitle=1mm]

\textbf{Practice Problem 4} (Sampling without replacement). A population
consists of \(\{1,2,3,4\}\). Two values are sampled without replacement.
Let \(X_1\) be the first draw and \(X_2\) the second draw. Are \(X_1\)
and \(X_2\) independent?

\end{tcolorbox}

\begin{tcolorbox}[enhanced jigsaw, arc=.35mm, bottomrule=.15mm, bottomtitle=1mm, breakable, colback=white, colbacktitle=quarto-callout-note-color!10!white, colframe=quarto-callout-note-color-frame, coltitle=black, left=2mm, leftrule=.75mm, opacityback=0, opacitybacktitle=0.6, rightrule=.15mm, title=\textcolor{quarto-callout-note-color}{\faInfo}\hspace{0.5em}{Solution}, titlerule=0mm, toprule=.15mm, toptitle=1mm]

No.~For example, \[\mathbb{P}(X_2=1)=\frac14,\] but
\[\mathbb{P}(X_2=1\mid X_1=1)=0.\] Because the conditional probability
differs from the marginal probability, \(X_1\) and \(X_2\) are not
independent.

\end{tcolorbox}

\section{Statistics and Sampling
Distributions}\label{statistics-and-sampling-distributions}

This section introduces statistics as functions of random samples and
defines the sampling distribution of a statistic.

\begin{tcolorbox}[enhanced jigsaw, arc=.35mm, bottomrule=.15mm, bottomtitle=1mm, breakable, colback=white, colbacktitle=quarto-callout-note-color!10!white, colframe=quarto-callout-note-color-frame, coltitle=black, left=2mm, leftrule=.75mm, opacityback=0, opacitybacktitle=0.6, rightrule=.15mm, title=\textcolor{quarto-callout-note-color}{\faInfo}\hspace{0.5em}{Definition}, titlerule=0mm, toprule=.15mm, toptitle=1mm]

\textbf{Definition 5} (Statistic). Let \(X_1,\ldots,X_n\) be a random
sample from a population. A \textbf{statistic} is any function of the
sample: \[Y=T(X_1,\ldots,X_n).\] The distribution of \(Y\) is called the
\textbf{sampling distribution} of the statistic \(T\).

\end{tcolorbox}

After the data are observed, we use lower-case letters
\(x_1,\ldots,x_n\) for the observed values. The statistic
\(T(X_1,\ldots,X_n)\) is random before the data are observed, while
\(T(x_1,\ldots,x_n)\) is the realized numerical value after observation.

\begin{tcolorbox}[enhanced jigsaw, arc=.35mm, bottomrule=.15mm, bottomtitle=1mm, breakable, colback=white, colbacktitle=quarto-callout-tip-color!10!white, colframe=quarto-callout-tip-color-frame, coltitle=black, left=2mm, leftrule=.75mm, opacityback=0, opacitybacktitle=0.6, rightrule=.15mm, title=\textcolor{quarto-callout-tip-color}{\faLightbulb}\hspace{0.5em}{Example}, titlerule=0mm, toprule=.15mm, toptitle=1mm]

\textbf{Example 6} (Statistics). Each of the following is a statistic:

\begin{enumerate}
\def\labelenumi{\arabic{enumi}.}
\item
  \(T=\max{X_1,\ldots,X_n}\);
\item
  \(T=\min{X_1,\ldots,X_n}\);
\item
  \(T=\overline X\);
\item
  \(T=S^2\);
\item
  \(T=3\).
\end{enumerate}

The last example is a strange statistic because it ignores the data, but
it is still a function of the sample.

\end{tcolorbox}

\begin{tcolorbox}[enhanced jigsaw, arc=.35mm, bottomrule=.15mm, bottomtitle=1mm, breakable, colback=white, colbacktitle=quarto-callout-note-color!10!white, colframe=quarto-callout-note-color-frame, coltitle=black, left=2mm, leftrule=.75mm, opacityback=0, opacitybacktitle=0.6, rightrule=.15mm, title=\textcolor{quarto-callout-note-color}{\faInfo}\hspace{0.5em}{Solution}, titlerule=0mm, toprule=.15mm, toptitle=1mm]

A statistic only needs to be a function of the random sample. It does
not need to be useful. Therefore a constant such as \(T=3\) is
technically a statistic, although it contains no information about the
population.

\end{tcolorbox}

\subsection{Sample mean and sample
variance}\label{sample-mean-and-sample-variance}

This subsection introduces the two most common statistics in the course.

\begin{tcolorbox}[enhanced jigsaw, arc=.35mm, bottomrule=.15mm, bottomtitle=1mm, breakable, colback=white, colbacktitle=quarto-callout-note-color!10!white, colframe=quarto-callout-note-color-frame, coltitle=black, left=2mm, leftrule=.75mm, opacityback=0, opacitybacktitle=0.6, rightrule=.15mm, title=\textcolor{quarto-callout-note-color}{\faInfo}\hspace{0.5em}{Definition}, titlerule=0mm, toprule=.15mm, toptitle=1mm]

\textbf{Definition 7} (Sample mean). The \textbf{sample mean} is
\[\overline X=\frac{X_1+\cdots+X_n}{n}=\frac1n\sum_{i=1}^n X_i.\] For
observed values \(x_1,\ldots,x_n\), the observed sample mean is
\[\overline x=\frac1n\sum_{i=1}^n x_i.\]

\end{tcolorbox}

\begin{tcolorbox}[enhanced jigsaw, arc=.35mm, bottomrule=.15mm, bottomtitle=1mm, breakable, colback=white, colbacktitle=quarto-callout-note-color!10!white, colframe=quarto-callout-note-color-frame, coltitle=black, left=2mm, leftrule=.75mm, opacityback=0, opacitybacktitle=0.6, rightrule=.15mm, title=\textcolor{quarto-callout-note-color}{\faInfo}\hspace{0.5em}{Definition}, titlerule=0mm, toprule=.15mm, toptitle=1mm]

\textbf{Definition 8} (Sample variance and sample standard deviation).
The \textbf{sample variance} is
\[S^2=\frac{1}{n-1}\sum_{i=1}^n (X_i-\overline X)^2.\] The
\textbf{sample standard deviation} is \(S=\sqrt{S^2}\). Observed values
are denoted by \(s^2\) and \(s\).

\end{tcolorbox}

Two useful algebraic formulas for observed values are
\[\min_{a\in\mathbb{R}}\sum_{i=1}^n (x_i-a)^2=\sum_{i=1}^n (x_i-\overline x)^2,\]
and \[(n-1)s^2=\sum_{i=1}^n (x_i-\overline x)^2
=\sum_{i=1}^n x_i^2-n\overline x^{\,2}.\]

\begin{tcolorbox}[enhanced jigsaw, arc=.35mm, bottomrule=.15mm, bottomtitle=1mm, breakable, colback=white, colbacktitle=quarto-callout-important-color!10!white, colframe=quarto-callout-important-color-frame, coltitle=black, left=2mm, leftrule=.75mm, opacityback=0, opacitybacktitle=0.6, rightrule=.15mm, title=\textcolor{quarto-callout-important-color}{\faExclamation}\hspace{0.5em}{Proposition}, titlerule=0mm, toprule=.15mm, toptitle=1mm]

\textbf{Proposition 9} (Least-squares property of the sample mean).
\emph{For fixed observed values \(x_1,\ldots,x_n\), the function
\[Q(a)=\sum_{i=1}^n (x_i-a)^2\] is minimized at \(a=\overline x\).}

\end{tcolorbox}

\begin{tcolorbox}[enhanced jigsaw, arc=.35mm, bottomrule=.15mm, bottomtitle=1mm, breakable, colback=white, colbacktitle=quarto-callout-note-color!10!white, colframe=quarto-callout-note-color-frame, coltitle=black, left=2mm, leftrule=.75mm, opacityback=0, opacitybacktitle=0.6, rightrule=.15mm, title=\textcolor{quarto-callout-note-color}{\faInfo}\hspace{0.5em}{Proof}, titlerule=0mm, toprule=.15mm, toptitle=1mm]

\emph{Proof.} Expand and differentiate:
\[Q(a)=\sum_{i=1}^n x_i^2-2a\sum_{i=1}^n x_i+na^2.\] Then
\[Q'(a)=-2\sum_{i=1}^n x_i+2na.\] Setting \(Q'(a)=0\) gives
\[a=\frac1n\sum_{i=1}^n x_i=\overline x.\] Also \(Q''(a)=2n>0\), so this
critical point is a minimum.~◻

\end{tcolorbox}

\begin{tcolorbox}[enhanced jigsaw, arc=.35mm, bottomrule=.15mm, bottomtitle=1mm, breakable, colback=white, colbacktitle=quarto-callout-important-color!10!white, colframe=quarto-callout-important-color-frame, coltitle=black, left=2mm, leftrule=.75mm, opacityback=0, opacitybacktitle=0.6, rightrule=.15mm, title=\textcolor{quarto-callout-important-color}{\faExclamation}\hspace{0.5em}{Proposition}, titlerule=0mm, toprule=.15mm, toptitle=1mm]

\textbf{Proposition 10} (Computational formula for sample variance).
\emph{For observed data \(x_1,\ldots,x_n\),
\[\sum_{i=1}^n (x_i-\overline x)^2=\sum_{i=1}^n x_i^2-n\overline x^{\,2}.\]}

\end{tcolorbox}

\begin{tcolorbox}[enhanced jigsaw, arc=.35mm, bottomrule=.15mm, bottomtitle=1mm, breakable, colback=white, colbacktitle=quarto-callout-note-color!10!white, colframe=quarto-callout-note-color-frame, coltitle=black, left=2mm, leftrule=.75mm, opacityback=0, opacitybacktitle=0.6, rightrule=.15mm, title=\textcolor{quarto-callout-note-color}{\faInfo}\hspace{0.5em}{Proof}, titlerule=0mm, toprule=.15mm, toptitle=1mm]

\emph{Proof.} Using \(\sum_{i=1}^n x_i=n\overline x\), \[\begin{aligned}
\sum_{i=1}^n (x_i-\overline x)^2
&=\sum_{i=1}^n \left(x_i^2-2\overline x x_i+\overline x^{\,2}\right)\\
&=\sum_{i=1}^n x_i^2-2\overline x\sum_{i=1}^n x_i+n\overline x^{\,2}\\
&=\sum_{i=1}^n x_i^2-2n\overline x^{\,2}+n\overline x^{\,2}\\
&=\sum_{i=1}^n x_i^2-n\overline x^{\,2}.
\end{aligned}\]~◻

\end{tcolorbox}

\subsection{Mean and variance of the sample
mean}\label{mean-and-variance-of-the-sample-mean}

This subsection explains why \(\overline X\) estimates the population
mean and why its variance decreases as the sample size increases.

\begin{tcolorbox}[enhanced jigsaw, arc=.35mm, bottomrule=.15mm, bottomtitle=1mm, breakable, colback=white, colbacktitle=quarto-callout-important-color!10!white, colframe=quarto-callout-important-color-frame, coltitle=black, left=2mm, leftrule=.75mm, opacityback=0, opacitybacktitle=0.6, rightrule=.15mm, title=\textcolor{quarto-callout-important-color}{\faExclamation}\hspace{0.5em}{Theorem}, titlerule=0mm, toprule=.15mm, toptitle=1mm]

\textbf{Theorem 11} (Basic properties of \(\overline X\) and \(S^2\)).
\emph{Let \(X_1,\ldots,X_n\) be a random sample from a population with
mean \(\mu\) and variance \(\sigma^2\). Then
\[\mathbb{E}[\overline X]=\mu,
\qquad
\operatorname{Var}(\overline X)=\frac{\sigma^2}{n},
\qquad
\mathbb{E}[S^2]=\sigma^2.\] Thus \(\overline X\) is an unbiased
estimator of \(\mu\), and \(S^2\) is an unbiased estimator of
\(\sigma^2\).}

\end{tcolorbox}

\begin{tcolorbox}[enhanced jigsaw, arc=.35mm, bottomrule=.15mm, bottomtitle=1mm, breakable, colback=white, colbacktitle=quarto-callout-note-color!10!white, colframe=quarto-callout-note-color-frame, coltitle=black, left=2mm, leftrule=.75mm, opacityback=0, opacitybacktitle=0.6, rightrule=.15mm, title=\textcolor{quarto-callout-note-color}{\faInfo}\hspace{0.5em}{Proof}, titlerule=0mm, toprule=.15mm, toptitle=1mm]

\emph{Proof.} By linearity of expectation, \[\mathbb{E}[\overline X]
=\mathbb{E}\left[\frac1n\sum_{i=1}^n X_i\right]
=\frac1n\sum_{i=1}^n \mathbb{E}[X_i]
=\frac1n(n\mu)=\mu.\] Since the observations are independent,
\[\operatorname{Var}(\overline X)
=\operatorname{Var}\left(\frac1n\sum_{i=1}^n X_i\right)
=\frac1{n^2}\sum_{i=1}^n\operatorname{Var}(X_i)
=\frac{n\sigma^2}{n^2}=\frac{\sigma^2}{n}.\] For the sample variance,
use
\[\sum_{i=1}^n (X_i-\overline X)^2=\sum_{i=1}^n (X_i-\mu)^2-n(\overline X-\mu)^2.\]
Taking expectations gives
\[\mathbb{E}\left[\sum_{i=1}^n (X_i-\overline X)^2\right]
=n\sigma^2-n\operatorname{Var}(\overline X)
=n\sigma^2-n\frac{\sigma^2}{n}
=(n-1)\sigma^2.\] Therefore \[\mathbb{E}[S^2]
=\mathbb{E}\left[\frac1{n-1}\sum_{i=1}^n (X_i-\overline X)^2\right]
=\sigma^2.\]~◻

\end{tcolorbox}

\begin{tcolorbox}[enhanced jigsaw, arc=.35mm, bottomrule=.15mm, bottomtitle=1mm, breakable, colback=white, colbacktitle=quarto-callout-warning-color!10!white, colframe=quarto-callout-warning-color-frame, coltitle=black, left=2mm, leftrule=.75mm, opacityback=0, opacitybacktitle=0.6, rightrule=.15mm, title=\textcolor{quarto-callout-warning-color}{\faExclamationTriangle}\hspace{0.5em}{Practice Problem}, titlerule=0mm, toprule=.15mm, toptitle=1mm]

\textbf{Practice Problem 12} (Unbiasedness check). Let
\(X_1,\ldots,X_n\) be iid with mean \(10\) and variance \(4\). Find
\(\mathbb{E}[\overline X]\), \(\operatorname{Var}(\overline X)\), and
\(\mathbb{E}[S^2]\).

\end{tcolorbox}

\begin{tcolorbox}[enhanced jigsaw, arc=.35mm, bottomrule=.15mm, bottomtitle=1mm, breakable, colback=white, colbacktitle=quarto-callout-note-color!10!white, colframe=quarto-callout-note-color-frame, coltitle=black, left=2mm, leftrule=.75mm, opacityback=0, opacitybacktitle=0.6, rightrule=.15mm, title=\textcolor{quarto-callout-note-color}{\faInfo}\hspace{0.5em}{Solution}, titlerule=0mm, toprule=.15mm, toptitle=1mm]

Using the theorem, \[\mathbb{E}[\overline X]=10,
\qquad
\operatorname{Var}(\overline X)=\frac{4}{n},
\qquad
\mathbb{E}[S^2]=4.\]

\end{tcolorbox}

\section{Sums of Random Samples}\label{sums-of-random-samples}

This section studies how the distribution of a sum, and therefore the
sample mean, can be derived from the distribution of the original
sample.

Let \[Y=X_1+\cdots+X_n,
\qquad
\overline X=\frac{Y}{n}.\] If the pdf of \(Y\) is \(f_Y\), then the pdf
of \(\overline X=Y/n\) is \[f_{\overline X}(x)=n f_Y(nx).\] This follows
from the one-dimensional change-of-variables formula.

\subsection{MGF method}\label{mgf-method}

This subsection uses moment generating functions to identify the
distribution of sums and sample means.

\begin{tcolorbox}[enhanced jigsaw, arc=.35mm, bottomrule=.15mm, bottomtitle=1mm, breakable, colback=white, colbacktitle=quarto-callout-important-color!10!white, colframe=quarto-callout-important-color-frame, coltitle=black, left=2mm, leftrule=.75mm, opacityback=0, opacitybacktitle=0.6, rightrule=.15mm, title=\textcolor{quarto-callout-important-color}{\faExclamation}\hspace{0.5em}{Theorem}, titlerule=0mm, toprule=.15mm, toptitle=1mm]

\textbf{Theorem 13} (MGF of the sample mean). \emph{Let
\(X_1,\ldots,X_n\) be a random sample from a population with moment
generating function \(M_X(t)\). Then
\[M_{\overline X}(t)=\left[M_X\left(\frac{t}{n}\right)\right]^n.\]
Equivalently, for the sum \(Y=X_1+\cdots+X_n\),
\[M_Y(t)=\left[M_X(t)\right]^n.\]}

\end{tcolorbox}

\begin{tcolorbox}[enhanced jigsaw, arc=.35mm, bottomrule=.15mm, bottomtitle=1mm, breakable, colback=white, colbacktitle=quarto-callout-note-color!10!white, colframe=quarto-callout-note-color-frame, coltitle=black, left=2mm, leftrule=.75mm, opacityback=0, opacitybacktitle=0.6, rightrule=.15mm, title=\textcolor{quarto-callout-note-color}{\faInfo}\hspace{0.5em}{Proof}, titlerule=0mm, toprule=.15mm, toptitle=1mm]

\emph{Proof.} By independence,
\[M_Y(t)=\mathbb{E}[e^{t(X_1+\cdots+X_n)}]
=\prod_{i=1}^n \mathbb{E}[e^{tX_i}]
=\left[M_X(t)\right]^n.\] Since \(\overline X=Y/n\),
\[M_{\overline X}(t)
=\mathbb{E}[e^{tY/n}]
=M_Y(t/n)
=\left[M_X(t/n)\right]^n.\]~◻

\end{tcolorbox}

\begin{tcolorbox}[enhanced jigsaw, arc=.35mm, bottomrule=.15mm, bottomtitle=1mm, breakable, colback=white, colbacktitle=quarto-callout-tip-color!10!white, colframe=quarto-callout-tip-color-frame, coltitle=black, left=2mm, leftrule=.75mm, opacityback=0, opacitybacktitle=0.6, rightrule=.15mm, title=\textcolor{quarto-callout-tip-color}{\faLightbulb}\hspace{0.5em}{Example}, titlerule=0mm, toprule=.15mm, toptitle=1mm]

\textbf{Example 14} (Normal sample mean). Suppose
\(X_i\overset{\text{iid}}{\sim}\operatorname{Normal}(\mu,\sigma^2)\) for
\(i=1,\ldots,n\). Find the distribution of \(\overline X\).

\end{tcolorbox}

\begin{tcolorbox}[enhanced jigsaw, arc=.35mm, bottomrule=.15mm, bottomtitle=1mm, breakable, colback=white, colbacktitle=quarto-callout-note-color!10!white, colframe=quarto-callout-note-color-frame, coltitle=black, left=2mm, leftrule=.75mm, opacityback=0, opacitybacktitle=0.6, rightrule=.15mm, title=\textcolor{quarto-callout-note-color}{\faInfo}\hspace{0.5em}{Solution}, titlerule=0mm, toprule=.15mm, toptitle=1mm]

The MGF of \(X_i\) is
\[M_X(t)=\exp\left(\mu t+\frac{\sigma^2t^2}{2}\right).\] Thus
\[\begin{aligned}
M_{\overline X}(t)
&=\left[M_X(t/n)\right]^n\\
&=\left[\exp\left(\mu\frac{t}{n}+\frac{\sigma^2(t/n)^2}{2}\right)\right]^n\\
&=\exp\left(\mu t+\frac{\sigma^2 t^2}{2n}\right).
\end{aligned}\] This is the MGF of
\(\operatorname{Normal}(\mu,\sigma^2/n)\). Therefore
\[\overline X\sim \operatorname{Normal}\left(\mu,\frac{\sigma^2}{n}\right).\]

\end{tcolorbox}

\begin{tcolorbox}[enhanced jigsaw, arc=.35mm, bottomrule=.15mm, bottomtitle=1mm, breakable, colback=white, colbacktitle=quarto-callout-tip-color!10!white, colframe=quarto-callout-tip-color-frame, coltitle=black, left=2mm, leftrule=.75mm, opacityback=0, opacitybacktitle=0.6, rightrule=.15mm, title=\textcolor{quarto-callout-tip-color}{\faLightbulb}\hspace{0.5em}{Example}, titlerule=0mm, toprule=.15mm, toptitle=1mm]

\textbf{Example 15} (Gamma sample mean). Suppose
\(X_i\overset{\text{iid}}{\sim}\operatorname{Gamma}(\alpha,\beta)\),
where \(\beta\) is the rate parameter, so
\[M_X(t)=\left(\frac{\beta}{\beta-t}\right)^\alpha,
\qquad t<\beta.\] Find the distribution of \(\overline X\).

\end{tcolorbox}

\begin{tcolorbox}[enhanced jigsaw, arc=.35mm, bottomrule=.15mm, bottomtitle=1mm, breakable, colback=white, colbacktitle=quarto-callout-note-color!10!white, colframe=quarto-callout-note-color-frame, coltitle=black, left=2mm, leftrule=.75mm, opacityback=0, opacitybacktitle=0.6, rightrule=.15mm, title=\textcolor{quarto-callout-note-color}{\faInfo}\hspace{0.5em}{Solution}, titlerule=0mm, toprule=.15mm, toptitle=1mm]

First, \[Y=\sum_{i=1}^n X_i\sim \operatorname{Gamma}(n\alpha,\beta),\]
because the MGF of \(Y\) is
\[M_Y(t)=\left(\frac{\beta}{\beta-t}\right)^{n\alpha}.\] Since
\(\overline X=Y/n\), \[M_{\overline X}(t)=M_Y(t/n)
=\left(\frac{\beta}{\beta-t/n}\right)^{n\alpha}
=\left(\frac{n\beta}{n\beta-t}\right)^{n\alpha}.\] Therefore
\[\overline X\sim \operatorname{Gamma}(n\alpha,n\beta)\] under the rate
parameterization. Under a scale parameterization \(\theta=1/\beta\),
this is \[\overline X\sim \operatorname{Gamma}(n\alpha,\theta/n).\]

\end{tcolorbox}

\subsection{Convolution method}\label{convolution-method}

This subsection gives the general integral formula for the density of a
sum of independent continuous random variables.

\begin{tcolorbox}[enhanced jigsaw, arc=.35mm, bottomrule=.15mm, bottomtitle=1mm, breakable, colback=white, colbacktitle=quarto-callout-important-color!10!white, colframe=quarto-callout-important-color-frame, coltitle=black, left=2mm, leftrule=.75mm, opacityback=0, opacitybacktitle=0.6, rightrule=.15mm, title=\textcolor{quarto-callout-important-color}{\faExclamation}\hspace{0.5em}{Theorem}, titlerule=0mm, toprule=.15mm, toptitle=1mm]

\textbf{Theorem 16} (Convolution formula). \emph{If \(X\) and \(Y\) are
independent continuous random variables with pdfs \(f_X\) and \(f_Y\),
then the pdf of \(Z=X+Y\) is
\[f_Z(z)=\int_{-\infty}^{\infty} f_X(w)f_Y(z-w)\,dw.\]}

\end{tcolorbox}

\begin{tcolorbox}[enhanced jigsaw, arc=.35mm, bottomrule=.15mm, bottomtitle=1mm, breakable, colback=white, colbacktitle=quarto-callout-note-color!10!white, colframe=quarto-callout-note-color-frame, coltitle=black, left=2mm, leftrule=.75mm, opacityback=0, opacitybacktitle=0.6, rightrule=.15mm, title=\textcolor{quarto-callout-note-color}{\faInfo}\hspace{0.5em}{Proof}, titlerule=0mm, toprule=.15mm, toptitle=1mm]

\emph{Proof.} Let \[Z=X+Y,
\qquad
W=X.\] Then \[X=W,
\qquad
Y=Z-W.\] The Jacobian determinant of the inverse transformation
\((z,w)\mapsto (x,y)=(w,z-w)\) has absolute value \(1\). Hence
\[f_{Z,W}(z,w)=f_{X,Y}(w,z-w)=f_X(w)f_Y(z-w).\] Integrating out \(w\)
gives \[f_Z(z)=\int_{-\infty}^{\infty}f_X(w)f_Y(z-w)\,dw.\]~◻

\end{tcolorbox}

\begin{tcolorbox}[enhanced jigsaw, arc=.35mm, bottomrule=.15mm, bottomtitle=1mm, breakable, colback=white, colbacktitle=quarto-callout-tip-color!10!white, colframe=quarto-callout-tip-color-frame, coltitle=black, left=2mm, leftrule=.75mm, opacityback=0, opacitybacktitle=0.6, rightrule=.15mm, title=\textcolor{quarto-callout-tip-color}{\faLightbulb}\hspace{0.5em}{Example}, titlerule=0mm, toprule=.15mm, toptitle=1mm]

\textbf{Example 17} (Sum of two exponential random variables). Let
\(X,Y\overset{\text{iid}}{\sim}\operatorname{Exp}(1)\), so
\(f_X(x)=e^{-x}\) for \(x\geq0\). Find the pdf of \(Z=X+Y\).

\end{tcolorbox}

\begin{tcolorbox}[enhanced jigsaw, arc=.35mm, bottomrule=.15mm, bottomtitle=1mm, breakable, colback=white, colbacktitle=quarto-callout-note-color!10!white, colframe=quarto-callout-note-color-frame, coltitle=black, left=2mm, leftrule=.75mm, opacityback=0, opacitybacktitle=0.6, rightrule=.15mm, title=\textcolor{quarto-callout-note-color}{\faInfo}\hspace{0.5em}{Solution}, titlerule=0mm, toprule=.15mm, toptitle=1mm]

For \(z>0\), \[f_Z(z)=\int_{0}^{z}e^{-w}e^{-(z-w)}\,dw
=\int_0^z e^{-z}\,dw
=ze^{-z}.\] For \(z\leq0\), \(f_Z(z)=0\). Therefore
\[f_Z(z)=ze^{-z}\mathbf{1}_{\{z>0\}},\] which is a gamma density with
shape \(2\) and rate \(1\).

\end{tcolorbox}

\begin{tcolorbox}[enhanced jigsaw, arc=.35mm, bottomrule=.15mm, bottomtitle=1mm, breakable, colback=white, colbacktitle=quarto-callout-tip-color!10!white, colframe=quarto-callout-tip-color-frame, coltitle=black, left=2mm, leftrule=.75mm, opacityback=0, opacitybacktitle=0.6, rightrule=.15mm, title=\textcolor{quarto-callout-tip-color}{\faLightbulb}\hspace{0.5em}{Example}, titlerule=0mm, toprule=.15mm, toptitle=1mm]

\textbf{Example 18} (Sum of Cauchy random variables). Suppose \(U\) and
\(V\) are independent Cauchy random variables:
\[U\sim \operatorname{Cauchy}(0,\sigma),
\qquad
V\sim \operatorname{Cauchy}(0,\tau),\] with pdfs
\[f_U(u)=\frac{1}{\pi\sigma}\frac{1}{1+(u/\sigma)^2},
\qquad
f_V(v)=\frac{1}{\pi\tau}\frac{1}{1+(v/\tau)^2}.\] Find the distribution
of \(Z=U+V\).

\end{tcolorbox}

\begin{tcolorbox}[enhanced jigsaw, arc=.35mm, bottomrule=.15mm, bottomtitle=1mm, breakable, colback=white, colbacktitle=quarto-callout-note-color!10!white, colframe=quarto-callout-note-color-frame, coltitle=black, left=2mm, leftrule=.75mm, opacityback=0, opacitybacktitle=0.6, rightrule=.15mm, title=\textcolor{quarto-callout-note-color}{\faInfo}\hspace{0.5em}{Solution}, titlerule=0mm, toprule=.15mm, toptitle=1mm]

By convolution, \[f_Z(z)=\int_{-\infty}^{\infty}
\frac{1}{\pi\sigma}\frac{1}{1+(w/\sigma)^2}
\frac{1}{\pi\tau}\frac{1}{1+((z-w)/\tau)^2}\,dw.\] A standard Cauchy
convolution calculation gives
\[f_Z(z)=\frac{1}{\pi(\sigma+\tau)}\frac{1}{1+(z/(\sigma+\tau))^2}.\]
Thus \[Z=U+V\sim \operatorname{Cauchy}(0,\sigma+\tau).\] This is an
important example showing that averaging Cauchy variables does not
produce a distribution with smaller spread in the usual variance sense;
the Cauchy distribution does not have a finite mean or variance.

\end{tcolorbox}

\begin{tcolorbox}[enhanced jigsaw, arc=.35mm, bottomrule=.15mm, bottomtitle=1mm, breakable, colback=white, colbacktitle=quarto-callout-warning-color!10!white, colframe=quarto-callout-warning-color-frame, coltitle=black, left=2mm, leftrule=.75mm, opacityback=0, opacitybacktitle=0.6, rightrule=.15mm, title=\textcolor{quarto-callout-warning-color}{\faExclamationTriangle}\hspace{0.5em}{Practice Problem}, titlerule=0mm, toprule=.15mm, toptitle=1mm]

\textbf{Practice Problem 19} (Convolution practice). Let
\(X,Y\overset{\text{iid}}{\sim}\operatorname{Uniform}(0,1)\). Use
convolution to find the pdf of \(Z=X+Y\).

\end{tcolorbox}

\begin{tcolorbox}[enhanced jigsaw, arc=.35mm, bottomrule=.15mm, bottomtitle=1mm, breakable, colback=white, colbacktitle=quarto-callout-note-color!10!white, colframe=quarto-callout-note-color-frame, coltitle=black, left=2mm, leftrule=.75mm, opacityback=0, opacitybacktitle=0.6, rightrule=.15mm, title=\textcolor{quarto-callout-note-color}{\faInfo}\hspace{0.5em}{Solution}, titlerule=0mm, toprule=.15mm, toptitle=1mm]

For \(0<z<1\), \[f_Z(z)=\int_0^z 1\,dw=z.\] For \(1\leq z<2\),
\[f_Z(z)=\int_{z-1}^{1}1\,dw=2-z.\] Thus \[f_Z(z)=
\begin{cases}
z, &0<z<1,\\
2-z, &1\leq z<2,\\
0, &\text{otherwise}.
\end{cases}\]

\end{tcolorbox}

\section{Sampling from the Normal
Distribution}\label{sampling-from-the-normal-distribution}

This section presents the central normal-sampling facts that lead to the
chi-square, Student's \(t\), and \(F\) distributions.

\subsection{Normal sample mean and sample
variance}\label{normal-sample-mean-and-sample-variance}

This subsection states the special independence and chi-square results
that hold only for normal samples.

\begin{tcolorbox}[enhanced jigsaw, arc=.35mm, bottomrule=.15mm, bottomtitle=1mm, breakable, colback=white, colbacktitle=quarto-callout-important-color!10!white, colframe=quarto-callout-important-color-frame, coltitle=black, left=2mm, leftrule=.75mm, opacityback=0, opacitybacktitle=0.6, rightrule=.15mm, title=\textcolor{quarto-callout-important-color}{\faExclamation}\hspace{0.5em}{Theorem}, titlerule=0mm, toprule=.15mm, toptitle=1mm]

\textbf{Theorem 20} (Normal sampling theorem). \emph{Let
\(X_1,\ldots,X_n\) be a random sample from
\(\operatorname{Normal}(\mu,\sigma^2)\). Then}

\begin{enumerate}
\def\labelenumi{\arabic{enumi}.}
\item
  \emph{\(\overline X\sim \operatorname{Normal}\left(\mu,\sigma^2/n\right)\);}
\item
  \emph{\(\overline X\) and \(S^2\) are independent;}
\item
  \emph{\[\frac{(n-1)S^2}{\sigma^2}\sim \chi^2_{n-1}.\]}
\end{enumerate}

\end{tcolorbox}

\begin{tcolorbox}[enhanced jigsaw, arc=.35mm, bottomrule=.15mm, bottomtitle=1mm, breakable, colback=white, colbacktitle=quarto-callout-tip-color!10!white, colframe=quarto-callout-tip-color-frame, coltitle=black, left=2mm, leftrule=.75mm, opacityback=0, opacitybacktitle=0.6, rightrule=.15mm, title=\textcolor{quarto-callout-tip-color}{\faLightbulb}\hspace{0.5em}{Why this theorem matters}, titlerule=0mm, toprule=.15mm, toptitle=1mm]

The theorem says that for normal samples, uncertainty about the mean and
uncertainty about the variance separate cleanly. This is the starting
point for the \(t\) distribution, confidence intervals, and many
classical tests.

\end{tcolorbox}

\subsection{Chi-square random
variables}\label{chi-square-random-variables}

This subsection reviews the chi-square distribution and its relationship
with squared standard normals.

\begin{tcolorbox}[enhanced jigsaw, arc=.35mm, bottomrule=.15mm, bottomtitle=1mm, breakable, colback=white, colbacktitle=quarto-callout-important-color!10!white, colframe=quarto-callout-important-color-frame, coltitle=black, left=2mm, leftrule=.75mm, opacityback=0, opacitybacktitle=0.6, rightrule=.15mm, title=\textcolor{quarto-callout-important-color}{\faExclamation}\hspace{0.5em}{Theorem}, titlerule=0mm, toprule=.15mm, toptitle=1mm]

\textbf{Theorem 21} (Chi-square facts). \emph{The chi-square
distribution with \(k\) degrees of freedom is a gamma distribution:
\[\chi_k^2\sim \operatorname{Gamma}\left(\alpha=\frac{k}{2},\theta=2\right),\]
where \(\theta\) is the scale parameter. Also:}

\begin{enumerate}
\def\labelenumi{\arabic{enumi}.}
\item
  \emph{If \(Z\sim \operatorname{Normal}(0,1)\), then
  \(Z^2\sim \chi^2_1\).}
\item
  \emph{If \(X_i\sim\chi^2_{p_i}\) are independent, then
  \[X_1+\cdots+X_n\sim \chi^2_{p_1+\cdots+p_n}.\]}
\end{enumerate}

\end{tcolorbox}

\begin{tcolorbox}[enhanced jigsaw, arc=.35mm, bottomrule=.15mm, bottomtitle=1mm, breakable, colback=white, colbacktitle=quarto-callout-note-color!10!white, colframe=quarto-callout-note-color-frame, coltitle=black, left=2mm, leftrule=.75mm, opacityback=0, opacitybacktitle=0.6, rightrule=.15mm, title=\textcolor{quarto-callout-note-color}{\faInfo}\hspace{0.5em}{Proof}, titlerule=0mm, toprule=.15mm, toptitle=1mm]

\emph{Proof sketch.} If \(Z\sim\operatorname{Normal}(0,1)\), then by the
transformation \(Y=Z^2\), the density of \(Y\) is
\[f_Y(y)=\frac{1}{\sqrt{2\pi y}}e^{-y/2},\qquad y>0,\] which is
\(\operatorname{Gamma}(1/2,2)\), hence \(\chi_1^2\). The sum property
follows from multiplying MGFs. The MGF of \(\chi_k^2\) is
\((1-2t)^{-k/2}\), and products of such MGFs add the degrees of
freedom.~◻

\end{tcolorbox}

\begin{tcolorbox}[enhanced jigsaw, arc=.35mm, bottomrule=.15mm, bottomtitle=1mm, breakable, colback=white, colbacktitle=quarto-callout-tip-color!10!white, colframe=quarto-callout-tip-color-frame, coltitle=black, left=2mm, leftrule=.75mm, opacityback=0, opacitybacktitle=0.6, rightrule=.15mm, title=\textcolor{quarto-callout-tip-color}{\faLightbulb}\hspace{0.5em}{Example}, titlerule=0mm, toprule=.15mm, toptitle=1mm]

\textbf{Example 22} (Normal sample variance). Let \(X_1,\ldots,X_{10}\)
be a random sample from \(\operatorname{Normal}(\mu,\sigma^2)\). What is
the distribution of \(9S^2/\sigma^2\)?

\end{tcolorbox}

\begin{tcolorbox}[enhanced jigsaw, arc=.35mm, bottomrule=.15mm, bottomtitle=1mm, breakable, colback=white, colbacktitle=quarto-callout-note-color!10!white, colframe=quarto-callout-note-color-frame, coltitle=black, left=2mm, leftrule=.75mm, opacityback=0, opacitybacktitle=0.6, rightrule=.15mm, title=\textcolor{quarto-callout-note-color}{\faInfo}\hspace{0.5em}{Solution}, titlerule=0mm, toprule=.15mm, toptitle=1mm]

By the normal sampling theorem,
\[\frac{(n-1)S^2}{\sigma^2}\sim \chi^2_{n-1}.\] With \(n=10\),
\[\frac{9S^2}{\sigma^2}\sim \chi^2_9.\]

\end{tcolorbox}

\subsection{\texorpdfstring{Student's \(t\)
distribution}{Student's t distribution}}\label{students-t-distribution}

This subsection introduces the distribution used when the population
variance is unknown.

This topic was first addressed by W. S. Gosset, who published under the
pseudonym \emph{Student}, in the early 1900s.

If \(X_1,\ldots,X_n\) is a random sample from a normal population with
mean \(\mu\) and variance \(\sigma^2\), and if \(\sigma^2\) is known,
then
\[\frac{\overline X-\mu}{\sigma/\sqrt n}\sim \operatorname{Normal}(0,1).\]
When \(\sigma^2\) is unknown, we replace \(\sigma\) by the sample
standard deviation \(S\) and study
\[T=\frac{\overline X-\mu}{S/\sqrt n}.\]

\begin{tcolorbox}[enhanced jigsaw, arc=.35mm, bottomrule=.15mm, bottomtitle=1mm, breakable, colback=white, colbacktitle=quarto-callout-note-color!10!white, colframe=quarto-callout-note-color-frame, coltitle=black, left=2mm, leftrule=.75mm, opacityback=0, opacitybacktitle=0.6, rightrule=.15mm, title=\textcolor{quarto-callout-note-color}{\faInfo}\hspace{0.5em}{Definition}, titlerule=0mm, toprule=.15mm, toptitle=1mm]

\textbf{Definition 23} (Student's \(t\) distribution). Let
\(U\sim\operatorname{Normal}(0,1)\) and \(V\sim \chi^2_p\) be
independent. Then \[T=\frac{U}{\sqrt{V/p}}\] has a \textbf{Student's
\(t\) distribution} with \(p\) degrees of freedom, denoted \(t_p\).

\end{tcolorbox}

For a normal sample, \[\frac{\overline X-\mu}{S/\sqrt n}
=\frac{(\overline X-\mu)/(\sigma/\sqrt n)}{\sqrt{{(n-1)S^2/\sigma^2}/(n-1)}}
\sim t_{n-1}.\]

The pdf of \(T\sim t_p\) is
\[f_T(t)=\frac{\Gamma\left(\frac{p+1}{2}\right)}{\sqrt{p\pi}\,\Gamma\left(\frac p2\right)}
\left(1+\frac{t^2}{p}\right)^{-(p+1)/2},
\qquad -\infty<t<\infty.\]

\begin{tcolorbox}[enhanced jigsaw, arc=.35mm, bottomrule=.15mm, bottomtitle=1mm, breakable, colback=white, colbacktitle=quarto-callout-note-color!10!white, colframe=quarto-callout-note-color-frame, coltitle=black, left=2mm, leftrule=.75mm, opacityback=0, opacitybacktitle=0.6, rightrule=.15mm, title=\textcolor{quarto-callout-note-color}{\faInfo}\hspace{0.5em}{Remark}, titlerule=0mm, toprule=.15mm, toptitle=1mm]

\emph{Remark 24}. If \(p=1\), the Student's \(t\) distribution becomes
the Cauchy distribution. The \(t\) distribution has heavier tails than
the normal distribution. It has no MGF because it does not have moments
of all orders. If \(T\sim t_p\), then only moments of order less than
\(p\) exist. In particular, \[\mathbb{E}[T]=0 \quad \text{for } p>1,
\qquad
\operatorname{Var}(T)=\frac{p}{p-2}\quad \text{for } p>2.\]

\end{tcolorbox}

\begin{tcolorbox}[enhanced jigsaw, arc=.35mm, bottomrule=.15mm, bottomtitle=1mm, breakable, colback=white, colbacktitle=quarto-callout-tip-color!10!white, colframe=quarto-callout-tip-color-frame, coltitle=black, left=2mm, leftrule=.75mm, opacityback=0, opacitybacktitle=0.6, rightrule=.15mm, title=\textcolor{quarto-callout-tip-color}{\faLightbulb}\hspace{0.5em}{Example}, titlerule=0mm, toprule=.15mm, toptitle=1mm]

\textbf{Example 25} (Constructing a \(t\) statistic). Suppose
\(X_1,\ldots,X_{16}\) are iid \(\operatorname{Normal}(\mu,\sigma^2)\).
Find the distribution of \[T=\frac{\overline X-\mu}{S/4}.\]

\end{tcolorbox}

\begin{tcolorbox}[enhanced jigsaw, arc=.35mm, bottomrule=.15mm, bottomtitle=1mm, breakable, colback=white, colbacktitle=quarto-callout-note-color!10!white, colframe=quarto-callout-note-color-frame, coltitle=black, left=2mm, leftrule=.75mm, opacityback=0, opacitybacktitle=0.6, rightrule=.15mm, title=\textcolor{quarto-callout-note-color}{\faInfo}\hspace{0.5em}{Solution}, titlerule=0mm, toprule=.15mm, toptitle=1mm]

Since \(n=16\), \(\sqrt n=4\). Therefore
\[T=\frac{\overline X-\mu}{S/\sqrt n}\sim t_{n-1}=t_{15}.\]

\end{tcolorbox}

\subsection{\texorpdfstring{\(F\)
distribution}{F distribution}}\label{f-distribution}

This subsection introduces the distribution of a ratio of scaled
independent chi-square random variables, which is used to compare
variances.

Let \(X_1,\ldots,X_n\) be a random sample from
\(\operatorname{Normal}(\mu_X,\sigma_X^2)\), and let \(Y_1,\ldots,Y_m\)
be an independent random sample from
\(\operatorname{Normal}(\mu_Y,\sigma_Y^2)\). Then
\[\frac{(n-1)S_X^2}{\sigma_X^2}\sim \chi^2_{n-1},
\qquad
\frac{(m-1)S_Y^2}{\sigma_Y^2}\sim \chi^2_{m-1}.\]

\begin{tcolorbox}[enhanced jigsaw, arc=.35mm, bottomrule=.15mm, bottomtitle=1mm, breakable, colback=white, colbacktitle=quarto-callout-note-color!10!white, colframe=quarto-callout-note-color-frame, coltitle=black, left=2mm, leftrule=.75mm, opacityback=0, opacitybacktitle=0.6, rightrule=.15mm, title=\textcolor{quarto-callout-note-color}{\faInfo}\hspace{0.5em}{Definition}, titlerule=0mm, toprule=.15mm, toptitle=1mm]

\textbf{Definition 26} (\(F\) distribution). If \(U\sim\chi^2_p\) and
\(V\sim\chi^2_q\) are independent, then \[F=\frac{U/p}{V/q}\] has an
\textbf{\(F\) distribution} with \(p\) and \(q\) degrees of freedom,
denoted \(F_{p,q}\).

\end{tcolorbox}

Thus \[\frac{S_X^2/\sigma_X^2}{S_Y^2/\sigma_Y^2}
\sim F_{n-1,m-1}.\] The pdf of \(F\sim F_{p,q}\) is, for \(x>0\),
\[f_F(x)=
\frac{1}{B(p/2,q/2)}
\left(\frac{p}{q}\right)^{p/2}
\frac{x^{p/2-1}}{\left(1+\frac{p}{q}x\right)^{(p+q)/2}}.\]

\begin{tcolorbox}[enhanced jigsaw, arc=.35mm, bottomrule=.15mm, bottomtitle=1mm, breakable, colback=white, colbacktitle=quarto-callout-important-color!10!white, colframe=quarto-callout-important-color-frame, coltitle=black, left=2mm, leftrule=.75mm, opacityback=0, opacitybacktitle=0.6, rightrule=.15mm, title=\textcolor{quarto-callout-important-color}{\faExclamation}\hspace{0.5em}{Theorem}, titlerule=0mm, toprule=.15mm, toptitle=1mm]

\textbf{Theorem 27} (Basic relationships for the \(F\) distribution).

\begin{enumerate}
\def\labelenumi{\arabic{enumi}.}
\item
  \emph{If \(X\sim F_{p,q}\), then \(1/X\sim F_{q,p}\).}
\item
  \emph{If \(T\sim t_q\), then \(T^2\sim F_{1,q}\).}
\item
  \emph{If \(X\sim F_{p,q}\), then
  \[\frac{(p/q)X}{1+(p/q)X}\sim \operatorname{Beta}\left(\frac p2,\frac q2\right).\]}
\end{enumerate}

\end{tcolorbox}

\begin{tcolorbox}[enhanced jigsaw, arc=.35mm, bottomrule=.15mm, bottomtitle=1mm, breakable, colback=white, colbacktitle=quarto-callout-tip-color!10!white, colframe=quarto-callout-tip-color-frame, coltitle=black, left=2mm, leftrule=.75mm, opacityback=0, opacitybacktitle=0.6, rightrule=.15mm, title=\textcolor{quarto-callout-tip-color}{\faLightbulb}\hspace{0.5em}{Example}, titlerule=0mm, toprule=.15mm, toptitle=1mm]

\textbf{Example 28} (Expected variance ratio). Let
\[F=\frac{S_X^2/\sigma_X^2}{S_Y^2/\sigma_Y^2}\sim F_{n-1,m-1}.\] Assume
\(m>3\). Find \(\mathbb{E}[F]\).

\end{tcolorbox}

\begin{tcolorbox}[enhanced jigsaw, arc=.35mm, bottomrule=.15mm, bottomtitle=1mm, breakable, colback=white, colbacktitle=quarto-callout-note-color!10!white, colframe=quarto-callout-note-color-frame, coltitle=black, left=2mm, leftrule=.75mm, opacityback=0, opacitybacktitle=0.6, rightrule=.15mm, title=\textcolor{quarto-callout-note-color}{\faInfo}\hspace{0.5em}{Solution}, titlerule=0mm, toprule=.15mm, toptitle=1mm]

Write \[F=\frac{\chi^2_{n-1}/(n-1)}{\chi^2_{m-1}/(m-1)},\] where the
numerator and denominator are independent. Since
\[\mathbb{E}\left[\frac{\chi^2_{n-1}}{n-1}\right]=1\] and, for
\(m-1>2\),
\[\mathbb{E}\left[\frac{m-1}{\chi^2_{m-1}}\right]=\frac{m-1}{m-3},\] we
get \[\mathbb{E}[F]=\frac{m-1}{m-3}.\] For large \(m\),
\[\mathbb{E}[F]=\frac{m-1}{m-3}\approx 1.\]

\end{tcolorbox}

\begin{tcolorbox}[enhanced jigsaw, arc=.35mm, bottomrule=.15mm, bottomtitle=1mm, breakable, colback=white, colbacktitle=quarto-callout-warning-color!10!white, colframe=quarto-callout-warning-color-frame, coltitle=black, left=2mm, leftrule=.75mm, opacityback=0, opacitybacktitle=0.6, rightrule=.15mm, title=\textcolor{quarto-callout-warning-color}{\faExclamationTriangle}\hspace{0.5em}{Practice Problem}, titlerule=0mm, toprule=.15mm, toptitle=1mm]

\textbf{Practice Problem 29} (\(t\) and \(F\)). If \(T\sim t_{12}\),
what is the distribution of \(T^2\)?

\end{tcolorbox}

\begin{tcolorbox}[enhanced jigsaw, arc=.35mm, bottomrule=.15mm, bottomtitle=1mm, breakable, colback=white, colbacktitle=quarto-callout-note-color!10!white, colframe=quarto-callout-note-color-frame, coltitle=black, left=2mm, leftrule=.75mm, opacityback=0, opacitybacktitle=0.6, rightrule=.15mm, title=\textcolor{quarto-callout-note-color}{\faInfo}\hspace{0.5em}{Solution}, titlerule=0mm, toprule=.15mm, toptitle=1mm]

By the relationship between \(t\) and \(F\) distributions,
\[T^2\sim F_{1,12}.\]

\end{tcolorbox}

\section{Order Statistics}\label{order-statistics}

This section studies the sorted values of a random sample, which are
used to describe extremes, quantiles, ranges, and medians.

\begin{tcolorbox}[enhanced jigsaw, arc=.35mm, bottomrule=.15mm, bottomtitle=1mm, breakable, colback=white, colbacktitle=quarto-callout-note-color!10!white, colframe=quarto-callout-note-color-frame, coltitle=black, left=2mm, leftrule=.75mm, opacityback=0, opacitybacktitle=0.6, rightrule=.15mm, title=\textcolor{quarto-callout-note-color}{\faInfo}\hspace{0.5em}{Definition}, titlerule=0mm, toprule=.15mm, toptitle=1mm]

\textbf{Definition 30} (Order statistics). Let \(X_1,\ldots,X_n\) be a
random sample. The \textbf{order statistics} are the ordered values
\[X_{(1)}<X_{(2)}<\cdots<X_{(n)},\] where \(X_{(j)}\) is the \(j\)-th
smallest value among \(X_1,\ldots,X_n\). In particular,
\[X_{(1)}=\min{X_1,\ldots,X_n},
\qquad
X_{(n)}=\max{X_1,\ldots,X_n}.\]

\end{tcolorbox}

If ties are possible, as in discrete distributions, one may write
\[X_{(1)}\leq X_{(2)}\leq\cdots\leq X_{(n)}.\] For continuous
distributions, ties have probability zero, so strict ordering is usually
safe.

\subsection{Discrete order statistics}\label{discrete-order-statistics}

This subsection gives the distribution of an order statistic when the
population is discrete.

\begin{tcolorbox}[enhanced jigsaw, arc=.35mm, bottomrule=.15mm, bottomtitle=1mm, breakable, colback=white, colbacktitle=quarto-callout-important-color!10!white, colframe=quarto-callout-important-color-frame, coltitle=black, left=2mm, leftrule=.75mm, opacityback=0, opacitybacktitle=0.6, rightrule=.15mm, title=\textcolor{quarto-callout-important-color}{\faExclamation}\hspace{0.5em}{Theorem}, titlerule=0mm, toprule=.15mm, toptitle=1mm]

\textbf{Theorem 31} (Discrete order statistics). \emph{Let
\(X_1,\ldots,X_n\) be a random sample from a discrete distribution with
pmf \[f_X(x_i)=p_i,
\qquad x_1<x_2<\cdots.\] Define \[P_0=0,
\qquad
P_i=p_1+\cdots+p_i.\] Then
\[\mathbb{P}(X_{(j)}\leq x_i)=\sum_{k=j}^n {n\choose k}P_i^k(1-P_i)^{n-k},\]
and \[\mathbb{P}(X_{(j)}=x_i)
=\sum_{k=j}^n {n\choose k}
\left[P_i^k(1-P_i)^{n-k}-P_{i-1}^k(1-P_{i-1})^{n-k}\right].\]}

\end{tcolorbox}

\begin{tcolorbox}[enhanced jigsaw, arc=.35mm, bottomrule=.15mm, bottomtitle=1mm, breakable, colback=white, colbacktitle=quarto-callout-note-color!10!white, colframe=quarto-callout-note-color-frame, coltitle=black, left=2mm, leftrule=.75mm, opacityback=0, opacitybacktitle=0.6, rightrule=.15mm, title=\textcolor{quarto-callout-note-color}{\faInfo}\hspace{0.5em}{Proof}, titlerule=0mm, toprule=.15mm, toptitle=1mm]

\emph{Proof.} The event \({X_{(j)}\leq x_i}\) occurs exactly when at
least \(j\) of the \(n\) observations are less than or equal to \(x_i\).
Each observation has probability \(P_i\) of being less than or equal to
\(x_i\). Therefore the number of observations less than or equal to
\(x_i\) is \(\operatorname{Binomial}(n,P_i)\), giving
\[\mathbb{P}(X_{(j)}\leq x_i)=\sum_{k=j}^n {n\choose k}P_i^k(1-P_i)^{n-k}.\]
Then \[\mathbb{P}(X_{(j)}=x_i)
=\mathbb{P}(X_{(j)}\leq x_i)-\mathbb{P}(X_{(j)}\leq x_{i-1}),\] which
gives the stated formula.~◻

\end{tcolorbox}

\begin{tcolorbox}[enhanced jigsaw, arc=.35mm, bottomrule=.15mm, bottomtitle=1mm, breakable, colback=white, colbacktitle=quarto-callout-tip-color!10!white, colframe=quarto-callout-tip-color-frame, coltitle=black, left=2mm, leftrule=.75mm, opacityback=0, opacitybacktitle=0.6, rightrule=.15mm, title=\textcolor{quarto-callout-tip-color}{\faLightbulb}\hspace{0.5em}{Example}, titlerule=0mm, toprule=.15mm, toptitle=1mm]

\textbf{Example 32} (Median of three fair die rolls). Roll a fair
six-sided die three times. Let \(X_{(2)}\) be the sample median. Find
\(\mathbb{P}(X_{(2)}\leq 4)\).

\end{tcolorbox}

\begin{tcolorbox}[enhanced jigsaw, arc=.35mm, bottomrule=.15mm, bottomtitle=1mm, breakable, colback=white, colbacktitle=quarto-callout-note-color!10!white, colframe=quarto-callout-note-color-frame, coltitle=black, left=2mm, leftrule=.75mm, opacityback=0, opacitybacktitle=0.6, rightrule=.15mm, title=\textcolor{quarto-callout-note-color}{\faInfo}\hspace{0.5em}{Solution}, titlerule=0mm, toprule=.15mm, toptitle=1mm]

For a fair die, \[P_4=\mathbb{P}(X\leq 4)=\frac46=\frac23.\] The event
\(X_{(2)}\leq4\) means at least two of the three observations are at
most \(4\). Hence \[\begin{aligned}
\mathbb{P}(X_{(2)}\leq4)
&={3\choose2}\left(\frac23\right)^2\left(\frac13\right)
+{3\choose3}\left(\frac23\right)^3\\
&=3\cdot\frac49\cdot\frac13+\frac8{27}
=\frac{12}{27}+\frac{8}{27}
=\frac{20}{27}.
\end{aligned}\]

\end{tcolorbox}

\subsection{Continuous order
statistics}\label{continuous-order-statistics}

This subsection gives the density formulas for continuous order
statistics.

\begin{tcolorbox}[enhanced jigsaw, arc=.35mm, bottomrule=.15mm, bottomtitle=1mm, breakable, colback=white, colbacktitle=quarto-callout-important-color!10!white, colframe=quarto-callout-important-color-frame, coltitle=black, left=2mm, leftrule=.75mm, opacityback=0, opacitybacktitle=0.6, rightrule=.15mm, title=\textcolor{quarto-callout-important-color}{\faExclamation}\hspace{0.5em}{Theorem}, titlerule=0mm, toprule=.15mm, toptitle=1mm]

\textbf{Theorem 33} (Pdf of a continuous order statistic). \emph{Let
\(X_1,\ldots,X_n\) be a random sample from a continuous population with
cdf \(F_X(x)\) and pdf \(f_X(x)\). Then the pdf of \(X_{(j)}\) is
\[f_{X_{(j)}}(x)
=\frac{n!}{(j-1)!(n-j)!}
 f_X(x)[F_X(x)]^{j-1}[1-F_X(x)]^{n-j}.\]}

\end{tcolorbox}

\begin{tcolorbox}[enhanced jigsaw, arc=.35mm, bottomrule=.15mm, bottomtitle=1mm, breakable, colback=white, colbacktitle=quarto-callout-note-color!10!white, colframe=quarto-callout-note-color-frame, coltitle=black, left=2mm, leftrule=.75mm, opacityback=0, opacitybacktitle=0.6, rightrule=.15mm, title=\textcolor{quarto-callout-note-color}{\faInfo}\hspace{0.5em}{Proof}, titlerule=0mm, toprule=.15mm, toptitle=1mm]

\emph{Proof.} For \(X_{(j)}\) to fall in a small interval near \(x\),
one observation must fall near \(x\), exactly \(j-1\) observations must
be below \(x\), and \(n-j\) observations must be above \(x\). The
combinatorial factor counts which observations play these roles:
\[\frac{n!}{(j-1)!1!(n-j)!}.\] The probability contribution is
approximately \[[F_X(x)]^{j-1} f_X(x)\,dx [1-F_X(x)]^{n-j}.\] Dividing
by \(dx\) gives the density.~◻

\end{tcolorbox}

\begin{tcolorbox}[enhanced jigsaw, arc=.35mm, bottomrule=.15mm, bottomtitle=1mm, breakable, colback=white, colbacktitle=quarto-callout-important-color!10!white, colframe=quarto-callout-important-color-frame, coltitle=black, left=2mm, leftrule=.75mm, opacityback=0, opacitybacktitle=0.6, rightrule=.15mm, title=\textcolor{quarto-callout-important-color}{\faExclamation}\hspace{0.5em}{Theorem}, titlerule=0mm, toprule=.15mm, toptitle=1mm]

\textbf{Theorem 34} (Joint pdf of two continuous order statistics).
\emph{Let \(1\leq i<j\leq n\). The joint pdf of \(X_{(i)}\) and
\(X_{(j)}\) is, for \(u<v\), \[\begin{aligned}
f_{X_{(i)},X_{(j)}}(u,v)
&=\frac{n!}{(i-1)!(j-1-i)!(n-j)!}
 f_X(u)f_X(v)[F_X(u)]^{i-1}\\
&\qquad \times [F_X(v)-F_X(u)]^{j-1-i}[1-F_X(v)]^{n-j}.
\end{aligned}\] It is \(0\) otherwise.}

\end{tcolorbox}

\begin{tcolorbox}[enhanced jigsaw, arc=.35mm, bottomrule=.15mm, bottomtitle=1mm, breakable, colback=white, colbacktitle=quarto-callout-important-color!10!white, colframe=quarto-callout-important-color-frame, coltitle=black, left=2mm, leftrule=.75mm, opacityback=0, opacitybacktitle=0.6, rightrule=.15mm, title=\textcolor{quarto-callout-important-color}{\faExclamation}\hspace{0.5em}{Theorem}, titlerule=0mm, toprule=.15mm, toptitle=1mm]

\textbf{Theorem 35} (Joint pdf of all continuous order statistics).
\emph{The joint pdf of \(X_{(1)},X_{(2)},\ldots,X_{(n)}\) is
\[f_{X_{(1)},\ldots,X_{(n)}}(x_1,\ldots,x_n)
=
\begin{cases}
n! f_X(x_1)\cdots f_X(x_n), & -\infty<x_1<\cdots<x_n<\infty,\\
0, &\text{otherwise}.
\end{cases}\]}

\end{tcolorbox}

\begin{tcolorbox}[enhanced jigsaw, arc=.35mm, bottomrule=.15mm, bottomtitle=1mm, breakable, colback=white, colbacktitle=quarto-callout-tip-color!10!white, colframe=quarto-callout-tip-color-frame, coltitle=black, left=2mm, leftrule=.75mm, opacityback=0, opacitybacktitle=0.6, rightrule=.15mm, title=\textcolor{quarto-callout-tip-color}{\faLightbulb}\hspace{0.5em}{Example}, titlerule=0mm, toprule=.15mm, toptitle=1mm]

\textbf{Example 36} (Uniform order statistic). Let \(X_1,\ldots,X_n\) be
a random sample from \(\operatorname{Uniform}(0,1)\). Find the
distribution of \(X_{(j)}\).

\end{tcolorbox}

\begin{tcolorbox}[enhanced jigsaw, arc=.35mm, bottomrule=.15mm, bottomtitle=1mm, breakable, colback=white, colbacktitle=quarto-callout-note-color!10!white, colframe=quarto-callout-note-color-frame, coltitle=black, left=2mm, leftrule=.75mm, opacityback=0, opacitybacktitle=0.6, rightrule=.15mm, title=\textcolor{quarto-callout-note-color}{\faInfo}\hspace{0.5em}{Solution}, titlerule=0mm, toprule=.15mm, toptitle=1mm]

For \(\operatorname{Uniform}(0,1)\), \[f_X(x)=1,
\qquad
F_X(x)=x,
\qquad 0<x<1.\] Thus \[f_{X_{(j)}}(x)
=\frac{n!}{(j-1)!(n-j)!}x^{j-1}(1-x)^{n-j},
\qquad 0<x<1.\] This is the beta density with parameters \(j\) and
\(n-j+1\). Therefore \[X_{(j)}\sim \operatorname{Beta}(j,n-j+1).\]

\end{tcolorbox}

\begin{tcolorbox}[enhanced jigsaw, arc=.35mm, bottomrule=.15mm, bottomtitle=1mm, breakable, colback=white, colbacktitle=quarto-callout-warning-color!10!white, colframe=quarto-callout-warning-color-frame, coltitle=black, left=2mm, leftrule=.75mm, opacityback=0, opacitybacktitle=0.6, rightrule=.15mm, title=\textcolor{quarto-callout-warning-color}{\faExclamationTriangle}\hspace{0.5em}{Practice Problem}, titlerule=0mm, toprule=.15mm, toptitle=1mm]

\textbf{Practice Problem 37} (Maximum of uniforms). Let
\(X_1,\ldots,X_n\overset{\text{iid}}{\sim}\operatorname{Uniform}(0,1)\).
Find the cdf and pdf of \(X_{(n)}=\max{X_1,\ldots,X_n}\).

\end{tcolorbox}

\begin{tcolorbox}[enhanced jigsaw, arc=.35mm, bottomrule=.15mm, bottomtitle=1mm, breakable, colback=white, colbacktitle=quarto-callout-note-color!10!white, colframe=quarto-callout-note-color-frame, coltitle=black, left=2mm, leftrule=.75mm, opacityback=0, opacitybacktitle=0.6, rightrule=.15mm, title=\textcolor{quarto-callout-note-color}{\faInfo}\hspace{0.5em}{Solution}, titlerule=0mm, toprule=.15mm, toptitle=1mm]

For \(0<x<1\),
\[F_{X_{(n)}}(x)=\mathbb{P}(X_{(n)}\leq x)=\mathbb{P}(X_1\leq x,\ldots,X_n\leq x)=x^n.\]
Therefore \[f_{X_{(n)}}(x)=nx^{n-1},\qquad 0<x<1.\] This agrees with
\(\operatorname{Beta}(n,1)\).

\end{tcolorbox}

\subsection{Sample range and median}\label{sample-range-and-median}

This subsection introduces range, median, and quartiles as statistics
built from order statistics.

\begin{tcolorbox}[enhanced jigsaw, arc=.35mm, bottomrule=.15mm, bottomtitle=1mm, breakable, colback=white, colbacktitle=quarto-callout-note-color!10!white, colframe=quarto-callout-note-color-frame, coltitle=black, left=2mm, leftrule=.75mm, opacityback=0, opacitybacktitle=0.6, rightrule=.15mm, title=\textcolor{quarto-callout-note-color}{\faInfo}\hspace{0.5em}{Definition}, titlerule=0mm, toprule=.15mm, toptitle=1mm]

\textbf{Definition 38} (Sample range). The \textbf{sample range} is
\[R=X_{(n)}-X_{(1)}.\] It measures the distance between the smallest and
largest observations.

\end{tcolorbox}

\begin{tcolorbox}[enhanced jigsaw, arc=.35mm, bottomrule=.15mm, bottomtitle=1mm, breakable, colback=white, colbacktitle=quarto-callout-note-color!10!white, colframe=quarto-callout-note-color-frame, coltitle=black, left=2mm, leftrule=.75mm, opacityback=0, opacitybacktitle=0.6, rightrule=.15mm, title=\textcolor{quarto-callout-note-color}{\faInfo}\hspace{0.5em}{Definition}, titlerule=0mm, toprule=.15mm, toptitle=1mm]

\textbf{Definition 39} (Sample median and quartiles). The \textbf{sample
median} is a number \(M\) such that approximately one-half of the
observations are below \(M\) and one-half are above \(M\). The lower
quartile is the 25th percentile, and the upper quartile is the 75th
percentile.

\end{tcolorbox}

\begin{tcolorbox}[enhanced jigsaw, arc=.35mm, bottomrule=.15mm, bottomtitle=1mm, breakable, colback=white, colbacktitle=quarto-callout-tip-color!10!white, colframe=quarto-callout-tip-color-frame, coltitle=black, left=2mm, leftrule=.75mm, opacityback=0, opacitybacktitle=0.6, rightrule=.15mm, title=\textcolor{quarto-callout-tip-color}{\faLightbulb}\hspace{0.5em}{Example}, titlerule=0mm, toprule=.15mm, toptitle=1mm]

\textbf{Example 40} (Range from a uniform sample). Let
\(X_1,\ldots,X_n\) be a random sample from
\(\operatorname{Uniform}(0,a)\). Find the distribution of the sample
range \[R=X_{(n)}-X_{(1)}.\]

\end{tcolorbox}

\begin{tcolorbox}[enhanced jigsaw, arc=.35mm, bottomrule=.15mm, bottomtitle=1mm, breakable, colback=white, colbacktitle=quarto-callout-note-color!10!white, colframe=quarto-callout-note-color-frame, coltitle=black, left=2mm, leftrule=.75mm, opacityback=0, opacitybacktitle=0.6, rightrule=.15mm, title=\textcolor{quarto-callout-note-color}{\faInfo}\hspace{0.5em}{Solution}, titlerule=0mm, toprule=.15mm, toptitle=1mm]

The joint pdf of the minimum and maximum is
\[f_{X_{(1)},X_{(n)}}(x_1,x_n)
=\frac{n(n-1)}{a^n}(x_n-x_1)^{n-2},
\qquad 0<x_1<x_n<a.\] Define \[V=\frac{X_{(n)}+X_{(1)}}{2},
\qquad
R=X_{(n)}-X_{(1)}.\] Then \[X_{(1)}=V-\frac{R}{2},
\qquad
X_{(n)}=V+\frac{R}{2}.\] The absolute value of the Jacobian is \(1\).
The support is \[0<r<a,
\qquad
\frac r2<v<a-\frac r2.\] Thus \[f_{R,V}(r,v)=\frac{n(n-1)r^{n-2}}{a^n},
\qquad 0<r<a,\quad \frac r2<v<a-\frac r2.\] Integrating out \(v\) gives
\[\begin{aligned}
f_R(r)
&=\int_{r/2}^{a-r/2}\frac{n(n-1)r^{n-2}}{a^n}\,dv\\
&=\frac{n(n-1)r^{n-2}(a-r)}{a^n},
\qquad 0<r<a.
\end{aligned}\] If \(a=1\), then \[R\sim \operatorname{Beta}(n-1,2).\]
For arbitrary \(a\), \[\frac{R}{a}\sim \operatorname{Beta}(n-1,2).\]

\end{tcolorbox}

\begin{tcolorbox}[enhanced jigsaw, arc=.35mm, bottomrule=.15mm, bottomtitle=1mm, breakable, colback=white, colbacktitle=quarto-callout-warning-color!10!white, colframe=quarto-callout-warning-color-frame, coltitle=black, left=2mm, leftrule=.75mm, opacityback=0, opacitybacktitle=0.6, rightrule=.15mm, title=\textcolor{quarto-callout-warning-color}{\faExclamationTriangle}\hspace{0.5em}{Practice Problem}, titlerule=0mm, toprule=.15mm, toptitle=1mm]

\textbf{Practice Problem 41} (Expected range from uniforms). Let
\(X_1,\ldots,X_n\overset{\text{iid}}{\sim}\operatorname{Uniform}(0,1)\).
Use the distribution of \(R\) to compute \(\mathbb{E}[R]\).

\end{tcolorbox}

\begin{tcolorbox}[enhanced jigsaw, arc=.35mm, bottomrule=.15mm, bottomtitle=1mm, breakable, colback=white, colbacktitle=quarto-callout-note-color!10!white, colframe=quarto-callout-note-color-frame, coltitle=black, left=2mm, leftrule=.75mm, opacityback=0, opacitybacktitle=0.6, rightrule=.15mm, title=\textcolor{quarto-callout-note-color}{\faInfo}\hspace{0.5em}{Solution}, titlerule=0mm, toprule=.15mm, toptitle=1mm]

For \(a=1\), \[f_R(r)=n(n-1)r^{n-2}(1-r),
\qquad 0<r<1.\] Therefore \[\begin{aligned}
\mathbb{E}[R]
&=\int_0^1 r\,n(n-1)r^{n-2}(1-r)\,dr\\
&=n(n-1)\int_0^1 (r^{n-1}-r^n)\,dr\\
&=n(n-1)\left(\frac1n-\frac1{n+1}\right)\\
&=\frac{n-1}{n+1}.
\end{aligned}\]

\end{tcolorbox}

\section{Summary}\label{summary-4}

This section summarizes the main formulas and ideas from sampling and
order statistics.

\begin{tcolorbox}[enhanced jigsaw, arc=.35mm, bottomrule=.15mm, bottomtitle=1mm, breakable, colback=white, colbacktitle=quarto-callout-tip-color!10!white, colframe=quarto-callout-tip-color-frame, coltitle=black, left=2mm, leftrule=.75mm, opacityback=0, opacitybacktitle=0.6, rightrule=.15mm, title=\textcolor{quarto-callout-tip-color}{\faLightbulb}\hspace{0.5em}{Key formulas}, titlerule=0mm, toprule=.15mm, toptitle=1mm]

For a random sample \(X_1,\ldots,X_n\) with mean \(\mu\) and variance
\(\sigma^2\), \[\mathbb{E}[\overline X]=\mu,
\qquad
\operatorname{Var}(\overline X)=\frac{\sigma^2}{n},
\qquad
\mathbb{E}[S^2]=\sigma^2.\] If the sample is normal,
\[\overline X\sim \operatorname{Normal}\left(\mu,\frac{\sigma^2}{n}\right),
\qquad
\frac{(n-1)S^2}{\sigma^2}\sim \chi^2_{n-1},
\qquad
\frac{\overline X-\mu}{S/\sqrt n}\sim t_{n-1}.\] For continuous order
statistics, \[f_{X_{(j)}}(x)
=\frac{n!}{(j-1)!(n-j)!}f_X(x)[F_X(x)]^{j-1}[1-F_X(x)]^{n-j}.\]

\end{tcolorbox}

\begin{tcolorbox}[enhanced jigsaw, arc=.35mm, bottomrule=.15mm, bottomtitle=1mm, breakable, colback=white, colbacktitle=quarto-callout-tip-color!10!white, colframe=quarto-callout-tip-color-frame, coltitle=black, left=2mm, leftrule=.75mm, opacityback=0, opacitybacktitle=0.6, rightrule=.15mm, title=\textcolor{quarto-callout-tip-color}{\faLightbulb}\hspace{0.5em}{Conceptual summary}, titlerule=0mm, toprule=.15mm, toptitle=1mm]

\begin{itemize}
\item
  Random samples are iid random variables from a population
  distribution.
\item
  Statistics are functions of the random sample.
\item
  The sampling distribution describes how a statistic varies from sample
  to sample.
\item
  Normal samples produce especially simple and important sampling
  distributions.
\item
  Order statistics describe sorted data and are used for medians,
  quantiles, minima, maxima, and ranges.
\end{itemize}

\end{tcolorbox}

\section{References}\label{references}

These notes follow the course lecture material and the standard
references:

\begin{itemize}
\item
  Casella and Berger, \emph{Statistical Inference}, 2nd edition,
  Sections 5.1--5.4.
\item
  Larry Wasserman, \emph{All of Statistics}.
\item
  C. M. Grinstead and J. L. Snell, \emph{Introduction to Probability}.
\item
  Sheldon Ross, \emph{Introduction to Probability Models}.
\item
  Online resource: \url{https://www.probabilitycourse.com/}.
\end{itemize}

\bookmarksetup{startatroot}

\chapter{Chapter 9: Convergence Theory and Limit
Theorems}\label{chapter-9-convergence-theory-and-limit-theorems}

This chapter develops the mathematical language for describing the
limiting behavior of random variables, sample averages, and estimators.
It connects convergence modes with the weak law of large numbers, the
strong law of large numbers, the central limit theorem, continuous
mapping, Slutsky's theorem, and concentration inequalities.

\textbf{Topics.} Convergence of numbers and functions; sequences of
random variables; convergence in distribution, probability, mean, and
almost surely; relationships among modes of convergence; weak and strong
laws of large numbers; central limit theorem; continuous mapping
theorem; Slutsky's theorem; concentration inequalities and
high-dimensional proportion estimation.

\bookmarksetup{startatroot}

\chapter{Overview}\label{overview-6}

This section develops the language needed to describe limiting behavior
of random variables, sample averages, and estimators.

In deterministic calculus, a sequence of numbers may converge to a fixed
limit. In probability and statistics, the objects that converge are
often random variables. Since random variables are functions on an
underlying sample space, there are several useful ways for a sequence of
random variables to converge: by the shape of the distribution, by
probability of large deviation, by average error, or by
almost-everywhere sample-path behavior.

\begin{tcolorbox}[enhanced jigsaw, arc=.35mm, bottomrule=.15mm, bottomtitle=1mm, breakable, colback=white, colbacktitle=quarto-callout-tip-color!10!white, colframe=quarto-callout-tip-color-frame, coltitle=black, left=2mm, leftrule=.75mm, opacityback=0, opacitybacktitle=0.6, rightrule=.15mm, title=\textcolor{quarto-callout-tip-color}{\faLightbulb}\hspace{0.5em}{Key idea}, titlerule=0mm, toprule=.15mm, toptitle=1mm]

The main question is not only whether \(X_n\) approaches \(X\), but also
\emph{in what sense} it approaches \(X\): \[X_n \xrightarrow{D} X,\qquad
X_n \xrightarrow{P} X,\qquad
X_n \xrightarrow{L^r} X,\qquad
X_n \xrightarrow{a.s.} X.\] These four arrows encode different
mathematical meanings and are used for different statistical purposes.

\end{tcolorbox}

The major results in this section are the weak law of large numbers, the
strong law of large numbers, the central limit theorem, the continuous
mapping theorem, Slutsky's theorem, and concentration inequalities.

\bookmarksetup{startatroot}

\chapter{Convergence of Numbers and
Functions}\label{convergence-of-numbers-and-functions}

This section reviews ordinary deterministic convergence before
introducing convergence of random variables.

\section{Convergence of a sequence of
numbers}\label{convergence-of-a-sequence-of-numbers}

A sequence of numbers converges when, eventually, all terms remain
arbitrarily close to a fixed limiting value.

\begin{tcolorbox}[enhanced jigsaw, arc=.35mm, bottomrule=.15mm, bottomtitle=1mm, breakable, colback=white, colbacktitle=quarto-callout-note-color!10!white, colframe=quarto-callout-note-color-frame, coltitle=black, left=2mm, leftrule=.75mm, opacityback=0, opacitybacktitle=0.6, rightrule=.15mm, title=\textcolor{quarto-callout-note-color}{\faInfo}\hspace{0.5em}{Definition}, titlerule=0mm, toprule=.15mm, toptitle=1mm]

\textbf{Definition 1} (Convergence of a numerical sequence). A sequence
of numbers \(a_1,a_2,a_3,\ldots\) converges to a limit \(a\) if
\[\lim_{n\to\infty} a_n=a.\] More precisely, for every \(\epsilon>0\),
there exists \(N\in\mathbb{N}\) such that
\[|a_n-a|<\epsilon \qquad \text{for all } n>N.\]

\end{tcolorbox}

\begin{tcolorbox}[enhanced jigsaw, arc=.35mm, bottomrule=.15mm, bottomtitle=1mm, breakable, colback=white, colbacktitle=quarto-callout-tip-color!10!white, colframe=quarto-callout-tip-color-frame, coltitle=black, left=2mm, leftrule=.75mm, opacityback=0, opacitybacktitle=0.6, rightrule=.15mm, title=\textcolor{quarto-callout-tip-color}{\faLightbulb}\hspace{0.5em}{Example}, titlerule=0mm, toprule=.15mm, toptitle=1mm]

\textbf{Example 2} (A basic convergent sequence). Let
\[a_n=\frac{n}{n+1}.\] Show that \(a_n\to 1\).

\end{tcolorbox}

\begin{tcolorbox}[enhanced jigsaw, arc=.35mm, bottomrule=.15mm, bottomtitle=1mm, breakable, colback=white, colbacktitle=quarto-callout-note-color!10!white, colframe=quarto-callout-note-color-frame, coltitle=black, left=2mm, leftrule=.75mm, opacityback=0, opacitybacktitle=0.6, rightrule=.15mm, title=\textcolor{quarto-callout-note-color}{\faInfo}\hspace{0.5em}{Solution}, titlerule=0mm, toprule=.15mm, toptitle=1mm]

We have
\[\left|\frac{n}{n+1}-1\right|=\left|\frac{n-(n+1)}{n+1}\right|=\frac{1}{n+1}.\]
Given \(\epsilon>0\), choose \(N>1/\epsilon-1\). Then for all \(n>N\),
\[\frac{1}{n+1}<\epsilon.\] Therefore \(a_n\to 1\).

\end{tcolorbox}

\section{Epsilon-delta definition of a
limit}\label{epsilon-delta-definition-of-a-limit}

The epsilon-delta definition is the formal way to state that the values
of a function become close to a limit when the input becomes close to a
point.

\begin{tcolorbox}[enhanced jigsaw, arc=.35mm, bottomrule=.15mm, bottomtitle=1mm, breakable, colback=white, colbacktitle=quarto-callout-note-color!10!white, colframe=quarto-callout-note-color-frame, coltitle=black, left=2mm, leftrule=.75mm, opacityback=0, opacitybacktitle=0.6, rightrule=.15mm, title=\textcolor{quarto-callout-note-color}{\faInfo}\hspace{0.5em}{Definition}, titlerule=0mm, toprule=.15mm, toptitle=1mm]

\textbf{Definition 3} (Limit of a function). We say
\[\lim_{x\to a} f(x)=L\] if for every \(\epsilon>0\), there exists
\(\delta>0\) such that whenever \[0<|x-a|<\delta,\] it follows that
\[|f(x)-L|<\epsilon.\] Here \(\epsilon\) measures how close \(f(x)\)
must be to \(L\), and \(\delta\) measures how close \(x\) must be to
\(a\).

\end{tcolorbox}

\begin{tcolorbox}[enhanced jigsaw, arc=.35mm, bottomrule=.15mm, bottomtitle=1mm, breakable, colback=white, colbacktitle=quarto-callout-note-color!10!white, colframe=quarto-callout-note-color-frame, coltitle=black, left=2mm, leftrule=.75mm, opacityback=0, opacitybacktitle=0.6, rightrule=.15mm, title=\textcolor{quarto-callout-note-color}{\faInfo}\hspace{0.5em}{Remark}, titlerule=0mm, toprule=.15mm, toptitle=1mm]

\emph{Remark 4}. The same idea appears repeatedly in probability. In
convergence in probability, for example, the tolerance \(\epsilon\)
controls how far \(X_n\) can be from \(X\), and the probability of
violating that tolerance must go to zero.

\end{tcolorbox}

\begin{tcolorbox}[enhanced jigsaw, arc=.35mm, bottomrule=.15mm, bottomtitle=1mm, breakable, colback=white, colbacktitle=quarto-callout-warning-color!10!white, colframe=quarto-callout-warning-color-frame, coltitle=black, left=2mm, leftrule=.75mm, opacityback=0, opacitybacktitle=0.6, rightrule=.15mm, title=\textcolor{quarto-callout-warning-color}{\faExclamationTriangle}\hspace{0.5em}{Practice Problem}, titlerule=0mm, toprule=.15mm, toptitle=1mm]

\textbf{Practice Problem 5} (A quick limit check). Use the epsilon-delta
idea to show that \(f(x)=2x+1\) satisfies \[\lim_{x\to 3} f(x)=7.\]

\end{tcolorbox}

\begin{tcolorbox}[enhanced jigsaw, arc=.35mm, bottomrule=.15mm, bottomtitle=1mm, breakable, colback=white, colbacktitle=quarto-callout-note-color!10!white, colframe=quarto-callout-note-color-frame, coltitle=black, left=2mm, leftrule=.75mm, opacityback=0, opacitybacktitle=0.6, rightrule=.15mm, title=\textcolor{quarto-callout-note-color}{\faInfo}\hspace{0.5em}{Solution}, titlerule=0mm, toprule=.15mm, toptitle=1mm]

We compute \[|f(x)-7|=|2x+1-7|=2|x-3|.\] Given \(\epsilon>0\), choose
\(\delta=\epsilon/2\). If \(0<|x-3|<\delta\), then
\[|f(x)-7|=2|x-3|<2\delta=\epsilon.\] Therefore
\(\lim_{x\to 3}(2x+1)=7\).

\end{tcolorbox}

\bookmarksetup{startatroot}

\chapter{Sequences of Random
Variables}\label{sequences-of-random-variables}

This section explains why convergence of random variables is richer than
convergence of numbers.

Recall that a random variable is a function \[X:S\to \mathbb{R}\] from a
sample space \(S\) to the real numbers. If \(S=\{s_1,\ldots,s_k\}\) is
finite, then a random variable is determined by its values
\[X(s_i)=x_i,\qquad i=1,
\ldots,k.\] A sequence of random variables \[X_1,X_2,X_3,\ldots\] is
therefore a sequence of functions \[X_n:S\to \mathbb{R}.\] For each
outcome \(s\in S\), the sequence \(X_1(s),X_2(s),\ldots\) is an ordinary
sequence of numbers. But random variables also have distributions,
probabilities, expectations, and variances. This creates several
different modes of convergence.

\begin{tcolorbox}[enhanced jigsaw, arc=.35mm, bottomrule=.15mm, bottomtitle=1mm, breakable, colback=white, colbacktitle=quarto-callout-tip-color!10!white, colframe=quarto-callout-tip-color-frame, coltitle=black, left=2mm, leftrule=.75mm, opacityback=0, opacitybacktitle=0.6, rightrule=.15mm, title=\textcolor{quarto-callout-tip-color}{\faLightbulb}\hspace{0.5em}{Example}, titlerule=0mm, toprule=.15mm, toptitle=1mm]

\textbf{Example 6} (A sequence of random variables on a finite sample
space). Let \(S=\{s_1,s_2\}\) with
\(\mathbb{P}(s_1)=\mathbb{P}(s_2)=1/2\), and define
\[X_n(s_1)=\frac{1}{n},\qquad X_n(s_2)=1+\frac{1}{n}.\] Find the
pointwise limit \(X(s)=\lim_{n\to\infty}X_n(s)\).

\end{tcolorbox}

\begin{tcolorbox}[enhanced jigsaw, arc=.35mm, bottomrule=.15mm, bottomtitle=1mm, breakable, colback=white, colbacktitle=quarto-callout-note-color!10!white, colframe=quarto-callout-note-color-frame, coltitle=black, left=2mm, leftrule=.75mm, opacityback=0, opacitybacktitle=0.6, rightrule=.15mm, title=\textcolor{quarto-callout-note-color}{\faInfo}\hspace{0.5em}{Solution}, titlerule=0mm, toprule=.15mm, toptitle=1mm]

For \(s_1\), \[X_n(s_1)=\frac{1}{n}\to 0.\] For \(s_2\),
\[X_n(s_2)=1+\frac{1}{n}\to 1.\] Thus the pointwise limiting random
variable is \[X(s_1)=0,
\qquad
X(s_2)=1.\]

\end{tcolorbox}

\bookmarksetup{startatroot}

\chapter{Convergence in Distribution}\label{convergence-in-distribution}

This section introduces the weakest and most distribution-centered mode
of convergence.

\begin{tcolorbox}[enhanced jigsaw, arc=.35mm, bottomrule=.15mm, bottomtitle=1mm, breakable, colback=white, colbacktitle=quarto-callout-note-color!10!white, colframe=quarto-callout-note-color-frame, coltitle=black, left=2mm, leftrule=.75mm, opacityback=0, opacitybacktitle=0.6, rightrule=.15mm, title=\textcolor{quarto-callout-note-color}{\faInfo}\hspace{0.5em}{Definition}, titlerule=0mm, toprule=.15mm, toptitle=1mm]

\textbf{Definition 7} (Convergence in distribution). Let \(X_n\) have
CDF \(F_n\) and let \(X\) have CDF \(F\). We say that \(X_n\)
\textbf{converges in distribution} to \(X\), also called convergence
weakly or convergence in law, if \[\lim_{n\to\infty}F_n(x)=F(x)\] for
every continuity point \(x\) of \(F\). We write
\[X_n \xrightarrow{D} X.\]

\end{tcolorbox}

\begin{tcolorbox}[enhanced jigsaw, arc=.35mm, bottomrule=.15mm, bottomtitle=1mm, breakable, colback=white, colbacktitle=quarto-callout-tip-color!10!white, colframe=quarto-callout-tip-color-frame, coltitle=black, left=2mm, leftrule=.75mm, opacityback=0, opacitybacktitle=0.6, rightrule=.15mm, title=\textcolor{quarto-callout-tip-color}{\faLightbulb}\hspace{0.5em}{Key idea}, titlerule=0mm, toprule=.15mm, toptitle=1mm]

Convergence in distribution means that the overall shape of the
distribution of \(X_n\) converges to the overall shape of the
distribution of \(X\). It does not necessarily mean that \(X_n\) and
\(X\) are close on the same sample space.

\end{tcolorbox}

\begin{tcolorbox}[enhanced jigsaw, arc=.35mm, bottomrule=.15mm, bottomtitle=1mm, breakable, colback=white, colbacktitle=quarto-callout-tip-color!10!white, colframe=quarto-callout-tip-color-frame, coltitle=black, left=2mm, leftrule=.75mm, opacityback=0, opacitybacktitle=0.6, rightrule=.15mm, title=\textcolor{quarto-callout-tip-color}{\faLightbulb}\hspace{0.5em}{Example}, titlerule=0mm, toprule=.15mm, toptitle=1mm]

\textbf{Example 8} (Normal distributions with variances converging to
one). Let \[X_n\sim \operatorname{Normal}\left(0,1+\frac{1}{n}\right).\]
Show that \[X_n\xrightarrow{D}\operatorname{Normal}(0,1).\]

\end{tcolorbox}

\begin{tcolorbox}[enhanced jigsaw, arc=.35mm, bottomrule=.15mm, bottomtitle=1mm, breakable, colback=white, colbacktitle=quarto-callout-note-color!10!white, colframe=quarto-callout-note-color-frame, coltitle=black, left=2mm, leftrule=.75mm, opacityback=0, opacitybacktitle=0.6, rightrule=.15mm, title=\textcolor{quarto-callout-note-color}{\faInfo}\hspace{0.5em}{Solution}, titlerule=0mm, toprule=.15mm, toptitle=1mm]

The CDF of \(X_n\) is
\[F_n(x)=\Phi\left(\frac{x}{\sqrt{1+1/n}}\right),\] where \(\Phi\) is
the standard normal CDF. Since \[\frac{x}{\sqrt{1+1/n}}\to x\] and
\(\Phi\) is continuous,
\[F_n(x)=\Phi\left(\frac{x}{\sqrt{1+1/n}}\right)\to \Phi(x).\] Therefore
\(X_n\xrightarrow{D}N(0,1)\).

\end{tcolorbox}

\begin{figure}

\centering{

\pandocbounded{\includegraphics[keepaspectratio]{chapters/chapter-09-convergence-theory_files/figure-pdf/fig-normal-convergence-ch9-output-1.pdf}}

}

\caption{\label{fig-normal-convergence-ch9}Densities N(0,2), N(0,1.2),
and N(0,1) illustrating convergence in distribution.}

\end{figure}%

\begin{tcolorbox}[enhanced jigsaw, arc=.35mm, bottomrule=.15mm, bottomtitle=1mm, breakable, colback=white, colbacktitle=quarto-callout-warning-color!10!white, colframe=quarto-callout-warning-color-frame, coltitle=black, left=2mm, leftrule=.75mm, opacityback=0, opacitybacktitle=0.6, rightrule=.15mm, title=\textcolor{quarto-callout-warning-color}{\faExclamationTriangle}\hspace{0.5em}{Practice Problem}, titlerule=0mm, toprule=.15mm, toptitle=1mm]

\textbf{Practice Problem 9} (CDF convergence). Let
\(X_n\sim \operatorname{Uniform}(0,1+1/n)\). Show that
\(X_n\xrightarrow{D}\operatorname{Uniform}(0,1)\).

\end{tcolorbox}

\begin{tcolorbox}[enhanced jigsaw, arc=.35mm, bottomrule=.15mm, bottomtitle=1mm, breakable, colback=white, colbacktitle=quarto-callout-note-color!10!white, colframe=quarto-callout-note-color-frame, coltitle=black, left=2mm, leftrule=.75mm, opacityback=0, opacitybacktitle=0.6, rightrule=.15mm, title=\textcolor{quarto-callout-note-color}{\faInfo}\hspace{0.5em}{Solution}, titlerule=0mm, toprule=.15mm, toptitle=1mm]

For \(0<x<1\), \[F_n(x)=\frac{x}{1+1/n}\to x.\] For \(x\le 0\),
\(F_n(x)=0\to 0\). For \(x>1\), eventually \(x>1+1/n\), so
\(F_n(x)=1\to 1\). These are exactly the values of the CDF of
\(\operatorname{Uniform}(0,1)\) at its continuity points. Therefore
\(X_n\xrightarrow{D}\operatorname{Uniform}(0,1)\).

\end{tcolorbox}

\bookmarksetup{startatroot}

\chapter{Convergence in Probability}\label{convergence-in-probability}

This section introduces the mode of convergence most closely connected
to statistical consistency.

\begin{tcolorbox}[enhanced jigsaw, arc=.35mm, bottomrule=.15mm, bottomtitle=1mm, breakable, colback=white, colbacktitle=quarto-callout-note-color!10!white, colframe=quarto-callout-note-color-frame, coltitle=black, left=2mm, leftrule=.75mm, opacityback=0, opacitybacktitle=0.6, rightrule=.15mm, title=\textcolor{quarto-callout-note-color}{\faInfo}\hspace{0.5em}{Definition}, titlerule=0mm, toprule=.15mm, toptitle=1mm]

\textbf{Definition 10} (Convergence in probability). A sequence \(X_n\)
converges in probability to \(X\) if for every \(\epsilon>0\),
\[\lim_{n\to\infty}\mathbb{P}(|X_n-X|>\epsilon)=0.\] We write
\[X_n\xrightarrow{P}X.\]

\end{tcolorbox}

\begin{tcolorbox}[enhanced jigsaw, arc=.35mm, bottomrule=.15mm, bottomtitle=1mm, breakable, colback=white, colbacktitle=quarto-callout-tip-color!10!white, colframe=quarto-callout-tip-color-frame, coltitle=black, left=2mm, leftrule=.75mm, opacityback=0, opacitybacktitle=0.6, rightrule=.15mm, title=\textcolor{quarto-callout-tip-color}{\faLightbulb}\hspace{0.5em}{Key idea}, titlerule=0mm, toprule=.15mm, toptitle=1mm]

Convergence in probability means that large deviations from the limiting
random variable become rare.

\end{tcolorbox}

\begin{tcolorbox}[enhanced jigsaw, arc=.35mm, bottomrule=.15mm, bottomtitle=1mm, breakable, colback=white, colbacktitle=quarto-callout-tip-color!10!white, colframe=quarto-callout-tip-color-frame, coltitle=black, left=2mm, leftrule=.75mm, opacityback=0, opacitybacktitle=0.6, rightrule=.15mm, title=\textcolor{quarto-callout-tip-color}{\faLightbulb}\hspace{0.5em}{Example}, titlerule=0mm, toprule=.15mm, toptitle=1mm]

\textbf{Example 11} (A shrinking normal sequence). Let
\[X_n\sim \operatorname{Normal}\left(0,\frac{1}{n}\right).\] Show that
\(X_n\xrightarrow{P}0\), and identify the distributional limit of
\(\sqrt nX_n\).

\end{tcolorbox}

\begin{tcolorbox}[enhanced jigsaw, arc=.35mm, bottomrule=.15mm, bottomtitle=1mm, breakable, colback=white, colbacktitle=quarto-callout-note-color!10!white, colframe=quarto-callout-note-color-frame, coltitle=black, left=2mm, leftrule=.75mm, opacityback=0, opacitybacktitle=0.6, rightrule=.15mm, title=\textcolor{quarto-callout-note-color}{\faInfo}\hspace{0.5em}{Solution}, titlerule=0mm, toprule=.15mm, toptitle=1mm]

For any \(\epsilon>0\), Chebyshev's inequality gives
\[\mathbb{P}(|X_n-0|>\epsilon)
\le \frac{\operatorname{Var}(X_n)}{\epsilon^2}
=\frac{1}{n\epsilon^2}\to 0.\] Thus \(X_n\xrightarrow{P}0\).

Also, since \(X_n\sim N(0,1/n)\), \[\sqrt n X_n\sim N(0,1).\] Therefore
\[\sqrt nX_n\xrightarrow{D}N(0,1).\] This example shows that a sequence
can converge to \(0\) in probability while a rescaled version has a
non-degenerate limiting distribution.

\end{tcolorbox}

\begin{tcolorbox}[enhanced jigsaw, arc=.35mm, bottomrule=.15mm, bottomtitle=1mm, breakable, colback=white, colbacktitle=quarto-callout-tip-color!10!white, colframe=quarto-callout-tip-color-frame, coltitle=black, left=2mm, leftrule=.75mm, opacityback=0, opacitybacktitle=0.6, rightrule=.15mm, title=\textcolor{quarto-callout-tip-color}{\faLightbulb}\hspace{0.5em}{Example}, titlerule=0mm, toprule=.15mm, toptitle=1mm]

\textbf{Example 12} (Binary sequence converging in probability). Let
\[\mathbb{P}(X_n=0)=1-\frac1n,
\qquad
\mathbb{P}(X_n=1)=\frac1n.\] Show that \(X_n\xrightarrow{P}0\).

\end{tcolorbox}

\begin{tcolorbox}[enhanced jigsaw, arc=.35mm, bottomrule=.15mm, bottomtitle=1mm, breakable, colback=white, colbacktitle=quarto-callout-note-color!10!white, colframe=quarto-callout-note-color-frame, coltitle=black, left=2mm, leftrule=.75mm, opacityback=0, opacitybacktitle=0.6, rightrule=.15mm, title=\textcolor{quarto-callout-note-color}{\faInfo}\hspace{0.5em}{Solution}, titlerule=0mm, toprule=.15mm, toptitle=1mm]

Let \(0<\epsilon<1\). Then
\[\mathbb{P}(|X_n-0|>\epsilon)=\mathbb{P}(X_n=1)=\frac1n\to 0.\] For
\(\epsilon\ge 1\), the probability is already \(0\). Thus
\(X_n\xrightarrow{P}0\).

\end{tcolorbox}

\section{Convergence in probability implies convergence in
distribution}\label{convergence-in-probability-implies-convergence-in-distribution}

This subsection states a key implication and shows why the converse is
not true in general.

\begin{tcolorbox}[enhanced jigsaw, arc=.35mm, bottomrule=.15mm, bottomtitle=1mm, breakable, colback=white, colbacktitle=quarto-callout-important-color!10!white, colframe=quarto-callout-important-color-frame, coltitle=black, left=2mm, leftrule=.75mm, opacityback=0, opacitybacktitle=0.6, rightrule=.15mm, title=\textcolor{quarto-callout-important-color}{\faExclamation}\hspace{0.5em}{Theorem}, titlerule=0mm, toprule=.15mm, toptitle=1mm]

\textbf{Theorem 13} (Convergence in probability implies convergence in
distribution). \emph{If \(X_n\xrightarrow{P}X\), then
\[X_n\xrightarrow{D}X.\]}

\end{tcolorbox}

\begin{tcolorbox}[enhanced jigsaw, arc=.35mm, bottomrule=.15mm, bottomtitle=1mm, breakable, colback=white, colbacktitle=quarto-callout-tip-color!10!white, colframe=quarto-callout-tip-color-frame, coltitle=black, left=2mm, leftrule=.75mm, opacityback=0, opacitybacktitle=0.6, rightrule=.15mm, title=\textcolor{quarto-callout-tip-color}{\faLightbulb}\hspace{0.5em}{Example}, titlerule=0mm, toprule=.15mm, toptitle=1mm]

\textbf{Example 14} (Same distribution does not imply convergence in
probability). Let \(X\sim N(0,1)\) and define
\[X_n=-X\qquad \text{for all } n.\] Then \(X_n\xrightarrow{D}X\), but
\(X_n\) does not converge in probability to \(X\).

\end{tcolorbox}

\begin{tcolorbox}[enhanced jigsaw, arc=.35mm, bottomrule=.15mm, bottomtitle=1mm, breakable, colback=white, colbacktitle=quarto-callout-note-color!10!white, colframe=quarto-callout-note-color-frame, coltitle=black, left=2mm, leftrule=.75mm, opacityback=0, opacitybacktitle=0.6, rightrule=.15mm, title=\textcolor{quarto-callout-note-color}{\faInfo}\hspace{0.5em}{Solution}, titlerule=0mm, toprule=.15mm, toptitle=1mm]

Since \(X\) and \(-X\) have the same standard normal distribution, every
\(X_n=-X\) has the same distribution as \(X\). Therefore
\(X_n\xrightarrow{D}X\).

However, \[|X_n-X|=|-X-X|=2|X|.\] Hence
\[\mathbb{P}(|X_n-X|>1)=\mathbb{P}(2|X|>1)=\mathbb{P}(|X|>1/2),\] which
is a positive constant and does not go to \(0\). Therefore \(X_n\) does
not converge in probability to \(X\).

\end{tcolorbox}

\begin{tcolorbox}[enhanced jigsaw, arc=.35mm, bottomrule=.15mm, bottomtitle=1mm, breakable, colback=white, colbacktitle=quarto-callout-important-color!10!white, colframe=quarto-callout-important-color-frame, coltitle=black, left=2mm, leftrule=.75mm, opacityback=0, opacitybacktitle=0.6, rightrule=.15mm, title=\textcolor{quarto-callout-important-color}{\faExclamation}\hspace{0.5em}{Theorem}, titlerule=0mm, toprule=.15mm, toptitle=1mm]

\textbf{Theorem 15} (Distributional convergence to a constant). \emph{If
\(X_n\xrightarrow{D}c\), where \(c\) is a constant, then
\[X_n\xrightarrow{P}c.\]}

\end{tcolorbox}

\begin{tcolorbox}[enhanced jigsaw, arc=.35mm, bottomrule=.15mm, bottomtitle=1mm, breakable, colback=white, colbacktitle=quarto-callout-note-color!10!white, colframe=quarto-callout-note-color-frame, coltitle=black, left=2mm, leftrule=.75mm, opacityback=0, opacitybacktitle=0.6, rightrule=.15mm, title=\textcolor{quarto-callout-note-color}{\faInfo}\hspace{0.5em}{Remark}, titlerule=0mm, toprule=.15mm, toptitle=1mm]

\emph{Remark 16}. The theorem is special to the case where the limiting
random variable is constant. When the limit is random, convergence in
distribution alone does not imply convergence in probability.

\end{tcolorbox}

\bookmarksetup{startatroot}

\chapter{Convergence in Mean}\label{convergence-in-mean}

This section introduces convergence based on the expected size of the
error.

\begin{tcolorbox}[enhanced jigsaw, arc=.35mm, bottomrule=.15mm, bottomtitle=1mm, breakable, colback=white, colbacktitle=quarto-callout-note-color!10!white, colframe=quarto-callout-note-color-frame, coltitle=black, left=2mm, leftrule=.75mm, opacityback=0, opacitybacktitle=0.6, rightrule=.15mm, title=\textcolor{quarto-callout-note-color}{\faInfo}\hspace{0.5em}{Definition}, titlerule=0mm, toprule=.15mm, toptitle=1mm]

\textbf{Definition 17} (Convergence in \(r\)-th mean). Let \(r\ge 1\) be
fixed. We say that \(X_n\) converges in the \(r\)-th mean to \(X\) if
\[\lim_{n\to\infty}\mathbb{E}|X_n-X|^r=0.\] We write
\[X_n\xrightarrow{L^r}X.\] The case \(r=1\) is convergence in mean, and
the case \(r=2\) is convergence in mean square.

\end{tcolorbox}

\begin{tcolorbox}[enhanced jigsaw, arc=.35mm, bottomrule=.15mm, bottomtitle=1mm, breakable, colback=white, colbacktitle=quarto-callout-tip-color!10!white, colframe=quarto-callout-tip-color-frame, coltitle=black, left=2mm, leftrule=.75mm, opacityback=0, opacitybacktitle=0.6, rightrule=.15mm, title=\textcolor{quarto-callout-tip-color}{\faLightbulb}\hspace{0.5em}{Example}, titlerule=0mm, toprule=.15mm, toptitle=1mm]

\textbf{Example 18} (Uniform sequence converging in \(L^r\)). Let
\[X_n\sim \operatorname{Uniform}\left(0,\frac1n\right).\] Show that
\(X_n\xrightarrow{L^r}0\) for every fixed \(r\ge 1\).

\end{tcolorbox}

\begin{tcolorbox}[enhanced jigsaw, arc=.35mm, bottomrule=.15mm, bottomtitle=1mm, breakable, colback=white, colbacktitle=quarto-callout-note-color!10!white, colframe=quarto-callout-note-color-frame, coltitle=black, left=2mm, leftrule=.75mm, opacityback=0, opacitybacktitle=0.6, rightrule=.15mm, title=\textcolor{quarto-callout-note-color}{\faInfo}\hspace{0.5em}{Solution}, titlerule=0mm, toprule=.15mm, toptitle=1mm]

The density of \(X_n\) is \(n\) on \((0,1/n)\). Thus
\[\mathbb{E}|X_n-0|^r
=\int_0^{1/n} x^r n\,dx
= n\cdot \frac{(1/n)^{r+1}}{r+1}
=\frac{1}{(r+1)n^r}\to 0.\] Therefore \(X_n\xrightarrow{L^r}0\).

\end{tcolorbox}

\section{Mean convergence implies probability
convergence}\label{mean-convergence-implies-probability-convergence}

This subsection uses Markov's inequality to connect convergence in mean
with convergence in probability.

\begin{tcolorbox}[enhanced jigsaw, arc=.35mm, bottomrule=.15mm, bottomtitle=1mm, breakable, colback=white, colbacktitle=quarto-callout-important-color!10!white, colframe=quarto-callout-important-color-frame, coltitle=black, left=2mm, leftrule=.75mm, opacityback=0, opacitybacktitle=0.6, rightrule=.15mm, title=\textcolor{quarto-callout-important-color}{\faExclamation}\hspace{0.5em}{Theorem}, titlerule=0mm, toprule=.15mm, toptitle=1mm]

\textbf{Theorem 19} (\(L^r\) convergence implies convergence in
probability). \emph{If \(X_n\xrightarrow{L^r}X\) for some \(r\ge 1\),
then \[X_n\xrightarrow{P}X.\]}

\end{tcolorbox}

\begin{tcolorbox}[enhanced jigsaw, arc=.35mm, bottomrule=.15mm, bottomtitle=1mm, breakable, colback=white, colbacktitle=quarto-callout-note-color!10!white, colframe=quarto-callout-note-color-frame, coltitle=black, left=2mm, leftrule=.75mm, opacityback=0, opacitybacktitle=0.6, rightrule=.15mm, title=\textcolor{quarto-callout-note-color}{\faInfo}\hspace{0.5em}{Proof}, titlerule=0mm, toprule=.15mm, toptitle=1mm]

\emph{Proof.} For any \(\epsilon>0\), Markov's inequality applied to
\(|X_n-X|^r\) gives \[\mathbb{P}(|X_n-X|>\epsilon)
=\mathbb{P}(|X_n-X|^r>\epsilon^r)
\le \frac{\mathbb{E}|X_n-X|^r}{\epsilon^r}.\] If
\(\mathbb{E}|X_n-X|^r\to 0\), then the right-hand side goes to \(0\).
Hence \(X_n\xrightarrow{P}X\). ◻

\end{tcolorbox}

\begin{tcolorbox}[enhanced jigsaw, arc=.35mm, bottomrule=.15mm, bottomtitle=1mm, breakable, colback=white, colbacktitle=quarto-callout-tip-color!10!white, colframe=quarto-callout-tip-color-frame, coltitle=black, left=2mm, leftrule=.75mm, opacityback=0, opacitybacktitle=0.6, rightrule=.15mm, title=\textcolor{quarto-callout-tip-color}{\faLightbulb}\hspace{0.5em}{Example}, titlerule=0mm, toprule=.15mm, toptitle=1mm]

\textbf{Example 20} (Convergence in probability without convergence in
mean). Define \(X_n\) by \[\mathbb{P}(X_n=0)=1-\frac1n,
\qquad
\mathbb{P}(X_n=n^2)=\frac1n.\] Show that \(X_n\xrightarrow{P}0\) but
\(\mathbb{E}X_n\to\infty\).

\end{tcolorbox}

\begin{tcolorbox}[enhanced jigsaw, arc=.35mm, bottomrule=.15mm, bottomtitle=1mm, breakable, colback=white, colbacktitle=quarto-callout-note-color!10!white, colframe=quarto-callout-note-color-frame, coltitle=black, left=2mm, leftrule=.75mm, opacityback=0, opacitybacktitle=0.6, rightrule=.15mm, title=\textcolor{quarto-callout-note-color}{\faInfo}\hspace{0.5em}{Solution}, titlerule=0mm, toprule=.15mm, toptitle=1mm]

For any fixed \(\epsilon>0\), once \(n^2>\epsilon\),
\[\mathbb{P}(|X_n-0|>\epsilon)=\mathbb{P}(X_n=n^2)=\frac1n\to 0.\] Thus
\(X_n\xrightarrow{P}0\).

However,
\[\mathbb{E}X_n=0\left(1-\frac1n\right)+n^2\cdot \frac1n=n\to\infty.\]
Therefore convergence in probability does not imply convergence in mean.

\end{tcolorbox}

\bookmarksetup{startatroot}

\chapter{Almost Sure Convergence}\label{almost-sure-convergence}

This section gives the strongest sample-path mode of convergence used in
this course.

\begin{tcolorbox}[enhanced jigsaw, arc=.35mm, bottomrule=.15mm, bottomtitle=1mm, breakable, colback=white, colbacktitle=quarto-callout-note-color!10!white, colframe=quarto-callout-note-color-frame, coltitle=black, left=2mm, leftrule=.75mm, opacityback=0, opacitybacktitle=0.6, rightrule=.15mm, title=\textcolor{quarto-callout-note-color}{\faInfo}\hspace{0.5em}{Definition}, titlerule=0mm, toprule=.15mm, toptitle=1mm]

\textbf{Definition 21} (Almost sure convergence). The sequence \(X_n\)
converges almost surely to \(X\) if
\[\mathbb{P}\left(\lim_{n\to\infty}X_n=X\right)=1.\] Equivalently,
\[\mathbb{P}\left(\omega:\lim_{n\to\infty}X_n(\omega)=X(\omega)\right)=1.\]
We write \[X_n\xrightarrow{a.s.}X.\]

\end{tcolorbox}

\begin{tcolorbox}[enhanced jigsaw, arc=.35mm, bottomrule=.15mm, bottomtitle=1mm, breakable, colback=white, colbacktitle=quarto-callout-tip-color!10!white, colframe=quarto-callout-tip-color-frame, coltitle=black, left=2mm, leftrule=.75mm, opacityback=0, opacitybacktitle=0.6, rightrule=.15mm, title=\textcolor{quarto-callout-tip-color}{\faLightbulb}\hspace{0.5em}{Key idea}, titlerule=0mm, toprule=.15mm, toptitle=1mm]

Interpretation Almost sure convergence means that, with probability
\(1\), the actual sample path \(X_n(\omega)\) converges pointwise to
\(X(\omega)\).

\end{tcolorbox}

\begin{tcolorbox}[enhanced jigsaw, arc=.35mm, bottomrule=.15mm, bottomtitle=1mm, breakable, colback=white, colbacktitle=quarto-callout-important-color!10!white, colframe=quarto-callout-important-color-frame, coltitle=black, left=2mm, leftrule=.75mm, opacityback=0, opacitybacktitle=0.6, rightrule=.15mm, title=\textcolor{quarto-callout-important-color}{\faExclamation}\hspace{0.5em}{Theorem}, titlerule=0mm, toprule=.15mm, toptitle=1mm]

\textbf{Theorem 22} (Almost sure convergence implies convergence in
probability). \emph{If \(X_n\xrightarrow{a.s.}X\), then
\[X_n\xrightarrow{P}X.\]}

\end{tcolorbox}

\begin{tcolorbox}[enhanced jigsaw, arc=.35mm, bottomrule=.15mm, bottomtitle=1mm, breakable, colback=white, colbacktitle=quarto-callout-note-color!10!white, colframe=quarto-callout-note-color-frame, coltitle=black, left=2mm, leftrule=.75mm, opacityback=0, opacitybacktitle=0.6, rightrule=.15mm, title=\textcolor{quarto-callout-note-color}{\faInfo}\hspace{0.5em}{Proof}, titlerule=0mm, toprule=.15mm, toptitle=1mm]

\emph{Reason.} Almost sure convergence says that for almost every
\(\omega\), the numerical sequence \(X_n(\omega)\) converges to
\(X(\omega)\). Therefore, for any \(\epsilon>0\), the event
\[A_{n,\epsilon}=\{\omega: |X_n(\omega)-X(\omega)|>\epsilon\}\] must
eventually become rare. This implies
\[\mathbb{P}(A_{n,\epsilon})\to 0,\] which is precisely convergence in
probability. ◻

\end{tcolorbox}

\includegraphics[width=5.83in,height=2.9in]{chapters/chapter-09-convergence-theory_files/figure-latex/mermaid-figure-1.png}

\begin{tcolorbox}[enhanced jigsaw, arc=.35mm, bottomrule=.15mm, bottomtitle=1mm, breakable, colback=white, colbacktitle=quarto-callout-warning-color!10!white, colframe=quarto-callout-warning-color-frame, coltitle=black, left=2mm, leftrule=.75mm, opacityback=0, opacitybacktitle=0.6, rightrule=.15mm, title=\textcolor{quarto-callout-warning-color}{\faExclamationTriangle}\hspace{0.5em}{Common mistake}, titlerule=0mm, toprule=.15mm, toptitle=1mm]

Common mistake The arrows above are one-way in general. For example,
convergence in distribution does not usually imply convergence in
probability, and convergence in probability does not usually imply
convergence in mean.

\end{tcolorbox}

\begin{tcolorbox}[enhanced jigsaw, arc=.35mm, bottomrule=.15mm, bottomtitle=1mm, breakable, colback=white, colbacktitle=quarto-callout-warning-color!10!white, colframe=quarto-callout-warning-color-frame, coltitle=black, left=2mm, leftrule=.75mm, opacityback=0, opacitybacktitle=0.6, rightrule=.15mm, title=\textcolor{quarto-callout-warning-color}{\faExclamationTriangle}\hspace{0.5em}{Practice Problem}, titlerule=0mm, toprule=.15mm, toptitle=1mm]

\textbf{Practice Problem 23} (Identifying modes of convergence). Suppose
\(X_n=Z/n\), where \(Z\) is a random variable with
\(\mathbb{E}Z^2<\infty\). Show that \(X_n\to 0\) almost surely and in
\(L^2\).

\end{tcolorbox}

\begin{tcolorbox}[enhanced jigsaw, arc=.35mm, bottomrule=.15mm, bottomtitle=1mm, breakable, colback=white, colbacktitle=quarto-callout-note-color!10!white, colframe=quarto-callout-note-color-frame, coltitle=black, left=2mm, leftrule=.75mm, opacityback=0, opacitybacktitle=0.6, rightrule=.15mm, title=\textcolor{quarto-callout-note-color}{\faInfo}\hspace{0.5em}{Solution}, titlerule=0mm, toprule=.15mm, toptitle=1mm]

For every outcome \(\omega\) such that \(Z(\omega)\) is finite,
\[X_n(\omega)=\frac{Z(\omega)}{n}\to 0.\] Since \(Z\) is finite almost
surely, \(X_n\xrightarrow{a.s.}0\).

Also,
\[\mathbb{E}|X_n|^2=\mathbb{E}\left(\frac{Z^2}{n^2}\right)=\frac{\mathbb{E}Z^2}{n^2}\to 0.\]
Therefore \(X_n\xrightarrow{L^2}0\).

\end{tcolorbox}

\bookmarksetup{startatroot}

\chapter{Weak Law of Large Numbers}\label{weak-law-of-large-numbers}

This section states and proves the law saying that sample averages
converge to the population mean in probability.

Let \(X_1,\ldots,X_n\) be iid random variables from a distribution
\(F\). The sample average is \[\overline X_n=\frac1n\sum_{i=1}^n X_i.\]

\begin{tcolorbox}[enhanced jigsaw, arc=.35mm, bottomrule=.15mm, bottomtitle=1mm, breakable, colback=white, colbacktitle=quarto-callout-important-color!10!white, colframe=quarto-callout-important-color-frame, coltitle=black, left=2mm, leftrule=.75mm, opacityback=0, opacitybacktitle=0.6, rightrule=.15mm, title=\textcolor{quarto-callout-important-color}{\faExclamation}\hspace{0.5em}{Theorem}, titlerule=0mm, toprule=.15mm, toptitle=1mm]

\textbf{Theorem 24} (Weak law of large numbers). \emph{Suppose
\(X_1,X_2,\ldots\) are iid with mean \[\mu=\mathbb{E}X_i<\infty\] and
finite variance \[\operatorname{Var}(X_i)=\sigma^2<\infty.\] Then
\[\overline X_n\xrightarrow{P}\mu.\] That is, for every \(\epsilon>0\),
\[\lim_{n\to\infty}\mathbb{P}(|\overline X_n-\mu|>\epsilon)=0.\]}

\end{tcolorbox}

\begin{tcolorbox}[enhanced jigsaw, arc=.35mm, bottomrule=.15mm, bottomtitle=1mm, breakable, colback=white, colbacktitle=quarto-callout-note-color!10!white, colframe=quarto-callout-note-color-frame, coltitle=black, left=2mm, leftrule=.75mm, opacityback=0, opacitybacktitle=0.6, rightrule=.15mm, title=\textcolor{quarto-callout-note-color}{\faInfo}\hspace{0.5em}{Proof}, titlerule=0mm, toprule=.15mm, toptitle=1mm]

\emph{Proof.} We know \[\mathbb{E}\overline X_n=\mu,
\qquad
\operatorname{Var}(\overline X_n)=\frac{\sigma^2}{n}.\] By Chebyshev's
inequality, \[\mathbb{P}(|\overline X_n-\mu|\ge \epsilon)
\le \frac{\operatorname{Var}(\overline X_n)}{\epsilon^2}
=\frac{\sigma^2}{n\epsilon^2}\to 0.\] Therefore
\(\overline X_n\xrightarrow{P}\mu\). ◻

\end{tcolorbox}

\begin{tcolorbox}[enhanced jigsaw, arc=.35mm, bottomrule=.15mm, bottomtitle=1mm, breakable, colback=white, colbacktitle=quarto-callout-tip-color!10!white, colframe=quarto-callout-tip-color-frame, coltitle=black, left=2mm, leftrule=.75mm, opacityback=0, opacitybacktitle=0.6, rightrule=.15mm, title=\textcolor{quarto-callout-tip-color}{\faLightbulb}\hspace{0.5em}{Example}, titlerule=0mm, toprule=.15mm, toptitle=1mm]

\textbf{Example 25} (Coin flips). Let
\(X_i\sim \operatorname{Bernoulli}(p)\) be iid, where \(X_i=1\) means
heads and \(X_i=0\) means tails. Let
\[\overline X_n=\frac1n\sum_{i=1}^n X_i.\] Interpret the weak law of
large numbers.

\end{tcolorbox}

\begin{tcolorbox}[enhanced jigsaw, arc=.35mm, bottomrule=.15mm, bottomtitle=1mm, breakable, colback=white, colbacktitle=quarto-callout-note-color!10!white, colframe=quarto-callout-note-color-frame, coltitle=black, left=2mm, leftrule=.75mm, opacityback=0, opacitybacktitle=0.6, rightrule=.15mm, title=\textcolor{quarto-callout-note-color}{\faInfo}\hspace{0.5em}{Solution}, titlerule=0mm, toprule=.15mm, toptitle=1mm]

Here \(\mathbb{E}X_i=p\) and \(\operatorname{Var}(X_i)=p(1-p)\). The
WLLN gives \[\overline X_n\xrightarrow{P}p.\] The sample proportion of
heads converges in probability to the true probability of heads.

\end{tcolorbox}

\bookmarksetup{startatroot}

\chapter{Strong Law of Large Numbers}\label{strong-law-of-large-numbers}

This section states the stronger sample-path version of the law of large
numbers.

\begin{tcolorbox}[enhanced jigsaw, arc=.35mm, bottomrule=.15mm, bottomtitle=1mm, breakable, colback=white, colbacktitle=quarto-callout-important-color!10!white, colframe=quarto-callout-important-color-frame, coltitle=black, left=2mm, leftrule=.75mm, opacityback=0, opacitybacktitle=0.6, rightrule=.15mm, title=\textcolor{quarto-callout-important-color}{\faExclamation}\hspace{0.5em}{Theorem}, titlerule=0mm, toprule=.15mm, toptitle=1mm]

\textbf{Theorem 26} (Strong law of large numbers). \emph{Suppose
\(X_1,X_2,\ldots\) are iid with finite mean \[\mu=\mathbb{E}X_i<\infty\]
and finite variance \[\operatorname{Var}(X_i)=\sigma^2<\infty.\] Then
\[\overline X_n\xrightarrow{a.s.}\mu.\] Equivalently,
\[\mathbb{P}\left(\lim_{n\to\infty}\overline X_n=\mu\right)=1.\]}

\end{tcolorbox}

\begin{tcolorbox}[enhanced jigsaw, arc=.35mm, bottomrule=.15mm, bottomtitle=1mm, breakable, colback=white, colbacktitle=quarto-callout-note-color!10!white, colframe=quarto-callout-note-color-frame, coltitle=black, left=2mm, leftrule=.75mm, opacityback=0, opacitybacktitle=0.6, rightrule=.15mm, title=\textcolor{quarto-callout-note-color}{\faInfo}\hspace{0.5em}{Remark}, titlerule=0mm, toprule=.15mm, toptitle=1mm]

\emph{Remark 27}. The weak law says the probability of a large error
goes to \(0\). The strong law says that, with probability \(1\), the
sequence of observed sample averages itself converges to \(\mu\).

\end{tcolorbox}

\begin{tcolorbox}[enhanced jigsaw, arc=.35mm, bottomrule=.15mm, bottomtitle=1mm, breakable, colback=white, colbacktitle=quarto-callout-tip-color!10!white, colframe=quarto-callout-tip-color-frame, coltitle=black, left=2mm, leftrule=.75mm, opacityback=0, opacitybacktitle=0.6, rightrule=.15mm, title=\textcolor{quarto-callout-tip-color}{\faLightbulb}\hspace{0.5em}{Example}, titlerule=0mm, toprule=.15mm, toptitle=1mm]

\textbf{Example 28} (Long-run frequency). In repeated iid Bernoulli
trials with success probability \(p\), the sample proportion of
successes satisfies \[\overline X_n\xrightarrow{a.s.}p.\] This
formalizes the long-run frequency interpretation of probability.

\end{tcolorbox}

\begin{tcolorbox}[enhanced jigsaw, arc=.35mm, bottomrule=.15mm, bottomtitle=1mm, breakable, colback=white, colbacktitle=quarto-callout-note-color!10!white, colframe=quarto-callout-note-color-frame, coltitle=black, left=2mm, leftrule=.75mm, opacityback=0, opacitybacktitle=0.6, rightrule=.15mm, title=\textcolor{quarto-callout-note-color}{\faInfo}\hspace{0.5em}{Solution}, titlerule=0mm, toprule=.15mm, toptitle=1mm]

Since \(X_i\sim \operatorname{Bernoulli}(p)\) have finite mean \(p\) and
finite variance \(p(1-p)\), the strong law applies directly:
\[\mathbb{P}\left(\lim_{n\to\infty}\overline X_n=p\right)=1.\]

\end{tcolorbox}

\bookmarksetup{startatroot}

\chapter{Central Limit Theorem}\label{central-limit-theorem}

This section states the central approximation result explaining why
normal distributions appear throughout statistics.

\begin{tcolorbox}[enhanced jigsaw, arc=.35mm, bottomrule=.15mm, bottomtitle=1mm, breakable, colback=white, colbacktitle=quarto-callout-important-color!10!white, colframe=quarto-callout-important-color-frame, coltitle=black, left=2mm, leftrule=.75mm, opacityback=0, opacitybacktitle=0.6, rightrule=.15mm, title=\textcolor{quarto-callout-important-color}{\faExclamation}\hspace{0.5em}{Theorem}, titlerule=0mm, toprule=.15mm, toptitle=1mm]

\textbf{Theorem 29} (Central limit theorem). \emph{Suppose
\(X_1,X_2,\ldots\) are iid with mean \[\mu=\mathbb{E}X_i<\infty\] and
variance \[\operatorname{Var}(X_i)=\sigma^2<\infty.\] Then
\[\frac{\overline X_n-\mu}{\sigma/\sqrt n}\xrightarrow{D}N(0,1).\]
Equivalently,
\[\overline X_n \approx N\left(\mu,\frac{\sigma^2}{n}\right)\] for large
\(n\), and \[X_1+\cdots+X_n \approx N(n\mu,n\sigma^2).\]}

\end{tcolorbox}

\begin{tcolorbox}[enhanced jigsaw, arc=.35mm, bottomrule=.15mm, bottomtitle=1mm, breakable, colback=white, colbacktitle=quarto-callout-warning-color!10!white, colframe=quarto-callout-warning-color-frame, coltitle=black, left=2mm, leftrule=.75mm, opacityback=0, opacitybacktitle=0.6, rightrule=.15mm, title=\textcolor{quarto-callout-warning-color}{\faExclamationTriangle}\hspace{0.5em}{Common mistake}, titlerule=0mm, toprule=.15mm, toptitle=1mm]

Precise interpretation The expression
\[\overline X_n \approx N\left(\mu,\frac{\sigma^2}{n}\right)\] is an
approximation for large \(n\). The exact convergence statement is the
standardized convergence
\[\frac{\overline X_n-\mu}{\sigma/\sqrt n}\xrightarrow{D}N(0,1).\]

\end{tcolorbox}

\begin{tcolorbox}[enhanced jigsaw, arc=.35mm, bottomrule=.15mm, bottomtitle=1mm, breakable, colback=white, colbacktitle=quarto-callout-tip-color!10!white, colframe=quarto-callout-tip-color-frame, coltitle=black, left=2mm, leftrule=.75mm, opacityback=0, opacitybacktitle=0.6, rightrule=.15mm, title=\textcolor{quarto-callout-tip-color}{\faLightbulb}\hspace{0.5em}{Example}, titlerule=0mm, toprule=.15mm, toptitle=1mm]

\textbf{Example 30} (Binomial approximation). Let
\(X\sim \operatorname{Binomial}(n,p)\). Use the CLT to approximate the
distribution of \(X\) for large \(n\).

\end{tcolorbox}

\begin{tcolorbox}[enhanced jigsaw, arc=.35mm, bottomrule=.15mm, bottomtitle=1mm, breakable, colback=white, colbacktitle=quarto-callout-note-color!10!white, colframe=quarto-callout-note-color-frame, coltitle=black, left=2mm, leftrule=.75mm, opacityback=0, opacitybacktitle=0.6, rightrule=.15mm, title=\textcolor{quarto-callout-note-color}{\faInfo}\hspace{0.5em}{Solution}, titlerule=0mm, toprule=.15mm, toptitle=1mm]

Write \[X=X_1+\cdots+X_n,\] where
\(X_i\sim \operatorname{Bernoulli}(p)\) are iid. Then \(\mu=p\) and
\(\sigma^2=p(1-p)\). By the CLT,
\[\frac{X-np}{\sqrt{np(1-p)}}\xrightarrow{D}N(0,1).\] Thus for large
\(n\), \[X\approx N(np,np(1-p)).\]

\end{tcolorbox}

\begin{tcolorbox}[enhanced jigsaw, arc=.35mm, bottomrule=.15mm, bottomtitle=1mm, breakable, colback=white, colbacktitle=quarto-callout-warning-color!10!white, colframe=quarto-callout-warning-color-frame, coltitle=black, left=2mm, leftrule=.75mm, opacityback=0, opacitybacktitle=0.6, rightrule=.15mm, title=\textcolor{quarto-callout-warning-color}{\faExclamationTriangle}\hspace{0.5em}{Practice Problem}, titlerule=0mm, toprule=.15mm, toptitle=1mm]

\textbf{Practice Problem 31} (Standardized sample mean). Suppose
\(X_1,\ldots,X_n\) are iid with \(\mathbb{E}X_i=10\) and
\(\operatorname{Var}(X_i)=25\). Find the approximate distribution of
\(\overline X_n\) for large \(n\).

\end{tcolorbox}

\begin{tcolorbox}[enhanced jigsaw, arc=.35mm, bottomrule=.15mm, bottomtitle=1mm, breakable, colback=white, colbacktitle=quarto-callout-note-color!10!white, colframe=quarto-callout-note-color-frame, coltitle=black, left=2mm, leftrule=.75mm, opacityback=0, opacitybacktitle=0.6, rightrule=.15mm, title=\textcolor{quarto-callout-note-color}{\faInfo}\hspace{0.5em}{Solution}, titlerule=0mm, toprule=.15mm, toptitle=1mm]

Here \(\mu=10\) and \(\sigma^2=25\). The CLT gives
\[\frac{\overline X_n-10}{5/\sqrt n}\xrightarrow{D}N(0,1),\] so for
large \(n\), \[\overline X_n\approx N\left(10,\frac{25}{n}\right).\]

\end{tcolorbox}

\bookmarksetup{startatroot}

\chapter{Continuous Mapping and Slutsky's
Theorem}\label{continuous-mapping-and-slutskys-theorem}

This section gives two tools for transforming convergent random
variables.

\section{Continuous mapping theorem}\label{continuous-mapping-theorem}

The continuous mapping theorem says that continuous transformations
preserve convergence.

\begin{tcolorbox}[enhanced jigsaw, arc=.35mm, bottomrule=.15mm, bottomtitle=1mm, breakable, colback=white, colbacktitle=quarto-callout-important-color!10!white, colframe=quarto-callout-important-color-frame, coltitle=black, left=2mm, leftrule=.75mm, opacityback=0, opacitybacktitle=0.6, rightrule=.15mm, title=\textcolor{quarto-callout-important-color}{\faExclamation}\hspace{0.5em}{Theorem}, titlerule=0mm, toprule=.15mm, toptitle=1mm]

\textbf{Theorem 32} (Continuous mapping theorem). \emph{Let \(g\) be a
continuous function.}

\begin{enumerate}
\def\labelenumi{\arabic{enumi}.}
\item
  \emph{If \(X_n\xrightarrow{P}X\), then \[g(X_n)\xrightarrow{P}g(X).\]}
\item
  \emph{If \(X_n\xrightarrow{D}X\), then \[g(X_n)\xrightarrow{D}g(X).\]}
\end{enumerate}

\end{tcolorbox}

\begin{tcolorbox}[enhanced jigsaw, arc=.35mm, bottomrule=.15mm, bottomtitle=1mm, breakable, colback=white, colbacktitle=quarto-callout-tip-color!10!white, colframe=quarto-callout-tip-color-frame, coltitle=black, left=2mm, leftrule=.75mm, opacityback=0, opacitybacktitle=0.6, rightrule=.15mm, title=\textcolor{quarto-callout-tip-color}{\faLightbulb}\hspace{0.5em}{Example}, titlerule=0mm, toprule=.15mm, toptitle=1mm]

\textbf{Example 33} (Squaring a consistent estimator). Suppose
\(X_n\xrightarrow{P}\theta\). Show that
\(X_n^2\xrightarrow{P}\theta^2\).

\end{tcolorbox}

\begin{tcolorbox}[enhanced jigsaw, arc=.35mm, bottomrule=.15mm, bottomtitle=1mm, breakable, colback=white, colbacktitle=quarto-callout-note-color!10!white, colframe=quarto-callout-note-color-frame, coltitle=black, left=2mm, leftrule=.75mm, opacityback=0, opacitybacktitle=0.6, rightrule=.15mm, title=\textcolor{quarto-callout-note-color}{\faInfo}\hspace{0.5em}{Solution}, titlerule=0mm, toprule=.15mm, toptitle=1mm]

Use the continuous mapping theorem with \(g(x)=x^2\), which is
continuous on \(\mathbb{R}\). Therefore
\[X_n^2=g(X_n)\xrightarrow{P}g(\theta)=\theta^2.\]

\end{tcolorbox}

\section{Slutsky's theorem}\label{slutskys-theorem}

Slutsky's theorem is essential for asymptotic statistics because it lets
us replace unknown constants by consistent estimates.

\begin{tcolorbox}[enhanced jigsaw, arc=.35mm, bottomrule=.15mm, bottomtitle=1mm, breakable, colback=white, colbacktitle=quarto-callout-important-color!10!white, colframe=quarto-callout-important-color-frame, coltitle=black, left=2mm, leftrule=.75mm, opacityback=0, opacitybacktitle=0.6, rightrule=.15mm, title=\textcolor{quarto-callout-important-color}{\faExclamation}\hspace{0.5em}{Theorem}, titlerule=0mm, toprule=.15mm, toptitle=1mm]

\textbf{Theorem 34} (Slutsky's theorem). \emph{If \[X_n\xrightarrow{D}X
\qquad \text{and} \qquad
Y_n\xrightarrow{P}c,\] where \(c\) is a constant, then
\[X_n+Y_n\xrightarrow{D}X+c,\] \[X_nY_n\xrightarrow{D}cX,\] and if
\(c\ne 0\), \[\frac{X_n}{Y_n}\xrightarrow{D}\frac{X}{c}.\]}

\end{tcolorbox}

\begin{tcolorbox}[enhanced jigsaw, arc=.35mm, bottomrule=.15mm, bottomtitle=1mm, breakable, colback=white, colbacktitle=quarto-callout-tip-color!10!white, colframe=quarto-callout-tip-color-frame, coltitle=black, left=2mm, leftrule=.75mm, opacityback=0, opacitybacktitle=0.6, rightrule=.15mm, title=\textcolor{quarto-callout-tip-color}{\faLightbulb}\hspace{0.5em}{Example}, titlerule=0mm, toprule=.15mm, toptitle=1mm]

\textbf{Example 35} (Replacing variance by a consistent estimator).
Suppose
\[\frac{\sqrt n(\overline X_n-\mu)}{\sigma}\xrightarrow{D}N(0,1)\] and
suppose \(S_n\xrightarrow{P}\sigma\), with \(\sigma>0\). Show that
\[\frac{\sqrt n(\overline X_n-\mu)}{S_n}\xrightarrow{D}N(0,1).\]

\end{tcolorbox}

\begin{tcolorbox}[enhanced jigsaw, arc=.35mm, bottomrule=.15mm, bottomtitle=1mm, breakable, colback=white, colbacktitle=quarto-callout-note-color!10!white, colframe=quarto-callout-note-color-frame, coltitle=black, left=2mm, leftrule=.75mm, opacityback=0, opacitybacktitle=0.6, rightrule=.15mm, title=\textcolor{quarto-callout-note-color}{\faInfo}\hspace{0.5em}{Solution}, titlerule=0mm, toprule=.15mm, toptitle=1mm]

Write \[\frac{\sqrt n(\overline X_n-\mu)}{S_n}
=\frac{\sqrt n(\overline X_n-\mu)}{\sigma}\cdot \frac{\sigma}{S_n}.\]
The first factor converges in distribution to \(N(0,1)\). Since
\(S_n\xrightarrow{P}\sigma\), the second factor satisfies
\[\frac{\sigma}{S_n}\xrightarrow{P}1.\] By Slutsky's theorem, the
product converges in distribution to \[N(0,1)\cdot 1=N(0,1).\]

\end{tcolorbox}

\bookmarksetup{startatroot}

\chapter{Statistical Meaning of
Convergence}\label{statistical-meaning-of-convergence}

This section connects convergence theory to estimation and inference.

\begin{itemize}
\item
  Convergence in probability is related to \textbf{statistical
  consistency}. An estimator \(\widehat\theta_n\) is consistent for
  \(\theta\) if \[\widehat\theta_n\xrightarrow{P}\theta.\]
\item
  Convergence in distribution is often used to construct maximum
  likelihood approximations, confidence intervals, and hypothesis tests.
\item
  The CLT often provides the limiting distribution of a properly
  standardized estimator.
\end{itemize}

\begin{tcolorbox}[enhanced jigsaw, arc=.35mm, bottomrule=.15mm, bottomtitle=1mm, breakable, colback=white, colbacktitle=quarto-callout-tip-color!10!white, colframe=quarto-callout-tip-color-frame, coltitle=black, left=2mm, leftrule=.75mm, opacityback=0, opacitybacktitle=0.6, rightrule=.15mm, title=\textcolor{quarto-callout-tip-color}{\faLightbulb}\hspace{0.5em}{Example}, titlerule=0mm, toprule=.15mm, toptitle=1mm]

\textbf{Example 36} (Sample mean as a consistent estimator). Let
\(X_1,\ldots,X_n\) be iid with mean \(\mu\) and finite variance. Show
that \(\overline X_n\) is a consistent estimator of \(\mu\).

\end{tcolorbox}

\begin{tcolorbox}[enhanced jigsaw, arc=.35mm, bottomrule=.15mm, bottomtitle=1mm, breakable, colback=white, colbacktitle=quarto-callout-note-color!10!white, colframe=quarto-callout-note-color-frame, coltitle=black, left=2mm, leftrule=.75mm, opacityback=0, opacitybacktitle=0.6, rightrule=.15mm, title=\textcolor{quarto-callout-note-color}{\faInfo}\hspace{0.5em}{Solution}, titlerule=0mm, toprule=.15mm, toptitle=1mm]

By the weak law of large numbers, \[\overline X_n\xrightarrow{P}\mu.\]
Therefore \(\overline X_n\) is consistent for \(\mu\).

\end{tcolorbox}

\bookmarksetup{startatroot}

\chapter{Concentration Inequalities and
Rates}\label{concentration-inequalities-and-rates}

This section explains how concentration inequalities strengthen
convergence statements by providing explicit finite-sample bounds.

A concentration inequality aims to find a function
\(\phi_n(\epsilon)\to 0\) such that
\[\mathbb{P}(|X_n-\mathbb{E}X_n|>\epsilon)\le \phi_n(\epsilon).\] This
automatically gives convergence in probability. The rate at which
\(\phi_n(\epsilon)\) goes to zero describes how fast \(X_n\)
concentrates around its expectation.

\section{Concentration of a Gaussian
mean}\label{concentration-of-a-gaussian-mean}

This subsection derives a Gaussian-rate tail bound for the sample mean
of normal observations.

\begin{tcolorbox}[enhanced jigsaw, arc=.35mm, bottomrule=.15mm, bottomtitle=1mm, breakable, colback=white, colbacktitle=quarto-callout-tip-color!10!white, colframe=quarto-callout-tip-color-frame, coltitle=black, left=2mm, leftrule=.75mm, opacityback=0, opacitybacktitle=0.6, rightrule=.15mm, title=\textcolor{quarto-callout-tip-color}{\faLightbulb}\hspace{0.5em}{Example}, titlerule=0mm, toprule=.15mm, toptitle=1mm]

\textbf{Example 37} (Gaussian sample mean concentration). Let
\[X_1,\ldots,X_n\overset{\text{iid}}{\sim}N(0,\sigma^2)\] and let
\(\overline X_n\) be the sample mean. Show that
\[\mathbb{P}(|\overline X_n|>\epsilon)
\le 2\exp\left(-\frac{n\epsilon^2}{2\sigma^2}\right).\]

\end{tcolorbox}

\begin{tcolorbox}[enhanced jigsaw, arc=.35mm, bottomrule=.15mm, bottomtitle=1mm, breakable, colback=white, colbacktitle=quarto-callout-note-color!10!white, colframe=quarto-callout-note-color-frame, coltitle=black, left=2mm, leftrule=.75mm, opacityback=0, opacitybacktitle=0.6, rightrule=.15mm, title=\textcolor{quarto-callout-note-color}{\faInfo}\hspace{0.5em}{Solution}, titlerule=0mm, toprule=.15mm, toptitle=1mm]

Since \(X_i\) are normal,
\[\overline X_n\sim N\left(0,\frac{\sigma^2}{n}\right).\] Let
\(Z\sim N(0,1)\). Then \[\mathbb{P}(|\overline X_n|>\epsilon)
=\mathbb{P}\left(|Z|>\frac{\sqrt n\epsilon}{\sigma}\right).\] Using the
standard normal tail bound \[\mathbb{P}(|Z|>t)\le 2e^{-t^2/2},\] we
obtain \[\mathbb{P}(|\overline X_n|>\epsilon)
\le 2\exp\left(-\frac{n\epsilon^2}{2\sigma^2}\right).\]

\end{tcolorbox}

\section{Concentration of a maximum}\label{concentration-of-a-maximum}

This subsection shows how a union bound converts a one-variable tail
bound into a maximum tail bound.

\begin{tcolorbox}[enhanced jigsaw, arc=.35mm, bottomrule=.15mm, bottomtitle=1mm, breakable, colback=white, colbacktitle=quarto-callout-tip-color!10!white, colframe=quarto-callout-tip-color-frame, coltitle=black, left=2mm, leftrule=.75mm, opacityback=0, opacitybacktitle=0.6, rightrule=.15mm, title=\textcolor{quarto-callout-tip-color}{\faLightbulb}\hspace{0.5em}{Example}, titlerule=0mm, toprule=.15mm, toptitle=1mm]

\textbf{Example 38} (Maximum of Gaussian variables). Let
\[X_1,\ldots,X_n\overset{\text{iid}}{\sim}N(0,\sigma^2)\] and define
\[Z_n=\max\{|X_1|,\ldots,|X_n|\}.\] Show that \[\mathbb{P}(Z_n>\epsilon)
\le 2n\exp\left(-\frac{\epsilon^2}{2\sigma^2}\right).\]

\end{tcolorbox}

\begin{tcolorbox}[enhanced jigsaw, arc=.35mm, bottomrule=.15mm, bottomtitle=1mm, breakable, colback=white, colbacktitle=quarto-callout-note-color!10!white, colframe=quarto-callout-note-color-frame, coltitle=black, left=2mm, leftrule=.75mm, opacityback=0, opacitybacktitle=0.6, rightrule=.15mm, title=\textcolor{quarto-callout-note-color}{\faInfo}\hspace{0.5em}{Solution}, titlerule=0mm, toprule=.15mm, toptitle=1mm]

By the union bound, \[\mathbb{P}(Z_n>\epsilon)
=\mathbb{P}\left(\bigcup_{i=1}^n\{|X_i|>\epsilon\}\right)
\le \sum_{i=1}^n \mathbb{P}(|X_i|>\epsilon).\] For each \(i\),
\[\mathbb{P}(|X_i|>\epsilon)
\le 2\exp\left(-\frac{\epsilon^2}{2\sigma^2}\right).\] Therefore
\[\mathbb{P}(Z_n>\epsilon)
\le 2n\exp\left(-\frac{\epsilon^2}{2\sigma^2}\right).\]

\end{tcolorbox}

\section{Chebyshev versus Hoeffding for sample
means}\label{chebyshev-versus-hoeffding-for-sample-means}

This subsection compares polynomial and exponential rates for the sample
mean.

Recall Chebyshev's inequality: \[\mathbb{P}(|X-\mathbb{E}X|\ge \epsilon)
\le \frac{\operatorname{Var}(X)}{\epsilon^2}.\] If \(X_1,\ldots,X_n\)
are iid with variance \(\sigma^2\), then
\[\mathbb{P}(|\overline X_n-\mathbb{E}\overline X_n|\ge \epsilon)
\le \frac{\sigma^2}{n\epsilon^2}.\] This gives a rate of order \(1/n\).

\begin{tcolorbox}[enhanced jigsaw, arc=.35mm, bottomrule=.15mm, bottomtitle=1mm, breakable, colback=white, colbacktitle=quarto-callout-important-color!10!white, colframe=quarto-callout-important-color-frame, coltitle=black, left=2mm, leftrule=.75mm, opacityback=0, opacitybacktitle=0.6, rightrule=.15mm, title=\textcolor{quarto-callout-important-color}{\faExclamation}\hspace{0.5em}{Theorem}, titlerule=0mm, toprule=.15mm, toptitle=1mm]

\textbf{Theorem 39} (Hoeffding's inequality). \emph{Suppose \(X_1,
\ldots,X_n\) are independent and \(0\le X_i\le 1\). Then for any
\(\epsilon>0\),
\[\mathbb{P}(|\overline X_n-\mathbb{E}\overline X_n|\ge \epsilon)
\le 2e^{-2n\epsilon^2}.\] More generally, if \(a\le X_i\le b\), then the
result can be scaled accordingly.}

\end{tcolorbox}

\begin{tcolorbox}[enhanced jigsaw, arc=.35mm, bottomrule=.15mm, bottomtitle=1mm, breakable, colback=white, colbacktitle=quarto-callout-note-color!10!white, colframe=quarto-callout-note-color-frame, coltitle=black, left=2mm, leftrule=.75mm, opacityback=0, opacitybacktitle=0.6, rightrule=.15mm, title=\textcolor{quarto-callout-note-color}{\faInfo}\hspace{0.5em}{Remark}, titlerule=0mm, toprule=.15mm, toptitle=1mm]

\emph{Remark 40}. Hoeffding's inequality gives exponential
concentration. This is much faster than the rate given by Chebyshev's
inequality. Such exponential rates are important in high-dimensional
statistics, nonparametric inference, semiparametric inference, and
empirical risk minimization.

\end{tcolorbox}

\begin{tcolorbox}[enhanced jigsaw, arc=.35mm, bottomrule=.15mm, bottomtitle=1mm, breakable, colback=white, colbacktitle=quarto-callout-warning-color!10!white, colframe=quarto-callout-warning-color-frame, coltitle=black, left=2mm, leftrule=.75mm, opacityback=0, opacitybacktitle=0.6, rightrule=.15mm, title=\textcolor{quarto-callout-warning-color}{\faExclamationTriangle}\hspace{0.5em}{Practice Problem}, titlerule=0mm, toprule=.15mm, toptitle=1mm]

\textbf{Practice Problem 41} (Comparing bounds). Suppose
\(X_1,\ldots,X_n\overset{\text{iid}}{\sim}\operatorname{Bernoulli}(p)\)
and \(\overline X_n\) is the sample proportion. Give a Chebyshev bound
and a Hoeffding bound for \[\mathbb{P}(|\overline X_n-p|\ge \epsilon).\]

\end{tcolorbox}

\begin{tcolorbox}[enhanced jigsaw, arc=.35mm, bottomrule=.15mm, bottomtitle=1mm, breakable, colback=white, colbacktitle=quarto-callout-note-color!10!white, colframe=quarto-callout-note-color-frame, coltitle=black, left=2mm, leftrule=.75mm, opacityback=0, opacitybacktitle=0.6, rightrule=.15mm, title=\textcolor{quarto-callout-note-color}{\faInfo}\hspace{0.5em}{Solution}, titlerule=0mm, toprule=.15mm, toptitle=1mm]

Since \(\operatorname{Var}(X_i)=p(1-p)\), Chebyshev gives
\[\mathbb{P}(|\overline X_n-p|\ge \epsilon)
\le \frac{p(1-p)}{n\epsilon^2}.\] Since Bernoulli variables lie in
\([0,1]\), Hoeffding gives \[\mathbb{P}(|\overline X_n-p|\ge \epsilon)
\le 2e^{-2n\epsilon^2}.\] The Hoeffding bound decays exponentially in
\(n\), while the Chebyshev bound decays only like \(1/n\).

\end{tcolorbox}

\section{Consistency of estimating a high-dimensional
proportion}\label{consistency-of-estimating-a-high-dimensional-proportion}

This subsection illustrates why concentration inequalities are central
in modern statistics.

Consider iid binary observations \[X_1,
\ldots,X_n\in \{0,1\}^d.\] Write \[X_i=(X_{i1},\ldots,X_{id}),\] and
suppose \[\mathbb{E}X_{ij}=p_j,
\qquad j=1,\ldots,d.\] The sample proportion estimator for coordinate
\(j\) is \[\widehat p_j=\frac1n\sum_{i=1}^n X_{ij}.\] We want all \(d\)
coordinates to be close to their targets simultaneously:
\[\max_{1\le j\le d}|\widehat p_j-p_j|.\]

\begin{tcolorbox}[enhanced jigsaw, arc=.35mm, bottomrule=.15mm, bottomtitle=1mm, breakable, colback=white, colbacktitle=quarto-callout-important-color!10!white, colframe=quarto-callout-important-color-frame, coltitle=black, left=2mm, leftrule=.75mm, opacityback=0, opacitybacktitle=0.6, rightrule=.15mm, title=\textcolor{quarto-callout-important-color}{\faExclamation}\hspace{0.5em}{Theorem}, titlerule=0mm, toprule=.15mm, toptitle=1mm]

\textbf{Theorem 42} (Uniform concentration for binary proportions).
\emph{For any \(\epsilon>0\),
\[\mathbb{P}\left(\max_{1\le j\le d}|\widehat p_j-p_j|\ge \epsilon\right)
\le 2d e^{-2n\epsilon^2}.\]}

\end{tcolorbox}

\begin{tcolorbox}[enhanced jigsaw, arc=.35mm, bottomrule=.15mm, bottomtitle=1mm, breakable, colback=white, colbacktitle=quarto-callout-note-color!10!white, colframe=quarto-callout-note-color-frame, coltitle=black, left=2mm, leftrule=.75mm, opacityback=0, opacitybacktitle=0.6, rightrule=.15mm, title=\textcolor{quarto-callout-note-color}{\faInfo}\hspace{0.5em}{Proof}, titlerule=0mm, toprule=.15mm, toptitle=1mm]

\emph{Proof.} For a fixed coordinate \(j\), Hoeffding's inequality gives
\[\mathbb{P}(|\widehat p_j-p_j|\ge \epsilon)
\le 2e^{-2n\epsilon^2}.\] Using the union bound over \(j=1,
\ldots,d\),
\[\mathbb{P}\left(\max_{1\le j\le d}|\widehat p_j-p_j|\ge \epsilon\right)
=\mathbb{P}\left(\bigcup_{j=1}^d\{|
\widehat p_j-p_j|\ge \epsilon\}\right)
\le \sum_{j=1}^d 2e^{-2n\epsilon^2}.\] Therefore
\[\mathbb{P}\left(\max_{1\le j\le d}|\widehat p_j-p_j|\ge \epsilon\right)
\le 2d e^{-2n\epsilon^2}.\] ◻

\end{tcolorbox}

\begin{tcolorbox}[enhanced jigsaw, arc=.35mm, bottomrule=.15mm, bottomtitle=1mm, breakable, colback=white, colbacktitle=quarto-callout-tip-color!10!white, colframe=quarto-callout-tip-color-frame, coltitle=black, left=2mm, leftrule=.75mm, opacityback=0, opacitybacktitle=0.6, rightrule=.15mm, title=\textcolor{quarto-callout-tip-color}{\faLightbulb}\hspace{0.5em}{Key idea}, titlerule=0mm, toprule=.15mm, toptitle=1mm]

The bound \[2d e^{-2n\epsilon^2}\] shows that the sample size \(n\) must
grow at least on the order of \(\log d\) to control all \(d\)
coordinates simultaneously. This logarithmic dependence on dimension is
a key theme in high-dimensional probability and statistics.

\end{tcolorbox}

\begin{tcolorbox}[enhanced jigsaw, arc=.35mm, bottomrule=.15mm, bottomtitle=1mm, breakable, colback=white, colbacktitle=quarto-callout-warning-color!10!white, colframe=quarto-callout-warning-color-frame, coltitle=black, left=2mm, leftrule=.75mm, opacityback=0, opacitybacktitle=0.6, rightrule=.15mm, title=\textcolor{quarto-callout-warning-color}{\faExclamationTriangle}\hspace{0.5em}{Practice Problem}, titlerule=0mm, toprule=.15mm, toptitle=1mm]

\textbf{Practice Problem 43} (Sample size from a concentration bound).
Find a sufficient sample size \(n\) so that
\[\mathbb{P}\left(\max_{1\le j\le d}|\widehat p_j-p_j|\ge \epsilon\right)
\le \alpha.\]

\end{tcolorbox}

\begin{tcolorbox}[enhanced jigsaw, arc=.35mm, bottomrule=.15mm, bottomtitle=1mm, breakable, colback=white, colbacktitle=quarto-callout-note-color!10!white, colframe=quarto-callout-note-color-frame, coltitle=black, left=2mm, leftrule=.75mm, opacityback=0, opacitybacktitle=0.6, rightrule=.15mm, title=\textcolor{quarto-callout-note-color}{\faInfo}\hspace{0.5em}{Solution}, titlerule=0mm, toprule=.15mm, toptitle=1mm]

It is enough to require \[2d e^{-2n\epsilon^2}\le \alpha.\] Taking
logarithms, \[\log(2d)-2n\epsilon^2\le \log\alpha.\] Equivalently,
\[2n\epsilon^2\ge \log\left(\frac{2d}{\alpha}\right).\] Thus a
sufficient condition is
\[n\ge \frac{1}{2\epsilon^2}\log\left(\frac{2d}{\alpha}\right).\]

\end{tcolorbox}

\bookmarksetup{startatroot}

\chapter{Summary of Modes of
Convergence}\label{summary-of-modes-of-convergence}

This section summarizes the four modes of convergence and their main
relationships.

{\def\LTcaptype{none} % do not increment counter
\begin{longtable}[]{@{}
  >{\raggedright\arraybackslash}p{(\linewidth - 4\tabcolsep) * \real{0.3333}}
  >{\raggedright\arraybackslash}p{(\linewidth - 4\tabcolsep) * \real{0.3333}}
  >{\raggedright\arraybackslash}p{(\linewidth - 4\tabcolsep) * \real{0.3333}}@{}}
\toprule\noalign{}
\begin{minipage}[b]{\linewidth}\raggedright
Mode
\end{minipage} & \begin{minipage}[b]{\linewidth}\raggedright
Definition
\end{minipage} & \begin{minipage}[b]{\linewidth}\raggedright
Interpretation
\end{minipage} \\
\midrule\noalign{}
\endhead
\bottomrule\noalign{}
\endlastfoot
Distribution & \(F_n(x)\to F(x)\) at continuity points of \(F\) &
Distributional shape converges \\
Probability & \(\mathbb{P}(|X_n-X|>\epsilon)\to 0\) & Large errors
become rare \\
\(r\)-th mean & \(\mathbb{E}|X_n-X|^r\to 0\) & Average \(r\)-th power
error vanishes \\
Almost surely & \(\mathbb{P}(\lim X_n=X)=1\) & Sample paths converge
with probability \(1\) \\
\end{longtable}
}

\begin{tcolorbox}[enhanced jigsaw, arc=.35mm, bottomrule=.15mm, bottomtitle=1mm, breakable, colback=white, colbacktitle=quarto-callout-tip-color!10!white, colframe=quarto-callout-tip-color-frame, coltitle=black, left=2mm, leftrule=.75mm, opacityback=0, opacitybacktitle=0.6, rightrule=.15mm, title=\textcolor{quarto-callout-tip-color}{\faLightbulb}\hspace{0.5em}{Key idea}, titlerule=0mm, toprule=.15mm, toptitle=1mm]

\[X_n\xrightarrow{a.s.}X \Longrightarrow X_n\xrightarrow{P}X \Longrightarrow X_n\xrightarrow{D}X.\]
Also, \[X_n\xrightarrow{L^r}X \Longrightarrow X_n\xrightarrow{P}X.\] The
converses are generally false without additional assumptions.

\end{tcolorbox}

\begin{tcolorbox}[enhanced jigsaw, arc=.35mm, bottomrule=.15mm, bottomtitle=1mm, breakable, colback=white, colbacktitle=quarto-callout-warning-color!10!white, colframe=quarto-callout-warning-color-frame, coltitle=black, left=2mm, leftrule=.75mm, opacityback=0, opacitybacktitle=0.6, rightrule=.15mm, title=\textcolor{quarto-callout-warning-color}{\faExclamationTriangle}\hspace{0.5em}{Practice Problem}, titlerule=0mm, toprule=.15mm, toptitle=1mm]

\textbf{Practice Problem 44} (Choosing the right theorem). For each
goal, identify the most relevant theorem.

\begin{enumerate}
\def\labelenumi{\arabic{enumi}.}
\item
  Show that a sample mean is close to \(\mu\) for large \(n\).
\item
  Approximate the distribution of a standardized sample mean.
\item
  Replace an unknown \(\sigma\) by a consistent estimator \(S_n\) in an
  asymptotic normal statement.
\item
  Show that \(g(\widehat\theta_n)\) is consistent if
  \(\widehat\theta_n\) is consistent.
\end{enumerate}

\end{tcolorbox}

\begin{tcolorbox}[enhanced jigsaw, arc=.35mm, bottomrule=.15mm, bottomtitle=1mm, breakable, colback=white, colbacktitle=quarto-callout-note-color!10!white, colframe=quarto-callout-note-color-frame, coltitle=black, left=2mm, leftrule=.75mm, opacityback=0, opacitybacktitle=0.6, rightrule=.15mm, title=\textcolor{quarto-callout-note-color}{\faInfo}\hspace{0.5em}{Solution}, titlerule=0mm, toprule=.15mm, toptitle=1mm]

\begin{enumerate}
\def\labelenumi{\arabic{enumi}.}
\item
  Use the weak law of large numbers, or a concentration inequality if a
  finite-sample rate is needed.
\item
  Use the central limit theorem.
\item
  Use Slutsky's theorem.
\item
  Use the continuous mapping theorem.
\end{enumerate}

\end{tcolorbox}

\bookmarksetup{startatroot}

\chapter{Chapter 10: Monte Carlo
Sampling}\label{chapter-10-monte-carlo-sampling}

This chapter introduces Monte Carlo methods as a practical bridge
between probability theory, statistical inference, numerical
integration, simulation, Bayesian computation, and optimization. The
main principle is to rewrite a difficult quantity as a probability or
expectation, simulate random samples, and approximate the target with a
sample average.

\begin{tcolorbox}[enhanced jigsaw, arc=.35mm, bottomrule=.15mm, bottomtitle=1mm, breakable, colback=white, colbacktitle=quarto-callout-note-color!10!white, colframe=quarto-callout-note-color-frame, coltitle=black, left=2mm, leftrule=.75mm, opacityback=0, opacitybacktitle=0.6, rightrule=.15mm, title=\textcolor{quarto-callout-note-color}{\faInfo}\hspace{0.5em}{Topics}, titlerule=0mm, toprule=.15mm, toptitle=1mm]

Introduction and applications; simulation; Monte Carlo integration; Law
of Large Numbers foundation; standard random-number generation;
inverse-transform sampling; Box--Muller normal sampling; rejection
sampling; importance sampling; sampling-importance-resampling;
advantages and limitations.

\end{tcolorbox}

\bookmarksetup{startatroot}

\chapter{Overview}\label{overview-7}

This section introduces Monte Carlo methods as a systematic way to use
randomness to approximate quantities that are difficult to compute
exactly.

Monte Carlo methods are algorithms that approximate numerical answers by
repeated random sampling. The key idea is simple: if a quantity can be
written as an expectation or probability, then it can often be
approximated by an average over simulated random samples.

\begin{tcolorbox}[enhanced jigsaw, arc=.35mm, bottomrule=.15mm, bottomtitle=1mm, breakable, colback=white, colbacktitle=quarto-callout-tip-color!10!white, colframe=quarto-callout-tip-color-frame, coltitle=black, left=2mm, leftrule=.75mm, opacityback=0, opacitybacktitle=0.6, rightrule=.15mm, title=\textcolor{quarto-callout-tip-color}{\faLightbulb}\hspace{0.5em}{Key idea}, titlerule=0mm, toprule=.15mm, toptitle=1mm]

Main message Monte Carlo sampling replaces a difficult probability,
expectation, integral, sum, or optimization problem by a random
experiment that can be repeated many times:
\[\text{target quantity} \approx \text{sample average from simulated draws}.\]
The accuracy improves as the number of samples increases, by the Law of
Large Numbers.

\end{tcolorbox}

This section emphasizes four central ideas:

\begin{enumerate}
\def\labelenumi{\arabic{enumi}.}
\item
  using simulation to estimate averages and probabilities;
\item
  using Monte Carlo integration to approximate integrals;
\item
  generating random variables from desired distributions;
\item
  using rejection sampling and importance sampling when direct sampling
  is hard.
\end{enumerate}

\bookmarksetup{startatroot}

\chapter{Background and History}\label{background-and-history}

This section places Monte Carlo sampling in its historical and
conceptual context.

\section{Historical background}\label{historical-background}

Monte Carlo methods are named after Monte Carlo, the famous casino
district in Monaco. The name reflects the central role of chance in
these methods, much like chance drives games of roulette.

Historically, several ideas contributed to Monte Carlo methods:

\begin{itemize}
\item
  \textbf{Buffon's needle problem} was studied by Georges-Louis Leclerc,
  Comte de Buffon, in the eighteenth century and gave an early example
  of using random experiments to estimate \(\pi\).
\item
  \textbf{Markov chains} were introduced by Andrey Markov to describe
  systems whose future evolution depends on the present state rather
  than the full past.
\item
  \textbf{Ulam and von Neumann} used Monte Carlo methods in the 1940s
  during the Manhattan Project, for example in neutron diffusion
  simulations.
\item
  \textbf{Markov chain Monte Carlo} later developed through algorithms
  such as Metropolis--Hastings, Gibbs sampling, and Hamiltonian Monte
  Carlo.
\end{itemize}

\begin{tcolorbox}[enhanced jigsaw, arc=.35mm, bottomrule=.15mm, bottomtitle=1mm, breakable, colback=white, colbacktitle=quarto-callout-tip-color!10!white, colframe=quarto-callout-tip-color-frame, coltitle=black, left=2mm, leftrule=.75mm, opacityback=0, opacitybacktitle=0.6, rightrule=.15mm, title=\textcolor{quarto-callout-tip-color}{\faLightbulb}\hspace{0.5em}{Key idea}, titlerule=0mm, toprule=.15mm, toptitle=1mm]

Historical timeline

\end{tcolorbox}

\section{Markov models and Monte Carlo
sampling}\label{markov-models-and-monte-carlo-sampling}

A Markov model describes a system that evolves over time according to a
memoryless rule.

\begin{tcolorbox}[enhanced jigsaw, arc=.35mm, bottomrule=.15mm, bottomtitle=1mm, breakable, colback=white, colbacktitle=quarto-callout-note-color!10!white, colframe=quarto-callout-note-color-frame, coltitle=black, left=2mm, leftrule=.75mm, opacityback=0, opacitybacktitle=0.6, rightrule=.15mm, title=\textcolor{quarto-callout-note-color}{\faInfo}\hspace{0.5em}{Definition}, titlerule=0mm, toprule=.15mm, toptitle=1mm]

\textbf{Definition 1} (Markov chain). A Markov chain is a sequence of
random states in which the probability of the next state depends only on
the current state, not on the full history.

\end{tcolorbox}

Monte Carlo sampling is broader than Markov chains. It refers to methods
that sample from probability distributions in order to estimate desired
quantities. Markov chain Monte Carlo, or MCMC, is used when direct
independent sampling is difficult, but constructing a Markov chain with
the desired stationary distribution is possible.

\begin{tcolorbox}[enhanced jigsaw, arc=.35mm, bottomrule=.15mm, bottomtitle=1mm, breakable, colback=white, colbacktitle=quarto-callout-important-color!10!white, colframe=quarto-callout-important-color-frame, coltitle=black, left=2mm, leftrule=.75mm, opacityback=0, opacitybacktitle=0.6, rightrule=.15mm, title=\textcolor{quarto-callout-important-color}{\faExclamation}\hspace{0.5em}{Remark}, titlerule=0mm, toprule=.15mm, toptitle=1mm]

\emph{Remark 2}. Monte Carlo sampling has limitations when the target
distribution is high-dimensional, sharply concentrated, multimodal, or
known only up to a normalizing constant. These limitations motivate
rejection sampling, importance sampling, and MCMC.

\end{tcolorbox}

\bookmarksetup{startatroot}

\chapter{How Monte Carlo Works}\label{how-monte-carlo-works}

This section explains the basic Monte Carlo principle: repeated random
sampling produces numerical approximations.

\section{The basic Monte Carlo idea}\label{the-basic-monte-carlo-idea}

Monte Carlo methods use random sampling and statistical modeling to
estimate mathematical functions or mimic complex systems.

A typical Monte Carlo procedure has the following form:

\begin{enumerate}
\def\labelenumi{\arabic{enumi}.}
\item
  Choose a probability model for the uncertain quantities.
\item
  Generate random samples from that probability model.
\item
  Evaluate the quantity of interest for each sample.
\item
  Average the results.
\item
  Use more samples to improve accuracy.
\end{enumerate}

\begin{tcolorbox}[enhanced jigsaw, arc=.35mm, bottomrule=.15mm, bottomtitle=1mm, breakable, colback=white, colbacktitle=quarto-callout-tip-color!10!white, colframe=quarto-callout-tip-color-frame, coltitle=black, left=2mm, leftrule=.75mm, opacityback=0, opacitybacktitle=0.6, rightrule=.15mm, title=\textcolor{quarto-callout-tip-color}{\faLightbulb}\hspace{0.5em}{Key idea}, titlerule=0mm, toprule=.15mm, toptitle=1mm]

Monte Carlo principle If
\(X_1,\ldots,X_N\overset{\text{iid}}{\sim}p(x)\) and \(g\) is a
function, then \[\mathbb{E}[g(X)] = \int g(x)p(x)\,\,dx\] can be
approximated by \[\widehat{I}_N = \frac{1}{N}\sum_{i=1}^N g(X_i).\] By
the Law of Large Numbers,
\[\widehat{I}_N \xrightarrow{P} \mathbb{E}[g(X)].\]

\end{tcolorbox}

\section{Applications}\label{applications}

Monte Carlo methods appear in many areas because many scientific and
statistical questions can be written as expectations, probabilities, or
high-dimensional integrals.

\begin{itemize}
\item
  \textbf{Simulation.} We learn about a random object by generating and
  observing many realizations of it. Examples include risk assessment
  and the spread of disease.
\item
  \textbf{Monte Carlo integration.} We estimate an integral by writing
  it as an expectation of a random variable.
\item
  \textbf{Statistical sampling.} MCMC and related methods are used in
  Bayesian inference, posterior computation, model fitting, and
  parameter estimation.
\item
  \textbf{Optimization.} Random search and simulated annealing use
  random sampling to explore complicated objective functions.
\item
  \textbf{Physical simulation.} Monte Carlo methods are fundamental in
  nuclear physics, statistical mechanics, computer graphics, and
  reinforcement learning.
\end{itemize}

\bookmarksetup{startatroot}

\chapter{Foundation: Law of Large
Numbers}\label{foundation-law-of-large-numbers}

This section reviews the probabilistic foundation that makes Monte Carlo
possible.

\section{Motivating example: average
height}\label{motivating-example-average-height}

Suppose we want to know the average height of students in a class. If
the class is small, we can measure everyone and compute the ordinary
average.

Now suppose we want to estimate the average height of all people in
Boston. Let \(\mathcal{B}\) denote the population and let \(f(x)\) be
the height of person \(x\). The population average is
\[\mathbb{E}_{x\in\mathcal{B}}[f(x)] = \frac{1}{N}\sum_{x\in\mathcal{B}} f(x),\]
which may be infeasible to compute exactly. Instead, we randomly survey
\(S\) people \(x^{(1)},\ldots,x^{(S)}\) and use
\[\frac{1}{S}\sum_{s=1}^S f(x^{(s)}).\]

\begin{tcolorbox}[enhanced jigsaw, arc=.35mm, bottomrule=.15mm, bottomtitle=1mm, breakable, colback=white, colbacktitle=quarto-callout-tip-color!10!white, colframe=quarto-callout-tip-color-frame, coltitle=black, left=2mm, leftrule=.75mm, opacityback=0, opacitybacktitle=0.6, rightrule=.15mm, title=\textcolor{quarto-callout-tip-color}{\faLightbulb}\hspace{0.5em}{Key idea}, titlerule=0mm, toprule=.15mm, toptitle=1mm]

Simulation summary Use a sample mean to approximate a true mean.

\end{tcolorbox}

\section{Weak Law of Large Numbers}\label{weak-law-of-large-numbers-1}

The Law of Large Numbers explains why sample averages approximate
expectations.

\begin{tcolorbox}[enhanced jigsaw, arc=.35mm, bottomrule=.15mm, bottomtitle=1mm, breakable, colback=white, colbacktitle=quarto-callout-note-color!10!white, colframe=quarto-callout-note-color-frame, coltitle=black, left=2mm, leftrule=.75mm, opacityback=0, opacitybacktitle=0.6, rightrule=.15mm, title=\textcolor{quarto-callout-note-color}{\faInfo}\hspace{0.5em}{Theorem}, titlerule=0mm, toprule=.15mm, toptitle=1mm]

\textbf{Theorem 3} (Weak Law of Large Numbers). \emph{Let
\(X_1,X_2,\ldots\) be IID random variables with common finite mean
\(\mu\) and variance \(\sigma^2\). Define the sample mean
\[\overline{X}_n = \frac{1}{n}(X_1+\cdots+X_n).\] Then
\[\overline{X}_n \xrightarrow{P} \mu.\] Equivalently, for every
\(\epsilon>0\),
\[\lim_{n\to\infty}\mathbb{P}\big(|\overline{X}_n-\mu|\ge \epsilon\big)=0.\]}

\end{tcolorbox}

\begin{tcolorbox}[enhanced jigsaw, arc=.35mm, bottomrule=.15mm, bottomtitle=1mm, breakable, colback=white, colbacktitle=quarto-callout-note-color!10!white, colframe=quarto-callout-note-color-frame, coltitle=black, left=2mm, leftrule=.75mm, opacityback=0, opacitybacktitle=0.6, rightrule=.15mm, title=\textcolor{quarto-callout-note-color}{\faInfo}\hspace{0.5em}{Solution}, titlerule=0mm, toprule=.15mm, toptitle=1mm]

The proof uses Chebyshev's inequality. Since the \(X_i\) are
independent,
\[\mathbb{E}[\overline{X}_n]=\mu,\qquad \operatorname{Var}(\overline{X}_n)=\frac{\sigma^2}{n}.\]
Therefore, \[\mathbb{P}(|\overline{X}_n-\mu|\ge \epsilon)
\le \frac{\operatorname{Var}(\overline{X}_n)}{\epsilon^2}
=\frac{\sigma^2}{n\epsilon^2}\to 0.\] This proves convergence in
probability.

\end{tcolorbox}

\begin{tcolorbox}[enhanced jigsaw, arc=.35mm, bottomrule=.15mm, bottomtitle=1mm, breakable, colback=white, colbacktitle=quarto-callout-important-color!10!white, colframe=quarto-callout-important-color-frame, coltitle=black, left=2mm, leftrule=.75mm, opacityback=0, opacitybacktitle=0.6, rightrule=.15mm, title=\textcolor{quarto-callout-important-color}{\faExclamation}\hspace{0.5em}{Remark}, titlerule=0mm, toprule=.15mm, toptitle=1mm]

\emph{Remark 4} (Monte Carlo convergence rate). For many Monte Carlo
averages, the standard error is proportional to \[\frac{1}{\sqrt{N}},\]
where \(N\) is the number of samples. This convergence is reliable but
may be slow.

\end{tcolorbox}

\bookmarksetup{startatroot}

\chapter{\texorpdfstring{Application I: Simulation and Estimating
\(\pi\)}{Application I: Simulation and Estimating \textbackslash pi}}\label{application-i-simulation-and-estimating-pi}

This section shows how a probability can be estimated by simulating
random points.

\section{Uniform raindrops in a
square}\label{uniform-raindrops-in-a-square}

Suppose raindrops fall uniformly at random over a square \([-1,1]^2\).
Inside the square is the unit circle \[x^2+y^2\le 1.\] The square has
area \(4\), and the circle has area \(\pi\). Therefore, if \(p\) is the
probability that a random point lands in the circle, then
\[p=\frac{\text{area of circle}}{\text{area of square}}=\frac{\pi}{4}.\]
Thus \[\pi = 4p.\]

If we generate random points \((X_i,Y_i)\) uniformly from \([-1,1]^2\),
then
\[\widehat{p}_N = \frac{1}{N}\sum_{i=1}^N \mathbf{1}\{X_i^2+Y_i^2\le 1\}\]
estimates \(p\), so \[\widehat{\pi}_N = 4\widehat{p}_N
=4\frac{\#\{\text{points inside circle}\}}{\#\{\text{points inside square}\}}.\]

\begin{tcolorbox}[enhanced jigsaw, arc=.35mm, bottomrule=.15mm, bottomtitle=1mm, breakable, colback=white, colbacktitle=quarto-callout-tip-color!10!white, colframe=quarto-callout-tip-color-frame, coltitle=black, left=2mm, leftrule=.75mm, opacityback=0, opacitybacktitle=0.6, rightrule=.15mm, title=\textcolor{quarto-callout-tip-color}{\faLightbulb}\hspace{0.5em}{Example}, titlerule=0mm, toprule=.15mm, toptitle=1mm]

\textbf{Example 5} (Monte Carlo estimate of \(\pi\)). Use uniformly
sampled points in \([-1,1]^2\) to estimate \(\pi\).

\end{tcolorbox}

\begin{tcolorbox}[enhanced jigsaw, arc=.35mm, bottomrule=.15mm, bottomtitle=1mm, breakable, colback=white, colbacktitle=quarto-callout-note-color!10!white, colframe=quarto-callout-note-color-frame, coltitle=black, left=2mm, leftrule=.75mm, opacityback=0, opacitybacktitle=0.6, rightrule=.15mm, title=\textcolor{quarto-callout-note-color}{\faInfo}\hspace{0.5em}{Solution}, titlerule=0mm, toprule=.15mm, toptitle=1mm]

Let \[I_i=\mathbf{1}\{X_i^2+Y_i^2\le 1\}.\] Then
\(I_i\sim\operatorname{Bernoulli}(p)\) with \(p=\pi/4\). The sample mean
\[\overline{I}_N=\frac{1}{N}\sum_{i=1}^N I_i\] converges in probability
to \(p\). Therefore \[4\overline{I}_N \xrightarrow{P} 4p = \pi.\] For
example, using more samples should move the estimate closer to \(\pi\)
on average. Typical simulations may give values such as \(3.2\) for
\(500\) samples, \(3.1312\) for \(10{,}000\) samples, and \(3.142208\)
for \(10^6\) samples.

\end{tcolorbox}

\section{Python simulation}\label{python-simulation}

The following Python code estimates \(\pi\) by random sampling.

\begin{Shaded}
\begin{Highlighting}[]
\ImportTok{import}\NormalTok{ random}

\KeywordTok{def}\NormalTok{ estimate\_pi(num\_samples):}
\NormalTok{    inside\_circle }\OperatorTok{=} \DecValTok{0}

    \ControlFlowTok{for}\NormalTok{ \_ }\KeywordTok{in} \BuiltInTok{range}\NormalTok{(num\_samples):}
\NormalTok{        x }\OperatorTok{=}\NormalTok{ random.uniform(}\OperatorTok{{-}}\DecValTok{1}\NormalTok{, }\DecValTok{1}\NormalTok{)}
\NormalTok{        y }\OperatorTok{=}\NormalTok{ random.uniform(}\OperatorTok{{-}}\DecValTok{1}\NormalTok{, }\DecValTok{1}\NormalTok{)}

        \ControlFlowTok{if}\NormalTok{ x}\OperatorTok{**}\DecValTok{2} \OperatorTok{+}\NormalTok{ y}\OperatorTok{**}\DecValTok{2} \OperatorTok{\textless{}=} \DecValTok{1}\NormalTok{:}
\NormalTok{            inside\_circle }\OperatorTok{+=} \DecValTok{1}

    \ControlFlowTok{return} \DecValTok{4} \OperatorTok{*}\NormalTok{ (inside\_circle }\OperatorTok{/}\NormalTok{ num\_samples)}

\NormalTok{num\_samples }\OperatorTok{=} \DecValTok{1\_000\_000}
\NormalTok{pi\_estimate }\OperatorTok{=}\NormalTok{ estimate\_pi(num\_samples)}
\BuiltInTok{print}\NormalTok{(}\StringTok{"Estimated value of pi:"}\NormalTok{, pi\_estimate)}
\end{Highlighting}
\end{Shaded}

\section{MATLAB simulation}\label{matlab-simulation}

The same idea can be implemented in MATLAB.

\begin{Shaded}
\begin{Highlighting}[]
\KeywordTok{function} \VariableTok{estimated\_pi} \OperatorTok{=} \VariableTok{estimate\_pi}\NormalTok{(}\VariableTok{num\_samples}\NormalTok{)}
    \VariableTok{inside\_circle} \OperatorTok{=} \FloatTok{0}\OperatorTok{;}
    \KeywordTok{for} \VariableTok{i} \OperatorTok{=} \FloatTok{1}\OperatorTok{:}\VariableTok{num\_samples}
        \VariableTok{x} \OperatorTok{=} \FloatTok{2}\OperatorTok{*}\VariableTok{rand}\NormalTok{() }\OperatorTok{{-}} \FloatTok{1}\OperatorTok{;}
        \VariableTok{y} \OperatorTok{=} \FloatTok{2}\OperatorTok{*}\VariableTok{rand}\NormalTok{() }\OperatorTok{{-}} \FloatTok{1}\OperatorTok{;}
        \VariableTok{distance} \OperatorTok{=} \VariableTok{x}\OperatorTok{\^{}}\FloatTok{2} \OperatorTok{+} \VariableTok{y}\OperatorTok{\^{}}\FloatTok{2}\OperatorTok{;}
        \KeywordTok{if} \VariableTok{distance} \OperatorTok{\textless{}=} \FloatTok{1}
            \VariableTok{inside\_circle} \OperatorTok{=} \VariableTok{inside\_circle} \OperatorTok{+} \FloatTok{1}\OperatorTok{;}
        \KeywordTok{end}
    \KeywordTok{end}
    \VariableTok{estimated\_pi} \OperatorTok{=} \FloatTok{4} \OperatorTok{*} \VariableTok{inside\_circle} \OperatorTok{/} \VariableTok{num\_samples}\OperatorTok{;}
\KeywordTok{end}

\CommentTok{\% Example usage}
\VariableTok{num\_samples} \OperatorTok{=} \FloatTok{1000000}\OperatorTok{;}
\VariableTok{pi\_estimate} \OperatorTok{=} \VariableTok{estimate\_pi}\NormalTok{(}\VariableTok{num\_samples}\NormalTok{)}\OperatorTok{;}
\VariableTok{disp}\NormalTok{([}\SpecialStringTok{\textquotesingle{}Estimated value of pi: \textquotesingle{}}\OperatorTok{,} \VariableTok{num2str}\NormalTok{(}\VariableTok{pi\_estimate}\NormalTok{)])}\OperatorTok{;}
\end{Highlighting}
\end{Shaded}

\begin{tcolorbox}[enhanced jigsaw, arc=.35mm, bottomrule=.15mm, bottomtitle=1mm, breakable, colback=white, colbacktitle=quarto-callout-warning-color!10!white, colframe=quarto-callout-warning-color-frame, coltitle=black, left=2mm, leftrule=.75mm, opacityback=0, opacitybacktitle=0.6, rightrule=.15mm, title=\textcolor{quarto-callout-warning-color}{\faExclamationTriangle}\hspace{0.5em}{Warning}, titlerule=0mm, toprule=.15mm, toptitle=1mm]

Common coding detail If the square is \([-1,1]^2\), use
\(x=2\operatorname{rand}()-1\) and \(y=2\operatorname{rand}()-1\). If
instead \(x,y\sim\operatorname{Uniform}(0,1)\), then the simulation
estimates a quarter circle rather than the full square setup.

\end{tcolorbox}

\bookmarksetup{startatroot}

\chapter{Buffon's Needle Problem}\label{buffons-needle-problem}

This section gives a classical Monte Carlo method for estimating \(\pi\)
using random line crossings.

\section{Geometric setup}\label{geometric-setup}

In Buffon's needle problem, parallel lines are separated by distance
\(D\). A needle of length \(L\) is dropped at random, with \(L\le D\).
Let \(P\) be the probability that the needle crosses a line.

If \(\theta\) is the angle between the needle and the lines, and if
\(y\) is the distance from the needle center to the nearest line, then
the needle crosses a line if the vertical projection of half the needle
is at least \(y\): \[y \le \frac{L}{2}\sin\theta.\] The classical
formula is \[P=\frac{2L}{\pi D}.\] Therefore, \[\pi = \frac{2L}{DP}.\]
If \(k\) out of \(n\) dropped needles cross a line, then
\(P\approx k/n\) and \[\widehat{\pi}=\frac{2Ln}{Dk}.\]

\begin{tcolorbox}[enhanced jigsaw, arc=.35mm, bottomrule=.15mm, bottomtitle=1mm, breakable, colback=white, colbacktitle=quarto-callout-tip-color!10!white, colframe=quarto-callout-tip-color-frame, coltitle=black, left=2mm, leftrule=.75mm, opacityback=0, opacitybacktitle=0.6, rightrule=.15mm, title=\textcolor{quarto-callout-tip-color}{\faLightbulb}\hspace{0.5em}{Example}, titlerule=0mm, toprule=.15mm, toptitle=1mm]

\textbf{Example 6} (Buffon's needle estimate). Suppose \(L=1\), \(D=2\),
and \(n\) needles are dropped. If \(k\) needles cross a line, estimate
\(\pi\).

\end{tcolorbox}

\begin{tcolorbox}[enhanced jigsaw, arc=.35mm, bottomrule=.15mm, bottomtitle=1mm, breakable, colback=white, colbacktitle=quarto-callout-note-color!10!white, colframe=quarto-callout-note-color-frame, coltitle=black, left=2mm, leftrule=.75mm, opacityback=0, opacitybacktitle=0.6, rightrule=.15mm, title=\textcolor{quarto-callout-note-color}{\faInfo}\hspace{0.5em}{Solution}, titlerule=0mm, toprule=.15mm, toptitle=1mm]

Here \[P=\frac{2L}{\pi D}=\frac{1}{\pi}.\] Since \(P\approx k/n\), we
obtain \[\frac{k}{n}\approx \frac{1}{\pi},
\qquad\text{so}\qquad
\widehat{\pi}=\frac{n}{k}.\] This estimator is intuitive, but
convergence is slow. Around \(100{,}000\) trials may be needed for only
a few correct digits.

\end{tcolorbox}

\begin{tcolorbox}[enhanced jigsaw, arc=.35mm, bottomrule=.15mm, bottomtitle=1mm, breakable, colback=white, colbacktitle=quarto-callout-important-color!10!white, colframe=quarto-callout-important-color-frame, coltitle=black, left=2mm, leftrule=.75mm, opacityback=0, opacitybacktitle=0.6, rightrule=.15mm, title=\textcolor{quarto-callout-important-color}{\faExclamation}\hspace{0.5em}{Remark}, titlerule=0mm, toprule=.15mm, toptitle=1mm]

\emph{Remark 7}. Buffon's needle is historically important because it
makes the connection between geometry, probability, and numerical
estimation very transparent. It is not computationally efficient
compared with modern numerical methods.

\end{tcolorbox}

\bookmarksetup{startatroot}

\chapter{Application II: Monte Carlo
Integration}\label{application-ii-monte-carlo-integration}

This section explains how to approximate integrals by rewriting them as
expectations.

\section{Expectation as integration}\label{expectation-as-integration}

Let \(X\) be a continuous random variable with density \(p(x)\). Then
\[\mathbb{E}[g(X)] = \int_{-\infty}^{\infty} g(x)p(x)\,\,dx.\] If
\(X_1,\ldots,X_N\overset{\text{iid}}{\sim}p(x)\), then
\[\int_{-\infty}^{\infty} g(x)p(x)\,\,dx
\approx \frac{1}{N}\sum_{i=1}^N g(X_i).\]

\begin{tcolorbox}[enhanced jigsaw, arc=.35mm, bottomrule=.15mm, bottomtitle=1mm, breakable, colback=white, colbacktitle=quarto-callout-tip-color!10!white, colframe=quarto-callout-tip-color-frame, coltitle=black, left=2mm, leftrule=.75mm, opacityback=0, opacitybacktitle=0.6, rightrule=.15mm, title=\textcolor{quarto-callout-tip-color}{\faLightbulb}\hspace{0.5em}{Key idea}, titlerule=0mm, toprule=.15mm, toptitle=1mm]

Monte Carlo integration A difficult integral becomes a sample average
once it is written as an expectation.

\end{tcolorbox}

\section{Example: an integral without a simple closed
form}\label{example-an-integral-without-a-simple-closed-form}

Consider \[I=\int_0^1 e^{-x^3}\,\,dx.\] Let
\(U\sim\operatorname{Uniform}(0,1)\). Since the density of \(U\) is
\(1\) on \([0,1]\), \[I=\mathbb{E}\left[e^{-U^3}\right].\] Generate
\(U_1,\ldots,U_N\overset{\text{iid}}{\sim}\operatorname{Uniform}(0,1)\)
and define \[W_i=e^{-U_i^3}.\] Then
\[\widehat{I}_N=\frac{1}{N}\sum_{i=1}^N e^{-U_i^3}\] estimates \(I\).

\begin{tcolorbox}[enhanced jigsaw, arc=.35mm, bottomrule=.15mm, bottomtitle=1mm, breakable, colback=white, colbacktitle=quarto-callout-tip-color!10!white, colframe=quarto-callout-tip-color-frame, coltitle=black, left=2mm, leftrule=.75mm, opacityback=0, opacitybacktitle=0.6, rightrule=.15mm, title=\textcolor{quarto-callout-tip-color}{\faLightbulb}\hspace{0.5em}{Example}, titlerule=0mm, toprule=.15mm, toptitle=1mm]

\textbf{Example 8} (Monte Carlo integration). Design a Monte Carlo
estimator for \[\int_0^1 e^{-x^3}\,\,dx.\]

\end{tcolorbox}

\begin{tcolorbox}[enhanced jigsaw, arc=.35mm, bottomrule=.15mm, bottomtitle=1mm, breakable, colback=white, colbacktitle=quarto-callout-note-color!10!white, colframe=quarto-callout-note-color-frame, coltitle=black, left=2mm, leftrule=.75mm, opacityback=0, opacitybacktitle=0.6, rightrule=.15mm, title=\textcolor{quarto-callout-note-color}{\faInfo}\hspace{0.5em}{Solution}, titlerule=0mm, toprule=.15mm, toptitle=1mm]

Let
\(U_1,\ldots,U_N\overset{\text{iid}}{\sim}\operatorname{Uniform}(0,1)\).
Use \[\widehat{I}_N=\frac{1}{N}\sum_{i=1}^N e^{-U_i^3}.\] By the Law of
Large Numbers, \[\widehat{I}_N \xrightarrow{P} \mathbb{E}[e^{-U^3}]
=\int_0^1 e^{-x^3}\,\,dx.\]

\end{tcolorbox}

\section{General Monte Carlo
integration}\label{general-monte-carlo-integration}

More generally, suppose \[I=\int f(x)p(x)\,\,dx,\] where \(p(x)\) is a
probability density. If \(X_i\overset{\text{iid}}{\sim}p\), then
\[\widehat{I}_N=\frac{1}{N}\sum_{i=1}^N f(X_i)\] is the basic Monte
Carlo estimator.

\begin{tcolorbox}[enhanced jigsaw, arc=.35mm, bottomrule=.15mm, bottomtitle=1mm, breakable, colback=white, colbacktitle=quarto-callout-important-color!10!white, colframe=quarto-callout-important-color-frame, coltitle=black, left=2mm, leftrule=.75mm, opacityback=0, opacitybacktitle=0.6, rightrule=.15mm, title=\textcolor{quarto-callout-important-color}{\faExclamation}\hspace{0.5em}{Remark}, titlerule=0mm, toprule=.15mm, toptitle=1mm]

\emph{Remark 9}. When we change the sampling distribution \(p\), the
function being averaged changes. This observation leads directly to
importance sampling.

\end{tcolorbox}

\begin{tcolorbox}[enhanced jigsaw, arc=.35mm, bottomrule=.15mm, bottomtitle=1mm, breakable, colback=white, colbacktitle=quarto-callout-warning-color!10!white, colframe=quarto-callout-warning-color-frame, coltitle=black, left=2mm, leftrule=.75mm, opacityback=0, opacitybacktitle=0.6, rightrule=.15mm, title=\textcolor{quarto-callout-warning-color}{\faExclamationTriangle}\hspace{0.5em}{Practice Problem}, titlerule=0mm, toprule=.15mm, toptitle=1mm]

\textbf{Practice Problem 10} (Uniform Monte Carlo integration). Use
Monte Carlo to estimate \[\int_a^b g(x)\,\,dx.\]

\end{tcolorbox}

\begin{tcolorbox}[enhanced jigsaw, arc=.35mm, bottomrule=.15mm, bottomtitle=1mm, breakable, colback=white, colbacktitle=quarto-callout-note-color!10!white, colframe=quarto-callout-note-color-frame, coltitle=black, left=2mm, leftrule=.75mm, opacityback=0, opacitybacktitle=0.6, rightrule=.15mm, title=\textcolor{quarto-callout-note-color}{\faInfo}\hspace{0.5em}{Solution}, titlerule=0mm, toprule=.15mm, toptitle=1mm]

Let \(X\sim\operatorname{Uniform}(a,b)\). Then \(p(x)=1/(b-a)\) on
\([a,b]\), and \[\mathbb{E}[g(X)] = \int_a^b g(x)\frac{1}{b-a}\,\,dx.\]
Thus \[\int_a^b g(x)\,\,dx = (b-a)\mathbb{E}[g(X)].\] Generate
\(X_1,\ldots,X_N\overset{\text{iid}}{\sim}\operatorname{Uniform}(a,b)\)
and use \[\widehat{I}_N=(b-a)\frac{1}{N}\sum_{i=1}^N g(X_i).\]

\end{tcolorbox}

\begin{tcolorbox}[enhanced jigsaw, arc=.35mm, bottomrule=.15mm, bottomtitle=1mm, breakable, colback=white, colbacktitle=quarto-callout-warning-color!10!white, colframe=quarto-callout-warning-color-frame, coltitle=black, left=2mm, leftrule=.75mm, opacityback=0, opacitybacktitle=0.6, rightrule=.15mm, title=\textcolor{quarto-callout-warning-color}{\faExclamationTriangle}\hspace{0.5em}{Practice Problem}, titlerule=0mm, toprule=.15mm, toptitle=1mm]

\textbf{Practice Problem 11} (Practice integral). Use the Monte Carlo
method to estimate \[I=\int_{-1}^{2} \cos^2(x)(x^3+1)\,\,dx.\]

\end{tcolorbox}

\begin{tcolorbox}[enhanced jigsaw, arc=.35mm, bottomrule=.15mm, bottomtitle=1mm, breakable, colback=white, colbacktitle=quarto-callout-note-color!10!white, colframe=quarto-callout-note-color-frame, coltitle=black, left=2mm, leftrule=.75mm, opacityback=0, opacitybacktitle=0.6, rightrule=.15mm, title=\textcolor{quarto-callout-note-color}{\faInfo}\hspace{0.5em}{Solution}, titlerule=0mm, toprule=.15mm, toptitle=1mm]

Let \(X\sim\operatorname{Uniform}(-1,2)\). Then \(b-a=3\), and
\[I=3\mathbb{E}\left[\cos^2(X)(X^3+1)\right].\] Generate
\(X_1,\ldots,X_N\overset{\text{iid}}{\sim}\operatorname{Uniform}(-1,2)\)
and compute
\[\widehat{I}_N=\frac{3}{N}\sum_{i=1}^N \cos^2(X_i)(X_i^3+1).\] For
comparison, a symbolic calculation gives
\[I=\frac{3\sin 4}{2}+\frac{21\cos 4}{16}-\frac{3\cos 2}{16}+\frac{3\sin 2}{8}+\frac{27}{8}
\approx 1.8009.\]

\end{tcolorbox}

\begin{tcolorbox}[enhanced jigsaw, arc=.35mm, bottomrule=.15mm, bottomtitle=1mm, breakable, colback=white, colbacktitle=quarto-callout-warning-color!10!white, colframe=quarto-callout-warning-color-frame, coltitle=black, left=2mm, leftrule=.75mm, opacityback=0, opacitybacktitle=0.6, rightrule=.15mm, title=\textcolor{quarto-callout-warning-color}{\faExclamationTriangle}\hspace{0.5em}{Practice Problem}, titlerule=0mm, toprule=.15mm, toptitle=1mm]

\textbf{Practice Problem 12} (Practice integral). Use Monte Carlo to
estimate \[I=\int_0^1 \frac{e^x-1}{e-1}\,\,dx.\]

\end{tcolorbox}

\begin{tcolorbox}[enhanced jigsaw, arc=.35mm, bottomrule=.15mm, bottomtitle=1mm, breakable, colback=white, colbacktitle=quarto-callout-note-color!10!white, colframe=quarto-callout-note-color-frame, coltitle=black, left=2mm, leftrule=.75mm, opacityback=0, opacitybacktitle=0.6, rightrule=.15mm, title=\textcolor{quarto-callout-note-color}{\faInfo}\hspace{0.5em}{Solution}, titlerule=0mm, toprule=.15mm, toptitle=1mm]

Let \(U\sim\operatorname{Uniform}(0,1)\). Then
\[I=\mathbb{E}\left[\frac{e^U-1}{e-1}\right].\] The Monte Carlo
estimator is
\[\widehat{I}_N=\frac{1}{N}\sum_{i=1}^N \frac{e^{U_i}-1}{e-1},
\qquad U_i\overset{\text{iid}}{\sim}\operatorname{Uniform}(0,1).\] The
exact value is \[I=\frac{1}{e-1}\int_0^1(e^x-1)\,\,dx
=\frac{e-2}{e-1}.\]

\end{tcolorbox}

\bookmarksetup{startatroot}

\chapter{Application III: Bayesian Inference and
Learning}\label{application-iii-bayesian-inference-and-learning}

This section shows why Monte Carlo methods are central in Bayesian
statistics.

Let \(X\in\mathcal{X}\) be an unknown variable and let \(Y\) be observed
data. Bayesian inference often requires integrals that are analytically
intractable.

\section{Normalization}\label{normalization}

Bayes' theorem gives \[p(X\mid Y)=\frac{p(Y\mid X)p(X)}{p(Y)}.\] The
normalizing constant is
\[p(Y)=\int_{\mathcal{X}}p(Y\mid X')p(X')\,\,dX'.\] This integral may be
difficult in high dimensions.

If \(x^{(1)},\ldots,x^{(S)}\) are sampled from the prior \(p(X)\), then
\[p(Y)=\mathbb{E}_{X\sim p(X)}[p(Y\mid X)]
\approx \frac{1}{S}\sum_{s=1}^S p(Y\mid x^{(s)}).\]

\section{Marginalization}\label{marginalization-1}

If a model contains latent variables \(Z\), the marginal posterior may
require \[p(X\mid Y)=\int p(X,Z\mid Y)\,\,dZ.\] Monte Carlo methods
approximate this integral using samples from a suitable distribution.

\section{Posterior expectations}\label{posterior-expectations}

Many Bayesian summaries are posterior expectations:
\[\mathbb{E}_{p(X\mid Y)}[f(X)]
=\int_{\mathcal{X}} f(X)p(X\mid Y)\,\,dX.\] If
\(X^{(1)},\ldots,X^{(S)}\) are samples from \(p(X\mid Y)\), then
\[\mathbb{E}_{p(X\mid Y)}[f(X)]\approx \frac{1}{S}\sum_{s=1}^S f(X^{(s)}).\]

\begin{tcolorbox}[enhanced jigsaw, arc=.35mm, bottomrule=.15mm, bottomtitle=1mm, breakable, colback=white, colbacktitle=quarto-callout-important-color!10!white, colframe=quarto-callout-important-color-frame, coltitle=black, left=2mm, leftrule=.75mm, opacityback=0, opacitybacktitle=0.6, rightrule=.15mm, title=\textcolor{quarto-callout-important-color}{\faExclamation}\hspace{0.5em}{Remark}, titlerule=0mm, toprule=.15mm, toptitle=1mm]

\emph{Remark 13}. In Bayesian inference, Monte Carlo is often needed
because the posterior distribution is known only up to a normalizing
constant.

\end{tcolorbox}

\bookmarksetup{startatroot}

\chapter{Application IV:
Optimization}\label{application-iv-optimization}

This section shows how random sampling can be used to search for good
solutions of optimization problems.

\section{Random search}\label{random-search}

Optimization problems ask us to minimize or maximize an objective
function, often in many dimensions. Random search is a simple Monte
Carlo optimization method.

\begin{Shaded}
\begin{Highlighting}[]
\ImportTok{import}\NormalTok{ numpy }\ImportTok{as}\NormalTok{ np}

\KeywordTok{def}\NormalTok{ objective\_function(x):}
    \ControlFlowTok{return}\NormalTok{ x}\OperatorTok{**}\DecValTok{2}

\KeywordTok{def}\NormalTok{ random\_search(objective\_function, bounds, max\_iter}\OperatorTok{=}\DecValTok{1000}\NormalTok{):}
\NormalTok{    best\_solution }\OperatorTok{=} \VariableTok{None}
\NormalTok{    best\_score }\OperatorTok{=} \BuiltInTok{float}\NormalTok{(}\StringTok{\textquotesingle{}inf\textquotesingle{}}\NormalTok{)}

    \ControlFlowTok{for}\NormalTok{ \_ }\KeywordTok{in} \BuiltInTok{range}\NormalTok{(max\_iter):}
\NormalTok{        solution }\OperatorTok{=}\NormalTok{ np.random.uniform(bounds[}\DecValTok{0}\NormalTok{], bounds[}\DecValTok{1}\NormalTok{])}
\NormalTok{        score }\OperatorTok{=}\NormalTok{ objective\_function(solution)}

        \ControlFlowTok{if}\NormalTok{ score }\OperatorTok{\textless{}}\NormalTok{ best\_score:}
\NormalTok{            best\_solution }\OperatorTok{=}\NormalTok{ solution}
\NormalTok{            best\_score }\OperatorTok{=}\NormalTok{ score}

    \ControlFlowTok{return}\NormalTok{ best\_solution, best\_score}

\NormalTok{bounds }\OperatorTok{=}\NormalTok{ (}\OperatorTok{{-}}\DecValTok{10}\NormalTok{, }\DecValTok{10}\NormalTok{)}
\NormalTok{best\_solution, best\_score }\OperatorTok{=}\NormalTok{ random\_search(objective\_function, bounds)}
\BuiltInTok{print}\NormalTok{(}\StringTok{"Best solution:"}\NormalTok{, best\_solution)}
\BuiltInTok{print}\NormalTok{(}\StringTok{"Best score:"}\NormalTok{, best\_score)}
\end{Highlighting}
\end{Shaded}

\begin{tcolorbox}[enhanced jigsaw, arc=.35mm, bottomrule=.15mm, bottomtitle=1mm, breakable, colback=white, colbacktitle=quarto-callout-tip-color!10!white, colframe=quarto-callout-tip-color-frame, coltitle=black, left=2mm, leftrule=.75mm, opacityback=0, opacitybacktitle=0.6, rightrule=.15mm, title=\textcolor{quarto-callout-tip-color}{\faLightbulb}\hspace{0.5em}{Example}, titlerule=0mm, toprule=.15mm, toptitle=1mm]

\textbf{Example 14} (Random search). Explain why random search can be
useful even though it seems inefficient.

\end{tcolorbox}

\begin{tcolorbox}[enhanced jigsaw, arc=.35mm, bottomrule=.15mm, bottomtitle=1mm, breakable, colback=white, colbacktitle=quarto-callout-note-color!10!white, colframe=quarto-callout-note-color-frame, coltitle=black, left=2mm, leftrule=.75mm, opacityback=0, opacitybacktitle=0.6, rightrule=.15mm, title=\textcolor{quarto-callout-note-color}{\faInfo}\hspace{0.5em}{Solution}, titlerule=0mm, toprule=.15mm, toptitle=1mm]

Random search does not follow a deterministic gradient direction.
Instead, it samples many candidate points from the search space. This
can be useful when the objective function is non-convex, discontinuous,
noisy, or has many local minima. It is also easy to implement and can
work surprisingly well in high-dimensional hyperparameter search when
only a small number of dimensions are truly important.

\end{tcolorbox}

\section{Simulated annealing}\label{simulated-annealing}

Simulated annealing is a Monte Carlo method for global optimization. It
is designed for complicated, multimodal, or non-convex objective
functions.

The method proposes random neighboring states. A better state is
accepted automatically, while a worse state may still be accepted with
probability depending on a temperature parameter. The temperature is
gradually decreased.

\begin{Shaded}
\begin{Highlighting}[]
\ImportTok{import}\NormalTok{ numpy }\ImportTok{as}\NormalTok{ np}

\KeywordTok{def}\NormalTok{ objective\_function(x):}
    \ControlFlowTok{return}\NormalTok{ np.sin(x) }\OperatorTok{+}\NormalTok{ np.sin(}\DecValTok{2}\OperatorTok{*}\NormalTok{x)}

\KeywordTok{def}\NormalTok{ simulated\_annealing(objective\_function, bounds,}
\NormalTok{                        initial\_temperature}\OperatorTok{=}\DecValTok{100}\NormalTok{,}
\NormalTok{                        cooling\_rate}\OperatorTok{=}\FloatTok{0.95}\NormalTok{,}
\NormalTok{                        num\_iterations}\OperatorTok{=}\DecValTok{1000}\NormalTok{):}
\NormalTok{    current\_solution }\OperatorTok{=}\NormalTok{ np.random.uniform(bounds[}\DecValTok{0}\NormalTok{], bounds[}\DecValTok{1}\NormalTok{])}
\NormalTok{    current\_score }\OperatorTok{=}\NormalTok{ objective\_function(current\_solution)}
\NormalTok{    best\_solution }\OperatorTok{=}\NormalTok{ current\_solution}
\NormalTok{    best\_score }\OperatorTok{=}\NormalTok{ current\_score}
\NormalTok{    temperature }\OperatorTok{=}\NormalTok{ initial\_temperature}
\NormalTok{    solution\_history }\OperatorTok{=}\NormalTok{ [(current\_solution, current\_score)]}

    \ControlFlowTok{for}\NormalTok{ \_ }\KeywordTok{in} \BuiltInTok{range}\NormalTok{(num\_iterations):}
\NormalTok{        neighbor\_solution }\OperatorTok{=}\NormalTok{ current\_solution }\OperatorTok{+}\NormalTok{ np.random.uniform(}\OperatorTok{{-}}\FloatTok{0.1}\NormalTok{, }\FloatTok{0.1}\NormalTok{)}
\NormalTok{        neighbor\_score }\OperatorTok{=}\NormalTok{ objective\_function(neighbor\_solution)}

        \ControlFlowTok{if}\NormalTok{ (neighbor\_score }\OperatorTok{\textless{}}\NormalTok{ current\_score }\KeywordTok{or}
\NormalTok{            np.random.rand() }\OperatorTok{\textless{}}\NormalTok{ np.exp((current\_score }\OperatorTok{{-}}\NormalTok{ neighbor\_score) }\OperatorTok{/}\NormalTok{ temperature)):}
\NormalTok{            current\_solution }\OperatorTok{=}\NormalTok{ neighbor\_solution}
\NormalTok{            current\_score }\OperatorTok{=}\NormalTok{ neighbor\_score}

        \ControlFlowTok{if}\NormalTok{ current\_score }\OperatorTok{\textless{}}\NormalTok{ best\_score:}
\NormalTok{            best\_solution }\OperatorTok{=}\NormalTok{ current\_solution}
\NormalTok{            best\_score }\OperatorTok{=}\NormalTok{ current\_score}

\NormalTok{        temperature }\OperatorTok{*=}\NormalTok{ cooling\_rate}
\NormalTok{        solution\_history.append((current\_solution, current\_score))}

    \ControlFlowTok{return}\NormalTok{ best\_solution, best\_score, solution\_history}
\end{Highlighting}
\end{Shaded}

\bookmarksetup{startatroot}

\chapter{Common Distributions Used in Monte
Carlo}\label{common-distributions-used-in-monte-carlo}

This section reviews probability distributions frequently used in Monte
Carlo modeling.

Common probability distributions include:

\begin{itemize}
\item
  \textbf{Uniform distribution:} models equally likely values over an
  interval.
\item
  \textbf{Normal distribution:} models symmetric noise and appears
  naturally through the Central Limit Theorem.
\item
  \textbf{Exponential distribution:} models waiting times between
  independent events.
\item
  \textbf{Poisson distribution:} models event counts over time or space.
\item
  \textbf{Discrete distributions:} model categorical outcomes or finite
  sample spaces.
\item
  \textbf{Lognormal distribution:} models positive variables whose
  logarithm is normal.
\item
  \textbf{Triangular distribution:} a continuous distribution on
  \([a,b]\) with mode \(c\), where \(a<b\) and \(a\le c\le b\).
\end{itemize}

\begin{tcolorbox}[enhanced jigsaw, arc=.35mm, bottomrule=.15mm, bottomtitle=1mm, breakable, colback=white, colbacktitle=quarto-callout-note-color!10!white, colframe=quarto-callout-note-color-frame, coltitle=black, left=2mm, leftrule=.75mm, opacityback=0, opacitybacktitle=0.6, rightrule=.15mm, title=\textcolor{quarto-callout-note-color}{\faInfo}\hspace{0.5em}{Definition}, titlerule=0mm, toprule=.15mm, toptitle=1mm]

\textbf{Definition 15} (Triangular distribution). A triangular
distribution with lower limit \(a\), upper limit \(b\), and mode \(c\)
has density \[f(x)=
\begin{cases}
\dfrac{2(x-a)}{(b-a)(c-a)}, & a\le x\le c,\\[6pt]
\dfrac{2(b-x)}{(b-a)(b-c)}, & c<x\le b,\\[6pt]
0, & \text{otherwise}.
\end{cases}\]

\end{tcolorbox}

\bookmarksetup{startatroot}

\chapter{Random Number Generation}\label{random-number-generation}

This section discusses the starting point of Monte Carlo: generating
random numbers.

\section{Uniform random variables}\label{uniform-random-variables}

Most random number generation begins with samples from a uniform
distribution on \([0,1]\). The general workflow is:

\begin{enumerate}
\def\labelenumi{\arabic{enumi}.}
\item
  Generate IID uniform random numbers
  \(U_1,U_2,\ldots\sim\operatorname{Uniform}(0,1)\).
\item
  Transform those uniform random numbers into samples from another
  desired distribution.
\end{enumerate}

If a target density \(f\) has CDF \[F(x)=\int_{-\infty}^x f(t)\,\,dt,\]
then the inverse-transform method uses \(F^{-1}\) to transform uniform
draws into samples from \(F\).

\section{Pseudorandom number
generators}\label{pseudorandom-number-generators}

In practice, computers generate pseudorandom numbers. A basic example is
a linear congruential generator.

\begin{tcolorbox}[enhanced jigsaw, arc=.35mm, bottomrule=.15mm, bottomtitle=1mm, breakable, colback=white, colbacktitle=quarto-callout-note-color!10!white, colframe=quarto-callout-note-color-frame, coltitle=black, left=2mm, leftrule=.75mm, opacityback=0, opacitybacktitle=0.6, rightrule=.15mm, title=\textcolor{quarto-callout-note-color}{\faInfo}\hspace{0.5em}{Definition}, titlerule=0mm, toprule=.15mm, toptitle=1mm]

\textbf{Definition 16} (Linear congruential generator). Choose a
starting seed \(X_0\) and constants \(a,c,m\). Generate
\[X_{n+1}=(aX_n+c)\mod m.\] Then \[U_n=\frac{X_n}{m}\] can be treated as
an approximate uniform number in \([0,1]\).

\end{tcolorbox}

\begin{tcolorbox}[enhanced jigsaw, arc=.35mm, bottomrule=.15mm, bottomtitle=1mm, breakable, colback=white, colbacktitle=quarto-callout-important-color!10!white, colframe=quarto-callout-important-color-frame, coltitle=black, left=2mm, leftrule=.75mm, opacityback=0, opacitybacktitle=0.6, rightrule=.15mm, title=\textcolor{quarto-callout-important-color}{\faExclamation}\hspace{0.5em}{Remark}, titlerule=0mm, toprule=.15mm, toptitle=1mm]

\emph{Remark 17}. The mathematical details of pseudorandom number
generation are outside the scope of this course, but the quality of the
random number generator can affect Monte Carlo accuracy.

\end{tcolorbox}

\bookmarksetup{startatroot}

\chapter{Standard Distributions: Inverse-Transform
Method}\label{standard-distributions-inverse-transform-method}

This section develops the inverse-transform method for sampling from
nonuniform distributions.

\section{General inverse-transform
method}\label{general-inverse-transform-method}

Suppose \(Z\sim\operatorname{Uniform}(0,1)\) and suppose the target
random variable \(Y\) has CDF \(h(y)\). Define \[Y=h^{-1}(Z).\] Then
\(Y\) has CDF \(h\).

\begin{tcolorbox}[enhanced jigsaw, arc=.35mm, bottomrule=.15mm, bottomtitle=1mm, breakable, colback=white, colbacktitle=quarto-callout-note-color!10!white, colframe=quarto-callout-note-color-frame, coltitle=black, left=2mm, leftrule=.75mm, opacityback=0, opacitybacktitle=0.6, rightrule=.15mm, title=\textcolor{quarto-callout-note-color}{\faInfo}\hspace{0.5em}{Theorem}, titlerule=0mm, toprule=.15mm, toptitle=1mm]

\textbf{Theorem 18} (Inverse-transform sampling). \emph{Let
\(U\sim\operatorname{Uniform}(0,1)\) and let \(F\) be a continuous,
strictly increasing CDF. Define \[Y=F^{-1}(U).\] Then \(Y\) has CDF
\(F\).}

\end{tcolorbox}

\begin{tcolorbox}[enhanced jigsaw, arc=.35mm, bottomrule=.15mm, bottomtitle=1mm, breakable, colback=white, colbacktitle=quarto-callout-note-color!10!white, colframe=quarto-callout-note-color-frame, coltitle=black, left=2mm, leftrule=.75mm, opacityback=0, opacitybacktitle=0.6, rightrule=.15mm, title=\textcolor{quarto-callout-note-color}{\faInfo}\hspace{0.5em}{Solution}, titlerule=0mm, toprule=.15mm, toptitle=1mm]

For any \(y\), \[\mathbb{P}(Y\le y)=\mathbb{P}(F^{-1}(U)\le y).\] Since
\(F\) is increasing, this event is equivalent to \(U\le F(y)\).
Therefore, \[\mathbb{P}(Y\le y)=\mathbb{P}(U\le F(y))=F(y),\] because
\(U\) is uniform on \([0,1]\).

\end{tcolorbox}

\section{Example: exponential
distribution}\label{example-exponential-distribution}

Let \(Y\sim\operatorname{Exp}(\lambda)\) with density
\[p(y)=\lambda e^{-\lambda y},\qquad y\ge 0.\] Its CDF is
\[h(y)=1-e^{-\lambda y}.\] Set \(z=h(y)\). Then \[z=1-e^{-\lambda y}
\quad\Longrightarrow\quad
1-z=e^{-\lambda y}
\quad\Longrightarrow\quad
 y=-\frac{1}{\lambda}\log(1-z).\] Thus, if
\(Z\sim\operatorname{Uniform}(0,1)\), then
\[Y=-\frac{1}{\lambda}\log(1-Z)\sim\operatorname{Exp}(\lambda).\] Since
\(1-Z\sim\operatorname{Uniform}(0,1)\), one often writes
\[Y=-\frac{1}{\lambda}\log U.\]

\begin{tcolorbox}[enhanced jigsaw, arc=.35mm, bottomrule=.15mm, bottomtitle=1mm, breakable, colback=white, colbacktitle=quarto-callout-tip-color!10!white, colframe=quarto-callout-tip-color-frame, coltitle=black, left=2mm, leftrule=.75mm, opacityback=0, opacitybacktitle=0.6, rightrule=.15mm, title=\textcolor{quarto-callout-tip-color}{\faLightbulb}\hspace{0.5em}{Example}, titlerule=0mm, toprule=.15mm, toptitle=1mm]

\textbf{Example 19} (Exponential inverse transform). Generate samples
from \(\operatorname{Exp}(2)\) using uniform random variables.

\end{tcolorbox}

\begin{tcolorbox}[enhanced jigsaw, arc=.35mm, bottomrule=.15mm, bottomtitle=1mm, breakable, colback=white, colbacktitle=quarto-callout-note-color!10!white, colframe=quarto-callout-note-color-frame, coltitle=black, left=2mm, leftrule=.75mm, opacityback=0, opacitybacktitle=0.6, rightrule=.15mm, title=\textcolor{quarto-callout-note-color}{\faInfo}\hspace{0.5em}{Solution}, titlerule=0mm, toprule=.15mm, toptitle=1mm]

Let \(U\sim\operatorname{Uniform}(0,1)\). Since \(\lambda=2\),
\[Y=-\frac{1}{2}\log(1-U)\] has the exponential distribution with rate
\(2\).

\end{tcolorbox}

\section{Example: Cauchy
distribution}\label{example-cauchy-distribution}

The Cauchy distribution with location \(x_0\) and scale \(\gamma\) has
density
\[p(x;x_0,\gamma)=\frac{1}{\pi\gamma}\frac{1}{1+\left(\frac{x-x_0}{\gamma}\right)^2}.\]
Its CDF is
\[F(x;x_0,\gamma)=\frac{1}{\pi}\tan^{-1}\left(\frac{x-x_0}{\gamma}\right)+\frac{1}{2}.\]
Solving \(z=F(x)\) gives the inverse CDF
\[x=F^{-1}(z)=x_0+\gamma\tan\left(\pi\left(z-\frac{1}{2}\right)\right).\]
Therefore, if \(Z\sim\operatorname{Uniform}(0,1)\), then
\[X=x_0+\gamma\tan\left(\pi\left(Z-\frac{1}{2}\right)\right)\] follows a
\(\operatorname{Cauchy}(x_0,\gamma)\) distribution.

\begin{tcolorbox}[enhanced jigsaw, arc=.35mm, bottomrule=.15mm, bottomtitle=1mm, breakable, colback=white, colbacktitle=quarto-callout-tip-color!10!white, colframe=quarto-callout-tip-color-frame, coltitle=black, left=2mm, leftrule=.75mm, opacityback=0, opacitybacktitle=0.6, rightrule=.15mm, title=\textcolor{quarto-callout-tip-color}{\faLightbulb}\hspace{0.5em}{Example}, titlerule=0mm, toprule=.15mm, toptitle=1mm]

\textbf{Example 20} (Cauchy inverse transform). Generate samples from
the standard Cauchy distribution.

\end{tcolorbox}

\begin{tcolorbox}[enhanced jigsaw, arc=.35mm, bottomrule=.15mm, bottomtitle=1mm, breakable, colback=white, colbacktitle=quarto-callout-note-color!10!white, colframe=quarto-callout-note-color-frame, coltitle=black, left=2mm, leftrule=.75mm, opacityback=0, opacitybacktitle=0.6, rightrule=.15mm, title=\textcolor{quarto-callout-note-color}{\faInfo}\hspace{0.5em}{Solution}, titlerule=0mm, toprule=.15mm, toptitle=1mm]

For the standard Cauchy distribution, \(x_0=0\) and \(\gamma=1\). If
\(U\sim\operatorname{Uniform}(0,1)\), then
\[X=\tan\left(\pi\left(U-\frac12\right)\right)\] has the standard Cauchy
distribution.

\end{tcolorbox}

\begin{Shaded}
\begin{Highlighting}[]
\ImportTok{import}\NormalTok{ numpy }\ImportTok{as}\NormalTok{ np}
\ImportTok{import}\NormalTok{ matplotlib.pyplot }\ImportTok{as}\NormalTok{ plt}

\KeywordTok{def}\NormalTok{ cauchy\_inverse\_cdf(p, x0}\OperatorTok{=}\DecValTok{0}\NormalTok{, gamma}\OperatorTok{=}\DecValTok{1}\NormalTok{):}
    \ControlFlowTok{return}\NormalTok{ x0 }\OperatorTok{+}\NormalTok{ gamma }\OperatorTok{*}\NormalTok{ np.tan(np.pi }\OperatorTok{*}\NormalTok{ (p }\OperatorTok{{-}} \FloatTok{0.5}\NormalTok{))}

\NormalTok{p\_values }\OperatorTok{=}\NormalTok{ np.linspace(}\FloatTok{0.001}\NormalTok{, }\FloatTok{0.999}\NormalTok{, }\DecValTok{1000}\NormalTok{)}
\NormalTok{x\_values }\OperatorTok{=}\NormalTok{ cauchy\_inverse\_cdf(p\_values)}

\NormalTok{plt.figure(figsize}\OperatorTok{=}\NormalTok{(}\DecValTok{8}\NormalTok{, }\DecValTok{5}\NormalTok{))}
\NormalTok{plt.plot(p\_values, x\_values, label}\OperatorTok{=}\VerbatimStringTok{r"}\DecValTok{$}\VerbatimStringTok{x=x\_0}\OperatorTok{+}\ErrorTok{\textbackslash{}}\VerbatimStringTok{gamma}\CharTok{\textbackslash{}t}\VerbatimStringTok{an}\KeywordTok{(}\ErrorTok{\textbackslash{}}\VerbatimStringTok{pi}\KeywordTok{(}\VerbatimStringTok{p{-}0}\DecValTok{.}\VerbatimStringTok{5}\KeywordTok{))}\DecValTok{$}\VerbatimStringTok{"}\NormalTok{)}
\NormalTok{plt.axhline(}\DecValTok{0}\NormalTok{, linestyle}\OperatorTok{=}\StringTok{\textquotesingle{}{-}{-}\textquotesingle{}}\NormalTok{, linewidth}\OperatorTok{=}\FloatTok{0.8}\NormalTok{)}
\NormalTok{plt.axvline(}\FloatTok{0.5}\NormalTok{, linestyle}\OperatorTok{=}\StringTok{\textquotesingle{}{-}{-}\textquotesingle{}}\NormalTok{, linewidth}\OperatorTok{=}\FloatTok{0.8}\NormalTok{)}
\NormalTok{plt.xlabel(}\StringTok{"Probability p"}\NormalTok{)}
\NormalTok{plt.ylabel(}\StringTok{"Quantile x"}\NormalTok{)}
\NormalTok{plt.title(}\StringTok{"Inverse CDF of the Cauchy Distribution"}\NormalTok{)}
\NormalTok{plt.legend()}
\NormalTok{plt.grid()}
\NormalTok{plt.show()}
\end{Highlighting}
\end{Shaded}

\bookmarksetup{startatroot}

\chapter{Multiple Variables and the Box--Muller
Method}\label{multiple-variables-and-the-boxmuller-method}

This section extends random variable transformations to multiple
variables.

\section{Jacobian transformation}\label{jacobian-transformation}

Suppose \(u_1,
\ldots,u_M\sim\operatorname{Uniform}(0,1)\) and we define a
transformation from \(u\) variables to \(y\) variables. The density
transformation is controlled by the Jacobian:
\[p(y_1,\ldots,y_M)=p(u_1,\ldots,u_M)
\left|\frac{\partial(u_1,\ldots,u_M)}{\partial(y_1,\ldots,y_M)}\right|.\]

\section{Polar transformation of two normal
variables}\label{polar-transformation-of-two-normal-variables}

Let \(Z_1,Z_2\) be independent standard normal variables. Their joint
density is
\[f(z_1,z_2)=\frac{1}{2\pi}\exp\left(-\frac{z_1^2+z_2^2}{2}\right).\]
Transform to polar coordinates: \[Z_1=R\cos\Theta,
\qquad
Z_2=R\sin\Theta.\] Since \(R^2=Z_1^2+Z_2^2\) and the Jacobian is \(R\),
the joint density of \((R,\Theta)\) is
\[f_{R,\Theta}(r,\theta)=\frac{1}{2\pi}e^{-r^2/2}r,
\qquad r>0,\\ 0\le \theta<2\pi.\] Therefore,
\[\Theta\sim\operatorname{Uniform}(0,2\pi),\] and
\[f_R(r)=re^{-r^2/2},\qquad r>0.\] The CDF of \(R\) is
\[F_R(r)=1-e^{-r^2/2}.\] Thus, if
\(U_1,U_2\overset{\text{iid}}{\sim}\operatorname{Uniform}(0,1)\),
\[R=\sqrt{-2\log U_1},
\qquad
\Theta=2\pi U_2.\]

\section{Box--Muller method}\label{boxmuller-method}

The Box--Muller method converts two independent uniform random variables
into two independent standard normal random variables.

\begin{tcolorbox}[enhanced jigsaw, arc=.35mm, bottomrule=.15mm, bottomtitle=1mm, breakable, colback=white, colbacktitle=quarto-callout-note-color!10!white, colframe=quarto-callout-note-color-frame, coltitle=black, left=2mm, leftrule=.75mm, opacityback=0, opacitybacktitle=0.6, rightrule=.15mm, title=\textcolor{quarto-callout-note-color}{\faInfo}\hspace{0.5em}{Theorem}, titlerule=0mm, toprule=.15mm, toptitle=1mm]

\textbf{Theorem 21} (Box--Muller method). \emph{Let
\(U_1,U_2\overset{\text{iid}}{\sim}\operatorname{Uniform}(0,1)\). Define
\[Z_1=\sqrt{-2\log U_1}\cos(2\pi U_2),\]
\[Z_2=\sqrt{-2\log U_1}\sin(2\pi U_2).\] Then \(Z_1\) and \(Z_2\) are
independent standard normal random variables.}

\end{tcolorbox}

\begin{tcolorbox}[enhanced jigsaw, arc=.35mm, bottomrule=.15mm, bottomtitle=1mm, breakable, colback=white, colbacktitle=quarto-callout-note-color!10!white, colframe=quarto-callout-note-color-frame, coltitle=black, left=2mm, leftrule=.75mm, opacityback=0, opacitybacktitle=0.6, rightrule=.15mm, title=\textcolor{quarto-callout-note-color}{\faInfo}\hspace{0.5em}{Solution}, titlerule=0mm, toprule=.15mm, toptitle=1mm]

The construction uses polar coordinates. The angle \(\Theta=2\pi U_2\)
is uniform on \([0,2\pi]\). The radius \(R=\sqrt{-2\log U_1}\) has CDF
\[\mathbb{P}(R\le r)=\mathbb{P}\left(\sqrt{-2\log U_1}\le r\right)
=\mathbb{P}(U_1\ge e^{-r^2/2})=1-e^{-r^2/2}.\] This is exactly the
Rayleigh distribution needed for the radius of a two-dimensional
standard normal vector. Therefore \((Z_1,Z_2)\) has joint density
\[\frac{1}{2\pi}e^{-(z_1^2+z_2^2)/2},\] which factors into two standard
normal densities.

\end{tcolorbox}

\bookmarksetup{startatroot}

\chapter{Rejection Sampling}\label{rejection-sampling}

This section introduces rejection sampling, a method for sampling from a
complex target distribution using a simpler proposal distribution.

\section{Setup}\label{setup}

Suppose we want to sample from a target density \(p(z)\), but direct
sampling is difficult. Choose a proposal density \(q(z)\) from which it
is easy to sample. Assume there exists a constant \(k<\infty\) such that
\[p(z)\le kq(z)
\qquad\text{for all }z.\] The function \(kq(z)\) forms an envelope above
\(p(z)\).

\begin{tcolorbox}[enhanced jigsaw, arc=.35mm, bottomrule=.15mm, bottomtitle=1mm, breakable, colback=white, colbacktitle=quarto-callout-note-color!10!white, colframe=quarto-callout-note-color-frame, coltitle=black, left=2mm, leftrule=.75mm, opacityback=0, opacitybacktitle=0.6, rightrule=.15mm, title=\textcolor{quarto-callout-note-color}{\faInfo}\hspace{0.5em}{Definition}, titlerule=0mm, toprule=.15mm, toptitle=1mm]

\textbf{Definition 22} (Proposal distribution). In rejection sampling,
\(q(z)\) is called the proposal distribution because it proposes
candidate samples.

\end{tcolorbox}

\section{Algorithm}\label{algorithm}

~

\begin{tcolorbox}[enhanced jigsaw, arc=.35mm, bottomrule=.15mm, bottomtitle=1mm, breakable, colback=white, colbacktitle=quarto-callout-tip-color!10!white, colframe=quarto-callout-tip-color-frame, coltitle=black, left=2mm, leftrule=.75mm, opacityback=0, opacitybacktitle=0.6, rightrule=.15mm, title=\textcolor{quarto-callout-tip-color}{\faLightbulb}\hspace{0.5em}{Key idea}, titlerule=0mm, toprule=.15mm, toptitle=1mm]

Rejection sampling algorithm Repeat the following steps:

\begin{enumerate}
\def\labelenumi{\arabic{enumi}.}
\item
  Sample \(Z\sim q(z)\).
\item
  Sample \(U\sim\operatorname{Uniform}(0,1)\) independently.
\item
  Accept \(Z\) if \[U\le \frac{p(Z)}{kq(Z)}.\] Otherwise reject \(Z\)
  and try again.
\end{enumerate}

The accepted samples have density \(p(z)\).

\end{tcolorbox}

\begin{tcolorbox}[enhanced jigsaw, arc=.35mm, bottomrule=.15mm, bottomtitle=1mm, breakable, colback=white, colbacktitle=quarto-callout-tip-color!10!white, colframe=quarto-callout-tip-color-frame, coltitle=black, left=2mm, leftrule=.75mm, opacityback=0, opacitybacktitle=0.6, rightrule=.15mm, title=\textcolor{quarto-callout-tip-color}{\faLightbulb}\hspace{0.5em}{Example}, titlerule=0mm, toprule=.15mm, toptitle=1mm]

\textbf{Example 23} (Acceptance probability). Show that the overall
probability of acceptance is \(1/k\).

\end{tcolorbox}

\begin{tcolorbox}[enhanced jigsaw, arc=.35mm, bottomrule=.15mm, bottomtitle=1mm, breakable, colback=white, colbacktitle=quarto-callout-note-color!10!white, colframe=quarto-callout-note-color-frame, coltitle=black, left=2mm, leftrule=.75mm, opacityback=0, opacitybacktitle=0.6, rightrule=.15mm, title=\textcolor{quarto-callout-note-color}{\faInfo}\hspace{0.5em}{Solution}, titlerule=0mm, toprule=.15mm, toptitle=1mm]

Conditioning on the proposed value \(Z=z\),
\[\mathbb{P}(\text{accept}\mid Z=z)=\frac{p(z)}{kq(z)}.\] Therefore
\[\mathbb{P}(\text{accept})
=\int \mathbb{P}(\text{accept}\mid Z=z)q(z)\,\,dz
=\int \frac{p(z)}{kq(z)}q(z)\,\,dz
=\frac{1}{k}\int p(z)\,\,dz
=\frac{1}{k}.\]

\end{tcolorbox}

\section{Limitations}\label{limitations}

Rejection sampling is simple and exact in principle, but it may be
inefficient.

\begin{itemize}
\item
  It may be hard to find a finite and useful \(k\) such that
  \(p(z)\le kq(z)\) everywhere.
\item
  If \(k\) is too large, the acceptance probability \(1/k\) is too
  small.
\item
  In high-dimensional spaces, most proposed points may be rejected,
  making the method exponentially slow.
\end{itemize}

\bookmarksetup{startatroot}

\chapter{Importance Sampling}\label{importance-sampling}

This section introduces importance sampling, a method for estimating
expectations by sampling from a convenient proposal distribution.

\section{Basic importance sampling
identity}\label{basic-importance-sampling-identity}

Suppose we want \[\mathbb{E}_p[f(Z)] = \int f(z)p(z)\,\,dz,\] but
sampling from \(p\) is difficult. Choose a proposal density \(q\) such
that \(q(z)>0\) whenever \(p(z)f(z)\ne 0\). Then \[\mathbb{E}_p[f(Z)]
=\int f(z)\frac{p(z)}{q(z)}q(z)\,\,dz.\] If
\(Z^{(1)},\ldots,Z^{(L)}\overset{\text{iid}}{\sim}q\), then
\[\widehat{I}_L
=\frac{1}{L}\sum_{\ell=1}^L
\frac{p(Z^{(\ell)})}{q(Z^{(\ell)})}f(Z^{(\ell)})\] estimates
\(\mathbb{E}_p[f(Z)]\).

\begin{tcolorbox}[enhanced jigsaw, arc=.35mm, bottomrule=.15mm, bottomtitle=1mm, breakable, colback=white, colbacktitle=quarto-callout-note-color!10!white, colframe=quarto-callout-note-color-frame, coltitle=black, left=2mm, leftrule=.75mm, opacityback=0, opacitybacktitle=0.6, rightrule=.15mm, title=\textcolor{quarto-callout-note-color}{\faInfo}\hspace{0.5em}{Definition}, titlerule=0mm, toprule=.15mm, toptitle=1mm]

\textbf{Definition 24} (Importance weights). The ratios
\[r_\ell=\frac{p(Z^{(\ell)})}{q(Z^{(\ell)})}\] are called importance
weights. They correct the bias introduced by sampling from \(q\) instead
of \(p\).

\end{tcolorbox}

\begin{tcolorbox}[enhanced jigsaw, arc=.35mm, bottomrule=.15mm, bottomtitle=1mm, breakable, colback=white, colbacktitle=quarto-callout-important-color!10!white, colframe=quarto-callout-important-color-frame, coltitle=black, left=2mm, leftrule=.75mm, opacityback=0, opacitybacktitle=0.6, rightrule=.15mm, title=\textcolor{quarto-callout-important-color}{\faExclamation}\hspace{0.5em}{Remark}, titlerule=0mm, toprule=.15mm, toptitle=1mm]

\emph{Remark 25}. Unlike rejection sampling, importance sampling keeps
all generated samples. However, the estimator can have high variance if
the weights are unstable.

\end{tcolorbox}

\section{Self-normalized importance
sampling}\label{self-normalized-importance-sampling}

Often, especially in Bayesian statistics, the target density is known
only up to a normalizing constant:
\[p(z)=\frac{\widetilde{p}(z)}{Z_p}.\] Similarly, we may write
\[q(z)=\frac{\widetilde{q}(z)}{Z_q}.\] Then \[\mathbb{E}_p[f]
=\frac{Z_q}{Z_p}\int f(z)\frac{\widetilde{p}(z)}{\widetilde{q}(z)}q(z)\,\,dz.\]
The ratio of normalizing constants can be estimated by \[\frac{Z_p}{Z_q}
=\int \frac{\widetilde{p}(z)}{\widetilde{q}(z)}q(z)\,\,dz
\approx \frac{1}{L}\sum_{\ell=1}^L \widetilde{r}_\ell,\] where
\[\widetilde{r}_\ell=\frac{\widetilde{p}(Z^{(\ell)})}{\widetilde{q}(Z^{(\ell)})}.\]
Thus the self-normalized estimator is \[\widehat{I}_{SN}
=\sum_{\ell=1}^L w_\ell f(Z^{(\ell)}),\] where
\[w_\ell=\frac{\widetilde{r}_\ell}{\sum_{m=1}^L \widetilde{r}_m}
=\frac{\widetilde{p}(Z^{(\ell)})/\widetilde{q}(Z^{(\ell)})}{\sum_{m=1}^L \widetilde{p}(Z^{(m)})/\widetilde{q}(Z^{(m)})}.\]

\begin{tcolorbox}[enhanced jigsaw, arc=.35mm, bottomrule=.15mm, bottomtitle=1mm, breakable, colback=white, colbacktitle=quarto-callout-tip-color!10!white, colframe=quarto-callout-tip-color-frame, coltitle=black, left=2mm, leftrule=.75mm, opacityback=0, opacitybacktitle=0.6, rightrule=.15mm, title=\textcolor{quarto-callout-tip-color}{\faLightbulb}\hspace{0.5em}{Example}, titlerule=0mm, toprule=.15mm, toptitle=1mm]

\textbf{Example 26} (Importance sampling for rare events). Suppose \(p\)
puts very small probability in a region that contributes a lot to the
expectation. Explain why a proposal \(q\) can help.

\end{tcolorbox}

\begin{tcolorbox}[enhanced jigsaw, arc=.35mm, bottomrule=.15mm, bottomtitle=1mm, breakable, colback=white, colbacktitle=quarto-callout-note-color!10!white, colframe=quarto-callout-note-color-frame, coltitle=black, left=2mm, leftrule=.75mm, opacityback=0, opacitybacktitle=0.6, rightrule=.15mm, title=\textcolor{quarto-callout-note-color}{\faInfo}\hspace{0.5em}{Solution}, titlerule=0mm, toprule=.15mm, toptitle=1mm]

If direct samples from \(p\) rarely visit the important region, the
ordinary Monte Carlo average will have high variance. A proposal \(q\)
can intentionally sample more often from the important region. The
weight \(p(z)/q(z)\) then corrects the estimator so it still targets the
original expectation under \(p\).

\end{tcolorbox}

\bookmarksetup{startatroot}

\chapter{Sampling-Importance-Resampling}\label{sampling-importance-resampling}

This section explains how importance weights can be used to
approximately draw samples from a target distribution.

Rejection sampling requires a constant \(k\) satisfying
\(p(z)\le kq(z)\), which may be hard to find.
Sampling-importance-resampling avoids this envelope constant.

\begin{tcolorbox}[enhanced jigsaw, arc=.35mm, bottomrule=.15mm, bottomtitle=1mm, breakable, colback=white, colbacktitle=quarto-callout-tip-color!10!white, colframe=quarto-callout-tip-color-frame, coltitle=black, left=2mm, leftrule=.75mm, opacityback=0, opacitybacktitle=0.6, rightrule=.15mm, title=\textcolor{quarto-callout-tip-color}{\faLightbulb}\hspace{0.5em}{Key idea}, titlerule=0mm, toprule=.15mm, toptitle=1mm]

Sampling-importance-resampling (SIR)

\begin{enumerate}
\def\labelenumi{\arabic{enumi}.}
\item
  Draw proposal samples \(z^{(1)},\ldots,z^{(L)}\sim q(z)\).
\item
  Compute weights \[w_\ell\propto \frac{p(z^{(\ell)})}{q(z^{(\ell)})},
  \qquad \sum_{\ell=1}^L w_\ell=1.\]
\item
  Resample from the discrete set \(\{z^{(1)},\ldots,z^{(L)}\}\) with
  probabilities \(w_1,\ldots,w_L\).
\end{enumerate}

\end{tcolorbox}

The second-stage resampled values are approximate samples from the
target distribution \(p\).

\begin{tcolorbox}[enhanced jigsaw, arc=.35mm, bottomrule=.15mm, bottomtitle=1mm, breakable, colback=white, colbacktitle=quarto-callout-important-color!10!white, colframe=quarto-callout-important-color-frame, coltitle=black, left=2mm, leftrule=.75mm, opacityback=0, opacitybacktitle=0.6, rightrule=.15mm, title=\textcolor{quarto-callout-important-color}{\faExclamation}\hspace{0.5em}{Remark}, titlerule=0mm, toprule=.15mm, toptitle=1mm]

\emph{Remark 27}. SIR is conceptually useful because it transforms
weighted samples into unweighted samples, but it can suffer from weight
degeneracy if a small number of weights dominate.

\end{tcolorbox}

\bookmarksetup{startatroot}

\chapter{Summary, Advantages, and
Limitations}\label{summary-advantages-and-limitations}

This section summarizes the main Monte Carlo methods and their practical
strengths and weaknesses.

\section{Summary of methods}\label{summary-of-methods}

{\def\LTcaptype{none} % do not increment counter
\begin{longtable}[]{@{}ll@{}}
\toprule\noalign{}
\textbf{Method} & \textbf{Main idea} \\
\midrule\noalign{}
\endhead
\bottomrule\noalign{}
\endlastfoot
Monte Carlo simulation & Approximate probabilities or expectations by
repeated random experiments. \\
Monte Carlo integration & Rewrite an integral as an expectation and
approximate it by a sample average. \\
Inverse-transform sampling & Generate samples by applying an inverse CDF
to uniform random variables. \\
Box--Muller method & Generate two independent normal random variables
from two independent uniforms. \\
Rejection sampling & Sample from a proposal and accept/reject candidates
to obtain target samples. \\
Importance sampling & Sample from a proposal and correct using
weights. \\
Sampling-importance-resampling & Use importance weights to resample
approximate draws from the target. \\
\end{longtable}
}

\section{Advantages}\label{advantages}

Monte Carlo methods are widely used because they are often easy to
implement and broadly applicable.

\begin{itemize}
\item
  They can solve complex problems using simple random sampling.
\item
  They turn difficult integrals into averages.
\item
  Randomness becomes a tool rather than an obstacle.
\item
  They scale conceptually to high-dimensional statistical models.
\end{itemize}

\section{Limitations}\label{limitations-1}

Monte Carlo methods also have important limitations.

\begin{itemize}
\item
  They may require substantial computational power.
\item
  Accuracy depends on the number of samples and the variance of the
  estimator.
\item
  Results can be sensitive to the chosen sampling distribution.
\item
  Basic Monte Carlo can perform poorly for high-dimensional integration
  or rare-event problems.
\item
  Rejection sampling may be inefficient when acceptance rates are small.
\item
  Importance sampling may be unstable when weights have high variance.
\end{itemize}

\begin{tcolorbox}[enhanced jigsaw, arc=.35mm, bottomrule=.15mm, bottomtitle=1mm, breakable, colback=white, colbacktitle=quarto-callout-warning-color!10!white, colframe=quarto-callout-warning-color-frame, coltitle=black, left=2mm, leftrule=.75mm, opacityback=0, opacitybacktitle=0.6, rightrule=.15mm, title=\textcolor{quarto-callout-warning-color}{\faExclamationTriangle}\hspace{0.5em}{Warning}, titlerule=0mm, toprule=.15mm, toptitle=1mm]

Alan Sokal's warning ``Monte Carlo is an extremely bad method; it should
be used only when all alternative methods are worse.''

This quote emphasizes that Monte Carlo is powerful, but not magic. It is
most useful when exact calculation, deterministic integration, or direct
optimization is impractical.

\end{tcolorbox}

\bookmarksetup{startatroot}

\chapter{Additional Practice
Problems}\label{additional-practice-problems-2}

This section gives additional exercises to reinforce the main ideas of
Monte Carlo sampling.

\begin{tcolorbox}[enhanced jigsaw, arc=.35mm, bottomrule=.15mm, bottomtitle=1mm, breakable, colback=white, colbacktitle=quarto-callout-warning-color!10!white, colframe=quarto-callout-warning-color-frame, coltitle=black, left=2mm, leftrule=.75mm, opacityback=0, opacitybacktitle=0.6, rightrule=.15mm, title=\textcolor{quarto-callout-warning-color}{\faExclamationTriangle}\hspace{0.5em}{Practice Problem}, titlerule=0mm, toprule=.15mm, toptitle=1mm]

\textbf{Practice Problem 28} (Variance of the Monte Carlo estimator).
Let \(X_1,\ldots,X_N\overset{\text{iid}}{\sim}p\) and define
\[\widehat{I}_N=\frac{1}{N}\sum_{i=1}^N f(X_i).\] Assume
\(\operatorname{Var}(f(X))<\infty\). Find \(\mathbb{E}[\widehat{I}_N]\)
and \(\operatorname{Var}(\widehat{I}_N)\).

\end{tcolorbox}

\begin{tcolorbox}[enhanced jigsaw, arc=.35mm, bottomrule=.15mm, bottomtitle=1mm, breakable, colback=white, colbacktitle=quarto-callout-note-color!10!white, colframe=quarto-callout-note-color-frame, coltitle=black, left=2mm, leftrule=.75mm, opacityback=0, opacitybacktitle=0.6, rightrule=.15mm, title=\textcolor{quarto-callout-note-color}{\faInfo}\hspace{0.5em}{Solution}, titlerule=0mm, toprule=.15mm, toptitle=1mm]

By linearity of expectation,
\[\mathbb{E}[\widehat{I}_N]=\frac{1}{N}\sum_{i=1}^N \mathbb{E}[f(X_i)]=\mathbb{E}[f(X)].\]
Thus the estimator is unbiased for \(I=\mathbb{E}[f(X)]\).

Because the samples are independent,
\[\operatorname{Var}(\widehat{I}_N)=\frac{1}{N^2}\sum_{i=1}^N \operatorname{Var}(f(X_i))
=\frac{\operatorname{Var}(f(X))}{N}.\] Therefore the standard deviation
of the estimator decreases at rate \(1/\sqrt{N}\).

\end{tcolorbox}

\begin{tcolorbox}[enhanced jigsaw, arc=.35mm, bottomrule=.15mm, bottomtitle=1mm, breakable, colback=white, colbacktitle=quarto-callout-warning-color!10!white, colframe=quarto-callout-warning-color-frame, coltitle=black, left=2mm, leftrule=.75mm, opacityback=0, opacitybacktitle=0.6, rightrule=.15mm, title=\textcolor{quarto-callout-warning-color}{\faExclamationTriangle}\hspace{0.5em}{Practice Problem}, titlerule=0mm, toprule=.15mm, toptitle=1mm]

\textbf{Practice Problem 29} (Inverse transform for a power
distribution). Suppose \(Y\) has CDF \(F(y)=y^\alpha\) for
\(0\le y\le 1\), where \(\alpha>0\). Use inverse-transform sampling to
generate \(Y\).

\end{tcolorbox}

\begin{tcolorbox}[enhanced jigsaw, arc=.35mm, bottomrule=.15mm, bottomtitle=1mm, breakable, colback=white, colbacktitle=quarto-callout-note-color!10!white, colframe=quarto-callout-note-color-frame, coltitle=black, left=2mm, leftrule=.75mm, opacityback=0, opacitybacktitle=0.6, rightrule=.15mm, title=\textcolor{quarto-callout-note-color}{\faInfo}\hspace{0.5em}{Solution}, titlerule=0mm, toprule=.15mm, toptitle=1mm]

Let \(U\sim\operatorname{Uniform}(0,1)\). Solve \[U=F(y)=y^\alpha.\]
Then \[y=U^{1/\alpha}.\] Therefore \[Y=U^{1/\alpha}\] has CDF
\(F(y)=y^\alpha\) on \([0,1]\).

\end{tcolorbox}

\begin{tcolorbox}[enhanced jigsaw, arc=.35mm, bottomrule=.15mm, bottomtitle=1mm, breakable, colback=white, colbacktitle=quarto-callout-warning-color!10!white, colframe=quarto-callout-warning-color-frame, coltitle=black, left=2mm, leftrule=.75mm, opacityback=0, opacitybacktitle=0.6, rightrule=.15mm, title=\textcolor{quarto-callout-warning-color}{\faExclamationTriangle}\hspace{0.5em}{Practice Problem}, titlerule=0mm, toprule=.15mm, toptitle=1mm]

\textbf{Practice Problem 30} (Importance sampling identity). Let \(p\)
and \(q\) be densities with \(p(z)>0\) only where \(q(z)>0\). Show that
\[\int f(z)p(z)\,\,dz
=\mathbb{E}_q\left[f(Z)\frac{p(Z)}{q(Z)}\right].\]

\end{tcolorbox}

\begin{tcolorbox}[enhanced jigsaw, arc=.35mm, bottomrule=.15mm, bottomtitle=1mm, breakable, colback=white, colbacktitle=quarto-callout-note-color!10!white, colframe=quarto-callout-note-color-frame, coltitle=black, left=2mm, leftrule=.75mm, opacityback=0, opacitybacktitle=0.6, rightrule=.15mm, title=\textcolor{quarto-callout-note-color}{\faInfo}\hspace{0.5em}{Solution}, titlerule=0mm, toprule=.15mm, toptitle=1mm]

By definition of expectation under \(q\),
\[\mathbb{E}_q\left[f(Z)\frac{p(Z)}{q(Z)}\right]
=\int f(z)\frac{p(z)}{q(z)}q(z)\,\,dz
=\int f(z)p(z)\,\,dz.\] This is the basic identity behind importance
sampling.

\end{tcolorbox}

\begin{tcolorbox}[enhanced jigsaw, arc=.35mm, bottomrule=.15mm, bottomtitle=1mm, breakable, colback=white, colbacktitle=quarto-callout-warning-color!10!white, colframe=quarto-callout-warning-color-frame, coltitle=black, left=2mm, leftrule=.75mm, opacityback=0, opacitybacktitle=0.6, rightrule=.15mm, title=\textcolor{quarto-callout-warning-color}{\faExclamationTriangle}\hspace{0.5em}{Practice Problem}, titlerule=0mm, toprule=.15mm, toptitle=1mm]

\textbf{Practice Problem 31} (Acceptance rate). In rejection sampling,
suppose \(p(z)\le kq(z)\) and \(k=5\). What is the expected acceptance
probability? Approximately how many proposal samples are needed to
obtain \(1000\) accepted samples?

\end{tcolorbox}

\begin{tcolorbox}[enhanced jigsaw, arc=.35mm, bottomrule=.15mm, bottomtitle=1mm, breakable, colback=white, colbacktitle=quarto-callout-note-color!10!white, colframe=quarto-callout-note-color-frame, coltitle=black, left=2mm, leftrule=.75mm, opacityback=0, opacitybacktitle=0.6, rightrule=.15mm, title=\textcolor{quarto-callout-note-color}{\faInfo}\hspace{0.5em}{Solution}, titlerule=0mm, toprule=.15mm, toptitle=1mm]

The acceptance probability is
\[\mathbb{P}(\text{accept})=\frac{1}{k}=\frac{1}{5}=0.2.\] To obtain
about \(1000\) accepted samples, we need approximately
\[\frac{1000}{0.2}=5000\] proposal samples.

\end{tcolorbox}

\bookmarksetup{startatroot}

\chapter{Chapter 11: Sufficient Statistics and the Likelihood
Principle}\label{chapter-11-sufficient-statistics-and-the-likelihood-principle}

This chapter introduces sufficient statistics as a formal data-reduction
principle for statistical inference. The main goal is to identify when a
statistic keeps all information in the sample about an unknown
parameter. The chapter also introduces the likelihood function and the
likelihood principle.

\begin{tcolorbox}[enhanced jigsaw, arc=.35mm, bottomrule=.15mm, bottomtitle=1mm, breakable, colback=white, colbacktitle=quarto-callout-note-color!10!white, colframe=quarto-callout-note-color-frame, coltitle=black, left=2mm, leftrule=.75mm, opacityback=0, opacitybacktitle=0.6, rightrule=.15mm, title=\textcolor{quarto-callout-note-color}{\faInfo}\hspace{0.5em}{Topics}, titlerule=0mm, toprule=.15mm, toptitle=1mm]

Statistics and estimators; heuristic and mathematical definitions of
sufficient statistics; the factorization theorem; Bernoulli, normal,
uniform, and gamma examples; minimal and complete sufficient statistics;
likelihood functions; likelihood principle; binomial versus negative
binomial coin-tossing example.

\end{tcolorbox}

\bookmarksetup{startatroot}

\chapter{Overview}\label{overview-8}

This section introduces sufficient statistics, which formalize the idea
of reducing a data set without losing information about an unknown
parameter.

In statistical inference, the data set may contain many observations,
but often the information relevant to a parameter can be summarized by a
lower-dimensional statistic. For example, for Bernoulli data, the total
number of successes contains all the information about the success
probability. For normal data with known variance, the sample mean
contains all the information about the unknown mean.

\begin{tcolorbox}[enhanced jigsaw, arc=.35mm, bottomrule=.15mm, bottomtitle=1mm, breakable, colback=white, colbacktitle=quarto-callout-tip-color!10!white, colframe=quarto-callout-tip-color-frame, coltitle=black, left=2mm, leftrule=.75mm, opacityback=0, opacitybacktitle=0.6, rightrule=.15mm, title=\textcolor{quarto-callout-tip-color}{\faLightbulb}\hspace{0.5em}{Key idea}, titlerule=0mm, toprule=.15mm, toptitle=1mm]

A statistic is \emph{sufficient} for a parameter if, after the statistic
is known, the remaining details of the sample no longer contain
information about the parameter.
\[\text{full data} \quad \longrightarrow \quad \text{sufficient statistic} \quad \longrightarrow \quad \text{inference about } \theta.\]

\end{tcolorbox}

The main tool in this section is the \textbf{factorization theorem}. It
gives an efficient way to check sufficiency by factoring the joint
density or joint mass function into two parts: one part involving the
parameter only through the statistic, and one part independent of the
parameter.

\bookmarksetup{startatroot}

\chapter{Statistics and Estimators}\label{statistics-and-estimators}

This section recalls the basic language of statistics before defining
sufficient statistics.

\begin{tcolorbox}[enhanced jigsaw, arc=.35mm, bottomrule=.15mm, bottomtitle=1mm, breakable, colback=white, colbacktitle=quarto-callout-note-color!10!white, colframe=quarto-callout-note-color-frame, coltitle=black, left=2mm, leftrule=.75mm, opacityback=0, opacitybacktitle=0.6, rightrule=.15mm, title=\textcolor{quarto-callout-note-color}{\faInfo}\hspace{0.5em}{Definition}, titlerule=0mm, toprule=.15mm, toptitle=1mm]

\textbf{Definition 1} (Statistic). Let \(X_1,\ldots,X_n\) be a random
sample from a population. A \textbf{statistic} is any function of the
sample, \[T=T(X_1,\ldots,X_n),\] that does not depend on unknown
parameters.

\end{tcolorbox}

A statistic is itself a random variable, because it is computed from
random data. After the data are observed, the statistic takes a
numerical value.

\begin{tcolorbox}[enhanced jigsaw, arc=.35mm, bottomrule=.15mm, bottomtitle=1mm, breakable, colback=white, colbacktitle=quarto-callout-note-color!10!white, colframe=quarto-callout-note-color-frame, coltitle=black, left=2mm, leftrule=.75mm, opacityback=0, opacitybacktitle=0.6, rightrule=.15mm, title=\textcolor{quarto-callout-note-color}{\faInfo}\hspace{0.5em}{Example}, titlerule=0mm, toprule=.15mm, toptitle=1mm]

\textbf{Example 2} (Common statistics). Let \(X_1,\ldots,X_n\) be a
random sample.

\begin{enumerate}
\def\labelenumi{\arabic{enumi}.}
\item
  The sample mean \[\overline X=\frac{1}{n}\sum_{i=1}^n X_i\] is a
  statistic.
\item
  The sample variance
  \[S^2=\frac{1}{n-1}\sum_{i=1}^n (X_i-\overline X)^2\] is a statistic.
\item
  The maximum \[X_{(n)}=\max\{X_1,\ldots,X_n\}\] is a statistic.
\item
  The constant statistic \(T=3\) is also technically a statistic,
  although it ignores the data.
\end{enumerate}

\end{tcolorbox}

\begin{tcolorbox}[enhanced jigsaw, arc=.35mm, bottomrule=.15mm, bottomtitle=1mm, breakable, colback=white, colbacktitle=quarto-callout-tip-color!10!white, colframe=quarto-callout-tip-color-frame, coltitle=black, left=2mm, leftrule=.75mm, opacityback=0, opacitybacktitle=0.6, rightrule=.15mm, title=\textcolor{quarto-callout-tip-color}{\faLightbulb}\hspace{0.5em}{Solution}, titlerule=0mm, toprule=.15mm, toptitle=1mm]

Each quantity listed is a function of \(X_1,\ldots,X_n\) and does not
use an unknown parameter. Therefore each is a statistic. The constant
statistic is mathematically allowed, but it is usually not useful
because it does not summarize any sample information.

\end{tcolorbox}

\begin{tcolorbox}[enhanced jigsaw, arc=.35mm, bottomrule=.15mm, bottomtitle=1mm, breakable, colback=white, colbacktitle=quarto-callout-note-color!10!white, colframe=quarto-callout-note-color-frame, coltitle=black, left=2mm, leftrule=.75mm, opacityback=0, opacitybacktitle=0.6, rightrule=.15mm, title=\textcolor{quarto-callout-note-color}{\faInfo}\hspace{0.5em}{Definition}, titlerule=0mm, toprule=.15mm, toptitle=1mm]

\textbf{Definition 3} (Estimator). A statistic \(T(X_1,\ldots,X_n)\) is
called an \textbf{estimator} of a population parameter \(\theta\) if it
is used to estimate \(\theta\).

\end{tcolorbox}

For example, \(\overline X\) is commonly used to estimate the population
mean \(\mu\), and \(S^2\) is commonly used to estimate the population
variance \(\sigma^2\).

\bookmarksetup{startatroot}

\chapter{Heuristic Meaning of Sufficient
Statistics}\label{heuristic-meaning-of-sufficient-statistics}

This section explains the intuitive meaning of sufficiency before giving
the formal definition.

A sufficient statistic is a function of the data that captures all the
information needed to estimate a parameter \(\theta\). Once the
sufficient statistic is known, the rest of the data provides no
additional information about \(\theta\).

\begin{tcolorbox}[enhanced jigsaw, arc=.35mm, bottomrule=.15mm, bottomtitle=1mm, breakable, colback=white, colbacktitle=quarto-callout-note-color!10!white, colframe=quarto-callout-note-color-frame, coltitle=black, left=2mm, leftrule=.75mm, opacityback=0, opacitybacktitle=0.6, rightrule=.15mm, title=\textcolor{quarto-callout-note-color}{\faInfo}\hspace{0.5em}{Definition}, titlerule=0mm, toprule=.15mm, toptitle=1mm]

\textbf{Definition 4} (Heuristic definition of sufficiency). Let
\(X_1,\ldots,X_n\) be a random sample from a model with unknown
parameter \(\theta\). A statistic \[T=T(X_1,\ldots,X_n)\] is
\textbf{sufficient} for \(\theta\) if a statistician who knows only
\(T\) can do just as well for inference about \(\theta\) as a
statistician who knows the full data set \((X_1,\ldots,X_n)\).

\end{tcolorbox}

\begin{tcolorbox}[enhanced jigsaw, arc=.35mm, bottomrule=.15mm, bottomtitle=1mm, breakable, colback=white, colbacktitle=quarto-callout-tip-color!10!white, colframe=quarto-callout-tip-color-frame, coltitle=black, left=2mm, leftrule=.75mm, opacityback=0, opacitybacktitle=0.6, rightrule=.15mm, title=\textcolor{quarto-callout-tip-color}{\faLightbulb}\hspace{0.5em}{Key idea}, titlerule=0mm, toprule=.15mm, toptitle=1mm]

Sufficiency is a data-reduction principle:
\[(X_1,\ldots,X_n) \quad \leadsto \quad T(X_1,\ldots,X_n).\] If \(T\) is
sufficient, then this reduction does not lose information about
\(\theta\).

\end{tcolorbox}

\begin{tcolorbox}[enhanced jigsaw, arc=.35mm, bottomrule=.15mm, bottomtitle=1mm, breakable, colback=white, colbacktitle=quarto-callout-note-color!10!white, colframe=quarto-callout-note-color-frame, coltitle=black, left=2mm, leftrule=.75mm, opacityback=0, opacitybacktitle=0.6, rightrule=.15mm, title=\textcolor{quarto-callout-note-color}{\faInfo}\hspace{0.5em}{Example}, titlerule=0mm, toprule=.15mm, toptitle=1mm]

\textbf{Example 5} (Bernoulli intuition). Suppose
\(X_1,\ldots,X_n \overset{\text{iid}}{\sim}\operatorname{Bernoulli}(p)\).
Each \(X_i\) is \(1\) for success and \(0\) for failure. If the goal is
to learn \(p\), the order of successes and failures does not matter. The
total number of successes \[T=\sum_{i=1}^n X_i\] contains all
information about \(p\).

\end{tcolorbox}

\begin{tcolorbox}[enhanced jigsaw, arc=.35mm, bottomrule=.15mm, bottomtitle=1mm, breakable, colback=white, colbacktitle=quarto-callout-tip-color!10!white, colframe=quarto-callout-tip-color-frame, coltitle=black, left=2mm, leftrule=.75mm, opacityback=0, opacitybacktitle=0.6, rightrule=.15mm, title=\textcolor{quarto-callout-tip-color}{\faLightbulb}\hspace{0.5em}{Solution}, titlerule=0mm, toprule=.15mm, toptitle=1mm]

For example, the sequences
\[(1,1,0,1,0) \quad \text{and} \quad (0,1,1,0,1)\] both have three
successes out of five trials. Their probabilities under the Bernoulli
model are both proportional to \[p^3(1-p)^2.\] For inference about
\(p\), the relevant information is the number of successes, not the
order in which they occurred.

\end{tcolorbox}

\bookmarksetup{startatroot}

\chapter{Mathematical Definition of Sufficient
Statistics}\label{mathematical-definition-of-sufficient-statistics}

This section gives the formal conditional-distribution definition of
sufficiency.

\begin{tcolorbox}[enhanced jigsaw, arc=.35mm, bottomrule=.15mm, bottomtitle=1mm, breakable, colback=white, colbacktitle=quarto-callout-note-color!10!white, colframe=quarto-callout-note-color-frame, coltitle=black, left=2mm, leftrule=.75mm, opacityback=0, opacitybacktitle=0.6, rightrule=.15mm, title=\textcolor{quarto-callout-note-color}{\faInfo}\hspace{0.5em}{Definition}, titlerule=0mm, toprule=.15mm, toptitle=1mm]

\textbf{Definition 6} (Sufficient statistic). Let \(X_1,\ldots,X_n\) be
a random sample from a distribution with parameter \(\theta\), and let
\[T=T(X_1,\ldots,X_n)\] be a statistic. We say that \(T\) is
\textbf{sufficient} for \(\theta\) if, for every possible value \(t\),
the conditional distribution of the full sample
\[(X_1,\ldots,X_n) \mid T=t\] does not depend on \(\theta\).

\end{tcolorbox}

In words, once \(T=t\) is known, the remaining randomness in the sample
does not involve the unknown parameter. Therefore the remaining details
of the sample cannot help us learn about \(\theta\).

\begin{tcolorbox}[enhanced jigsaw, arc=.35mm, bottomrule=.15mm, bottomtitle=1mm, breakable, colback=white, colbacktitle=quarto-callout-important-color!10!white, colframe=quarto-callout-important-color-frame, coltitle=black, left=2mm, leftrule=.75mm, opacityback=0, opacitybacktitle=0.6, rightrule=.15mm, title=\textcolor{quarto-callout-important-color}{\faExclamation}\hspace{0.5em}{Remark}, titlerule=0mm, toprule=.15mm, toptitle=1mm]

\emph{Remark 7}. The conditional-distribution definition is conceptually
important, but it can be difficult to check directly. The factorization
theorem gives a much easier method.

\end{tcolorbox}

\bookmarksetup{startatroot}

\chapter{The Factorization Theorem}\label{the-factorization-theorem}

This section introduces the main theorem for checking sufficiency.

\begin{tcolorbox}[enhanced jigsaw, arc=.35mm, bottomrule=.15mm, bottomtitle=1mm, breakable, colback=white, colbacktitle=quarto-callout-important-color!10!white, colframe=quarto-callout-important-color-frame, coltitle=black, left=2mm, leftrule=.75mm, opacityback=0, opacitybacktitle=0.6, rightrule=.15mm, title=\textcolor{quarto-callout-important-color}{\faExclamation}\hspace{0.5em}{Theorem}, titlerule=0mm, toprule=.15mm, toptitle=1mm]

\textbf{Theorem 8} (Neyman--Fisher factorization theorem). \emph{Let
\(X_1,\ldots,X_n\) be a random sample with joint density or joint mass
function \[f(x_1,\ldots,x_n\mid\theta).\] A statistic
\[T=T(X_1,\ldots,X_n)\] is sufficient for \(\theta\) if and only if
there exist functions \(g\) and \(h\) such that, for all sample points
\(x=(x_1,\ldots,x_n)\) and all parameter values \(\theta\),
\[f(x_1,\ldots,x_n\mid\theta)
= g(T(x_1,\ldots,x_n),\theta)\,h(x_1,\ldots,x_n).\] Here:}

\begin{itemize}
\item
  \emph{\(g(T(x),\theta)\) depends on the data only through \(T(x)\) and
  may depend on \(\theta\);}
\item
  \emph{\(h(x)\) may depend on the full data but does not depend on
  \(\theta\).}
\end{itemize}

\end{tcolorbox}

\begin{tcolorbox}[enhanced jigsaw, arc=.35mm, bottomrule=.15mm, bottomtitle=1mm, breakable, colback=white, colbacktitle=quarto-callout-tip-color!10!white, colframe=quarto-callout-tip-color-frame, coltitle=black, left=2mm, leftrule=.75mm, opacityback=0, opacitybacktitle=0.6, rightrule=.15mm, title=\textcolor{quarto-callout-tip-color}{\faLightbulb}\hspace{0.5em}{Key idea}, titlerule=0mm, toprule=.15mm, toptitle=1mm]

To prove that \(T\) is sufficient:

\begin{enumerate}
\def\labelenumi{\arabic{enumi}.}
\item
  Write the joint density or joint pmf of the sample.
\item
  Rearrange it so that every parameter-dependent term involving the data
  appears only through \(T\).
\item
  Put the remaining parameter-free data terms into \(h(x)\).
\end{enumerate}

\end{tcolorbox}

\bookmarksetup{startatroot}

\chapter{Example: Bernoulli
Distribution}\label{example-bernoulli-distribution}

This section shows the most basic example: the number of successes is
sufficient for the Bernoulli parameter.

\begin{tcolorbox}[enhanced jigsaw, arc=.35mm, bottomrule=.15mm, bottomtitle=1mm, breakable, colback=white, colbacktitle=quarto-callout-note-color!10!white, colframe=quarto-callout-note-color-frame, coltitle=black, left=2mm, leftrule=.75mm, opacityback=0, opacitybacktitle=0.6, rightrule=.15mm, title=\textcolor{quarto-callout-note-color}{\faInfo}\hspace{0.5em}{Example}, titlerule=0mm, toprule=.15mm, toptitle=1mm]

\textbf{Example 9} (Bernoulli distribution). Let
\[X_1,\ldots,X_n \overset{\text{iid}}{\sim}\operatorname{Bernoulli}(p), \qquad 0<p<1.\]
Show that \[T(X)=\sum_{i=1}^n X_i\] is sufficient for \(p\).

\end{tcolorbox}

\begin{tcolorbox}[enhanced jigsaw, arc=.35mm, bottomrule=.15mm, bottomtitle=1mm, breakable, colback=white, colbacktitle=quarto-callout-tip-color!10!white, colframe=quarto-callout-tip-color-frame, coltitle=black, left=2mm, leftrule=.75mm, opacityback=0, opacitybacktitle=0.6, rightrule=.15mm, title=\textcolor{quarto-callout-tip-color}{\faLightbulb}\hspace{0.5em}{Solution}, titlerule=0mm, toprule=.15mm, toptitle=1mm]

The joint pmf is \[f(x_1,\ldots,x_n\mid p)
=\prod_{i=1}^n p^{x_i}(1-p)^{1-x_i},\] where each \(x_i\in\{0,1\}\).
Therefore \[f(x_1,\ldots,x_n\mid p)
=p^{\sum_{i=1}^n x_i}(1-p)^{n-\sum_{i=1}^n x_i}.\] Let
\[T(x)=\sum_{i=1}^n x_i.\] Then \[f(x_1,\ldots,x_n\mid p)
=\underbrace{p^{T(x)}(1-p)^{n-T(x)}}_{g(T(x),p)}\underbrace{1}_{h(x)}.\]
By the factorization theorem, \(T(X)=\sum_{i=1}^n X_i\) is sufficient
for \(p\).

\end{tcolorbox}

\begin{tcolorbox}[enhanced jigsaw, arc=.35mm, bottomrule=.15mm, bottomtitle=1mm, breakable, colback=white, colbacktitle=quarto-callout-important-color!10!white, colframe=quarto-callout-important-color-frame, coltitle=black, left=2mm, leftrule=.75mm, opacityback=0, opacitybacktitle=0.6, rightrule=.15mm, title=\textcolor{quarto-callout-important-color}{\faExclamation}\hspace{0.5em}{Remark}, titlerule=0mm, toprule=.15mm, toptitle=1mm]

\emph{Remark 10}. The sample mean \(\overline X=T/n\) is also sufficient
for \(p\), because it contains exactly the same information as \(T\).

\end{tcolorbox}

\bookmarksetup{startatroot}

\chapter{Example: Normal Distribution with Known
Variance}\label{example-normal-distribution-with-known-variance}

This section shows that, for normal data with known variance, the sample
mean is sufficient for the unknown mean.

\begin{tcolorbox}[enhanced jigsaw, arc=.35mm, bottomrule=.15mm, bottomtitle=1mm, breakable, colback=white, colbacktitle=quarto-callout-note-color!10!white, colframe=quarto-callout-note-color-frame, coltitle=black, left=2mm, leftrule=.75mm, opacityback=0, opacitybacktitle=0.6, rightrule=.15mm, title=\textcolor{quarto-callout-note-color}{\faInfo}\hspace{0.5em}{Example}, titlerule=0mm, toprule=.15mm, toptitle=1mm]

\textbf{Example 11} (Normal distribution with known variance). Let
\[X_1,\ldots,X_n \overset{\text{iid}}{\sim}\operatorname{Normal}(\mu,\sigma^2),\]
where \(\sigma^2\) is known and \(\mu\) is unknown. Show that
\[T(X)=\sum_{i=1}^n X_i\] is sufficient for \(\mu\).

\end{tcolorbox}

\begin{tcolorbox}[enhanced jigsaw, arc=.35mm, bottomrule=.15mm, bottomtitle=1mm, breakable, colback=white, colbacktitle=quarto-callout-tip-color!10!white, colframe=quarto-callout-tip-color-frame, coltitle=black, left=2mm, leftrule=.75mm, opacityback=0, opacitybacktitle=0.6, rightrule=.15mm, title=\textcolor{quarto-callout-tip-color}{\faLightbulb}\hspace{0.5em}{Solution}, titlerule=0mm, toprule=.15mm, toptitle=1mm]

The joint density is \[f(x_1,\ldots,x_n\mid\mu)
=\prod_{i=1}^n \frac{1}{\sqrt{2\pi\sigma^2}}\exp\left\{-\frac{(x_i-\mu)^2}{2\sigma^2}\right\}.\]
Thus \[f(x_1,\ldots,x_n\mid\mu)
=(2\pi)^{-n/2}\sigma^{-n}
\exp\left\{-\frac{1}{2\sigma^2}\sum_{i=1}^n (x_i-\mu)^2\right\}.\]
Expand the quadratic term: \[\sum_{i=1}^n (x_i-\mu)^2
=\sum_{i=1}^n x_i^2-2\mu\sum_{i=1}^n x_i+n\mu^2.\] Therefore
\[\begin{aligned}
f(x_1,\ldots,x_n\mid\mu)
&=(2\pi)^{-n/2}\sigma^{-n}
\exp\left\{-\frac{1}{2\sigma^2}\sum_{i=1}^n x_i^2
+\frac{\mu}{\sigma^2}\sum_{i=1}^n x_i
-\frac{n\mu^2}{2\sigma^2}\right\} \\
&=\underbrace{\exp\left\{\frac{\mu}{\sigma^2}T(x)-\frac{n\mu^2}{2\sigma^2}\right\}}_{g(T(x),\mu)}
\underbrace{(2\pi)^{-n/2}\sigma^{-n}\exp\left\{-\frac{1}{2\sigma^2}\sum_{i=1}^n x_i^2\right\}}_{h(x)}.
\end{aligned}\] The function \(h(x)\) does not depend on \(\mu\), and
the function \(g\) depends on the data only through \(T(x)=\sum_i x_i\).
Hence \(T(X)=\sum_i X_i\) is sufficient for \(\mu\).

Since \(\overline X=T/n\) is a one-to-one function of \(T\), the sample
mean \(\overline X\) is also sufficient for \(\mu\).

\end{tcolorbox}

\bookmarksetup{startatroot}

\chapter{Example: Uniform Population with Unknown Upper
Bound}\label{example-uniform-population-with-unknown-upper-bound}

This section illustrates an important point: the sufficient statistic
may not be the sample mean.

\begin{tcolorbox}[enhanced jigsaw, arc=.35mm, bottomrule=.15mm, bottomtitle=1mm, breakable, colback=white, colbacktitle=quarto-callout-note-color!10!white, colframe=quarto-callout-note-color-frame, coltitle=black, left=2mm, leftrule=.75mm, opacityback=0, opacitybacktitle=0.6, rightrule=.15mm, title=\textcolor{quarto-callout-note-color}{\faInfo}\hspace{0.5em}{Example}, titlerule=0mm, toprule=.15mm, toptitle=1mm]

\textbf{Example 12} (Uniform population). Suppose
\[X_1,\ldots,X_n \overset{\text{iid}}{\sim}\operatorname{Uniform}(0,\theta),\]
where \(\theta>0\) is unknown. Show that
\[T(X)=\max\{X_1,\ldots,X_n\}=X_{(n)}\] is sufficient for \(\theta\).

\end{tcolorbox}

\begin{tcolorbox}[enhanced jigsaw, arc=.35mm, bottomrule=.15mm, bottomtitle=1mm, breakable, colback=white, colbacktitle=quarto-callout-tip-color!10!white, colframe=quarto-callout-tip-color-frame, coltitle=black, left=2mm, leftrule=.75mm, opacityback=0, opacitybacktitle=0.6, rightrule=.15mm, title=\textcolor{quarto-callout-tip-color}{\faLightbulb}\hspace{0.5em}{Solution}, titlerule=0mm, toprule=.15mm, toptitle=1mm]

The density of one observation is
\[f(x\mid\theta)=\frac{1}{\theta}\mathbbm{1}(0\le x\le \theta).\] The
joint density is \[f(x_1,\ldots,x_n\mid\theta)
=\theta^{-n}\prod_{i=1}^n \mathbbm{1}(0\le x_i\le \theta).\] The
condition \(x_i\le \theta\) for all \(i\) is equivalent to
\[\max\{x_1,\ldots,x_n\}\le \theta.\] Also, the condition \(0\le x_i\)
for all \(i\) does not involve \(\theta\). Therefore
\[f(x_1,\ldots,x_n\mid\theta)
=\underbrace{\theta^{-n}\mathbbm{1}\bigl(\max\{x_1,\ldots,x_n\}\le \theta\bigr)}_{g(T(x),\theta)}
\underbrace{\prod_{i=1}^n \mathbbm{1}(x_i\ge 0)}_{h(x)}.\] By the
factorization theorem, \(T(X)=X_{(n)}\) is sufficient for \(\theta\).

\end{tcolorbox}

\begin{tcolorbox}[enhanced jigsaw, arc=.35mm, bottomrule=.15mm, bottomtitle=1mm, breakable, colback=white, colbacktitle=quarto-callout-warning-color!10!white, colframe=quarto-callout-warning-color-frame, coltitle=black, left=2mm, leftrule=.75mm, opacityback=0, opacitybacktitle=0.6, rightrule=.15mm, title=\textcolor{quarto-callout-warning-color}{\faExclamationTriangle}\hspace{0.5em}{Question}, titlerule=0mm, toprule=.15mm, toptitle=1mm]

\textbf{Question 13} (Is the sample mean sufficient?). For the uniform
model \(X_i\overset{\text{iid}}{\sim}\operatorname{Uniform}(0,\theta)\),
is the sample mean \(\overline X\) sufficient for \(\theta\)?

\end{tcolorbox}

\begin{tcolorbox}[enhanced jigsaw, arc=.35mm, bottomrule=.15mm, bottomtitle=1mm, breakable, colback=white, colbacktitle=quarto-callout-tip-color!10!white, colframe=quarto-callout-tip-color-frame, coltitle=black, left=2mm, leftrule=.75mm, opacityback=0, opacitybacktitle=0.6, rightrule=.15mm, title=\textcolor{quarto-callout-tip-color}{\faLightbulb}\hspace{0.5em}{Solution}, titlerule=0mm, toprule=.15mm, toptitle=1mm]

No.~By the factorization theorem, sufficiency would require the
indicator \[\mathbbm{1}\bigl(\max\{x_1,\ldots,x_n\}\le \theta\bigr)\] to
be expressible as a function of only \(\overline x\) and \(\theta\).
This is impossible because two samples can have the same sample mean but
different maxima.

For example, with \(n=2\), \[(0.5,0.5) \quad \text{and} \quad (0,1)\]
have the same sample mean \(0.5\), but their maxima are \(0.5\) and
\(1\). If \(\theta=0.75\), the first sample is possible under
\(\operatorname{Uniform}(0,\theta)\), but the second is not. Therefore
the sample mean does not contain all information about \(\theta\).

\end{tcolorbox}

\bookmarksetup{startatroot}

\chapter{Example: Gamma Population with Unknown
Shape}\label{example-gamma-population-with-unknown-shape}

This section gives an example where the sufficient statistic is a sum of
logarithms rather than a sum of observations.

\begin{tcolorbox}[enhanced jigsaw, arc=.35mm, bottomrule=.15mm, bottomtitle=1mm, breakable, colback=white, colbacktitle=quarto-callout-note-color!10!white, colframe=quarto-callout-note-color-frame, coltitle=black, left=2mm, leftrule=.75mm, opacityback=0, opacitybacktitle=0.6, rightrule=.15mm, title=\textcolor{quarto-callout-note-color}{\faInfo}\hspace{0.5em}{Example}, titlerule=0mm, toprule=.15mm, toptitle=1mm]

\textbf{Example 14} (Gamma population: shape unknown, rate known).
Suppose the population has a gamma distribution with density
\[f(x\mid\alpha)=\frac{\beta^\alpha}{\Gamma(\alpha)}x^{\alpha-1}e^{-\beta x}, \qquad x>0,\]
where \(\beta\) is known and \(\alpha\) is unknown. Let
\(X_1,\ldots,X_n\) be a random sample. Show that
\[T(X)=\sum_{i=1}^n \log X_i\] is sufficient for \(\alpha\).

\end{tcolorbox}

\begin{tcolorbox}[enhanced jigsaw, arc=.35mm, bottomrule=.15mm, bottomtitle=1mm, breakable, colback=white, colbacktitle=quarto-callout-tip-color!10!white, colframe=quarto-callout-tip-color-frame, coltitle=black, left=2mm, leftrule=.75mm, opacityback=0, opacitybacktitle=0.6, rightrule=.15mm, title=\textcolor{quarto-callout-tip-color}{\faLightbulb}\hspace{0.5em}{Solution}, titlerule=0mm, toprule=.15mm, toptitle=1mm]

The joint density is \[\begin{aligned}
f(x_1,\ldots,x_n\mid\alpha)
&=\prod_{i=1}^n \frac{\beta^\alpha}{\Gamma(\alpha)}x_i^{\alpha-1}e^{-\beta x_i} \\
&=\frac{\beta^{n\alpha}}{\Gamma(\alpha)^n}
\left(\prod_{i=1}^n x_i^{\alpha-1}\right)
\exp\left\{-\beta\sum_{i=1}^n x_i\right\}.
\end{aligned}\] Now \[\prod_{i=1}^n x_i^{\alpha-1}
=\exp\left\{(\alpha-1)\sum_{i=1}^n \log x_i\right\}.\] Hence
\[f(x_1,\ldots,x_n\mid\alpha)
=\underbrace{\frac{\beta^{n\alpha}}{\Gamma(\alpha)^n}
\exp\left\{(\alpha-1)T(x)\right\}}_{g(T(x),\alpha)}
\underbrace{\exp\left\{-\beta\sum_{i=1}^n x_i\right\}}_{h(x)}.\] Because
\(h(x)\) does not depend on \(\alpha\) and \(g\) depends on the data
only through \(T(x)=\sum_i \log x_i\), the statistic
\(T(X)=\sum_i \log X_i\) is sufficient for \(\alpha\).

\end{tcolorbox}

\begin{tcolorbox}[enhanced jigsaw, arc=.35mm, bottomrule=.15mm, bottomtitle=1mm, breakable, colback=white, colbacktitle=quarto-callout-important-color!10!white, colframe=quarto-callout-important-color-frame, coltitle=black, left=2mm, leftrule=.75mm, opacityback=0, opacitybacktitle=0.6, rightrule=.15mm, title=\textcolor{quarto-callout-important-color}{\faExclamation}\hspace{0.5em}{Remark}, titlerule=0mm, toprule=.15mm, toptitle=1mm]

\emph{Remark 15}. The statistic \[\exp(T)=\prod_{i=1}^n X_i\] is also
sufficient, because it is a one-to-one transformation of \(T\) on the
positive sample space. In this model, the sample mean \(\overline X\) is
not sufficient for the unknown shape parameter \(\alpha\).

\end{tcolorbox}

\bookmarksetup{startatroot}

\chapter{Minimal and Complete Sufficient
Statistics}\label{minimal-and-complete-sufficient-statistics}

This section introduces two refinements of sufficiency that are
important in statistical theory.

\begin{tcolorbox}[enhanced jigsaw, arc=.35mm, bottomrule=.15mm, bottomtitle=1mm, breakable, colback=white, colbacktitle=quarto-callout-note-color!10!white, colframe=quarto-callout-note-color-frame, coltitle=black, left=2mm, leftrule=.75mm, opacityback=0, opacitybacktitle=0.6, rightrule=.15mm, title=\textcolor{quarto-callout-note-color}{\faInfo}\hspace{0.5em}{Definition}, titlerule=0mm, toprule=.15mm, toptitle=1mm]

\textbf{Definition 16} (Minimal sufficient statistic). A sufficient
statistic \(T\) is called \textbf{minimal sufficient} if it is the
smallest sufficient statistic in the sense that every other sufficient
statistic is a function of it.

\end{tcolorbox}

Minimal sufficiency means that the statistic gives the most compressed
version of the data that still preserves all information about the
parameter.

\begin{tcolorbox}[enhanced jigsaw, arc=.35mm, bottomrule=.15mm, bottomtitle=1mm, breakable, colback=white, colbacktitle=quarto-callout-note-color!10!white, colframe=quarto-callout-note-color-frame, coltitle=black, left=2mm, leftrule=.75mm, opacityback=0, opacitybacktitle=0.6, rightrule=.15mm, title=\textcolor{quarto-callout-note-color}{\faInfo}\hspace{0.5em}{Definition}, titlerule=0mm, toprule=.15mm, toptitle=1mm]

\textbf{Definition 17} (Complete sufficient statistic). A statistic
\(T\) is \textbf{complete and sufficient} for \(\theta\) if it is
sufficient and has the following property:
\[\mathbb{E}_\theta[g(T)]=0 \text{ for all } \theta
\quad \Longrightarrow \quad
\mathbb{P}_\theta(g(T)=0)=1 \text{ for all } \theta.\]

\end{tcolorbox}

Completeness is a technical condition that is useful for proving
uniqueness and optimality results for unbiased estimators.

\begin{tcolorbox}[enhanced jigsaw, arc=.35mm, bottomrule=.15mm, bottomtitle=1mm, breakable, colback=white, colbacktitle=quarto-callout-tip-color!10!white, colframe=quarto-callout-tip-color-frame, coltitle=black, left=2mm, leftrule=.75mm, opacityback=0, opacitybacktitle=0.6, rightrule=.15mm, title=\textcolor{quarto-callout-tip-color}{\faLightbulb}\hspace{0.5em}{Key idea}, titlerule=0mm, toprule=.15mm, toptitle=1mm]

Why sufficiency matters Sufficient statistics are important because they
support:

\begin{itemize}
\item
  \textbf{data reduction}: compress the sample without losing
  information about \(\theta\);
\item
  \textbf{maximum likelihood estimation}: MLEs often depend only on
  sufficient statistics;
\item
  \textbf{Bayesian inference}: posterior distributions often depend on
  the data only through sufficient statistics;
\item
  \textbf{efficient inference}: sufficient statistics help identify
  estimators with good theoretical properties.
\end{itemize}

\end{tcolorbox}

\bookmarksetup{startatroot}

\chapter{The Likelihood Function}\label{the-likelihood-function}

This section introduces the likelihood function, which is the central
object connecting observed data to inference about parameters.

\begin{tcolorbox}[enhanced jigsaw, arc=.35mm, bottomrule=.15mm, bottomtitle=1mm, breakable, colback=white, colbacktitle=quarto-callout-note-color!10!white, colframe=quarto-callout-note-color-frame, coltitle=black, left=2mm, leftrule=.75mm, opacityback=0, opacitybacktitle=0.6, rightrule=.15mm, title=\textcolor{quarto-callout-note-color}{\faInfo}\hspace{0.5em}{Definition}, titlerule=0mm, toprule=.15mm, toptitle=1mm]

\textbf{Definition 18} (Likelihood function). Let \(f(x\mid\theta)\) be
the joint density or joint pmf of the sample \(X=(X_1,\ldots,X_n)\).
After observing \[X=x=(x_1,\ldots,x_n),\] the \textbf{likelihood
function} is the function of \(\theta\) defined by
\[L(\theta\mid x)=f(x\mid\theta).\]

\end{tcolorbox}

The joint density and the likelihood function have the same algebraic
expression, but they are viewed differently:

\begin{itemize}
\item
  as a density or pmf, \(x\) varies and \(\theta\) is fixed;
\item
  as a likelihood, \(x\) is observed and fixed, while \(\theta\) varies.
\end{itemize}

\begin{tcolorbox}[enhanced jigsaw, arc=.35mm, bottomrule=.15mm, bottomtitle=1mm, breakable, colback=white, colbacktitle=quarto-callout-note-color!10!white, colframe=quarto-callout-note-color-frame, coltitle=black, left=2mm, leftrule=.75mm, opacityback=0, opacitybacktitle=0.6, rightrule=.15mm, title=\textcolor{quarto-callout-note-color}{\faInfo}\hspace{0.5em}{Example}, titlerule=0mm, toprule=.15mm, toptitle=1mm]

\textbf{Example 19} (Bernoulli likelihood). Let
\(X_1,\ldots,X_n\overset{\text{iid}}{\sim}\operatorname{Bernoulli}(p)\),
and suppose the observed data have \(t=\sum_i x_i\) successes. The
likelihood is \[L(p\mid x)=p^t(1-p)^{n-t}, \qquad 0<p<1.\]

\end{tcolorbox}

\begin{tcolorbox}[enhanced jigsaw, arc=.35mm, bottomrule=.15mm, bottomtitle=1mm, breakable, colback=white, colbacktitle=quarto-callout-tip-color!10!white, colframe=quarto-callout-tip-color-frame, coltitle=black, left=2mm, leftrule=.75mm, opacityback=0, opacitybacktitle=0.6, rightrule=.15mm, title=\textcolor{quarto-callout-tip-color}{\faLightbulb}\hspace{0.5em}{Solution}, titlerule=0mm, toprule=.15mm, toptitle=1mm]

The joint pmf is \[\prod_{i=1}^n p^{x_i}(1-p)^{1-x_i}
=p^{\sum_i x_i}(1-p)^{n-\sum_i x_i}.\] After the data are observed,
\(t=\sum_i x_i\) is fixed, so the likelihood as a function of \(p\) is
\[L(p\mid x)=p^t(1-p)^{n-t}.\]

\end{tcolorbox}

\begin{tcolorbox}[enhanced jigsaw, arc=.35mm, bottomrule=.15mm, bottomtitle=1mm, breakable, colback=white, colbacktitle=quarto-callout-important-color!10!white, colframe=quarto-callout-important-color-frame, coltitle=black, left=2mm, leftrule=.75mm, opacityback=0, opacitybacktitle=0.6, rightrule=.15mm, title=\textcolor{quarto-callout-important-color}{\faExclamation}\hspace{0.5em}{Remark}, titlerule=0mm, toprule=.15mm, toptitle=1mm]

\emph{Remark 20} (Connection with sufficiency). The factorization
theorem says that if \(T\) is sufficient, then the likelihood can be
written so that all parameter-dependent information from the data passes
through \(T\).

\end{tcolorbox}

\bookmarksetup{startatroot}

\chapter{The Likelihood Principle}\label{the-likelihood-principle}

This section states the likelihood principle, a foundational idea
especially important in Bayesian statistics.

\begin{tcolorbox}[enhanced jigsaw, arc=.35mm, bottomrule=.15mm, bottomtitle=1mm, breakable, colback=white, colbacktitle=quarto-callout-note-color!10!white, colframe=quarto-callout-note-color-frame, coltitle=black, left=2mm, leftrule=.75mm, opacityback=0, opacitybacktitle=0.6, rightrule=.15mm, title=\textcolor{quarto-callout-note-color}{\faInfo}\hspace{0.5em}{Definition}, titlerule=0mm, toprule=.15mm, toptitle=1mm]

\textbf{Definition 21} (Likelihood principle). If two experiments
produce likelihood functions that are proportional as functions of
\(\theta\), then they provide the same evidence about \(\theta\) and
should lead to the same inference about \(\theta\).

More precisely, if two observed data sets \(x\) and \(y\) have
likelihoods satisfying \[L_1(\theta\mid x)=C(x,y)L_2(\theta\mid y)
\quad \text{for all } \theta,\] where \(C(x,y)\) does not depend on
\(\theta\), then the two data sets have the same likelihood information
about \(\theta\).

\end{tcolorbox}

\begin{tcolorbox}[enhanced jigsaw, arc=.35mm, bottomrule=.15mm, bottomtitle=1mm, breakable, colback=white, colbacktitle=quarto-callout-tip-color!10!white, colframe=quarto-callout-tip-color-frame, coltitle=black, left=2mm, leftrule=.75mm, opacityback=0, opacitybacktitle=0.6, rightrule=.15mm, title=\textcolor{quarto-callout-tip-color}{\faLightbulb}\hspace{0.5em}{Key idea}, titlerule=0mm, toprule=.15mm, toptitle=1mm]

Interpretation The likelihood principle says that once the data are
observed, the evidential content about \(\theta\) is contained in the
likelihood function. How the data were collected matters only through
its effect on the likelihood.

\end{tcolorbox}

This principle is central in Bayesian inference, where the posterior is
proportional to likelihood times prior:
\[\pi(\theta\mid x)\propto L(\theta\mid x)\pi(\theta).\] If two
likelihoods are proportional and the same prior is used, then the
posterior distributions are identical.

\bookmarksetup{startatroot}

\chapter{Coin-Tossing Example: Binomial versus Negative
Binomial}\label{coin-tossing-example-binomial-versus-negative-binomial}

This section illustrates the likelihood principle using two coin-tossing
experiments.

\begin{tcolorbox}[enhanced jigsaw, arc=.35mm, bottomrule=.15mm, bottomtitle=1mm, breakable, colback=white, colbacktitle=quarto-callout-note-color!10!white, colframe=quarto-callout-note-color-frame, coltitle=black, left=2mm, leftrule=.75mm, opacityback=0, opacitybacktitle=0.6, rightrule=.15mm, title=\textcolor{quarto-callout-note-color}{\faInfo}\hspace{0.5em}{Example}, titlerule=0mm, toprule=.15mm, toptitle=1mm]

\textbf{Example 22} (Coin tossing and the likelihood principle).
Consider two different experiments involving a coin with unknown
probability \(p\) of heads.

\textbf{Scenario 1: Binomial experiment.} Toss the coin \(n=10\) times
and observe \(7\) heads. Then
\[X=7, \qquad X\sim \operatorname{Binomial}(10,p),\] and the likelihood
is \[L_1(p)=\binom{10}{7}p^7(1-p)^3.\]

\textbf{Scenario 2: Negative binomial experiment.} Toss the coin until
\(7\) heads appear. Suppose it takes \(10\) tosses. Then the first \(9\)
tosses contain \(6\) heads and \(3\) tails, and the tenth toss is heads.
The likelihood is \[L_2(p)=\binom{9}{6}p^7(1-p)^3.\]

Explain why the likelihood principle says these two experiments provide
the same evidence about \(p\).

\end{tcolorbox}

\begin{tcolorbox}[enhanced jigsaw, arc=.35mm, bottomrule=.15mm, bottomtitle=1mm, breakable, colback=white, colbacktitle=quarto-callout-tip-color!10!white, colframe=quarto-callout-tip-color-frame, coltitle=black, left=2mm, leftrule=.75mm, opacityback=0, opacitybacktitle=0.6, rightrule=.15mm, title=\textcolor{quarto-callout-tip-color}{\faLightbulb}\hspace{0.5em}{Solution}, titlerule=0mm, toprule=.15mm, toptitle=1mm]

The two likelihoods are \[L_1(p)=\binom{10}{7}p^7(1-p)^3\] and
\[L_2(p)=\binom{9}{6}p^7(1-p)^3.\] They differ only by a multiplicative
constant that does not depend on \(p\):
\[L_1(p)=\frac{\binom{10}{7}}{\binom{9}{6}}L_2(p).\] Thus the likelihood
functions are proportional as functions of \(p\). By the likelihood
principle, the two experiments provide the same evidence about \(p\),
even though the stopping rules were different.

\end{tcolorbox}

\begin{tcolorbox}[enhanced jigsaw, arc=.35mm, bottomrule=.15mm, bottomtitle=1mm, breakable, colback=white, colbacktitle=quarto-callout-important-color!10!white, colframe=quarto-callout-important-color-frame, coltitle=black, left=2mm, leftrule=.75mm, opacityback=0, opacitybacktitle=0.6, rightrule=.15mm, title=\textcolor{quarto-callout-important-color}{\faExclamation}\hspace{0.5em}{Remark}, titlerule=0mm, toprule=.15mm, toptitle=1mm]

\emph{Remark 23}. Frequentist procedures sometimes depend on the
sampling plan or stopping rule, while likelihood-based and Bayesian
procedures focus on the likelihood after the data are observed. This is
one reason the likelihood principle is philosophically important.

\end{tcolorbox}

\bookmarksetup{startatroot}

\chapter{Practice Problems}\label{practice-problems-2}

This section gives additional practice with the factorization theorem
and likelihood principle.

\begin{tcolorbox}[enhanced jigsaw, arc=.35mm, bottomrule=.15mm, bottomtitle=1mm, breakable, colback=white, colbacktitle=quarto-callout-warning-color!10!white, colframe=quarto-callout-warning-color-frame, coltitle=black, left=2mm, leftrule=.75mm, opacityback=0, opacitybacktitle=0.6, rightrule=.15mm, title=\textcolor{quarto-callout-warning-color}{\faExclamationTriangle}\hspace{0.5em}{Practice Problem}, titlerule=0mm, toprule=.15mm, toptitle=1mm]

\textbf{Practice Problem 24} (Poisson sufficient statistic). Let
\(X_1,\ldots,X_n\overset{\text{iid}}{\sim}\operatorname{Poisson}(\lambda)\).
Find a sufficient statistic for \(\lambda\).

\end{tcolorbox}

\begin{tcolorbox}[enhanced jigsaw, arc=.35mm, bottomrule=.15mm, bottomtitle=1mm, breakable, colback=white, colbacktitle=quarto-callout-tip-color!10!white, colframe=quarto-callout-tip-color-frame, coltitle=black, left=2mm, leftrule=.75mm, opacityback=0, opacitybacktitle=0.6, rightrule=.15mm, title=\textcolor{quarto-callout-tip-color}{\faLightbulb}\hspace{0.5em}{Solution}, titlerule=0mm, toprule=.15mm, toptitle=1mm]

The joint pmf is \[f(x_1,\ldots,x_n\mid\lambda)
=\prod_{i=1}^n \frac{e^{-\lambda}\lambda^{x_i}}{x_i!}
=e^{-n\lambda}\lambda^{\sum_i x_i}\prod_{i=1}^n \frac{1}{x_i!}.\] Let
\[T(X)=\sum_{i=1}^n X_i.\] Then
\[f(x\mid\lambda)=\underbrace{e^{-n\lambda}\lambda^{T(x)}}_{g(T(x),\lambda)}
\underbrace{\prod_{i=1}^n \frac{1}{x_i!}}_{h(x)}.\] By the factorization
theorem, \(T(X)=\sum_i X_i\) is sufficient for \(\lambda\).

\end{tcolorbox}

\begin{tcolorbox}[enhanced jigsaw, arc=.35mm, bottomrule=.15mm, bottomtitle=1mm, breakable, colback=white, colbacktitle=quarto-callout-warning-color!10!white, colframe=quarto-callout-warning-color-frame, coltitle=black, left=2mm, leftrule=.75mm, opacityback=0, opacitybacktitle=0.6, rightrule=.15mm, title=\textcolor{quarto-callout-warning-color}{\faExclamationTriangle}\hspace{0.5em}{Practice Problem}, titlerule=0mm, toprule=.15mm, toptitle=1mm]

\textbf{Practice Problem 25} (Exponential sufficient statistic). Let
\(X_1,\ldots,X_n \overset{\text{iid}}{\sim}Exp(\lambda)\) with density
\[f(x\mid\lambda)=\lambda e^{-\lambda x},\qquad x>0.\] Find a sufficient
statistic for \(\lambda\).

\end{tcolorbox}

\begin{tcolorbox}[enhanced jigsaw, arc=.35mm, bottomrule=.15mm, bottomtitle=1mm, breakable, colback=white, colbacktitle=quarto-callout-tip-color!10!white, colframe=quarto-callout-tip-color-frame, coltitle=black, left=2mm, leftrule=.75mm, opacityback=0, opacitybacktitle=0.6, rightrule=.15mm, title=\textcolor{quarto-callout-tip-color}{\faLightbulb}\hspace{0.5em}{Solution}, titlerule=0mm, toprule=.15mm, toptitle=1mm]

The joint density is \[f(x_1,\ldots,x_n\mid\lambda)
=\lambda^n\exp\left\{-\lambda\sum_{i=1}^n x_i\right\}
\prod_{i=1}^n \mathbbm{1}(x_i>0).\] Let \[T(X)=\sum_{i=1}^n X_i.\] Then
\[f(x\mid\lambda)
=\underbrace{\lambda^n e^{-\lambda T(x)}}_{g(T(x),\lambda)}
\underbrace{\prod_{i=1}^n \mathbbm{1}(x_i>0)}_{h(x)}.\] Therefore
\(T(X)=\sum_i X_i\) is sufficient for \(\lambda\).

\end{tcolorbox}

\begin{tcolorbox}[enhanced jigsaw, arc=.35mm, bottomrule=.15mm, bottomtitle=1mm, breakable, colback=white, colbacktitle=quarto-callout-warning-color!10!white, colframe=quarto-callout-warning-color-frame, coltitle=black, left=2mm, leftrule=.75mm, opacityback=0, opacitybacktitle=0.6, rightrule=.15mm, title=\textcolor{quarto-callout-warning-color}{\faExclamationTriangle}\hspace{0.5em}{Practice Problem}, titlerule=0mm, toprule=.15mm, toptitle=1mm]

\textbf{Practice Problem 26} (Normal distribution with both parameters
unknown). Let
\(X_1,\ldots,X_n\overset{\text{iid}}{\sim}\operatorname{Normal}(\mu,\sigma^2)\),
where both \(\mu\) and \(\sigma^2\) are unknown. Find a sufficient
statistic for \((\mu,\sigma^2)\).

\end{tcolorbox}

\begin{tcolorbox}[enhanced jigsaw, arc=.35mm, bottomrule=.15mm, bottomtitle=1mm, breakable, colback=white, colbacktitle=quarto-callout-tip-color!10!white, colframe=quarto-callout-tip-color-frame, coltitle=black, left=2mm, leftrule=.75mm, opacityback=0, opacitybacktitle=0.6, rightrule=.15mm, title=\textcolor{quarto-callout-tip-color}{\faLightbulb}\hspace{0.5em}{Solution}, titlerule=0mm, toprule=.15mm, toptitle=1mm]

The joint density is \[f(x\mid\mu,\sigma^2)
=(2\pi\sigma^2)^{-n/2}
\exp\left\{-\frac{1}{2\sigma^2}\sum_{i=1}^n (x_i-\mu)^2\right\}.\]
Expand: \[\sum_{i=1}^n (x_i-\mu)^2
=\sum_{i=1}^n x_i^2-2\mu\sum_{i=1}^n x_i+n\mu^2.\] Thus the joint
density depends on the data through \[\sum_{i=1}^n x_i
\quad \text{and} \quad
\sum_{i=1}^n x_i^2.\] Therefore
\[T(X)=\left(\sum_{i=1}^n X_i,\sum_{i=1}^n X_i^2\right)\] is sufficient
for \((\mu,\sigma^2)\) by the factorization theorem.

\end{tcolorbox}

\begin{tcolorbox}[enhanced jigsaw, arc=.35mm, bottomrule=.15mm, bottomtitle=1mm, breakable, colback=white, colbacktitle=quarto-callout-warning-color!10!white, colframe=quarto-callout-warning-color-frame, coltitle=black, left=2mm, leftrule=.75mm, opacityback=0, opacitybacktitle=0.6, rightrule=.15mm, title=\textcolor{quarto-callout-warning-color}{\faExclamationTriangle}\hspace{0.5em}{Practice Problem}, titlerule=0mm, toprule=.15mm, toptitle=1mm]

\textbf{Practice Problem 27} (Uniform lower and upper endpoints). Let
\(X_1,\ldots,X_n\overset{\text{iid}}{\sim}\operatorname{Uniform}(a,b)\),
where both \(a\) and \(b\) are unknown. Find a sufficient statistic for
\((a,b)\).

\end{tcolorbox}

\begin{tcolorbox}[enhanced jigsaw, arc=.35mm, bottomrule=.15mm, bottomtitle=1mm, breakable, colback=white, colbacktitle=quarto-callout-tip-color!10!white, colframe=quarto-callout-tip-color-frame, coltitle=black, left=2mm, leftrule=.75mm, opacityback=0, opacitybacktitle=0.6, rightrule=.15mm, title=\textcolor{quarto-callout-tip-color}{\faLightbulb}\hspace{0.5em}{Solution}, titlerule=0mm, toprule=.15mm, toptitle=1mm]

The density of one observation is
\[f(x\mid a,b)=\frac{1}{b-a}\mathbbm{1}(a\le x\le b).\] The joint
density is \[f(x_1,\ldots,x_n\mid a,b)
=(b-a)^{-n}\prod_{i=1}^n \mathbbm{1}(a\le x_i\le b).\] The condition
\(a\le x_i\le b\) for all \(i\) is equivalent to \[a\le \min_i x_i
\quad \text{and} \quad
\max_i x_i\le b.\] Hence \[f(x\mid a,b)
=(b-a)^{-n}\mathbbm{1}(a\le x_{(1)})\mathbbm{1}(x_{(n)}\le b).\] This
depends on the data only through \[T(X)=(X_{(1)},X_{(n)}).\] Therefore
\((X_{(1)},X_{(n)})\) is sufficient for \((a,b)\).

\end{tcolorbox}

\begin{tcolorbox}[enhanced jigsaw, arc=.35mm, bottomrule=.15mm, bottomtitle=1mm, breakable, colback=white, colbacktitle=quarto-callout-warning-color!10!white, colframe=quarto-callout-warning-color-frame, coltitle=black, left=2mm, leftrule=.75mm, opacityback=0, opacitybacktitle=0.6, rightrule=.15mm, title=\textcolor{quarto-callout-warning-color}{\faExclamationTriangle}\hspace{0.5em}{Practice Problem}, titlerule=0mm, toprule=.15mm, toptitle=1mm]

\textbf{Practice Problem 28} (Likelihood principle). Suppose two
experiments give likelihoods \[L_1(\theta)=12\theta^4(1-\theta)^6,
\qquad
L_2(\theta)=3\theta^4(1-\theta)^6,\] for \(0<\theta<1\). Do they provide
the same likelihood information about \(\theta\)?

\end{tcolorbox}

\begin{tcolorbox}[enhanced jigsaw, arc=.35mm, bottomrule=.15mm, bottomtitle=1mm, breakable, colback=white, colbacktitle=quarto-callout-tip-color!10!white, colframe=quarto-callout-tip-color-frame, coltitle=black, left=2mm, leftrule=.75mm, opacityback=0, opacitybacktitle=0.6, rightrule=.15mm, title=\textcolor{quarto-callout-tip-color}{\faLightbulb}\hspace{0.5em}{Solution}, titlerule=0mm, toprule=.15mm, toptitle=1mm]

Yes. We have \[L_1(\theta)=4L_2(\theta),\] and the constant \(4\) does
not depend on \(\theta\). Therefore the likelihoods are proportional as
functions of \(\theta\). By the likelihood principle, they provide the
same likelihood information about \(\theta\).

\end{tcolorbox}

\bookmarksetup{startatroot}

\chapter{Summary}\label{summary-5}

This section is about reducing data without losing information about
parameters.

\begin{tcolorbox}[enhanced jigsaw, arc=.35mm, bottomrule=.15mm, bottomtitle=1mm, breakable, colback=white, colbacktitle=quarto-callout-tip-color!10!white, colframe=quarto-callout-tip-color-frame, coltitle=black, left=2mm, leftrule=.75mm, opacityback=0, opacitybacktitle=0.6, rightrule=.15mm, title=\textcolor{quarto-callout-tip-color}{\faLightbulb}\hspace{0.5em}{Key idea}, titlerule=0mm, toprule=.15mm, toptitle=1mm]

\begin{enumerate}
\def\labelenumi{\arabic{enumi}.}
\item
  A statistic is a function of the random sample.
\item
  A statistic is sufficient for \(\theta\) if the conditional
  distribution of the full sample given the statistic does not depend on
  \(\theta\).
\item
  The factorization theorem is the main practical tool for proving
  sufficiency.
\item
  For Bernoulli, Poisson, exponential, and normal-with-known-variance
  models, sums often appear as sufficient statistics.
\item
  For a uniform model with unknown endpoint, the maximum is sufficient;
  the sample mean is not.
\item
  The likelihood function is the joint density or pmf viewed as a
  function of the parameter after observing data.
\item
  The likelihood principle says proportional likelihoods provide the
  same evidence about the parameter.
\end{enumerate}

\end{tcolorbox}

\bookmarksetup{startatroot}

\chapter{Chapter 12: Point Estimation I --- Finding
Estimators}\label{chapter-12-point-estimation-i-finding-estimators}

This chapter begins the statistical inference part of the course. The
main goal is to construct point estimators for unknown parameters using
three major principles: method of moments, maximum likelihood
estimation, and Bayesian estimation.

\begin{tcolorbox}[enhanced jigsaw, arc=.35mm, bottomrule=.15mm, bottomtitle=1mm, breakable, colback=white, colbacktitle=quarto-callout-note-color!10!white, colframe=quarto-callout-note-color-frame, coltitle=black, left=2mm, leftrule=.75mm, opacityback=0, opacitybacktitle=0.6, rightrule=.15mm, title=\textcolor{quarto-callout-note-color}{\faInfo}\hspace{0.5em}{Topics}, titlerule=0mm, toprule=.15mm, toptitle=1mm]

Point estimators and estimates; method of moments; maximum likelihood
estimation; log-likelihood optimization; Hessian and second derivative
test; normal, log-normal, Bernoulli, binomial, and uniform examples;
invariance property of MLE; Bayesian estimators; beta-binomial
conjugacy; posterior mean, variance, mode, credible probabilities;
conjugate priors; MAP versus MLE; normal-normal MAP.

\end{tcolorbox}

\bookmarksetup{startatroot}

\chapter{Overview}\label{overview-9}

This section introduces three main methods for finding point estimators:
the method of moments, maximum likelihood estimation, and Bayesian
estimation.

In point estimation, we observe a random sample \(X_1,\ldots,X_n\) from
a population model with density or mass function \(f(x\mid \theta)\).
The unknown parameter \(\theta\), or a function \(g(\theta)\),
determines important features of the population. The goal is to
construct a statistic from the data that gives a reasonable numerical
estimate of the unknown quantity.

\begin{tcolorbox}[enhanced jigsaw, arc=.35mm, bottomrule=.15mm, bottomtitle=1mm, breakable, colback=white, colbacktitle=quarto-callout-tip-color!10!white, colframe=quarto-callout-tip-color-frame, coltitle=black, left=2mm, leftrule=.75mm, opacityback=0, opacitybacktitle=0.6, rightrule=.15mm, title=\textcolor{quarto-callout-tip-color}{\faLightbulb}\hspace{0.5em}{Key idea}, titlerule=0mm, toprule=.15mm, toptitle=1mm]

A point estimator is a statistic used to estimate an unknown population
parameter. This section studies three major construction principles:
\[\text{match moments}, \qquad \text{maximize likelihood}, \qquad \text{update prior information by Bayes' rule}.\]

\end{tcolorbox}

The method of moments is often simple and intuitive. Maximum likelihood
is usually more efficient and is the most widely used method in
statistical modeling. Bayesian estimation treats the parameter as
uncertain and combines prior information with the likelihood from the
observed data.

\bookmarksetup{startatroot}

\chapter{Point Estimators and
Estimates}\label{point-estimators-and-estimates}

This section sets up the language of estimators and estimates.

\begin{tcolorbox}[enhanced jigsaw, arc=.35mm, bottomrule=.15mm, bottomtitle=1mm, breakable, colback=white, colbacktitle=quarto-callout-note-color!10!white, colframe=quarto-callout-note-color-frame, coltitle=black, left=2mm, leftrule=.75mm, opacityback=0, opacitybacktitle=0.6, rightrule=.15mm, title=\textcolor{quarto-callout-note-color}{\faInfo}\hspace{0.5em}{Definition}, titlerule=0mm, toprule=.15mm, toptitle=1mm]

\textbf{Definition 1} (Point estimator). Let \(X_1,\ldots,X_n\) be a
random sample from a population distribution with pdf or pmf
\(f(x\mid \theta)\). A \textbf{point estimator} of \(\theta\) or
\(g(\theta)\) is any statistic \[W=W(X_1,\ldots,X_n)\] used to estimate
the unknown parameter or function of the parameter.

\end{tcolorbox}

\begin{tcolorbox}[enhanced jigsaw, arc=.35mm, bottomrule=.15mm, bottomtitle=1mm, breakable, colback=white, colbacktitle=quarto-callout-note-color!10!white, colframe=quarto-callout-note-color-frame, coltitle=black, left=2mm, leftrule=.75mm, opacityback=0, opacitybacktitle=0.6, rightrule=.15mm, title=\textcolor{quarto-callout-note-color}{\faInfo}\hspace{0.5em}{Definition}, titlerule=0mm, toprule=.15mm, toptitle=1mm]

\textbf{Definition 2} (Estimate). After observing the sample values
\[(X_1,\ldots,X_n)=(x_1,\ldots,x_n),\] the number \[W(x_1,\ldots,x_n)\]
is called an \textbf{estimate}.

\end{tcolorbox}

The distinction is important: an estimator is a random variable before
observing data; an estimate is the realized numerical value after
observing data.

\begin{tcolorbox}[enhanced jigsaw, arc=.35mm, bottomrule=.15mm, bottomtitle=1mm, breakable, colback=white, colbacktitle=quarto-callout-note-color!10!white, colframe=quarto-callout-note-color-frame, coltitle=black, left=2mm, leftrule=.75mm, opacityback=0, opacitybacktitle=0.6, rightrule=.15mm, title=\textcolor{quarto-callout-note-color}{\faInfo}\hspace{0.5em}{Example}, titlerule=0mm, toprule=.15mm, toptitle=1mm]

\textbf{Example 3} (Common point estimators). For a random sample
\(X_1,\ldots,X_n\), common point estimators include:

\begin{enumerate}
\def\labelenumi{\arabic{enumi}.}
\item
  the sample mean \[\overline X=\frac{1}{n}\sum_{i=1}^n X_i,\] which
  estimates the population mean \(\mu\);
\item
  the sample variance
  \[S^2=\frac{1}{n-1}\sum_{i=1}^n (X_i-\overline X)^2,\] which estimates
  the population variance \(\sigma^2\);
\item
  the sample proportion \[\widehat p=\frac{1}{n}\sum_{i=1}^n X_i\] for
  Bernoulli data, which estimates the success probability \(p\).
\end{enumerate}

\end{tcolorbox}

\begin{tcolorbox}[enhanced jigsaw, arc=.35mm, bottomrule=.15mm, bottomtitle=1mm, breakable, colback=white, colbacktitle=quarto-callout-tip-color!10!white, colframe=quarto-callout-tip-color-frame, coltitle=black, left=2mm, leftrule=.75mm, opacityback=0, opacitybacktitle=0.6, rightrule=.15mm, title=\textcolor{quarto-callout-tip-color}{\faLightbulb}\hspace{0.5em}{Solution}, titlerule=0mm, toprule=.15mm, toptitle=1mm]

Each quantity is a function of the random sample and does not depend on
the unknown parameter. Therefore each is a statistic. When used to
estimate a population quantity, it is called a point estimator.

\end{tcolorbox}

\bookmarksetup{startatroot}

\chapter{Method of Moments}\label{method-of-moments}

This section introduces the method of moments, also called moment
matching.

Suppose \(X_1,\ldots,X_n\) is a sample from a population distribution
with pdf or pmf
\[f(x\mid \boldsymbol\theta), \qquad \boldsymbol\theta=(\theta_1,\ldots,\theta_k).\]
The method of moments estimates the unknown parameters by equating
sample moments with population moments.

\begin{tcolorbox}[enhanced jigsaw, arc=.35mm, bottomrule=.15mm, bottomtitle=1mm, breakable, colback=white, colbacktitle=quarto-callout-note-color!10!white, colframe=quarto-callout-note-color-frame, coltitle=black, left=2mm, leftrule=.75mm, opacityback=0, opacitybacktitle=0.6, rightrule=.15mm, title=\textcolor{quarto-callout-note-color}{\faInfo}\hspace{0.5em}{Definition}, titlerule=0mm, toprule=.15mm, toptitle=1mm]

\textbf{Definition 4} (Sample and population moments). The first \(k\)
sample moments are \[m_1=\frac{1}{n}\sum_{i=1}^n X_i, \qquad
m_2=\frac{1}{n}\sum_{i=1}^n X_i^2, \qquad \ldots, \qquad
m_k=\frac{1}{n}\sum_{i=1}^n X_i^k.\] The corresponding population
moments are \[\mu_1'=\mathbb{E}[X], \qquad
\mu_2'=\mathbb{E}[X^2], \qquad \ldots, \qquad
\mu_k'=\mathbb{E}[X^k].\] Usually the population moments depend on the
unknown parameter vector \(\boldsymbol\theta\).

\end{tcolorbox}

\begin{tcolorbox}[enhanced jigsaw, arc=.35mm, bottomrule=.15mm, bottomtitle=1mm, breakable, colback=white, colbacktitle=quarto-callout-note-color!10!white, colframe=quarto-callout-note-color-frame, coltitle=black, left=2mm, leftrule=.75mm, opacityback=0, opacitybacktitle=0.6, rightrule=.15mm, title=\textcolor{quarto-callout-note-color}{\faInfo}\hspace{0.5em}{Definition}, titlerule=0mm, toprule=.15mm, toptitle=1mm]

\textbf{Definition 5} (Method of moments estimator). The \textbf{method
of moments estimators} are obtained by solving the system
\[m_1=\mu_1'(\boldsymbol\theta), \qquad
m_2=\mu_2'(\boldsymbol\theta), \qquad \ldots, \qquad
m_k=\mu_k'(\boldsymbol\theta)\] for \(\theta_1,\ldots,\theta_k\).

\end{tcolorbox}

\begin{tcolorbox}[enhanced jigsaw, arc=.35mm, bottomrule=.15mm, bottomtitle=1mm, breakable, colback=white, colbacktitle=quarto-callout-tip-color!10!white, colframe=quarto-callout-tip-color-frame, coltitle=black, left=2mm, leftrule=.75mm, opacityback=0, opacitybacktitle=0.6, rightrule=.15mm, title=\textcolor{quarto-callout-tip-color}{\faLightbulb}\hspace{0.5em}{Key idea}, titlerule=0mm, toprule=.15mm, toptitle=1mm]

How to use method of moments

\begin{enumerate}
\def\labelenumi{\arabic{enumi}.}
\item
  Count the number of unknown parameters.
\item
  Write the same number of population moment equations.
\item
  Replace population moments by sample moments.
\item
  Solve the resulting equations for the parameters.
\end{enumerate}

\end{tcolorbox}

\section{Method of Moments for the Normal
Distribution}\label{method-of-moments-for-the-normal-distribution}

This subsection shows how method of moments estimates the parameters of
a normal distribution.

\begin{tcolorbox}[enhanced jigsaw, arc=.35mm, bottomrule=.15mm, bottomtitle=1mm, breakable, colback=white, colbacktitle=quarto-callout-note-color!10!white, colframe=quarto-callout-note-color-frame, coltitle=black, left=2mm, leftrule=.75mm, opacityback=0, opacitybacktitle=0.6, rightrule=.15mm, title=\textcolor{quarto-callout-note-color}{\faInfo}\hspace{0.5em}{Example}, titlerule=0mm, toprule=.15mm, toptitle=1mm]

\textbf{Example 6} (Normal distribution). Suppose
\[X_1,\ldots,X_n \overset{\text{iid}}{\sim}\operatorname{Normal}(\mu,\sigma^2).\]
Find the method of moments estimators of \(\mu\) and \(\sigma^2\).

\end{tcolorbox}

\begin{tcolorbox}[enhanced jigsaw, arc=.35mm, bottomrule=.15mm, bottomtitle=1mm, breakable, colback=white, colbacktitle=quarto-callout-tip-color!10!white, colframe=quarto-callout-tip-color-frame, coltitle=black, left=2mm, leftrule=.75mm, opacityback=0, opacitybacktitle=0.6, rightrule=.15mm, title=\textcolor{quarto-callout-tip-color}{\faLightbulb}\hspace{0.5em}{Solution}, titlerule=0mm, toprule=.15mm, toptitle=1mm]

There are two unknown parameters, so we match the first two moments. Let
\[\theta_1=\mu, \qquad \theta_2=\sigma^2.\] The first sample moment is
\[m_1=\frac{1}{n}\sum_{i=1}^n X_i=\overline X.\] The second sample
moment is \[m_2=\frac{1}{n}\sum_{i=1}^n X_i^2.\] For a normal random
variable, \[\mu_1'=\mathbb{E}[X]=\mu,\] and
\[\mu_2'=\mathbb{E}[X^2]=\operatorname{Var}(X)+\{\mathbb{E}[X]\}^2=\sigma^2+\mu^2.\]
The method of moments equations are
\[m_1=\mu, \qquad m_2=\mu^2+\sigma^2.\] Therefore
\[\widetilde\mu_{\mathrm{MOM}}=\overline X\] and
\[\widetilde\sigma^2_{\mathrm{MOM}}=m_2-m_1^2
=\frac{1}{n}\sum_{i=1}^n X_i^2-\overline X^{\,2}
=\frac{1}{n}\sum_{i=1}^n (X_i-\overline X)^2.\]

\end{tcolorbox}

\section{Method of Moments for the Binomial
Distribution}\label{method-of-moments-for-the-binomial-distribution}

This subsection illustrates that method of moments may produce
estimators that are algebraically valid but practically imperfect.

\begin{tcolorbox}[enhanced jigsaw, arc=.35mm, bottomrule=.15mm, bottomtitle=1mm, breakable, colback=white, colbacktitle=quarto-callout-note-color!10!white, colframe=quarto-callout-note-color-frame, coltitle=black, left=2mm, leftrule=.75mm, opacityback=0, opacitybacktitle=0.6, rightrule=.15mm, title=\textcolor{quarto-callout-note-color}{\faInfo}\hspace{0.5em}{Example}, titlerule=0mm, toprule=.15mm, toptitle=1mm]

\textbf{Example 7} (Binomial distribution with two unknown parameters).
Suppose
\[X_1,\ldots,X_n \overset{\text{iid}}{\sim}\operatorname{Binomial}(k,p),\]
where both \(k\) and \(p\) are unknown. Find the method of moments
estimators of \(k\) and \(p\).

\end{tcolorbox}

\begin{tcolorbox}[enhanced jigsaw, arc=.35mm, bottomrule=.15mm, bottomtitle=1mm, breakable, colback=white, colbacktitle=quarto-callout-tip-color!10!white, colframe=quarto-callout-tip-color-frame, coltitle=black, left=2mm, leftrule=.75mm, opacityback=0, opacitybacktitle=0.6, rightrule=.15mm, title=\textcolor{quarto-callout-tip-color}{\faLightbulb}\hspace{0.5em}{Solution}, titlerule=0mm, toprule=.15mm, toptitle=1mm]

For \(X\sim \operatorname{Binomial}(k,p)\), \[\mathbb{E}[X]=kp,\] and
\[\operatorname{Var}(X)=kp(1-p).\] Therefore
\[\mathbb{E}[X^2]=\operatorname{Var}(X)+\{\mathbb{E}[X]\}^2=kp(1-p)+k^2p^2.\]
Let \[m_1=\overline X, \qquad m_2=\frac{1}{n}\sum_{i=1}^n X_i^2.\] The
method of moments equations are \[m_1=kp,\] and \[m_2=kp(1-p)+k^2p^2.\]
Using \(kp=m_1\), we get \[m_2=m_1(1-p)+m_1^2.\] Thus
\[p=1-\frac{m_2-m_1^2}{m_1}.\] The method of moments estimator of \(p\)
is \[\widetilde p_{\mathrm{MOM}}
=1-\frac{m_2-m_1^2}{m_1}.\] Then \[k=\frac{m_1}{p},\] so
\[\widetilde k_{\mathrm{MOM}}
=\frac{m_1}{\widetilde p_{\mathrm{MOM}}}.\]

\end{tcolorbox}

\begin{tcolorbox}[enhanced jigsaw, arc=.35mm, bottomrule=.15mm, bottomtitle=1mm, breakable, colback=white, colbacktitle=quarto-callout-warning-color!10!white, colframe=quarto-callout-warning-color-frame, coltitle=black, left=2mm, leftrule=.75mm, opacityback=0, opacitybacktitle=0.6, rightrule=.15mm, title=\textcolor{quarto-callout-warning-color}{\faExclamationTriangle}\hspace{0.5em}{Warning}, titlerule=0mm, toprule=.15mm, toptitle=1mm]

Important remark For the normal distribution, method of moments gives
estimators that agree with intuition. For the binomial model with both
\(k\) and \(p\) unknown, the method of moments estimators may behave
poorly. For some data sets, the formulas may even produce impossible
parameter values, such as negative estimates or estimates outside the
valid parameter space.

\end{tcolorbox}

Another famous use of moment matching is the Satterthwaite
approximation, where a complicated random quantity is approximated by a
scaled chi-square distribution by matching moments.

\bookmarksetup{startatroot}

\chapter{Maximum Likelihood
Estimation}\label{maximum-likelihood-estimation}

This section introduces maximum likelihood estimation, the most widely
used general method for deriving point estimators.

\begin{tcolorbox}[enhanced jigsaw, arc=.35mm, bottomrule=.15mm, bottomtitle=1mm, breakable, colback=white, colbacktitle=quarto-callout-note-color!10!white, colframe=quarto-callout-note-color-frame, coltitle=black, left=2mm, leftrule=.75mm, opacityback=0, opacitybacktitle=0.6, rightrule=.15mm, title=\textcolor{quarto-callout-note-color}{\faInfo}\hspace{0.5em}{Definition}, titlerule=0mm, toprule=.15mm, toptitle=1mm]

\textbf{Definition 8} (Likelihood function). Let \(X_1,\ldots,X_n\) be a
random sample from a population distribution with pdf or pmf
\(f(x\mid \boldsymbol\theta)\), where
\[\boldsymbol\theta=(\theta_1,\ldots,\theta_k).\] Given observed data
\(x_1,\ldots,x_n\), the \textbf{likelihood function} is
\[L(\boldsymbol\theta\mid \mathbf{x})
=\prod_{i=1}^n f(x_i\mid \boldsymbol\theta)
=f(x_1,\ldots,x_n\mid \theta_1,\ldots,\theta_k).\]

\end{tcolorbox}

The likelihood function treats the observed data as fixed and the
parameter as the variable.

\begin{tcolorbox}[enhanced jigsaw, arc=.35mm, bottomrule=.15mm, bottomtitle=1mm, breakable, colback=white, colbacktitle=quarto-callout-note-color!10!white, colframe=quarto-callout-note-color-frame, coltitle=black, left=2mm, leftrule=.75mm, opacityback=0, opacitybacktitle=0.6, rightrule=.15mm, title=\textcolor{quarto-callout-note-color}{\faInfo}\hspace{0.5em}{Definition}, titlerule=0mm, toprule=.15mm, toptitle=1mm]

\textbf{Definition 9} (Maximum likelihood estimator). A \textbf{maximum
likelihood estimator} is a value of the parameter that maximizes the
likelihood function:
\[\widehat{\boldsymbol\theta}_{\mathrm{MLE}}(\mathbf{x})
=\arg\max_{\boldsymbol\theta} L(\boldsymbol\theta\mid \mathbf{x}).\]

\end{tcolorbox}

Since logarithm is an increasing function, maximizing the likelihood is
equivalent to maximizing the log-likelihood.

\begin{tcolorbox}[enhanced jigsaw, arc=.35mm, bottomrule=.15mm, bottomtitle=1mm, breakable, colback=white, colbacktitle=quarto-callout-note-color!10!white, colframe=quarto-callout-note-color-frame, coltitle=black, left=2mm, leftrule=.75mm, opacityback=0, opacitybacktitle=0.6, rightrule=.15mm, title=\textcolor{quarto-callout-note-color}{\faInfo}\hspace{0.5em}{Definition}, titlerule=0mm, toprule=.15mm, toptitle=1mm]

\textbf{Definition 10} (Log-likelihood). The \textbf{log-likelihood
function} is \[\ell(\boldsymbol\theta\mid \mathbf{x})
=\log L(\boldsymbol\theta\mid \mathbf{x}).\] Then
\[\arg\max_{\boldsymbol\theta} L(\boldsymbol\theta\mid \mathbf{x})
=
\arg\max_{\boldsymbol\theta} \ell(\boldsymbol\theta\mid \mathbf{x}).\]

\end{tcolorbox}

\begin{tcolorbox}[enhanced jigsaw, arc=.35mm, bottomrule=.15mm, bottomtitle=1mm, breakable, colback=white, colbacktitle=quarto-callout-tip-color!10!white, colframe=quarto-callout-tip-color-frame, coltitle=black, left=2mm, leftrule=.75mm, opacityback=0, opacitybacktitle=0.6, rightrule=.15mm, title=\textcolor{quarto-callout-tip-color}{\faLightbulb}\hspace{0.5em}{Key idea}, titlerule=0mm, toprule=.15mm, toptitle=1mm]

Calculus method for MLE When the parameter space is continuous and the
maximum occurs in the interior, solve
\[\frac{\partial}{\partial \theta_i}\ell(\boldsymbol\theta\mid \mathbf{x})=0,
\qquad i=1,\ldots,k.\] Then check whether the critical point gives a
local or global maximum.

\end{tcolorbox}

\section{Review: Hessian and the Second Derivative
Test}\label{review-hessian-and-the-second-derivative-test}

This subsection recalls the multivariable calculus tool used to classify
critical points.

\begin{tcolorbox}[enhanced jigsaw, arc=.35mm, bottomrule=.15mm, bottomtitle=1mm, breakable, colback=white, colbacktitle=quarto-callout-note-color!10!white, colframe=quarto-callout-note-color-frame, coltitle=black, left=2mm, leftrule=.75mm, opacityback=0, opacitybacktitle=0.6, rightrule=.15mm, title=\textcolor{quarto-callout-note-color}{\faInfo}\hspace{0.5em}{Definition}, titlerule=0mm, toprule=.15mm, toptitle=1mm]

\textbf{Definition 11} (Hessian matrix). For a smooth function
\(f:\mathbb{R}^n\to \mathbb{R}\), the \textbf{Hessian matrix} is
\[Hf(\mathbf{x})=
\begin{pmatrix}
\frac{\partial^2 f}{\partial x_1^2} & \cdots & \frac{\partial^2 f}{\partial x_1\partial x_n}\\
\vdots & \ddots & \vdots\\
\frac{\partial^2 f}{\partial x_n\partial x_1} & \cdots & \frac{\partial^2 f}{\partial x_n^2}
\end{pmatrix}.\]

\end{tcolorbox}

\begin{tcolorbox}[enhanced jigsaw, arc=.35mm, bottomrule=.15mm, bottomtitle=1mm, breakable, colback=white, colbacktitle=quarto-callout-important-color!10!white, colframe=quarto-callout-important-color-frame, coltitle=black, left=2mm, leftrule=.75mm, opacityback=0, opacitybacktitle=0.6, rightrule=.15mm, title=\textcolor{quarto-callout-important-color}{\faExclamation}\hspace{0.5em}{Theorem}, titlerule=0mm, toprule=.15mm, toptitle=1mm]

\textbf{Theorem 12} (Second derivative test). \emph{Let
\(f:\mathbb{R}^n\to \mathbb{R}\) be smooth and suppose
\(\nabla f(\mathbf{a})=0\).}

\begin{enumerate}
\def\labelenumi{\arabic{enumi}.}
\item
  \emph{If \(Hf(\mathbf{a})\) is positive definite, then \(\mathbf{a}\)
  is a local minimum.}
\item
  \emph{If \(Hf(\mathbf{a})\) is negative definite, then \(\mathbf{a}\)
  is a local maximum.}
\item
  \emph{If \(Hf(\mathbf{a})\) has both positive and negative
  eigenvalues, then \(\mathbf{a}\) is a saddle point.}
\item
  \emph{Other cases require additional analysis.}
\end{enumerate}

\end{tcolorbox}

\section{MLE for the Normal
Distribution}\label{mle-for-the-normal-distribution}

This subsection derives the MLE for the mean and variance of a normal
distribution.

\begin{tcolorbox}[enhanced jigsaw, arc=.35mm, bottomrule=.15mm, bottomtitle=1mm, breakable, colback=white, colbacktitle=quarto-callout-note-color!10!white, colframe=quarto-callout-note-color-frame, coltitle=black, left=2mm, leftrule=.75mm, opacityback=0, opacitybacktitle=0.6, rightrule=.15mm, title=\textcolor{quarto-callout-note-color}{\faInfo}\hspace{0.5em}{Example}, titlerule=0mm, toprule=.15mm, toptitle=1mm]

\textbf{Example 13} (Normal distribution). Suppose
\[X_1,\ldots,X_n \overset{\text{iid}}{\sim}\operatorname{Normal}(\mu,\sigma^2).\]
Find the MLEs of \(\mu\) and \(\sigma^2\).

\end{tcolorbox}

\begin{tcolorbox}[enhanced jigsaw, arc=.35mm, bottomrule=.15mm, bottomtitle=1mm, breakable, colback=white, colbacktitle=quarto-callout-tip-color!10!white, colframe=quarto-callout-tip-color-frame, coltitle=black, left=2mm, leftrule=.75mm, opacityback=0, opacitybacktitle=0.6, rightrule=.15mm, title=\textcolor{quarto-callout-tip-color}{\faLightbulb}\hspace{0.5em}{Solution}, titlerule=0mm, toprule=.15mm, toptitle=1mm]

\textbf{Step 1: Write the likelihood.} The density is
\[f(x_i\mid \mu,\sigma^2)=\frac{1}{\sqrt{2\pi\sigma^2}}
\exp\left\{-\frac{(x_i-\mu)^2}{2\sigma^2}\right\}.\] Thus
\[L(\mu,\sigma^2\mid \mathbf{x})
=\prod_{i=1}^n \frac{1}{\sqrt{2\pi\sigma^2}}
\exp\left\{-\frac{(x_i-\mu)^2}{2\sigma^2}\right\}.\]

\textbf{Step 2: Write the log-likelihood.}
\[\ell(\mu,\sigma^2\mid \mathbf{x})
= -\frac{n}{2}\log(2\pi)-\frac{n}{2}\log(\sigma^2)
-\frac{1}{2\sigma^2}\sum_{i=1}^n (x_i-\mu)^2.\]

\textbf{Step 3: Differentiate with respect to \(\mu\).}
\[\frac{\partial}{\partial \mu}\ell(\mu,\sigma^2\mid \mathbf{x})
=\frac{1}{\sigma^2}\sum_{i=1}^n (x_i-\mu)
=\frac{n}{\sigma^2}(\overline x-\mu).\] Setting this equal to zero gives
\[\widehat\mu_{\mathrm{MLE}}=\overline x=\frac{1}{n}\sum_{i=1}^n x_i.\]

\textbf{Step 4: Differentiate with respect to \(\sigma^2\).} Treat
\(\sigma^2\) as the parameter. Then
\[\frac{\partial}{\partial \sigma^2}\ell(\mu,\sigma^2\mid \mathbf{x})
= -\frac{n}{2\sigma^2}+\frac{1}{2(\sigma^2)^2}\sum_{i=1}^n (x_i-\mu)^2.\]
Setting this equal to zero gives
\[\sigma^2=\frac{1}{n}\sum_{i=1}^n (x_i-\mu)^2.\] Replacing \(\mu\) by
\(\widehat\mu_{\mathrm{MLE}}=\overline x\), we obtain
\[\widehat\sigma^2_{\mathrm{MLE}}
=\frac{1}{n}\sum_{i=1}^n (x_i-\overline x)^2.\] This estimator divides
by \(n\), not by \(n-1\), so it is biased for \(\sigma^2\).

\end{tcolorbox}

\section{Exercise: Log-Normal
Distribution}\label{exercise-log-normal-distribution}

This subsection applies the normal MLE calculation after a logarithmic
transformation.

\begin{tcolorbox}[enhanced jigsaw, arc=.35mm, bottomrule=.15mm, bottomtitle=1mm, breakable, colback=white, colbacktitle=quarto-callout-warning-color!10!white, colframe=quarto-callout-warning-color-frame, coltitle=black, left=2mm, leftrule=.75mm, opacityback=0, opacitybacktitle=0.6, rightrule=.15mm, title=\textcolor{quarto-callout-warning-color}{\faExclamationTriangle}\hspace{0.5em}{Practice Problem}, titlerule=0mm, toprule=.15mm, toptitle=1mm]

\textbf{Practice Problem 14} (Log-normal MLE). Suppose
\[X_1,\ldots,X_n \overset{\text{iid}}{\sim}\operatorname{LogNormal}(\mu,\sigma^2),\]
so that
\[Y_i=\ln X_i \overset{\text{iid}}{\sim}\operatorname{Normal}(\mu,\sigma^2).\]
The density of \(X_i\) is \[f_X(x_i)=\frac{1}{x_i\sqrt{2\pi\sigma^2}}
\exp\left\{-\frac{(\ln x_i-\mu)^2}{2\sigma^2}\right\},
\qquad x_i>0.\] Find the MLEs of \(\mu\) and \(\sigma^2\).

\end{tcolorbox}

\begin{tcolorbox}[enhanced jigsaw, arc=.35mm, bottomrule=.15mm, bottomtitle=1mm, breakable, colback=white, colbacktitle=quarto-callout-tip-color!10!white, colframe=quarto-callout-tip-color-frame, coltitle=black, left=2mm, leftrule=.75mm, opacityback=0, opacitybacktitle=0.6, rightrule=.15mm, title=\textcolor{quarto-callout-tip-color}{\faLightbulb}\hspace{0.5em}{Solution}, titlerule=0mm, toprule=.15mm, toptitle=1mm]

Since \(Y_i=\ln X_i\) is normal, we can apply the normal MLE result to
the transformed data \[y_i=\ln x_i.\] Thus
\[\widehat\mu_{\mathrm{MLE}}=\overline y=\frac{1}{n}\sum_{i=1}^n \ln x_i,\]
and \[\widehat\sigma^2_{\mathrm{MLE}}
=\frac{1}{n}\sum_{i=1}^n (\ln x_i-\overline y)^2.\] The extra factor
\(1/x_i\) in the log-normal density does not depend on \(\mu\) or
\(\sigma^2\), so it does not change the maximizer.

\end{tcolorbox}

\section{MLE for Bernoulli and Binomial
Models}\label{mle-for-bernoulli-and-binomial-models}

This subsection derives MLEs for common discrete models.

\begin{tcolorbox}[enhanced jigsaw, arc=.35mm, bottomrule=.15mm, bottomtitle=1mm, breakable, colback=white, colbacktitle=quarto-callout-note-color!10!white, colframe=quarto-callout-note-color-frame, coltitle=black, left=2mm, leftrule=.75mm, opacityback=0, opacitybacktitle=0.6, rightrule=.15mm, title=\textcolor{quarto-callout-note-color}{\faInfo}\hspace{0.5em}{Example}, titlerule=0mm, toprule=.15mm, toptitle=1mm]

\textbf{Example 15} (Bernoulli distribution). Suppose
\[X_1,\ldots,X_n \overset{\text{iid}}{\sim}\operatorname{Bernoulli}(p).\]
Find the MLE of \(p\).

\end{tcolorbox}

\begin{tcolorbox}[enhanced jigsaw, arc=.35mm, bottomrule=.15mm, bottomtitle=1mm, breakable, colback=white, colbacktitle=quarto-callout-tip-color!10!white, colframe=quarto-callout-tip-color-frame, coltitle=black, left=2mm, leftrule=.75mm, opacityback=0, opacitybacktitle=0.6, rightrule=.15mm, title=\textcolor{quarto-callout-tip-color}{\faLightbulb}\hspace{0.5em}{Solution}, titlerule=0mm, toprule=.15mm, toptitle=1mm]

\textbf{Step 1: Likelihood.}
\[L(p\mid \mathbf{x})=\prod_{i=1}^n p^{x_i}(1-p)^{1-x_i}.\] Let
\[S=\sum_{i=1}^n x_i\] be the total number of successes. Then
\[L(p\mid \mathbf{x})=p^S(1-p)^{n-S}.\]

\textbf{Step 2: Log-likelihood.}
\[\ell(p\mid \mathbf{x})=S\log p+(n-S)\log(1-p).\]

\textbf{Step 3: Differentiate and solve.}
\[\frac{\partial}{\partial p}\ell(p\mid \mathbf{x})
=\frac{S}{p}-\frac{n-S}{1-p}.\] Setting this equal to zero gives
\[S(1-p)=p(n-S).\] Hence
\[\widehat p_{\mathrm{MLE}}=\frac{S}{n}=\frac{1}{n}\sum_{i=1}^n x_i=\overline x.\]

\end{tcolorbox}

\begin{tcolorbox}[enhanced jigsaw, arc=.35mm, bottomrule=.15mm, bottomtitle=1mm, breakable, colback=white, colbacktitle=quarto-callout-note-color!10!white, colframe=quarto-callout-note-color-frame, coltitle=black, left=2mm, leftrule=.75mm, opacityback=0, opacitybacktitle=0.6, rightrule=.15mm, title=\textcolor{quarto-callout-note-color}{\faInfo}\hspace{0.5em}{Example}, titlerule=0mm, toprule=.15mm, toptitle=1mm]

\textbf{Example 16} (Binomial distribution with known number of trials).
Suppose
\[X_1,\ldots,X_n \overset{\text{iid}}{\sim}\operatorname{Binomial}(k,p),\]
where \(k\) is known and \(p\) is unknown. Find the MLE of \(p\).

\end{tcolorbox}

\begin{tcolorbox}[enhanced jigsaw, arc=.35mm, bottomrule=.15mm, bottomtitle=1mm, breakable, colback=white, colbacktitle=quarto-callout-tip-color!10!white, colframe=quarto-callout-tip-color-frame, coltitle=black, left=2mm, leftrule=.75mm, opacityback=0, opacitybacktitle=0.6, rightrule=.15mm, title=\textcolor{quarto-callout-tip-color}{\faLightbulb}\hspace{0.5em}{Solution}, titlerule=0mm, toprule=.15mm, toptitle=1mm]

The likelihood is proportional to
\[L(p\mid \mathbf{x})\propto \prod_{i=1}^n p^{x_i}(1-p)^{k-x_i}
=p^{\sum x_i}(1-p)^{kn-\sum x_i}.\] The log-likelihood is
\[\ell(p\mid \mathbf{x})=S\log p+(kn-S)\log(1-p)+\text{constant},\]
where \(S=\sum_{i=1}^n x_i\). Differentiating gives
\[\frac{S}{p}-\frac{kn-S}{1-p}=0.\] Therefore
\[\widehat p_{\mathrm{MLE}}=\frac{S}{kn}=\frac{\sum_{i=1}^n x_i}{kn}.\]
This is the total number of observed successes divided by the total
number of trials.

\end{tcolorbox}

\section{Binomial MLE with Unknown Number of
Trials}\label{binomial-mle-with-unknown-number-of-trials}

This subsection describes a harder binomial MLE problem where the
parameter is discrete.

\begin{tcolorbox}[enhanced jigsaw, arc=.35mm, bottomrule=.15mm, bottomtitle=1mm, breakable, colback=white, colbacktitle=quarto-callout-note-color!10!white, colframe=quarto-callout-note-color-frame, coltitle=black, left=2mm, leftrule=.75mm, opacityback=0, opacitybacktitle=0.6, rightrule=.15mm, title=\textcolor{quarto-callout-note-color}{\faInfo}\hspace{0.5em}{Example}, titlerule=0mm, toprule=.15mm, toptitle=1mm]

\textbf{Example 17} (Binomial model with known \(p\) and unknown \(k\)).
Suppose
\[X_1,\ldots,X_n \overset{\text{iid}}{\sim}\operatorname{Binomial}(k,p),\]
where \(p\) is known and \(k\) is unknown. The likelihood is
\[L(k\mid \mathbf{x},p)=\prod_{i=1}^n {k\choose x_i}p^{x_i}(1-p)^{k-x_i}.\]
Explain why this MLE problem is not solved by ordinary differentiation.

\end{tcolorbox}

\begin{tcolorbox}[enhanced jigsaw, arc=.35mm, bottomrule=.15mm, bottomtitle=1mm, breakable, colback=white, colbacktitle=quarto-callout-tip-color!10!white, colframe=quarto-callout-tip-color-frame, coltitle=black, left=2mm, leftrule=.75mm, opacityback=0, opacitybacktitle=0.6, rightrule=.15mm, title=\textcolor{quarto-callout-tip-color}{\faLightbulb}\hspace{0.5em}{Solution}, titlerule=0mm, toprule=.15mm, toptitle=1mm]

The parameter \(k\) is a positive integer, not a continuous real number.
The likelihood contains binomial coefficients involving factorials,
\[{k\choose x_i}=\frac{k!}{x_i!(k-x_i)!},\] so ordinary calculus with
respect to \(k\) is not the natural tool. Instead, we compare likelihood
values at neighboring integers. A common approach is to study the
likelihood ratio
\[\frac{L(k\mid \mathbf{x},p)}{L(k-1\mid \mathbf{x},p)}.\] The
likelihood increases while this ratio is at least \(1\) and decreases
after it falls below \(1\). Therefore the MLE is found by discrete
optimization over integer values satisfying \[k\geq \max_i x_i.\]

\end{tcolorbox}

The lecture notes also mention an equivalent transformation \(z=1/k\),
which turns the likelihood condition into an equation in \(z\) on the
interval \[0<z<\frac{1}{\max_i x_i}.\] The relevant function is strictly
decreasing, so a unique solution \(\widehat z\) exists, and the
corresponding estimate is
\[\widehat k_{\mathrm{MLE}}=\frac{1}{\widehat z}.\] In practice, because
\(k\) must be an integer, one checks the nearest admissible integer
values.

\section{MLE for a Scale Uniform
Distribution}\label{mle-for-a-scale-uniform-distribution}

This subsection gives an example where the maximum occurs at the
boundary of the parameter space.

\begin{tcolorbox}[enhanced jigsaw, arc=.35mm, bottomrule=.15mm, bottomtitle=1mm, breakable, colback=white, colbacktitle=quarto-callout-note-color!10!white, colframe=quarto-callout-note-color-frame, coltitle=black, left=2mm, leftrule=.75mm, opacityback=0, opacitybacktitle=0.6, rightrule=.15mm, title=\textcolor{quarto-callout-note-color}{\faInfo}\hspace{0.5em}{Example}, titlerule=0mm, toprule=.15mm, toptitle=1mm]

\textbf{Example 18} (Uniform distribution on \((0,\theta)\) ). Suppose
\[X_1,\ldots,X_n \overset{\text{iid}}{\sim}\operatorname{Uniform}(0,\theta),\]
where \(\theta>0\) is unknown. Find the MLE of \(\theta\).

\end{tcolorbox}

\begin{tcolorbox}[enhanced jigsaw, arc=.35mm, bottomrule=.15mm, bottomtitle=1mm, breakable, colback=white, colbacktitle=quarto-callout-tip-color!10!white, colframe=quarto-callout-tip-color-frame, coltitle=black, left=2mm, leftrule=.75mm, opacityback=0, opacitybacktitle=0.6, rightrule=.15mm, title=\textcolor{quarto-callout-tip-color}{\faLightbulb}\hspace{0.5em}{Solution}, titlerule=0mm, toprule=.15mm, toptitle=1mm]

The density is
\[f(x\mid \theta)=\frac{1}{\theta}, \qquad 0\leq x\leq \theta.\] The
likelihood is \[L(\theta\mid \mathbf{x})
=\prod_{i=1}^n \frac{1}{\theta}\mathbbm{1}\{0\leq x_i\leq \theta\}
=\theta^{-n}\mathbbm{1}\{\theta\geq x_{(n)}\},\] where
\[x_{(n)}=\max\{x_1,\ldots,x_n\}.\] For all feasible
\(\theta\geq x_{(n)}\), the function \(\theta^{-n}\) is decreasing in
\(\theta\). Therefore it is maximized by choosing the smallest feasible
value: \[\widehat\theta_{\mathrm{MLE}}=x_{(n)}=\max\{x_1,\ldots,x_n\}.\]

\end{tcolorbox}

\begin{tcolorbox}[enhanced jigsaw, arc=.35mm, bottomrule=.15mm, bottomtitle=1mm, breakable, colback=white, colbacktitle=quarto-callout-warning-color!10!white, colframe=quarto-callout-warning-color-frame, coltitle=black, left=2mm, leftrule=.75mm, opacityback=0, opacitybacktitle=0.6, rightrule=.15mm, title=\textcolor{quarto-callout-warning-color}{\faExclamationTriangle}\hspace{0.5em}{Warning}, titlerule=0mm, toprule=.15mm, toptitle=1mm]

Boundary maximum In the uniform example, the derivative method does not
find the answer because the maximum occurs at the boundary
\(\theta=x_{(n)}\), not at an interior critical point.

\end{tcolorbox}

\bookmarksetup{startatroot}

\chapter{Invariance Property of Maximum Likelihood
Estimators}\label{invariance-property-of-maximum-likelihood-estimators}

This section explains how to estimate a function of a parameter once the
MLE of the parameter is known.

\begin{tcolorbox}[enhanced jigsaw, arc=.35mm, bottomrule=.15mm, bottomtitle=1mm, breakable, colback=white, colbacktitle=quarto-callout-important-color!10!white, colframe=quarto-callout-important-color-frame, coltitle=black, left=2mm, leftrule=.75mm, opacityback=0, opacitybacktitle=0.6, rightrule=.15mm, title=\textcolor{quarto-callout-important-color}{\faExclamation}\hspace{0.5em}{Theorem}, titlerule=0mm, toprule=.15mm, toptitle=1mm]

\textbf{Theorem 19} (Invariance property of MLE). \emph{Suppose
\(\widehat\theta\) is the MLE of \(\theta\). Then, for any function
\(g(\theta)\), the MLE of \(g(\theta)\) is \[g(\widehat\theta).\]}

\end{tcolorbox}

\begin{tcolorbox}[enhanced jigsaw, arc=.35mm, bottomrule=.15mm, bottomtitle=1mm, breakable, colback=white, colbacktitle=quarto-callout-note-color!10!white, colframe=quarto-callout-note-color-frame, coltitle=black, left=2mm, leftrule=.75mm, opacityback=0, opacitybacktitle=0.6, rightrule=.15mm, title=\textcolor{quarto-callout-note-color}{\faInfo}\hspace{0.5em}{Proof}, titlerule=0mm, toprule=.15mm, toptitle=1mm]

\emph{Proof.} The MLE chooses the value of \(\theta\) that maximizes the
likelihood. Reparametrizing by \(\phi=g(\theta)\) changes the label of
the parameter but not the height of the likelihood curve. Therefore the
maximizing value of the transformed parameter is the transformation of
the maximizing value of the original parameter.~◻

\end{tcolorbox}

\begin{tcolorbox}[enhanced jigsaw, arc=.35mm, bottomrule=.15mm, bottomtitle=1mm, breakable, colback=white, colbacktitle=quarto-callout-note-color!10!white, colframe=quarto-callout-note-color-frame, coltitle=black, left=2mm, leftrule=.75mm, opacityback=0, opacitybacktitle=0.6, rightrule=.15mm, title=\textcolor{quarto-callout-note-color}{\faInfo}\hspace{0.5em}{Example}, titlerule=0mm, toprule=.15mm, toptitle=1mm]

\textbf{Example 20} (MLE of the normal standard deviation). Suppose
\[X_1,\ldots,X_n\overset{\text{iid}}{\sim}\operatorname{Normal}(\mu,\sigma^2).\]
We already found
\[\widehat\sigma^2_{\mathrm{MLE}}=\frac{1}{n}\sum_{i=1}^n (x_i-\overline x)^2.\]
Find the MLE of \(\sigma\).

\end{tcolorbox}

\begin{tcolorbox}[enhanced jigsaw, arc=.35mm, bottomrule=.15mm, bottomtitle=1mm, breakable, colback=white, colbacktitle=quarto-callout-tip-color!10!white, colframe=quarto-callout-tip-color-frame, coltitle=black, left=2mm, leftrule=.75mm, opacityback=0, opacitybacktitle=0.6, rightrule=.15mm, title=\textcolor{quarto-callout-tip-color}{\faLightbulb}\hspace{0.5em}{Solution}, titlerule=0mm, toprule=.15mm, toptitle=1mm]

Since \(\sigma=g(\sigma^2)=\sqrt{\sigma^2}\), the invariance property
gives \[\widehat\sigma_{\mathrm{MLE}}
=\sqrt{\widehat\sigma^2_{\mathrm{MLE}}}
=\sqrt{\frac{1}{n}\sum_{i=1}^n (x_i-\overline x)^2}.\]

\end{tcolorbox}

\bookmarksetup{startatroot}

\chapter{Bayesian Estimators}\label{bayesian-estimators}

This section introduces the Bayesian approach to point estimation.

In the method of moments and maximum likelihood estimation, the
parameter \(\theta\) is considered unknown but fixed. In the Bayesian
approach, \(\theta\) is treated as an uncertain quantity described by a
probability distribution.

\begin{tcolorbox}[enhanced jigsaw, arc=.35mm, bottomrule=.15mm, bottomtitle=1mm, breakable, colback=white, colbacktitle=quarto-callout-note-color!10!white, colframe=quarto-callout-note-color-frame, coltitle=black, left=2mm, leftrule=.75mm, opacityback=0, opacitybacktitle=0.6, rightrule=.15mm, title=\textcolor{quarto-callout-note-color}{\faInfo}\hspace{0.5em}{Definition}, titlerule=0mm, toprule=.15mm, toptitle=1mm]

\textbf{Definition 21} (Prior and posterior). Let \(X_1,\ldots,X_n\) be
sampled from a population distribution with pdf or pmf
\(f(x\mid \theta)\). In Bayesian inference:

\begin{enumerate}
\def\labelenumi{\arabic{enumi}.}
\item
  The \textbf{prior distribution} \(\pi(\theta)\) represents belief
  about \(\theta\) before seeing the data.
\item
  The \textbf{posterior distribution} \(\pi(\theta\mid \mathbf{x})\)
  represents updated belief about \(\theta\) after observing the data
  \(\mathbf{x}=(x_1,\ldots,x_n)\).
\end{enumerate}

By Bayes' rule, \[\pi(\theta\mid \mathbf{x})
=\frac{f(\mathbf{x}\mid \theta)\pi(\theta)}{m(\mathbf{x})},\] where
\[m(\mathbf{x})=\int f(\mathbf{x}\mid \theta)\pi(\theta)\,d\theta\] is
the marginal distribution, or evidence, of the data.

\end{tcolorbox}

\begin{tcolorbox}[enhanced jigsaw, arc=.35mm, bottomrule=.15mm, bottomtitle=1mm, breakable, colback=white, colbacktitle=quarto-callout-tip-color!10!white, colframe=quarto-callout-tip-color-frame, coltitle=black, left=2mm, leftrule=.75mm, opacityback=0, opacitybacktitle=0.6, rightrule=.15mm, title=\textcolor{quarto-callout-tip-color}{\faLightbulb}\hspace{0.5em}{Key idea}, titlerule=0mm, toprule=.15mm, toptitle=1mm]

Bayesian update
\[\text{posterior} \; \propto \; \text{likelihood} \times \text{prior}.\]
Equivalently, \[P(\boldsymbol\theta\mid \mathcal D)
=\frac{P(\mathcal D\mid \boldsymbol\theta)P(\boldsymbol\theta)}{P(\mathcal D)}.\]

\end{tcolorbox}

Once the posterior distribution has been derived, several point
estimates are possible:

\begin{enumerate}
\def\labelenumi{\arabic{enumi}.}
\item
  posterior mean: \(\mathbb{E}[\theta\mid \mathbf{x}]\);
\item
  posterior median;
\item
  posterior mode, also called the MAP estimate.
\end{enumerate}

\begin{tcolorbox}[enhanced jigsaw, arc=.35mm, bottomrule=.15mm, bottomtitle=1mm, breakable, colback=white, colbacktitle=quarto-callout-note-color!10!white, colframe=quarto-callout-note-color-frame, coltitle=black, left=2mm, leftrule=.75mm, opacityback=0, opacitybacktitle=0.6, rightrule=.15mm, title=\textcolor{quarto-callout-note-color}{\faInfo}\hspace{0.5em}{Definition}, titlerule=0mm, toprule=.15mm, toptitle=1mm]

\textbf{Definition 22} (MAP estimator). The \textbf{maximum a
posteriori} estimator is
\[\widehat\theta_{\mathrm{MAP}}=\arg\max_\theta \pi(\theta\mid \mathbf{x}).\]

\end{tcolorbox}

\section{Bayesian Coin Tossing: Beta-Binomial
Model}\label{bayesian-coin-tossing-beta-binomial-model}

This subsection studies the standard Bayesian model for an unknown coin
probability.

\begin{tcolorbox}[enhanced jigsaw, arc=.35mm, bottomrule=.15mm, bottomtitle=1mm, breakable, colback=white, colbacktitle=quarto-callout-note-color!10!white, colframe=quarto-callout-note-color-frame, coltitle=black, left=2mm, leftrule=.75mm, opacityback=0, opacitybacktitle=0.6, rightrule=.15mm, title=\textcolor{quarto-callout-note-color}{\faInfo}\hspace{0.5em}{Example}, titlerule=0mm, toprule=.15mm, toptitle=1mm]

\textbf{Example 23} (Bayesian estimation for a coin). Suppose a coin has
unknown probability \(\theta\) of heads. We toss the coin \(n\) times
and observe \(x\) heads. The likelihood is
\[P(x\mid n,\theta)={n\choose x}\theta^x(1-\theta)^{n-x}.\] Use a beta
prior \[\theta\sim \operatorname{Beta}(\alpha,\beta).\] Find the
posterior distribution.

\end{tcolorbox}

\begin{tcolorbox}[enhanced jigsaw, arc=.35mm, bottomrule=.15mm, bottomtitle=1mm, breakable, colback=white, colbacktitle=quarto-callout-tip-color!10!white, colframe=quarto-callout-tip-color-frame, coltitle=black, left=2mm, leftrule=.75mm, opacityback=0, opacitybacktitle=0.6, rightrule=.15mm, title=\textcolor{quarto-callout-tip-color}{\faLightbulb}\hspace{0.5em}{Solution}, titlerule=0mm, toprule=.15mm, toptitle=1mm]

The beta prior has density \[p(\theta\mid \alpha,\beta)
=\frac{1}{B(\alpha,\beta)}\theta^{\alpha-1}(1-\theta)^{\beta-1},
\qquad 0<\theta<1,\] where
\[B(\alpha,\beta)=\int_0^1 \theta^{\alpha-1}(1-\theta)^{\beta-1}\,d\theta\]
and equivalently
\[B(\alpha,\beta)=\frac{\Gamma(\alpha)\Gamma(\beta)}{\Gamma(\alpha+\beta)}.\]
Using Bayes' rule, \[p(\theta\mid x,n,\alpha,\beta)
\propto P(x\mid n,\theta)p(\theta\mid \alpha,\beta).\] Therefore
\[p(\theta\mid x,n,\alpha,\beta)
\propto
\theta^x(1-\theta)^{n-x}\theta^{\alpha-1}(1-\theta)^{\beta-1}.\]
Combining powers gives \[p(\theta\mid x,n,\alpha,\beta)
\propto
\theta^{\alpha+x-1}(1-\theta)^{\beta+n-x-1}.\] Thus
\[\theta\mid x,n,\alpha,\beta \sim \operatorname{Beta}(\alpha+x,\beta+n-x).\]
The normalized posterior density is \[p(\theta\mid x,n,\alpha,\beta)
=\frac{1}{B(\alpha+x,\beta+n-x)}
\theta^{\alpha+x-1}(1-\theta)^{\beta+n-x-1}.\]

\end{tcolorbox}

\begin{tcolorbox}[enhanced jigsaw, arc=.35mm, bottomrule=.15mm, bottomtitle=1mm, breakable, colback=white, colbacktitle=quarto-callout-tip-color!10!white, colframe=quarto-callout-tip-color-frame, coltitle=black, left=2mm, leftrule=.75mm, opacityback=0, opacitybacktitle=0.6, rightrule=.15mm, title=\textcolor{quarto-callout-tip-color}{\faLightbulb}\hspace{0.5em}{Key idea}, titlerule=0mm, toprule=.15mm, toptitle=1mm]

Pseudocount interpretation The beta hyperparameters \(\alpha\) and
\(\beta\) act like prior pseudocounts. After observing \(x\) heads and
\(n-x\) tails, the posterior parameters become
\[\alpha \longmapsto \alpha+x,
\qquad
\beta \longmapsto \beta+n-x.\]

\end{tcolorbox}

\begin{tcolorbox}[enhanced jigsaw, arc=.35mm, bottomrule=.15mm, bottomtitle=1mm, breakable, colback=white, colbacktitle=quarto-callout-note-color!10!white, colframe=quarto-callout-note-color-frame, coltitle=black, left=2mm, leftrule=.75mm, opacityback=0, opacitybacktitle=0.6, rightrule=.15mm, title=\textcolor{quarto-callout-note-color}{\faInfo}\hspace{0.5em}{Example}, titlerule=0mm, toprule=.15mm, toptitle=1mm]

\textbf{Example 24} (Prior, likelihood, posterior). Suppose
\[(\alpha,\beta)=(3,5), \qquad (x,n)=(5,6).\] Find the posterior
distribution.

\end{tcolorbox}

\begin{tcolorbox}[enhanced jigsaw, arc=.35mm, bottomrule=.15mm, bottomtitle=1mm, breakable, colback=white, colbacktitle=quarto-callout-tip-color!10!white, colframe=quarto-callout-tip-color-frame, coltitle=black, left=2mm, leftrule=.75mm, opacityback=0, opacitybacktitle=0.6, rightrule=.15mm, title=\textcolor{quarto-callout-tip-color}{\faLightbulb}\hspace{0.5em}{Solution}, titlerule=0mm, toprule=.15mm, toptitle=1mm]

The posterior is
\[\theta\mid x,n,\alpha,\beta \sim \operatorname{Beta}(\alpha+x,\beta+n-x).\]
Substituting the values gives
\[\theta\mid \mathcal D \sim \operatorname{Beta}(3+5,5+6-5)=\operatorname{Beta}(8,6).\]
The posterior distribution is shifted toward larger values of \(\theta\)
because \(5\) heads were observed in \(6\) tosses.

\end{tcolorbox}

\section{Posterior Mean, Variance, Mode, and Credible
Probability}\label{posterior-mean-variance-mode-and-credible-probability}

This subsection extracts useful point estimates and uncertainty
summaries from the beta posterior.

\begin{tcolorbox}[enhanced jigsaw, arc=.35mm, bottomrule=.15mm, bottomtitle=1mm, breakable, colback=white, colbacktitle=quarto-callout-important-color!10!white, colframe=quarto-callout-important-color-frame, coltitle=black, left=2mm, leftrule=.75mm, opacityback=0, opacitybacktitle=0.6, rightrule=.15mm, title=\textcolor{quarto-callout-important-color}{\faExclamation}\hspace{0.5em}{Theorem}, titlerule=0mm, toprule=.15mm, toptitle=1mm]

\textbf{Theorem 25} (Beta-binomial posterior summaries). \emph{If
\[\theta\mid \mathcal D\sim \operatorname{Beta}(\alpha+x,\beta+n-x),\]
then \[\mathbb{E}[\theta\mid \mathcal D]
=\frac{\alpha+x}{\alpha+\beta+n},\]
\[\operatorname{Var}(\theta\mid \mathcal D)
=\frac{(\alpha+x)(\beta+n-x)}{(\alpha+\beta+n)^2(\alpha+\beta+n+1)},\]
and, when \(\alpha+x>1\) and \(\beta+n-x>1\), the posterior mode is
\[\operatorname{Mode}(\theta\mid \mathcal D)
=\frac{\alpha+x-1}{\alpha+\beta+n-2}.\]}

\end{tcolorbox}

\begin{tcolorbox}[enhanced jigsaw, arc=.35mm, bottomrule=.15mm, bottomtitle=1mm, breakable, colback=white, colbacktitle=quarto-callout-note-color!10!white, colframe=quarto-callout-note-color-frame, coltitle=black, left=2mm, leftrule=.75mm, opacityback=0, opacitybacktitle=0.6, rightrule=.15mm, title=\textcolor{quarto-callout-note-color}{\faInfo}\hspace{0.5em}{Example}, titlerule=0mm, toprule=.15mm, toptitle=1mm]

\textbf{Example 26} (Posterior probability as a Bayesian test). In the
coin example, compute the posterior probability that the coin is biased
toward tails:
\[\mathbb{P}\left(\theta<\frac12\mid x,n,\alpha,\beta\right).\]

\end{tcolorbox}

\begin{tcolorbox}[enhanced jigsaw, arc=.35mm, bottomrule=.15mm, bottomtitle=1mm, breakable, colback=white, colbacktitle=quarto-callout-tip-color!10!white, colframe=quarto-callout-tip-color-frame, coltitle=black, left=2mm, leftrule=.75mm, opacityback=0, opacitybacktitle=0.6, rightrule=.15mm, title=\textcolor{quarto-callout-tip-color}{\faLightbulb}\hspace{0.5em}{Solution}, titlerule=0mm, toprule=.15mm, toptitle=1mm]

The posterior distribution is
\[\theta\mid x,n,\alpha,\beta\sim \operatorname{Beta}(\alpha+x,\beta+n-x).\]
Therefore \[\mathbb{P}\left(\theta<\frac12\mid x,n,\alpha,\beta\right)
=\int_0^{1/2}
\frac{1}{B(\alpha+x,\beta+n-x)}
\theta^{\alpha+x-1}(1-\theta)^{\beta+n-x-1}
\,d\theta.\] This is a posterior probability. It plays a role similar to
a hypothesis test, but its interpretation is directly probabilistic
under the Bayesian model.

\end{tcolorbox}

\bookmarksetup{startatroot}

\chapter{Conjugate Priors}\label{conjugate-priors}

This section explains why the beta prior is especially convenient for
binomial data.

\begin{tcolorbox}[enhanced jigsaw, arc=.35mm, bottomrule=.15mm, bottomtitle=1mm, breakable, colback=white, colbacktitle=quarto-callout-note-color!10!white, colframe=quarto-callout-note-color-frame, coltitle=black, left=2mm, leftrule=.75mm, opacityback=0, opacitybacktitle=0.6, rightrule=.15mm, title=\textcolor{quarto-callout-note-color}{\faInfo}\hspace{0.5em}{Definition}, titlerule=0mm, toprule=.15mm, toptitle=1mm]

\textbf{Definition 27} (Conjugate prior). For a likelihood
\(P(\mathcal D\mid \boldsymbol\theta)\), a prior
\(P(\boldsymbol\theta)\) is called a \textbf{conjugate prior} if the
posterior distribution \(P(\boldsymbol\theta\mid \mathcal D)\) belongs
to the same distribution family as the prior.

\end{tcolorbox}

\begin{tcolorbox}[enhanced jigsaw, arc=.35mm, bottomrule=.15mm, bottomtitle=1mm, breakable, colback=white, colbacktitle=quarto-callout-note-color!10!white, colframe=quarto-callout-note-color-frame, coltitle=black, left=2mm, leftrule=.75mm, opacityback=0, opacitybacktitle=0.6, rightrule=.15mm, title=\textcolor{quarto-callout-note-color}{\faInfo}\hspace{0.5em}{Example}, titlerule=0mm, toprule=.15mm, toptitle=1mm]

\textbf{Example 28} (Beta-binomial conjugacy). For binomial data,
\[P(x\mid n,\theta)={n\choose x}\theta^x(1-\theta)^{n-x},\] the beta
prior \[\theta\sim \operatorname{Beta}(\alpha,\beta)\] is conjugate
because the posterior is again beta:
\[\theta\mid x,n\sim \operatorname{Beta}(\alpha+x,\beta+n-x).\]

\end{tcolorbox}

\begin{tcolorbox}[enhanced jigsaw, arc=.35mm, bottomrule=.15mm, bottomtitle=1mm, breakable, colback=white, colbacktitle=quarto-callout-tip-color!10!white, colframe=quarto-callout-tip-color-frame, coltitle=black, left=2mm, leftrule=.75mm, opacityback=0, opacitybacktitle=0.6, rightrule=.15mm, title=\textcolor{quarto-callout-tip-color}{\faLightbulb}\hspace{0.5em}{Solution}, titlerule=0mm, toprule=.15mm, toptitle=1mm]

The posterior is obtained by multiplying the likelihood and prior:
\[\theta^x(1-\theta)^{n-x}\theta^{\alpha-1}(1-\theta)^{\beta-1}
=\theta^{\alpha+x-1}(1-\theta)^{\beta+n-x-1}.\] This has the beta form,
so the beta prior is conjugate to the binomial likelihood.

\end{tcolorbox}

\begin{tcolorbox}[enhanced jigsaw, arc=.35mm, bottomrule=.15mm, bottomtitle=1mm, breakable, colback=white, colbacktitle=quarto-callout-tip-color!10!white, colframe=quarto-callout-tip-color-frame, coltitle=black, left=2mm, leftrule=.75mm, opacityback=0, opacitybacktitle=0.6, rightrule=.15mm, title=\textcolor{quarto-callout-tip-color}{\faLightbulb}\hspace{0.5em}{Key idea}, titlerule=0mm, toprule=.15mm, toptitle=1mm]

Why conjugacy helps When the prior is conjugate, we can update
parameters without recomputing the full integral every time. In the
beta-binomial case,
\[(\alpha,\beta) \longrightarrow (\alpha+x,\beta+n-x).\] This makes
Bayesian updating fast and interpretable.

\end{tcolorbox}

For many likelihoods in the exponential family, conjugate priors also
exist. For example, the conjugate prior for the mean of a normal
distribution with known variance is also normal.

\bookmarksetup{startatroot}

\chapter{MAP versus MLE}\label{map-versus-mle}

This section compares maximum likelihood estimation with maximum a
posteriori estimation.

Suppose \(\boldsymbol\theta\) denotes model parameters and
\[\mathcal D=\{(\mathbf{x}_i,\mathbf{y}_i):i=1,\ldots,N\}\] denotes
observed data. The likelihood is
\[P(\mathcal D\mid \boldsymbol\theta).\] The MLE is
\[\widehat{\boldsymbol\theta}_{\mathrm{MLE}}
=\arg\max_{\boldsymbol\theta}P(\mathcal D\mid \boldsymbol\theta)
=\arg\max_{\boldsymbol\theta}\log P(\mathcal D\mid \boldsymbol\theta).\]

In Bayesian statistics, after choosing a prior \(P(\boldsymbol\theta)\),
the posterior is \[P(\boldsymbol\theta\mid \mathcal D)
=\frac{P(\mathcal D\mid \boldsymbol\theta)P(\boldsymbol\theta)}{P(\mathcal D)}.\]
The MAP estimator is \[\widehat{\boldsymbol\theta}_{\mathrm{MAP}}
=\arg\max_{\boldsymbol\theta} P(\boldsymbol\theta\mid \mathcal D).\]
Since \(P(\mathcal D)\) does not depend on \(\boldsymbol\theta\),
\[\widehat{\boldsymbol\theta}_{\mathrm{MAP}}
=\arg\max_{\boldsymbol\theta} P(\mathcal D\mid \boldsymbol\theta)P(\boldsymbol\theta)
=\arg\max_{\boldsymbol\theta}\left\{\log P(\mathcal D\mid \boldsymbol\theta)+\log P(\boldsymbol\theta)\right\}.\]

\begin{tcolorbox}[enhanced jigsaw, arc=.35mm, bottomrule=.15mm, bottomtitle=1mm, breakable, colback=white, colbacktitle=quarto-callout-tip-color!10!white, colframe=quarto-callout-tip-color-frame, coltitle=black, left=2mm, leftrule=.75mm, opacityback=0, opacitybacktitle=0.6, rightrule=.15mm, title=\textcolor{quarto-callout-tip-color}{\faLightbulb}\hspace{0.5em}{Key idea}, titlerule=0mm, toprule=.15mm, toptitle=1mm]

Comparison \[\text{MLE: maximize likelihood only.}\]
\[\text{MAP: maximize likelihood plus prior information.}\] The prior
term \(\log P(\boldsymbol\theta)\) can be viewed as a regularization
term.

\end{tcolorbox}

\section{MAP for the Mean of a Normal
Distribution}\label{map-for-the-mean-of-a-normal-distribution}

This subsection derives a normal-normal Bayesian estimator.

\begin{tcolorbox}[enhanced jigsaw, arc=.35mm, bottomrule=.15mm, bottomtitle=1mm, breakable, colback=white, colbacktitle=quarto-callout-note-color!10!white, colframe=quarto-callout-note-color-frame, coltitle=black, left=2mm, leftrule=.75mm, opacityback=0, opacitybacktitle=0.6, rightrule=.15mm, title=\textcolor{quarto-callout-note-color}{\faInfo}\hspace{0.5em}{Example}, titlerule=0mm, toprule=.15mm, toptitle=1mm]

\textbf{Example 29} (MAP for \(\mu\) with known variance). Suppose
\[x_1,\ldots,x_N\] are observed from
\[X_i\mid \mu \sim \operatorname{Normal}(\mu,\sigma^2),\] where
\(\sigma^2\) is known. Assume the prior
\[\mu\sim \operatorname{Normal}(\mu_0,\sigma_0^2).\] Find the MAP
estimate of \(\mu\).

\end{tcolorbox}

\begin{tcolorbox}[enhanced jigsaw, arc=.35mm, bottomrule=.15mm, bottomtitle=1mm, breakable, colback=white, colbacktitle=quarto-callout-tip-color!10!white, colframe=quarto-callout-tip-color-frame, coltitle=black, left=2mm, leftrule=.75mm, opacityback=0, opacitybacktitle=0.6, rightrule=.15mm, title=\textcolor{quarto-callout-tip-color}{\faLightbulb}\hspace{0.5em}{Solution}, titlerule=0mm, toprule=.15mm, toptitle=1mm]

The likelihood is \[P(\mathcal D\mid \mu)
=\prod_{i=1}^N \frac{1}{\sqrt{2\pi}\sigma}
\exp\left\{-\frac{(x_i-\mu)^2}{2\sigma^2}\right\}.\] The prior density
is \[P(\mu)=\frac{1}{\sqrt{2\pi}\sigma_0}
\exp\left\{-\frac{(\mu-\mu_0)^2}{2\sigma_0^2}\right\}.\] The MAP
estimate maximizes the posterior, equivalently minimizes the negative
log-posterior: \[\widehat\mu_{\mathrm{MAP}}
=\arg\min_\mu \left\{
\frac{1}{2\sigma_0^2}(\mu-\mu_0)^2
+\sum_{i=1}^N \frac{1}{2\sigma^2}(x_i-\mu)^2
\right\}.\] Differentiate with respect to \(\mu\):
\[\frac{\mu-\mu_0}{\sigma_0^2}
+\sum_{i=1}^N \frac{\mu-x_i}{\sigma^2}=0.\] Thus
\[\mu\left(\frac{1}{\sigma_0^2}+\frac{N}{\sigma^2}\right)
=\frac{\mu_0}{\sigma_0^2}+\frac{\sum_{i=1}^N x_i}{\sigma^2}.\] Solving
gives \[\widehat\mu_{\mathrm{MAP}}
=\frac{\sigma_0^2\sum_{i=1}^N x_i+\sigma^2\mu_0}{N\sigma_0^2+\sigma^2}.\]
Equivalently, \[\widehat\mu_{\mathrm{MAP}}
=\frac{N\sigma_0^2}{N\sigma_0^2+\sigma^2}\overline x
+\frac{\sigma^2}{N\sigma_0^2+\sigma^2}\mu_0.\] So the MAP estimate is a
weighted average of the sample mean and the prior mean.

The posterior variance is \[\operatorname{Var}(\mu\mid \mathbf{x})
=\left(\frac{1}{\sigma_0^2}+\frac{N}{\sigma^2}\right)^{-1}
=\frac{\sigma^2\sigma_0^2}{\sigma^2+N\sigma_0^2}.\] As
\(\sigma_0^2\to \infty\), the prior becomes very diffuse and
\[\widehat\mu_{\mathrm{MAP}}\to \overline x=\widehat\mu_{\mathrm{MLE}}.\]

\end{tcolorbox}

\bookmarksetup{startatroot}

\chapter{Practice Problems}\label{practice-problems-3}

This section gives additional problems that reinforce the main
estimation methods.

\begin{tcolorbox}[enhanced jigsaw, arc=.35mm, bottomrule=.15mm, bottomtitle=1mm, breakable, colback=white, colbacktitle=quarto-callout-warning-color!10!white, colframe=quarto-callout-warning-color-frame, coltitle=black, left=2mm, leftrule=.75mm, opacityback=0, opacitybacktitle=0.6, rightrule=.15mm, title=\textcolor{quarto-callout-warning-color}{\faExclamationTriangle}\hspace{0.5em}{Practice Problem}, titlerule=0mm, toprule=.15mm, toptitle=1mm]

\textbf{Practice Problem 30} (Method of moments for exponential data).
Suppose
\[X_1,\ldots,X_n\overset{\text{iid}}{\sim}\operatorname{Exponential}(\lambda),\]
with density \(f(x\mid \lambda)=\lambda e^{-\lambda x}\) for
\(x\geq 0\). Find the method of moments estimator of \(\lambda\).

\end{tcolorbox}

\begin{tcolorbox}[enhanced jigsaw, arc=.35mm, bottomrule=.15mm, bottomtitle=1mm, breakable, colback=white, colbacktitle=quarto-callout-tip-color!10!white, colframe=quarto-callout-tip-color-frame, coltitle=black, left=2mm, leftrule=.75mm, opacityback=0, opacitybacktitle=0.6, rightrule=.15mm, title=\textcolor{quarto-callout-tip-color}{\faLightbulb}\hspace{0.5em}{Solution}, titlerule=0mm, toprule=.15mm, toptitle=1mm]

For \(X\sim \operatorname{Exponential}(\lambda)\),
\[\mathbb{E}[X]=\frac{1}{\lambda}.\] The first sample moment is
\(m_1=\overline X\). Equating sample and population moments gives
\[\overline X=\frac{1}{\lambda}.\] Thus
\[\widetilde\lambda_{\mathrm{MOM}}=\frac{1}{\overline X}.\]

\end{tcolorbox}

\begin{tcolorbox}[enhanced jigsaw, arc=.35mm, bottomrule=.15mm, bottomtitle=1mm, breakable, colback=white, colbacktitle=quarto-callout-warning-color!10!white, colframe=quarto-callout-warning-color-frame, coltitle=black, left=2mm, leftrule=.75mm, opacityback=0, opacitybacktitle=0.6, rightrule=.15mm, title=\textcolor{quarto-callout-warning-color}{\faExclamationTriangle}\hspace{0.5em}{Practice Problem}, titlerule=0mm, toprule=.15mm, toptitle=1mm]

\textbf{Practice Problem 31} (MLE for exponential data). For the same
exponential model,
\[X_1,\ldots,X_n\overset{\text{iid}}{\sim}\operatorname{Exponential}(\lambda),\]
find the MLE of \(\lambda\).

\end{tcolorbox}

\begin{tcolorbox}[enhanced jigsaw, arc=.35mm, bottomrule=.15mm, bottomtitle=1mm, breakable, colback=white, colbacktitle=quarto-callout-tip-color!10!white, colframe=quarto-callout-tip-color-frame, coltitle=black, left=2mm, leftrule=.75mm, opacityback=0, opacitybacktitle=0.6, rightrule=.15mm, title=\textcolor{quarto-callout-tip-color}{\faLightbulb}\hspace{0.5em}{Solution}, titlerule=0mm, toprule=.15mm, toptitle=1mm]

The likelihood is
\[L(\lambda\mid \mathbf{x})=\prod_{i=1}^n \lambda e^{-\lambda x_i}
=\lambda^n e^{-\lambda\sum_{i=1}^n x_i}.\] The log-likelihood is
\[\ell(\lambda\mid \mathbf{x})=n\log\lambda-\lambda\sum_{i=1}^n x_i.\]
Differentiating gives
\[\frac{\partial \ell}{\partial \lambda}=\frac{n}{\lambda}-\sum_{i=1}^n x_i.\]
Setting this equal to zero gives
\[\widehat\lambda_{\mathrm{MLE}}=\frac{n}{\sum_{i=1}^n x_i}=\frac{1}{\overline x}.\]
In this model, the MOM estimator and MLE coincide.

\end{tcolorbox}

\begin{tcolorbox}[enhanced jigsaw, arc=.35mm, bottomrule=.15mm, bottomtitle=1mm, breakable, colback=white, colbacktitle=quarto-callout-warning-color!10!white, colframe=quarto-callout-warning-color-frame, coltitle=black, left=2mm, leftrule=.75mm, opacityback=0, opacitybacktitle=0.6, rightrule=.15mm, title=\textcolor{quarto-callout-warning-color}{\faExclamationTriangle}\hspace{0.5em}{Practice Problem}, titlerule=0mm, toprule=.15mm, toptitle=1mm]

\textbf{Practice Problem 32} (MLE invariance). Suppose
\(\widehat\lambda_{\mathrm{MLE}}=1/\overline X\) is the MLE for the
exponential rate \(\lambda\). Find the MLE for the mean
\(\theta=1/\lambda\).

\end{tcolorbox}

\begin{tcolorbox}[enhanced jigsaw, arc=.35mm, bottomrule=.15mm, bottomtitle=1mm, breakable, colback=white, colbacktitle=quarto-callout-tip-color!10!white, colframe=quarto-callout-tip-color-frame, coltitle=black, left=2mm, leftrule=.75mm, opacityback=0, opacitybacktitle=0.6, rightrule=.15mm, title=\textcolor{quarto-callout-tip-color}{\faLightbulb}\hspace{0.5em}{Solution}, titlerule=0mm, toprule=.15mm, toptitle=1mm]

By the invariance property,
\[\widehat\theta_{\mathrm{MLE}}=\frac{1}{\widehat\lambda_{\mathrm{MLE}}}
=\overline X.\]

\end{tcolorbox}

\begin{tcolorbox}[enhanced jigsaw, arc=.35mm, bottomrule=.15mm, bottomtitle=1mm, breakable, colback=white, colbacktitle=quarto-callout-warning-color!10!white, colframe=quarto-callout-warning-color-frame, coltitle=black, left=2mm, leftrule=.75mm, opacityback=0, opacitybacktitle=0.6, rightrule=.15mm, title=\textcolor{quarto-callout-warning-color}{\faExclamationTriangle}\hspace{0.5em}{Practice Problem}, titlerule=0mm, toprule=.15mm, toptitle=1mm]

\textbf{Practice Problem 33} (Bayesian beta-binomial update). Suppose
\(\theta\sim \operatorname{Beta}(2,2)\) and then \(n=10\) coin tosses
produce \(x=7\) heads. Find the posterior distribution, posterior mean,
and posterior mode.

\end{tcolorbox}

\begin{tcolorbox}[enhanced jigsaw, arc=.35mm, bottomrule=.15mm, bottomtitle=1mm, breakable, colback=white, colbacktitle=quarto-callout-tip-color!10!white, colframe=quarto-callout-tip-color-frame, coltitle=black, left=2mm, leftrule=.75mm, opacityback=0, opacitybacktitle=0.6, rightrule=.15mm, title=\textcolor{quarto-callout-tip-color}{\faLightbulb}\hspace{0.5em}{Solution}, titlerule=0mm, toprule=.15mm, toptitle=1mm]

The posterior is
\[\theta\mid \mathcal D\sim \operatorname{Beta}(2+7,2+10-7)=\operatorname{Beta}(9,5).\]
The posterior mean is
\[\mathbb{E}[\theta\mid \mathcal D]=\frac{9}{9+5}=\frac{9}{14}.\] The
posterior mode is \[\frac{9-1}{9+5-2}=\frac{8}{12}=\frac{2}{3}.\]

\end{tcolorbox}

\bookmarksetup{startatroot}

\chapter{Summary}\label{summary-6}

This section developed the first set of tools for finding point
estimators.

\begin{tcolorbox}[enhanced jigsaw, arc=.35mm, bottomrule=.15mm, bottomtitle=1mm, breakable, colback=white, colbacktitle=quarto-callout-tip-color!10!white, colframe=quarto-callout-tip-color-frame, coltitle=black, left=2mm, leftrule=.75mm, opacityback=0, opacitybacktitle=0.6, rightrule=.15mm, title=\textcolor{quarto-callout-tip-color}{\faLightbulb}\hspace{0.5em}{Key idea}, titlerule=0mm, toprule=.15mm, toptitle=1mm]

\begin{enumerate}
\def\labelenumi{\arabic{enumi}.}
\item
  A point estimator is a statistic used to estimate an unknown
  parameter.
\item
  Method of moments estimates parameters by matching sample moments to
  population moments.
\item
  Maximum likelihood estimation chooses the parameter value that makes
  the observed data most likely.
\item
  The log-likelihood is usually easier to optimize than the likelihood.
\item
  MLEs satisfy the invariance property: the MLE of \(g(\theta)\) is
  \(g(\widehat\theta)\).
\item
  Bayesian estimation combines prior information with the likelihood to
  form a posterior distribution.
\item
  MAP estimation maximizes the posterior; posterior mean and posterior
  mode are both common Bayesian point estimators.
\item
  Conjugate priors make Bayesian updating algebraically simple, as in
  the beta-binomial model.
\end{enumerate}

\end{tcolorbox}

\bookmarksetup{startatroot}

\chapter{Chapter 13: Point Estimation II --- Evaluating
Estimators}\label{chapter-13-point-estimation-ii-evaluating-estimators}

This chapter continues point estimation. Chapter 12 focused on how to
construct estimators. This chapter focuses on how to evaluate and
compare estimators using mean squared error, bias, variance, Fisher
information, Cramér--Rao lower bounds, Rao--Blackwellization, UMVUEs,
loss functions, risk functions, and Bayes risk.

\begin{tcolorbox}[enhanced jigsaw, arc=.35mm, bottomrule=.15mm, bottomtitle=1mm, breakable, colback=white, colbacktitle=quarto-callout-note-color!10!white, colframe=quarto-callout-note-color-frame, coltitle=black, left=2mm, leftrule=.75mm, opacityback=0, opacitybacktitle=0.6, rightrule=.15mm, title=\textcolor{quarto-callout-note-color}{\faInfo}\hspace{0.5em}{Topics}, titlerule=0mm, toprule=.15mm, toptitle=1mm]

Evaluating and comparing point estimators; mean squared error; bias and
variance; unbiasedness; best unbiased estimators; Cramer-Rao inequality;
Fisher information; Rao-Blackwell theorem; Lehmann-Scheffe theorem; loss
functions; risk functions; Bayes risk; examples for normal, Bernoulli,
binomial, and Poisson models.

\end{tcolorbox}

\bookmarksetup{startatroot}

\chapter{Overview: Evaluating and Comparing
Estimators}\label{overview-evaluating-and-comparing-estimators}

This section studies how to choose among different point estimators when
several construction methods are available.

In Section 12, we discussed ways to \emph{find} estimators: method of
moments, maximum likelihood estimation, and Bayesian methods. In this
section, the question changes from construction to evaluation.

\begin{tcolorbox}[enhanced jigsaw, arc=.35mm, bottomrule=.15mm, bottomtitle=1mm, breakable, colback=white, colbacktitle=quarto-callout-tip-color!10!white, colframe=quarto-callout-tip-color-frame, coltitle=black, left=2mm, leftrule=.75mm, opacityback=0, opacitybacktitle=0.6, rightrule=.15mm, title=\textcolor{quarto-callout-tip-color}{\faLightbulb}\hspace{0.5em}{Key idea}, titlerule=0mm, toprule=.15mm, toptitle=1mm]

Suppose several estimators are available for the same parameter. Which
estimator should we use?

\end{tcolorbox}

Common criteria for a good estimator include:

\begin{itemize}
\item
  small mean squared error;
\item
  unbiasedness;
\item
  small variance among unbiased estimators;
\item
  efficiency relative to a lower bound;
\item
  use of sufficient statistics;
\item
  consistency as the sample size grows;
\item
  robustness under small model deviations.
\end{itemize}

The main theme is that there is usually no single universal meaning of
``best.'' The answer depends on the loss function, the parameter space,
and whether we care about finite-sample or large-sample performance.

\bookmarksetup{startatroot}

\chapter{Mean Squared Error and Bias}\label{mean-squared-error-and-bias}

Mean squared error is one of the most useful criteria because it
combines variance and systematic bias into a single number.

\section{Mean squared error}\label{mean-squared-error}

We begin with the most common risk measure for point estimation.

\begin{tcolorbox}[enhanced jigsaw, arc=.35mm, bottomrule=.15mm, bottomtitle=1mm, breakable, colback=white, colbacktitle=quarto-callout-note-color!10!white, colframe=quarto-callout-note-color-frame, coltitle=black, left=2mm, leftrule=.75mm, opacityback=0, opacitybacktitle=0.6, rightrule=.15mm, title=\textcolor{quarto-callout-note-color}{\faInfo}\hspace{0.5em}{Definition}, titlerule=0mm, toprule=.15mm, toptitle=1mm]

\textbf{Definition 1} (Mean squared error). Let \(W=W(X_1,\ldots,X_n)\)
be an estimator of a scalar parameter \(\theta\). The \textbf{mean
squared error} of \(W\) is
\[\operatorname{MSE}_\theta(W)=\mathbb{E}_\theta\bigl[(W-\theta)^2\bigr].\]

\end{tcolorbox}

The MSE measures the average squared distance between the random
estimator and the unknown true parameter. It depends on \(\theta\), so
it is a function on the parameter space.

\section{Bias and unbiasedness}\label{bias-and-unbiasedness}

Bias measures whether the estimator systematically overestimates or
underestimates the target.

\begin{tcolorbox}[enhanced jigsaw, arc=.35mm, bottomrule=.15mm, bottomtitle=1mm, breakable, colback=white, colbacktitle=quarto-callout-note-color!10!white, colframe=quarto-callout-note-color-frame, coltitle=black, left=2mm, leftrule=.75mm, opacityback=0, opacitybacktitle=0.6, rightrule=.15mm, title=\textcolor{quarto-callout-note-color}{\faInfo}\hspace{0.5em}{Definition}, titlerule=0mm, toprule=.15mm, toptitle=1mm]

\textbf{Definition 2} (Bias). The \textbf{bias} of a point estimator
\(W\) of \(\theta\) is
\[\operatorname{Bias}_\theta(W)=\mathbb{E}_\theta[W]-\theta.\] The
estimator \(W\) is \textbf{unbiased} for \(\theta\) if
\[\mathbb{E}_\theta[W]=\theta\] for all possible values of \(\theta\).

\end{tcolorbox}

\begin{tcolorbox}[enhanced jigsaw, arc=.35mm, bottomrule=.15mm, bottomtitle=1mm, breakable, colback=white, colbacktitle=quarto-callout-important-color!10!white, colframe=quarto-callout-important-color-frame, coltitle=black, left=2mm, leftrule=.75mm, opacityback=0, opacitybacktitle=0.6, rightrule=.15mm, title=\textcolor{quarto-callout-important-color}{\faExclamation}\hspace{0.5em}{Remark}, titlerule=0mm, toprule=.15mm, toptitle=1mm]

\emph{Remark 3}. Unbiasedness is a natural criterion, but it is not the
only criterion. A biased estimator can have smaller MSE than an unbiased
estimator if the reduction in variance is large enough.

\end{tcolorbox}

\section{Bias-variance decomposition}\label{bias-variance-decomposition}

The key reason MSE is convenient is that it decomposes into variance
plus squared bias.

\begin{tcolorbox}[enhanced jigsaw, arc=.35mm, bottomrule=.15mm, bottomtitle=1mm, breakable, colback=white, colbacktitle=quarto-callout-important-color!10!white, colframe=quarto-callout-important-color-frame, coltitle=black, left=2mm, leftrule=.75mm, opacityback=0, opacitybacktitle=0.6, rightrule=.15mm, title=\textcolor{quarto-callout-important-color}{\faExclamation}\hspace{0.5em}{Theorem}, titlerule=0mm, toprule=.15mm, toptitle=1mm]

\textbf{Theorem 4} (Bias-variance decomposition). \emph{For any
estimator \(W\) of \(\theta\), \[\operatorname{MSE}_\theta(W)
    =\mathbb{E}_\theta[(W-\theta)^2]
    =\operatorname{Var}_\theta(W)+\bigl(\mathbb{E}_\theta[W]-\theta\bigr)^2.\]
Equivalently,
\[\operatorname{MSE}_\theta(W)=\operatorname{Var}_\theta(W)+\operatorname{Bias}_\theta(W)^2.\]}

\end{tcolorbox}

\begin{tcolorbox}[enhanced jigsaw, arc=.35mm, bottomrule=.15mm, bottomtitle=1mm, breakable, colback=white, colbacktitle=quarto-callout-note-color!10!white, colframe=quarto-callout-note-color-frame, coltitle=black, left=2mm, leftrule=.75mm, opacityback=0, opacitybacktitle=0.6, rightrule=.15mm, title=\textcolor{quarto-callout-note-color}{\faInfo}\hspace{0.5em}{Proof}, titlerule=0mm, toprule=.15mm, toptitle=1mm]

\emph{Proof.} Write
\[W-\theta=(W-\mathbb{E}_\theta W)+(\mathbb{E}_\theta W-\theta).\] Then
\[\begin{aligned}
\mathbb{E}_\theta[(W-\theta)^2]
&=\mathbb{E}_\theta\left[(W-\mathbb{E}_\theta W)^2\right]
  +2(\mathbb{E}_\theta W-\theta)\mathbb{E}_\theta[W-\mathbb{E}_\theta W]
  +(\mathbb{E}_\theta W-\theta)^2\\
&=\operatorname{Var}_\theta(W)+(\mathbb{E}_\theta W-\theta)^2,
\end{aligned}\] because the middle term is zero.~◻

\end{tcolorbox}

\begin{tcolorbox}[enhanced jigsaw, arc=.35mm, bottomrule=.15mm, bottomtitle=1mm, breakable, colback=white, colbacktitle=quarto-callout-important-color!10!white, colframe=quarto-callout-important-color-frame, coltitle=black, left=2mm, leftrule=.75mm, opacityback=0, opacitybacktitle=0.6, rightrule=.15mm, title=\textcolor{quarto-callout-important-color}{\faExclamation}\hspace{0.5em}{Remark}, titlerule=0mm, toprule=.15mm, toptitle=1mm]

\emph{Remark 5}. It is reasonable to consider other errors, such as
\(\mathbb{E}_\theta|W-\theta|\). The MSE is especially common because it
is differentiable and has the bias-variance decomposition.

\end{tcolorbox}

\bookmarksetup{startatroot}

\chapter{Examples of MSE}\label{examples-of-mse}

These examples show how MSE is computed and how bias can sometimes
improve MSE.

\section{Sample mean and sample variance under a normal
model}\label{sample-mean-and-sample-variance-under-a-normal-model}

We first recall basic facts about the sample mean and sample variance.

\begin{tcolorbox}[enhanced jigsaw, arc=.35mm, bottomrule=.15mm, bottomtitle=1mm, breakable, colback=white, colbacktitle=quarto-callout-note-color!10!white, colframe=quarto-callout-note-color-frame, coltitle=black, left=2mm, leftrule=.75mm, opacityback=0, opacitybacktitle=0.6, rightrule=.15mm, title=\textcolor{quarto-callout-note-color}{\faInfo}\hspace{0.5em}{Example}, titlerule=0mm, toprule=.15mm, toptitle=1mm]

\textbf{Example 6} (Normal distribution). Let \(X_1,\ldots,X_n\) be a
random sample from a population with mean \(\mu\) and variance
\(\sigma^2\). Define \[\bar X=\frac1n\sum_{i=1}^n X_i,
    \qquad
    S^2=\frac1{n-1}\sum_{i=1}^n (X_i-\bar X)^2.\] Then
\[\mathbb{E}(\bar X)=\mu,\qquad
    \operatorname{Var}(\bar X)=\frac{\sigma^2}{n},\qquad
    \mathbb{E}(S^2)=\sigma^2.\] Thus \(\bar X\) is unbiased for \(\mu\),
and \(S^2\) is unbiased for \(\sigma^2\).

\end{tcolorbox}

\begin{tcolorbox}[enhanced jigsaw, arc=.35mm, bottomrule=.15mm, bottomtitle=1mm, breakable, colback=white, colbacktitle=quarto-callout-tip-color!10!white, colframe=quarto-callout-tip-color-frame, coltitle=black, left=2mm, leftrule=.75mm, opacityback=0, opacitybacktitle=0.6, rightrule=.15mm, title=\textcolor{quarto-callout-tip-color}{\faLightbulb}\hspace{0.5em}{Solution}, titlerule=0mm, toprule=.15mm, toptitle=1mm]

Since \(X_1,\ldots,X_n\) are iid,
\[\mathbb{E}(\bar X)=\frac1n\sum_{i=1}^n \mathbb{E}(X_i)=\mu\] and
\[\operatorname{Var}(\bar X)=\frac1{n^2}\sum_{i=1}^n \operatorname{Var}(X_i)=\frac{\sigma^2}{n}.\]
The identity \(\mathbb{E}(S^2)=\sigma^2\) was proved in the sampling
section. Therefore
\[\operatorname{MSE}(\bar X)=\operatorname{Var}(\bar X)=\frac{\sigma^2}{n}\]
when estimating \(\mu\). If
\(X_i\sim \operatorname{Normal}(\mu,\sigma^2)\), then
\[\frac{(n-1)S^2}{\sigma^2}\sim \chi^2_{n-1},\] so
\[\operatorname{Var}(S^2)=\frac{2\sigma^4}{n-1}.\] Since \(S^2\) is
unbiased for \(\sigma^2\),
\[\operatorname{MSE}(S^2)=\frac{2\sigma^4}{n-1}.\]

\end{tcolorbox}

\section{The MLE of normal variance}\label{the-mle-of-normal-variance}

The MLE for the normal variance is biased, but it can have smaller MSE
than the unbiased sample variance.

\begin{tcolorbox}[enhanced jigsaw, arc=.35mm, bottomrule=.15mm, bottomtitle=1mm, breakable, colback=white, colbacktitle=quarto-callout-note-color!10!white, colframe=quarto-callout-note-color-frame, coltitle=black, left=2mm, leftrule=.75mm, opacityback=0, opacitybacktitle=0.6, rightrule=.15mm, title=\textcolor{quarto-callout-note-color}{\faInfo}\hspace{0.5em}{Example}, titlerule=0mm, toprule=.15mm, toptitle=1mm]

\textbf{Example 7} (Normal variance MLE). Suppose
\(X_1,\ldots,X_n\overset{\text{iid}}{\sim}\operatorname{Normal}(\mu,\sigma^2)\).
The MLE of \(\sigma^2\) is \[\widehat\sigma^2_{\mathrm{MLE}}
    =\frac1n\sum_{i=1}^n (X_i-\bar X)^2
    =\frac{n-1}{n}S^2.\] Find its bias, variance, and MSE.

\end{tcolorbox}

\begin{tcolorbox}[enhanced jigsaw, arc=.35mm, bottomrule=.15mm, bottomtitle=1mm, breakable, colback=white, colbacktitle=quarto-callout-tip-color!10!white, colframe=quarto-callout-tip-color-frame, coltitle=black, left=2mm, leftrule=.75mm, opacityback=0, opacitybacktitle=0.6, rightrule=.15mm, title=\textcolor{quarto-callout-tip-color}{\faLightbulb}\hspace{0.5em}{Solution}, titlerule=0mm, toprule=.15mm, toptitle=1mm]

Since \(\mathbb{E}(S^2)=\sigma^2\),
\[\mathbb{E}\bigl(\widehat\sigma^2_{\mathrm{MLE}}\bigr)
    =\frac{n-1}{n}\sigma^2.\] Thus
\[\operatorname{Bias}\bigl(\widehat\sigma^2_{\mathrm{MLE}}\bigr)
    =-\frac{\sigma^2}{n}.\] Also,
\[\operatorname{Var}\bigl(\widehat\sigma^2_{\mathrm{MLE}}\bigr)
    =\left(\frac{n-1}{n}\right)^2\operatorname{Var}(S^2)
    =\left(\frac{n-1}{n}\right)^2\frac{2\sigma^4}{n-1}
    =\frac{2(n-1)\sigma^4}{n^2}.\] Therefore
\[\operatorname{MSE}\bigl(\widehat\sigma^2_{\mathrm{MLE}}\bigr)
    =\frac{2(n-1)\sigma^4}{n^2}+\frac{\sigma^4}{n^2}
    =\frac{(2n-1)\sigma^4}{n^2}.\] Compare this with
\[\operatorname{MSE}(S^2)=\frac{2\sigma^4}{n-1}.\] For \(n\ge 2\),
\[\frac{(2n-1)}{n^2}<\frac{2}{n-1},\] so
\[\operatorname{MSE}\bigl(\widehat\sigma^2_{\mathrm{MLE}}\bigr)<\operatorname{MSE}(S^2).\]
The MLE trades a small bias for a larger variance reduction.

\end{tcolorbox}

\begin{tcolorbox}[enhanced jigsaw, arc=.35mm, bottomrule=.15mm, bottomtitle=1mm, breakable, colback=white, colbacktitle=quarto-callout-tip-color!10!white, colframe=quarto-callout-tip-color-frame, coltitle=black, left=2mm, leftrule=.75mm, opacityback=0, opacitybacktitle=0.6, rightrule=.15mm, title=\textcolor{quarto-callout-tip-color}{\faLightbulb}\hspace{0.5em}{Key idea}, titlerule=0mm, toprule=.15mm, toptitle=1mm]

Bias-variance tradeoff An unbiased estimator is not automatically best
under MSE. Sometimes a biased estimator has smaller MSE because its
variance is much smaller.

\end{tcolorbox}

\section{A Bayes estimator for Bernoulli
data}\label{a-bayes-estimator-for-bernoulli-data}

This example compares the usual sample proportion with a Bayesian
shrinkage estimator.

\begin{tcolorbox}[enhanced jigsaw, arc=.35mm, bottomrule=.15mm, bottomtitle=1mm, breakable, colback=white, colbacktitle=quarto-callout-note-color!10!white, colframe=quarto-callout-note-color-frame, coltitle=black, left=2mm, leftrule=.75mm, opacityback=0, opacitybacktitle=0.6, rightrule=.15mm, title=\textcolor{quarto-callout-note-color}{\faInfo}\hspace{0.5em}{Example}, titlerule=0mm, toprule=.15mm, toptitle=1mm]

\textbf{Example 8} (Binomial Bayes estimator). Suppose
\(X_1,\ldots,X_n\overset{\text{iid}}{\sim}\operatorname{Bernoulli}(p)\),
where \(p\) is unknown. Let \[S=X_1+\cdots+X_n.\] The MLE is
\[\widehat p_{\mathrm{MLE}}=\bar X=\frac{S}{n}.\] If the prior is
\(p\sim \operatorname{Beta}(\alpha,\beta)\), then the posterior mean
estimator is \[\widehat p_{\mathrm{B}}
    =\mathbb{E}[p\mid S]
    =\frac{\alpha+S}{\alpha+\beta+n}.\] Compute the MSE of \(\bar X\)
and \(\widehat p_{\mathrm{B}}\) as functions of \(p\).

\end{tcolorbox}

\begin{tcolorbox}[enhanced jigsaw, arc=.35mm, bottomrule=.15mm, bottomtitle=1mm, breakable, colback=white, colbacktitle=quarto-callout-tip-color!10!white, colframe=quarto-callout-tip-color-frame, coltitle=black, left=2mm, leftrule=.75mm, opacityback=0, opacitybacktitle=0.6, rightrule=.15mm, title=\textcolor{quarto-callout-tip-color}{\faLightbulb}\hspace{0.5em}{Solution}, titlerule=0mm, toprule=.15mm, toptitle=1mm]

For the MLE, \[\mathbb{E}_p(\bar X)=p,
    \qquad
    \operatorname{Var}_p(\bar X)=\frac{p(1-p)}{n}.\] Thus
\[\operatorname{MSE}_p(\bar X)=\frac{p(1-p)}{n}.\] For the Bayes
estimator, \[\widehat p_{\mathrm{B}}=\frac{\alpha+S}{\alpha+\beta+n}.\]
Since \(S\sim \operatorname{Binomial}(n,p)\), \[\mathbb{E}_p(S)=np,
    \qquad
    \operatorname{Var}_p(S)=np(1-p).\] Therefore
\[\mathbb{E}_p(\widehat p_{\mathrm{B}})
    =\frac{\alpha+np}{\alpha+\beta+n},\]
\[\operatorname{Var}_p(\widehat p_{\mathrm{B}})
    =\frac{np(1-p)}{(\alpha+\beta+n)^2},\] and
\[\operatorname{Bias}_p(\widehat p_{\mathrm{B}})
    =\frac{\alpha+np}{\alpha+\beta+n}-p
    =\frac{\alpha-p(\alpha+\beta)}{\alpha+\beta+n}.\] Hence
\[\operatorname{MSE}_p(\widehat p_{\mathrm{B}})
    =\frac{np(1-p)}{(\alpha+\beta+n)^2}
    +\left(\frac{\alpha-p(\alpha+\beta)}{\alpha+\beta+n}\right)^2.\] If
we choose \[\alpha=\beta=\frac{\sqrt n}{4},\] then this estimator
shrinks \(\bar X\) toward \(1/2\). It can have better MSE for small
\(n\) or when one strongly believes \(p\) is close to \(1/2\). For large
\(n\), the MLE becomes very strong because the data dominates.

\end{tcolorbox}

\begin{tcolorbox}[enhanced jigsaw, arc=.35mm, bottomrule=.15mm, bottomtitle=1mm, breakable, colback=white, colbacktitle=quarto-callout-important-color!10!white, colframe=quarto-callout-important-color-frame, coltitle=black, left=2mm, leftrule=.75mm, opacityback=0, opacitybacktitle=0.6, rightrule=.15mm, title=\textcolor{quarto-callout-important-color}{\faExclamation}\hspace{0.5em}{Remark}, titlerule=0mm, toprule=.15mm, toptitle=1mm]

\emph{Remark 9}. The lecture slide emphasizes a special choice of
\(\alpha\) and \(\beta\) that makes the MSE curve flatter. The main
statistical idea is shrinkage: the estimator sacrifices unbiasedness to
reduce risk near values favored by the prior.

\end{tcolorbox}

\bookmarksetup{startatroot}

\chapter{Best Unbiased Estimators}\label{best-unbiased-estimators}

Unbiasedness alone is not enough, so we compare unbiased estimators by
their variance.

\section{Definition and uniqueness}\label{definition-and-uniqueness}

The best unbiased estimator has the smallest variance among all unbiased
estimators.

\begin{tcolorbox}[enhanced jigsaw, arc=.35mm, bottomrule=.15mm, bottomtitle=1mm, breakable, colback=white, colbacktitle=quarto-callout-note-color!10!white, colframe=quarto-callout-note-color-frame, coltitle=black, left=2mm, leftrule=.75mm, opacityback=0, opacitybacktitle=0.6, rightrule=.15mm, title=\textcolor{quarto-callout-note-color}{\faInfo}\hspace{0.5em}{Definition}, titlerule=0mm, toprule=.15mm, toptitle=1mm]

\textbf{Definition 10} (Best unbiased estimator and UMVUE). An estimator
\(\widehat\theta\) of \(\theta\) is called a \textbf{best unbiased
estimator} if:

\begin{enumerate}
\def\labelenumi{\arabic{enumi}.}
\item
  it is unbiased: \(\mathbb{E}_\theta[\widehat\theta]=\theta\) for all
  \(\theta\);
\item
  it has minimum variance among all unbiased estimators of \(\theta\)
  for every \(\theta\).
\end{enumerate}

It is also called a \textbf{uniform minimum variance unbiased
estimator}, or \textbf{UMVUE}.

\end{tcolorbox}

The definition extends naturally to estimating a function \(g(\theta)\).

\begin{tcolorbox}[enhanced jigsaw, arc=.35mm, bottomrule=.15mm, bottomtitle=1mm, breakable, colback=white, colbacktitle=quarto-callout-important-color!10!white, colframe=quarto-callout-important-color-frame, coltitle=black, left=2mm, leftrule=.75mm, opacityback=0, opacitybacktitle=0.6, rightrule=.15mm, title=\textcolor{quarto-callout-important-color}{\faExclamation}\hspace{0.5em}{Theorem}, titlerule=0mm, toprule=.15mm, toptitle=1mm]

\textbf{Theorem 11} (Uniqueness). \emph{If \(W\) is a best unbiased
estimator of \(g(\theta)\), then \(W\) is unique almost surely.}

\end{tcolorbox}

\begin{tcolorbox}[enhanced jigsaw, arc=.35mm, bottomrule=.15mm, bottomtitle=1mm, breakable, colback=white, colbacktitle=quarto-callout-note-color!10!white, colframe=quarto-callout-note-color-frame, coltitle=black, left=2mm, leftrule=.75mm, opacityback=0, opacitybacktitle=0.6, rightrule=.15mm, title=\textcolor{quarto-callout-note-color}{\faInfo}\hspace{0.5em}{Proof}, titlerule=0mm, toprule=.15mm, toptitle=1mm]

\emph{Proof.} Suppose \(W_1\) and \(W_2\) are both best unbiased
estimators. Then \[W_3=\frac{W_1+W_2}{2}\] is also unbiased. Since
\(W_1\) and \(W_2\) are both minimum-variance unbiased estimators,
\(W_3\) cannot have smaller variance. But
\[\operatorname{Var}(W_1-W_2)\ge 0\] and the usual variance identity
implies that averaging two different unbiased estimators would strictly
reduce variance unless \(W_1=W_2\) almost surely. Therefore \(W_1=W_2\)
almost surely.~◻

\end{tcolorbox}

\section{Poisson unbiased estimation}\label{poisson-unbiased-estimation}

This example shows that many unbiased estimators can exist, but one may
be clearly better.

\begin{tcolorbox}[enhanced jigsaw, arc=.35mm, bottomrule=.15mm, bottomtitle=1mm, breakable, colback=white, colbacktitle=quarto-callout-note-color!10!white, colframe=quarto-callout-note-color-frame, coltitle=black, left=2mm, leftrule=.75mm, opacityback=0, opacitybacktitle=0.6, rightrule=.15mm, title=\textcolor{quarto-callout-note-color}{\faInfo}\hspace{0.5em}{Example}, titlerule=0mm, toprule=.15mm, toptitle=1mm]

\textbf{Example 12} (Poisson unbiased estimators). Suppose
\(X_1,\ldots,X_n\overset{\text{iid}}{\sim}\operatorname{Poisson}(\lambda)\).
Then both \(\bar X\) and \(S^2\) are unbiased estimators of \(\lambda\).
Compare their MSEs.

\end{tcolorbox}

\begin{tcolorbox}[enhanced jigsaw, arc=.35mm, bottomrule=.15mm, bottomtitle=1mm, breakable, colback=white, colbacktitle=quarto-callout-tip-color!10!white, colframe=quarto-callout-tip-color-frame, coltitle=black, left=2mm, leftrule=.75mm, opacityback=0, opacitybacktitle=0.6, rightrule=.15mm, title=\textcolor{quarto-callout-tip-color}{\faLightbulb}\hspace{0.5em}{Solution}, titlerule=0mm, toprule=.15mm, toptitle=1mm]

For the sample mean, \[\mathbb{E}_\lambda(\bar X)=\lambda,
    \qquad
    \operatorname{Var}_\lambda(\bar X)=\frac{\lambda}{n}.\] Thus
\[\operatorname{MSE}(\bar X)=\frac{\lambda}{n}.\] The sample variance
\(S^2\) is also unbiased for the population variance, and for a Poisson
random variable the population variance equals \(\lambda\). Hence
\(\mathbb{E}(S^2)=\lambda\).

A general formula for the variance of the sample variance is
\[\operatorname{Var}(S^2)=\frac1n\left(\mu_4-\frac{n-3}{n-1}\sigma^4\right),\]
where \(\mu_4=\mathbb{E}[(X-\mu)^4]\) and
\(\sigma^2=\operatorname{Var}(X)\). For
\(X\sim \operatorname{Poisson}(\lambda)\), \[\mu=\lambda,
    \qquad
    \sigma^2=\lambda,
    \qquad
    \mu_4=\lambda+3\lambda^2.\] Therefore \[\begin{aligned}
    \operatorname{Var}(S^2)
    &=\frac1n\left(\lambda+3\lambda^2-\frac{n-3}{n-1}\lambda^2\right)\\
    &=\frac{\lambda}{n}+\frac{2\lambda^2}{n-1}.
\end{aligned}\] Since \(S^2\) is unbiased,
\[\operatorname{MSE}(S^2)=\frac{\lambda}{n}+\frac{2\lambda^2}{n-1}.\]
Thus \[\operatorname{MSE}(\bar X)<\operatorname{MSE}(S^2)\] for
\(\lambda>0\). This shows \(\bar X\) is better than \(S^2\) among these
two unbiased estimators, but by itself it does not prove \(\bar X\) is
the best among \emph{all} unbiased estimators.

\end{tcolorbox}

\bookmarksetup{startatroot}

\chapter{Cramer-Rao Inequality and Fisher
Information}\label{cramer-rao-inequality-and-fisher-information}

The Cramer-Rao inequality gives a lower bound on the variance of
unbiased estimators.

\section{The inequality}\label{the-inequality}

The bound is one of the main tools for proving that an unbiased
estimator is best.

\begin{tcolorbox}[enhanced jigsaw, arc=.35mm, bottomrule=.15mm, bottomtitle=1mm, breakable, colback=white, colbacktitle=quarto-callout-important-color!10!white, colframe=quarto-callout-important-color-frame, coltitle=black, left=2mm, leftrule=.75mm, opacityback=0, opacitybacktitle=0.6, rightrule=.15mm, title=\textcolor{quarto-callout-important-color}{\faExclamation}\hspace{0.5em}{Theorem}, titlerule=0mm, toprule=.15mm, toptitle=1mm]

\textbf{Theorem 13} (Cramer-Rao inequality). \emph{Let
\(X=(X_1,\ldots,X_n)\) have joint density or mass function
\(f(x\mid \theta)\). Let \(W(X)\) be an estimator of \(g(\theta)\)
satisfying the usual regularity conditions that allow differentiation
under the integral sign. Then \[\operatorname{Var}_\theta(W)
    \ge
    \frac{\left(\dfrac{d}{d\theta}\mathbb{E}_\theta[W]\right)^2}
    {\mathbb{E}_\theta\left[\left(\dfrac{\partial}{\partial\theta}\log f(X\mid\theta)\right)^2\right]}.\]
If \(W\) is unbiased for \(g(\theta)\), then
\(\mathbb{E}_\theta[W]=g(\theta)\) and \[\operatorname{Var}_\theta(W)
    \ge
    \frac{\left(g'(\theta)\right)^2}{I_X(\theta)},\] where
\[I_X(\theta)=\mathbb{E}_\theta\left[\left(\frac{\partial}{\partial\theta}\log f(X\mid\theta)\right)^2\right]\]
is the Fisher information in the sample.}

\end{tcolorbox}

\begin{tcolorbox}[enhanced jigsaw, arc=.35mm, bottomrule=.15mm, bottomtitle=1mm, breakable, colback=white, colbacktitle=quarto-callout-note-color!10!white, colframe=quarto-callout-note-color-frame, coltitle=black, left=2mm, leftrule=.75mm, opacityback=0, opacitybacktitle=0.6, rightrule=.15mm, title=\textcolor{quarto-callout-note-color}{\faInfo}\hspace{0.5em}{Proof}, titlerule=0mm, toprule=.15mm, toptitle=1mm]

\emph{Proof idea.} Let
\[U_\theta(X)=\frac{\partial}{\partial\theta}\log f(X\mid \theta)\] be
the score function. Under regularity conditions,
\[\mathbb{E}_\theta[U_\theta(X)]=0.\] Also, \[\begin{aligned}
\frac{d}{d\theta}\mathbb{E}_\theta[W]
&=\frac{d}{d\theta}\int W(x)f(x\mid \theta)\,dx\\
&=\int W(x)\frac{\partial}{\partial\theta}f(x\mid\theta)\,dx\\
&=\int W(x)\left(\frac{\partial}{\partial\theta}\log f(x\mid\theta)\right)f(x\mid\theta)\,dx\\
&=\mathbb{E}_\theta[W U_\theta(X)]\\
&=\operatorname{Cov}_\theta(W,U_\theta(X)).
\end{aligned}\] By Cauchy-Schwarz,
\[\left(\frac{d}{d\theta}\mathbb{E}_\theta[W]\right)^2
    \le \operatorname{Var}_\theta(W)\operatorname{Var}_\theta(U_\theta(X)).\]
Since
\(\operatorname{Var}_\theta(U_\theta)=\mathbb{E}_\theta[U_\theta^2]=I_X(\theta)\),
rearranging gives the result.~◻

\end{tcolorbox}

\section{Equality condition}\label{equality-condition}

The equality condition explains when the lower bound is attained.

\begin{tcolorbox}[enhanced jigsaw, arc=.35mm, bottomrule=.15mm, bottomtitle=1mm, breakable, colback=white, colbacktitle=quarto-callout-important-color!10!white, colframe=quarto-callout-important-color-frame, coltitle=black, left=2mm, leftrule=.75mm, opacityback=0, opacitybacktitle=0.6, rightrule=.15mm, title=\textcolor{quarto-callout-important-color}{\faExclamation}\hspace{0.5em}{Theorem}, titlerule=0mm, toprule=.15mm, toptitle=1mm]

\textbf{Theorem 14} (Equality condition). \emph{Equality in the
Cramer-Rao inequality holds if and only if there exists a function
\(h(\theta)\) such that \[h(\theta)\{W(x)-g(\theta)\}
    =\frac{\partial}{\partial\theta}\log L(\theta\mid x),\] where
\(L(\theta\mid x)\) is the likelihood function.}

\end{tcolorbox}

This condition is useful because if an unbiased estimator attains the
Cramer-Rao lower bound, then it is a best unbiased estimator.

\section{IID form and information
number}\label{iid-form-and-information-number}

For independent samples, the Fisher information adds.

\begin{tcolorbox}[enhanced jigsaw, arc=.35mm, bottomrule=.15mm, bottomtitle=1mm, breakable, colback=white, colbacktitle=quarto-callout-important-color!10!white, colframe=quarto-callout-important-color-frame, coltitle=black, left=2mm, leftrule=.75mm, opacityback=0, opacitybacktitle=0.6, rightrule=.15mm, title=\textcolor{quarto-callout-important-color}{\faExclamation}\hspace{0.5em}{Corollary}, titlerule=0mm, toprule=.15mm, toptitle=1mm]

\textbf{Corollary 15} (IID sample). \emph{Suppose \(X_1,\ldots,X_n\) are
iid with density or mass function \(f(x\mid\theta)\). Then \[I_X(\theta)
    =n I_1(\theta),
    \qquad
    I_1(\theta)=\mathbb{E}_\theta\left[\left(\frac{\partial}{\partial\theta}\log f(X\mid\theta)\right)^2\right].\]
Hence \[\operatorname{Var}_\theta(W)
    \ge
    \frac{\left(\dfrac{d}{d\theta}\mathbb{E}_\theta[W]\right)^2}
    {nI_1(\theta)}.\]}

\end{tcolorbox}

\begin{tcolorbox}[enhanced jigsaw, arc=.35mm, bottomrule=.15mm, bottomtitle=1mm, breakable, colback=white, colbacktitle=quarto-callout-important-color!10!white, colframe=quarto-callout-important-color-frame, coltitle=black, left=2mm, leftrule=.75mm, opacityback=0, opacitybacktitle=0.6, rightrule=.15mm, title=\textcolor{quarto-callout-important-color}{\faExclamation}\hspace{0.5em}{Lemma}, titlerule=0mm, toprule=.15mm, toptitle=1mm]

\textbf{Lemma 16} (Computing Fisher information). \emph{Under suitable
regularity conditions, often satisfied for exponential-family models,
\[I_1(\theta)
    =\mathbb{E}_\theta\left[\left(\frac{\partial}{\partial\theta}\log f(X\mid\theta)\right)^2\right]
    =-\mathbb{E}_\theta\left[\frac{\partial^2}{\partial\theta^2}\log f(X\mid\theta)\right].\]}

\end{tcolorbox}

\section{Poisson conclusion}\label{poisson-conclusion}

We now return to the Poisson example and prove that \(\bar X\) is best
unbiased.

\begin{tcolorbox}[enhanced jigsaw, arc=.35mm, bottomrule=.15mm, bottomtitle=1mm, breakable, colback=white, colbacktitle=quarto-callout-note-color!10!white, colframe=quarto-callout-note-color-frame, coltitle=black, left=2mm, leftrule=.75mm, opacityback=0, opacitybacktitle=0.6, rightrule=.15mm, title=\textcolor{quarto-callout-note-color}{\faInfo}\hspace{0.5em}{Example}, titlerule=0mm, toprule=.15mm, toptitle=1mm]

\textbf{Example 17} (Poisson UMVUE). Let
\(X_1,\ldots,X_n\overset{\text{iid}}{\sim}\operatorname{Poisson}(\lambda)\).
Show that \(\bar X\) is the best unbiased estimator of \(\lambda\).

\end{tcolorbox}

\begin{tcolorbox}[enhanced jigsaw, arc=.35mm, bottomrule=.15mm, bottomtitle=1mm, breakable, colback=white, colbacktitle=quarto-callout-tip-color!10!white, colframe=quarto-callout-tip-color-frame, coltitle=black, left=2mm, leftrule=.75mm, opacityback=0, opacitybacktitle=0.6, rightrule=.15mm, title=\textcolor{quarto-callout-tip-color}{\faLightbulb}\hspace{0.5em}{Solution}, titlerule=0mm, toprule=.15mm, toptitle=1mm]

For one observation,
\[f(x\mid\lambda)=e^{-\lambda}\frac{\lambda^x}{x!},\] so
\[\log f(x\mid\lambda)=-\lambda+x\log\lambda-\log(x!).\] Then
\[\frac{\partial^2}{\partial\lambda^2}\log f(x\mid\lambda)
    =-\frac{x}{\lambda^2}.\] Therefore \[I_1(\lambda)
    =-\mathbb{E}_\lambda\left[-\frac{X}{\lambda^2}\right]
    =\frac{\lambda}{\lambda^2}
    =\frac1\lambda.\] For \(n\) iid observations,
\[I_X(\lambda)=\frac{n}{\lambda}.\] If \(W\) is unbiased for
\(\lambda\), then \(g(\lambda)=\lambda\) and \(g'(\lambda)=1\). The
Cramer-Rao lower bound gives
\[\operatorname{Var}_\lambda(W)\ge \frac{1}{n/\lambda}=\frac{\lambda}{n}.\]
But \[\mathbb{E}_\lambda(\bar X)=\lambda,
    \qquad
    \operatorname{Var}_\lambda(\bar X)=\frac{\lambda}{n}.\] Thus
\(\bar X\) is unbiased and attains the lower bound. Hence \(\bar X\) is
the best unbiased estimator of \(\lambda\).

\end{tcolorbox}

\section{Normal variance example}\label{normal-variance-example}

The Cramer-Rao bound is a lower bound, but not every unbiased estimator
attains it.

\begin{tcolorbox}[enhanced jigsaw, arc=.35mm, bottomrule=.15mm, bottomtitle=1mm, breakable, colback=white, colbacktitle=quarto-callout-note-color!10!white, colframe=quarto-callout-note-color-frame, coltitle=black, left=2mm, leftrule=.75mm, opacityback=0, opacitybacktitle=0.6, rightrule=.15mm, title=\textcolor{quarto-callout-note-color}{\faInfo}\hspace{0.5em}{Example}, titlerule=0mm, toprule=.15mm, toptitle=1mm]

\textbf{Example 18} (Normal variance). Let
\(X_1,\ldots,X_n\overset{\text{iid}}{\sim}\operatorname{Normal}(\mu,\sigma^2)\).
Consider estimating \(\sigma^2\). If we use the parameter
\(\theta=\sigma^2\), then the normal log-density for one observation is
\[\log f(x\mid \mu,\theta)
    =-\frac12\log(2\pi\theta)-\frac{(x-\mu)^2}{2\theta}.\] When \(\mu\)
is treated as known in this direct one-parameter calculation,
\[-\mathbb{E}\left[\frac{\partial^2}{\partial\theta^2}\log f(X\mid\mu,\theta)\right]
    =\frac{1}{2\theta^2}
    =\frac{1}{2\sigma^4}.\] Thus the Cramer-Rao lower bound for unbiased
estimators of \(\sigma^2\) is
\[\operatorname{Var}(W)\ge \frac{2\sigma^4}{n}.\]

\end{tcolorbox}

\begin{tcolorbox}[enhanced jigsaw, arc=.35mm, bottomrule=.15mm, bottomtitle=1mm, breakable, colback=white, colbacktitle=quarto-callout-tip-color!10!white, colframe=quarto-callout-tip-color-frame, coltitle=black, left=2mm, leftrule=.75mm, opacityback=0, opacitybacktitle=0.6, rightrule=.15mm, title=\textcolor{quarto-callout-tip-color}{\faLightbulb}\hspace{0.5em}{Solution}, titlerule=0mm, toprule=.15mm, toptitle=1mm]

Compute derivatives:
\[\frac{\partial}{\partial\theta}\log f(x\mid\mu,\theta)
    =-\frac1{2\theta}+\frac{(x-\mu)^2}{2\theta^2},\]
\[\frac{\partial^2}{\partial\theta^2}\log f(x\mid\mu,\theta)
    =\frac1{2\theta^2}-\frac{(x-\mu)^2}{\theta^3}.\] Since
\(\mathbb{E}[(X-\mu)^2]=\theta\),
\[-\mathbb{E}\left[\frac{\partial^2}{\partial\theta^2}\log f(X\mid\mu,\theta)\right]
    =-\left(\frac1{2\theta^2}-\frac{\theta}{\theta^3}\right)
    =\frac1{2\theta^2}.\] Thus \(I_X(\theta)=n/(2\theta^2)\), so
\[\operatorname{Var}(W)\ge \frac{1}{n/(2\theta^2)}=\frac{2\theta^2}{n}=\frac{2\sigma^4}{n}.\]
The usual unbiased sample variance satisfies
\[\operatorname{Var}(S^2)=\frac{2\sigma^4}{n-1},\] which is larger than
\(2\sigma^4/n\). Therefore this simple Cramer-Rao calculation does not
prove that \(S^2\) attains the lower bound.

\end{tcolorbox}

\begin{tcolorbox}[enhanced jigsaw, arc=.35mm, bottomrule=.15mm, bottomtitle=1mm, breakable, colback=white, colbacktitle=quarto-callout-important-color!10!white, colframe=quarto-callout-important-color-frame, coltitle=black, left=2mm, leftrule=.75mm, opacityback=0, opacitybacktitle=0.6, rightrule=.15mm, title=\textcolor{quarto-callout-important-color}{\faExclamation}\hspace{0.5em}{Remark}, titlerule=0mm, toprule=.15mm, toptitle=1mm]

\emph{Remark 19}. When nuisance parameters are present, such as unknown
\(\mu\) while estimating \(\sigma^2\), one must be careful with the
precise form of Fisher information. The lecture emphasizes the main
message: the Cramer-Rao lower bound is a benchmark, and not every
natural unbiased estimator attains it.

\end{tcolorbox}

\bookmarksetup{startatroot}

\chapter{Sufficiency and
Unbiasedness}\label{sufficiency-and-unbiasedness}

Sufficient statistics can improve estimators by conditioning away
irrelevant randomness.

\section{Rao-Blackwell theorem}\label{rao-blackwell-theorem}

The Rao-Blackwell theorem is a systematic estimator-improvement
principle.

\begin{tcolorbox}[enhanced jigsaw, arc=.35mm, bottomrule=.15mm, bottomtitle=1mm, breakable, colback=white, colbacktitle=quarto-callout-important-color!10!white, colframe=quarto-callout-important-color-frame, coltitle=black, left=2mm, leftrule=.75mm, opacityback=0, opacitybacktitle=0.6, rightrule=.15mm, title=\textcolor{quarto-callout-important-color}{\faExclamation}\hspace{0.5em}{Theorem}, titlerule=0mm, toprule=.15mm, toptitle=1mm]

\textbf{Theorem 20} (Rao-Blackwell theorem). \emph{Let
\(X_1,\ldots,X_n\) be a sample from a population distribution with
density or mass function \(f(x\mid\theta)\). Let \(W(X)\) be an unbiased
estimator of \(g(\theta)\), and let \(T(X)\) be a sufficient statistic
for \(\theta\). Define \[\phi(T)=\mathbb{E}[W(X)\mid T].\] Then
\[\mathbb{E}_\theta[\phi(T)]=g(\theta)\] and
\[\operatorname{Var}_\theta(\phi(T))\le \operatorname{Var}_\theta(W)\]
for every \(\theta\).}

\end{tcolorbox}

\begin{tcolorbox}[enhanced jigsaw, arc=.35mm, bottomrule=.15mm, bottomtitle=1mm, breakable, colback=white, colbacktitle=quarto-callout-note-color!10!white, colframe=quarto-callout-note-color-frame, coltitle=black, left=2mm, leftrule=.75mm, opacityback=0, opacitybacktitle=0.6, rightrule=.15mm, title=\textcolor{quarto-callout-note-color}{\faInfo}\hspace{0.5em}{Proof}, titlerule=0mm, toprule=.15mm, toptitle=1mm]

\emph{Proof.} For unbiasedness, use the law of total expectation:
\[\mathbb{E}_\theta[\phi(T)]
    =\mathbb{E}_\theta\bigl[\mathbb{E}_\theta(W\mid T)\bigr]
    =\mathbb{E}_\theta(W)=g(\theta).\] For variance, use the law of
total variance: \[\operatorname{Var}_\theta(W)
    =\operatorname{Var}_\theta\bigl(\mathbb{E}_\theta(W\mid T)\bigr)
     +\mathbb{E}_\theta\bigl[\operatorname{Var}_\theta(W\mid T)\bigr].\]
Since the second term is nonnegative,
\[\operatorname{Var}_\theta(W)\ge \operatorname{Var}_\theta(\phi(T)).\]
Thus \(\phi(T)\) is uniformly at least as good as \(W\) among unbiased
estimators.~◻

\end{tcolorbox}

\begin{tcolorbox}[enhanced jigsaw, arc=.35mm, bottomrule=.15mm, bottomtitle=1mm, breakable, colback=white, colbacktitle=quarto-callout-tip-color!10!white, colframe=quarto-callout-tip-color-frame, coltitle=black, left=2mm, leftrule=.75mm, opacityback=0, opacitybacktitle=0.6, rightrule=.15mm, title=\textcolor{quarto-callout-tip-color}{\faLightbulb}\hspace{0.5em}{Key idea}, titlerule=0mm, toprule=.15mm, toptitle=1mm]

Rao-Blackwellization replaces an estimator by its conditional
expectation given a sufficient statistic. This preserves unbiasedness
and reduces variance.

\end{tcolorbox}

\section{Lehmann-Scheffe theorem}\label{lehmann-scheffe-theorem}

Completeness turns the Rao-Blackwell estimator into the unique best
unbiased estimator.

\begin{tcolorbox}[enhanced jigsaw, arc=.35mm, bottomrule=.15mm, bottomtitle=1mm, breakable, colback=white, colbacktitle=quarto-callout-note-color!10!white, colframe=quarto-callout-note-color-frame, coltitle=black, left=2mm, leftrule=.75mm, opacityback=0, opacitybacktitle=0.6, rightrule=.15mm, title=\textcolor{quarto-callout-note-color}{\faInfo}\hspace{0.5em}{Definition}, titlerule=0mm, toprule=.15mm, toptitle=1mm]

\textbf{Definition 21} (Complete sufficient statistic). A statistic
\(T\) is \textbf{complete} for a family of distributions if
\[\mathbb{E}_\theta[g(T)]=0 \quad \text{for all } \theta
    \quad \Longrightarrow \quad
    g(T)=0 \text{ almost surely}.\] If \(T\) is both sufficient and
complete, it is called a complete sufficient statistic.

\end{tcolorbox}

\begin{tcolorbox}[enhanced jigsaw, arc=.35mm, bottomrule=.15mm, bottomtitle=1mm, breakable, colback=white, colbacktitle=quarto-callout-important-color!10!white, colframe=quarto-callout-important-color-frame, coltitle=black, left=2mm, leftrule=.75mm, opacityback=0, opacitybacktitle=0.6, rightrule=.15mm, title=\textcolor{quarto-callout-important-color}{\faExclamation}\hspace{0.5em}{Theorem}, titlerule=0mm, toprule=.15mm, toptitle=1mm]

\textbf{Theorem 22} (Lehmann-Scheffe theorem). \emph{Let \(T(X)\) be a
complete sufficient statistic for \(\theta\). If \(\phi(T)\) is unbiased
for \(g(\theta)\), then \(\phi(T)\) is the unique best unbiased
estimator of \(g(\theta)\).}

\end{tcolorbox}

\begin{tcolorbox}[enhanced jigsaw, arc=.35mm, bottomrule=.15mm, bottomtitle=1mm, breakable, colback=white, colbacktitle=quarto-callout-important-color!10!white, colframe=quarto-callout-important-color-frame, coltitle=black, left=2mm, leftrule=.75mm, opacityback=0, opacitybacktitle=0.6, rightrule=.15mm, title=\textcolor{quarto-callout-important-color}{\faExclamation}\hspace{0.5em}{Remark}, titlerule=0mm, toprule=.15mm, toptitle=1mm]

\emph{Remark 23}. A related characterization says that an unbiased
estimator is best unbiased if and only if it is uncorrelated with every
unbiased estimator of zero.

\end{tcolorbox}

\section{Common UMVUEs}\label{common-umvues}

The following table summarizes common best unbiased estimators discussed
in the course.

{\def\LTcaptype{none} % do not increment counter
\begin{longtable}[]{@{}lll@{}}
\toprule\noalign{}
Parameter & Distribution / model & UMVUE \\
\midrule\noalign{}
\endhead
\bottomrule\noalign{}
\endlastfoot
Mean \(\mu\) & \(\operatorname{Normal}(\mu,\sigma^2)\), known
\(\sigma^2\) & \(\bar X\) \\
Mean \(\mu\) & \(\operatorname{Normal}(\mu,\sigma^2)\), unknown
\(\sigma^2\) & \(\bar X\) \\
Variance \(\sigma^2\) & \(\operatorname{Normal}(\mu,\sigma^2)\), unknown
\(\mu\) & \(\displaystyle \frac1{n-1}\sum_{i=1}^n (X_i-\bar X)^2\) \\
Parameter \(p\) & Bernoulli / binomial & sample proportion
\(\widehat p\) \\
Parameter \(\lambda\) & Poisson & \(\bar X\) \\
Parameter \(\theta\) & \(\operatorname{Uniform}[0,\theta]\) &
\(\displaystyle \frac{n+1}{n}\max_i X_i\) \\
\end{longtable}
}

\begin{tcolorbox}[enhanced jigsaw, arc=.35mm, bottomrule=.15mm, bottomtitle=1mm, breakable, colback=white, colbacktitle=quarto-callout-note-color!10!white, colframe=quarto-callout-note-color-frame, coltitle=black, left=2mm, leftrule=.75mm, opacityback=0, opacitybacktitle=0.6, rightrule=.15mm, title=\textcolor{quarto-callout-note-color}{\faInfo}\hspace{0.5em}{Example}, titlerule=0mm, toprule=.15mm, toptitle=1mm]

\textbf{Example 24} (Uniform UMVUE). Let
\(X_1,\ldots,X_n\overset{\text{iid}}{\sim}\operatorname{Uniform}(0,\theta)\).
Show that \[\widehat\theta=\frac{n+1}{n}X_{(n)}\] is unbiased for
\(\theta\), where \(X_{(n)}=\max_i X_i\).

\end{tcolorbox}

\begin{tcolorbox}[enhanced jigsaw, arc=.35mm, bottomrule=.15mm, bottomtitle=1mm, breakable, colback=white, colbacktitle=quarto-callout-tip-color!10!white, colframe=quarto-callout-tip-color-frame, coltitle=black, left=2mm, leftrule=.75mm, opacityback=0, opacitybacktitle=0.6, rightrule=.15mm, title=\textcolor{quarto-callout-tip-color}{\faLightbulb}\hspace{0.5em}{Solution}, titlerule=0mm, toprule=.15mm, toptitle=1mm]

The CDF of \(X_{(n)}\) is
\[\mathbb{P}(X_{(n)}\le x)=\left(\frac{x}{\theta}\right)^n,
    \qquad 0\le x\le \theta.\] Thus the density is
\[f_{X_{(n)}}(x)=\frac{n x^{n-1}}{\theta^n},
    \qquad 0\le x\le \theta.\] Then \[\mathbb{E}[X_{(n)}]
    =\int_0^\theta x\frac{n x^{n-1}}{\theta^n}\,dx
    =\frac{n}{\theta^n}\cdot \frac{\theta^{n+1}}{n+1}
    =\frac{n}{n+1}\theta.\] Therefore
\[\mathbb{E}\left[\frac{n+1}{n}X_{(n)}\right]=\theta.\]

\end{tcolorbox}

\bookmarksetup{startatroot}

\chapter{Loss Functions and Risk}\label{loss-functions-and-risk}

Decision theory evaluates estimators by the loss incurred after making a
decision.

\section{Loss functions}\label{loss-functions}

A loss function measures the cost of estimating \(\theta\) by an action
\(a\).

\begin{tcolorbox}[enhanced jigsaw, arc=.35mm, bottomrule=.15mm, bottomtitle=1mm, breakable, colback=white, colbacktitle=quarto-callout-note-color!10!white, colframe=quarto-callout-note-color-frame, coltitle=black, left=2mm, leftrule=.75mm, opacityback=0, opacitybacktitle=0.6, rightrule=.15mm, title=\textcolor{quarto-callout-note-color}{\faInfo}\hspace{0.5em}{Definition}, titlerule=0mm, toprule=.15mm, toptitle=1mm]

\textbf{Definition 25} (Loss function). Let \(\Theta\) be the parameter
space and \(\mathcal A\) be the action space. A \textbf{loss function}
is a function \[L:\Theta\times\mathcal A\to \mathbb{R}\] where
\(L(\theta,a)\) measures the cost of taking action \(a\) when the true
parameter is \(\theta\).

\end{tcolorbox}

For point estimation, the action space is often \(\mathcal A=\Theta\).
Common losses include:

\begin{itemize}
\item
  absolute error loss: \[L(\theta,a)=|a-\theta|;\]
\item
  squared error loss: \[L(\theta,a)=(a-\theta)^2;\]
\item
  zero-one loss for exact decisions: \[L(\theta,a)=
          \begin{cases}
          0, & a=\theta,\\
          1, & a\ne \theta.
          \end{cases}\]
\end{itemize}

\section{Risk function}\label{risk-function}

The risk is the expected loss when a decision rule is used.

\begin{tcolorbox}[enhanced jigsaw, arc=.35mm, bottomrule=.15mm, bottomtitle=1mm, breakable, colback=white, colbacktitle=quarto-callout-note-color!10!white, colframe=quarto-callout-note-color-frame, coltitle=black, left=2mm, leftrule=.75mm, opacityback=0, opacitybacktitle=0.6, rightrule=.15mm, title=\textcolor{quarto-callout-note-color}{\faInfo}\hspace{0.5em}{Definition}, titlerule=0mm, toprule=.15mm, toptitle=1mm]

\textbf{Definition 26} (Decision rule and risk). A decision rule is a
function \[\delta:\mathcal X\to \mathcal A\] that selects an action
based on the observed data. Its risk function is
\[R(\theta,\delta)=\mathbb{E}_\theta[L(\theta,\delta(X))].\]

\end{tcolorbox}

For squared error loss,
\[R(\theta,\delta)=\mathbb{E}_\theta[(\delta(X)-\theta)^2]=\operatorname{MSE}_\theta(\delta).\]
Thus MSE is a special case of risk.

\begin{tcolorbox}[enhanced jigsaw, arc=.35mm, bottomrule=.15mm, bottomtitle=1mm, breakable, colback=white, colbacktitle=quarto-callout-tip-color!10!white, colframe=quarto-callout-tip-color-frame, coltitle=black, left=2mm, leftrule=.75mm, opacityback=0, opacitybacktitle=0.6, rightrule=.15mm, title=\textcolor{quarto-callout-tip-color}{\faLightbulb}\hspace{0.5em}{Key idea}, titlerule=0mm, toprule=.15mm, toptitle=1mm]

Risk comparison Given two estimators \(\delta_1\) and \(\delta_2\), if
\[R(\theta,\delta_1)<R(\theta,\delta_2)\] for all \(\theta\), then
\(\delta_1\) is uniformly better under the chosen loss.

\end{tcolorbox}

\bookmarksetup{startatroot}

\chapter{Risk Examples}\label{risk-examples}

The risk function can lead to conclusions different from unbiasedness.

\section{Bernoulli shrinkage
estimator}\label{bernoulli-shrinkage-estimator}

The Bernoulli example illustrates how shrinkage changes risk.

\begin{tcolorbox}[enhanced jigsaw, arc=.35mm, bottomrule=.15mm, bottomtitle=1mm, breakable, colback=white, colbacktitle=quarto-callout-note-color!10!white, colframe=quarto-callout-note-color-frame, coltitle=black, left=2mm, leftrule=.75mm, opacityback=0, opacitybacktitle=0.6, rightrule=.15mm, title=\textcolor{quarto-callout-note-color}{\faInfo}\hspace{0.5em}{Example}, titlerule=0mm, toprule=.15mm, toptitle=1mm]

\textbf{Example 27} (Bernoulli risk comparison). Suppose
\(X_1,\ldots,X_n\overset{\text{iid}}{\sim}\operatorname{Bernoulli}(p)\).
Compare \[\bar X=\frac1n\sum_{i=1}^n X_i\] with the shrinkage estimator
\[\widehat p_B=\frac{\sum_{i=1}^n X_i+\sqrt n/4}{n+\sqrt n}.\] Use
squared error loss.

\end{tcolorbox}

\begin{tcolorbox}[enhanced jigsaw, arc=.35mm, bottomrule=.15mm, bottomtitle=1mm, breakable, colback=white, colbacktitle=quarto-callout-tip-color!10!white, colframe=quarto-callout-tip-color-frame, coltitle=black, left=2mm, leftrule=.75mm, opacityback=0, opacitybacktitle=0.6, rightrule=.15mm, title=\textcolor{quarto-callout-tip-color}{\faLightbulb}\hspace{0.5em}{Solution}, titlerule=0mm, toprule=.15mm, toptitle=1mm]

Let \(S=\sum_{i=1}^n X_i\). Then
\[\operatorname{MSE}_p(\bar X)=\frac{p(1-p)}{n}.\] For
\[\widehat p_B=\frac{S+c}{n+d},
    \qquad c=\frac{\sqrt n}{4},\quad d=\sqrt n,\] we have
\[\mathbb{E}_p(\widehat p_B)=\frac{np+c}{n+d},
    \qquad
    \operatorname{Var}_p(\widehat p_B)=\frac{np(1-p)}{(n+d)^2}.\] Thus
\[\operatorname{MSE}_p(\widehat p_B)
    =\frac{np(1-p)}{(n+d)^2}
    +\left(\frac{np+c}{n+d}-p\right)^2.\] Since \(d=\sqrt n\) and
\(c=\sqrt n/4\), \[\frac{np+c}{n+d}-p
    =\frac{\sqrt n(1/4-p)}{n+\sqrt n}.\] Therefore
\[\operatorname{MSE}_p(\widehat p_B)
    =\frac{np(1-p)}{(n+\sqrt n)^2}
    +\frac{n(1/4-p)^2}{(n+\sqrt n)^2}.\] This estimator shrinks toward
\(1/4\) in this parametrization. Its risk may be smaller than that of
\(\bar X\) in some parts of the parameter space and larger in others.
The lesson is that the preferred estimator depends on the loss function
and the values of \(p\) considered important.

\end{tcolorbox}

\begin{tcolorbox}[enhanced jigsaw, arc=.35mm, bottomrule=.15mm, bottomtitle=1mm, breakable, colback=white, colbacktitle=quarto-callout-important-color!10!white, colframe=quarto-callout-important-color-frame, coltitle=black, left=2mm, leftrule=.75mm, opacityback=0, opacitybacktitle=0.6, rightrule=.15mm, title=\textcolor{quarto-callout-important-color}{\faExclamation}\hspace{0.5em}{Remark}, titlerule=0mm, toprule=.15mm, toptitle=1mm]

\emph{Remark 28}. The lecture slide displays risk curves for \(n=4\) and
\(n=400\). The qualitative message is that shrinkage estimators can
improve risk in small samples or near the shrinkage target, while the
sample proportion becomes very strong as \(n\) grows.

\end{tcolorbox}

\section{\texorpdfstring{Normal variance: choosing a constant multiple
of
\(S^2\)}{Normal variance: choosing a constant multiple of S\^{}2}}\label{normal-variance-choosing-a-constant-multiple-of-s2}

This example shows how MSE can favor a biased multiple of the unbiased
sample variance.

\begin{tcolorbox}[enhanced jigsaw, arc=.35mm, bottomrule=.15mm, bottomtitle=1mm, breakable, colback=white, colbacktitle=quarto-callout-note-color!10!white, colframe=quarto-callout-note-color-frame, coltitle=black, left=2mm, leftrule=.75mm, opacityback=0, opacitybacktitle=0.6, rightrule=.15mm, title=\textcolor{quarto-callout-note-color}{\faInfo}\hspace{0.5em}{Example}, titlerule=0mm, toprule=.15mm, toptitle=1mm]

\textbf{Example 29} (Normal variance risk). Let
\(X_1,\ldots,X_n\overset{\text{iid}}{\sim}\operatorname{Normal}(\mu,\sigma^2)\)
and consider estimators of the form \[\delta_b(X)=bS^2,
    \qquad
    S^2=\frac1{n-1}\sum_{i=1}^n (X_i-\bar X)^2.\] Find the value of
\(b\) that minimizes MSE for estimating \(\sigma^2\).

\end{tcolorbox}

\begin{tcolorbox}[enhanced jigsaw, arc=.35mm, bottomrule=.15mm, bottomtitle=1mm, breakable, colback=white, colbacktitle=quarto-callout-tip-color!10!white, colframe=quarto-callout-tip-color-frame, coltitle=black, left=2mm, leftrule=.75mm, opacityback=0, opacitybacktitle=0.6, rightrule=.15mm, title=\textcolor{quarto-callout-tip-color}{\faLightbulb}\hspace{0.5em}{Solution}, titlerule=0mm, toprule=.15mm, toptitle=1mm]

We know \[\mathbb{E}(S^2)=\sigma^2,
    \qquad
    \operatorname{Var}(S^2)=\frac{2\sigma^4}{n-1}.\] For
\(\delta_b=bS^2\), \[\mathbb{E}(\delta_b)=b\sigma^2,
    \qquad
    \operatorname{Var}(\delta_b)=b^2\frac{2\sigma^4}{n-1}.\] Thus
\[\begin{aligned}
    R(\sigma^2,\delta_b)
    &=\operatorname{MSE}(\delta_b)\\
    &=b^2\frac{2\sigma^4}{n-1}+(b\sigma^2-\sigma^2)^2\\
    &=\sigma^4\left(\frac{2b^2}{n-1}+(b-1)^2\right).
\end{aligned}\] Differentiate with respect to \(b\):
\[\frac{d}{db}\left(\frac{2b^2}{n-1}+(b-1)^2\right)
    =\frac{4b}{n-1}+2(b-1).\] Set equal to zero:
\[\frac{4b}{n-1}+2b-2=0.\] Hence \[b\left(\frac{4}{n-1}+2\right)=2,
    \qquad
    b=\frac{n-1}{n+1}.\] Therefore the MSE-minimizing estimator in this
class is \[\widetilde S^2
    =\frac{n-1}{n+1}S^2
    =\frac1{n+1}\sum_{i=1}^n (X_i-\bar X)^2.\]

\end{tcolorbox}

\begin{tcolorbox}[enhanced jigsaw, arc=.35mm, bottomrule=.15mm, bottomtitle=1mm, breakable, colback=white, colbacktitle=quarto-callout-warning-color!10!white, colframe=quarto-callout-warning-color-frame, coltitle=black, left=2mm, leftrule=.75mm, opacityback=0, opacitybacktitle=0.6, rightrule=.15mm, title=\textcolor{quarto-callout-warning-color}{\faExclamationTriangle}\hspace{0.5em}{Warning}, titlerule=0mm, toprule=.15mm, toptitle=1mm]

Correction note The calculation above gives \(b=(n-1)/(n+1)\) under
squared error loss. This estimator is biased but improves MSE within the
class \(bS^2\).

\end{tcolorbox}

\section{Risk curves for variance
estimators}\label{risk-curves-for-variance-estimators}

The three common estimators of \(\sigma^2\) are
\[S^2=\frac1{n-1}\sum_{i=1}^n (X_i-\bar X)^2,\]
\[\widehat\sigma^2_{\mathrm{MLE}}=\frac1n\sum_{i=1}^n (X_i-\bar X)^2=\frac{n-1}{n}S^2,\]
\[\widetilde S^2=\frac1{n+1}\sum_{i=1}^n (X_i-\bar X)^2=\frac{n-1}{n+1}S^2.\]
Under squared error loss, their risks are proportional to \(\sigma^4\):
\[R(\sigma^2,S^2)=\frac{2\sigma^4}{n-1},\]
\[R(\sigma^2,\widehat\sigma^2_{\mathrm{MLE}})=\frac{(2n-1)\sigma^4}{n^2},\]
\[R(\sigma^2,\widetilde S^2)=\sigma^4\left(\frac{2}{n-1}\left(\frac{n-1}{n+1}\right)^2+\frac{4}{(n+1)^2}\right).\]
The MSE-minimizing constant multiple \(\widetilde S^2\) has the smallest
risk among estimators of the form \(bS^2\).

\section{Different loss functions: Stein's
loss}\label{different-loss-functions-steins-loss}

Changing the loss function can change which estimator is preferred.

\begin{tcolorbox}[enhanced jigsaw, arc=.35mm, bottomrule=.15mm, bottomtitle=1mm, breakable, colback=white, colbacktitle=quarto-callout-note-color!10!white, colframe=quarto-callout-note-color-frame, coltitle=black, left=2mm, leftrule=.75mm, opacityback=0, opacitybacktitle=0.6, rightrule=.15mm, title=\textcolor{quarto-callout-note-color}{\faInfo}\hspace{0.5em}{Definition}, titlerule=0mm, toprule=.15mm, toptitle=1mm]

\textbf{Definition 30} (Stein's loss for variance estimation). For
estimating \(\sigma^2\) by an action \(a>0\), Stein's loss is
\[L(\sigma^2,a)=\frac{a}{\sigma^2}-1-\log\left(\frac{a}{\sigma^2}\right).\]

\end{tcolorbox}

\begin{tcolorbox}[enhanced jigsaw, arc=.35mm, bottomrule=.15mm, bottomtitle=1mm, breakable, colback=white, colbacktitle=quarto-callout-note-color!10!white, colframe=quarto-callout-note-color-frame, coltitle=black, left=2mm, leftrule=.75mm, opacityback=0, opacitybacktitle=0.6, rightrule=.15mm, title=\textcolor{quarto-callout-note-color}{\faInfo}\hspace{0.5em}{Example}, titlerule=0mm, toprule=.15mm, toptitle=1mm]

\textbf{Example 31} (Stein's loss for \(bS^2\)). For estimators
\(\delta_b(X)=bS^2\), compute the risk under Stein's loss up to terms
independent of \(b\).

\end{tcolorbox}

\begin{tcolorbox}[enhanced jigsaw, arc=.35mm, bottomrule=.15mm, bottomtitle=1mm, breakable, colback=white, colbacktitle=quarto-callout-tip-color!10!white, colframe=quarto-callout-tip-color-frame, coltitle=black, left=2mm, leftrule=.75mm, opacityback=0, opacitybacktitle=0.6, rightrule=.15mm, title=\textcolor{quarto-callout-tip-color}{\faLightbulb}\hspace{0.5em}{Solution}, titlerule=0mm, toprule=.15mm, toptitle=1mm]

The risk is \[\begin{aligned}
    R(\sigma^2,\delta_b)
    &=\mathbb{E}\left[\frac{bS^2}{\sigma^2}-1-\log\left(\frac{bS^2}{\sigma^2}\right)\right]\\
    &=b\mathbb{E}\left[\frac{S^2}{\sigma^2}\right]-1-\log b-\mathbb{E}\left[\log\left(\frac{S^2}{\sigma^2}\right)\right].
\end{aligned}\] Since \(\mathbb{E}[S^2/\sigma^2]=1\),
\[R(\sigma^2,\delta_b)=b-\log b-1-C,\] where
\[C=\mathbb{E}\left[\log\left(\frac{S^2}{\sigma^2}\right)\right]\] does
not depend on \(b\). Differentiating \(b-\log b\) gives \[1-\frac1b=0,\]
so \(b=1\). Under Stein's loss, the best choice in this class is
\(\delta_b=S^2\).

\end{tcolorbox}

\begin{tcolorbox}[enhanced jigsaw, arc=.35mm, bottomrule=.15mm, bottomtitle=1mm, breakable, colback=white, colbacktitle=quarto-callout-tip-color!10!white, colframe=quarto-callout-tip-color-frame, coltitle=black, left=2mm, leftrule=.75mm, opacityback=0, opacitybacktitle=0.6, rightrule=.15mm, title=\textcolor{quarto-callout-tip-color}{\faLightbulb}\hspace{0.5em}{Key idea}, titlerule=0mm, toprule=.15mm, toptitle=1mm]

Important lesson The word ``best'' depends on the loss function. Under
squared error loss, a biased multiple of \(S^2\) can be preferred; under
Stein's loss, \(S^2\) is preferred within the class \(bS^2\).

\end{tcolorbox}

\bookmarksetup{startatroot}

\chapter{Bayes Risk}\label{bayes-risk}

Bayes risk averages the frequentist risk over a prior distribution on
the parameter.

\section{Definition}\label{definition-3}

When a prior distribution is specified, we can compare estimators by
their average risk under the prior.

\begin{tcolorbox}[enhanced jigsaw, arc=.35mm, bottomrule=.15mm, bottomtitle=1mm, breakable, colback=white, colbacktitle=quarto-callout-note-color!10!white, colframe=quarto-callout-note-color-frame, coltitle=black, left=2mm, leftrule=.75mm, opacityback=0, opacitybacktitle=0.6, rightrule=.15mm, title=\textcolor{quarto-callout-note-color}{\faInfo}\hspace{0.5em}{Definition}, titlerule=0mm, toprule=.15mm, toptitle=1mm]

\textbf{Definition 32} (Bayes risk). Given a prior distribution
\(\pi(\theta)\), the Bayes risk of a decision rule \(\delta\) is
\[R_B(\delta)=\int_\Theta R(\theta,\delta)\pi(\theta)\,d\theta.\]
Equivalently,
\[R_B(\delta)=\int_\Theta\int_{\mathcal X} L(\theta,\delta(x))f(x\mid\theta)\pi(\theta)\,dx\,d\theta.\]

\end{tcolorbox}

Using Bayes' rule, this can also be written as \[R_B(\delta)
    =\int_{\mathcal X}\left[\int_\Theta L(\theta,\delta(x))\pi(\theta\mid x)\,d\theta\right]m(x)\,dx,\]
where \(\pi(\theta\mid x)\) is the posterior distribution and
\[m(x)=\int f(x\mid\theta)\pi(\theta)\,d\theta\] is the marginal
distribution of the data.

\section{Posterior expected loss}\label{posterior-expected-loss}

For a fixed observed data value \(x\), the posterior expected loss is
\[\int_\Theta L(\theta,a)\pi(\theta\mid x)\,d\theta.\] A Bayes estimator
chooses the action \(a\) minimizing this posterior expected loss.

\begin{tcolorbox}[enhanced jigsaw, arc=.35mm, bottomrule=.15mm, bottomtitle=1mm, breakable, colback=white, colbacktitle=quarto-callout-note-color!10!white, colframe=quarto-callout-note-color-frame, coltitle=black, left=2mm, leftrule=.75mm, opacityback=0, opacitybacktitle=0.6, rightrule=.15mm, title=\textcolor{quarto-callout-note-color}{\faInfo}\hspace{0.5em}{Example}, titlerule=0mm, toprule=.15mm, toptitle=1mm]

\textbf{Example 33} (Bayes estimator under squared error loss). Suppose
\(L(\theta,a)=(a-\theta)^2\). Show that the Bayes estimator is the
posterior mean.

\end{tcolorbox}

\begin{tcolorbox}[enhanced jigsaw, arc=.35mm, bottomrule=.15mm, bottomtitle=1mm, breakable, colback=white, colbacktitle=quarto-callout-tip-color!10!white, colframe=quarto-callout-tip-color-frame, coltitle=black, left=2mm, leftrule=.75mm, opacityback=0, opacitybacktitle=0.6, rightrule=.15mm, title=\textcolor{quarto-callout-tip-color}{\faLightbulb}\hspace{0.5em}{Solution}, titlerule=0mm, toprule=.15mm, toptitle=1mm]

For fixed data \(x\), minimize \[\mathbb{E}[(a-\theta)^2\mid x]\] with
respect to \(a\). Expand: \[\mathbb{E}[(a-\theta)^2\mid x]
    =a^2-2a\mathbb{E}[\theta\mid x]+\mathbb{E}[\theta^2\mid x].\]
Differentiating with respect to \(a\) gives
\[2a-2\mathbb{E}[\theta\mid x]=0.\] Thus \[a=\mathbb{E}[\theta\mid x].\]
So the Bayes estimator under squared error loss is the posterior mean.

\end{tcolorbox}

\begin{tcolorbox}[enhanced jigsaw, arc=.35mm, bottomrule=.15mm, bottomtitle=1mm, breakable, colback=white, colbacktitle=quarto-callout-note-color!10!white, colframe=quarto-callout-note-color-frame, coltitle=black, left=2mm, leftrule=.75mm, opacityback=0, opacitybacktitle=0.6, rightrule=.15mm, title=\textcolor{quarto-callout-note-color}{\faInfo}\hspace{0.5em}{Example}, titlerule=0mm, toprule=.15mm, toptitle=1mm]

\textbf{Example 34} (Bayes estimator under absolute error loss). Suppose
\(L(\theta,a)=|a-\theta|\). The Bayes estimator is a posterior median.

\end{tcolorbox}

\begin{tcolorbox}[enhanced jigsaw, arc=.35mm, bottomrule=.15mm, bottomtitle=1mm, breakable, colback=white, colbacktitle=quarto-callout-tip-color!10!white, colframe=quarto-callout-tip-color-frame, coltitle=black, left=2mm, leftrule=.75mm, opacityback=0, opacitybacktitle=0.6, rightrule=.15mm, title=\textcolor{quarto-callout-tip-color}{\faLightbulb}\hspace{0.5em}{Solution}, titlerule=0mm, toprule=.15mm, toptitle=1mm]

For fixed data \(x\), minimize \[\mathbb{E}[|a-\theta|\mid x]
    =\int |a-\theta|\pi(\theta\mid x)\,d\theta.\] A standard
derivative/subgradient argument shows that a minimizer satisfies
\[\mathbb{P}(\theta\le a\mid x)\ge \frac12,
    \qquad
    \mathbb{P}(\theta\ge a\mid x)\ge \frac12.\] Thus any posterior
median is a Bayes estimator under absolute error loss.

\end{tcolorbox}

\bookmarksetup{startatroot}

\chapter{Summary of Estimator
Criteria}\label{summary-of-estimator-criteria}

This section closes with a comparison table for the main
estimator-evaluation criteria.

{\def\LTcaptype{none} % do not increment counter
\begin{longtable}[]{@{}
  >{\raggedright\arraybackslash}p{(\linewidth - 6\tabcolsep) * \real{0.0704}}
  >{\raggedright\arraybackslash}p{(\linewidth - 6\tabcolsep) * \real{0.3474}}
  >{\raggedright\arraybackslash}p{(\linewidth - 6\tabcolsep) * \real{0.2723}}
  >{\raggedright\arraybackslash}p{(\linewidth - 6\tabcolsep) * \real{0.3005}}@{}}
\toprule\noalign{}
\begin{minipage}[b]{\linewidth}\raggedright
Criterion
\end{minipage} & \begin{minipage}[b]{\linewidth}\raggedright
Definition
\end{minipage} & \begin{minipage}[b]{\linewidth}\raggedright
Desirable property
\end{minipage} & \begin{minipage}[b]{\linewidth}\raggedright
Notes
\end{minipage} \\
\midrule\noalign{}
\endhead
\bottomrule\noalign{}
\endlastfoot
Unbiasedness & \(\mathbb{E}[\widehat\theta]=\theta\) & Expected value
equals the target & Does not guarantee minimum MSE \\
Bias &
\(\operatorname{Bias}(\widehat\theta)=\mathbb{E}[\widehat\theta]-\theta\)
& Bias close to zero & Can be positive or negative \\
Variance & \(\operatorname{Var}(\widehat\theta)\) & Smaller is better,
especially among unbiased estimators & Measures sampling variability \\
MSE & \(\mathbb{E}[(\widehat\theta-\theta)^2]\) & Smaller is better &
\(\operatorname{MSE}=\operatorname{Var}+\operatorname{Bias}^2\) \\
Efficiency & Variance relative to best possible lower bound & Close to
Cramer-Rao lower bound & Usually for unbiased estimators \\
Consistency & \(\widehat\theta_n\xrightarrow{P}\theta\) & Converges to
target as \(n\to\infty\) & Large-sample property \\
Sufficiency & Statistic captures all information about parameter & No
information loss about \(\theta\) & Linked to factorization theorem and
Rao-Blackwell \\
Robustness & Stability under small departures from assumptions & Less
sensitive to outliers/model error & Important in applications \\
\end{longtable}
}

\begin{tcolorbox}[enhanced jigsaw, arc=.35mm, bottomrule=.15mm, bottomtitle=1mm, breakable, colback=white, colbacktitle=quarto-callout-tip-color!10!white, colframe=quarto-callout-tip-color-frame, coltitle=black, left=2mm, leftrule=.75mm, opacityback=0, opacitybacktitle=0.6, rightrule=.15mm, title=\textcolor{quarto-callout-tip-color}{\faLightbulb}\hspace{0.5em}{Key idea}, titlerule=0mm, toprule=.15mm, toptitle=1mm]

Concept map

\end{tcolorbox}

\bookmarksetup{startatroot}

\chapter{Practice Problems}\label{practice-problems-4}

The practice problems below reinforce the main methods for evaluating
estimators.

\begin{tcolorbox}[enhanced jigsaw, arc=.35mm, bottomrule=.15mm, bottomtitle=1mm, breakable, colback=white, colbacktitle=quarto-callout-warning-color!10!white, colframe=quarto-callout-warning-color-frame, coltitle=black, left=2mm, leftrule=.75mm, opacityback=0, opacitybacktitle=0.6, rightrule=.15mm, title=\textcolor{quarto-callout-warning-color}{\faExclamationTriangle}\hspace{0.5em}{Practice Problem}, titlerule=0mm, toprule=.15mm, toptitle=1mm]

\textbf{Practice Problem 35} (Bias-variance decomposition). Let \(W\) be
an estimator of \(\theta\). Prove that
\[\mathbb{E}_\theta[(W-\theta)^2]=\operatorname{Var}_\theta(W)+\operatorname{Bias}_\theta(W)^2.\]

\end{tcolorbox}

\begin{tcolorbox}[enhanced jigsaw, arc=.35mm, bottomrule=.15mm, bottomtitle=1mm, breakable, colback=white, colbacktitle=quarto-callout-tip-color!10!white, colframe=quarto-callout-tip-color-frame, coltitle=black, left=2mm, leftrule=.75mm, opacityback=0, opacitybacktitle=0.6, rightrule=.15mm, title=\textcolor{quarto-callout-tip-color}{\faLightbulb}\hspace{0.5em}{Solution}, titlerule=0mm, toprule=.15mm, toptitle=1mm]

Use \[W-\theta=(W-\mathbb{E}W)+(\mathbb{E}W-\theta).\] Squaring and
taking expectations gives \[\mathbb{E}[(W-\theta)^2]
    =\mathbb{E}[(W-\mathbb{E}W)^2]+2(\mathbb{E}W-\theta)\mathbb{E}[W-\mathbb{E}W]+(\mathbb{E}W-\theta)^2.\]
The middle term is zero, so
\[\mathbb{E}[(W-\theta)^2]=\operatorname{Var}(W)+\operatorname{Bias}(W)^2.\]

\end{tcolorbox}

\begin{tcolorbox}[enhanced jigsaw, arc=.35mm, bottomrule=.15mm, bottomtitle=1mm, breakable, colback=white, colbacktitle=quarto-callout-warning-color!10!white, colframe=quarto-callout-warning-color-frame, coltitle=black, left=2mm, leftrule=.75mm, opacityback=0, opacitybacktitle=0.6, rightrule=.15mm, title=\textcolor{quarto-callout-warning-color}{\faExclamationTriangle}\hspace{0.5em}{Practice Problem}, titlerule=0mm, toprule=.15mm, toptitle=1mm]

\textbf{Practice Problem 36} (MSE for a biased variance estimator). Let
\(X_1,\ldots,X_n\overset{\text{iid}}{\sim}\operatorname{Normal}(\mu,\sigma^2)\)
and let \[\widehat\sigma^2_c=cS^2.\] Find
\(\operatorname{MSE}(\widehat\sigma^2_c)\) as a function of \(c\).

\end{tcolorbox}

\begin{tcolorbox}[enhanced jigsaw, arc=.35mm, bottomrule=.15mm, bottomtitle=1mm, breakable, colback=white, colbacktitle=quarto-callout-tip-color!10!white, colframe=quarto-callout-tip-color-frame, coltitle=black, left=2mm, leftrule=.75mm, opacityback=0, opacitybacktitle=0.6, rightrule=.15mm, title=\textcolor{quarto-callout-tip-color}{\faLightbulb}\hspace{0.5em}{Solution}, titlerule=0mm, toprule=.15mm, toptitle=1mm]

Since \(\mathbb{E}(S^2)=\sigma^2\) and
\(\operatorname{Var}(S^2)=2\sigma^4/(n-1)\),
\[\mathbb{E}(cS^2)=c\sigma^2,
    \qquad
    \operatorname{Var}(cS^2)=c^2\frac{2\sigma^4}{n-1}.\] Therefore
\[\operatorname{MSE}(cS^2)
    =c^2\frac{2\sigma^4}{n-1}+(c\sigma^2-\sigma^2)^2
    =\sigma^4\left(\frac{2c^2}{n-1}+(c-1)^2\right).\]

\end{tcolorbox}

\begin{tcolorbox}[enhanced jigsaw, arc=.35mm, bottomrule=.15mm, bottomtitle=1mm, breakable, colback=white, colbacktitle=quarto-callout-warning-color!10!white, colframe=quarto-callout-warning-color-frame, coltitle=black, left=2mm, leftrule=.75mm, opacityback=0, opacitybacktitle=0.6, rightrule=.15mm, title=\textcolor{quarto-callout-warning-color}{\faExclamationTriangle}\hspace{0.5em}{Practice Problem}, titlerule=0mm, toprule=.15mm, toptitle=1mm]

\textbf{Practice Problem 37} (Optimal constant multiple). Using the
previous problem, find the value of \(c\) that minimizes
\(\operatorname{MSE}(cS^2)\).

\end{tcolorbox}

\begin{tcolorbox}[enhanced jigsaw, arc=.35mm, bottomrule=.15mm, bottomtitle=1mm, breakable, colback=white, colbacktitle=quarto-callout-tip-color!10!white, colframe=quarto-callout-tip-color-frame, coltitle=black, left=2mm, leftrule=.75mm, opacityback=0, opacitybacktitle=0.6, rightrule=.15mm, title=\textcolor{quarto-callout-tip-color}{\faLightbulb}\hspace{0.5em}{Solution}, titlerule=0mm, toprule=.15mm, toptitle=1mm]

Minimize \[h(c)=\frac{2c^2}{n-1}+(c-1)^2.\] Differentiate:
\[h'(c)=\frac{4c}{n-1}+2(c-1).\] Set \(h'(c)=0\):
\[\frac{4c}{n-1}+2c-2=0.\] Thus \[c\left(\frac{4}{n-1}+2\right)=2,
    \qquad
    c=\frac{n-1}{n+1}.\]

\end{tcolorbox}

\begin{tcolorbox}[enhanced jigsaw, arc=.35mm, bottomrule=.15mm, bottomtitle=1mm, breakable, colback=white, colbacktitle=quarto-callout-warning-color!10!white, colframe=quarto-callout-warning-color-frame, coltitle=black, left=2mm, leftrule=.75mm, opacityback=0, opacitybacktitle=0.6, rightrule=.15mm, title=\textcolor{quarto-callout-warning-color}{\faExclamationTriangle}\hspace{0.5em}{Practice Problem}, titlerule=0mm, toprule=.15mm, toptitle=1mm]

\textbf{Practice Problem 38} (Cramer-Rao for Bernoulli). Let
\(X_1,\ldots,X_n\overset{\text{iid}}{\sim}\operatorname{Bernoulli}(p)\).
Find the Fisher information and the Cramer-Rao lower bound for unbiased
estimators of \(p\).

\end{tcolorbox}

\begin{tcolorbox}[enhanced jigsaw, arc=.35mm, bottomrule=.15mm, bottomtitle=1mm, breakable, colback=white, colbacktitle=quarto-callout-tip-color!10!white, colframe=quarto-callout-tip-color-frame, coltitle=black, left=2mm, leftrule=.75mm, opacityback=0, opacitybacktitle=0.6, rightrule=.15mm, title=\textcolor{quarto-callout-tip-color}{\faLightbulb}\hspace{0.5em}{Solution}, titlerule=0mm, toprule=.15mm, toptitle=1mm]

For one observation, \[f(x\mid p)=p^x(1-p)^{1-x},
    \qquad x\in\{0,1\}.\] The log-likelihood for one observation is
\[\ell(p)=X\log p+(1-X)\log(1-p).\] Then
\[\ell'(p)=\frac{X}{p}-\frac{1-X}{1-p}
    =\frac{X-p}{p(1-p)}.\] Thus
\[I_1(p)=\mathbb{E}\left[\left(\frac{X-p}{p(1-p)}\right)^2\right]
    =\frac{\operatorname{Var}(X)}{p^2(1-p)^2}
    =\frac{p(1-p)}{p^2(1-p)^2}
    =\frac1{p(1-p)}.\] For \(n\) iid observations,
\[I_n(p)=\frac{n}{p(1-p)}.\] For unbiased estimators of \(p\),
\[\operatorname{Var}(W)\ge \frac{1}{I_n(p)}=\frac{p(1-p)}{n}.\] Since
\(\bar X\) is unbiased and has variance \(p(1-p)/n\), it attains the
lower bound.

\end{tcolorbox}

\begin{tcolorbox}[enhanced jigsaw, arc=.35mm, bottomrule=.15mm, bottomtitle=1mm, breakable, colback=white, colbacktitle=quarto-callout-warning-color!10!white, colframe=quarto-callout-warning-color-frame, coltitle=black, left=2mm, leftrule=.75mm, opacityback=0, opacitybacktitle=0.6, rightrule=.15mm, title=\textcolor{quarto-callout-warning-color}{\faExclamationTriangle}\hspace{0.5em}{Practice Problem}, titlerule=0mm, toprule=.15mm, toptitle=1mm]

\textbf{Practice Problem 39} (Rao-Blackwell improvement). Suppose \(W\)
is unbiased for \(g(\theta)\) and \(T\) is sufficient. Let
\(\phi(T)=\mathbb{E}(W\mid T)\). Prove that \(\phi(T)\) is unbiased and
has variance no larger than \(W\).

\end{tcolorbox}

\begin{tcolorbox}[enhanced jigsaw, arc=.35mm, bottomrule=.15mm, bottomtitle=1mm, breakable, colback=white, colbacktitle=quarto-callout-tip-color!10!white, colframe=quarto-callout-tip-color-frame, coltitle=black, left=2mm, leftrule=.75mm, opacityback=0, opacitybacktitle=0.6, rightrule=.15mm, title=\textcolor{quarto-callout-tip-color}{\faLightbulb}\hspace{0.5em}{Solution}, titlerule=0mm, toprule=.15mm, toptitle=1mm]

By the law of total expectation,
\[\mathbb{E}[\phi(T)]=\mathbb{E}[\mathbb{E}(W\mid T)]=\mathbb{E}(W)=g(\theta).\]
By the law of total variance,
\[\operatorname{Var}(W)=\operatorname{Var}(\mathbb{E}(W\mid T))+\mathbb{E}[\operatorname{Var}(W\mid T)]
    =\operatorname{Var}(\phi(T))+\mathbb{E}[\operatorname{Var}(W\mid T)].\]
The second term is nonnegative, so
\[\operatorname{Var}(\phi(T))\le \operatorname{Var}(W).\]

\end{tcolorbox}

\begin{tcolorbox}[enhanced jigsaw, arc=.35mm, bottomrule=.15mm, bottomtitle=1mm, breakable, colback=white, colbacktitle=quarto-callout-warning-color!10!white, colframe=quarto-callout-warning-color-frame, coltitle=black, left=2mm, leftrule=.75mm, opacityback=0, opacitybacktitle=0.6, rightrule=.15mm, title=\textcolor{quarto-callout-warning-color}{\faExclamationTriangle}\hspace{0.5em}{Practice Problem}, titlerule=0mm, toprule=.15mm, toptitle=1mm]

\textbf{Practice Problem 40} (Bayes estimator under squared error). Let
\(\pi(\theta\mid x)\) be a posterior density. Show that the posterior
mean minimizes posterior expected squared error loss.

\end{tcolorbox}

\begin{tcolorbox}[enhanced jigsaw, arc=.35mm, bottomrule=.15mm, bottomtitle=1mm, breakable, colback=white, colbacktitle=quarto-callout-tip-color!10!white, colframe=quarto-callout-tip-color-frame, coltitle=black, left=2mm, leftrule=.75mm, opacityback=0, opacitybacktitle=0.6, rightrule=.15mm, title=\textcolor{quarto-callout-tip-color}{\faLightbulb}\hspace{0.5em}{Solution}, titlerule=0mm, toprule=.15mm, toptitle=1mm]

For fixed \(x\), minimize
\[Q(a)=\int (a-\theta)^2\pi(\theta\mid x)\,d\theta.\] Expanding,
\[Q(a)=a^2-2a\mathbb{E}(\theta\mid x)+\mathbb{E}(\theta^2\mid x).\]
Differentiate: \[Q'(a)=2a-2\mathbb{E}(\theta\mid x).\] Thus the
minimizer is \[a=\mathbb{E}(\theta\mid x).\]

\end{tcolorbox}

\bookmarksetup{startatroot}

\chapter{Chapter 14: Hypothesis Tests I --- Methods of Finding
Tests}\label{chapter-14-hypothesis-tests-i-methods-of-finding-tests}

This chapter introduces the foundations of hypothesis testing and
several general methods for constructing tests. The central ideas are
null and alternative hypotheses, test statistics, rejection regions,
Type I and Type II errors, significance level, power, \(p\)-values,
likelihood ratio tests, Bayesian tests, union-intersection tests,
intersection-union tests, and the Neyman--Pearson lemma.

\begin{tcolorbox}[enhanced jigsaw, arc=.35mm, bottomrule=.15mm, bottomtitle=1mm, breakable, colback=white, colbacktitle=quarto-callout-note-color!10!white, colframe=quarto-callout-note-color-frame, coltitle=black, left=2mm, leftrule=.75mm, opacityback=0, opacitybacktitle=0.6, rightrule=.15mm, title=\textcolor{quarto-callout-note-color}{\faInfo}\hspace{0.5em}{Topics}, titlerule=0mm, toprule=.15mm, toptitle=1mm]

Hypothesis testing; null and alternative hypotheses; Type I and Type II
errors; significance level; power; \(p\)-values; likelihood ratio tests;
Bayesian tests; union-intersection tests; intersection-union tests;
radar detection; coin testing; normal mean tests; sufficient statistics
and LRTs.

\end{tcolorbox}

\bookmarksetup{startatroot}

\chapter{Introduction to Hypothesis
Testing}\label{introduction-to-hypothesis-testing}

This section develops the mathematical theory behind hypothesis tests
and explains how classical tests arise from general testing principles.

In introductory statistics, students often learn concrete tests such as
the \(Z\)-test, \(t\)-test, chi-square test, two-proportion test,
two-sample mean test, and \(F\)-test. Here we study the theoretical
principles behind such procedures.

\begin{tcolorbox}[enhanced jigsaw, arc=.35mm, bottomrule=.15mm, bottomtitle=1mm, breakable, colback=white, colbacktitle=quarto-callout-tip-color!10!white, colframe=quarto-callout-tip-color-frame, coltitle=black, left=2mm, leftrule=.75mm, opacityback=0, opacitybacktitle=0.6, rightrule=.15mm, title=\textcolor{quarto-callout-tip-color}{\faLightbulb}\hspace{0.5em}{Key idea}, titlerule=0mm, toprule=.15mm, toptitle=1mm]

Main goal A hypothesis test uses observed data to decide between two
competing statements about a population parameter.

\end{tcolorbox}

For example, a pharmaceutical company may want to know whether a new
drug is effective in treating a disease. A natural pair of hypotheses is
\[H_0: \text{the drug is not effective},
\qquad
H_1: \text{the drug is effective}.\] The null hypothesis \(H_0\)
represents the default or no-effect claim, while the alternative
hypothesis \(H_1\) represents the new effect or departure from the
default.

\section{The testing decision}\label{the-testing-decision}

A hypothesis test is a rule that specifies which sample values lead us
to reject \(H_0\) and which sample values lead us not to reject \(H_0\).

\begin{tcolorbox}[enhanced jigsaw, arc=.35mm, bottomrule=.15mm, bottomtitle=1mm, breakable, colback=white, colbacktitle=quarto-callout-note-color!10!white, colframe=quarto-callout-note-color-frame, coltitle=black, left=2mm, leftrule=.75mm, opacityback=0, opacitybacktitle=0.6, rightrule=.15mm, title=\textcolor{quarto-callout-note-color}{\faInfo}\hspace{0.5em}{Definition}, titlerule=0mm, toprule=.15mm, toptitle=1mm]

\textbf{Definition 1} (Hypothesis test). Let \(X=(X_1,\ldots,X_n)\) be a
sample from a population distribution depending on a parameter
\(\theta\). A hypothesis test is a rule that divides the sample space
into two regions:

\begin{itemize}
\item
  a \textbf{rejection region} \(R\), where we reject \(H_0\);
\item
  an \textbf{acceptance region} \(A=R^c\), where we fail to reject
  \(H_0\).
\end{itemize}

Equivalently, a test may be specified by a test statistic \(W(X)\) and a
threshold rule.

\end{tcolorbox}

\begin{tcolorbox}[enhanced jigsaw, arc=.35mm, bottomrule=.15mm, bottomtitle=1mm, breakable, colback=white, colbacktitle=quarto-callout-important-color!10!white, colframe=quarto-callout-important-color-frame, coltitle=black, left=2mm, leftrule=.75mm, opacityback=0, opacitybacktitle=0.6, rightrule=.15mm, title=\textcolor{quarto-callout-important-color}{\faExclamation}\hspace{0.5em}{Remark}, titlerule=0mm, toprule=.15mm, toptitle=1mm]

\emph{Remark 2} (Terminology). In modern statistical language, one often
says ``fail to reject \(H_0\)'' instead of ``accept \(H_0\).'' This
emphasizes that lack of evidence against \(H_0\) is not the same as
proof that \(H_0\) is true.

\end{tcolorbox}

\section{Mathematical formulation}\label{mathematical-formulation}

The mathematical formulation of testing is based on partitioning the
parameter space.

\begin{tcolorbox}[enhanced jigsaw, arc=.35mm, bottomrule=.15mm, bottomtitle=1mm, breakable, colback=white, colbacktitle=quarto-callout-note-color!10!white, colframe=quarto-callout-note-color-frame, coltitle=black, left=2mm, leftrule=.75mm, opacityback=0, opacitybacktitle=0.6, rightrule=.15mm, title=\textcolor{quarto-callout-note-color}{\faInfo}\hspace{0.5em}{Definition}, titlerule=0mm, toprule=.15mm, toptitle=1mm]

\textbf{Definition 3} (Null and alternative hypotheses). Let \(\Theta\)
be the parameter space. A hypothesis test compares
\[H_0: \theta\in \Theta_0
\qquad \text{versus} \qquad
H_1: \theta\in \Theta_0^c.\] The set \(\Theta_0\) is the null parameter
space, and \(\Theta_0^c\) is the alternative parameter space.

\end{tcolorbox}

Examples include
\[H_0:\mu=100 \quad \text{versus}\quad H_1:\mu\ne 100,\] or
\[H_0:\mu\ge 100 \quad \text{versus}\quad H_1:\mu<100.\]

\bookmarksetup{startatroot}

\chapter{A First Example: Testing Whether a Coin Is
Fair}\label{a-first-example-testing-whether-a-coin-is-fair}

This example introduces the main ingredients of hypothesis testing
through the familiar problem of checking whether a coin is fair.

Suppose \(\theta\) is the probability of heads. We want to test
\[H_0:\theta=\frac12
\qquad \text{versus}\qquad
H_1:\theta\ne \frac12.\] Let
\(X_i\sim \operatorname{Bernoulli}(\theta)\) represent the result of the
\(i\)th toss, where \(X_i=1\) for heads and \(X_i=0\) for tails. Let
\[X=X_1+\cdots+X_n.\] Then \[X\sim \operatorname{Binomial}(n,\theta).\]

\begin{tcolorbox}[enhanced jigsaw, arc=.35mm, bottomrule=.15mm, bottomtitle=1mm, breakable, colback=white, colbacktitle=quarto-callout-note-color!10!white, colframe=quarto-callout-note-color-frame, coltitle=black, left=2mm, leftrule=.75mm, opacityback=0, opacitybacktitle=0.6, rightrule=.15mm, title=\textcolor{quarto-callout-note-color}{\faInfo}\hspace{0.5em}{Example}, titlerule=0mm, toprule=.15mm, toptitle=1mm]

\textbf{Example 4} (Coin test with 100 tosses). Suppose the coin is
tossed \(n=100\) times. Under \(H_0\), \(\theta=1/2\), so
\[X\sim \operatorname{Binomial}\left(100,\frac12\right),
\qquad
\mathbb{E}_{H_0}[X]=50,
\qquad
\operatorname{Var}_{H_0}(X)=25.\] A natural test rejects \(H_0\) when
\(X\) is too far from \(50\).

\end{tcolorbox}

\begin{tcolorbox}[enhanced jigsaw, arc=.35mm, bottomrule=.15mm, bottomtitle=1mm, breakable, colback=white, colbacktitle=quarto-callout-tip-color!10!white, colframe=quarto-callout-tip-color-frame, coltitle=black, left=2mm, leftrule=.75mm, opacityback=0, opacitybacktitle=0.6, rightrule=.15mm, title=\textcolor{quarto-callout-tip-color}{\faLightbulb}\hspace{0.5em}{Solution}, titlerule=0mm, toprule=.15mm, toptitle=1mm]

Let \(t>0\) be a threshold. The test is
\[\text{fail to reject }H_0 \text{ if } |X-50|\le t,
\qquad
\text{reject }H_0 \text{ if } |X-50|>t.\] The threshold should control
the probability of rejecting a fair coin:
\[\mathbb{P}(\text{Type I error})=
\mathbb{P}_{H_0}(|X-50|>t).\] Using the central limit theorem, under
\(H_0\), \[Y=\frac{X-n\theta_0}{\sqrt{n\theta_0(1-\theta_0)}}
=\frac{X-50}{5}
\approx \operatorname{Normal}(0,1).\] For significance level
\(\alpha=0.05\), the two-sided critical value is approximately \(1.96\).
Thus we reject \(H_0\) when \[\left|\frac{X-50}{5}\right|>1.96.\]
Equivalently, \[|X-50|>9.8.\] Using integer values, we fail to reject
\(H_0\) approximately when \[X\in\{41,42,\ldots,59\},\] and reject
\(H_0\) otherwise.

\end{tcolorbox}

\section{Type I and Type II errors}\label{type-i-and-type-ii-errors}

Every hypothesis test can make two kinds of mistakes.

\begin{tcolorbox}[enhanced jigsaw, arc=.35mm, bottomrule=.15mm, bottomtitle=1mm, breakable, colback=white, colbacktitle=quarto-callout-note-color!10!white, colframe=quarto-callout-note-color-frame, coltitle=black, left=2mm, leftrule=.75mm, opacityback=0, opacitybacktitle=0.6, rightrule=.15mm, title=\textcolor{quarto-callout-note-color}{\faInfo}\hspace{0.5em}{Definition}, titlerule=0mm, toprule=.15mm, toptitle=1mm]

\textbf{Definition 5} (Type I and Type II errors). For a hypothesis test
of \(H_0\) versus \(H_1\):

\begin{itemize}
\item
  A \textbf{Type I error} occurs when \(H_0\) is true but we reject
  \(H_0\).
\item
  A \textbf{Type II error} occurs when \(H_1\) is true but we fail to
  reject \(H_0\).
\end{itemize}

The Type I error probability is usually denoted by \(\alpha\), and the
Type II error probability is usually denoted by \(\beta\).

\end{tcolorbox}

For the coin example, if \(\theta_1\ne 1/2\), then
\[\beta(\theta_1)=\mathbb{P}_{\theta_1}(\text{fail to reject }H_0)
=\mathbb{P}_{\theta_1}(41\le X\le 59).\] The power function is
\[\operatorname{Power}(\theta)=\mathbb{P}_\theta(\text{reject }H_0)=1-\beta(\theta)\]
for \(\theta\) in the alternative.

{\def\LTcaptype{none} % do not increment counter
\begin{longtable}[]{@{}ccc@{}}
\toprule\noalign{}
Decision / Truth & \(H_0\) true & \(H_1\) true \\
\midrule\noalign{}
\endhead
\bottomrule\noalign{}
\endlastfoot
Reject \(H_0\) & Type I error & Correct decision \\
Fail to reject \(H_0\) & Correct decision & Type II error \\
\end{longtable}
}

\begin{tcolorbox}[enhanced jigsaw, arc=.35mm, bottomrule=.15mm, bottomtitle=1mm, breakable, colback=white, colbacktitle=quarto-callout-warning-color!10!white, colframe=quarto-callout-warning-color-frame, coltitle=black, left=2mm, leftrule=.75mm, opacityback=0, opacitybacktitle=0.6, rightrule=.15mm, title=\textcolor{quarto-callout-warning-color}{\faExclamationTriangle}\hspace{0.5em}{Warning}, titlerule=0mm, toprule=.15mm, toptitle=1mm]

Trade-off For a fixed sample size, making \(\alpha\) smaller often makes
\(\beta\) larger. Reducing both errors usually requires increasing the
sample size or using a more informative statistic.

\end{tcolorbox}

\bookmarksetup{startatroot}

\chapter{\texorpdfstring{Significance Level, Critical Values, and
\(p\)-Values}{Significance Level, Critical Values, and p-Values}}\label{significance-level-critical-values-and-p-values}

This section explains the practical language of hypothesis testing:
significance level, rejection region, critical value, and \(p\)-value.

\section{Significance level}\label{significance-level}

The significance level controls the probability of falsely rejecting the
null hypothesis.

\begin{tcolorbox}[enhanced jigsaw, arc=.35mm, bottomrule=.15mm, bottomtitle=1mm, breakable, colback=white, colbacktitle=quarto-callout-note-color!10!white, colframe=quarto-callout-note-color-frame, coltitle=black, left=2mm, leftrule=.75mm, opacityback=0, opacitybacktitle=0.6, rightrule=.15mm, title=\textcolor{quarto-callout-note-color}{\faInfo}\hspace{0.5em}{Definition}, titlerule=0mm, toprule=.15mm, toptitle=1mm]

\textbf{Definition 6} (Level-\(\alpha\) test). A test has significance
level \(\alpha\) if
\[\sup_{\theta\in\Theta_0}\mathbb{P}_\theta(\text{reject }H_0)\le \alpha.\]
For a simple null hypothesis \(H_0:\theta=\theta_0\), this reduces to
\[\mathbb{P}_{\theta_0}(\text{reject }H_0)\le \alpha.\]

\end{tcolorbox}

\section{\texorpdfstring{\(p\)-value}{p-value}}\label{p-value}

The \(p\)-value measures how extreme the observed statistic is under the
null model.

\begin{tcolorbox}[enhanced jigsaw, arc=.35mm, bottomrule=.15mm, bottomtitle=1mm, breakable, colback=white, colbacktitle=quarto-callout-note-color!10!white, colframe=quarto-callout-note-color-frame, coltitle=black, left=2mm, leftrule=.75mm, opacityback=0, opacitybacktitle=0.6, rightrule=.15mm, title=\textcolor{quarto-callout-note-color}{\faInfo}\hspace{0.5em}{Definition}, titlerule=0mm, toprule=.15mm, toptitle=1mm]

\textbf{Definition 7} (\(p\)-value). For an observed statistic
\(W(x_1,\ldots,x_n)\), the \(p\)-value is the probability, computed
under \(H_0\), of observing a test statistic at least as extreme as the
one observed. For a right-sided test, this is often
\[p\text{-value}=\mathbb{P}_{H_0}\left(W(X_1,\ldots,X_n)\ge W(x_1,\ldots,x_n)\right).\]
For a left-sided test, this is often
\[p\text{-value}=\mathbb{P}_{H_0}\left(W(X_1,\ldots,X_n)\le W(x_1,\ldots,x_n)\right).\]

\end{tcolorbox}

\begin{tcolorbox}[enhanced jigsaw, arc=.35mm, bottomrule=.15mm, bottomtitle=1mm, breakable, colback=white, colbacktitle=quarto-callout-tip-color!10!white, colframe=quarto-callout-tip-color-frame, coltitle=black, left=2mm, leftrule=.75mm, opacityback=0, opacitybacktitle=0.6, rightrule=.15mm, title=\textcolor{quarto-callout-tip-color}{\faLightbulb}\hspace{0.5em}{Key idea}, titlerule=0mm, toprule=.15mm, toptitle=1mm]

Interpretation The \(p\)-value is the smallest significance level at
which the observed data would lead to rejection of \(H_0\).

\end{tcolorbox}

A small \(p\)-value indicates that the observed data are unusual under
\(H_0\), so the data provide evidence against \(H_0\).

\bookmarksetup{startatroot}

\chapter{Example: Radar Aircraft
Detection}\label{example-radar-aircraft-detection}

This example shows how hypothesis testing appears in signal detection
and illustrates the trade-off between Type I and Type II errors.

A radar system receives a signal \(X\). If no aircraft is present, then
\[X=W.\] If an aircraft is present, then \[X=1+W.\] Here
\[W\sim \operatorname{Normal}\left(0,\sigma^2\right),
\qquad
\sigma^2=\frac19.\] Equivalently, \[X=\theta+W,
\qquad
\theta=\begin{cases}
0, & \text{no aircraft is present},\\
1, & \text{an aircraft is present}.
\end{cases}\] We test \[H_0:\theta=0
\qquad \text{versus}\qquad
H_1:\theta=1.\]

\begin{tcolorbox}[enhanced jigsaw, arc=.35mm, bottomrule=.15mm, bottomtitle=1mm, breakable, colback=white, colbacktitle=quarto-callout-note-color!10!white, colframe=quarto-callout-note-color-frame, coltitle=black, left=2mm, leftrule=.75mm, opacityback=0, opacitybacktitle=0.6, rightrule=.15mm, title=\textcolor{quarto-callout-note-color}{\faInfo}\hspace{0.5em}{Example}, titlerule=0mm, toprule=.15mm, toptitle=1mm]

\textbf{Example 8} (Level \(0.05\) radar test). Construct a level
\(\alpha=0.05\) test that rejects \(H_0\) when \(X>c\).

\end{tcolorbox}

\begin{tcolorbox}[enhanced jigsaw, arc=.35mm, bottomrule=.15mm, bottomtitle=1mm, breakable, colback=white, colbacktitle=quarto-callout-tip-color!10!white, colframe=quarto-callout-tip-color-frame, coltitle=black, left=2mm, leftrule=.75mm, opacityback=0, opacitybacktitle=0.6, rightrule=.15mm, title=\textcolor{quarto-callout-tip-color}{\faLightbulb}\hspace{0.5em}{Solution}, titlerule=0mm, toprule=.15mm, toptitle=1mm]

Under \(H_0\), \(X=W\sim \operatorname{Normal}(0,1/9)\). Thus
\[\mathbb{P}_{H_0}(X>c)=\mathbb{P}(3X>3c)=1-\Phi(3c).\] To make this
probability equal to \(0.05\), choose \[1-\Phi(3c)=0.05.\] Therefore
\[3c=z_{0.95}\approx 1.645,
\qquad
c\approx \frac{1.645}{3}=0.5483.\] The level-\(0.05\) test is
\[\text{reject }H_0 \quad \text{if} \quad X>0.5483.\]

\end{tcolorbox}

\begin{tcolorbox}[enhanced jigsaw, arc=.35mm, bottomrule=.15mm, bottomtitle=1mm, breakable, colback=white, colbacktitle=quarto-callout-note-color!10!white, colframe=quarto-callout-note-color-frame, coltitle=black, left=2mm, leftrule=.75mm, opacityback=0, opacitybacktitle=0.6, rightrule=.15mm, title=\textcolor{quarto-callout-note-color}{\faInfo}\hspace{0.5em}{Example}, titlerule=0mm, toprule=.15mm, toptitle=1mm]

\textbf{Example 9} (Type II error for the radar test). For the level
\(0.05\) radar test above, compute the Type II error probability.

\end{tcolorbox}

\begin{tcolorbox}[enhanced jigsaw, arc=.35mm, bottomrule=.15mm, bottomtitle=1mm, breakable, colback=white, colbacktitle=quarto-callout-tip-color!10!white, colframe=quarto-callout-tip-color-frame, coltitle=black, left=2mm, leftrule=.75mm, opacityback=0, opacitybacktitle=0.6, rightrule=.15mm, title=\textcolor{quarto-callout-tip-color}{\faLightbulb}\hspace{0.5em}{Solution}, titlerule=0mm, toprule=.15mm, toptitle=1mm]

Under \(H_1\), \(X=1+W\) with \(W\sim \operatorname{Normal}(0,1/9)\).
The Type II error probability is
\[\beta=\mathbb{P}_{H_1}(X\le c)=\mathbb{P}(1+W\le c)=\mathbb{P}(W\le c-1).\]
Standardizing gives \[\beta=\Phi(3(c-1)).\] Using
\(c=1.645/3\approx 0.5483\), \[3(c-1)\approx -1.355,\] so
\[\beta\approx \Phi(-1.355)\approx 0.0877.\] Thus the probability of
missing a present aircraft is about \(8.77\%\).

\end{tcolorbox}

\begin{tcolorbox}[enhanced jigsaw, arc=.35mm, bottomrule=.15mm, bottomtitle=1mm, breakable, colback=white, colbacktitle=quarto-callout-note-color!10!white, colframe=quarto-callout-note-color-frame, coltitle=black, left=2mm, leftrule=.75mm, opacityback=0, opacitybacktitle=0.6, rightrule=.15mm, title=\textcolor{quarto-callout-note-color}{\faInfo}\hspace{0.5em}{Example}, titlerule=0mm, toprule=.15mm, toptitle=1mm]

\textbf{Example 10} (Evidence check at level \(0.01\)). Suppose the
observed signal is \(X=0.6\). Determine whether there is sufficient
evidence to reject \(H_0\) at significance level \(\alpha=0.01\).

\end{tcolorbox}

\begin{tcolorbox}[enhanced jigsaw, arc=.35mm, bottomrule=.15mm, bottomtitle=1mm, breakable, colback=white, colbacktitle=quarto-callout-tip-color!10!white, colframe=quarto-callout-tip-color-frame, coltitle=black, left=2mm, leftrule=.75mm, opacityback=0, opacitybacktitle=0.6, rightrule=.15mm, title=\textcolor{quarto-callout-tip-color}{\faLightbulb}\hspace{0.5em}{Solution}, titlerule=0mm, toprule=.15mm, toptitle=1mm]

For a right-sided level-\(0.01\) test, \[\mathbb{P}_{H_0}(X>c)=0.01.\]
Thus \[3c=z_{0.99}\approx 2.326,
\qquad
c\approx \frac{2.326}{3}=0.7753.\] Since the observed value is
\[0.6<0.7753,\] we do not reject \(H_0\) at the \(0.01\) level.

\end{tcolorbox}

\begin{tcolorbox}[enhanced jigsaw, arc=.35mm, bottomrule=.15mm, bottomtitle=1mm, breakable, colback=white, colbacktitle=quarto-callout-note-color!10!white, colframe=quarto-callout-note-color-frame, coltitle=black, left=2mm, leftrule=.75mm, opacityback=0, opacitybacktitle=0.6, rightrule=.15mm, title=\textcolor{quarto-callout-note-color}{\faInfo}\hspace{0.5em}{Example}, titlerule=0mm, toprule=.15mm, toptitle=1mm]

\textbf{Example 11} (Power constraint). Find a critical value \(c\) so
that the probability of missing a present aircraft is less than \(5\%\).
Then compute the resulting significance level.

\end{tcolorbox}

\begin{tcolorbox}[enhanced jigsaw, arc=.35mm, bottomrule=.15mm, bottomtitle=1mm, breakable, colback=white, colbacktitle=quarto-callout-tip-color!10!white, colframe=quarto-callout-tip-color-frame, coltitle=black, left=2mm, leftrule=.75mm, opacityback=0, opacitybacktitle=0.6, rightrule=.15mm, title=\textcolor{quarto-callout-tip-color}{\faLightbulb}\hspace{0.5em}{Solution}, titlerule=0mm, toprule=.15mm, toptitle=1mm]

We want \[\beta=\mathbb{P}_{H_1}(X\le c)=0.05.\] Since under \(H_1\),
\(X\sim \operatorname{Normal}(1,1/9)\),
\[\mathbb{P}(3(X-1)\le 3(c-1))=0.05.\] Thus
\[3(c-1)=z_{0.05}\approx -1.645,\] so
\[c=1-\frac{1.645}{3}\approx 0.4517.\] The resulting significance level
is \[\alpha=\mathbb{P}_{H_0}(X>c)=1-\Phi(3c).\] Since
\(3c\approx 1.355\), \[\alpha\approx 1-\Phi(1.355)\approx 0.0877.\] To
reduce the miss probability to \(5\%\), the false alarm probability must
increase to about \(8.77\%\).

\end{tcolorbox}

\begin{tcolorbox}[enhanced jigsaw, arc=.35mm, bottomrule=.15mm, bottomtitle=1mm, breakable, colback=white, colbacktitle=quarto-callout-note-color!10!white, colframe=quarto-callout-note-color-frame, coltitle=black, left=2mm, leftrule=.75mm, opacityback=0, opacitybacktitle=0.6, rightrule=.15mm, title=\textcolor{quarto-callout-note-color}{\faInfo}\hspace{0.5em}{Example}, titlerule=0mm, toprule=.15mm, toptitle=1mm]

\textbf{Example 12} (\(p\)-value for the radar test). For the observed
value \(X_0=0.6\), compute the \(p\)-value for the right-sided radar
test.

\end{tcolorbox}

\begin{tcolorbox}[enhanced jigsaw, arc=.35mm, bottomrule=.15mm, bottomtitle=1mm, breakable, colback=white, colbacktitle=quarto-callout-tip-color!10!white, colframe=quarto-callout-tip-color-frame, coltitle=black, left=2mm, leftrule=.75mm, opacityback=0, opacitybacktitle=0.6, rightrule=.15mm, title=\textcolor{quarto-callout-tip-color}{\faLightbulb}\hspace{0.5em}{Solution}, titlerule=0mm, toprule=.15mm, toptitle=1mm]

Under \(H_0\), \(X\sim \operatorname{Normal}(0,1/9)\). The standardized
observed statistic is \[Z_0=\frac{0.6}{1/3}=1.8.\] For a right-sided
test,
\[p\text{-value}=\mathbb{P}_{H_0}(X\ge 0.6)=\mathbb{P}(Z\ge 1.8)=1-\Phi(1.8).\]
Numerically, \[p\text{-value}\approx 0.0359.\] Therefore:

\begin{itemize}
\item
  at \(\alpha=0.05\), reject \(H_0\);
\item
  at \(\alpha=0.01\), do not reject \(H_0\).
\end{itemize}

\end{tcolorbox}

\bookmarksetup{startatroot}

\chapter{Classical Test for a Normal
Mean}\label{classical-test-for-a-normal-mean}

This section reviews how a familiar normal mean test fits into the
general framework of test statistics and rejection regions.

Suppose \[X_1,\ldots,X_n\sim \operatorname{Normal}(\mu,\sigma^2),\]
where \(\sigma^2\) is known. We test \[H_0:\mu=\mu_0
\qquad \text{versus}\qquad
H_1:\mu\ne \mu_0.\]

Under \(H_0\),
\[\bar X\sim \operatorname{Normal}\left(\mu_0,\frac{\sigma^2}{n}\right).\]
Therefore
\[Z=\frac{\bar X-\mu_0}{\sigma/\sqrt n}\sim \operatorname{Normal}(0,1).\]

A two-sided level-\(\alpha\) test rejects \(H_0\) when
\[|Z|>z_{1-\alpha/2}.\] Equivalently,
\[\left|\bar X-\mu_0\right|> z_{1-\alpha/2}\frac{\sigma}{\sqrt n}.\]

When \(\sigma^2\) is unknown and the data are normal, the classical test
statistic is \[t=\frac{\bar X-\mu_0}{S/\sqrt n},
\qquad
S^2=\frac{1}{n-1}\sum_{i=1}^n (X_i-\bar X)^2,\] and under \(H_0\),
\[t\sim t_{n-1}.\] The two-sided level-\(\alpha\) test rejects when
\[|t|>t_{n-1,1-\alpha/2}.\]

\bookmarksetup{startatroot}

\chapter{Methods of Finding Tests}\label{methods-of-finding-tests}

This section introduces several general methods for constructing
hypothesis tests.

The main methods discussed in this section are:

\begin{enumerate}
\def\labelenumi{\arabic{enumi}.}
\item
  likelihood ratio tests;
\item
  Bayesian tests;
\item
  union-intersection tests;
\item
  intersection-union tests;
\item
  the Neyman-Pearson lemma for simple hypotheses.
\end{enumerate}

\begin{tcolorbox}[enhanced jigsaw, arc=.35mm, bottomrule=.15mm, bottomtitle=1mm, breakable, colback=white, colbacktitle=quarto-callout-tip-color!10!white, colframe=quarto-callout-tip-color-frame, coltitle=black, left=2mm, leftrule=.75mm, opacityback=0, opacitybacktitle=0.6, rightrule=.15mm, title=\textcolor{quarto-callout-tip-color}{\faLightbulb}\hspace{0.5em}{Key idea}, titlerule=0mm, toprule=.15mm, toptitle=1mm]

Guiding principle Different testing methods correspond to different
statistical philosophies. Likelihood ratio tests compare best
likelihoods; Bayesian tests use posterior probabilities;
union-intersection and intersection-union tests build complex tests from
simpler component tests.

\end{tcolorbox}

\bookmarksetup{startatroot}

\chapter{Likelihood Ratio Tests}\label{likelihood-ratio-tests}

Likelihood ratio tests compare how well the null parameter space
explains the data against how well the full parameter space explains the
data.

Let \(X_1,\ldots,X_n\) be a sample from a population distribution with
density or mass function \(f(x\mid\theta)\). For observed data
\(x=(x_1,\ldots,x_n)\), the likelihood function is
\[L(\theta\mid x)=f(x_1,\ldots,x_n\mid\theta)=\prod_{i=1}^n f(x_i\mid\theta).\]

\begin{tcolorbox}[enhanced jigsaw, arc=.35mm, bottomrule=.15mm, bottomtitle=1mm, breakable, colback=white, colbacktitle=quarto-callout-note-color!10!white, colframe=quarto-callout-note-color-frame, coltitle=black, left=2mm, leftrule=.75mm, opacityback=0, opacitybacktitle=0.6, rightrule=.15mm, title=\textcolor{quarto-callout-note-color}{\faInfo}\hspace{0.5em}{Definition}, titlerule=0mm, toprule=.15mm, toptitle=1mm]

\textbf{Definition 13} (Likelihood ratio statistic). For testing
\[H_0:\theta\in\Theta_0
\qquad \text{versus}\qquad
H_1:\theta\in\Theta_0^c,\] the likelihood ratio statistic is
\[\lambda(x)=\frac{\sup_{\theta\in\Theta_0}L(\theta\mid x)}{\sup_{\theta\in\Theta}L(\theta\mid x)}.\]
A likelihood ratio test rejects \(H_0\) for small values of
\(\lambda(x)\): \[R=\{x:\lambda(x)\le c\},
\qquad 0\le c\le 1.\]

\end{tcolorbox}

Since \(\Theta_0\subseteq \Theta\), we always have
\[0\le \lambda(x)\le 1.\] A small value of \(\lambda(x)\) means that the
null model fits much worse than the unrestricted model.

\section{Simple versus simple likelihood ratio
test}\label{simple-versus-simple-likelihood-ratio-test}

The simplest LRT compares two point hypotheses.

\begin{tcolorbox}[enhanced jigsaw, arc=.35mm, bottomrule=.15mm, bottomtitle=1mm, breakable, colback=white, colbacktitle=quarto-callout-note-color!10!white, colframe=quarto-callout-note-color-frame, coltitle=black, left=2mm, leftrule=.75mm, opacityback=0, opacitybacktitle=0.6, rightrule=.15mm, title=\textcolor{quarto-callout-note-color}{\faInfo}\hspace{0.5em}{Definition}, titlerule=0mm, toprule=.15mm, toptitle=1mm]

\textbf{Definition 14} (Simple likelihood ratio). For
\[H_0:\theta=\theta_0
\qquad \text{versus}\qquad
H_1:\theta=\theta_1,\] define
\[\lambda(x)=\frac{L(\theta_0\mid x)}{L(\theta_1\mid x)}.\] The
likelihood ratio test rejects \(H_0\) for small values of
\(\lambda(x)\).

\end{tcolorbox}

\bookmarksetup{startatroot}

\chapter{Likelihood Ratio Test: Radar
Example}\label{likelihood-ratio-test-radar-example}

This section rederives the radar test using the likelihood ratio method.

Recall that \[X=\theta+W,
\qquad
W\sim \operatorname{Normal}\left(0,\frac19\right),\] and we test
\[H_0:\theta=0
\qquad \text{versus}\qquad
H_1:\theta=1.\]

\begin{tcolorbox}[enhanced jigsaw, arc=.35mm, bottomrule=.15mm, bottomtitle=1mm, breakable, colback=white, colbacktitle=quarto-callout-note-color!10!white, colframe=quarto-callout-note-color-frame, coltitle=black, left=2mm, leftrule=.75mm, opacityback=0, opacitybacktitle=0.6, rightrule=.15mm, title=\textcolor{quarto-callout-note-color}{\faInfo}\hspace{0.5em}{Example}, titlerule=0mm, toprule=.15mm, toptitle=1mm]

\textbf{Example 15} (LRT for radar detection). Find the likelihood ratio
statistic and show that the LRT rejects for large values of \(X\).

\end{tcolorbox}

\begin{tcolorbox}[enhanced jigsaw, arc=.35mm, bottomrule=.15mm, bottomtitle=1mm, breakable, colback=white, colbacktitle=quarto-callout-tip-color!10!white, colframe=quarto-callout-tip-color-frame, coltitle=black, left=2mm, leftrule=.75mm, opacityback=0, opacitybacktitle=0.6, rightrule=.15mm, title=\textcolor{quarto-callout-tip-color}{\faLightbulb}\hspace{0.5em}{Solution}, titlerule=0mm, toprule=.15mm, toptitle=1mm]

The density under \(\theta=0\) is
\[L(0\mid x)=\frac{3}{\sqrt{2\pi}}\exp\left(-\frac{9x^2}{2}\right).\]
The density under \(\theta=1\) is
\[L(1\mid x)=\frac{3}{\sqrt{2\pi}}\exp\left(-\frac{9(x-1)^2}{2}\right).\]
Thus the likelihood ratio is \[\lambda(x)=\frac{L(0\mid x)}{L(1\mid x)}
=\exp\left(-\frac{9x^2}{2}+\frac{9(x-1)^2}{2}\right).\] Simplifying,
\[\lambda(x)=\exp\left(\frac{9(1-2x)}{2}\right).\] Since this is a
decreasing function of \(x\), rejecting for small \(\lambda(x)\) is
equivalent to rejecting for large \(x\). Therefore the LRT has the form
\[\text{reject }H_0 \quad \text{if}\quad x>c'.\] For a level \(0.05\)
test, \[\mathbb{P}_{H_0}(X>c')=0.05,\] so
\[c'=\frac{z_{0.95}}{3}\approx \frac{1.645}{3}=0.5483.\] This is the
same test constructed directly from Type I error control.

\end{tcolorbox}

\bookmarksetup{startatroot}

\chapter{LRT for a Normal Mean with Known
Variance}\label{lrt-for-a-normal-mean-with-known-variance}

This section shows that the classical two-sided \(Z\)-test is a
likelihood ratio test.

Suppose \[X_1,\ldots,X_n\sim \operatorname{Normal}(\mu,\sigma^2),\]
where \(\sigma^2\) is known. We test \[H_0:\mu=\mu_0
\qquad \text{versus}\qquad
H_1:\mu\ne\mu_0.\]

\begin{tcolorbox}[enhanced jigsaw, arc=.35mm, bottomrule=.15mm, bottomtitle=1mm, breakable, colback=white, colbacktitle=quarto-callout-note-color!10!white, colframe=quarto-callout-note-color-frame, coltitle=black, left=2mm, leftrule=.75mm, opacityback=0, opacitybacktitle=0.6, rightrule=.15mm, title=\textcolor{quarto-callout-note-color}{\faInfo}\hspace{0.5em}{Example}, titlerule=0mm, toprule=.15mm, toptitle=1mm]

\textbf{Example 16} (Normal mean LRT with known variance). Derive the
likelihood ratio statistic.

\end{tcolorbox}

\begin{tcolorbox}[enhanced jigsaw, arc=.35mm, bottomrule=.15mm, bottomtitle=1mm, breakable, colback=white, colbacktitle=quarto-callout-tip-color!10!white, colframe=quarto-callout-tip-color-frame, coltitle=black, left=2mm, leftrule=.75mm, opacityback=0, opacitybacktitle=0.6, rightrule=.15mm, title=\textcolor{quarto-callout-tip-color}{\faLightbulb}\hspace{0.5em}{Solution}, titlerule=0mm, toprule=.15mm, toptitle=1mm]

The likelihood is \[L(\mu\mid x)= (2\pi\sigma^2)^{-n/2}
\exp\left(-\frac{1}{2\sigma^2}\sum_{i=1}^n (x_i-\mu)^2\right).\] Under
\(H_0\), the best null value is fixed at \(\mu_0\). Under the full
parameter space, the MLE is \[\widehat\mu=\bar x.\] Therefore
\[\lambda(x)=\frac{L(\mu_0\mid x)}{L(\bar x\mid x)}.\] Using the
identity \[\sum_{i=1}^n (x_i-\mu_0)^2
=
\sum_{i=1}^n (x_i-\bar x)^2+n(\bar x-\mu_0)^2,\] we obtain
\[\lambda(x)=\exp\left(-\frac{n(\bar x-\mu_0)^2}{2\sigma^2}\right).\]
This statistic decreases as \(|\bar x-\mu_0|\) increases. Thus rejecting
for small \(\lambda(x)\) is equivalent to rejecting for large
\[\left|\frac{\bar X-\mu_0}{\sigma/\sqrt n}\right|.\] Therefore the
level-\(\alpha\) LRT rejects when
\[\left|\frac{\bar X-\mu_0}{\sigma/\sqrt n}\right|>z_{1-\alpha/2}.\]
This is the classical two-sided \(Z\)-test.

\end{tcolorbox}

The cutoff \(c\) in the likelihood-ratio form can be written as
\[c=\exp\left(-\frac{z_{1-\alpha/2}^2}{2}\right).\]

\bookmarksetup{startatroot}

\chapter{LRT for a Normal Mean with Unknown
Variance}\label{lrt-for-a-normal-mean-with-unknown-variance}

This section shows that the classical two-sided Student's \(t\)-test can
also be derived as a likelihood ratio test.

Suppose \[X_1,\ldots,X_n\sim \operatorname{Normal}(\mu,\sigma^2),\]
where both \(\mu\) and \(\sigma^2\) are unknown. We test \[H_0:\mu=\mu_0
\qquad \text{versus}\qquad
H_1:\mu\ne\mu_0.\]

\begin{tcolorbox}[enhanced jigsaw, arc=.35mm, bottomrule=.15mm, bottomtitle=1mm, breakable, colback=white, colbacktitle=quarto-callout-note-color!10!white, colframe=quarto-callout-note-color-frame, coltitle=black, left=2mm, leftrule=.75mm, opacityback=0, opacitybacktitle=0.6, rightrule=.15mm, title=\textcolor{quarto-callout-note-color}{\faInfo}\hspace{0.5em}{Example}, titlerule=0mm, toprule=.15mm, toptitle=1mm]

\textbf{Example 17} (Normal mean LRT with unknown variance). Derive the
LRT and connect it to the Student \(t\) statistic.

\end{tcolorbox}

\begin{tcolorbox}[enhanced jigsaw, arc=.35mm, bottomrule=.15mm, bottomtitle=1mm, breakable, colback=white, colbacktitle=quarto-callout-tip-color!10!white, colframe=quarto-callout-tip-color-frame, coltitle=black, left=2mm, leftrule=.75mm, opacityback=0, opacitybacktitle=0.6, rightrule=.15mm, title=\textcolor{quarto-callout-tip-color}{\faLightbulb}\hspace{0.5em}{Solution}, titlerule=0mm, toprule=.15mm, toptitle=1mm]

Under the full model, the MLEs are \[\widehat\mu=\bar x,
\qquad
\widehat\sigma^2=\frac1n\sum_{i=1}^n (x_i-\bar x)^2.\] Under \(H_0\),
\(\mu=\mu_0\), and the MLE of \(\sigma^2\) is
\[\widehat\sigma_0^2=\frac1n\sum_{i=1}^n (x_i-\mu_0)^2.\] After
substituting these into the normal likelihood, the likelihood ratio is
\[\lambda(x)=
\left(\frac{\widehat\sigma_0^2}{\widehat\sigma^2}\right)^{-n/2}.\] Using
\[\sum_{i=1}^n (x_i-\mu_0)^2
=
\sum_{i=1}^n (x_i-\bar x)^2+n(\bar x-\mu_0)^2,\] we get
\[\frac{\widehat\sigma_0^2}{\widehat\sigma^2}
=1+\frac{n(\bar x-\mu_0)^2}{\sum_{i=1}^n (x_i-\bar x)^2}.\] Let
\[S^2=\frac{1}{n-1}\sum_{i=1}^n (X_i-\bar X)^2\] and
\[t=\frac{\bar X-\mu_0}{S/\sqrt n}.\] Then
\[\frac{\widehat\sigma_0^2}{\widehat\sigma^2}=1+\frac{t^2}{n-1}.\]
Therefore \(\lambda(x)\) is small exactly when \(|t|\) is large. The LRT
rejects \(H_0\) when \[|t|>t_{n-1,1-\alpha/2}.\] This is the classical
two-sided Student's \(t\)-test with \(n-1\) degrees of freedom.

\end{tcolorbox}

\begin{tcolorbox}[enhanced jigsaw, arc=.35mm, bottomrule=.15mm, bottomtitle=1mm, breakable, colback=white, colbacktitle=quarto-callout-important-color!10!white, colframe=quarto-callout-important-color-frame, coltitle=black, left=2mm, leftrule=.75mm, opacityback=0, opacitybacktitle=0.6, rightrule=.15mm, title=\textcolor{quarto-callout-important-color}{\faExclamation}\hspace{0.5em}{Remark}, titlerule=0mm, toprule=.15mm, toptitle=1mm]

\emph{Remark 18}. The classical \(Z\)-, \(t\)-, chi-square, proportion,
pooled two-sample \(t\)-, two-proportion, and \(F\)-tests can all be
interpreted as special cases or asymptotic versions of likelihood ratio
tests.

\end{tcolorbox}

\bookmarksetup{startatroot}

\chapter{LRTs and Sufficient
Statistics}\label{lrts-and-sufficient-statistics}

This section explains why likelihood ratio tests can be computed using
sufficient statistics without losing information.

\begin{tcolorbox}[enhanced jigsaw, arc=.35mm, bottomrule=.15mm, bottomtitle=1mm, breakable, colback=white, colbacktitle=quarto-callout-important-color!10!white, colframe=quarto-callout-important-color-frame, coltitle=black, left=2mm, leftrule=.75mm, opacityback=0, opacitybacktitle=0.6, rightrule=.15mm, title=\textcolor{quarto-callout-important-color}{\faExclamation}\hspace{0.5em}{Theorem}, titlerule=0mm, toprule=.15mm, toptitle=1mm]

\textbf{Theorem 19} (LRT based on a sufficient statistic). \emph{Suppose
\(T(X)\) is sufficient for \(\theta\). Let \(\lambda(x)\) be the
likelihood ratio statistic based on the full data \(X\), and let
\(\lambda^*(T(x))\) be the likelihood ratio statistic based on the
statistic \(T\). Then \[\lambda^*(T(x))=\lambda(x).\]}

\end{tcolorbox}

\begin{tcolorbox}[enhanced jigsaw, arc=.35mm, bottomrule=.15mm, bottomtitle=1mm, breakable, colback=white, colbacktitle=quarto-callout-note-color!10!white, colframe=quarto-callout-note-color-frame, coltitle=black, left=2mm, leftrule=.75mm, opacityback=0, opacitybacktitle=0.6, rightrule=.15mm, title=\textcolor{quarto-callout-note-color}{\faInfo}\hspace{0.5em}{Proof}, titlerule=0mm, toprule=.15mm, toptitle=1mm]

\emph{Proof.} By the factorization theorem, the likelihood can be
written as \[L(\theta\mid x)=g(T(x),\theta)h(x),\] where \(h(x)\) does
not depend on \(\theta\). Then
\[\lambda(x)=\frac{\sup_{\theta\in\Theta_0}g(T(x),\theta)h(x)}{\sup_{\theta\in\Theta}g(T(x),\theta)h(x)}.\]
The factor \(h(x)\) cancels, so
\[\lambda(x)=\frac{\sup_{\theta\in\Theta_0}g(T(x),\theta)}{\sup_{\theta\in\Theta}g(T(x),\theta)},\]
which is exactly the likelihood ratio based on \(T(x)\).~◻

\end{tcolorbox}

\begin{tcolorbox}[enhanced jigsaw, arc=.35mm, bottomrule=.15mm, bottomtitle=1mm, breakable, colback=white, colbacktitle=quarto-callout-note-color!10!white, colframe=quarto-callout-note-color-frame, coltitle=black, left=2mm, leftrule=.75mm, opacityback=0, opacitybacktitle=0.6, rightrule=.15mm, title=\textcolor{quarto-callout-note-color}{\faInfo}\hspace{0.5em}{Example}, titlerule=0mm, toprule=.15mm, toptitle=1mm]

\textbf{Example 20} (Normal mean with known variance using \(\bar X\)).
Suppose \(X_1,\ldots,X_n\sim \operatorname{Normal}(\mu,\sigma^2)\) with
known \(\sigma^2\). Since \(\bar X\) is sufficient for \(\mu\), derive
the LRT using \(T=\bar X\).

\end{tcolorbox}

\begin{tcolorbox}[enhanced jigsaw, arc=.35mm, bottomrule=.15mm, bottomtitle=1mm, breakable, colback=white, colbacktitle=quarto-callout-tip-color!10!white, colframe=quarto-callout-tip-color-frame, coltitle=black, left=2mm, leftrule=.75mm, opacityback=0, opacitybacktitle=0.6, rightrule=.15mm, title=\textcolor{quarto-callout-tip-color}{\faLightbulb}\hspace{0.5em}{Solution}, titlerule=0mm, toprule=.15mm, toptitle=1mm]

The statistic \[T=\bar X\] has distribution
\[T\sim \operatorname{Normal}\left(\mu,\frac{\sigma^2}{n}\right).\] Thus
the likelihood based on \(T=t\) is
\[L(\mu\mid t)=\frac{1}{\sqrt{2\pi\sigma^2/n}}
\exp\left(-\frac{n(t-\mu)^2}{2\sigma^2}\right).\] For testing
\(H_0:\mu=\mu_0\) versus \(H_1:\mu\ne\mu_0\), the unrestricted MLE based
on \(t\) is \(\widehat\mu=t\). Therefore
\[\lambda^*(t)=\frac{L(\mu_0\mid t)}{L(t\mid t)}
=\exp\left(-\frac{n(t-\mu_0)^2}{2\sigma^2}\right).\] Substituting
\(t=\bar x\) gives
\[\lambda^*(\bar x)=\exp\left(-\frac{n(\bar x-\mu_0)^2}{2\sigma^2}\right),\]
which is the same likelihood ratio statistic obtained from the full
sample.

\end{tcolorbox}

\bookmarksetup{startatroot}

\chapter{Bayesian Tests}\label{bayesian-tests}

This section presents hypothesis testing from a Bayesian viewpoint,
where inference is based on posterior probabilities.

In Bayesian statistics, the parameter \(\theta\) is treated as a random
quantity with prior distribution \(\pi(\theta)\). Given data
\(x=(x_1,\ldots,x_n)\), the posterior distribution is
\[\pi(\theta\mid x)=\frac{f(x\mid\theta)\pi(\theta)}{m(x)},\] where
\[m(x)=\int f(x\mid\theta)\pi(\theta)\,d\theta\] is the marginal
distribution of the data.

\begin{tcolorbox}[enhanced jigsaw, arc=.35mm, bottomrule=.15mm, bottomtitle=1mm, breakable, colback=white, colbacktitle=quarto-callout-note-color!10!white, colframe=quarto-callout-note-color-frame, coltitle=black, left=2mm, leftrule=.75mm, opacityback=0, opacitybacktitle=0.6, rightrule=.15mm, title=\textcolor{quarto-callout-note-color}{\faInfo}\hspace{0.5em}{Definition}, titlerule=0mm, toprule=.15mm, toptitle=1mm]

\textbf{Definition 21} (Bayesian test by posterior probabilities). For
testing \[H_0:\theta\in\Theta_0
\qquad \text{versus}\qquad
H_1:\theta\in\Theta_0^c,\] compute the posterior probabilities
\[\mathbb{P}(\theta\in\Theta_0\mid x)
\quad \text{and}\quad
\mathbb{P}(\theta\in\Theta_0^c\mid x).\] A simple Bayesian rule rejects
\(H_0\) if
\[\mathbb{P}(\theta\in\Theta_0\mid x)<\mathbb{P}(\theta\in\Theta_0^c\mid x).\]
Equivalently, reject if \[\mathbb{P}(\theta\in\Theta_0\mid x)<\frac12.\]
A more conservative rule might reject only if
\[\mathbb{P}(\theta\in\Theta_0\mid x)<0.05.\]

\end{tcolorbox}

\section{Bayesian normal mean test}\label{bayesian-normal-mean-test}

We now compute a Bayesian test for a normal mean.

\begin{tcolorbox}[enhanced jigsaw, arc=.35mm, bottomrule=.15mm, bottomtitle=1mm, breakable, colback=white, colbacktitle=quarto-callout-note-color!10!white, colframe=quarto-callout-note-color-frame, coltitle=black, left=2mm, leftrule=.75mm, opacityback=0, opacitybacktitle=0.6, rightrule=.15mm, title=\textcolor{quarto-callout-note-color}{\faInfo}\hspace{0.5em}{Example}, titlerule=0mm, toprule=.15mm, toptitle=1mm]

\textbf{Example 22} (Bayesian test for a normal mean). Suppose
\[X_1,\ldots,X_n\mid\mu \sim \operatorname{Normal}(\mu,\sigma^2),\]
where \(\sigma^2\) is known. Suppose the prior is
\[\mu\sim \operatorname{Normal}(\theta,\tau^2).\] Test
\[H_0:\mu\le \mu_0
\qquad \text{versus}\qquad
H_1:\mu>\mu_0.\] Derive the posterior and the Bayesian decision rule
that rejects when the posterior probability of \(H_0\) is less than
\(1/2\).

\end{tcolorbox}

\begin{tcolorbox}[enhanced jigsaw, arc=.35mm, bottomrule=.15mm, bottomtitle=1mm, breakable, colback=white, colbacktitle=quarto-callout-tip-color!10!white, colframe=quarto-callout-tip-color-frame, coltitle=black, left=2mm, leftrule=.75mm, opacityback=0, opacitybacktitle=0.6, rightrule=.15mm, title=\textcolor{quarto-callout-tip-color}{\faLightbulb}\hspace{0.5em}{Solution}, titlerule=0mm, toprule=.15mm, toptitle=1mm]

The normal-normal conjugate posterior is normal:
\[\mu\mid x \sim \operatorname{Normal}(m_n,v_n),\] where
\[m_n=\frac{\tau^2\sum_{i=1}^n x_i+\sigma^2\theta}{n\tau^2+\sigma^2}\]
and \[v_n=\frac{\sigma^2\tau^2}{\sigma^2+n\tau^2}.\] We reject \(H_0\)
when \[\mathbb{P}(\mu\le \mu_0\mid x)<\frac12.\] Since the posterior
distribution is normal and symmetric, this condition is equivalent to
\[m_n>\mu_0.\] Thus the Bayesian test rejects when
\[\frac{\tau^2\sum_{i=1}^n x_i+\sigma^2\theta}{n\tau^2+\sigma^2}>\mu_0.\]
Equivalently, \[\bar x>\mu_0+\frac{\sigma^2}{n\tau^2}(\mu_0-\theta).\]
This shows how the prior mean \(\theta\) shifts the rejection threshold.

\end{tcolorbox}

\bookmarksetup{startatroot}

\chapter{Union-Intersection Tests}\label{union-intersection-tests}

This section explains how to construct a test for a complicated
alternative by combining simpler component tests.

The union-intersection test is useful when the null hypothesis can be
written as an intersection of simpler hypotheses.

Suppose
\[H_0:\theta\in\Theta_0=\bigcap_{\gamma\in\Gamma}\Theta_\gamma.\] Then
by De Morgan's law,
\[H_1:\theta\in\Theta_0^c=\bigcup_{\gamma\in\Gamma}\Theta_\gamma^c.\]
For each \(\gamma\), we test \[H_{0\gamma}:\theta\in\Theta_\gamma
\qquad \text{versus}\qquad
H_{1\gamma}:\theta\in\Theta_\gamma^c.\] Let the rejection region for the
\(\gamma\)th subtest be \[R_\gamma=\{x:T_\gamma(x)>c\}.\] The
union-intersection test rejects if any component test rejects:
\[R=\bigcup_{\gamma\in\Gamma}R_\gamma.\] Equivalently,
\[R=\left\{x:\sup_{\gamma\in\Gamma}T_\gamma(x)>c\right\}.\] Thus the
combined test statistic is \[T(x)=\sup_{\gamma\in\Gamma}T_\gamma(x).\]

\begin{tcolorbox}[enhanced jigsaw, arc=.35mm, bottomrule=.15mm, bottomtitle=1mm, breakable, colback=white, colbacktitle=quarto-callout-note-color!10!white, colframe=quarto-callout-note-color-frame, coltitle=black, left=2mm, leftrule=.75mm, opacityback=0, opacitybacktitle=0.6, rightrule=.15mm, title=\textcolor{quarto-callout-note-color}{\faInfo}\hspace{0.5em}{Example}, titlerule=0mm, toprule=.15mm, toptitle=1mm]

\textbf{Example 23} (UIT for a two-sided normal mean test). Suppose
\(X_1,\ldots,X_n\) are normal and we want to test \[H_0:\mu=\mu_0
\qquad \text{versus}\qquad
H_1:\mu\ne \mu_0.\] Explain how this can be constructed from one-sided
tests.

\end{tcolorbox}

\begin{tcolorbox}[enhanced jigsaw, arc=.35mm, bottomrule=.15mm, bottomtitle=1mm, breakable, colback=white, colbacktitle=quarto-callout-tip-color!10!white, colframe=quarto-callout-tip-color-frame, coltitle=black, left=2mm, leftrule=.75mm, opacityback=0, opacitybacktitle=0.6, rightrule=.15mm, title=\textcolor{quarto-callout-tip-color}{\faLightbulb}\hspace{0.5em}{Solution}, titlerule=0mm, toprule=.15mm, toptitle=1mm]

The null hypothesis can be written as \[H_0:\mu=\mu_0
\quad \Longleftrightarrow \quad
\{\mu\le \mu_0\}\cap \{\mu\ge \mu_0\}.\] The alternative is
\[H_1:\mu\ne \mu_0
\quad \Longleftrightarrow \quad
\{\mu>\mu_0\}\cup \{\mu<\mu_0\}.\] Thus we combine two one-sided tests:

\begin{itemize}
\item
  reject \(H_{0L}:\mu\le \mu_0\) for large positive values of the
  standardized statistic;
\item
  reject \(H_{0R}:\mu\ge \mu_0\) for large negative values of the
  standardized statistic.
\end{itemize}

If \(\sigma^2\) is known, the standardized statistic is
\[Z=\frac{\bar X-\mu_0}{\sigma/\sqrt n}.\] The UIT rejects when either
tail is too extreme, which is equivalent to \[|Z|>z_{1-\alpha/2}.\] This
is the classical two-sided \(Z\)-test.

If \(\sigma^2\) is unknown, replace \(\sigma\) by \(S\):
\[t=\frac{\bar X-\mu_0}{S/\sqrt n},\] and reject when
\[|t|>t_{n-1,1-\alpha/2}.\] This is the classical two-sided \(t\)-test.

\end{tcolorbox}

\bookmarksetup{startatroot}

\chapter{Intersection-Union Tests}\label{intersection-union-tests}

This section explains the complementary construction, where we reject
only if all component tests reject.

The intersection-union test is useful when the alternative hypothesis
requires several conditions to hold simultaneously.

Suppose
\[H_0:\theta\in\Theta_0=\bigcup_{\gamma\in\Gamma}\Theta_\gamma.\] Then
by De Morgan's law,
\[H_1:\theta\in\Theta_0^c=\bigcap_{\gamma\in\Gamma}\Theta_\gamma^c.\]
For each component test, let \[R_\gamma=\{x:T_\gamma(x)>c\}.\] The IUT
rejects only when all component tests reject:
\[R=\bigcap_{\gamma\in\Gamma}R_\gamma.\] Equivalently,
\[R=\left\{x:\inf_{\gamma\in\Gamma}T_\gamma(x)>c\right\}.\] Thus the
combined test statistic is \[T(x)=\inf_{\gamma\in\Gamma}T_\gamma(x).\]

\begin{tcolorbox}[enhanced jigsaw, arc=.35mm, bottomrule=.15mm, bottomtitle=1mm, breakable, colback=white, colbacktitle=quarto-callout-note-color!10!white, colframe=quarto-callout-note-color-frame, coltitle=black, left=2mm, leftrule=.75mm, opacityback=0, opacitybacktitle=0.6, rightrule=.15mm, title=\textcolor{quarto-callout-note-color}{\faInfo}\hspace{0.5em}{Example}, titlerule=0mm, toprule=.15mm, toptitle=1mm]

\textbf{Example 24} (Acceptance sampling for upholstery fabric). A batch
of upholstery fabric has two quality parameters:

\begin{itemize}
\item
  \(\theta_1\): mean breaking strength, which should exceed \(50\)
  pounds;
\item
  \(\theta_2\): probability of passing a flammability test, which should
  exceed \(0.95\).
\end{itemize}

Set up an intersection-union test for determining whether a batch is
acceptable.

\end{tcolorbox}

\begin{tcolorbox}[enhanced jigsaw, arc=.35mm, bottomrule=.15mm, bottomtitle=1mm, breakable, colback=white, colbacktitle=quarto-callout-tip-color!10!white, colframe=quarto-callout-tip-color-frame, coltitle=black, left=2mm, leftrule=.75mm, opacityback=0, opacitybacktitle=0.6, rightrule=.15mm, title=\textcolor{quarto-callout-tip-color}{\faLightbulb}\hspace{0.5em}{Solution}, titlerule=0mm, toprule=.15mm, toptitle=1mm]

The batch is acceptable only if both standards are met:
\[H_1:\theta_1>50 \quad \text{and}\quad \theta_2>0.95.\] Thus the null
hypothesis is that at least one standard fails:
\[H_0:\theta_1\le 50 \quad \text{or}\quad \theta_2\le 0.95.\]
Equivalently, \[H_0=\{\theta_1\le 50\}\cup \{\theta_2\le 0.95\}\] and
\[H_1=\{\theta_1>50\}\cap \{\theta_2>0.95\}.\] We collect two types of
data:

\begin{itemize}
\item
  breaking strengths \(X_1,\ldots,X_n\), often modeled as normal with
  mean \(\theta_1\);
\item
  flammability indicators \(Y_1,\ldots,Y_m\), where \(Y_i=1\) if the
  item passes and \(0\) otherwise, often modeled as Bernoulli with
  probability \(\theta_2\).
\end{itemize}

We test the two component null hypotheses \[H_{01}:\theta_1\le 50
\qquad \text{and}\qquad
H_{02}:\theta_2\le 0.95.\] The intersection-union rule rejects the
overall null \(H_0\) only if both component null hypotheses are
rejected. In words, we declare the batch acceptable only if the strength
requirement and the flammability requirement both pass their respective
tests.

\end{tcolorbox}

\bookmarksetup{startatroot}

\chapter{Neyman-Pearson Lemma: Simple
Hypotheses}\label{neyman-pearson-lemma-simple-hypotheses}

This section records the fundamental optimality result behind likelihood
ratio tests for simple hypotheses.

Although the detailed proof is usually saved for a more advanced
treatment, the main message is essential: for simple null and simple
alternative hypotheses, the most powerful test at a fixed significance
level is a likelihood ratio test.

\begin{tcolorbox}[enhanced jigsaw, arc=.35mm, bottomrule=.15mm, bottomtitle=1mm, breakable, colback=white, colbacktitle=quarto-callout-important-color!10!white, colframe=quarto-callout-important-color-frame, coltitle=black, left=2mm, leftrule=.75mm, opacityback=0, opacitybacktitle=0.6, rightrule=.15mm, title=\textcolor{quarto-callout-important-color}{\faExclamation}\hspace{0.5em}{Theorem}, titlerule=0mm, toprule=.15mm, toptitle=1mm]

\textbf{Theorem 25} (Neyman-Pearson lemma). \emph{Consider testing two
simple hypotheses \[H_0:\theta=\theta_0
\qquad \text{versus}\qquad
H_1:\theta=\theta_1.\] Among all tests with Type I error probability at
most \(\alpha\), the most powerful test rejects \(H_0\) for sufficiently
small values of \[\frac{L(\theta_0\mid X)}{L(\theta_1\mid X)}.\]
Equivalently, it rejects for sufficiently large values of
\[\frac{L(\theta_1\mid X)}{L(\theta_0\mid X)}.\]}

\end{tcolorbox}

\begin{tcolorbox}[enhanced jigsaw, arc=.35mm, bottomrule=.15mm, bottomtitle=1mm, breakable, colback=white, colbacktitle=quarto-callout-important-color!10!white, colframe=quarto-callout-important-color-frame, coltitle=black, left=2mm, leftrule=.75mm, opacityback=0, opacitybacktitle=0.6, rightrule=.15mm, title=\textcolor{quarto-callout-important-color}{\faExclamation}\hspace{0.5em}{Remark}, titlerule=0mm, toprule=.15mm, toptitle=1mm]

\emph{Remark 26}. The Neyman-Pearson lemma explains why likelihood ratio
tests are not only natural, but optimal, for simple-versus-simple
testing problems.

\end{tcolorbox}

\bookmarksetup{startatroot}

\chapter{Practice Problems}\label{practice-problems-5}

This section gives practice problems that reinforce the main ideas of
hypothesis tests, likelihood ratios, \(p\)-values, and combined tests.

\begin{tcolorbox}[enhanced jigsaw, arc=.35mm, bottomrule=.15mm, bottomtitle=1mm, breakable, colback=white, colbacktitle=quarto-callout-warning-color!10!white, colframe=quarto-callout-warning-color-frame, coltitle=black, left=2mm, leftrule=.75mm, opacityback=0, opacitybacktitle=0.6, rightrule=.15mm, title=\textcolor{quarto-callout-warning-color}{\faExclamationTriangle}\hspace{0.5em}{Practice Problem}, titlerule=0mm, toprule=.15mm, toptitle=1mm]

\textbf{Practice Problem 27} (Coin test). A coin is tossed \(100\) times
and \(62\) heads are observed. Test \[H_0:p=\frac12
\qquad \text{versus}\qquad
H_1:p\ne \frac12\] using the normal approximation at level
\(\alpha=0.05\).

\end{tcolorbox}

\begin{tcolorbox}[enhanced jigsaw, arc=.35mm, bottomrule=.15mm, bottomtitle=1mm, breakable, colback=white, colbacktitle=quarto-callout-tip-color!10!white, colframe=quarto-callout-tip-color-frame, coltitle=black, left=2mm, leftrule=.75mm, opacityback=0, opacitybacktitle=0.6, rightrule=.15mm, title=\textcolor{quarto-callout-tip-color}{\faLightbulb}\hspace{0.5em}{Solution}, titlerule=0mm, toprule=.15mm, toptitle=1mm]

Under \(H_0\), \[X\sim \operatorname{Binomial}\left(100,\frac12\right),
\qquad
\mathbb{E}[X]=50,
\qquad
\operatorname{sd}(X)=5.\] The standardized statistic is
\[Z=\frac{62-50}{5}=2.4.\] For a two-sided level-\(0.05\) test, the
critical value is \(1.96\). Since \[|2.4|>1.96,\] we reject \(H_0\). The
approximate two-sided \(p\)-value is \[2(1-\Phi(2.4))\approx 0.0164.\]

\end{tcolorbox}

\begin{tcolorbox}[enhanced jigsaw, arc=.35mm, bottomrule=.15mm, bottomtitle=1mm, breakable, colback=white, colbacktitle=quarto-callout-warning-color!10!white, colframe=quarto-callout-warning-color-frame, coltitle=black, left=2mm, leftrule=.75mm, opacityback=0, opacitybacktitle=0.6, rightrule=.15mm, title=\textcolor{quarto-callout-warning-color}{\faExclamationTriangle}\hspace{0.5em}{Practice Problem}, titlerule=0mm, toprule=.15mm, toptitle=1mm]

\textbf{Practice Problem 28} (Radar \(p\)-value). In the radar example,
suppose \(X=0.4\). Compute the right-sided \(p\)-value under \(H_0\).

\end{tcolorbox}

\begin{tcolorbox}[enhanced jigsaw, arc=.35mm, bottomrule=.15mm, bottomtitle=1mm, breakable, colback=white, colbacktitle=quarto-callout-tip-color!10!white, colframe=quarto-callout-tip-color-frame, coltitle=black, left=2mm, leftrule=.75mm, opacityback=0, opacitybacktitle=0.6, rightrule=.15mm, title=\textcolor{quarto-callout-tip-color}{\faLightbulb}\hspace{0.5em}{Solution}, titlerule=0mm, toprule=.15mm, toptitle=1mm]

Under \(H_0\), \(X\sim \operatorname{Normal}(0,1/9)\). Standardizing
gives \[Z=3X=3(0.4)=1.2.\] Thus the right-sided \(p\)-value is
\[p=\mathbb{P}(Z\ge 1.2)=1-\Phi(1.2)\approx 0.1151.\] At level \(0.05\),
we would not reject \(H_0\).

\end{tcolorbox}

\begin{tcolorbox}[enhanced jigsaw, arc=.35mm, bottomrule=.15mm, bottomtitle=1mm, breakable, colback=white, colbacktitle=quarto-callout-warning-color!10!white, colframe=quarto-callout-warning-color-frame, coltitle=black, left=2mm, leftrule=.75mm, opacityback=0, opacitybacktitle=0.6, rightrule=.15mm, title=\textcolor{quarto-callout-warning-color}{\faExclamationTriangle}\hspace{0.5em}{Practice Problem}, titlerule=0mm, toprule=.15mm, toptitle=1mm]

\textbf{Practice Problem 29} (Normal mean, known variance). Suppose
\(X_1,\ldots,X_{25}\sim \operatorname{Normal}(\mu,16)\) and
\(\bar x=53\). Test \[H_0:\mu=50
\qquad \text{versus}\qquad
H_1:\mu\ne 50\] at level \(0.05\).

\end{tcolorbox}

\begin{tcolorbox}[enhanced jigsaw, arc=.35mm, bottomrule=.15mm, bottomtitle=1mm, breakable, colback=white, colbacktitle=quarto-callout-tip-color!10!white, colframe=quarto-callout-tip-color-frame, coltitle=black, left=2mm, leftrule=.75mm, opacityback=0, opacitybacktitle=0.6, rightrule=.15mm, title=\textcolor{quarto-callout-tip-color}{\faLightbulb}\hspace{0.5em}{Solution}, titlerule=0mm, toprule=.15mm, toptitle=1mm]

Here \(\sigma=4\) and \(n=25\), so
\[Z=\frac{\bar x-\mu_0}{\sigma/\sqrt n}
=\frac{53-50}{4/5}=\frac{3}{0.8}=3.75.\] Since \[|3.75|>1.96,\] we
reject \(H_0\) at level \(0.05\).

\end{tcolorbox}

\begin{tcolorbox}[enhanced jigsaw, arc=.35mm, bottomrule=.15mm, bottomtitle=1mm, breakable, colback=white, colbacktitle=quarto-callout-warning-color!10!white, colframe=quarto-callout-warning-color-frame, coltitle=black, left=2mm, leftrule=.75mm, opacityback=0, opacitybacktitle=0.6, rightrule=.15mm, title=\textcolor{quarto-callout-warning-color}{\faExclamationTriangle}\hspace{0.5em}{Practice Problem}, titlerule=0mm, toprule=.15mm, toptitle=1mm]

\textbf{Practice Problem 30} (Normal mean, unknown variance). Suppose
\(X_1,\ldots,X_{16}\) are normal, \(\bar x=10.8\), and \(s=2.4\). Test
\[H_0:\mu=10
\qquad \text{versus}\qquad
H_1:\mu\ne 10\] at level \(0.05\).

\end{tcolorbox}

\begin{tcolorbox}[enhanced jigsaw, arc=.35mm, bottomrule=.15mm, bottomtitle=1mm, breakable, colback=white, colbacktitle=quarto-callout-tip-color!10!white, colframe=quarto-callout-tip-color-frame, coltitle=black, left=2mm, leftrule=.75mm, opacityback=0, opacitybacktitle=0.6, rightrule=.15mm, title=\textcolor{quarto-callout-tip-color}{\faLightbulb}\hspace{0.5em}{Solution}, titlerule=0mm, toprule=.15mm, toptitle=1mm]

The test statistic is \[t=\frac{\bar x-\mu_0}{s/\sqrt n}
=\frac{10.8-10}{2.4/4}
=\frac{0.8}{0.6}=1.333.\] There are \(n-1=15\) degrees of freedom. The
two-sided \(0.05\) critical value is about
\[t_{15,0.975}\approx 2.131.\] Since \[|1.333|<2.131,\] we fail to
reject \(H_0\).

\end{tcolorbox}

\begin{tcolorbox}[enhanced jigsaw, arc=.35mm, bottomrule=.15mm, bottomtitle=1mm, breakable, colback=white, colbacktitle=quarto-callout-warning-color!10!white, colframe=quarto-callout-warning-color-frame, coltitle=black, left=2mm, leftrule=.75mm, opacityback=0, opacitybacktitle=0.6, rightrule=.15mm, title=\textcolor{quarto-callout-warning-color}{\faExclamationTriangle}\hspace{0.5em}{Practice Problem}, titlerule=0mm, toprule=.15mm, toptitle=1mm]

\textbf{Practice Problem 31} (Likelihood ratio for Bernoulli simple
hypotheses). Let \(X_1,\ldots,X_n\sim \operatorname{Bernoulli}(p)\).
Test \[H_0:p=p_0
\qquad \text{versus}\qquad
H_1:p=p_1,\] where \(p_1>p_0\). Find the likelihood ratio and describe
the rejection region.

\end{tcolorbox}

\begin{tcolorbox}[enhanced jigsaw, arc=.35mm, bottomrule=.15mm, bottomtitle=1mm, breakable, colback=white, colbacktitle=quarto-callout-tip-color!10!white, colframe=quarto-callout-tip-color-frame, coltitle=black, left=2mm, leftrule=.75mm, opacityback=0, opacitybacktitle=0.6, rightrule=.15mm, title=\textcolor{quarto-callout-tip-color}{\faLightbulb}\hspace{0.5em}{Solution}, titlerule=0mm, toprule=.15mm, toptitle=1mm]

Let \[S=\sum_{i=1}^n X_i.\] The likelihood is
\[L(p\mid x)=p^S(1-p)^{n-S}.\] Thus
\[\lambda(x)=\frac{L(p_0\mid x)}{L(p_1\mid x)}
=\left(\frac{p_0}{p_1}\right)^S
\left(\frac{1-p_0}{1-p_1}\right)^{n-S}.\] Taking logs,
\[\log\lambda(x)=S\log\left(\frac{p_0}{p_1}\right)
+(n-S)\log\left(\frac{1-p_0}{1-p_1}\right).\] Because \(p_1>p_0\), the
likelihood ratio decreases as \(S\) increases. Therefore the LRT rejects
\(H_0\) for large values of \(S\), i.e.~\[\sum_{i=1}^n X_i \ge k\] for a
threshold \(k\) chosen to control the Type I error probability.

\end{tcolorbox}

\begin{tcolorbox}[enhanced jigsaw, arc=.35mm, bottomrule=.15mm, bottomtitle=1mm, breakable, colback=white, colbacktitle=quarto-callout-warning-color!10!white, colframe=quarto-callout-warning-color-frame, coltitle=black, left=2mm, leftrule=.75mm, opacityback=0, opacitybacktitle=0.6, rightrule=.15mm, title=\textcolor{quarto-callout-warning-color}{\faExclamationTriangle}\hspace{0.5em}{Practice Problem}, titlerule=0mm, toprule=.15mm, toptitle=1mm]

\textbf{Practice Problem 32} (Bayesian posterior decision). Suppose
\(\mu\mid x\sim \operatorname{Normal}(4,1)\). Test \[H_0:\mu\le 3
\qquad \text{versus}\qquad
H_1:\mu>3.\] Using the Bayesian rule ``reject \(H_0\) if
\(\mathbb{P}(H_0\mid x)<1/2\),'' what is the decision?

\end{tcolorbox}

\begin{tcolorbox}[enhanced jigsaw, arc=.35mm, bottomrule=.15mm, bottomtitle=1mm, breakable, colback=white, colbacktitle=quarto-callout-tip-color!10!white, colframe=quarto-callout-tip-color-frame, coltitle=black, left=2mm, leftrule=.75mm, opacityback=0, opacitybacktitle=0.6, rightrule=.15mm, title=\textcolor{quarto-callout-tip-color}{\faLightbulb}\hspace{0.5em}{Solution}, titlerule=0mm, toprule=.15mm, toptitle=1mm]

Since the posterior is normal with mean \(4\) and variance \(1\),
\[\mathbb{P}(H_0\mid x)=\mathbb{P}(\mu\le 3\mid x)=\Phi\left(\frac{3-4}{1}\right)=\Phi(-1)\approx 0.1587.\]
Since \[0.1587<\frac12,\] we reject \(H_0\) using this Bayesian decision
rule.

\end{tcolorbox}

\begin{tcolorbox}[enhanced jigsaw, arc=.35mm, bottomrule=.15mm, bottomtitle=1mm, breakable, colback=white, colbacktitle=quarto-callout-warning-color!10!white, colframe=quarto-callout-warning-color-frame, coltitle=black, left=2mm, leftrule=.75mm, opacityback=0, opacitybacktitle=0.6, rightrule=.15mm, title=\textcolor{quarto-callout-warning-color}{\faExclamationTriangle}\hspace{0.5em}{Practice Problem}, titlerule=0mm, toprule=.15mm, toptitle=1mm]

\textbf{Practice Problem 33} (Intersection-union logic). A device is
acceptable only if its battery life is above \(10\) hours and its
failure probability is below \(0.01\). Write the null and alternative
hypotheses for an IUT.

\end{tcolorbox}

\begin{tcolorbox}[enhanced jigsaw, arc=.35mm, bottomrule=.15mm, bottomtitle=1mm, breakable, colback=white, colbacktitle=quarto-callout-tip-color!10!white, colframe=quarto-callout-tip-color-frame, coltitle=black, left=2mm, leftrule=.75mm, opacityback=0, opacitybacktitle=0.6, rightrule=.15mm, title=\textcolor{quarto-callout-tip-color}{\faLightbulb}\hspace{0.5em}{Solution}, titlerule=0mm, toprule=.15mm, toptitle=1mm]

Let \(\theta_1\) be the mean battery life and \(\theta_2\) be the
failure probability. The device is acceptable only if \[\theta_1>10
\qquad \text{and}\qquad
\theta_2<0.01.\] Thus the alternative is
\[H_1:\theta_1>10 \text{ and } \theta_2<0.01.\] The null is that at
least one requirement fails:
\[H_0:\theta_1\le 10 \text{ or } \theta_2\ge 0.01.\] An IUT rejects
\(H_0\) only if both component tests reject their corresponding
component null hypotheses.

\end{tcolorbox}

\bookmarksetup{startatroot}

\chapter{Summary}\label{summary-7}

This section introduced the foundations and construction methods for
hypothesis testing.

\begin{tcolorbox}[enhanced jigsaw, arc=.35mm, bottomrule=.15mm, bottomtitle=1mm, breakable, colback=white, colbacktitle=quarto-callout-tip-color!10!white, colframe=quarto-callout-tip-color-frame, coltitle=black, left=2mm, leftrule=.75mm, opacityback=0, opacitybacktitle=0.6, rightrule=.15mm, title=\textcolor{quarto-callout-tip-color}{\faLightbulb}\hspace{0.5em}{Key idea}, titlerule=0mm, toprule=.15mm, toptitle=1mm]

Key takeaways

\begin{itemize}
\item
  A hypothesis test divides the sample space into an acceptance region
  and a rejection region.
\item
  Type I error means rejecting a true null; Type II error means failing
  to reject a false null.
\item
  The significance level controls Type I error probability.
\item
  The \(p\)-value is the smallest significance level that would reject
  \(H_0\).
\item
  Likelihood ratio tests reject when the null model fits much worse than
  the unrestricted model.
\item
  The classical \(Z\)-test and \(t\)-test can be derived from likelihood
  ratio tests.
\item
  Bayesian tests use posterior probabilities of hypotheses.
\item
  UITs reject when any component test rejects; IUTs reject only when all
  component tests reject.
\end{itemize}

\end{tcolorbox}

\bookmarksetup{startatroot}

\chapter{Chapter 15: Hypothesis Tests II --- Evaluating
Tests}\label{chapter-15-hypothesis-tests-ii-evaluating-tests}

This chapter explains how to evaluate and compare hypothesis tests after
constructing them. The main ideas are Type I and Type II errors, power
functions, test size and level, unbiased tests, uniformly most powerful
tests, \(p\)-values, and risk functions.

\begin{tcolorbox}[enhanced jigsaw, arc=.35mm, bottomrule=.15mm, bottomtitle=1mm, breakable, colback=white, colbacktitle=quarto-callout-note-color!10!white, colframe=quarto-callout-note-color-frame, coltitle=black, left=2mm, leftrule=.75mm, opacityback=0, opacitybacktitle=0.6, rightrule=.15mm, title=\textcolor{quarto-callout-note-color}{\faInfo}\hspace{0.5em}{Topics}, titlerule=0mm, toprule=.15mm, toptitle=1mm]

Evaluating hypothesis tests; Type I and Type II errors; power function;
size and level; unbiased tests; uniformly most powerful tests;
Neyman--Pearson lemma; binomial and normal UMP tests; \(p\)-values; loss
functions; risk functions; decision-theoretic evaluation of tests.

\end{tcolorbox}

\bookmarksetup{startatroot}

\chapter{Why We Need to Evaluate
Tests}\label{why-we-need-to-evaluate-tests}

This section explains how to compare hypothesis tests after we have
learned how to construct them.

In hypothesis testing, different methods may lead to different tests. A
likelihood ratio test, a Bayesian test, a union-intersection test, and a
classical test statistic can all be reasonable procedures. The key
question is: \emph{how do we decide whether one test is better than
another?}

\begin{tcolorbox}[enhanced jigsaw, arc=.35mm, bottomrule=.15mm, bottomtitle=1mm, breakable, colback=white, colbacktitle=quarto-callout-tip-color!10!white, colframe=quarto-callout-tip-color-frame, coltitle=black, left=2mm, leftrule=.75mm, opacityback=0, opacitybacktitle=0.6, rightrule=.15mm, title=\textcolor{quarto-callout-tip-color}{\faLightbulb}\hspace{0.5em}{Key idea}, titlerule=0mm, toprule=.15mm, toptitle=1mm]

Main idea Hypothesis tests are evaluated and compared through their
probabilities of making mistakes. A good test should have small Type I
error probability and small Type II error probability, subject to the
unavoidable trade-off between the two.

\end{tcolorbox}

The main evaluation tools in this section are:

\begin{itemize}
\item
  the \textbf{power function};
\item
  the \textbf{size} or \textbf{significance level};
\item
  \textbf{unbiasedness} of tests;
\item
  \textbf{uniformly most powerful} tests;
\item
  \textbf{risk functions} and expected loss.
\end{itemize}

\bookmarksetup{startatroot}

\chapter{Two Types of Errors}\label{two-types-of-errors}

This section reviews the two fundamental errors in hypothesis testing
and introduces the power function as a unified way to measure them.

Consider a hypothesis test \[H_0: \theta\in \Theta_0
\qquad \text{versus} \qquad
H_1: \theta\in \Theta_0^c.\] Let \(R\) be the rejection region of the
test. The test rejects \(H_0\) when \(X\in R\) and fails to reject
\(H_0\) when \(X\in R^c\).

\section{Type I and Type II errors}\label{type-i-and-type-ii-errors-1}

This subsection defines the two ways a hypothesis test can be wrong.

\begin{tcolorbox}[enhanced jigsaw, arc=.35mm, bottomrule=.15mm, bottomtitle=1mm, breakable, colback=white, colbacktitle=quarto-callout-note-color!10!white, colframe=quarto-callout-note-color-frame, coltitle=black, left=2mm, leftrule=.75mm, opacityback=0, opacitybacktitle=0.6, rightrule=.15mm, title=\textcolor{quarto-callout-note-color}{\faInfo}\hspace{0.5em}{Definition}, titlerule=0mm, toprule=.15mm, toptitle=1mm]

\textbf{Definition 1} (Type I and Type II errors). For the test
\(H_0:\theta\in\Theta_0\) versus \(H_1:\theta\in\Theta_0^c\):

\begin{itemize}
\item
  A \textbf{Type I error} occurs when \(\theta\in\Theta_0\) but the test
  rejects \(H_0\).
\item
  A \textbf{Type II error} occurs when \(\theta\in\Theta_0^c\) but the
  test fails to reject \(H_0\).
\end{itemize}

\end{tcolorbox}

{\def\LTcaptype{none} % do not increment counter
\begin{longtable}[]{@{}lll@{}}
\toprule\noalign{}
Decision & \(H_0\) true & \(H_1\) true \\
\midrule\noalign{}
\endhead
\bottomrule\noalign{}
\endlastfoot
Reject \(H_0\) & Type I error & Correct decision \\
Fail to reject \(H_0\) & Correct decision & Type II error \\
\end{longtable}
}

For \(\theta\in\Theta_0\), the Type I error probability is
\[\mathbb{P}_\theta(X\in R).\] For \(\theta\in\Theta_0^c\), the Type II
error probability is
\[\mathbb{P}_\theta(X\in R^c)=1-\mathbb{P}_\theta(X\in R).\]

\section{Power function}\label{power-function}

This subsection introduces the main function used to evaluate a test
over all parameter values.

\begin{tcolorbox}[enhanced jigsaw, arc=.35mm, bottomrule=.15mm, bottomtitle=1mm, breakable, colback=white, colbacktitle=quarto-callout-note-color!10!white, colframe=quarto-callout-note-color-frame, coltitle=black, left=2mm, leftrule=.75mm, opacityback=0, opacitybacktitle=0.6, rightrule=.15mm, title=\textcolor{quarto-callout-note-color}{\faInfo}\hspace{0.5em}{Definition}, titlerule=0mm, toprule=.15mm, toptitle=1mm]

\textbf{Definition 2} (Power function). The \textbf{power function} of a
hypothesis test with rejection region \(R\) is
\[\beta(\theta):=\mathbb{P}_\theta(X\in R).\] Thus, \[\beta(\theta)=
\begin{cases}
\mathbb{P}_\theta(\text{Type I error}), & \theta\in\Theta_0,\\[3pt]
1-\mathbb{P}_\theta(\text{Type II error}), & \theta\in\Theta_0^c.
\end{cases}\]

\end{tcolorbox}

\begin{tcolorbox}[enhanced jigsaw, arc=.35mm, bottomrule=.15mm, bottomtitle=1mm, breakable, colback=white, colbacktitle=quarto-callout-tip-color!10!white, colframe=quarto-callout-tip-color-frame, coltitle=black, left=2mm, leftrule=.75mm, opacityback=0, opacitybacktitle=0.6, rightrule=.15mm, title=\textcolor{quarto-callout-tip-color}{\faLightbulb}\hspace{0.5em}{Key idea}, titlerule=0mm, toprule=.15mm, toptitle=1mm]

Ideal power function The ideal test would have \[\beta(\theta)=
\begin{cases}
0, & \theta\in\Theta_0,\\
1, & \theta\in\Theta_0^c.
\end{cases}\] This would mean no Type I errors and no Type II errors. In
real problems, such a perfect test is usually impossible.

\end{tcolorbox}

\bookmarksetup{startatroot}

\chapter{Power Function Examples}\label{power-function-examples}

This section computes power functions for binomial and normal tests and
shows the trade-off between Type I and Type II errors.

\section{Binomial power functions}\label{binomial-power-functions}

This subsection compares two tests for the same binomial hypothesis.

\begin{tcolorbox}[enhanced jigsaw, arc=.35mm, bottomrule=.15mm, bottomtitle=1mm, breakable, colback=white, colbacktitle=quarto-callout-tip-color!10!white, colframe=quarto-callout-tip-color-frame, coltitle=black, left=2mm, leftrule=.75mm, opacityback=0, opacitybacktitle=0.6, rightrule=.15mm, title=\textcolor{quarto-callout-tip-color}{\faLightbulb}\hspace{0.5em}{Example}, titlerule=0mm, toprule=.15mm, toptitle=1mm]

\textbf{Example 3} (Binomial power function: two tests). Let
\[X\sim \operatorname{Binomial}(n=5,\theta),\] and consider
\[H_0:\theta\le \frac12
\qquad \text{versus} \qquad
H_1:\theta>\frac12.\] Compare the following two tests.

\textbf{Test 1:} reject \(H_0\) if \(X=5\).

The rejection region is \(R_1=\{5\}\), so the power function is
\[\beta_1(\theta)=\mathbb{P}_\theta(X=5)=\theta^5.\] For
\(\theta\le 1/2\),
\[\beta_1(\theta)\le \left(\frac12\right)^5=\frac{1}{32}\approx 0.03125.\]
Thus Test 1 has a small Type I error probability, but for values of
\(\theta\) only moderately larger than \(1/2\), it has a large Type II
error probability.

\textbf{Test 2:} reject \(H_0\) if \(X=3,4,\) or \(5\).

The rejection region is \(R_2=\{3,4,5\}\), so the power function is
\[\beta_2(\theta)=\mathbb{P}_\theta(X=3,4,5)
=\binom53\theta^3(1-\theta)^2+\binom54\theta^4(1-\theta)+\binom55\theta^5.\]

\end{tcolorbox}

\begin{tcolorbox}[enhanced jigsaw, arc=.35mm, bottomrule=.15mm, bottomtitle=1mm, breakable, colback=white, colbacktitle=quarto-callout-note-color!10!white, colframe=quarto-callout-note-color-frame, coltitle=black, left=2mm, leftrule=.75mm, opacityback=0, opacitybacktitle=0.6, rightrule=.15mm, title=\textcolor{quarto-callout-note-color}{\faInfo}\hspace{0.5em}{Solution}, titlerule=0mm, toprule=.15mm, toptitle=1mm]

The first test is conservative: it rejects only when all five trials are
successes. Its Type I error probability is at most \(1/32\) under
\(H_0\).

The second test rejects more often, so it has larger power for
\(\theta>1/2\). However, it also has larger Type I error probability. At
the boundary \(\theta=1/2\), \[\beta_2\left(\frac12\right)
=\binom53\left(\frac12\right)^5+\binom54\left(\frac12\right)^5+\binom55\left(\frac12\right)^5
=\frac{10+5+1}{32}=\frac12.\] Thus Test 2 has much larger power but also
much larger size. This illustrates the Type I--Type II error trade-off.

\end{tcolorbox}

\begin{figure}

\centering{

\pandocbounded{\includegraphics[keepaspectratio]{chapters/chapter-15-hypothesis-tests2_files/figure-pdf/fig-binomial-power-comparison-output-1.pdf}}

}

\caption{\label{fig-binomial-power-comparison}Power functions for two
binomial tests with n = 5.}

\end{figure}%

\section{Normal power function}\label{normal-power-function}

This subsection derives the power function of a one-sided normal mean
test.

\begin{tcolorbox}[enhanced jigsaw, arc=.35mm, bottomrule=.15mm, bottomtitle=1mm, breakable, colback=white, colbacktitle=quarto-callout-tip-color!10!white, colframe=quarto-callout-tip-color-frame, coltitle=black, left=2mm, leftrule=.75mm, opacityback=0, opacitybacktitle=0.6, rightrule=.15mm, title=\textcolor{quarto-callout-tip-color}{\faLightbulb}\hspace{0.5em}{Example}, titlerule=0mm, toprule=.15mm, toptitle=1mm]

\textbf{Example 4} (Normal power function). Let
\[X_1,\ldots,X_n\sim \operatorname{Normal}(\theta,\sigma^2),\] where
\(\sigma^2\) is known. Test \[H_0:\theta\le \theta_0
\qquad \text{versus} \qquad
H_1:\theta>\theta_0.\] Consider the test that rejects \(H_0\) if
\[\frac{\bar X-\theta_0}{\sigma/\sqrt n}>c,\] where \(c>0\) is a
constant.

\end{tcolorbox}

\begin{tcolorbox}[enhanced jigsaw, arc=.35mm, bottomrule=.15mm, bottomtitle=1mm, breakable, colback=white, colbacktitle=quarto-callout-note-color!10!white, colframe=quarto-callout-note-color-frame, coltitle=black, left=2mm, leftrule=.75mm, opacityback=0, opacitybacktitle=0.6, rightrule=.15mm, title=\textcolor{quarto-callout-note-color}{\faInfo}\hspace{0.5em}{Solution}, titlerule=0mm, toprule=.15mm, toptitle=1mm]

For any fixed \(\theta\),
\[Z=\frac{\bar X-\theta}{\sigma/\sqrt n}\sim N(0,1).\] Thus the power
function is \[\begin{aligned}
\beta(\theta)
&=\mathbb{P}_\theta\left(\frac{\bar X-\theta_0}{\sigma/\sqrt n}>c\right)\\
&=\mathbb{P}_\theta\left(\frac{\bar X-\theta}{\sigma/\sqrt n}>c+\frac{\theta_0-\theta}{\sigma/\sqrt n}\right)\\
&=1-\Phi\left(c+\frac{\theta_0-\theta}{\sigma/\sqrt n}\right).
\end{aligned}\] Therefore, \[\beta(\theta_0)=1-\Phi(c).\] If \(c=1.28\),
then \(\mathbb{P}(Z>1.28)\approx 0.10\), so the test has about a
\(10\%\) significance level at the boundary \(\theta=\theta_0\).

\end{tcolorbox}

The power function has the following behavior:
\[\theta\to -\infty \implies \beta(\theta)\to 0,
\qquad
\theta\to +\infty \implies \beta(\theta)\to 1.\] For fixed sample size,
decreasing Type I error usually increases Type II error, and increasing
power usually increases Type I error.

\section{Choosing sample size from power
requirements}\label{choosing-sample-size-from-power-requirements}

This subsection shows how power calculations determine a sample size.

\begin{tcolorbox}[enhanced jigsaw, arc=.35mm, bottomrule=.15mm, bottomtitle=1mm, breakable, colback=white, colbacktitle=quarto-callout-tip-color!10!white, colframe=quarto-callout-tip-color-frame, coltitle=black, left=2mm, leftrule=.75mm, opacityback=0, opacitybacktitle=0.6, rightrule=.15mm, title=\textcolor{quarto-callout-tip-color}{\faLightbulb}\hspace{0.5em}{Example}, titlerule=0mm, toprule=.15mm, toptitle=1mm]

\textbf{Example 5} (Choosing \(c\) and \(n\)). Let
\[X_1,\ldots,X_n\sim \operatorname{Normal}(\theta,\sigma^2),\] with
known \(\sigma^2\). Test \[H_0:\theta\le \theta_0
\qquad \text{versus} \qquad
H_1:\theta>\theta_0\] using the rule
\[\frac{\bar X-\theta_0}{\sigma/\sqrt n}>c.\] Choose \(c\) and \(n\) so
that \[\mathbb{P}(\text{Type I error})\le 0.10,
\qquad
\mathbb{P}(\text{Type II error})\le 0.20\quad \text{if } \theta\ge \theta_0+\sigma.\]

\end{tcolorbox}

\begin{tcolorbox}[enhanced jigsaw, arc=.35mm, bottomrule=.15mm, bottomtitle=1mm, breakable, colback=white, colbacktitle=quarto-callout-note-color!10!white, colframe=quarto-callout-note-color-frame, coltitle=black, left=2mm, leftrule=.75mm, opacityback=0, opacitybacktitle=0.6, rightrule=.15mm, title=\textcolor{quarto-callout-note-color}{\faInfo}\hspace{0.5em}{Solution}, titlerule=0mm, toprule=.15mm, toptitle=1mm]

The Type I error is largest at the boundary \(\theta=\theta_0\). Hence
we require \[\beta(\theta_0)=\mathbb{P}(Z>c)=0.10.\] Thus
\[c=z_{0.10}\approx 1.28.\] Next, requiring Type II error at most
\(0.20\) when \(\theta=\theta_0+\sigma\) is equivalent to requiring
power at least \(0.80\): \[\beta(\theta_0+\sigma)=0.80.\] But
\[\beta(\theta_0+\sigma)
=\mathbb{P}\left(Z>c-\sqrt n\right).\] So
\[\mathbb{P}\left(Z>1.28-\sqrt n\right)=0.80.\] Equivalently,
\[1.28-\sqrt n=z_{0.80}^{\text{right tail}},\] where
\(\mathbb{P}(Z>z_{0.80}^{\text{right tail}})=0.80\). Since
\(z_{0.80}^{\text{right tail}}\approx -0.84\),
\[1.28-\sqrt n\approx -0.84,
\qquad
\sqrt n\approx 2.12,
\qquad
n\approx 4.49.\] Therefore, choose \[\boxed{n=5.}\]

\end{tcolorbox}

\bookmarksetup{startatroot}

\chapter{Size and Level of a Test}\label{size-and-level-of-a-test}

This section distinguishes the exact size of a test from the more
flexible idea of a level-\(\alpha\) test.

\section{Size and level}\label{size-and-level}

This subsection gives the formal definitions.

\begin{tcolorbox}[enhanced jigsaw, arc=.35mm, bottomrule=.15mm, bottomtitle=1mm, breakable, colback=white, colbacktitle=quarto-callout-note-color!10!white, colframe=quarto-callout-note-color-frame, coltitle=black, left=2mm, leftrule=.75mm, opacityback=0, opacitybacktitle=0.6, rightrule=.15mm, title=\textcolor{quarto-callout-note-color}{\faInfo}\hspace{0.5em}{Definition}, titlerule=0mm, toprule=.15mm, toptitle=1mm]

\textbf{Definition 6} (Size-\(\alpha\) test). For \(\alpha\in[0,1]\), a
test with power function \(\beta(\theta)\) is called a \textbf{size
\(\alpha\) test} if
\[\sup_{\theta\in\Theta_0}\mathbb{P}_\theta(\text{reject }H_0)
=\sup_{\theta\in\Theta_0}\beta(\theta)=\alpha.\]

\end{tcolorbox}

\begin{tcolorbox}[enhanced jigsaw, arc=.35mm, bottomrule=.15mm, bottomtitle=1mm, breakable, colback=white, colbacktitle=quarto-callout-note-color!10!white, colframe=quarto-callout-note-color-frame, coltitle=black, left=2mm, leftrule=.75mm, opacityback=0, opacitybacktitle=0.6, rightrule=.15mm, title=\textcolor{quarto-callout-note-color}{\faInfo}\hspace{0.5em}{Definition}, titlerule=0mm, toprule=.15mm, toptitle=1mm]

\textbf{Definition 7} (Level-\(\alpha\) test). A test is called a
\textbf{level \(\alpha\) test} if
\[\sup_{\theta\in\Theta_0}\beta(\theta)\le \alpha.\]

\end{tcolorbox}

\begin{tcolorbox}[enhanced jigsaw, arc=.35mm, bottomrule=.15mm, bottomtitle=1mm, breakable, colback=white, colbacktitle=quarto-callout-important-color!10!white, colframe=quarto-callout-important-color-frame, coltitle=black, left=2mm, leftrule=.75mm, opacityback=0, opacitybacktitle=0.6, rightrule=.15mm, title=\textcolor{quarto-callout-important-color}{\faExclamation}\hspace{0.5em}{Remark}, titlerule=0mm, toprule=.15mm, toptitle=1mm]

\emph{Remark 8}. A size-\(\alpha\) test uses the full Type I error
budget exactly. A level-\(\alpha\) test may be conservative and use less
than the allowed Type I error probability.

\end{tcolorbox}

\section{Size of likelihood ratio
tests}\label{size-of-likelihood-ratio-tests}

This subsection explains how cutoffs are chosen for likelihood ratio
tests.

For a likelihood ratio test with rejection region
\[R=\{x:\lambda(x)\le c\},\] the size is
\[\sup_{\theta\in\Theta_0}\mathbb{P}_\theta(\lambda(X)\le c).\] To
obtain a size-\(\alpha\) LRT, choose \(c\) so that
\[\sup_{\theta\in\Theta_0}\mathbb{P}_\theta(\lambda(X)\le c)=\alpha.\]

\begin{tcolorbox}[enhanced jigsaw, arc=.35mm, bottomrule=.15mm, bottomtitle=1mm, breakable, colback=white, colbacktitle=quarto-callout-tip-color!10!white, colframe=quarto-callout-tip-color-frame, coltitle=black, left=2mm, leftrule=.75mm, opacityback=0, opacitybacktitle=0.6, rightrule=.15mm, title=\textcolor{quarto-callout-tip-color}{\faLightbulb}\hspace{0.5em}{Example}, titlerule=0mm, toprule=.15mm, toptitle=1mm]

\textbf{Example 9} (Size-\(\alpha\) LRT for a normal mean). Suppose
\(X_1,
\ldots,X_n\sim N(\theta,\sigma^2)\) with known \(\sigma^2\), and test
\[H_0:\theta=\theta_0
\qquad \text{versus} \qquad
H_1:\theta\ne\theta_0.\] The LRT rejects for large values of
\[\left|\frac{\bar X-\theta_0}{\sigma/\sqrt n}\right|.\] A
size-\(\alpha\) test rejects \(H_0\) if
\[\left|\frac{\bar X-\theta_0}{\sigma/\sqrt n}\right|>z_{\alpha/2},\]
where \(\mathbb{P}(Z>z_{\alpha/2})=\alpha/2\) for \(Z\sim N(0,1)\).

\end{tcolorbox}

\begin{tcolorbox}[enhanced jigsaw, arc=.35mm, bottomrule=.15mm, bottomtitle=1mm, breakable, colback=white, colbacktitle=quarto-callout-note-color!10!white, colframe=quarto-callout-note-color-frame, coltitle=black, left=2mm, leftrule=.75mm, opacityback=0, opacitybacktitle=0.6, rightrule=.15mm, title=\textcolor{quarto-callout-note-color}{\faInfo}\hspace{0.5em}{Solution}, titlerule=0mm, toprule=.15mm, toptitle=1mm]

Under \(H_0\), \[Z=\frac{\bar X-\theta_0}{\sigma/\sqrt n}\sim N(0,1).\]
Thus \[\mathbb{P}_{\theta_0}(|Z|>z_{\alpha/2})=\alpha.\] This is the
classical two-sided \(Z\)-test. In LRT form, the cutoff is equivalent to
\[c=\exp\left(-\frac{z_{\alpha/2}^2}{2}\right).\]

\end{tcolorbox}

\begin{tcolorbox}[enhanced jigsaw, arc=.35mm, bottomrule=.15mm, bottomtitle=1mm, breakable, colback=white, colbacktitle=quarto-callout-tip-color!10!white, colframe=quarto-callout-tip-color-frame, coltitle=black, left=2mm, leftrule=.75mm, opacityback=0, opacitybacktitle=0.6, rightrule=.15mm, title=\textcolor{quarto-callout-tip-color}{\faLightbulb}\hspace{0.5em}{Key idea}, titlerule=0mm, toprule=.15mm, toptitle=1mm]

Common critical values The following notation is standard:
\[\begin{aligned}
&z_\alpha: &&\mathbb{P}(Z>z_\alpha)=\alpha,
\quad Z\sim N(0,1),\\
&t_{n-1,\alpha}: &&\mathbb{P}(T_{n-1}>t_{n-1,\alpha})=\alpha,\\
&\chi^2_{p,1-\alpha}: &&\mathbb{P}(\chi_p^2>\chi^2_{p,1-\alpha})=\alpha.
\end{aligned}\]

\end{tcolorbox}

\bookmarksetup{startatroot}

\chapter{Unbiased Tests}\label{unbiased-tests}

This section adapts the idea of unbiasedness from estimation to
hypothesis testing.

\begin{tcolorbox}[enhanced jigsaw, arc=.35mm, bottomrule=.15mm, bottomtitle=1mm, breakable, colback=white, colbacktitle=quarto-callout-note-color!10!white, colframe=quarto-callout-note-color-frame, coltitle=black, left=2mm, leftrule=.75mm, opacityback=0, opacitybacktitle=0.6, rightrule=.15mm, title=\textcolor{quarto-callout-note-color}{\faInfo}\hspace{0.5em}{Definition}, titlerule=0mm, toprule=.15mm, toptitle=1mm]

\textbf{Definition 10} (Unbiased test). A test with power function
\(\beta(\theta)\) is called \textbf{unbiased} if
\[\beta(\theta')\ge \beta(\theta'')\] for every \(\theta'\in\Theta_0^c\)
and every \(\theta''\in\Theta_0\).

\end{tcolorbox}

This means that the probability of rejecting \(H_0\) under the
alternative should be at least as large as the probability of rejecting
\(H_0\) under the null.

\begin{tcolorbox}[enhanced jigsaw, arc=.35mm, bottomrule=.15mm, bottomtitle=1mm, breakable, colback=white, colbacktitle=quarto-callout-tip-color!10!white, colframe=quarto-callout-tip-color-frame, coltitle=black, left=2mm, leftrule=.75mm, opacityback=0, opacitybacktitle=0.6, rightrule=.15mm, title=\textcolor{quarto-callout-tip-color}{\faLightbulb}\hspace{0.5em}{Key idea}, titlerule=0mm, toprule=.15mm, toptitle=1mm]

Interpretation An unbiased test should not be more likely to reject
\(H_0\) when \(H_0\) is true than when \(H_1\) is true. It should point
in the correct direction.

\end{tcolorbox}

\begin{tcolorbox}[enhanced jigsaw, arc=.35mm, bottomrule=.15mm, bottomtitle=1mm, breakable, colback=white, colbacktitle=quarto-callout-tip-color!10!white, colframe=quarto-callout-tip-color-frame, coltitle=black, left=2mm, leftrule=.75mm, opacityback=0, opacitybacktitle=0.6, rightrule=.15mm, title=\textcolor{quarto-callout-tip-color}{\faLightbulb}\hspace{0.5em}{Example}, titlerule=0mm, toprule=.15mm, toptitle=1mm]

\textbf{Example 11} (One-sided normal test is unbiased). Let
\(X_1,\ldots,X_n\sim N(\theta,\sigma^2)\) with known \(\sigma^2\). Test
\[H_0:\theta\le \theta_0
\qquad \text{versus} \qquad
H_1:\theta>\theta_0\] by rejecting \(H_0\) when
\[\frac{\bar X-\theta_0}{\sigma/\sqrt n}>c.\]

\end{tcolorbox}

\begin{tcolorbox}[enhanced jigsaw, arc=.35mm, bottomrule=.15mm, bottomtitle=1mm, breakable, colback=white, colbacktitle=quarto-callout-note-color!10!white, colframe=quarto-callout-note-color-frame, coltitle=black, left=2mm, leftrule=.75mm, opacityback=0, opacitybacktitle=0.6, rightrule=.15mm, title=\textcolor{quarto-callout-note-color}{\faInfo}\hspace{0.5em}{Solution}, titlerule=0mm, toprule=.15mm, toptitle=1mm]

The power function is
\[\beta(\theta)=1-\Phi\left(c+\frac{\theta_0-\theta}{\sigma/\sqrt n}\right).\]
This function is increasing in \(\theta\). Therefore, if
\(\theta'>\theta_0\) and \(\theta''\le\theta_0\), then
\[\beta(\theta')\ge \beta(\theta''),\] so the test is unbiased.

\end{tcolorbox}

\bookmarksetup{startatroot}

\chapter{Uniformly Most Powerful
Tests}\label{uniformly-most-powerful-tests}

This section introduces the strongest optimality criterion for
hypothesis tests in a class of tests.

\section{Most powerful and UMP tests}\label{most-powerful-and-ump-tests}

This subsection defines UMP tests.

A level-\(\alpha\) test guarantees that the Type I error probability is
at most \(\alpha\) for all parameter values in the null. After
controlling Type I error, we want to make Type II error small.
Equivalently, we want the power function to be large under the
alternative.

\begin{tcolorbox}[enhanced jigsaw, arc=.35mm, bottomrule=.15mm, bottomtitle=1mm, breakable, colback=white, colbacktitle=quarto-callout-note-color!10!white, colframe=quarto-callout-note-color-frame, coltitle=black, left=2mm, leftrule=.75mm, opacityback=0, opacitybacktitle=0.6, rightrule=.15mm, title=\textcolor{quarto-callout-note-color}{\faInfo}\hspace{0.5em}{Definition}, titlerule=0mm, toprule=.15mm, toptitle=1mm]

\textbf{Definition 12} (Uniformly most powerful test). Let
\(\mathcal C\) be a class of tests for \[H_0:\theta\in\Theta_0
\qquad \text{versus} \qquad
H_1:\theta\in\Theta_0^c.\] A test in \(\mathcal C\) with power function
\(\beta(\theta)\) is called a \textbf{uniformly most powerful} test in
\(\mathcal C\) if \[\beta(\theta)\ge \beta'(\theta)\] for every
\(\theta\in\Theta_0^c\) and for every competing power function
\(\beta'(\theta)\) corresponding to another test in \(\mathcal C\).

\end{tcolorbox}

\begin{tcolorbox}[enhanced jigsaw, arc=.35mm, bottomrule=.15mm, bottomtitle=1mm, breakable, colback=white, colbacktitle=quarto-callout-important-color!10!white, colframe=quarto-callout-important-color-frame, coltitle=black, left=2mm, leftrule=.75mm, opacityback=0, opacitybacktitle=0.6, rightrule=.15mm, title=\textcolor{quarto-callout-important-color}{\faExclamation}\hspace{0.5em}{Remark}, titlerule=0mm, toprule=.15mm, toptitle=1mm]

\emph{Remark 13}. UMP tests do not always exist, especially for
two-sided alternatives. When they do exist, they provide a very strong
optimality guarantee.

\end{tcolorbox}

\section{Neyman--Pearson lemma}\label{neymanpearson-lemma}

This subsection gives the fundamental result for finding most powerful
tests for simple hypotheses.

\begin{tcolorbox}[enhanced jigsaw, arc=.35mm, bottomrule=.15mm, bottomtitle=1mm, breakable, colback=white, colbacktitle=quarto-callout-important-color!10!white, colframe=quarto-callout-important-color-frame, coltitle=black, left=2mm, leftrule=.75mm, opacityback=0, opacitybacktitle=0.6, rightrule=.15mm, title=\textcolor{quarto-callout-important-color}{\faExclamation}\hspace{0.5em}{Theorem}, titlerule=0mm, toprule=.15mm, toptitle=1mm]

\textbf{Theorem 14} (Neyman--Pearson Lemma). \emph{Consider testing two
simple hypotheses \[H_0:\theta=\theta_0
\qquad \text{versus} \qquad
H_1:\theta=\theta_1.\] Let \(f(x\mid\theta_i)\) be the pdf or pmf of
\(X\) under \(\theta_i\). Define a rejection region \(R\) by
\[x\in R \quad \text{if} \quad f(x\mid\theta_1)>k f(x\mid\theta_0),\]
and
\[x\in R^c \quad \text{if} \quad f(x\mid\theta_1)<k f(x\mid\theta_0),\]
for some constant \(k\ge0\), chosen so that
\[\mathbb{P}_{\theta_0}(X\in R)=\alpha.\] Then any test with this
rejection region is a most powerful level-\(\alpha\) test for testing
\(H_0\) against \(H_1\).}

\end{tcolorbox}

\begin{tcolorbox}[enhanced jigsaw, arc=.35mm, bottomrule=.15mm, bottomtitle=1mm, breakable, colback=white, colbacktitle=quarto-callout-note-color!10!white, colframe=quarto-callout-note-color-frame, coltitle=black, left=2mm, leftrule=.75mm, opacityback=0, opacitybacktitle=0.6, rightrule=.15mm, title=\textcolor{quarto-callout-note-color}{\faInfo}\hspace{0.5em}{Proof}, titlerule=0mm, toprule=.15mm, toptitle=1mm]

\emph{Proof.} The Neyman--Pearson lemma says that for a fixed Type I
error probability, the best way to maximize power against the simple
alternative \(\theta_1\) is to reject where the likelihood under \(H_1\)
is large relative to the likelihood under \(H_0\). In other words, the
test rejects when \[\frac{f(x\mid\theta_1)}{f(x\mid\theta_0)}>k.\] This
is exactly the likelihood-ratio principle for simple hypotheses. The
full proof compares the power of any competing level-\(\alpha\) test
with the power of this likelihood-ratio test and uses the sign of
\[f(x\mid\theta_1)-k f(x\mid\theta_0)\] on the rejection and
non-rejection regions.~◻

\end{tcolorbox}

\begin{tcolorbox}[enhanced jigsaw, arc=.35mm, bottomrule=.15mm, bottomtitle=1mm, breakable, colback=white, colbacktitle=quarto-callout-tip-color!10!white, colframe=quarto-callout-tip-color-frame, coltitle=black, left=2mm, leftrule=.75mm, opacityback=0, opacitybacktitle=0.6, rightrule=.15mm, title=\textcolor{quarto-callout-tip-color}{\faLightbulb}\hspace{0.5em}{Key idea}, titlerule=0mm, toprule=.15mm, toptitle=1mm]

Likelihood-ratio form The Neyman--Pearson most powerful test rejects
\(H_0\) for large values of
\[\frac{f(x\mid\theta_1)}{f(x\mid\theta_0)}.\] This means that the
observed data are much more likely under \(H_1\) than under \(H_0\).

\end{tcolorbox}

\bookmarksetup{startatroot}

\chapter{UMP Examples}\label{ump-examples}

This section applies the Neyman--Pearson lemma to binomial and normal
models.

\section{UMP binomial test}\label{ump-binomial-test}

This subsection works out a discrete example where boundary cases may
require special attention.

\begin{tcolorbox}[enhanced jigsaw, arc=.35mm, bottomrule=.15mm, bottomtitle=1mm, breakable, colback=white, colbacktitle=quarto-callout-tip-color!10!white, colframe=quarto-callout-tip-color-frame, coltitle=black, left=2mm, leftrule=.75mm, opacityback=0, opacitybacktitle=0.6, rightrule=.15mm, title=\textcolor{quarto-callout-tip-color}{\faLightbulb}\hspace{0.5em}{Example}, titlerule=0mm, toprule=.15mm, toptitle=1mm]

\textbf{Example 15} (UMP binomial test). Let
\[X\sim \operatorname{Binomial}(2,\theta).\] Test
\[H_0:\theta=\theta_0=\frac12
\qquad \text{versus} \qquad
H_1:\theta=\theta_1=\frac34.\] Compute the likelihood ratios for
\(X\in\{0,1,2\}\).

\end{tcolorbox}

\begin{tcolorbox}[enhanced jigsaw, arc=.35mm, bottomrule=.15mm, bottomtitle=1mm, breakable, colback=white, colbacktitle=quarto-callout-note-color!10!white, colframe=quarto-callout-note-color-frame, coltitle=black, left=2mm, leftrule=.75mm, opacityback=0, opacitybacktitle=0.6, rightrule=.15mm, title=\textcolor{quarto-callout-note-color}{\faInfo}\hspace{0.5em}{Solution}, titlerule=0mm, toprule=.15mm, toptitle=1mm]

The pmf is \[f(x\mid\theta)=\binom{2}{x}\theta^x(1-\theta)^{2-x}.\] The
likelihood ratio is \[\frac{f(x\mid 3/4)}{f(x\mid 1/2)}.\] For \(x=0\),
\[\frac{f(0\mid 3/4)}{f(0\mid 1/2)}
=\frac{(1/4)^2}{(1/2)^2}=\frac14.\] For \(x=1\),
\[\frac{f(1\mid 3/4)}{f(1\mid 1/2)}
=\frac{2(3/4)(1/4)}{2(1/2)(1/2)}=\frac34.\] For \(x=2\),
\[\frac{f(2\mid 3/4)}{f(2\mid 1/2)}
=\frac{(3/4)^2}{(1/2)^2}=\frac94.\] Thus the likelihood ratios are
ordered as \[\frac14<\frac34<\frac94.\] By the Neyman--Pearson lemma,
the most powerful test rejects first at \(X=2\), then at \(X=1\), then
at \(X=0\), depending on the desired size.

For example:

\begin{itemize}
\item
  If \(3/4<k<9/4\), reject \(H_0\) when \(X=2\). The size is
  \[\alpha=\mathbb{P}_{1/2}(X=2)=\frac14.\]
\item
  If \(1/4<k<3/4\), reject \(H_0\) when \(X=1\) or \(X=2\). The size is
  \[\alpha=\mathbb{P}_{1/2}(X=1\text{ or }2)=\frac34.\]
\item
  If \(k>9/4\), never reject \(H_0\), giving size \(0\).
\item
  If \(k<1/4\), always reject \(H_0\), giving size \(1\).
\end{itemize}

\end{tcolorbox}

For the rejection rule \(R=\{2\}\), the power function is
\[\beta(\theta)=\mathbb{P}_\theta(X=2)=\theta^2.\] For the rejection
rule \(R=\{1,2\}\), the power function is
\[\beta(\theta)=\mathbb{P}_\theta(X\ge1)=1-(1-\theta)^2.\]

\section{UMP normal test}\label{ump-normal-test}

This subsection derives the one-sided most powerful normal test.

\begin{tcolorbox}[enhanced jigsaw, arc=.35mm, bottomrule=.15mm, bottomtitle=1mm, breakable, colback=white, colbacktitle=quarto-callout-tip-color!10!white, colframe=quarto-callout-tip-color-frame, coltitle=black, left=2mm, leftrule=.75mm, opacityback=0, opacitybacktitle=0.6, rightrule=.15mm, title=\textcolor{quarto-callout-tip-color}{\faLightbulb}\hspace{0.5em}{Example}, titlerule=0mm, toprule=.15mm, toptitle=1mm]

\textbf{Example 16} (UMP normal test). Suppose
\[X_1,\ldots,X_n\sim \operatorname{Normal}(\theta,\sigma^2),\] where
\(\sigma^2\) is known. Test \[H_0:\theta=\theta_0
\qquad \text{versus} \qquad
H_1:\theta=\theta_1,\] where \(\theta_1<\theta_0\).

\end{tcolorbox}

\begin{tcolorbox}[enhanced jigsaw, arc=.35mm, bottomrule=.15mm, bottomtitle=1mm, breakable, colback=white, colbacktitle=quarto-callout-note-color!10!white, colframe=quarto-callout-note-color-frame, coltitle=black, left=2mm, leftrule=.75mm, opacityback=0, opacitybacktitle=0.6, rightrule=.15mm, title=\textcolor{quarto-callout-note-color}{\faInfo}\hspace{0.5em}{Solution}, titlerule=0mm, toprule=.15mm, toptitle=1mm]

The sample mean \(\bar X\) is sufficient for \(\theta\), and
\[\bar X\sim \operatorname{Normal}\left(\theta,\frac{\sigma^2}{n}\right).\]
The density of \(\bar X=t\) is \[g(t\mid \theta)
=\frac{1}{\sqrt{2\pi\sigma^2/n}}
\exp\left\{-\frac{n(t-\theta)^2}{2\sigma^2}\right\}.\] The
Neyman--Pearson test rejects \(H_0\) when
\[g(\bar X\mid\theta_1)>k g(\bar X\mid\theta_0).\] Taking logarithms,
this inequality is equivalent to rejecting for small values of
\(\bar X\), because \(\theta_1<\theta_0\). Hence the most powerful test
has the form \[\text{Reject }H_0 \quad \text{if} \quad \bar X<c.\] To
make this a level-\(\alpha\) test, choose \(c\) so that
\[\alpha=\mathbb{P}_{\theta_0}(\bar X<c).\] Under \(H_0\),
\[\bar X\sim \operatorname{Normal}\left(\theta_0,\frac{\sigma^2}{n}\right),\]
so \[\boxed{c=\theta_0-\frac{\sigma}{\sqrt n}z_\alpha,}\] where
\(\mathbb{P}(Z>z_\alpha)=\alpha\).

\end{tcolorbox}

If instead \(\theta_1>\theta_0\), then the UMP test rejects for large
values of \(\bar X\):
\[\bar X>\theta_0+\frac{\sigma}{\sqrt n}z_\alpha.\]

\bookmarksetup{startatroot}

\chapter{\texorpdfstring{\(p\)-Values}{p-Values}}\label{p-values}

This section defines valid \(p\)-values and connects them to
significance levels and rejection rules.

\section{Definition and
interpretation}\label{definition-and-interpretation}

This subsection explains why \(p\)-values summarize the strength of
evidence against the null hypothesis.

A \(p\)-value is a test statistic that measures how surprising the
observed data are if the null hypothesis is true. Small \(p\)-values
provide evidence in favor of the alternative hypothesis.

\begin{tcolorbox}[enhanced jigsaw, arc=.35mm, bottomrule=.15mm, bottomtitle=1mm, breakable, colback=white, colbacktitle=quarto-callout-note-color!10!white, colframe=quarto-callout-note-color-frame, coltitle=black, left=2mm, leftrule=.75mm, opacityback=0, opacitybacktitle=0.6, rightrule=.15mm, title=\textcolor{quarto-callout-note-color}{\faInfo}\hspace{0.5em}{Definition}, titlerule=0mm, toprule=.15mm, toptitle=1mm]

\textbf{Definition 17} (Valid \(p\)-value). A statistic \(p(X)\) is a
valid \(p\)-value if, for every \(\theta\in\Theta_0\) and every
\(0\le \alpha\le 1\), \[\mathbb{P}_\theta(p(X)\le \alpha)\le \alpha.\]

\end{tcolorbox}

\begin{tcolorbox}[enhanced jigsaw, arc=.35mm, bottomrule=.15mm, bottomtitle=1mm, breakable, colback=white, colbacktitle=quarto-callout-important-color!10!white, colframe=quarto-callout-important-color-frame, coltitle=black, left=2mm, leftrule=.75mm, opacityback=0, opacitybacktitle=0.6, rightrule=.15mm, title=\textcolor{quarto-callout-important-color}{\faExclamation}\hspace{0.5em}{Remark}, titlerule=0mm, toprule=.15mm, toptitle=1mm]

\emph{Remark 18}. A \(p\)-value is not the probability that \(H_0\) is
true. It is a probability computed under the assumption that \(H_0\) is
true.

\end{tcolorbox}

Equivalently, the \(p\)-value can be understood as the smallest
significance level at which the observed test statistic would lead to
rejection:
\[p(x)=\inf\{\alpha:\text{ reject }H_0\text{ at level }\alpha\}.\]

\section{\texorpdfstring{General construction of a
\(p\)-value}{General construction of a p-value}}\label{general-construction-of-a-p-value}

This subsection gives a general formula for a valid \(p\)-value.

\begin{tcolorbox}[enhanced jigsaw, arc=.35mm, bottomrule=.15mm, bottomtitle=1mm, breakable, colback=white, colbacktitle=quarto-callout-important-color!10!white, colframe=quarto-callout-important-color-frame, coltitle=black, left=2mm, leftrule=.75mm, opacityback=0, opacitybacktitle=0.6, rightrule=.15mm, title=\textcolor{quarto-callout-important-color}{\faExclamation}\hspace{0.5em}{Theorem}, titlerule=0mm, toprule=.15mm, toptitle=1mm]

\textbf{Theorem 19} (General \(p\)-value formula). \emph{Let \(W(X)\) be
a test statistic such that large values of \(W\) provide evidence in
favor of \(H_1\). Then
\[p(x):=\sup_{\theta\in\Theta_0}\mathbb{P}_\theta\left(W(X)\ge W(x)\right)\]
defines a valid \(p\)-value.}

\end{tcolorbox}

\begin{tcolorbox}[enhanced jigsaw, arc=.35mm, bottomrule=.15mm, bottomtitle=1mm, breakable, colback=white, colbacktitle=quarto-callout-note-color!10!white, colframe=quarto-callout-note-color-frame, coltitle=black, left=2mm, leftrule=.75mm, opacityback=0, opacitybacktitle=0.6, rightrule=.15mm, title=\textcolor{quarto-callout-note-color}{\faInfo}\hspace{0.5em}{Proof}, titlerule=0mm, toprule=.15mm, toptitle=1mm]

\emph{Proof.} The quantity \(p(x)\) is the largest null probability of
observing a test statistic at least as extreme as the observed
statistic. Therefore, for any \(\theta\in\Theta_0\), the probability
that this tail probability is at most \(\alpha\) is no larger than
\(\alpha\). Hence \[\mathbb{P}_\theta(p(X)\le\alpha)\le\alpha,\] which
is the validity condition.~◻

\end{tcolorbox}

\begin{tcolorbox}[enhanced jigsaw, arc=.35mm, bottomrule=.15mm, bottomtitle=1mm, breakable, colback=white, colbacktitle=quarto-callout-tip-color!10!white, colframe=quarto-callout-tip-color-frame, coltitle=black, left=2mm, leftrule=.75mm, opacityback=0, opacitybacktitle=0.6, rightrule=.15mm, title=\textcolor{quarto-callout-tip-color}{\faLightbulb}\hspace{0.5em}{Key idea}, titlerule=0mm, toprule=.15mm, toptitle=1mm]

Interpretation A \(p\)-value answers the question:
\[\text{How surprising is my observed signal, if the null hypothesis were true?}\]
The smaller the \(p\)-value, the stronger the evidence against \(H_0\).

\end{tcolorbox}

\section{\texorpdfstring{One-sided normal
\(p\)-value}{One-sided normal p-value}}\label{one-sided-normal-p-value}

This subsection computes a standard \(p\)-value for a one-sided normal
test.

\begin{tcolorbox}[enhanced jigsaw, arc=.35mm, bottomrule=.15mm, bottomtitle=1mm, breakable, colback=white, colbacktitle=quarto-callout-tip-color!10!white, colframe=quarto-callout-tip-color-frame, coltitle=black, left=2mm, leftrule=.75mm, opacityback=0, opacitybacktitle=0.6, rightrule=.15mm, title=\textcolor{quarto-callout-tip-color}{\faLightbulb}\hspace{0.5em}{Example}, titlerule=0mm, toprule=.15mm, toptitle=1mm]

\textbf{Example 20} (One-sided normal \(p\)-value). Let
\(X_1,\ldots,X_n\sim \operatorname{Normal}(\mu,\sigma^2)\) with
\(\sigma\) known. Test \[H_0:\mu\le \mu_0
\qquad \text{versus} \qquad
H_1:\mu>\mu_0.\] Use the statistic
\[W(X)=Z=\frac{\bar X-\mu_0}{\sigma/\sqrt n}.\] Find the \(p\)-value.

\end{tcolorbox}

\begin{tcolorbox}[enhanced jigsaw, arc=.35mm, bottomrule=.15mm, bottomtitle=1mm, breakable, colback=white, colbacktitle=quarto-callout-note-color!10!white, colframe=quarto-callout-note-color-frame, coltitle=black, left=2mm, leftrule=.75mm, opacityback=0, opacitybacktitle=0.6, rightrule=.15mm, title=\textcolor{quarto-callout-note-color}{\faInfo}\hspace{0.5em}{Solution}, titlerule=0mm, toprule=.15mm, toptitle=1mm]

Large values of \(Z\) favor \(H_1\). For any \(\mu\le\mu_0\),
\[\begin{aligned}
\mathbb{P}_\mu(W\ge w)
&=\mathbb{P}_\mu\left(\frac{\bar X-\mu_0}{\sigma/\sqrt n}\ge w\right)\\
&=\mathbb{P}_\mu\left(\frac{\bar X-\mu}{\sigma/\sqrt n}\ge w+\frac{\mu_0-\mu}{\sigma/\sqrt n}\right).
\end{aligned}\] The right tail probability is largest when
\(\mu=\mu_0\). Hence the supremum over \(\Theta_0=\{\mu:\mu\le\mu_0\}\)
is attained at \(\mu=\mu_0\).

If \[z_{\mathrm{obs}}=\frac{\bar x-\mu_0}{\sigma/\sqrt n},\] then the
valid \(p\)-value is
\[\boxed{p(x)=\mathbb{P}_{\mu_0}(Z\ge z_{\mathrm{obs}})=1-\Phi(z_{\mathrm{obs}}).}\]

\end{tcolorbox}

\begin{tcolorbox}[enhanced jigsaw, arc=.35mm, bottomrule=.15mm, bottomtitle=1mm, breakable, colback=white, colbacktitle=quarto-callout-tip-color!10!white, colframe=quarto-callout-tip-color-frame, coltitle=black, left=2mm, leftrule=.75mm, opacityback=0, opacitybacktitle=0.6, rightrule=.15mm, title=\textcolor{quarto-callout-tip-color}{\faLightbulb}\hspace{0.5em}{Example}, titlerule=0mm, toprule=.15mm, toptitle=1mm]

\textbf{Example 21} (Numerical \(p\)-value). Suppose
\(z_{\mathrm{obs}}=2.1\) in the one-sided normal test above. Compute the
\(p\)-value and state the decision at levels \(\alpha=0.05\) and
\(\alpha=0.01\).

\end{tcolorbox}

\begin{tcolorbox}[enhanced jigsaw, arc=.35mm, bottomrule=.15mm, bottomtitle=1mm, breakable, colback=white, colbacktitle=quarto-callout-note-color!10!white, colframe=quarto-callout-note-color-frame, coltitle=black, left=2mm, leftrule=.75mm, opacityback=0, opacitybacktitle=0.6, rightrule=.15mm, title=\textcolor{quarto-callout-note-color}{\faInfo}\hspace{0.5em}{Solution}, titlerule=0mm, toprule=.15mm, toptitle=1mm]

The one-sided \(p\)-value is \[p=1-\Phi(2.1)\approx 0.0179.\] At
\(\alpha=0.05\), since \(0.0179<0.05\), reject \(H_0\).

At \(\alpha=0.01\), since \(0.0179>0.01\), fail to reject \(H_0\).

\end{tcolorbox}

\bookmarksetup{startatroot}

\chapter{Loss and Risk Functions for
Tests}\label{loss-and-risk-functions-for-tests}

This section treats hypothesis testing as a decision problem with
possible costs for different mistakes.

\section{Decision-theoretic
formulation}\label{decision-theoretic-formulation}

This subsection introduces actions, loss functions, and risk functions
for hypothesis testing.

In hypothesis testing, there are two possible actions:
\[a_0=\text{fail to reject }H_0,
\qquad
a_1=\text{reject }H_0.\] A decision rule \(\delta(x)\) selects an action
based on the observed data \(x\).

\begin{tcolorbox}[enhanced jigsaw, arc=.35mm, bottomrule=.15mm, bottomtitle=1mm, breakable, colback=white, colbacktitle=quarto-callout-note-color!10!white, colframe=quarto-callout-note-color-frame, coltitle=black, left=2mm, leftrule=.75mm, opacityback=0, opacitybacktitle=0.6, rightrule=.15mm, title=\textcolor{quarto-callout-note-color}{\faInfo}\hspace{0.5em}{Definition}, titlerule=0mm, toprule=.15mm, toptitle=1mm]

\textbf{Definition 22} (Loss function). A loss function
\[L:\Theta\times \mathcal A\to \mathbb{R}\] measures the penalty for
taking action \(a\in\mathcal A\) when the true parameter is \(\theta\).

\end{tcolorbox}

\begin{tcolorbox}[enhanced jigsaw, arc=.35mm, bottomrule=.15mm, bottomtitle=1mm, breakable, colback=white, colbacktitle=quarto-callout-tip-color!10!white, colframe=quarto-callout-tip-color-frame, coltitle=black, left=2mm, leftrule=.75mm, opacityback=0, opacitybacktitle=0.6, rightrule=.15mm, title=\textcolor{quarto-callout-tip-color}{\faLightbulb}\hspace{0.5em}{Example}, titlerule=0mm, toprule=.15mm, toptitle=1mm]

\textbf{Example 23} (0--1 loss). For testing \(H_0:\theta\in\Theta_0\)
versus \(H_1:\theta\in\Theta_0^c\), the ordinary 0--1 loss is
\[L(\theta,a)=
\begin{cases}
0, & \theta\in\Theta_0\text{ and }a=a_0,\text{ or }\theta\in\Theta_0^c\text{ and }a=a_1,\\
1, & \theta\in\Theta_0\text{ and }a=a_1,\text{ or }\theta\in\Theta_0^c\text{ and }a=a_0.
\end{cases}\]

\end{tcolorbox}

\begin{tcolorbox}[enhanced jigsaw, arc=.35mm, bottomrule=.15mm, bottomtitle=1mm, breakable, colback=white, colbacktitle=quarto-callout-note-color!10!white, colframe=quarto-callout-note-color-frame, coltitle=black, left=2mm, leftrule=.75mm, opacityback=0, opacitybacktitle=0.6, rightrule=.15mm, title=\textcolor{quarto-callout-note-color}{\faInfo}\hspace{0.5em}{Solution}, titlerule=0mm, toprule=.15mm, toptitle=1mm]

The loss is \(0\) for a correct decision and \(1\) for an incorrect
decision. Thus both Type I and Type II errors are treated as equally
costly.

\end{tcolorbox}

\section{Generalized 0--1 loss}\label{generalized-01-loss}

This subsection allows Type I and Type II errors to have different
costs.

\begin{tcolorbox}[enhanced jigsaw, arc=.35mm, bottomrule=.15mm, bottomtitle=1mm, breakable, colback=white, colbacktitle=quarto-callout-note-color!10!white, colframe=quarto-callout-note-color-frame, coltitle=black, left=2mm, leftrule=.75mm, opacityback=0, opacitybacktitle=0.6, rightrule=.15mm, title=\textcolor{quarto-callout-note-color}{\faInfo}\hspace{0.5em}{Definition}, titlerule=0mm, toprule=.15mm, toptitle=1mm]

\textbf{Definition 24} (Generalized 0--1 loss). Let \(c_I\) be the cost
of a Type I error and \(c_{II}\) be the cost of a Type II error. Define
\[L(\theta,a_0)=
\begin{cases}
0, & \theta\in\Theta_0,\\
c_{II}, & \theta\in\Theta_0^c,
\end{cases}
\qquad
L(\theta,a_1)=
\begin{cases}
c_I, & \theta\in\Theta_0,\\
0, & \theta\in\Theta_0^c.
\end{cases}\]

\end{tcolorbox}

If \(c_I=c_{II}=1\), this reduces to the ordinary 0--1 loss. If
\(c_I\ne c_{II}\), the test should account for the relative seriousness
of the two errors.

\section{Risk function}\label{risk-function-1}

This subsection defines risk as the expected loss of a test.

\begin{tcolorbox}[enhanced jigsaw, arc=.35mm, bottomrule=.15mm, bottomtitle=1mm, breakable, colback=white, colbacktitle=quarto-callout-note-color!10!white, colframe=quarto-callout-note-color-frame, coltitle=black, left=2mm, leftrule=.75mm, opacityback=0, opacitybacktitle=0.6, rightrule=.15mm, title=\textcolor{quarto-callout-note-color}{\faInfo}\hspace{0.5em}{Definition}, titlerule=0mm, toprule=.15mm, toptitle=1mm]

\textbf{Definition 25} (Risk function). The risk function of a decision
rule \(\delta\) is
\[R(\theta,\delta)=\mathbb{E}_\theta[L(\theta,\delta(X))].\]

\end{tcolorbox}

Let the power function be
\[\beta(\theta)=\mathbb{P}_\theta(\delta(X)=a_1),\] the probability of
rejecting \(H_0\) under parameter \(\theta\). Under generalized 0--1
loss, \[R(\theta,\delta)=
\begin{cases}
c_I\beta(\theta), & \theta\in\Theta_0,\\[3pt]
c_{II}\{1-\beta(\theta)\}, & \theta\in\Theta_0^c.
\end{cases}\]

\begin{tcolorbox}[enhanced jigsaw, arc=.35mm, bottomrule=.15mm, bottomtitle=1mm, breakable, colback=white, colbacktitle=quarto-callout-tip-color!10!white, colframe=quarto-callout-tip-color-frame, coltitle=black, left=2mm, leftrule=.75mm, opacityback=0, opacitybacktitle=0.6, rightrule=.15mm, title=\textcolor{quarto-callout-tip-color}{\faLightbulb}\hspace{0.5em}{Key idea}, titlerule=0mm, toprule=.15mm, toptitle=1mm]

Risk interpretation For \(\theta\in\Theta_0\), the risk is the Type I
error probability weighted by the cost \(c_I\). For
\(\theta\in\Theta_0^c\), the risk is the Type II error probability
weighted by the cost \(c_{II}\).

\end{tcolorbox}

\bookmarksetup{startatroot}

\chapter{Risk Function Example}\label{risk-function-example}

This section computes the risk function for a one-sided normal test.

\begin{tcolorbox}[enhanced jigsaw, arc=.35mm, bottomrule=.15mm, bottomtitle=1mm, breakable, colback=white, colbacktitle=quarto-callout-tip-color!10!white, colframe=quarto-callout-tip-color-frame, coltitle=black, left=2mm, leftrule=.75mm, opacityback=0, opacitybacktitle=0.6, rightrule=.15mm, title=\textcolor{quarto-callout-tip-color}{\faLightbulb}\hspace{0.5em}{Example}, titlerule=0mm, toprule=.15mm, toptitle=1mm]

\textbf{Example 26} (Normal one-sided test risk). Let
\[X_1,\ldots,X_n\stackrel{iid}{\sim} N(\theta,\sigma^2),
\qquad
\sigma=1,
\qquad
n=25.\] Test \[H_0:\theta\le 0
\qquad \text{versus} \qquad
H_1:\theta>0.\] Use the size \(\alpha=0.05\) test that rejects \(H_0\)
when \[\bar X>z_{1-\alpha}\frac{\sigma}{\sqrt n}.\] Take \(c_I=2\) and
\(c_{II}=1\). Compute the power function and risk function.

\end{tcolorbox}

\begin{tcolorbox}[enhanced jigsaw, arc=.35mm, bottomrule=.15mm, bottomtitle=1mm, breakable, colback=white, colbacktitle=quarto-callout-note-color!10!white, colframe=quarto-callout-note-color-frame, coltitle=black, left=2mm, leftrule=.75mm, opacityback=0, opacitybacktitle=0.6, rightrule=.15mm, title=\textcolor{quarto-callout-note-color}{\faInfo}\hspace{0.5em}{Solution}, titlerule=0mm, toprule=.15mm, toptitle=1mm]

Here \[z_{1-\alpha}=z_{0.95}\approx 1.645,
\qquad
\frac{\sigma}{\sqrt n}=\frac15.\] So the rejection rule is
\[\bar X>1.645\cdot \frac15=0.329.\] For a general \(\theta\),
\[\bar X\sim N\left(\theta,\frac1{25}\right).\] Thus the power function
is \[\begin{aligned}
\beta(\theta)
&=\mathbb{P}_\theta(\bar X>0.329)\\
&=\mathbb{P}\left(\frac{\bar X-\theta}{1/5}>\frac{0.329-\theta}{1/5}\right)\\
&=1-\Phi(1.645-5\theta).
\end{aligned}\] Under the generalized 0--1 loss with \(c_I=2\) and
\(c_{II}=1\), \[R(\theta)=
\begin{cases}
2\beta(\theta), & \theta\le 0,\\[3pt]
1-\beta(\theta), & \theta>0.
\end{cases}\]

\end{tcolorbox}

\begin{tcolorbox}[enhanced jigsaw, arc=.35mm, bottomrule=.15mm, bottomtitle=1mm, breakable, colback=white, colbacktitle=quarto-callout-tip-color!10!white, colframe=quarto-callout-tip-color-frame, coltitle=black, left=2mm, leftrule=.75mm, opacityback=0, opacitybacktitle=0.6, rightrule=.15mm, title=\textcolor{quarto-callout-tip-color}{\faLightbulb}\hspace{0.5em}{Example}, titlerule=0mm, toprule=.15mm, toptitle=1mm]

\textbf{Example 27} (Numerical risk values). Using the previous example,
compute the risk at \(\theta=0,-0.2,0.2,0.5\).

\end{tcolorbox}

\begin{tcolorbox}[enhanced jigsaw, arc=.35mm, bottomrule=.15mm, bottomtitle=1mm, breakable, colback=white, colbacktitle=quarto-callout-note-color!10!white, colframe=quarto-callout-note-color-frame, coltitle=black, left=2mm, leftrule=.75mm, opacityback=0, opacitybacktitle=0.6, rightrule=.15mm, title=\textcolor{quarto-callout-note-color}{\faInfo}\hspace{0.5em}{Solution}, titlerule=0mm, toprule=.15mm, toptitle=1mm]

The power function is \[\beta(\theta)=1-\Phi(1.645-5\theta).\] At the
boundary \(\theta=0\), \[\beta(0)=0.05,
\qquad
R(0)=2(0.05)=0.10.\] At \(\theta=-0.2\),
\[\beta(-0.2)=1-\Phi(1.645+1)=1-\Phi(2.645)\approx 0.0041,\] so
\[R(-0.2)=2(0.0041)=0.0082.\] At \(\theta=0.2\),
\[\beta(0.2)=1-\Phi(1.645-1)=1-\Phi(0.645)\approx 0.260.\] Since
\(\theta>0\), \[R(0.2)=1-\beta(0.2)\approx 0.740.\] At \(\theta=0.5\),
\[\beta(0.5)=1-\Phi(1.645-2.5)=1-\Phi(-0.855)=\Phi(0.855)\approx 0.803.\]
Thus \[R(0.5)=1-0.803=0.197.\]

\end{tcolorbox}

\begin{tcolorbox}[enhanced jigsaw, arc=.35mm, bottomrule=.15mm, bottomtitle=1mm, breakable, colback=white, colbacktitle=quarto-callout-warning-color!10!white, colframe=quarto-callout-warning-color-frame, coltitle=black, left=2mm, leftrule=.75mm, opacityback=0, opacitybacktitle=0.6, rightrule=.15mm, title=\textcolor{quarto-callout-warning-color}{\faExclamationTriangle}\hspace{0.5em}{Important lesson}, titlerule=0mm, toprule=.15mm, toptitle=1mm]

Important lesson The risk can be small under parts of the null and large
near the alternative boundary. The shape of the risk function depends on
both the power function and the relative costs \(c_I\) and \(c_{II}\).

\end{tcolorbox}

\bookmarksetup{startatroot}

\chapter{Practice Problems}\label{practice-problems-6}

This section gives practice problems that reinforce the main definitions
and calculations from the section.

\begin{tcolorbox}[enhanced jigsaw, arc=.35mm, bottomrule=.15mm, bottomtitle=1mm, breakable, colback=white, colbacktitle=quarto-callout-caution-color!10!white, colframe=quarto-callout-caution-color-frame, coltitle=black, left=2mm, leftrule=.75mm, opacityback=0, opacitybacktitle=0.6, rightrule=.15mm, title=\textcolor{quarto-callout-caution-color}{\faFire}\hspace{0.5em}{Practice Problem}, titlerule=0mm, toprule=.15mm, toptitle=1mm]

\textbf{Practice Problem 28} (Power function for a binomial test). Let
\(X\sim\operatorname{Binomial}(6,\theta)\) and test
\[H_0:\theta\le \frac12
\qquad \text{versus} \qquad
H_1:\theta>\frac12.\] Consider the test that rejects \(H_0\) if
\(X\ge5\).

\begin{enumerate}
\def\labelenumi{\arabic{enumi}.}
\item
  Find the power function.
\item
  Find the size of the test.
\item
  Find the Type II error probability when \(\theta=0.8\).
\end{enumerate}

\end{tcolorbox}

\begin{tcolorbox}[enhanced jigsaw, arc=.35mm, bottomrule=.15mm, bottomtitle=1mm, breakable, colback=white, colbacktitle=quarto-callout-note-color!10!white, colframe=quarto-callout-note-color-frame, coltitle=black, left=2mm, leftrule=.75mm, opacityback=0, opacitybacktitle=0.6, rightrule=.15mm, title=\textcolor{quarto-callout-note-color}{\faInfo}\hspace{0.5em}{Solution}, titlerule=0mm, toprule=.15mm, toptitle=1mm]

(a) The rejection region is \(R=\{5,6\}\), so
\[\beta(\theta)=\mathbb{P}_\theta(X\ge5)
=\binom65\theta^5(1-\theta)+\binom66\theta^6
=6\theta^5(1-\theta)+\theta^6.\] (b) Since \(\beta(\theta)\) is
increasing in \(\theta\), the size is attained at \(\theta=1/2\):
\[\alpha=\beta(1/2)=6\left(\frac12\right)^5\left(\frac12\right)+\left(\frac12\right)^6
=\frac{6+1}{64}=\frac7{64}.\] (c) At \(\theta=0.8\),
\[\beta(0.8)=6(0.8)^5(0.2)+(0.8)^6.\] The Type II error probability is
\[1-\beta(0.8)=1-\{6(0.8)^5(0.2)+(0.8)^6\}.\] Numerically, this is
approximately \(0.3446\).

\end{tcolorbox}

\begin{tcolorbox}[enhanced jigsaw, arc=.35mm, bottomrule=.15mm, bottomtitle=1mm, breakable, colback=white, colbacktitle=quarto-callout-caution-color!10!white, colframe=quarto-callout-caution-color-frame, coltitle=black, left=2mm, leftrule=.75mm, opacityback=0, opacitybacktitle=0.6, rightrule=.15mm, title=\textcolor{quarto-callout-caution-color}{\faFire}\hspace{0.5em}{Practice Problem}, titlerule=0mm, toprule=.15mm, toptitle=1mm]

\textbf{Practice Problem 29} (Normal sample size). Let
\(X_1,\ldots,X_n\sim N(\theta,\sigma^2)\) with known \(\sigma^2\). Test
\[H_0:\theta\le\theta_0
\qquad \text{versus} \qquad
H_1:\theta>\theta_0\] using a level \(0.05\) test. Find a formula for
the smallest \(n\) such that the power is at least \(0.90\) when
\(\theta=\theta_0+\Delta\), where \(\Delta>0\).

\end{tcolorbox}

\begin{tcolorbox}[enhanced jigsaw, arc=.35mm, bottomrule=.15mm, bottomtitle=1mm, breakable, colback=white, colbacktitle=quarto-callout-note-color!10!white, colframe=quarto-callout-note-color-frame, coltitle=black, left=2mm, leftrule=.75mm, opacityback=0, opacitybacktitle=0.6, rightrule=.15mm, title=\textcolor{quarto-callout-note-color}{\faInfo}\hspace{0.5em}{Solution}, titlerule=0mm, toprule=.15mm, toptitle=1mm]

The level \(0.05\) test rejects when
\[\frac{\bar X-\theta_0}{\sigma/\sqrt n}>z_{0.05}.\] At
\(\theta=\theta_0+\Delta\), the power is \[\beta(\theta_0+\Delta)
=1-\Phi\left(z_{0.05}-\frac{\Delta\sqrt n}{\sigma}\right).\] We require
\[1-\Phi\left(z_{0.05}-\frac{\Delta\sqrt n}{\sigma}\right)\ge 0.90.\]
Equivalently,
\[\Phi\left(z_{0.05}-\frac{\Delta\sqrt n}{\sigma}\right)\le0.10.\] Thus
\[z_{0.05}-\frac{\Delta\sqrt n}{\sigma}\le z_{0.90}^{\text{right tail}},\]
where \(\mathbb{P}(Z>z_{0.90}^{\text{right tail}})=0.90\), so
\(z_{0.90}^{\text{right tail}}\approx -1.2816\). Hence
\[\sqrt n\ge \frac{\sigma}{\Delta}(z_{0.05}+1.2816).\] Therefore
\[\boxed{n\ge \left(\frac{\sigma}{\Delta}(z_{0.05}+1.2816)\right)^2.}\]
The smallest integer \(n\) is the ceiling of the right-hand side.

\end{tcolorbox}

\begin{tcolorbox}[enhanced jigsaw, arc=.35mm, bottomrule=.15mm, bottomtitle=1mm, breakable, colback=white, colbacktitle=quarto-callout-caution-color!10!white, colframe=quarto-callout-caution-color-frame, coltitle=black, left=2mm, leftrule=.75mm, opacityback=0, opacitybacktitle=0.6, rightrule=.15mm, title=\textcolor{quarto-callout-caution-color}{\faFire}\hspace{0.5em}{Practice Problem}, titlerule=0mm, toprule=.15mm, toptitle=1mm]

\textbf{Practice Problem 30} (One-sided normal p-value). Let
\(X_1,\ldots,X_{25}\sim N(\mu,4)\) and test \[H_0:\mu\le 10
\qquad \text{versus} \qquad
H_1:\mu>10.\] Suppose \(\bar x=10.9\). Compute the \(p\)-value and state
the decision at \(\alpha=0.05\).

\end{tcolorbox}

\begin{tcolorbox}[enhanced jigsaw, arc=.35mm, bottomrule=.15mm, bottomtitle=1mm, breakable, colback=white, colbacktitle=quarto-callout-note-color!10!white, colframe=quarto-callout-note-color-frame, coltitle=black, left=2mm, leftrule=.75mm, opacityback=0, opacitybacktitle=0.6, rightrule=.15mm, title=\textcolor{quarto-callout-note-color}{\faInfo}\hspace{0.5em}{Solution}, titlerule=0mm, toprule=.15mm, toptitle=1mm]

Here \(\sigma=2\) and \(n=25\), so \(\sigma/\sqrt n=2/5=0.4\). The
observed test statistic is
\[z_{\mathrm{obs}}=\frac{10.9-10}{0.4}=2.25.\] The one-sided \(p\)-value
is \[p=1-\Phi(2.25)\approx 0.0122.\] Since \(0.0122<0.05\), reject
\(H_0\) at the \(5\%\) level.

\end{tcolorbox}

\begin{tcolorbox}[enhanced jigsaw, arc=.35mm, bottomrule=.15mm, bottomtitle=1mm, breakable, colback=white, colbacktitle=quarto-callout-caution-color!10!white, colframe=quarto-callout-caution-color-frame, coltitle=black, left=2mm, leftrule=.75mm, opacityback=0, opacitybacktitle=0.6, rightrule=.15mm, title=\textcolor{quarto-callout-caution-color}{\faFire}\hspace{0.5em}{Practice Problem}, titlerule=0mm, toprule=.15mm, toptitle=1mm]

\textbf{Practice Problem 31} (Risk function). For the one-sided normal
test in the previous problem, use generalized 0--1 loss with \(c_I=3\)
and \(c_{II}=1\). Write the risk function in terms of the power
function.

\end{tcolorbox}

\begin{tcolorbox}[enhanced jigsaw, arc=.35mm, bottomrule=.15mm, bottomtitle=1mm, breakable, colback=white, colbacktitle=quarto-callout-note-color!10!white, colframe=quarto-callout-note-color-frame, coltitle=black, left=2mm, leftrule=.75mm, opacityback=0, opacitybacktitle=0.6, rightrule=.15mm, title=\textcolor{quarto-callout-note-color}{\faInfo}\hspace{0.5em}{Solution}, titlerule=0mm, toprule=.15mm, toptitle=1mm]

The decision is to reject or fail to reject \(H_0:\mu\le10\). The risk
function is \[R(\mu,\delta)=
\begin{cases}
3\beta(\mu), & \mu\le 10,\\[3pt]
1-\beta(\mu), & \mu>10.
\end{cases}\] where \[\beta(\mu)=\mathbb{P}_\mu(\text{reject }H_0).\]
For a level-\(0.05\) test,
\[\beta(\mu)=1-\Phi\left(z_{0.05}+\frac{10-\mu}{2/5}\right).\]

\end{tcolorbox}

\bookmarksetup{startatroot}

\chapter{Summary}\label{summary-8}

This section summarizes the main concepts used to evaluate hypothesis
tests.

{\def\LTcaptype{none} % do not increment counter
\begin{longtable}[]{@{}lll@{}}
\toprule\noalign{}
\textbf{Concept} & \textbf{Definition} & \textbf{Meaning} \\
\midrule\noalign{}
\endhead
\bottomrule\noalign{}
\endlastfoot
Power function & \(\beta(\theta)=\mathbb{P}_\theta(X\in R)\) &
Probability of rejecting \(H_0\) at parameter value \(\theta\) \\
Type I error & \(\mathbb{P}_\theta(X\in R)\) for \(\theta\in\Theta_0\) &
Rejecting \(H_0\) when \(H_0\) is true \\
Type II error & \(1-\beta(\theta)\) for \(\theta\in\Theta_0^c\) &
Failing to reject \(H_0\) when \(H_1\) is true \\
Size & \(\sup_{\theta\in\Theta_0}\beta(\theta)\) & Worst-case Type I
error probability \\
Level \(\alpha\) & \(\sup_{\theta\in\Theta_0}\beta(\theta)\le\alpha\) &
Type I error controlled by \(\alpha\) \\
Unbiased test & \(\beta(\theta')\ge\beta(\theta'')\) for
\(\theta'\in\Theta_0^c\), \(\theta''\in\Theta_0\) & Rejects more often
under the alternative than under the null \\
UMP test & Maximizes power for all \(\theta\in\Theta_0^c\) within a
class & Strong optimality criterion \\
\(p\)-value & Smallest level at which \(H_0\) is rejected & Measures
evidence against \(H_0\) \\
Risk function &
\(R(\theta,\delta)=\mathbb{E}_\theta[L(\theta,\delta(X))]\) & Expected
loss of a test \\
\end{longtable}
}

\begin{tcolorbox}[enhanced jigsaw, arc=.35mm, bottomrule=.15mm, bottomtitle=1mm, breakable, colback=white, colbacktitle=quarto-callout-tip-color!10!white, colframe=quarto-callout-tip-color-frame, coltitle=black, left=2mm, leftrule=.75mm, opacityback=0, opacitybacktitle=0.6, rightrule=.15mm, title=\textcolor{quarto-callout-tip-color}{\faLightbulb}\hspace{0.5em}{Key idea}, titlerule=0mm, toprule=.15mm, toptitle=1mm]

Key message A hypothesis test is not evaluated by whether it is correct
on one dataset. It is evaluated by its long-run error probabilities, its
power function, and, when costs matter, its risk function.

\end{tcolorbox}

\bookmarksetup{startatroot}

\chapter{Chapter 16: Interval Estimation I --- Finding Interval
Estimators}\label{chapter-16-interval-estimation-i-finding-interval-estimators}

This chapter begins interval estimation. The main goal is to move from a
single point estimate to a random interval procedure with a controlled
coverage probability. We study confidence intervals built from
acceptance regions, pivotal quantities, CDF pivoting, and Bayesian
posterior distributions.

\begin{tcolorbox}[enhanced jigsaw, arc=.35mm, bottomrule=.15mm, bottomtitle=1mm, breakable, colback=white, colbacktitle=quarto-callout-note-color!10!white, colframe=quarto-callout-note-color-frame, coltitle=black, left=2mm, leftrule=.75mm, opacityback=0, opacitybacktitle=0.6, rightrule=.15mm, title=\textcolor{quarto-callout-note-color}{\faInfo}\hspace{0.5em}{Topics}, titlerule=0mm, toprule=.15mm, toptitle=1mm]

Interval estimators; coverage probability; confidence coefficient;
confidence intervals from test inversion; pivotal quantities; pivoting
the CDF; normal mean and variance intervals; proportion intervals;
uniform and exponential examples; Bayesian credible intervals;
beta-binomial, normal-normal, and gamma-Poisson examples.

\end{tcolorbox}

\bookmarksetup{startatroot}

\chapter{From Point Estimation to Interval
Estimation}\label{from-point-estimation-to-interval-estimation}

This section begins the study of interval estimators, which quantify
uncertainty by producing a range of plausible parameter values instead
of a single number.

In point estimation, a statistic \(W(X_1,\ldots,X_n)\) estimates an
unknown fixed parameter \(\theta\) by a single number. For example,
\(\bar X\) estimates a population mean \(\mu\), \(S^2\) estimates a
population variance \(\sigma^2\), and \(\hat p=X/n\) estimates a
binomial proportion \(p\).

However, a single estimate does not show how much uncertainty remains.
An interval estimator gives lower and upper random bounds. The goal is
not only to estimate the parameter, but also to describe how reliable
the estimate is.

\begin{tcolorbox}[enhanced jigsaw, arc=.35mm, bottomrule=.15mm, bottomtitle=1mm, breakable, colback=white, colbacktitle=quarto-callout-note-color!10!white, colframe=quarto-callout-note-color-frame, coltitle=black, left=2mm, leftrule=.75mm, opacityback=0, opacitybacktitle=0.6, rightrule=.15mm, title=\textcolor{quarto-callout-note-color}{\faInfo}\hspace{0.5em}{Definition}, titlerule=0mm, toprule=.15mm, toptitle=1mm]

\textbf{Definition 1} (Interval estimator). Let \(X=(X_1,\ldots,X_n)\)
be a random sample from a population distribution with pdf or pmf
\(f(x\mid \theta)\). An \textbf{interval estimator} of a real parameter
\(\theta\) is a pair of statistics \(L(X)\) and \(U(X)\), with
\(L(X)\le U(X)\), such that after observing \(X=x\) we report
\[L(x)\le \theta \le U(x).\] The random interval \([L(X),U(X)]\) is the
interval estimator, and the observed interval \([L(x),U(x)]\) is the
interval estimate.

\end{tcolorbox}

\begin{tcolorbox}[enhanced jigsaw, arc=.35mm, bottomrule=.15mm, bottomtitle=1mm, breakable, colback=white, colbacktitle=quarto-callout-tip-color!10!white, colframe=quarto-callout-tip-color-frame, coltitle=black, left=2mm, leftrule=.75mm, opacityback=0, opacitybacktitle=0.6, rightrule=.15mm, title=\textcolor{quarto-callout-tip-color}{\faLightbulb}\hspace{0.5em}{Key idea}, titlerule=0mm, toprule=.15mm, toptitle=1mm]

Random interval, fixed parameter In frequentist interval estimation, the
parameter \(\theta\) is treated as fixed but unknown. The interval
\([L(X),U(X)]\) is random because it depends on the random sample \(X\).

\end{tcolorbox}

\section{A first normal example}\label{a-first-normal-example}

This subsection introduces confidence intervals through the simplest
normal model.

\begin{tcolorbox}[enhanced jigsaw, arc=.35mm, bottomrule=.15mm, bottomtitle=1mm, breakable, colback=white, colbacktitle=quarto-callout-tip-color!10!white, colframe=quarto-callout-tip-color-frame, coltitle=black, left=2mm, leftrule=.75mm, opacityback=0, opacitybacktitle=0.6, rightrule=.15mm, title=\textcolor{quarto-callout-tip-color}{\faLightbulb}\hspace{0.5em}{Example}, titlerule=0mm, toprule=.15mm, toptitle=1mm]

\textbf{Example 2} (Normal mean with known variance). Suppose
\(X_1,\ldots,X_n\sim \operatorname{Normal}(\mu,\sigma^2)\)
independently, where \(\sigma^2\) is known and \(\mu\) is unknown. An
interval estimator for \(\mu\) is
\[\left[\bar X-k\frac{\sigma}{\sqrt n},\; \bar X+k\frac{\sigma}{\sqrt n}\right]\]
for a chosen constant \(k>0\).

\end{tcolorbox}

\begin{tcolorbox}[enhanced jigsaw, arc=.35mm, bottomrule=.15mm, bottomtitle=1mm, breakable, colback=white, colbacktitle=quarto-callout-note-color!10!white, colframe=quarto-callout-note-color-frame, coltitle=black, left=2mm, leftrule=.75mm, opacityback=0, opacitybacktitle=0.6, rightrule=.15mm, title=\textcolor{quarto-callout-note-color}{\faInfo}\hspace{0.5em}{Solution}, titlerule=0mm, toprule=.15mm, toptitle=1mm]

Since
\[Z=\frac{\bar X-\mu}{\sigma/\sqrt n}\sim \operatorname{Normal}(0,1),\]
we have \[\begin{aligned}
\mathbb{P}_\mu\left(\mu\in \left[\bar X-k\frac{\sigma}{\sqrt n},\bar X+k\frac{\sigma}{\sqrt n}\right]\right)
&=\mathbb{P}_\mu\left(-k\le \frac{\bar X-\mu}{\sigma/\sqrt n}\le k\right)\\
&=\mathbb{P}(-k\le Z\le k)\\
&=2\Phi(k)-1.
\end{aligned}\] Thus the constant \(k\) controls the probability that
the random interval covers the fixed parameter \(\mu\). For example,
choosing \(k=z_{1-\alpha/2}\) gives coverage \(1-\alpha\).

\end{tcolorbox}

\begin{tcolorbox}[enhanced jigsaw, arc=.35mm, bottomrule=.15mm, bottomtitle=1mm, breakable, colback=white, colbacktitle=quarto-callout-note-color!10!white, colframe=quarto-callout-note-color-frame, coltitle=black, left=2mm, leftrule=.75mm, opacityback=0, opacitybacktitle=0.6, rightrule=.15mm, title=\textcolor{quarto-callout-note-color}{\faInfo}\hspace{0.5em}{Remark}, titlerule=0mm, toprule=.15mm, toptitle=1mm]

\emph{Remark 3}. In the normal interval above, \(\bar X\) is random and
\(\mu\) is an unknown constant. After observing the data, the interval
becomes a fixed numerical interval.

\end{tcolorbox}

\bookmarksetup{startatroot}

\chapter{Coverage Probability and Confidence
Coefficient}\label{coverage-probability-and-confidence-coefficient}

This section defines the main frequentist criterion for interval
estimators: how often the random interval covers the true parameter
value.

\begin{tcolorbox}[enhanced jigsaw, arc=.35mm, bottomrule=.15mm, bottomtitle=1mm, breakable, colback=white, colbacktitle=quarto-callout-note-color!10!white, colframe=quarto-callout-note-color-frame, coltitle=black, left=2mm, leftrule=.75mm, opacityback=0, opacitybacktitle=0.6, rightrule=.15mm, title=\textcolor{quarto-callout-note-color}{\faInfo}\hspace{0.5em}{Definition}, titlerule=0mm, toprule=.15mm, toptitle=1mm]

\textbf{Definition 4} (Coverage probability). For an interval estimator
\([L(X),U(X)]\) of a parameter \(\theta\), the \textbf{coverage
probability} at \(\theta\) is
\[\mathbb{P}_\theta\bigl(\theta\in [L(X),U(X)]\bigr)
    =\mathbb{P}_\theta\bigl(L(X)\le \theta\le U(X)\bigr).\]

\end{tcolorbox}

\begin{tcolorbox}[enhanced jigsaw, arc=.35mm, bottomrule=.15mm, bottomtitle=1mm, breakable, colback=white, colbacktitle=quarto-callout-note-color!10!white, colframe=quarto-callout-note-color-frame, coltitle=black, left=2mm, leftrule=.75mm, opacityback=0, opacitybacktitle=0.6, rightrule=.15mm, title=\textcolor{quarto-callout-note-color}{\faInfo}\hspace{0.5em}{Definition}, titlerule=0mm, toprule=.15mm, toptitle=1mm]

\textbf{Definition 5} (Confidence coefficient). The \textbf{confidence
coefficient} or guaranteed coverage of an interval estimator is the
worst-case coverage over all possible parameter values:
\[\inf_{\theta}\mathbb{P}_\theta\bigl(\theta\in [L(X),U(X)]\bigr).\] If
this quantity is at least \(1-\alpha\), we call the interval a
\(100(1-\alpha)\%\) confidence interval.

\end{tcolorbox}

\begin{tcolorbox}[enhanced jigsaw, arc=.35mm, bottomrule=.15mm, bottomtitle=1mm, breakable, colback=white, colbacktitle=quarto-callout-tip-color!10!white, colframe=quarto-callout-tip-color-frame, coltitle=black, left=2mm, leftrule=.75mm, opacityback=0, opacitybacktitle=0.6, rightrule=.15mm, title=\textcolor{quarto-callout-tip-color}{\faLightbulb}\hspace{0.5em}{Key idea}, titlerule=0mm, toprule=.15mm, toptitle=1mm]

Interpretation A \(95\%\) confidence interval procedure is a random
procedure that covers the true fixed parameter in at least \(95\%\) of
repeated samples. It does not mean that, after observing one fixed
interval, there is a \(95\%\) frequentist probability that the fixed
parameter lies inside that fixed interval.

\end{tcolorbox}

\section{Location-equivariant
intervals}\label{location-equivariant-intervals}

This subsection explains why the usual normal confidence intervals have
coverage that does not depend on the unknown mean.

\begin{tcolorbox}[enhanced jigsaw, arc=.35mm, bottomrule=.15mm, bottomtitle=1mm, breakable, colback=white, colbacktitle=quarto-callout-tip-color!10!white, colframe=quarto-callout-tip-color-frame, coltitle=black, left=2mm, leftrule=.75mm, opacityback=0, opacitybacktitle=0.6, rightrule=.15mm, title=\textcolor{quarto-callout-tip-color}{\faLightbulb}\hspace{0.5em}{Example}, titlerule=0mm, toprule=.15mm, toptitle=1mm]

\textbf{Example 6} (A family of normal intervals). Suppose
\(X_1,\ldots,X_n\sim \operatorname{Normal}(\mu,\sigma^2)\) independently
with known \(\sigma^2\). Let \(c<d\) be fixed constants. Consider
\[I_1(X)=\left[\bar X+c\frac{\sigma}{\sqrt n},\; \bar X+d\frac{\sigma}{\sqrt n}\right].\]
Find its coverage probability.

\end{tcolorbox}

\begin{tcolorbox}[enhanced jigsaw, arc=.35mm, bottomrule=.15mm, bottomtitle=1mm, breakable, colback=white, colbacktitle=quarto-callout-note-color!10!white, colframe=quarto-callout-note-color-frame, coltitle=black, left=2mm, leftrule=.75mm, opacityback=0, opacitybacktitle=0.6, rightrule=.15mm, title=\textcolor{quarto-callout-note-color}{\faInfo}\hspace{0.5em}{Solution}, titlerule=0mm, toprule=.15mm, toptitle=1mm]

The event \(\mu\in I_1(X)\) is equivalent to
\[\bar X+c\frac{\sigma}{\sqrt n}\le \mu \le \bar X+d\frac{\sigma}{\sqrt n}.\]
Subtracting \(\bar X\) and multiplying by \(\sqrt n/\sigma\) gives
\[c\le \frac{\mu-\bar X}{\sigma/\sqrt n}\le d.\] Equivalently,
\[-d\le \frac{\bar X-\mu}{\sigma/\sqrt n}\le -c.\] Therefore
\[\mathbb{P}_\mu(\mu\in I_1(X))
=\mathbb{P}(-d\le Z\le -c)
=\Phi(-c)-\Phi(-d).\] Using symmetry of the standard normal
distribution, this can also be written as \[\Phi(d)-\Phi(c)\] if the
constants are defined through the equivalent form \(c\le Z\le d\). The
key point is that the coverage does not depend on \(\mu\).

\end{tcolorbox}

\begin{tcolorbox}[enhanced jigsaw, arc=.35mm, bottomrule=.15mm, bottomtitle=1mm, breakable, colback=white, colbacktitle=quarto-callout-important-color!10!white, colframe=quarto-callout-important-color-frame, coltitle=black, left=2mm, leftrule=.75mm, opacityback=0, opacitybacktitle=0.6, rightrule=.15mm, title=\textcolor{quarto-callout-important-color}{\faExclamation}\hspace{0.5em}{Corollary}, titlerule=0mm, toprule=.15mm, toptitle=1mm]

\textbf{Corollary 7} (Two-sided normal confidence interval).
\emph{Choosing \(c=-z_{1-\alpha/2}\) and \(d=z_{1-\alpha/2}\) yields the
usual two-sided \(100(1-\alpha)\%\) confidence interval
\[\left[\bar X-z_{1-\alpha/2}\frac{\sigma}{\sqrt n},\; \bar X+z_{1-\alpha/2}\frac{\sigma}{\sqrt n}\right].\]}

\end{tcolorbox}

\section{A bad interval for a location
parameter}\label{a-bad-interval-for-a-location-parameter}

This subsection shows that not every random interval is a good interval
estimator.

\begin{tcolorbox}[enhanced jigsaw, arc=.35mm, bottomrule=.15mm, bottomtitle=1mm, breakable, colback=white, colbacktitle=quarto-callout-tip-color!10!white, colframe=quarto-callout-tip-color-frame, coltitle=black, left=2mm, leftrule=.75mm, opacityback=0, opacitybacktitle=0.6, rightrule=.15mm, title=\textcolor{quarto-callout-tip-color}{\faLightbulb}\hspace{0.5em}{Example}, titlerule=0mm, toprule=.15mm, toptitle=1mm]

\textbf{Example 8} (Multiplicative interval for a normal mean). Suppose
\(X_1,\ldots,X_n\sim \operatorname{Normal}(\mu,\sigma^2)\) with known
\(\sigma^2\). Let \(0<a<1<b\) and consider \[I_2(X)=[a\bar X,b\bar X].\]
Show that this interval has bad coverage at \(\mu=0\).

\end{tcolorbox}

\begin{tcolorbox}[enhanced jigsaw, arc=.35mm, bottomrule=.15mm, bottomtitle=1mm, breakable, colback=white, colbacktitle=quarto-callout-note-color!10!white, colframe=quarto-callout-note-color-frame, coltitle=black, left=2mm, leftrule=.75mm, opacityback=0, opacitybacktitle=0.6, rightrule=.15mm, title=\textcolor{quarto-callout-note-color}{\faInfo}\hspace{0.5em}{Solution}, titlerule=0mm, toprule=.15mm, toptitle=1mm]

At \(\mu=0\), the interval covers \(\mu\) if
\[a\bar X\le 0\le b\bar X.\] Since \(a>0\) and \(b>0\), the first
inequality implies \(\bar X\le 0\), and the second inequality implies
\(\bar X\ge 0\). Hence coverage occurs only when \(\bar X=0\). Since
\(\bar X\) is a continuous normal random variable,
\[\mathbb{P}_0(\bar X=0)=0.\] Thus the coverage probability at \(\mu=0\)
is \(0\). This interval is inappropriate for a location family.

\end{tcolorbox}

\begin{tcolorbox}[enhanced jigsaw, arc=.35mm, bottomrule=.15mm, bottomtitle=1mm, breakable, colback=white, colbacktitle=quarto-callout-warning-color!10!white, colframe=quarto-callout-warning-color-frame, coltitle=black, left=2mm, leftrule=.75mm, opacityback=0, opacitybacktitle=0.6, rightrule=.15mm, title=\textcolor{quarto-callout-warning-color}{\faExclamationTriangle}\hspace{0.5em}{Warning}, titlerule=0mm, toprule=.15mm, toptitle=1mm]

Important lesson A confidence interval should be designed to have good
coverage uniformly over the parameter space. A formula that looks like
an interval may still have very poor coverage.

\end{tcolorbox}

\bookmarksetup{startatroot}

\chapter{Methods for Building Interval
Estimators}\label{methods-for-building-interval-estimators}

This section summarizes the main methods for constructing confidence
intervals and credible intervals.

There are several standard methods:

\begin{enumerate}
\def\labelenumi{\arabic{enumi}.}
\item
  inverting a test statistic;
\item
  using pivotal quantities;
\item
  pivoting the CDF;
\item
  Bayesian credible intervals.
\end{enumerate}

Each method starts from a probability statement involving the data and
parameter, then solves the statement for the parameter.

\bookmarksetup{startatroot}

\chapter{Method 1: Inverting a Test
Statistic}\label{method-1-inverting-a-test-statistic}

This section explains the deep connection between hypothesis tests and
confidence intervals.

The basic idea is simple: a confidence interval consists of all
parameter values that would not be rejected by a corresponding
hypothesis test.

\begin{tcolorbox}[enhanced jigsaw, arc=.35mm, bottomrule=.15mm, bottomtitle=1mm, breakable, colback=white, colbacktitle=quarto-callout-tip-color!10!white, colframe=quarto-callout-tip-color-frame, coltitle=black, left=2mm, leftrule=.75mm, opacityback=0, opacitybacktitle=0.6, rightrule=.15mm, title=\textcolor{quarto-callout-tip-color}{\faLightbulb}\hspace{0.5em}{Example}, titlerule=0mm, toprule=.15mm, toptitle=1mm]

\textbf{Example 9} (Inverting the two-sided normal \(z\)-test). Suppose
\(X_1,\ldots,X_n\sim \operatorname{Normal}(\mu,\sigma^2)\) independently
with known \(\sigma^2\). For each fixed value \(\mu_0\), test
\[H_0:\mu=\mu_0
    \qquad \text{versus}\qquad
    H_1:\mu\ne \mu_0.\] At level \(\alpha\), the usual two-sided test
rejects \(H_0\) when
\[\left|\frac{\bar X-\mu_0}{\sigma/\sqrt n}\right|>z_{1-\alpha/2}.\]
Find the confidence interval obtained by inverting this test.

\end{tcolorbox}

\begin{tcolorbox}[enhanced jigsaw, arc=.35mm, bottomrule=.15mm, bottomtitle=1mm, breakable, colback=white, colbacktitle=quarto-callout-note-color!10!white, colframe=quarto-callout-note-color-frame, coltitle=black, left=2mm, leftrule=.75mm, opacityback=0, opacitybacktitle=0.6, rightrule=.15mm, title=\textcolor{quarto-callout-note-color}{\faInfo}\hspace{0.5em}{Solution}, titlerule=0mm, toprule=.15mm, toptitle=1mm]

The test accepts, or fails to reject, \(H_0:\mu=\mu_0\) when
\[\left|\frac{\bar X-\mu_0}{\sigma/\sqrt n}\right|\le z_{1-\alpha/2}.\]
This is equivalent to
\[-z_{1-\alpha/2}\le \frac{\bar X-\mu_0}{\sigma/\sqrt n}\le z_{1-\alpha/2}.\]
Solving for \(\mu_0\) gives
\[\bar X-z_{1-\alpha/2}\frac{\sigma}{\sqrt n}
    \le \mu_0 \le
    \bar X+z_{1-\alpha/2}\frac{\sigma}{\sqrt n}.\] Therefore the set of
all values \(\mu_0\) not rejected by the test is
\[C(X)=\left[\bar X-z_{1-\alpha/2}\frac{\sigma}{\sqrt n},\;
    \bar X+z_{1-\alpha/2}\frac{\sigma}{\sqrt n}\right].\] Because the
test has level \(\alpha\), the confidence set has coverage \(1-\alpha\).

\end{tcolorbox}

\section{Acceptance regions and confidence
sets}\label{acceptance-regions-and-confidence-sets}

This subsection states the formal test-confidence interval duality.

For a test of \(H_0:\theta=\theta_0\), let \(A(\theta_0)\) be the
acceptance region. That is, \(x\in A(\theta_0)\) means the data \(x\) do
not lead us to reject the hypothesis \(\theta=\theta_0\).

The corresponding confidence set is
\[C(x)=\{\theta_0: x\in A(\theta_0)\}.\] So the confidence set is the
set of parameter values that are compatible with the observed data under
the testing rule.

\begin{tcolorbox}[enhanced jigsaw, arc=.35mm, bottomrule=.15mm, bottomtitle=1mm, breakable, colback=white, colbacktitle=quarto-callout-important-color!10!white, colframe=quarto-callout-important-color-frame, coltitle=black, left=2mm, leftrule=.75mm, opacityback=0, opacitybacktitle=0.6, rightrule=.15mm, title=\textcolor{quarto-callout-important-color}{\faExclamation}\hspace{0.5em}{Theorem}, titlerule=0mm, toprule=.15mm, toptitle=1mm]

\textbf{Theorem 10} (Tests and confidence sets). \emph{For each
\(\theta_0\), let \(A(\theta_0)\) be the acceptance region of a
level-\(\alpha\) test of \[H_0:\theta=\theta_0.\] Define
\[C(x)=\{\theta_0:x\in A(\theta_0)\}.\] Then \(C(X)\) is a
\(100(1-\alpha)\%\) confidence set:
\[\mathbb{P}_\theta(\theta\in C(X))\ge 1-\alpha.\] Conversely, given any
\(100(1-\alpha)\%\) confidence set \(C(X)\), define
\[A(\theta_0)=\{x:\theta_0\in C(x)\}.\] Then \(A(\theta_0)\) is the
acceptance region of a level-\(\alpha\) test of
\(H_0:\theta=\theta_0\).}

\end{tcolorbox}

\begin{tcolorbox}[enhanced jigsaw, arc=.35mm, bottomrule=.15mm, bottomtitle=1mm, breakable, colback=white, colbacktitle=quarto-callout-note-color!10!white, colframe=quarto-callout-note-color-frame, coltitle=black, left=2mm, leftrule=.75mm, opacityback=0, opacitybacktitle=0.6, rightrule=.15mm, title=\textcolor{quarto-callout-note-color}{\faInfo}\hspace{0.5em}{Proof}, titlerule=0mm, toprule=.15mm, toptitle=1mm]

\emph{Proof.} For the first direction,
\[\{\theta\in C(X)\}=\{X\in A(\theta)\}.\] Since \(A(\theta)\) is the
acceptance region of a level-\(\alpha\) test,
\[\mathbb{P}_\theta(X\notin A(\theta))\le \alpha,\] so
\[\mathbb{P}_\theta(\theta\in C(X))=\mathbb{P}_\theta(X\in A(\theta))\ge 1-\alpha.\]
For the converse, if \(C(X)\) has coverage at least \(1-\alpha\), then
\[\mathbb{P}_{\theta_0}(\theta_0\notin C(X))\le \alpha.\] But
\(\theta_0\notin C(X)\) is exactly the rejection event for the test with
acceptance region \(A(\theta_0)=\{x:\theta_0\in C(x)\}\). Therefore the
test has level \(\alpha\).~◻

\end{tcolorbox}

\bookmarksetup{startatroot}

\chapter{Method 2: Pivotal
Quantities}\label{method-2-pivotal-quantities}

This section introduces pivotal quantities, one of the most useful tools
for deriving confidence intervals.

\begin{tcolorbox}[enhanced jigsaw, arc=.35mm, bottomrule=.15mm, bottomtitle=1mm, breakable, colback=white, colbacktitle=quarto-callout-note-color!10!white, colframe=quarto-callout-note-color-frame, coltitle=black, left=2mm, leftrule=.75mm, opacityback=0, opacitybacktitle=0.6, rightrule=.15mm, title=\textcolor{quarto-callout-note-color}{\faInfo}\hspace{0.5em}{Definition}, titlerule=0mm, toprule=.15mm, toptitle=1mm]

\textbf{Definition 11} (Pivotal quantity). A \textbf{pivot} or
\textbf{pivotal quantity} is a function \(Q(X,\theta)\) of the data
\(X\) and the parameter \(\theta\) whose distribution does not depend on
any unknown parameter. That is, when \(X\sim f(x\mid\theta)\), the
distribution of \(Q(X,\theta)\) is the same for every \(\theta\).

\end{tcolorbox}

\begin{tcolorbox}[enhanced jigsaw, arc=.35mm, bottomrule=.15mm, bottomtitle=1mm, breakable, colback=white, colbacktitle=quarto-callout-tip-color!10!white, colframe=quarto-callout-tip-color-frame, coltitle=black, left=2mm, leftrule=.75mm, opacityback=0, opacitybacktitle=0.6, rightrule=.15mm, title=\textcolor{quarto-callout-tip-color}{\faLightbulb}\hspace{0.5em}{Key idea}, titlerule=0mm, toprule=.15mm, toptitle=1mm]

Pivot method If \(Q(X,\theta)\) is a pivot and \(\mathcal A\) is a set
such that \[\mathbb{P}_\theta(Q(X,\theta)\in \mathcal A)=1-\alpha,\]
then \[C(x)=\{\theta: Q(x,\theta)\in \mathcal A\}\] forms a
\(100(1-\alpha)\%\) confidence set.

\end{tcolorbox}

\section{Normal mean with known
variance}\label{normal-mean-with-known-variance}

This subsection derives the usual normal interval using a pivot.

\begin{tcolorbox}[enhanced jigsaw, arc=.35mm, bottomrule=.15mm, bottomtitle=1mm, breakable, colback=white, colbacktitle=quarto-callout-tip-color!10!white, colframe=quarto-callout-tip-color-frame, coltitle=black, left=2mm, leftrule=.75mm, opacityback=0, opacitybacktitle=0.6, rightrule=.15mm, title=\textcolor{quarto-callout-tip-color}{\faLightbulb}\hspace{0.5em}{Example}, titlerule=0mm, toprule=.15mm, toptitle=1mm]

\textbf{Example 12} (Normal mean, known variance). Suppose
\(X_1,\ldots,X_n\sim \operatorname{Normal}(\mu,\sigma^2)\) independently
with known \(\sigma^2\). Find a \(100(1-\alpha)\%\) confidence interval
for \(\mu\).

\end{tcolorbox}

\begin{tcolorbox}[enhanced jigsaw, arc=.35mm, bottomrule=.15mm, bottomtitle=1mm, breakable, colback=white, colbacktitle=quarto-callout-note-color!10!white, colframe=quarto-callout-note-color-frame, coltitle=black, left=2mm, leftrule=.75mm, opacityback=0, opacitybacktitle=0.6, rightrule=.15mm, title=\textcolor{quarto-callout-note-color}{\faInfo}\hspace{0.5em}{Solution}, titlerule=0mm, toprule=.15mm, toptitle=1mm]

The pivot is
\[Z=\frac{\bar X-\mu}{\sigma/\sqrt n}\sim \operatorname{Normal}(0,1).\]
Choose standard normal quantiles so that
\[\mathbb{P}\left(-z_{1-\alpha/2}\le Z\le z_{1-\alpha/2}\right)=1-\alpha.\]
Substitute the pivot:
\[\mathbb{P}\left(-z_{1-\alpha/2}\le \frac{\bar X-\mu}{\sigma/\sqrt n}\le z_{1-\alpha/2}\right)=1-\alpha.\]
Solving for \(\mu\) gives
\[\mu\in \left[\bar X-z_{1-\alpha/2}\frac{\sigma}{\sqrt n},\;
    \bar X+z_{1-\alpha/2}\frac{\sigma}{\sqrt n}\right].\]

\end{tcolorbox}

\section{Normal mean with unknown
variance}\label{normal-mean-with-unknown-variance}

This subsection adds the classical Student \(t\) interval, which is a
central example of the pivot method.

\begin{tcolorbox}[enhanced jigsaw, arc=.35mm, bottomrule=.15mm, bottomtitle=1mm, breakable, colback=white, colbacktitle=quarto-callout-tip-color!10!white, colframe=quarto-callout-tip-color-frame, coltitle=black, left=2mm, leftrule=.75mm, opacityback=0, opacitybacktitle=0.6, rightrule=.15mm, title=\textcolor{quarto-callout-tip-color}{\faLightbulb}\hspace{0.5em}{Example}, titlerule=0mm, toprule=.15mm, toptitle=1mm]

\textbf{Example 13} (Normal mean, unknown variance). Suppose
\(X_1,\ldots,X_n\sim \operatorname{Normal}(\mu,\sigma^2)\)
independently, where both \(\mu\) and \(\sigma^2\) are unknown. Find a
\(100(1-\alpha)\%\) confidence interval for \(\mu\).

\end{tcolorbox}

\begin{tcolorbox}[enhanced jigsaw, arc=.35mm, bottomrule=.15mm, bottomtitle=1mm, breakable, colback=white, colbacktitle=quarto-callout-note-color!10!white, colframe=quarto-callout-note-color-frame, coltitle=black, left=2mm, leftrule=.75mm, opacityback=0, opacitybacktitle=0.6, rightrule=.15mm, title=\textcolor{quarto-callout-note-color}{\faInfo}\hspace{0.5em}{Solution}, titlerule=0mm, toprule=.15mm, toptitle=1mm]

The sample standard deviation is
\[S=\sqrt{\frac{1}{n-1}\sum_{i=1}^n (X_i-\bar X)^2}.\] For a normal
sample, \[T=\frac{\bar X-\mu}{S/\sqrt n}\sim t_{n-1},\] which does not
depend on \(\mu\) or \(\sigma^2\). Thus \(T\) is a pivot. Choose
\(t_{n-1,1-\alpha/2}\) such that
\[\mathbb{P}(-t_{n-1,1-\alpha/2}\le T\le t_{n-1,1-\alpha/2})=1-\alpha.\]
Solving for \(\mu\) gives
\[\mu\in \left[\bar X-t_{n-1,1-\alpha/2}\frac{S}{\sqrt n},\;
    \bar X+t_{n-1,1-\alpha/2}\frac{S}{\sqrt n}\right].\] This is the
classical Student \(t\) confidence interval.

\end{tcolorbox}

\section{Variance of a normal
distribution}\label{variance-of-a-normal-distribution}

This subsection constructs an exact chi-square confidence interval for a
normal variance.

\begin{tcolorbox}[enhanced jigsaw, arc=.35mm, bottomrule=.15mm, bottomtitle=1mm, breakable, colback=white, colbacktitle=quarto-callout-tip-color!10!white, colframe=quarto-callout-tip-color-frame, coltitle=black, left=2mm, leftrule=.75mm, opacityback=0, opacitybacktitle=0.6, rightrule=.15mm, title=\textcolor{quarto-callout-tip-color}{\faLightbulb}\hspace{0.5em}{Example}, titlerule=0mm, toprule=.15mm, toptitle=1mm]

\textbf{Example 14} (Variance of a normal distribution). Suppose
\(X_1,\ldots,X_n\sim \operatorname{Normal}(\mu,\sigma^2)\)
independently. Find a \(100(1-\alpha)\%\) confidence interval for
\(\sigma^2\).

\end{tcolorbox}

\begin{tcolorbox}[enhanced jigsaw, arc=.35mm, bottomrule=.15mm, bottomtitle=1mm, breakable, colback=white, colbacktitle=quarto-callout-note-color!10!white, colframe=quarto-callout-note-color-frame, coltitle=black, left=2mm, leftrule=.75mm, opacityback=0, opacitybacktitle=0.6, rightrule=.15mm, title=\textcolor{quarto-callout-note-color}{\faInfo}\hspace{0.5em}{Solution}, titlerule=0mm, toprule=.15mm, toptitle=1mm]

For a normal sample, \[Q=\frac{(n-1)S^2}{\sigma^2}\sim \chi^2_{n-1}.\]
This is a pivot. Let \(\chi^2_{\nu,\gamma}\) denote the \(\gamma\)
quantile of the chi-square distribution with \(\nu\) degrees of freedom.
Then
\[\mathbb{P}\left(\chi^2_{n-1,\alpha/2}\le \frac{(n-1)S^2}{\sigma^2}
    \le \chi^2_{n-1,1-\alpha/2}\right)=1-\alpha.\] Now solve the
inequalities for \(\sigma^2\). Because \(\sigma^2\) appears in the
denominator, the endpoints reverse when solving: \[\sigma^2\in \left[
    \frac{(n-1)S^2}{\chi^2_{n-1,1-\alpha/2}},\;
    \frac{(n-1)S^2}{\chi^2_{n-1,\alpha/2}}
    \right].\]

\end{tcolorbox}

\section{Proportion using an approximate
pivot}\label{proportion-using-an-approximate-pivot}

This subsection derives the common Wald interval for a binomial
proportion.

\begin{tcolorbox}[enhanced jigsaw, arc=.35mm, bottomrule=.15mm, bottomtitle=1mm, breakable, colback=white, colbacktitle=quarto-callout-tip-color!10!white, colframe=quarto-callout-tip-color-frame, coltitle=black, left=2mm, leftrule=.75mm, opacityback=0, opacitybacktitle=0.6, rightrule=.15mm, title=\textcolor{quarto-callout-tip-color}{\faLightbulb}\hspace{0.5em}{Example}, titlerule=0mm, toprule=.15mm, toptitle=1mm]

\textbf{Example 15} (Proportion: approximate pivot). Let
\(X\sim \operatorname{Binomial}(n,p)\) and let \(\hat p=X/n\). Use the
central limit theorem to derive an approximate confidence interval for
\(p\).

\end{tcolorbox}

\begin{tcolorbox}[enhanced jigsaw, arc=.35mm, bottomrule=.15mm, bottomtitle=1mm, breakable, colback=white, colbacktitle=quarto-callout-note-color!10!white, colframe=quarto-callout-note-color-frame, coltitle=black, left=2mm, leftrule=.75mm, opacityback=0, opacitybacktitle=0.6, rightrule=.15mm, title=\textcolor{quarto-callout-note-color}{\faInfo}\hspace{0.5em}{Solution}, titlerule=0mm, toprule=.15mm, toptitle=1mm]

By the central limit theorem,
\[\frac{\hat p-p}{\sqrt{p(1-p)/n}}\approx \operatorname{Normal}(0,1).\]
Since the unknown \(p\) appears in the standard error, the Wald interval
replaces \(p\) by \(\hat p\):
\[\frac{\hat p-p}{\sqrt{\hat p(1-\hat p)/n}}\approx \operatorname{Normal}(0,1).\]
Thus an approximate \(100(1-\alpha)\%\) confidence interval is
\[p\in \left[
    \hat p-z_{1-\alpha/2}\sqrt{\frac{\hat p(1-\hat p)}{n}},\;
    \hat p+z_{1-\alpha/2}\sqrt{\frac{\hat p(1-\hat p)}{n}}
    \right].\] This interval is called the Wald confidence interval.

\end{tcolorbox}

\begin{tcolorbox}[enhanced jigsaw, arc=.35mm, bottomrule=.15mm, bottomtitle=1mm, breakable, colback=white, colbacktitle=quarto-callout-warning-color!10!white, colframe=quarto-callout-warning-color-frame, coltitle=black, left=2mm, leftrule=.75mm, opacityback=0, opacitybacktitle=0.6, rightrule=.15mm, title=\textcolor{quarto-callout-warning-color}{\faExclamationTriangle}\hspace{0.5em}{Warning}, titlerule=0mm, toprule=.15mm, toptitle=1mm]

Approximate intervals The Wald interval is easy to remember, but it can
perform poorly when \(n\) is small or when \(p\) is close to \(0\) or
\(1\). Its coverage is approximate, not exact.

\end{tcolorbox}

\bookmarksetup{startatroot}

\chapter{Method 3: Pivoting the CDF}\label{method-3-pivoting-the-cdf}

This section studies interval construction using the CDF of a statistic
whose distribution depends on the parameter.

Sometimes a simple pivot is not obvious, but we know the distribution of
a statistic \(T(X)\). Suppose
\[F_T(t;\theta)=\mathbb{P}_\theta(T(X)\le t).\] If we can find
inequalities involving \(T\) and \(\theta\) with probability
\(1-\alpha\), then we can solve those inequalities for \(\theta\).

\begin{tcolorbox}[enhanced jigsaw, arc=.35mm, bottomrule=.15mm, bottomtitle=1mm, breakable, colback=white, colbacktitle=quarto-callout-tip-color!10!white, colframe=quarto-callout-tip-color-frame, coltitle=black, left=2mm, leftrule=.75mm, opacityback=0, opacitybacktitle=0.6, rightrule=.15mm, title=\textcolor{quarto-callout-tip-color}{\faLightbulb}\hspace{0.5em}{Key idea}, titlerule=0mm, toprule=.15mm, toptitle=1mm]

CDF pivoting strategy Find functions \(a(\theta)\) and \(b(\theta)\)
such that
\[\mathbb{P}_\theta(a(\theta)\le T(X)\le b(\theta))=1-\alpha.\] After
observing \(T=t_0\), solve \[a(\theta)\le t_0\le b(\theta)\] for
\(\theta\). The solution is a confidence set.

\end{tcolorbox}

\section{Uniform distribution}\label{uniform-distribution-1}

This subsection uses an order statistic to build a confidence interval
for the endpoint of a uniform distribution.

\begin{tcolorbox}[enhanced jigsaw, arc=.35mm, bottomrule=.15mm, bottomtitle=1mm, breakable, colback=white, colbacktitle=quarto-callout-tip-color!10!white, colframe=quarto-callout-tip-color-frame, coltitle=black, left=2mm, leftrule=.75mm, opacityback=0, opacitybacktitle=0.6, rightrule=.15mm, title=\textcolor{quarto-callout-tip-color}{\faLightbulb}\hspace{0.5em}{Example}, titlerule=0mm, toprule=.15mm, toptitle=1mm]

\textbf{Example 16} (Uniform endpoint). Suppose
\(X_1,\ldots,X_n\sim \operatorname{Uniform}(0,\theta)\) independently.
Let \[M=\max_{1\le i\le n}X_i.\] Find a one-sided \(100(1-\alpha)\%\)
confidence interval for \(\theta\).

\end{tcolorbox}

\begin{tcolorbox}[enhanced jigsaw, arc=.35mm, bottomrule=.15mm, bottomtitle=1mm, breakable, colback=white, colbacktitle=quarto-callout-note-color!10!white, colframe=quarto-callout-note-color-frame, coltitle=black, left=2mm, leftrule=.75mm, opacityback=0, opacitybacktitle=0.6, rightrule=.15mm, title=\textcolor{quarto-callout-note-color}{\faInfo}\hspace{0.5em}{Solution}, titlerule=0mm, toprule=.15mm, toptitle=1mm]

The CDF of \(M\) is
\[F_M(m;\theta)=\mathbb{P}_\theta(M\le m)=\left(\frac{m}{\theta}\right)^n,
    \qquad 0<m<\theta.\] Since \(M\le \theta\) always, the lower
endpoint must be at least \(M\). Choose \(c\) so that
\[\mathbb{P}_\theta(M\ge c\theta)=1-\alpha.\] Now
\[\mathbb{P}_\theta(M<c\theta)=c^n.\] Thus choose \(c=\alpha^{1/n}\) if
we want lower-tail probability \(\alpha\), or choose
\(c=(1-\alpha)^{1/n}\) depending on the convention for the one-sided
interval.

Using the construction in which
\[\mathbb{P}_\theta\left(M\ge \theta(1-\alpha)^{1/n}\right)=\alpha,\] we
obtain, after observing \(M=m\),
\[m\le \theta\le \frac{m}{(1-\alpha)^{1/n}}.\] Thus an interval of the
form \[\left[m,\frac{m}{(1-\alpha)^{1/n}}\right]\] comes from solving
the CDF inequality.

A more common \(100(1-\alpha)\%\) one-sided upper confidence interval is
\[\left[M,\frac{M}{\alpha^{1/n}}\right],\] because
\[\mathbb{P}_\theta\left(\theta\le \frac{M}{\alpha^{1/n}}\right)=\mathbb{P}_\theta(M\ge \alpha^{1/n}\theta)=1-\alpha.\]
The exact endpoint depends on how the tail probability is allocated.

\end{tcolorbox}

\begin{tcolorbox}[enhanced jigsaw, arc=.35mm, bottomrule=.15mm, bottomtitle=1mm, breakable, colback=white, colbacktitle=quarto-callout-note-color!10!white, colframe=quarto-callout-note-color-frame, coltitle=black, left=2mm, leftrule=.75mm, opacityback=0, opacitybacktitle=0.6, rightrule=.15mm, title=\textcolor{quarto-callout-note-color}{\faInfo}\hspace{0.5em}{Remark}, titlerule=0mm, toprule=.15mm, toptitle=1mm]

\emph{Remark 17}. The statistic \(M=\max_i X_i\) is natural here because
all observations must be less than or equal to \(\theta\). This is also
the sufficient statistic for \(\theta\) in the uniform endpoint model.

\end{tcolorbox}

\section{Exponential distribution}\label{exponential-distribution-2}

This subsection derives an exact confidence interval for the exponential
rate and for the mean lifetime.

\begin{tcolorbox}[enhanced jigsaw, arc=.35mm, bottomrule=.15mm, bottomtitle=1mm, breakable, colback=white, colbacktitle=quarto-callout-tip-color!10!white, colframe=quarto-callout-tip-color-frame, coltitle=black, left=2mm, leftrule=.75mm, opacityback=0, opacitybacktitle=0.6, rightrule=.15mm, title=\textcolor{quarto-callout-tip-color}{\faLightbulb}\hspace{0.5em}{Example}, titlerule=0mm, toprule=.15mm, toptitle=1mm]

\textbf{Example 18} (Exponential rate and mean lifetime). Suppose
\(X_1,\ldots,X_n\sim \operatorname{Exp}(\lambda)\) independently, where
the density is \[f(x\mid \lambda)=\lambda e^{-\lambda x},\qquad x>0.\]
Let \[Y=\sum_{i=1}^n X_i.\] Find a confidence interval for \(\lambda\),
and then for the mean lifetime \(\theta=1/\lambda\).

\end{tcolorbox}

\begin{tcolorbox}[enhanced jigsaw, arc=.35mm, bottomrule=.15mm, bottomtitle=1mm, breakable, colback=white, colbacktitle=quarto-callout-note-color!10!white, colframe=quarto-callout-note-color-frame, coltitle=black, left=2mm, leftrule=.75mm, opacityback=0, opacitybacktitle=0.6, rightrule=.15mm, title=\textcolor{quarto-callout-note-color}{\faInfo}\hspace{0.5em}{Solution}, titlerule=0mm, toprule=.15mm, toptitle=1mm]

The sum satisfies \[Y\sim \operatorname{Gamma}(n,\lambda),\] where
\(\lambda\) is the rate parameter. Equivalently,
\[2\lambda Y\sim \chi^2_{2n}.\] Thus
\[\mathbb{P}\left(\chi^2_{2n,\alpha/2}\le 2\lambda Y\le \chi^2_{2n,1-\alpha/2}\right)=1-\alpha.\]
Solving for \(\lambda\) gives \[\lambda\in \left[
    \frac{\chi^2_{2n,\alpha/2}}{2Y},\;
    \frac{\chi^2_{2n,1-\alpha/2}}{2Y}
    \right].\] If \(\theta=1/\lambda\) is the mean lifetime, then taking
reciprocals reverses the endpoints: \[\theta\in \left[
    \frac{2Y}{\chi^2_{2n,1-\alpha/2}},\;
    \frac{2Y}{\chi^2_{2n,\alpha/2}}
    \right].\]

\end{tcolorbox}

\bookmarksetup{startatroot}

\chapter{Method 4: Bayesian Credible
Intervals}\label{method-4-bayesian-credible-intervals}

This section presents the Bayesian counterpart of confidence intervals.

In classical frequentist statistics, the parameter \(\theta\) is fixed
and the interval is random. Therefore, after observing a fixed interval
such as \([3,10]\), it is not correct in the frequentist sense to say
``there is a \(90\%\) probability that \(\theta\) lies in \([3,10]\).''
The frequentist statement is about long-run coverage of the procedure.

In Bayesian statistics, parameters are treated as random variables. We
place a prior distribution on \(\theta\) and update it to a posterior
distribution after observing data. Then it is meaningful to say that
there is a posterior probability that \(\theta\) lies in an interval.

\begin{tcolorbox}[enhanced jigsaw, arc=.35mm, bottomrule=.15mm, bottomtitle=1mm, breakable, colback=white, colbacktitle=quarto-callout-note-color!10!white, colframe=quarto-callout-note-color-frame, coltitle=black, left=2mm, leftrule=.75mm, opacityback=0, opacitybacktitle=0.6, rightrule=.15mm, title=\textcolor{quarto-callout-note-color}{\faInfo}\hspace{0.5em}{Definition}, titlerule=0mm, toprule=.15mm, toptitle=1mm]

\textbf{Definition 19} (Credible interval). Let \(\pi(\theta\mid x)\) be
the posterior density of \(\theta\) given data \(x\). A set
\(A\subseteq \Theta\) is a \(100(1-\alpha)\%\) \textbf{credible set} if
\[\mathbb{P}(\theta\in A\mid x)=\int_A \pi(\theta\mid x)\,d\theta=1-\alpha.\]
If \(A=[a,b]\), then \([a,b]\) is a credible interval.

\end{tcolorbox}

\section{General Bayesian procedure}\label{general-bayesian-procedure}

This subsection summarizes how to build Bayesian intervals from a
posterior distribution.

Choose a prior distribution \(\pi(\theta)\) and combine it with the
likelihood \(f(x\mid \theta)\). The posterior distribution is
\[\pi(\theta\mid x)=\frac{f(x\mid \theta)\pi(\theta)}{\int f(x\mid \theta')\pi(\theta')\,d\theta'}.\]
Then choose \(a\) and \(b\) so that
\[\int_a^b \pi(\theta\mid x)\,d\theta=1-\alpha.\] A common choice is the
equal-tail credible interval, where
\[\mathbb{P}(\theta<a\mid x)=\frac{\alpha}{2},
    \qquad
    \mathbb{P}(\theta>b\mid x)=\frac{\alpha}{2}.\]

\begin{tcolorbox}[enhanced jigsaw, arc=.35mm, bottomrule=.15mm, bottomtitle=1mm, breakable, colback=white, colbacktitle=quarto-callout-note-color!10!white, colframe=quarto-callout-note-color-frame, coltitle=black, left=2mm, leftrule=.75mm, opacityback=0, opacitybacktitle=0.6, rightrule=.15mm, title=\textcolor{quarto-callout-note-color}{\faInfo}\hspace{0.5em}{Remark}, titlerule=0mm, toprule=.15mm, toptitle=1mm]

\emph{Remark 20}. Another common Bayesian interval is the highest
posterior density interval, which contains points with the largest
posterior density. Equal-tail and highest posterior density intervals
may differ for skewed posterior distributions.

\end{tcolorbox}

\section{Binomial data with beta
prior}\label{binomial-data-with-beta-prior}

This subsection presents the beta-binomial credible interval.

\begin{tcolorbox}[enhanced jigsaw, arc=.35mm, bottomrule=.15mm, bottomtitle=1mm, breakable, colback=white, colbacktitle=quarto-callout-tip-color!10!white, colframe=quarto-callout-tip-color-frame, coltitle=black, left=2mm, leftrule=.75mm, opacityback=0, opacitybacktitle=0.6, rightrule=.15mm, title=\textcolor{quarto-callout-tip-color}{\faLightbulb}\hspace{0.5em}{Example}, titlerule=0mm, toprule=.15mm, toptitle=1mm]

\textbf{Example 21} (Binomial data with beta prior). Suppose
\[X\sim \operatorname{Binomial}(n,p),\] and use the prior
\[p\sim \operatorname{Beta}(\alpha,\beta).\] Find the posterior
distribution and describe a credible interval for \(p\).

\end{tcolorbox}

\begin{tcolorbox}[enhanced jigsaw, arc=.35mm, bottomrule=.15mm, bottomtitle=1mm, breakable, colback=white, colbacktitle=quarto-callout-note-color!10!white, colframe=quarto-callout-note-color-frame, coltitle=black, left=2mm, leftrule=.75mm, opacityback=0, opacitybacktitle=0.6, rightrule=.15mm, title=\textcolor{quarto-callout-note-color}{\faInfo}\hspace{0.5em}{Solution}, titlerule=0mm, toprule=.15mm, toptitle=1mm]

The likelihood is proportional to \[p^x(1-p)^{n-x}.\] The beta prior
density is proportional to \[p^{\alpha-1}(1-p)^{\beta-1}.\] Therefore
the posterior density is proportional to
\[p^{\alpha+x-1}(1-p)^{\beta+n-x-1}.\] Hence
\[p\mid X=x\sim \operatorname{Beta}(\alpha+x,\beta+n-x).\] A
\(100(1-\alpha_0)\%\) equal-tail credible interval is
\[\left[q_{\alpha_0/2},q_{1-\alpha_0/2}\right],\] where \(q_\gamma\) is
the \(\gamma\) quantile of the
\(\operatorname{Beta}(\alpha+x,\beta+n-x)\) distribution.

\end{tcolorbox}

\begin{tcolorbox}[enhanced jigsaw, arc=.35mm, bottomrule=.15mm, bottomtitle=1mm, breakable, colback=white, colbacktitle=quarto-callout-tip-color!10!white, colframe=quarto-callout-tip-color-frame, coltitle=black, left=2mm, leftrule=.75mm, opacityback=0, opacitybacktitle=0.6, rightrule=.15mm, title=\textcolor{quarto-callout-tip-color}{\faLightbulb}\hspace{0.5em}{Example}, titlerule=0mm, toprule=.15mm, toptitle=1mm]

\textbf{Example 22} (Numerical beta-binomial credible interval). Suppose
\(n=20\), \(x=12\), and the prior is \(p\sim \operatorname{Beta}(2,2)\).
Find the posterior distribution and state the approximate \(95\%\)
credible interval given in the lecture notes.

\end{tcolorbox}

\begin{tcolorbox}[enhanced jigsaw, arc=.35mm, bottomrule=.15mm, bottomtitle=1mm, breakable, colback=white, colbacktitle=quarto-callout-note-color!10!white, colframe=quarto-callout-note-color-frame, coltitle=black, left=2mm, leftrule=.75mm, opacityback=0, opacitybacktitle=0.6, rightrule=.15mm, title=\textcolor{quarto-callout-note-color}{\faInfo}\hspace{0.5em}{Solution}, titlerule=0mm, toprule=.15mm, toptitle=1mm]

The posterior is
\[p\mid X=12\sim \operatorname{Beta}(2+12,2+20-12)=\operatorname{Beta}(14,10).\]
The equal-tail \(95\%\) credible interval is given by the \(0.025\) and
\(0.975\) posterior quantiles. The lecture-note computation gives
approximately \[p\in [0.385,0.768].\] This means that, under the beta
prior and the observed data, the posterior probability that \(p\) lies
in this interval is \(0.95\).

\end{tcolorbox}

\section{Normal mean with known variance and normal
prior}\label{normal-mean-with-known-variance-and-normal-prior}

This subsection studies a conjugate Bayesian interval for a normal mean.

\begin{tcolorbox}[enhanced jigsaw, arc=.35mm, bottomrule=.15mm, bottomtitle=1mm, breakable, colback=white, colbacktitle=quarto-callout-tip-color!10!white, colframe=quarto-callout-tip-color-frame, coltitle=black, left=2mm, leftrule=.75mm, opacityback=0, opacitybacktitle=0.6, rightrule=.15mm, title=\textcolor{quarto-callout-tip-color}{\faLightbulb}\hspace{0.5em}{Example}, titlerule=0mm, toprule=.15mm, toptitle=1mm]

\textbf{Example 23} (Normal-normal model). Suppose
\(X_1,\ldots,X_n\sim \operatorname{Normal}(\mu,\sigma^2)\)
independently, where \(\sigma^2\) is known. Put a normal prior on
\(\mu\): \[\mu\sim \operatorname{Normal}(\theta,\tau^2).\] Find the
posterior mean and variance, and describe a credible interval.

\end{tcolorbox}

\begin{tcolorbox}[enhanced jigsaw, arc=.35mm, bottomrule=.15mm, bottomtitle=1mm, breakable, colback=white, colbacktitle=quarto-callout-note-color!10!white, colframe=quarto-callout-note-color-frame, coltitle=black, left=2mm, leftrule=.75mm, opacityback=0, opacitybacktitle=0.6, rightrule=.15mm, title=\textcolor{quarto-callout-note-color}{\faInfo}\hspace{0.5em}{Solution}, titlerule=0mm, toprule=.15mm, toptitle=1mm]

The likelihood is normal in \(\mu\), and the normal prior is conjugate.
The posterior is also normal:
\[\mu\mid x\sim \operatorname{Normal}(m_n,v_n),\] where
\[m_n=\frac{\tau^2\sum_{i=1}^n x_i+\sigma^2\theta}{n\tau^2+\sigma^2}\]
and \[v_n=\frac{\sigma^2\tau^2}{\sigma^2+n\tau^2}.\] Therefore a
\(100(1-\alpha)\%\) equal-tail credible interval is
\[\left[m_n-z_{1-\alpha/2}\sqrt{v_n},\;
    m_n+z_{1-\alpha/2}\sqrt{v_n}\right].\]

\end{tcolorbox}

\begin{tcolorbox}[enhanced jigsaw, arc=.35mm, bottomrule=.15mm, bottomtitle=1mm, breakable, colback=white, colbacktitle=quarto-callout-tip-color!10!white, colframe=quarto-callout-tip-color-frame, coltitle=black, left=2mm, leftrule=.75mm, opacityback=0, opacitybacktitle=0.6, rightrule=.15mm, title=\textcolor{quarto-callout-tip-color}{\faLightbulb}\hspace{0.5em}{Key idea}, titlerule=0mm, toprule=.15mm, toptitle=1mm]

Weighted average interpretation The posterior mean is a weighted average
of the prior mean \(\theta\) and the sample mean \(\bar x\):
\[m_n=\frac{n\tau^2}{n\tau^2+\sigma^2}\bar x+
    \frac{\sigma^2}{n\tau^2+\sigma^2}\theta.\] As \(n\) grows, the
sample mean receives more weight.

\end{tcolorbox}

\section{Poisson data with gamma
prior}\label{poisson-data-with-gamma-prior}

This subsection presents the gamma-Poisson credible interval.

\begin{tcolorbox}[enhanced jigsaw, arc=.35mm, bottomrule=.15mm, bottomtitle=1mm, breakable, colback=white, colbacktitle=quarto-callout-tip-color!10!white, colframe=quarto-callout-tip-color-frame, coltitle=black, left=2mm, leftrule=.75mm, opacityback=0, opacitybacktitle=0.6, rightrule=.15mm, title=\textcolor{quarto-callout-tip-color}{\faLightbulb}\hspace{0.5em}{Example}, titlerule=0mm, toprule=.15mm, toptitle=1mm]

\textbf{Example 24} (Poisson data). Suppose
\(X_1,\ldots,X_n\sim \operatorname{Poisson}(\lambda)\) independently.
Use the prior \[\lambda\sim \operatorname{Gamma}(\alpha,\beta),\] where
\(\beta\) is the rate parameter. Find the posterior distribution.

\end{tcolorbox}

\begin{tcolorbox}[enhanced jigsaw, arc=.35mm, bottomrule=.15mm, bottomtitle=1mm, breakable, colback=white, colbacktitle=quarto-callout-note-color!10!white, colframe=quarto-callout-note-color-frame, coltitle=black, left=2mm, leftrule=.75mm, opacityback=0, opacitybacktitle=0.6, rightrule=.15mm, title=\textcolor{quarto-callout-note-color}{\faInfo}\hspace{0.5em}{Solution}, titlerule=0mm, toprule=.15mm, toptitle=1mm]

The likelihood is proportional to \[e^{-n\lambda}\lambda^{\sum_i x_i}.\]
The gamma prior density is proportional to
\[\lambda^{\alpha-1}e^{-\beta\lambda}.\] Thus the posterior density is
proportional to \[\lambda^{\alpha+\\sum_i x_i-1}e^{-(\beta+n)\lambda}.\]
Therefore
\[\lambda\mid X\sim \operatorname{Gamma}\left(\alpha+\sum_{i=1}^n x_i,\; \beta+n\right).\]
A credible interval for \(\lambda\) is obtained from the corresponding
gamma posterior quantiles.

\end{tcolorbox}

\begin{tcolorbox}[enhanced jigsaw, arc=.35mm, bottomrule=.15mm, bottomtitle=1mm, breakable, colback=white, colbacktitle=quarto-callout-tip-color!10!white, colframe=quarto-callout-tip-color-frame, coltitle=black, left=2mm, leftrule=.75mm, opacityback=0, opacitybacktitle=0.6, rightrule=.15mm, title=\textcolor{quarto-callout-tip-color}{\faLightbulb}\hspace{0.5em}{Example}, titlerule=0mm, toprule=.15mm, toptitle=1mm]

\textbf{Example 25} (Numerical gamma-Poisson posterior). Suppose
\(n=10\), the observed counts sum to \(\sum_i x_i=26\), and the prior is
\[\lambda\sim \operatorname{Gamma}(2,1).\] Find the posterior
distribution.

\end{tcolorbox}

\begin{tcolorbox}[enhanced jigsaw, arc=.35mm, bottomrule=.15mm, bottomtitle=1mm, breakable, colback=white, colbacktitle=quarto-callout-note-color!10!white, colframe=quarto-callout-note-color-frame, coltitle=black, left=2mm, leftrule=.75mm, opacityback=0, opacitybacktitle=0.6, rightrule=.15mm, title=\textcolor{quarto-callout-note-color}{\faInfo}\hspace{0.5em}{Solution}, titlerule=0mm, toprule=.15mm, toptitle=1mm]

Using the gamma-Poisson update,
\[\lambda\mid X\sim \operatorname{Gamma}(2+26,1+10)=\operatorname{Gamma}(28,11).\]
A \(95\%\) credible interval is given by the \(0.025\) and \(0.975\)
quantiles of this gamma distribution.

\end{tcolorbox}

\bookmarksetup{startatroot}

\chapter{Comparing Confidence Intervals and Credible
Intervals}\label{comparing-confidence-intervals-and-credible-intervals}

This section clarifies the difference between frequentist and Bayesian
interval statements.

\begin{tcolorbox}[enhanced jigsaw, arc=.35mm, bottomrule=.15mm, bottomtitle=1mm, breakable, colback=white, colbacktitle=quarto-callout-note-color!10!white, colframe=quarto-callout-note-color-frame, coltitle=black, left=2mm, leftrule=.75mm, opacityback=0, opacitybacktitle=0.6, rightrule=.15mm, title=\textcolor{quarto-callout-note-color}{\faInfo}\hspace{0.5em}{Note}, titlerule=0mm, toprule=.15mm, toptitle=1mm]

{\def\LTcaptype{none} % do not increment counter
\begin{longtable}[]{@{}lll@{}}
\toprule\noalign{}
& Frequentist confidence interval & Bayesian credible interval \\
\midrule\noalign{}
\endhead
\bottomrule\noalign{}
\endlastfoot
Parameter & Fixed unknown constant & Random variable with
prior/posterior distribution \\
Interval & Random before data; fixed after data & Fixed after data, with
posterior probability statement \\
Probability statement & Long-run coverage of the procedure & Posterior
probability of parameter lying in interval \\
Inputs & Sampling distribution & Likelihood plus prior \\
Dependence on prior & No prior needed & Depends on chosen prior \\
\end{longtable}
}

\end{tcolorbox}

\begin{tcolorbox}[enhanced jigsaw, arc=.35mm, bottomrule=.15mm, bottomtitle=1mm, breakable, colback=white, colbacktitle=quarto-callout-warning-color!10!white, colframe=quarto-callout-warning-color-frame, coltitle=black, left=2mm, leftrule=.75mm, opacityback=0, opacitybacktitle=0.6, rightrule=.15mm, title=\textcolor{quarto-callout-warning-color}{\faExclamationTriangle}\hspace{0.5em}{Warning}, titlerule=0mm, toprule=.15mm, toptitle=1mm]

Common interpretation mistake For a frequentist \(90\%\) confidence
interval \([3,10]\), it is not correct to say that
\(\mathbb{P}(\theta\in[3,10])=0.90\) after observing the data. In the
Bayesian setting, a \(90\%\) credible interval does allow that posterior
probability statement, but it depends on the prior model.

\end{tcolorbox}

\bookmarksetup{startatroot}

\chapter{Practice Problems}\label{practice-problems-7}

This section gives practice problems that reinforce interval
construction by test inversion, pivots, CDF pivoting, and Bayesian
posterior intervals.

\begin{tcolorbox}[enhanced jigsaw, arc=.35mm, bottomrule=.15mm, bottomtitle=1mm, breakable, colback=white, colbacktitle=quarto-callout-caution-color!10!white, colframe=quarto-callout-caution-color-frame, coltitle=black, left=2mm, leftrule=.75mm, opacityback=0, opacitybacktitle=0.6, rightrule=.15mm, title=\textcolor{quarto-callout-caution-color}{\faFire}\hspace{0.5em}{Practice Problem}, titlerule=0mm, toprule=.15mm, toptitle=1mm]

\textbf{Practice Problem 26} (Normal mean, known variance). Suppose
\(X_1,\ldots,X_{25}\sim \operatorname{Normal}(\mu,9)\) independently and
\(\bar x=10.4\). Find a \(95\%\) confidence interval for \(\mu\).

\end{tcolorbox}

\begin{tcolorbox}[enhanced jigsaw, arc=.35mm, bottomrule=.15mm, bottomtitle=1mm, breakable, colback=white, colbacktitle=quarto-callout-note-color!10!white, colframe=quarto-callout-note-color-frame, coltitle=black, left=2mm, leftrule=.75mm, opacityback=0, opacitybacktitle=0.6, rightrule=.15mm, title=\textcolor{quarto-callout-note-color}{\faInfo}\hspace{0.5em}{Solution}, titlerule=0mm, toprule=.15mm, toptitle=1mm]

Here \(n=25\), \(\sigma=3\), and \(z_{0.975}=1.96\). The interval is
\[\bar x\pm z_{0.975}\frac{\sigma}{\sqrt n}
    =10.4\pm 1.96\frac{3}{5}
    =10.4\pm 1.176.\] Thus \[\mu\in [9.224,11.576].\]

\end{tcolorbox}

\begin{tcolorbox}[enhanced jigsaw, arc=.35mm, bottomrule=.15mm, bottomtitle=1mm, breakable, colback=white, colbacktitle=quarto-callout-caution-color!10!white, colframe=quarto-callout-caution-color-frame, coltitle=black, left=2mm, leftrule=.75mm, opacityback=0, opacitybacktitle=0.6, rightrule=.15mm, title=\textcolor{quarto-callout-caution-color}{\faFire}\hspace{0.5em}{Practice Problem}, titlerule=0mm, toprule=.15mm, toptitle=1mm]

\textbf{Practice Problem 27} (Normal mean, unknown variance). Suppose
\(X_1,\ldots,X_{16}\sim \operatorname{Normal}(\mu,\sigma^2)\),
\(\bar x=5.2\), and \(s=1.6\). Write the \(95\%\) confidence interval
for \(\mu\) in terms of the appropriate \(t\) quantile.

\end{tcolorbox}

\begin{tcolorbox}[enhanced jigsaw, arc=.35mm, bottomrule=.15mm, bottomtitle=1mm, breakable, colback=white, colbacktitle=quarto-callout-note-color!10!white, colframe=quarto-callout-note-color-frame, coltitle=black, left=2mm, leftrule=.75mm, opacityback=0, opacitybacktitle=0.6, rightrule=.15mm, title=\textcolor{quarto-callout-note-color}{\faInfo}\hspace{0.5em}{Solution}, titlerule=0mm, toprule=.15mm, toptitle=1mm]

The pivot is \[T=\frac{\bar X-\mu}{S/\sqrt n}\sim t_{15}.\] Therefore
the \(95\%\) confidence interval is
\[5.2\pm t_{15,0.975}\frac{1.6}{4}.\] That is,
\[\mu\in \left[5.2-0.4t_{15,0.975},\;5.2+0.4t_{15,0.975}\right].\]

\end{tcolorbox}

\begin{tcolorbox}[enhanced jigsaw, arc=.35mm, bottomrule=.15mm, bottomtitle=1mm, breakable, colback=white, colbacktitle=quarto-callout-caution-color!10!white, colframe=quarto-callout-caution-color-frame, coltitle=black, left=2mm, leftrule=.75mm, opacityback=0, opacitybacktitle=0.6, rightrule=.15mm, title=\textcolor{quarto-callout-caution-color}{\faFire}\hspace{0.5em}{Practice Problem}, titlerule=0mm, toprule=.15mm, toptitle=1mm]

\textbf{Practice Problem 28} (Normal variance). Suppose
\(X_1,\ldots,X_{10}\sim \operatorname{Normal}(\mu,\sigma^2)\) and the
observed sample variance is \(s^2=4\). Write a \(95\%\) confidence
interval for \(\sigma^2\).

\end{tcolorbox}

\begin{tcolorbox}[enhanced jigsaw, arc=.35mm, bottomrule=.15mm, bottomtitle=1mm, breakable, colback=white, colbacktitle=quarto-callout-note-color!10!white, colframe=quarto-callout-note-color-frame, coltitle=black, left=2mm, leftrule=.75mm, opacityback=0, opacitybacktitle=0.6, rightrule=.15mm, title=\textcolor{quarto-callout-note-color}{\faInfo}\hspace{0.5em}{Solution}, titlerule=0mm, toprule=.15mm, toptitle=1mm]

Here \(n=10\), so \(n-1=9\). The chi-square pivot is
\[\frac{(n-1)S^2}{\sigma^2}\sim \chi^2_9.\] The \(95\%\) interval is
\[\left[
    \frac{9\cdot 4}{\chi^2_{9,0.975}},\;
    \frac{9\cdot 4}{\chi^2_{9,0.025}}
    \right]
    =\left[
    \frac{36}{\chi^2_{9,0.975}},\;
    \frac{36}{\chi^2_{9,0.025}}
    \right].\]

\end{tcolorbox}

\begin{tcolorbox}[enhanced jigsaw, arc=.35mm, bottomrule=.15mm, bottomtitle=1mm, breakable, colback=white, colbacktitle=quarto-callout-caution-color!10!white, colframe=quarto-callout-caution-color-frame, coltitle=black, left=2mm, leftrule=.75mm, opacityback=0, opacitybacktitle=0.6, rightrule=.15mm, title=\textcolor{quarto-callout-caution-color}{\faFire}\hspace{0.5em}{Practice Problem}, titlerule=0mm, toprule=.15mm, toptitle=1mm]

\textbf{Practice Problem 29} (Approximate binomial proportion interval).
Suppose \(X\sim \operatorname{Binomial}(200,p)\) and \(x=126\). Use the
Wald method to construct an approximate \(95\%\) confidence interval for
\(p\).

\end{tcolorbox}

\begin{tcolorbox}[enhanced jigsaw, arc=.35mm, bottomrule=.15mm, bottomtitle=1mm, breakable, colback=white, colbacktitle=quarto-callout-note-color!10!white, colframe=quarto-callout-note-color-frame, coltitle=black, left=2mm, leftrule=.75mm, opacityback=0, opacitybacktitle=0.6, rightrule=.15mm, title=\textcolor{quarto-callout-note-color}{\faInfo}\hspace{0.5em}{Solution}, titlerule=0mm, toprule=.15mm, toptitle=1mm]

The sample proportion is \[\hat p=\frac{126}{200}=0.63.\] The estimated
standard error is \[\sqrt{\frac{\hat p(1-\hat p)}{n}}
    =\sqrt{\frac{0.63(0.37)}{200}}
    \approx 0.0341.\] Using \(z_{0.975}=1.96\), the margin of error is
\[1.96(0.0341)\approx 0.0668.\] Thus the approximate interval is
\[p\in [0.5632,0.6968].\]

\end{tcolorbox}

\begin{tcolorbox}[enhanced jigsaw, arc=.35mm, bottomrule=.15mm, bottomtitle=1mm, breakable, colback=white, colbacktitle=quarto-callout-caution-color!10!white, colframe=quarto-callout-caution-color-frame, coltitle=black, left=2mm, leftrule=.75mm, opacityback=0, opacitybacktitle=0.6, rightrule=.15mm, title=\textcolor{quarto-callout-caution-color}{\faFire}\hspace{0.5em}{Practice Problem}, titlerule=0mm, toprule=.15mm, toptitle=1mm]

\textbf{Practice Problem 30} (Uniform endpoint). Suppose
\(X_1,\ldots,X_8\sim \operatorname{Uniform}(0,\theta)\) and the observed
maximum is \(m=12\). Give a \(95\%\) one-sided upper confidence interval
of the form \([m,m/\alpha^{1/n}]\).

\end{tcolorbox}

\begin{tcolorbox}[enhanced jigsaw, arc=.35mm, bottomrule=.15mm, bottomtitle=1mm, breakable, colback=white, colbacktitle=quarto-callout-note-color!10!white, colframe=quarto-callout-note-color-frame, coltitle=black, left=2mm, leftrule=.75mm, opacityback=0, opacitybacktitle=0.6, rightrule=.15mm, title=\textcolor{quarto-callout-note-color}{\faInfo}\hspace{0.5em}{Solution}, titlerule=0mm, toprule=.15mm, toptitle=1mm]

Here \(n=8\) and \(\alpha=0.05\). The interval is
\[\left[12,\frac{12}{0.05^{1/8}}\right].\] This interval has coverage
\(0.95\) because
\[\mathbb{P}_\theta\left(\theta\le \frac{M}{0.05^{1/8}}\right)
    =\mathbb{P}_\theta(M\ge 0.05^{1/8}\theta)
    =1-0.05.\]

\end{tcolorbox}

\begin{tcolorbox}[enhanced jigsaw, arc=.35mm, bottomrule=.15mm, bottomtitle=1mm, breakable, colback=white, colbacktitle=quarto-callout-caution-color!10!white, colframe=quarto-callout-caution-color-frame, coltitle=black, left=2mm, leftrule=.75mm, opacityback=0, opacitybacktitle=0.6, rightrule=.15mm, title=\textcolor{quarto-callout-caution-color}{\faFire}\hspace{0.5em}{Practice Problem}, titlerule=0mm, toprule=.15mm, toptitle=1mm]

\textbf{Practice Problem 31} (Exponential mean lifetime). Suppose
\(X_1,\ldots,X_6\sim \operatorname{Exp}(\lambda)\) and the observed sum
is \(Y=18\). Write a \(90\%\) confidence interval for the mean lifetime
\(\theta=1/\lambda\).

\end{tcolorbox}

\begin{tcolorbox}[enhanced jigsaw, arc=.35mm, bottomrule=.15mm, bottomtitle=1mm, breakable, colback=white, colbacktitle=quarto-callout-note-color!10!white, colframe=quarto-callout-note-color-frame, coltitle=black, left=2mm, leftrule=.75mm, opacityback=0, opacitybacktitle=0.6, rightrule=.15mm, title=\textcolor{quarto-callout-note-color}{\faInfo}\hspace{0.5em}{Solution}, titlerule=0mm, toprule=.15mm, toptitle=1mm]

Here \(2n=12\) and \(\alpha=0.10\). From \[2\lambda Y\sim \chi^2_{12},\]
we obtain \[\theta=\frac{1}{\lambda}\in
    \left[
    \frac{2Y}{\chi^2_{12,1-\alpha/2}},\;
    \frac{2Y}{\chi^2_{12,\alpha/2}}
    \right].\] Thus \[\theta\in
    \left[
    \frac{36}{\chi^2_{12,0.95}},\;
    \frac{36}{\chi^2_{12,0.05}}
    \right].\]

\end{tcolorbox}

\begin{tcolorbox}[enhanced jigsaw, arc=.35mm, bottomrule=.15mm, bottomtitle=1mm, breakable, colback=white, colbacktitle=quarto-callout-caution-color!10!white, colframe=quarto-callout-caution-color-frame, coltitle=black, left=2mm, leftrule=.75mm, opacityback=0, opacitybacktitle=0.6, rightrule=.15mm, title=\textcolor{quarto-callout-caution-color}{\faFire}\hspace{0.5em}{Practice Problem}, titlerule=0mm, toprule=.15mm, toptitle=1mm]

\textbf{Practice Problem 32} (Beta-binomial credible interval). Suppose
\(X\sim \operatorname{Binomial}(30,p)\), \(x=18\), and
\(p\sim \operatorname{Beta}(3,3)\). Find the posterior distribution and
describe a \(95\%\) equal-tail credible interval.

\end{tcolorbox}

\begin{tcolorbox}[enhanced jigsaw, arc=.35mm, bottomrule=.15mm, bottomtitle=1mm, breakable, colback=white, colbacktitle=quarto-callout-note-color!10!white, colframe=quarto-callout-note-color-frame, coltitle=black, left=2mm, leftrule=.75mm, opacityback=0, opacitybacktitle=0.6, rightrule=.15mm, title=\textcolor{quarto-callout-note-color}{\faInfo}\hspace{0.5em}{Solution}, titlerule=0mm, toprule=.15mm, toptitle=1mm]

The posterior is
\[p\mid X=18\sim \operatorname{Beta}(3+18,3+30-18)=\operatorname{Beta}(21,15).\]
A \(95\%\) equal-tail credible interval is
\[\left[q_{0.025},q_{0.975}\right],\] where \(q_\gamma\) is the
\(\gamma\) quantile of the \(\operatorname{Beta}(21,15)\) distribution.

\end{tcolorbox}

\begin{tcolorbox}[enhanced jigsaw, arc=.35mm, bottomrule=.15mm, bottomtitle=1mm, breakable, colback=white, colbacktitle=quarto-callout-caution-color!10!white, colframe=quarto-callout-caution-color-frame, coltitle=black, left=2mm, leftrule=.75mm, opacityback=0, opacitybacktitle=0.6, rightrule=.15mm, title=\textcolor{quarto-callout-caution-color}{\faFire}\hspace{0.5em}{Practice Problem}, titlerule=0mm, toprule=.15mm, toptitle=1mm]

\textbf{Practice Problem 33} (Normal-normal credible interval). Suppose
\(X_1,\ldots,X_n\sim \operatorname{Normal}(\mu,\sigma^2)\) with known
\(\sigma^2\), and the prior is
\(\mu\sim \operatorname{Normal}(\theta,\tau^2)\). Show that the
posterior variance is smaller than both \(\tau^2\) and \(\sigma^2/n\).

\end{tcolorbox}

\begin{tcolorbox}[enhanced jigsaw, arc=.35mm, bottomrule=.15mm, bottomtitle=1mm, breakable, colback=white, colbacktitle=quarto-callout-note-color!10!white, colframe=quarto-callout-note-color-frame, coltitle=black, left=2mm, leftrule=.75mm, opacityback=0, opacitybacktitle=0.6, rightrule=.15mm, title=\textcolor{quarto-callout-note-color}{\faInfo}\hspace{0.5em}{Solution}, titlerule=0mm, toprule=.15mm, toptitle=1mm]

The posterior variance is
\[v_n=\frac{\sigma^2\tau^2}{\sigma^2+n\tau^2}.\] Since
\(\sigma^2+n\tau^2>\sigma^2\), \[v_n<\tau^2.\] Also, since
\(\sigma^2+n\tau^2>n\tau^2\),
\[v_n<\frac{\sigma^2\tau^2}{n\tau^2}=\frac{\sigma^2}{n}.\] Thus the
posterior variance is smaller than both the prior variance and the
sampling variance of \(\bar X\).

\end{tcolorbox}

\bookmarksetup{startatroot}

\chapter{Summary}\label{summary-9}

This section summarizes the main ideas of interval estimation.

\begin{tcolorbox}[enhanced jigsaw, arc=.35mm, bottomrule=.15mm, bottomtitle=1mm, breakable, colback=white, colbacktitle=quarto-callout-tip-color!10!white, colframe=quarto-callout-tip-color-frame, coltitle=black, left=2mm, leftrule=.75mm, opacityback=0, opacitybacktitle=0.6, rightrule=.15mm, title=\textcolor{quarto-callout-tip-color}{\faLightbulb}\hspace{0.5em}{Key idea}, titlerule=0mm, toprule=.15mm, toptitle=1mm]

Key takeaways

\begin{itemize}
\item
  A point estimator gives one number; an interval estimator gives a
  random range of plausible parameter values.
\item
  The coverage probability is
  \(\mathbb{P}_\theta(L(X)\le \theta\le U(X))\).
\item
  A confidence coefficient is the worst-case coverage over the parameter
  space.
\item
  Confidence intervals can be obtained by inverting hypothesis tests.
\item
  Pivotal quantities are functions of data and parameters with
  parameter-free distributions.
\item
  CDF pivoting is useful when the distribution of an order statistic or
  sufficient statistic is known.
\item
  Bayesian credible intervals are posterior probability statements and
  depend on the prior.
\end{itemize}

\end{tcolorbox}

\begin{tcolorbox}[enhanced jigsaw, arc=.35mm, bottomrule=.15mm, bottomtitle=1mm, breakable, colback=white, colbacktitle=quarto-callout-note-color!10!white, colframe=quarto-callout-note-color-frame, coltitle=black, left=2mm, leftrule=.75mm, opacityback=0, opacitybacktitle=0.6, rightrule=.15mm, title=\textcolor{quarto-callout-note-color}{\faInfo}\hspace{0.5em}{Note}, titlerule=0mm, toprule=.15mm, toptitle=1mm]

{\def\LTcaptype{none} % do not increment counter
\begin{longtable}[]{@{}lll@{}}
\toprule\noalign{}
Model & Key pivot/statistic & Interval idea \\
\midrule\noalign{}
\endhead
\bottomrule\noalign{}
\endlastfoot
Normal mean, known \(\sigma^2\) & \((\bar X-\mu)/(\sigma/\sqrt n)\) &
Normal \(z\) interval \\
Normal mean, unknown \(\sigma^2\) & \((\bar X-\mu)/(S/\sqrt n)\) &
Student \(t\) interval \\
Normal variance & \((n-1)S^2/\sigma^2\) & Chi-square interval \\
Binomial proportion & \((\hat p-p)/\sqrt{p(1-p)/n}\) & Approximate Wald
interval \\
Uniform endpoint & \(M=\max X_i\) & CDF pivoting \\
Exponential rate & \(2\lambda\sum X_i\) & Chi-square interval \\
Beta-binomial Bayes & Beta posterior & Posterior quantiles \\
Gamma-Poisson Bayes & Gamma posterior & Posterior quantiles \\
\end{longtable}
}

\end{tcolorbox}

\bookmarksetup{startatroot}

\chapter{Chapter 17: Interval Estimation II --- Evaluating Interval
Estimators}\label{chapter-17-interval-estimation-ii-evaluating-interval-estimators}

This chapter continues interval estimation. After constructing
confidence intervals and credible intervals, we now ask how to compare
them. The main theme is the tradeoff between \textbf{reliability} and
\textbf{precision}: a good interval should cover the true parameter with
high probability, but it should not be unnecessarily wide.

\begin{tcolorbox}[enhanced jigsaw, arc=.35mm, bottomrule=.15mm, bottomtitle=1mm, breakable, colback=white, colbacktitle=quarto-callout-note-color!10!white, colframe=quarto-callout-note-color-frame, coltitle=black, left=2mm, leftrule=.75mm, opacityback=0, opacitybacktitle=0.6, rightrule=.15mm, title=\textcolor{quarto-callout-note-color}{\faInfo}\hspace{0.5em}{Topics}, titlerule=0mm, toprule=.15mm, toptitle=1mm]

Coverage probability; interval length; expected length; shortest
intervals for unimodal densities; test-related optimality; false
coverage probability; uniformly most accurate confidence sets; unbiased
confidence sets; Bayesian HPD credible intervals; loss-function
optimality; risk of confidence sets.

\end{tcolorbox}

\bookmarksetup{startatroot}

\chapter{Overview}\label{overview-10}

This section moves from constructing confidence and credible intervals
to comparing them and deciding which interval is preferable.

In Section 16, we learned several ways to build interval estimators:
test inversion, pivotal quantities, pivoting the CDF, and Bayesian
credible intervals. In this section, we evaluate interval estimators
using four main ideas:

\begin{enumerate}
\def\labelenumi{\arabic{enumi}.}
\item
  size and coverage probability;
\item
  test-related optimality;
\item
  Bayesian optimality;
\item
  loss-function optimality.
\end{enumerate}

\begin{tcolorbox}[enhanced jigsaw, arc=.35mm, bottomrule=.15mm, bottomtitle=1mm, breakable, colback=white, colbacktitle=quarto-callout-tip-color!10!white, colframe=quarto-callout-tip-color-frame, coltitle=black, left=2mm, leftrule=.75mm, opacityback=0, opacitybacktitle=0.6, rightrule=.15mm, title=\textcolor{quarto-callout-tip-color}{\faLightbulb}\hspace{0.5em}{Key idea}, titlerule=0mm, toprule=.15mm, toptitle=1mm]

\textbf{Main idea.} A useful interval estimator should have high
probability of covering the true parameter, but it should not be
unnecessarily wide. Evaluation of interval estimators is therefore a
balance between \emph{reliability} and \emph{precision}.

\end{tcolorbox}

\bookmarksetup{startatroot}

\chapter{Size and Coverage
Probability}\label{size-and-coverage-probability}

This section introduces the two most basic numerical criteria for
evaluating an interval estimator: its coverage probability and its
length.

\section{Coverage probability}\label{coverage-probability}

Coverage probability measures how often a confidence interval procedure
contains the true parameter value in repeated sampling.

\begin{tcolorbox}[enhanced jigsaw, arc=.35mm, bottomrule=.15mm, bottomtitle=1mm, breakable, colback=white, colbacktitle=quarto-callout-note-color!10!white, colframe=quarto-callout-note-color-frame, coltitle=black, left=2mm, leftrule=.75mm, opacityback=0, opacitybacktitle=0.6, rightrule=.15mm, title=\textcolor{quarto-callout-note-color}{\faInfo}\hspace{0.5em}{Definition}, titlerule=0mm, toprule=.15mm, toptitle=1mm]

\textbf{Definition 1} (Coverage probability). Let \(C(X)\) be a
confidence set for a real parameter \(\theta\). The \textbf{coverage
probability} at \(\theta\) is \[\mathbb{P}_\theta\{\theta \in C(X)\}.\]
A confidence interval is designed so that
\[\mathbb{P}_\theta\{\theta \in C(X)\} \approx 1-\alpha.\] If the lower
endpoint is \(L(X)\) and the upper endpoint is \(U(X)\), then
\[C(X)=[L(X),U(X)],
    \qquad
    \mathbb{P}_\theta\{L(X)\leq \theta \leq U(X)\}\approx 1-\alpha.\]

\end{tcolorbox}

\begin{tcolorbox}[enhanced jigsaw, arc=.35mm, bottomrule=.15mm, bottomtitle=1mm, breakable, colback=white, colbacktitle=quarto-callout-note-color!10!white, colframe=quarto-callout-note-color-frame, coltitle=black, left=2mm, leftrule=.75mm, opacityback=0, opacitybacktitle=0.6, rightrule=.15mm, title=\textcolor{quarto-callout-note-color}{\faInfo}\hspace{0.5em}{Remark}, titlerule=0mm, toprule=.15mm, toptitle=1mm]

\emph{Remark 2}. In frequentist interval estimation, \(\theta\) is fixed
but unknown, and the interval \(C(X)\) is random because it depends on
the sample. Coverage probability is a long-run property of the
procedure, not a posterior probability statement about \(\theta\).

\end{tcolorbox}

\section{Size and expected length}\label{size-and-expected-length}

The size of an interval measures how much uncertainty remains after the
data have been observed.

\begin{tcolorbox}[enhanced jigsaw, arc=.35mm, bottomrule=.15mm, bottomtitle=1mm, breakable, colback=white, colbacktitle=quarto-callout-note-color!10!white, colframe=quarto-callout-note-color-frame, coltitle=black, left=2mm, leftrule=.75mm, opacityback=0, opacitybacktitle=0.6, rightrule=.15mm, title=\textcolor{quarto-callout-note-color}{\faInfo}\hspace{0.5em}{Definition}, titlerule=0mm, toprule=.15mm, toptitle=1mm]

\textbf{Definition 3} (Length and expected length). For an interval
estimator \(C(X)=[L(X),U(X)]\), the \textbf{length} is
\[\operatorname{Length}(C(X))=U(X)-L(X).\] The \textbf{expected length}
is \[\mathbb{E}_\theta[\operatorname{Length}(C(X))]
    =\mathbb{E}_\theta[U(X)-L(X)].\]

\end{tcolorbox}

\begin{tcolorbox}[enhanced jigsaw, arc=.35mm, bottomrule=.15mm, bottomtitle=1mm, breakable, colback=white, colbacktitle=quarto-callout-tip-color!10!white, colframe=quarto-callout-tip-color-frame, coltitle=black, left=2mm, leftrule=.75mm, opacityback=0, opacitybacktitle=0.6, rightrule=.15mm, title=\textcolor{quarto-callout-tip-color}{\faLightbulb}\hspace{0.5em}{Key idea}, titlerule=0mm, toprule=.15mm, toptitle=1mm]

\textbf{Precision principle.} Among confidence intervals with the same
coverage probability, the shorter interval is usually preferred because
it gives a more precise estimate of the parameter.

\end{tcolorbox}

\section{Normal mean example: length and
coverage}\label{normal-mean-example-length-and-coverage}

The normal mean example shows explicitly how coverage and length are
calculated.

\begin{tcolorbox}[enhanced jigsaw, arc=.35mm, bottomrule=.15mm, bottomtitle=1mm, breakable, colback=white, colbacktitle=quarto-callout-tip-color!10!white, colframe=quarto-callout-tip-color-frame, coltitle=black, left=2mm, leftrule=.75mm, opacityback=0, opacitybacktitle=0.6, rightrule=.15mm, title=\textcolor{quarto-callout-tip-color}{\faLightbulb}\hspace{0.5em}{Example}, titlerule=0mm, toprule=.15mm, toptitle=1mm]

\textbf{Example 4} (Normal confidence interval with known variance).
Suppose \[X_1,\ldots,X_n \sim \operatorname{Normal}(\mu,\sigma^2),\]
where \(\sigma^2\) is known. From the pivot construction,
\[Z=\frac{\bar X-\mu}{\sigma/\sqrt n}\sim \operatorname{Normal}(0,1).\]
For constants \(a<b\) satisfying \[\mathbb{P}(a\leq Z\leq b)=1-\alpha,\]
we obtain \[a \leq \frac{\bar X-\mu}{\sigma/\sqrt n}\leq b.\] Solving
for \(\mu\) gives
\[\bar X-b\frac{\sigma}{\sqrt n}\leq \mu\leq \bar X-a\frac{\sigma}{\sqrt n}.\]
Thus a \((1-\alpha)\) confidence interval is
\[C(X)=\left[\bar X-b\frac{\sigma}{\sqrt n},\; \bar X-a\frac{\sigma}{\sqrt n}\right].\]
Its length is
\[\operatorname{Length}(C(X))=(b-a)\frac{\sigma}{\sqrt n}.\]

\end{tcolorbox}

\begin{tcolorbox}[enhanced jigsaw, arc=.35mm, bottomrule=.15mm, bottomtitle=1mm, breakable, colback=white, colbacktitle=quarto-callout-note-color!10!white, colframe=quarto-callout-note-color-frame, coltitle=black, left=2mm, leftrule=.75mm, opacityback=0, opacitybacktitle=0.6, rightrule=.15mm, title=\textcolor{quarto-callout-note-color}{\faInfo}\hspace{0.5em}{Solution}, titlerule=0mm, toprule=.15mm, toptitle=1mm]

The coverage is \[\begin{aligned}
\mathbb{P}_\mu\{\mu\in C(X)\}
&=\mathbb{P}_\mu\left\{\bar X-b\frac{\sigma}{\sqrt n}\leq \mu\leq \bar X-a\frac{\sigma}{\sqrt n}\right\}\\
&=\mathbb{P}_\mu\left\{a\leq \frac{\bar X-\mu}{\sigma/\sqrt n}\leq b\right\}\\
&=\mathbb{P}(a\leq Z\leq b)=1-\alpha.
\end{aligned}\] Since \(a\) and \(b\) are constants, the length is
nonrandom: \[\operatorname{Length}(C(X))=(b-a)\frac{\sigma}{\sqrt n}.\]
Thus, for fixed \(\alpha\), \(\sigma\), and \(n\), minimizing the
interval length is equivalent to minimizing \(b-a\) subject to
\(\mathbb{P}(a\leq Z\leq b)=1-\alpha\).

\end{tcolorbox}

\section{Shortest interval for a normal
pivot}\label{shortest-interval-for-a-normal-pivot}

For a symmetric unimodal distribution such as the standard normal
distribution, the shortest interval with fixed probability is the
equal-tail central interval.

\begin{tcolorbox}[enhanced jigsaw, arc=.35mm, bottomrule=.15mm, bottomtitle=1mm, breakable, colback=white, colbacktitle=quarto-callout-tip-color!10!white, colframe=quarto-callout-tip-color-frame, coltitle=black, left=2mm, leftrule=.75mm, opacityback=0, opacitybacktitle=0.6, rightrule=.15mm, title=\textcolor{quarto-callout-tip-color}{\faLightbulb}\hspace{0.5em}{Example}, titlerule=0mm, toprule=.15mm, toptitle=1mm]

\textbf{Example 5} (Shortest normal interval). Suppose
\[X_1,\ldots,X_n\sim \operatorname{Normal}(\mu,\sigma^2),\] where
\(\sigma^2\) is known. The usual \(95\%\) confidence interval for
\(\mu\) is \[\bar X \pm z_{0.975}\frac{\sigma}{\sqrt n}.\] That is,
\[C(X)=\left[\bar X-z_{0.975}\frac{\sigma}{\sqrt n},\;\bar X+z_{0.975}\frac{\sigma}{\sqrt n}\right].\]
The coverage probability is exactly \(0.95\) for every \(\mu\), and the
length is \[2z_{0.975}\frac{\sigma}{\sqrt n}.\]

\end{tcolorbox}

\begin{tcolorbox}[enhanced jigsaw, arc=.35mm, bottomrule=.15mm, bottomtitle=1mm, breakable, colback=white, colbacktitle=quarto-callout-note-color!10!white, colframe=quarto-callout-note-color-frame, coltitle=black, left=2mm, leftrule=.75mm, opacityback=0, opacitybacktitle=0.6, rightrule=.15mm, title=\textcolor{quarto-callout-note-color}{\faInfo}\hspace{0.5em}{Solution}, titlerule=0mm, toprule=.15mm, toptitle=1mm]

Because
\[\frac{\bar X-\mu}{\sigma/\sqrt n}\sim \operatorname{Normal}(0,1),\] we
have
\[\mathbb{P}\left(-z_{0.975}\leq \frac{\bar X-\mu}{\sigma/\sqrt n}\leq z_{0.975}\right)=0.95.\]
Solving the inequalities for \(\mu\) gives the stated interval. Its
length is \[\left(\bar X+z_{0.975}\frac{\sigma}{\sqrt n}\right)
    -\left(\bar X-z_{0.975}\frac{\sigma}{\sqrt n}\right)
    =2z_{0.975}\frac{\sigma}{\sqrt n}.\] As \(n\) increases, this length
decreases like \(1/\sqrt n\), so larger samples give more precise
intervals.

\end{tcolorbox}

\bookmarksetup{startatroot}

\chapter{Shortest Intervals for a Unimodal
PDF}\label{shortest-intervals-for-a-unimodal-pdf}

This section gives a general rule for finding the shortest interval with
a prescribed probability when the density is unimodal.

\section{Equal boundary density
principle}\label{equal-boundary-density-principle}

For a unimodal density, the shortest interval containing a fixed
probability mass should cut the density at equal heights on the left and
right boundaries.

\begin{tcolorbox}[enhanced jigsaw, arc=.35mm, bottomrule=.15mm, bottomtitle=1mm, breakable, colback=white, colbacktitle=quarto-callout-important-color!10!white, colframe=quarto-callout-important-color-frame, coltitle=black, left=2mm, leftrule=.75mm, opacityback=0, opacitybacktitle=0.6, rightrule=.15mm, title=\textcolor{quarto-callout-important-color}{\faExclamation}\hspace{0.5em}{Theorem}, titlerule=0mm, toprule=.15mm, toptitle=1mm]

\textbf{Theorem 6} (Shortest interval with a unimodal density).
\emph{Let \(f(x)\) be a unimodal probability density function. Suppose
the mode is \(x^*\), and \(f\) is nondecreasing for \(x\leq x^*\) and
nonincreasing for \(x\geq x^*\).}

\emph{If an interval \([a,b]\) satisfies}

\begin{enumerate}
\def\labelenumi{\arabic{enumi}.}
\item
  \emph{\(\displaystyle \int_a^b f(x)\,dx=1-\alpha\);}
\item
  \emph{\(f(a)=f(b)>0\);}
\item
  \emph{\(a\leq x^*\leq b\);}
\end{enumerate}

\emph{then \([a,b]\) is the shortest interval among intervals having
probability \(1-\alpha\).}

\end{tcolorbox}

\begin{tcolorbox}[enhanced jigsaw, arc=.35mm, bottomrule=.15mm, bottomtitle=1mm, breakable, colback=white, colbacktitle=quarto-callout-note-color!10!white, colframe=quarto-callout-note-color-frame, coltitle=black, left=2mm, leftrule=.75mm, opacityback=0, opacitybacktitle=0.6, rightrule=.15mm, title=\textcolor{quarto-callout-note-color}{\faInfo}\hspace{0.5em}{Proof}, titlerule=0mm, toprule=.15mm, toptitle=1mm]

\emph{Idea of proof.} If the two boundary densities are not equal,
suppose for example that \(f(a)<f(b)\). Then we can move the left
endpoint slightly inward and move the right endpoint slightly outward to
preserve the same probability. Since probability is removed from a
lower-density region and added in a higher-density region, the right
endpoint needs to move less than the left endpoint moves. The total
length decreases. Therefore, at the shortest interval, the boundary
densities must be equal. The interval must also contain the mode;
otherwise shifting it toward the mode would increase included
probability without increasing length. ◻

\end{tcolorbox}

\begin{tcolorbox}[enhanced jigsaw, arc=.35mm, bottomrule=.15mm, bottomtitle=1mm, breakable, colback=white, colbacktitle=quarto-callout-note-color!10!white, colframe=quarto-callout-note-color-frame, coltitle=black, left=2mm, leftrule=.75mm, opacityback=0, opacitybacktitle=0.6, rightrule=.15mm, title=\textcolor{quarto-callout-note-color}{\faInfo}\hspace{0.5em}{Remark}, titlerule=0mm, toprule=.15mm, toptitle=1mm]

\emph{Remark 7}. For a symmetric unimodal density such as
\(\operatorname{Normal}(0,1)\), the equal boundary density condition
gives a symmetric interval \([-z,z]\). This recovers the usual
equal-tail normal interval.

\end{tcolorbox}

\begin{tcolorbox}[enhanced jigsaw, arc=.35mm, bottomrule=.15mm, bottomtitle=1mm, breakable, colback=white, colbacktitle=quarto-callout-warning-color!10!white, colframe=quarto-callout-warning-color-frame, coltitle=black, left=2mm, leftrule=.75mm, opacityback=0, opacitybacktitle=0.6, rightrule=.15mm, title=\textcolor{quarto-callout-warning-color}{\faExclamationTriangle}\hspace{0.5em}{Practice Problem}, titlerule=0mm, toprule=.15mm, toptitle=1mm]

\textbf{Practice Problem 8} (Shortest interval for a symmetric unimodal
density). Let \(Z\sim\operatorname{Normal}(0,1)\). Among all intervals
\([a,b]\) satisfying \(\mathbb{P}(a\leq Z\leq b)=0.95\), show that the
shortest interval is \([-z_{0.975},z_{0.975}]\).

\end{tcolorbox}

\begin{tcolorbox}[enhanced jigsaw, arc=.35mm, bottomrule=.15mm, bottomtitle=1mm, breakable, colback=white, colbacktitle=quarto-callout-note-color!10!white, colframe=quarto-callout-note-color-frame, coltitle=black, left=2mm, leftrule=.75mm, opacityback=0, opacitybacktitle=0.6, rightrule=.15mm, title=\textcolor{quarto-callout-note-color}{\faInfo}\hspace{0.5em}{Solution}, titlerule=0mm, toprule=.15mm, toptitle=1mm]

The standard normal density is symmetric about \(0\) and unimodal with
mode \(0\). By the shortest interval theorem, the shortest interval must
satisfy \[\phi(a)=\phi(b),\] and it must contain the mode \(0\). Since
\[\phi(x)=\frac{1}{\sqrt{2\pi}}e^{-x^2/2},\] \(\phi(a)=\phi(b)\) implies
\(a^2=b^2\). Because the interval contains \(0\) and \(a<b\), we must
have \(a=-b\). The coverage condition becomes
\[\mathbb{P}(-b\leq Z\leq b)=0.95,\] so \(b=z_{0.975}\) and
\(a=-z_{0.975}\).

\end{tcolorbox}

\bookmarksetup{startatroot}

\chapter{Test-Related Optimality}\label{test-related-optimality}

This section connects optimal confidence intervals with optimal
hypothesis tests.

\section{Confidence sets by inverting
tests}\label{confidence-sets-by-inverting-tests}

A confidence interval can be obtained by collecting all parameter values
that would not be rejected by a hypothesis test.

Suppose that for each \(\theta_0\), \(A(\theta_0)\) is the acceptance
region of a level-\(\alpha\) test of \[H_0:\theta=\theta_0.\] Then the
inverted confidence set is \[C(x)=\{\theta_0:x\in A(\theta_0)\}.\] This
set has coverage at least \(1-\alpha\) because
\[\mathbb{P}_\theta\{\theta\in C(X)\}
    =\mathbb{P}_\theta\{X\in A(\theta)\}\geq 1-\alpha.\]

\section{False coverage probability}\label{false-coverage-probability}

While ordinary coverage asks whether the confidence set contains the
true parameter, false coverage asks whether the confidence set contains
a wrong parameter value.

\begin{tcolorbox}[enhanced jigsaw, arc=.35mm, bottomrule=.15mm, bottomtitle=1mm, breakable, colback=white, colbacktitle=quarto-callout-note-color!10!white, colframe=quarto-callout-note-color-frame, coltitle=black, left=2mm, leftrule=.75mm, opacityback=0, opacitybacktitle=0.6, rightrule=.15mm, title=\textcolor{quarto-callout-note-color}{\faInfo}\hspace{0.5em}{Definition}, titlerule=0mm, toprule=.15mm, toptitle=1mm]

\textbf{Definition 9} (False coverage probability). Let \(C(X)\) be a
confidence set. For \(\theta'\neq \theta\), the \textbf{false coverage
probability} is \[\mathbb{P}_\theta\{\theta'\in C(X)\}.\] This is the
probability that the confidence set includes a false parameter value
\(\theta'\) when the true value is \(\theta\).

\end{tcolorbox}

\begin{tcolorbox}[enhanced jigsaw, arc=.35mm, bottomrule=.15mm, bottomtitle=1mm, breakable, colback=white, colbacktitle=quarto-callout-note-color!10!white, colframe=quarto-callout-note-color-frame, coltitle=black, left=2mm, leftrule=.75mm, opacityback=0, opacitybacktitle=0.6, rightrule=.15mm, title=\textcolor{quarto-callout-note-color}{\faInfo}\hspace{0.5em}{Definition}, titlerule=0mm, toprule=.15mm, toptitle=1mm]

\textbf{Definition 10} (Uniformly most accurate confidence set). A
\((1-\alpha)\) confidence set is called \textbf{uniformly most accurate}
(UMA) if it minimizes the probability of false coverage among all
confidence sets with the same coverage probability.

\end{tcolorbox}

\section{UMA intervals from UMP
tests}\label{uma-intervals-from-ump-tests}

Optimal tests lead to optimal confidence sets after inversion.

\begin{tcolorbox}[enhanced jigsaw, arc=.35mm, bottomrule=.15mm, bottomtitle=1mm, breakable, colback=white, colbacktitle=quarto-callout-important-color!10!white, colframe=quarto-callout-important-color-frame, coltitle=black, left=2mm, leftrule=.75mm, opacityback=0, opacitybacktitle=0.6, rightrule=.15mm, title=\textcolor{quarto-callout-important-color}{\faExclamation}\hspace{0.5em}{Theorem}, titlerule=0mm, toprule=.15mm, toptitle=1mm]

\textbf{Theorem 11} (UMP tests give UMA confidence sets). \emph{Let
\(X\sim f(x\mid \theta)\), where \(\theta\in\mathbb{R}\). Let
\(A^*(\theta_0)\) be the acceptance region of a UMP level-\(\alpha\)
test of \[H_0:\theta=\theta_0
    \qquad \text{versus}\qquad
    H_1:\theta>\theta_0.\] Let
\[C^*(X)=\{\theta_0:X\in A^*(\theta_0)\}\] be the confidence set
obtained by inverting these acceptance regions. Then for any other
\((1-\alpha)\) confidence set \(C\),
\[\mathbb{P}_\theta\{\theta'\in C^*(X)\}
    \leq
    \mathbb{P}_\theta\{\theta'\in C(X)\},
    \qquad \theta'<\theta.\] Thus \(C^*\) is uniformly most accurate
against false values below the true parameter.}

\end{tcolorbox}

\begin{tcolorbox}[enhanced jigsaw, arc=.35mm, bottomrule=.15mm, bottomtitle=1mm, breakable, colback=white, colbacktitle=quarto-callout-note-color!10!white, colframe=quarto-callout-note-color-frame, coltitle=black, left=2mm, leftrule=.75mm, opacityback=0, opacitybacktitle=0.6, rightrule=.15mm, title=\textcolor{quarto-callout-note-color}{\faInfo}\hspace{0.5em}{Proof}, titlerule=0mm, toprule=.15mm, toptitle=1mm]

\emph{Idea of proof.} For a fixed false value \(\theta'\), the event
\(\{\theta'\in C(X)\}\) is the same as accepting the test of
\(H_0:\theta=\theta'\) when the true parameter is \(\theta\). Minimizing
false coverage is therefore equivalent to maximizing the probability of
rejecting \(H_0:\theta=\theta'\) when the truth is \(\theta>\theta'\). A
UMP test maximizes this rejection probability, so its inverted
confidence set minimizes false coverage. ◻

\end{tcolorbox}

\section{Example: normal lower confidence
bound}\label{example-normal-lower-confidence-bound}

One-sided normal tests have UMP properties, so their inverted confidence
bounds have UMA properties.

\begin{tcolorbox}[enhanced jigsaw, arc=.35mm, bottomrule=.15mm, bottomtitle=1mm, breakable, colback=white, colbacktitle=quarto-callout-tip-color!10!white, colframe=quarto-callout-tip-color-frame, coltitle=black, left=2mm, leftrule=.75mm, opacityback=0, opacitybacktitle=0.6, rightrule=.15mm, title=\textcolor{quarto-callout-tip-color}{\faLightbulb}\hspace{0.5em}{Example}, titlerule=0mm, toprule=.15mm, toptitle=1mm]

\textbf{Example 12} (UMA lower confidence bound for a normal mean).
Suppose \[X_1,\ldots,X_n\sim \operatorname{Normal}(\mu,\sigma^2),\]
where \(\sigma\) is known. To build a lower confidence bound for
\(\mu\), invert the UMP level-\(\alpha\) test of \[H_0:\mu=\mu_0
    \qquad\text{versus}\qquad
    H_1:\mu>\mu_0.\] The UMP rejection rule is
\[\frac{\bar X-\mu_0}{\sigma/\sqrt n}>z_\alpha.\] Thus the acceptance
region is \[\frac{\bar X-\mu_0}{\sigma/\sqrt n}\leq z_\alpha.\] Solving
for \(\mu_0\) gives \[\mu_0\geq \bar X-z_\alpha\frac{\sigma}{\sqrt n}.\]
Therefore the inverted confidence set is
\[C(X)=\left[\bar X-z_\alpha\frac{\sigma}{\sqrt n},\infty\right).\]

\end{tcolorbox}

\begin{tcolorbox}[enhanced jigsaw, arc=.35mm, bottomrule=.15mm, bottomtitle=1mm, breakable, colback=white, colbacktitle=quarto-callout-note-color!10!white, colframe=quarto-callout-note-color-frame, coltitle=black, left=2mm, leftrule=.75mm, opacityback=0, opacitybacktitle=0.6, rightrule=.15mm, title=\textcolor{quarto-callout-note-color}{\faInfo}\hspace{0.5em}{Solution}, titlerule=0mm, toprule=.15mm, toptitle=1mm]

The coverage is \[\begin{aligned}
\mathbb{P}_\mu\left\{\mu\geq \bar X-z_\alpha\frac{\sigma}{\sqrt n}\right\}
&=\mathbb{P}_\mu\left\{\frac{\bar X-\mu}{\sigma/\sqrt n}\leq z_\alpha\right\}\\
&=\Phi(z_\alpha)=1-\alpha.
\end{aligned}\] Because it comes from inverting a UMP one-sided test, it
is UMA for excluding false parameter values below the true value.

\end{tcolorbox}

\begin{tcolorbox}[enhanced jigsaw, arc=.35mm, bottomrule=.15mm, bottomtitle=1mm, breakable, colback=white, colbacktitle=quarto-callout-note-color!10!white, colframe=quarto-callout-note-color-frame, coltitle=black, left=2mm, leftrule=.75mm, opacityback=0, opacitybacktitle=0.6, rightrule=.15mm, title=\textcolor{quarto-callout-note-color}{\faInfo}\hspace{0.5em}{Remark}, titlerule=0mm, toprule=.15mm, toptitle=1mm]

\emph{Remark 13} (Two-sided intervals). The usual two-sided interval
\[\left[\bar X-z_{\alpha/2}\frac{\sigma}{\sqrt n},\;
    \bar X+z_{\alpha/2}\frac{\sigma}{\sqrt n}\right]\] is not UMA in the
same one-sided sense, because for two-sided alternatives there is
generally no UMP test. This is why alternative optimality concepts, such
as unbiasedness, are useful.

\end{tcolorbox}

\section{Unbiased confidence sets}\label{unbiased-confidence-sets}

When UMA confidence sets do not exist, unbiasedness is often used as a
weaker but useful optimality condition.

\begin{tcolorbox}[enhanced jigsaw, arc=.35mm, bottomrule=.15mm, bottomtitle=1mm, breakable, colback=white, colbacktitle=quarto-callout-note-color!10!white, colframe=quarto-callout-note-color-frame, coltitle=black, left=2mm, leftrule=.75mm, opacityback=0, opacitybacktitle=0.6, rightrule=.15mm, title=\textcolor{quarto-callout-note-color}{\faInfo}\hspace{0.5em}{Definition}, titlerule=0mm, toprule=.15mm, toptitle=1mm]

\textbf{Definition 14} (Unbiased confidence set). A \((1-\alpha)\)
confidence set \(C(X)\) is called \textbf{unbiased} if
\[\mathbb{P}_\theta\{\theta'\in C(X)\}\leq 1-\alpha
    \qquad\text{for all }\theta\neq \theta'.\] Equivalently, the
interval should not cover a false parameter value more often than it
covers the true parameter value.

\end{tcolorbox}

\begin{tcolorbox}[enhanced jigsaw, arc=.35mm, bottomrule=.15mm, bottomtitle=1mm, breakable, colback=white, colbacktitle=quarto-callout-tip-color!10!white, colframe=quarto-callout-tip-color-frame, coltitle=black, left=2mm, leftrule=.75mm, opacityback=0, opacitybacktitle=0.6, rightrule=.15mm, title=\textcolor{quarto-callout-tip-color}{\faLightbulb}\hspace{0.5em}{Key idea}, titlerule=0mm, toprule=.15mm, toptitle=1mm]

Parallel with tests An unbiased test has power under the alternative at
least as large as its rejection probability under the null. An unbiased
confidence set has false coverage probability no larger than its true
coverage probability.

\end{tcolorbox}

\begin{tcolorbox}[enhanced jigsaw, arc=.35mm, bottomrule=.15mm, bottomtitle=1mm, breakable, colback=white, colbacktitle=quarto-callout-warning-color!10!white, colframe=quarto-callout-warning-color-frame, coltitle=black, left=2mm, leftrule=.75mm, opacityback=0, opacitybacktitle=0.6, rightrule=.15mm, title=\textcolor{quarto-callout-warning-color}{\faExclamationTriangle}\hspace{0.5em}{Practice Problem}, titlerule=0mm, toprule=.15mm, toptitle=1mm]

\textbf{Practice Problem 15} (Unbiasedness intuition). Explain why a
confidence set that frequently includes false parameter values more
often than the true value is undesirable.

\end{tcolorbox}

\begin{tcolorbox}[enhanced jigsaw, arc=.35mm, bottomrule=.15mm, bottomtitle=1mm, breakable, colback=white, colbacktitle=quarto-callout-note-color!10!white, colframe=quarto-callout-note-color-frame, coltitle=black, left=2mm, leftrule=.75mm, opacityback=0, opacitybacktitle=0.6, rightrule=.15mm, title=\textcolor{quarto-callout-note-color}{\faInfo}\hspace{0.5em}{Solution}, titlerule=0mm, toprule=.15mm, toptitle=1mm]

A confidence set is intended to localize the true parameter. If for some
false value \(\theta'\neq\theta\) we had
\[\mathbb{P}_\theta\{\theta'\in C(X)\}>\mathbb{P}_\theta\{\theta\in C(X)\},\]
then the procedure would include that wrong value more often than the
correct value. This means the procedure is biased toward false parameter
values. The unbiasedness condition prevents this behavior by requiring
\[\mathbb{P}_\theta\{\theta'\in C(X)\}\leq 1-\alpha.\]

\end{tcolorbox}

\bookmarksetup{startatroot}

\chapter{Bayesian Optimality}\label{bayesian-optimality}

This section evaluates Bayesian credible intervals by posterior
probability and posterior length.

\section{Shortest credible sets}\label{shortest-credible-sets}

In Bayesian inference, after observing the data, the posterior
distribution describes uncertainty about the parameter.

Let \(\pi(\theta\mid x)\) be the posterior density. A credible set
\(C(x)\) with credibility \(1-\alpha\) satisfies
\[\int_{C(x)} \pi(\theta\mid x)\,d\theta=1-\alpha.\] Among all credible
sets with posterior probability \(1-\alpha\), we often prefer the one
with the smallest length.

\begin{tcolorbox}[enhanced jigsaw, arc=.35mm, bottomrule=.15mm, bottomtitle=1mm, breakable, colback=white, colbacktitle=quarto-callout-note-color!10!white, colframe=quarto-callout-note-color-frame, coltitle=black, left=2mm, leftrule=.75mm, opacityback=0, opacitybacktitle=0.6, rightrule=.15mm, title=\textcolor{quarto-callout-note-color}{\faInfo}\hspace{0.5em}{Definition}, titlerule=0mm, toprule=.15mm, toptitle=1mm]

\textbf{Definition 16} (Highest posterior density region). A
\textbf{highest posterior density} (HPD) region has the form
\[C(x)=\{\theta:\pi(\theta\mid x)\geq k\},\] where \(k\) is chosen so
that \[\int_{C(x)}\pi(\theta\mid x)\,d\theta=1-\alpha.\]

\end{tcolorbox}

\begin{tcolorbox}[enhanced jigsaw, arc=.35mm, bottomrule=.15mm, bottomtitle=1mm, breakable, colback=white, colbacktitle=quarto-callout-important-color!10!white, colframe=quarto-callout-important-color-frame, coltitle=black, left=2mm, leftrule=.75mm, opacityback=0, opacitybacktitle=0.6, rightrule=.15mm, title=\textcolor{quarto-callout-important-color}{\faExclamation}\hspace{0.5em}{Corollary}, titlerule=0mm, toprule=.15mm, toptitle=1mm]

\textbf{Corollary 17} (HPD is shortest for unimodal posterior). \emph{If
\(\pi(\theta\mid x)\) is unimodal, then the shortest credible interval
with posterior probability \(1-\alpha\) is the HPD interval
\[C(x)=\{\theta:\pi(\theta\mid x)\geq k\}.\]}

\end{tcolorbox}

\begin{tcolorbox}[enhanced jigsaw, arc=.35mm, bottomrule=.15mm, bottomtitle=1mm, breakable, colback=white, colbacktitle=quarto-callout-note-color!10!white, colframe=quarto-callout-note-color-frame, coltitle=black, left=2mm, leftrule=.75mm, opacityback=0, opacitybacktitle=0.6, rightrule=.15mm, title=\textcolor{quarto-callout-note-color}{\faInfo}\hspace{0.5em}{Remark}, titlerule=0mm, toprule=.15mm, toptitle=1mm]

\emph{Remark 18}. For symmetric unimodal posteriors, the HPD interval
and the equal-tail credible interval are the same. For skewed
posteriors, such as many Gamma posteriors, the HPD interval is typically
shorter than the equal-tail interval.

\end{tcolorbox}

\section{Poisson HPD region}\label{poisson-hpd-region}

The Poisson-Gamma example illustrates the difference between an
equal-tail credible interval and an HPD credible interval.

\begin{tcolorbox}[enhanced jigsaw, arc=.35mm, bottomrule=.15mm, bottomtitle=1mm, breakable, colback=white, colbacktitle=quarto-callout-tip-color!10!white, colframe=quarto-callout-tip-color-frame, coltitle=black, left=2mm, leftrule=.75mm, opacityback=0, opacitybacktitle=0.6, rightrule=.15mm, title=\textcolor{quarto-callout-tip-color}{\faLightbulb}\hspace{0.5em}{Example}, titlerule=0mm, toprule=.15mm, toptitle=1mm]

\textbf{Example 19} (Poisson HPD region). Suppose
\[X_1,\ldots,X_n\sim \operatorname{Poisson}(\lambda).\] Use a conjugate
Gamma prior for \(\lambda\). With the Gamma prior parameterized by shape
\(a\) and scale \(b\), \[\lambda\sim \operatorname{Gamma}(a,b),\] the
posterior is \[\lambda\mid \sum_i x_i
    \sim
    \operatorname{Gamma}\left(a+\sum_i x_i,\frac{1}{n+1/b}\right).\] The
HPD credible region is
\[\left\{\lambda:\pi\left(\lambda\mid \sum_i x_i\right)\geq k\right\},\]
where \(k\) is chosen to satisfy
\[\int_{\{\lambda:\pi(\lambda\mid \sum_i x_i)\geq k\}}
    \pi\left(\lambda\mid \sum_i x_i\right)\,d\lambda
    =1-\alpha.\] For the specific case \[a=b=1,
    \qquad n=10,
    \qquad \sum_i x_i=6,\] the posterior is
\[\lambda\mid x\sim \operatorname{Gamma}\left(7,\frac{1}{11}\right).\] A
\(90\%\) HPD credible set is approximately \[[0.253,1.005].\]

\end{tcolorbox}

\begin{tcolorbox}[enhanced jigsaw, arc=.35mm, bottomrule=.15mm, bottomtitle=1mm, breakable, colback=white, colbacktitle=quarto-callout-note-color!10!white, colframe=quarto-callout-note-color-frame, coltitle=black, left=2mm, leftrule=.75mm, opacityback=0, opacitybacktitle=0.6, rightrule=.15mm, title=\textcolor{quarto-callout-note-color}{\faInfo}\hspace{0.5em}{Solution}, titlerule=0mm, toprule=.15mm, toptitle=1mm]

The likelihood from \(n\) independent Poisson observations is
\[L(\lambda\mid x)
    \propto
    \lambda^{\sum_i x_i}e^{-n\lambda}.\] The Gamma prior with shape
\(a\) and scale \(b\) has density proportional to
\[\lambda^{a-1}e^{-\lambda/b}.\] Multiplying prior and likelihood gives
\[\pi(\lambda\mid x)
    \propto
    \lambda^{a+
\sum_i x_i-1}e^{-(n+1/b)\lambda}.\] Therefore,
\[\lambda\mid x\sim \operatorname{Gamma}\left(a+\sum_i x_i,\frac{1}{n+1/b}\right).\]
For \(a=b=1\), \(n=10\), and \(\sum_i x_i=6\), this becomes
\[\operatorname{Gamma}\left(7,\frac{1}{11}\right).\] The HPD interval is
found by choosing a density threshold \(k\) so that the set where the
posterior density exceeds \(k\) contains \(90\%\) posterior probability.
Numerically, this gives approximately \([0.253,1.005]\).

\end{tcolorbox}

\begin{tcolorbox}[enhanced jigsaw, arc=.35mm, bottomrule=.15mm, bottomtitle=1mm, breakable, colback=white, colbacktitle=quarto-callout-tip-color!10!white, colframe=quarto-callout-tip-color-frame, coltitle=black, left=2mm, leftrule=.75mm, opacityback=0, opacitybacktitle=0.6, rightrule=.15mm, title=\textcolor{quarto-callout-tip-color}{\faLightbulb}\hspace{0.5em}{Example}, titlerule=0mm, toprule=.15mm, toptitle=1mm]

\textbf{Example 20} (Equal-tail versus HPD for the Poisson example). For
the posterior
\[\lambda\mid x\sim \operatorname{Gamma}\left(7,\frac{1}{11}\right),\]
an equal-tail \(90\%\) credible interval is approximately
\[[0.299,1.077],\] with length \[1.077-0.299=0.778.\] The HPD \(90\%\)
credible interval is approximately \[[0.247,1.000],\] with length
\[1.000-0.247=0.753.\]

\end{tcolorbox}

\begin{tcolorbox}[enhanced jigsaw, arc=.35mm, bottomrule=.15mm, bottomtitle=1mm, breakable, colback=white, colbacktitle=quarto-callout-note-color!10!white, colframe=quarto-callout-note-color-frame, coltitle=black, left=2mm, leftrule=.75mm, opacityback=0, opacitybacktitle=0.6, rightrule=.15mm, title=\textcolor{quarto-callout-note-color}{\faInfo}\hspace{0.5em}{Solution}, titlerule=0mm, toprule=.15mm, toptitle=1mm]

Both intervals contain \(90\%\) posterior probability. The equal-tail
interval puts \(5\%\) posterior probability in each tail. The HPD
interval instead includes the points with highest posterior density
until \(90\%\) probability is accumulated. Because the Gamma posterior
is skewed, the HPD interval is shorter: \[0.753<0.778.\] Thus the HPD
interval is preferable under the criterion of shortest posterior
credible set.

\end{tcolorbox}

{\def\LTcaptype{none} % do not increment counter
\begin{longtable}[]{@{}llll@{}}
\toprule\noalign{}
terval & Lower & Upper L & ength \\
\midrule\noalign{}
\endhead
\bottomrule\noalign{}
\endlastfoot
Equal-tail \(90\%\) & \(0.299\) & \(1.077\) & \(0.778\) \\
HPD (shortest) \(90\%\) & \(0.247\) & \(1.000\) & \(0.753\) \\
\end{longtable}
}

\begin{tcolorbox}[enhanced jigsaw, arc=.35mm, bottomrule=.15mm, bottomtitle=1mm, breakable, colback=white, colbacktitle=quarto-callout-warning-color!10!white, colframe=quarto-callout-warning-color-frame, coltitle=black, left=2mm, leftrule=.75mm, opacityback=0, opacitybacktitle=0.6, rightrule=.15mm, title=\textcolor{quarto-callout-warning-color}{\faExclamationTriangle}\hspace{0.5em}{Practice Problem}, titlerule=0mm, toprule=.15mm, toptitle=1mm]

\textbf{Practice Problem 21} (Posterior for a Poisson mean). Suppose
\(X_1,\ldots,X_n\sim\operatorname{Poisson}(\lambda)\) and the prior is
\(\lambda\sim\operatorname{Gamma}(a,b)\) with shape \(a\) and scale
\(b\). Derive the posterior distribution of \(\lambda\).

\end{tcolorbox}

\begin{tcolorbox}[enhanced jigsaw, arc=.35mm, bottomrule=.15mm, bottomtitle=1mm, breakable, colback=white, colbacktitle=quarto-callout-note-color!10!white, colframe=quarto-callout-note-color-frame, coltitle=black, left=2mm, leftrule=.75mm, opacityback=0, opacitybacktitle=0.6, rightrule=.15mm, title=\textcolor{quarto-callout-note-color}{\faInfo}\hspace{0.5em}{Solution}, titlerule=0mm, toprule=.15mm, toptitle=1mm]

The likelihood is \[L(\lambda\mid x)
    =\prod_{i=1}^n e^{-\lambda}\frac{\lambda^{x_i}}{x_i!}
    \propto
    e^{-n\lambda}\lambda^{\sum_i x_i}.\] The prior density is
\[\pi(\lambda)\propto \lambda^{a-1}e^{-\lambda/b}.\] Thus
\[\begin{aligned}
    \pi(\lambda\mid x)
    &\propto L(\lambda\mid x)\pi(\lambda)\\
    &\propto \lambda^{\sum_i x_i}e^{-n\lambda}\lambda^{a-1}e^{-\lambda/b}\\
    &=\lambda^{a+
\sum_i x_i-1}e^{-(n+1/b)\lambda}.
\end{aligned}\] Therefore
\[\lambda\mid x\sim\operatorname{Gamma}\left(a+\sum_i x_i,\frac{1}{n+1/b}\right).\]

\end{tcolorbox}

\bookmarksetup{startatroot}

\chapter{Loss-Function Optimality}\label{loss-function-optimality}

This section evaluates intervals using a decision-theoretic risk that
combines interval length and coverage.

\section{A loss function for confidence
sets}\label{a-loss-function-for-confidence-sets}

Loss-function optimality turns interval estimation into a decision
problem.

The action is to choose a confidence set \(C\). A simple loss function
is
\[L(\theta,C)=b\cdot \operatorname{Length}(C)-\mathbbm{1}\{\theta\in C\},\]
where \(b>0\) is a tuning constant.

\begin{itemize}
\item
  The term \(b\cdot \operatorname{Length}(C)\) penalizes long intervals.
\item
  The term \(-\mathbbm{1}\{\theta\in C\}\) rewards intervals that cover
  the true parameter.
\item
  Large \(b\) prioritizes shorter intervals.
\item
  Small \(b\) prioritizes coverage.
\end{itemize}

\begin{tcolorbox}[enhanced jigsaw, arc=.35mm, bottomrule=.15mm, bottomtitle=1mm, breakable, colback=white, colbacktitle=quarto-callout-note-color!10!white, colframe=quarto-callout-note-color-frame, coltitle=black, left=2mm, leftrule=.75mm, opacityback=0, opacitybacktitle=0.6, rightrule=.15mm, title=\textcolor{quarto-callout-note-color}{\faInfo}\hspace{0.5em}{Definition}, titlerule=0mm, toprule=.15mm, toptitle=1mm]

\textbf{Definition 22} (Risk of a confidence set). The risk of an
interval procedure \(C(X)\) is \[R(\theta,C)
    =\mathbb{E}_\theta[L(\theta,C(X))]
    =b\mathbb{E}_\theta[\operatorname{Length}(C(X))]-\mathbb{P}_\theta\{\theta\in C(X)\}.\]

\end{tcolorbox}

\begin{tcolorbox}[enhanced jigsaw, arc=.35mm, bottomrule=.15mm, bottomtitle=1mm, breakable, colback=white, colbacktitle=quarto-callout-tip-color!10!white, colframe=quarto-callout-tip-color-frame, coltitle=black, left=2mm, leftrule=.75mm, opacityback=0, opacitybacktitle=0.6, rightrule=.15mm, title=\textcolor{quarto-callout-tip-color}{\faLightbulb}\hspace{0.5em}{Key idea}, titlerule=0mm, toprule=.15mm, toptitle=1mm]

Interpretation The risk combines two competing goals:
\[\text{short expected length}
    \qquad\text{and}\qquad
    \text{large coverage probability}.\] A low-risk interval is short
but still covers the true parameter with high probability.

\end{tcolorbox}

\section{Normal example: optimizing interval
half-width}\label{normal-example-optimizing-interval-half-width}

The normal example shows how a loss function can determine an optimal
confidence level.

\begin{tcolorbox}[enhanced jigsaw, arc=.35mm, bottomrule=.15mm, bottomtitle=1mm, breakable, colback=white, colbacktitle=quarto-callout-tip-color!10!white, colframe=quarto-callout-tip-color-frame, coltitle=black, left=2mm, leftrule=.75mm, opacityback=0, opacitybacktitle=0.6, rightrule=.15mm, title=\textcolor{quarto-callout-tip-color}{\faLightbulb}\hspace{0.5em}{Example}, titlerule=0mm, toprule=.15mm, toptitle=1mm]

\textbf{Example 23} (Risk for symmetric normal intervals). Suppose
\[X\sim \operatorname{Normal}(\mu,\sigma^2),\] where \(\sigma^2\) is
known. Consider the class of symmetric intervals
\[C(X)=[X-c\sigma,X+c\sigma],\qquad c\geq 0.\] The length is
\[\operatorname{Length}(C)=2c\sigma.\] The coverage probability is
\[\begin{aligned}
    \mathbb{P}_\mu\{\mu\in C(X)\}
    &=\mathbb{P}_\mu\{X-c\sigma\leq \mu\leq X+c\sigma\}\\
    &=\mathbb{P}\left(-c\leq \frac{X-\mu}{\sigma}\leq c\right)\\
    &=2\Phi(c)-1.
\end{aligned}\] Therefore the risk is
\[R(\mu,C)=b(2c\sigma)-\{2\Phi(c)-1\}.\] This risk does not depend on
\(\mu\).

\end{tcolorbox}

\begin{tcolorbox}[enhanced jigsaw, arc=.35mm, bottomrule=.15mm, bottomtitle=1mm, breakable, colback=white, colbacktitle=quarto-callout-note-color!10!white, colframe=quarto-callout-note-color-frame, coltitle=black, left=2mm, leftrule=.75mm, opacityback=0, opacitybacktitle=0.6, rightrule=.15mm, title=\textcolor{quarto-callout-note-color}{\faInfo}\hspace{0.5em}{Solution}, titlerule=0mm, toprule=.15mm, toptitle=1mm]

Since \(Z=(X-\mu)/\sigma\sim\operatorname{Normal}(0,1)\),
\[\mathbb{P}_\mu\{\mu\in C(X)\}=\mathbb{P}(-c\leq Z\leq c)=2\Phi(c)-1.\]
The expected length is simply \(2c\sigma\) because the interval length
is nonrandom. Substituting into
\[R(\mu,C)=b\mathbb{E}_\mu[\operatorname{Length}(C)]-\mathbb{P}_\mu\{\mu\in C(X)\}\]
gives \[R(c)=2b\sigma c-2\Phi(c)+1.\]

\end{tcolorbox}

\begin{tcolorbox}[enhanced jigsaw, arc=.35mm, bottomrule=.15mm, bottomtitle=1mm, breakable, colback=white, colbacktitle=quarto-callout-important-color!10!white, colframe=quarto-callout-important-color-frame, coltitle=black, left=2mm, leftrule=.75mm, opacityback=0, opacitybacktitle=0.6, rightrule=.15mm, title=\textcolor{quarto-callout-important-color}{\faExclamation}\hspace{0.5em}{Proposition}, titlerule=0mm, toprule=.15mm, toptitle=1mm]

\textbf{Proposition 24} (Optimal half-width). \emph{For the normal
interval risk \[R(c)=2b\sigma c-2\Phi(c)+1,
    \qquad c\geq 0,\] we have:}

\begin{enumerate}
\def\labelenumi{\arabic{enumi}.}
\item
  \emph{If \(b\sigma>1/\sqrt{2\pi}\), the minimizing value is
  \(c^*=0\).}
\item
  \emph{If \(b\sigma\leq 1/\sqrt{2\pi}\), the interior minimizing value
  satisfies \[\phi(c^*)=b\sigma,\] so
  \[c^*=\sqrt{-2\log(b\sigma\sqrt{2\pi})}.\]}
\end{enumerate}

\end{tcolorbox}

\begin{tcolorbox}[enhanced jigsaw, arc=.35mm, bottomrule=.15mm, bottomtitle=1mm, breakable, colback=white, colbacktitle=quarto-callout-note-color!10!white, colframe=quarto-callout-note-color-frame, coltitle=black, left=2mm, leftrule=.75mm, opacityback=0, opacitybacktitle=0.6, rightrule=.15mm, title=\textcolor{quarto-callout-note-color}{\faInfo}\hspace{0.5em}{Proof}, titlerule=0mm, toprule=.15mm, toptitle=1mm]

\emph{Proof.} Differentiate \(R(c)\): \[R'(c)=2b\sigma-2\phi(c).\] At
\(c=0\), \(\phi(0)=1/\sqrt{2\pi}\). If
\[b\sigma>\frac{1}{\sqrt{2\pi}},\] then \(R'(0)>0\) and \(R'(c)>0\) for
all \(c\geq 0\), so the minimum occurs at \(c=0\).

If \[b\sigma\leq \frac{1}{\sqrt{2\pi}},\] then an interior solution
satisfies \[\phi(c)=b\sigma.\] Because
\[\phi(c)=\frac{1}{\sqrt{2\pi}}e^{-c^2/2},\] we solve
\[\frac{1}{\sqrt{2\pi}}e^{-c^2/2}=b\sigma.\] Thus
\[e^{-c^2/2}=b\sigma\sqrt{2\pi},\] and hence
\[c^*=\sqrt{-2\log(b\sigma\sqrt{2\pi})}.\] ◻

\end{tcolorbox}

\begin{tcolorbox}[enhanced jigsaw, arc=.35mm, bottomrule=.15mm, bottomtitle=1mm, breakable, colback=white, colbacktitle=quarto-callout-note-color!10!white, colframe=quarto-callout-note-color-frame, coltitle=black, left=2mm, leftrule=.75mm, opacityback=0, opacitybacktitle=0.6, rightrule=.15mm, title=\textcolor{quarto-callout-note-color}{\faInfo}\hspace{0.5em}{Remark}, titlerule=0mm, toprule=.15mm, toptitle=1mm]

\emph{Remark 25}. If we write \(c=z_{\alpha/2}\), then the optimal risk
corresponds to a standard two-sided confidence interval with confidence
level \[1-\alpha=2\Phi(c)-1.\] Thus the loss-function approach can be
interpreted as choosing the confidence level by balancing coverage
against length.

\end{tcolorbox}

\begin{tcolorbox}[enhanced jigsaw, arc=.35mm, bottomrule=.15mm, bottomtitle=1mm, breakable, colback=white, colbacktitle=quarto-callout-warning-color!10!white, colframe=quarto-callout-warning-color-frame, coltitle=black, left=2mm, leftrule=.75mm, opacityback=0, opacitybacktitle=0.6, rightrule=.15mm, title=\textcolor{quarto-callout-warning-color}{\faExclamationTriangle}\hspace{0.5em}{Practice Problem}, titlerule=0mm, toprule=.15mm, toptitle=1mm]

\textbf{Practice Problem 26} (Optimal half-width). Let
\(X\sim\operatorname{Normal}(\mu,1)\) and consider intervals
\(C(X)=[X-c,X+c]\). With loss
\[L(\mu,C)=b\operatorname{Length}(C)-\mathbbm{1}\{\mu\in C\},\] find the
optimal \(c\) when \(b=0.2\).

\end{tcolorbox}

\begin{tcolorbox}[enhanced jigsaw, arc=.35mm, bottomrule=.15mm, bottomtitle=1mm, breakable, colback=white, colbacktitle=quarto-callout-note-color!10!white, colframe=quarto-callout-note-color-frame, coltitle=black, left=2mm, leftrule=.75mm, opacityback=0, opacitybacktitle=0.6, rightrule=.15mm, title=\textcolor{quarto-callout-note-color}{\faInfo}\hspace{0.5em}{Solution}, titlerule=0mm, toprule=.15mm, toptitle=1mm]

Here \(\sigma=1\), so the interior solution satisfies \[\phi(c)=b=0.2.\]
Since \[\phi(c)=\frac{1}{\sqrt{2\pi}}e^{-c^2/2},\] we solve
\[c=\sqrt{-2\log(0.2\sqrt{2\pi})}.\] Numerically,
\[0.2\sqrt{2\pi}\approx 0.5013,\] so
\[c\approx \sqrt{-2\log(0.5013)}\approx 1.18.\] The optimal interval is
therefore approximately \[[X-1.18,X+1.18].\] The corresponding coverage
is \[2\Phi(1.18)-1\approx 0.762.\]

\end{tcolorbox}

\bookmarksetup{startatroot}

\chapter{Summary and Comparison of
Criteria}\label{summary-and-comparison-of-criteria}

This final section summarizes the main ways to evaluate interval
estimators.

{\def\LTcaptype{none} % do not increment counter
\begin{longtable}[]{@{}lll@{}}
\toprule\noalign{}
Criterion** ** & Main idea** ** & Typical goal** \\
\midrule\noalign{}
\endhead
\bottomrule\noalign{}
\endlastfoot
Coverage probability & Probability the interval contains the true
parameter & Achieve at least \(1-\alpha\) coverage \\
Length / expected length & Width of the interval & Prefer shorter
intervals among those with same coverage \\
Shortest unimodal interval & Choose interval with equal boundary density
& Minimize length for fixed probability \\
Test-related optimality & Invert optimal hypothesis tests & Obtain UMA
confidence sets when UMP tests exist \\
Unbiased confidence sets & Avoid covering false values too often &
Ensure false coverage is not larger than true coverage \\
Bayesian HPD & Use highest posterior density region & Shortest credible
set for fixed posterior probability \\
Loss-function optimality & Combine length and coverage in one risk &
Choose interval minimizing expected loss \\
\end{longtable}
}

\begin{tcolorbox}[enhanced jigsaw, arc=.35mm, bottomrule=.15mm, bottomtitle=1mm, breakable, colback=white, colbacktitle=quarto-callout-tip-color!10!white, colframe=quarto-callout-tip-color-frame, coltitle=black, left=2mm, leftrule=.75mm, opacityback=0, opacitybacktitle=0.6, rightrule=.15mm, title=\textcolor{quarto-callout-tip-color}{\faLightbulb}\hspace{0.5em}{Key idea}, titlerule=0mm, toprule=.15mm, toptitle=1mm]

\textbf{Key takeaways.}

\begin{enumerate}
\def\labelenumi{\arabic{enumi}.}
\item
  A confidence interval should have high coverage probability and small
  length.
\item
  For unimodal densities, the shortest probability interval has equal
  density at its endpoints.
\item
  Inverting UMP tests can produce uniformly most accurate confidence
  sets.
\item
  For Bayesian inference, HPD intervals are shortest credible sets for
  unimodal posterior distributions.
\item
  Loss functions provide a decision-theoretic way to balance coverage
  and interval length.
\end{enumerate}

\end{tcolorbox}

\bookmarksetup{startatroot}

\chapter{Chapter 18: Analysis of Variance I --- One-Way
ANOVA}\label{chapter-18-analysis-of-variance-i-one-way-anova}

This chapter introduces \textbf{one-way analysis of variance (ANOVA)}.
The goal is to compare more than two normal population means while
controlling the overall Type I error rate. We start from the pooled
two-sample \(t\)-test, introduce the cell-means ANOVA model, study
contrasts, and derive the classical ANOVA \(F\)-test and ANOVA table.

\begin{tcolorbox}[enhanced jigsaw, arc=.35mm, bottomrule=.15mm, bottomtitle=1mm, breakable, colback=white, colbacktitle=quarto-callout-note-color!10!white, colframe=quarto-callout-note-color-frame, coltitle=black, left=2mm, leftrule=.75mm, opacityback=0, opacitybacktitle=0.6, rightrule=.15mm, title=\textcolor{quarto-callout-note-color}{\faInfo}\hspace{0.5em}{Topics}, titlerule=0mm, toprule=.15mm, toptitle=1mm]

Pooled two-sample \(t\)-test; one-way ANOVA model assumptions; global
ANOVA hypothesis; treatment effects; identifiability; linear
combinations and contrasts; inference for contrasts; sums of squares;
ANOVA decomposition; chi-square theory; ANOVA \(F\)-test; one-way ANOVA
table.

\end{tcolorbox}

\bookmarksetup{startatroot}

\chapter{Overview}\label{overview-11}

This section explains how ANOVA extends the two-sample pooled \(t\)-test
to compare several treatment means at once.

The main topics are:

\begin{enumerate}
\def\labelenumi{\arabic{enumi}.}
\item
  review of the pooled two-sample \(t\)-test;
\item
  one-way ANOVA model assumptions;
\item
  the classical ANOVA hypothesis;
\item
  linear combinations and contrasts;
\item
  inference for contrasts;
\item
  the ANOVA \(F\)-test;
\item
  the one-way ANOVA table.
\end{enumerate}

\begin{tcolorbox}[enhanced jigsaw, arc=.35mm, bottomrule=.15mm, bottomtitle=1mm, breakable, colback=white, colbacktitle=quarto-callout-tip-color!10!white, colframe=quarto-callout-tip-color-frame, coltitle=black, left=2mm, leftrule=.75mm, opacityback=0, opacitybacktitle=0.6, rightrule=.15mm, title=\textcolor{quarto-callout-tip-color}{\faLightbulb}\hspace{0.5em}{Key idea}, titlerule=0mm, toprule=.15mm, toptitle=1mm]

\textbf{Main idea.} ANOVA compares variation \emph{between} group means
with variation \emph{within} groups. If the between-group variation is
large relative to the within-group variation, then the data provide
evidence that not all group means are equal.

\end{tcolorbox}

\bookmarksetup{startatroot}

\chapter{\texorpdfstring{Review: Two Independent Samples and the Pooled
\(t\)-Test}{Review: Two Independent Samples and the Pooled t-Test}}\label{review-two-independent-samples-and-the-pooled-t-test}

This section reviews the two-sample method that ANOVA generalizes.

\section{Two independent normal
samples}\label{two-independent-normal-samples}

Suppose we observe two independent samples
\[X_1,\ldots,X_n \sim \operatorname{Normal}(\mu_X,\sigma^2),
    \qquad
    Y_1,\ldots,Y_m \sim \operatorname{Normal}(\mu_Y,\sigma^2),\] where
the two population variances are assumed equal but unknown. The
classical pooled \(t\)-test tests \[H_0: \mu_X=\mu_Y
    \qquad \text{versus} \qquad
    H_1: \mu_X \neq \mu_Y,\] or one-sided alternatives such as
\(H_1: \mu_X>\mu_Y\).

\begin{tcolorbox}[enhanced jigsaw, arc=.35mm, bottomrule=.15mm, bottomtitle=1mm, breakable, colback=white, colbacktitle=quarto-callout-note-color!10!white, colframe=quarto-callout-note-color-frame, coltitle=black, left=2mm, leftrule=.75mm, opacityback=0, opacitybacktitle=0.6, rightrule=.15mm, title=\textcolor{quarto-callout-note-color}{\faInfo}\hspace{0.5em}{Definition}, titlerule=0mm, toprule=.15mm, toptitle=1mm]

\textbf{Definition 1} (Pooled variance estimator). The pooled variance
estimator is \[S_p^2
    =\frac{(n-1)S_X^2+(m-1)S_Y^2}{n+m-2},\] where
\[S_X^2=\frac{1}{n-1}\sum_{i=1}^n (X_i-\bar X)^2,
    \qquad
    S_Y^2=\frac{1}{m-1}\sum_{j=1}^m (Y_j-\bar Y)^2.\]

\end{tcolorbox}

\begin{tcolorbox}[enhanced jigsaw, arc=.35mm, bottomrule=.15mm, bottomtitle=1mm, breakable, colback=white, colbacktitle=quarto-callout-important-color!10!white, colframe=quarto-callout-important-color-frame, coltitle=black, left=2mm, leftrule=.75mm, opacityback=0, opacitybacktitle=0.6, rightrule=.15mm, title=\textcolor{quarto-callout-important-color}{\faExclamation}\hspace{0.5em}{Theorem}, titlerule=0mm, toprule=.15mm, toptitle=1mm]

\textbf{Theorem 2} (Pooled \(t\)-statistic). \emph{Under
\(H_0:\mu_X=\mu_Y\), the statistic \[T
    =\frac{\bar X-\bar Y}{S_p\sqrt{\frac{1}{n}+\frac{1}{m}}}\] has a
Student \(t\) distribution with \(n+m-2\) degrees of freedom:
\[T\sim t_{n+m-2}.\]}

\end{tcolorbox}

\section{Rejection rules}\label{rejection-rules}

The rejection rule depends on the alternative hypothesis.

{\def\LTcaptype{none} % do not increment counter
\begin{longtable}[]{@{}lll@{}}
\toprule\noalign{}
Alternative & Rejection region & Interpretation \\
\midrule\noalign{}
\endhead
\bottomrule\noalign{}
\endlastfoot
\(H_1:\mu_X>\mu_Y\) & \(T\geq t_{n+m-2,\alpha}\) & right-tailed test \\
\(H_1:\mu_X<\mu_Y\) & \(T\leq -t_{n+m-2,\alpha}\) & left-tailed test \\
\(H_1:\mu_X\neq \mu_Y\) & \(|T|\geq t_{n+m-2,\alpha/2}\) & two-tailed
test \\
\end{longtable}
}

\begin{tcolorbox}[enhanced jigsaw, arc=.35mm, bottomrule=.15mm, bottomtitle=1mm, breakable, colback=white, colbacktitle=quarto-callout-note-color!10!white, colframe=quarto-callout-note-color-frame, coltitle=black, left=2mm, leftrule=.75mm, opacityback=0, opacitybacktitle=0.6, rightrule=.15mm, title=\textcolor{quarto-callout-note-color}{\faInfo}\hspace{0.5em}{Remark}, titlerule=0mm, toprule=.15mm, toptitle=1mm]

\emph{Remark 3}. The pooled \(t\)-test is a test for the difference
between two means. One-way ANOVA asks the corresponding question for
\(k\) independent groups: are all \(k\) population means equal?

\end{tcolorbox}

\begin{tcolorbox}[enhanced jigsaw, arc=.35mm, bottomrule=.15mm, bottomtitle=1mm, breakable, colback=white, colbacktitle=quarto-callout-tip-color!10!white, colframe=quarto-callout-tip-color-frame, coltitle=black, left=2mm, leftrule=.75mm, opacityback=0, opacitybacktitle=0.6, rightrule=.15mm, title=\textcolor{quarto-callout-tip-color}{\faLightbulb}\hspace{0.5em}{Example}, titlerule=0mm, toprule=.15mm, toptitle=1mm]

\textbf{Example 4} (Pooled \(t\)-test as a special case of ANOVA).
Suppose there are only two treatment groups. Then one-way ANOVA compares
\(\theta_1\) and \(\theta_2\). The ANOVA null hypothesis
\[H_0:\theta_1=\theta_2\] is exactly the same null hypothesis as the
two-sample pooled \(t\)-test.

\end{tcolorbox}

\begin{tcolorbox}[enhanced jigsaw, arc=.35mm, bottomrule=.15mm, bottomtitle=1mm, breakable, colback=white, colbacktitle=quarto-callout-caution-color!10!white, colframe=quarto-callout-caution-color-frame, coltitle=black, left=2mm, leftrule=.75mm, opacityback=0, opacitybacktitle=0.6, rightrule=.15mm, title=\textcolor{quarto-callout-caution-color}{\faFire}\hspace{0.5em}{Solution}, titlerule=0mm, toprule=.15mm, toptitle=1mm]

When \(k=2\), the between-group variation has only one degree of
freedom. The one-way ANOVA \(F\) statistic satisfies \[F=T^2,\] where
\(T\) is the pooled two-sample \(t\) statistic. Therefore, for two
groups, the ANOVA \(F\)-test and the two-sided pooled \(t\)-test lead to
the same conclusion.

\end{tcolorbox}

\bookmarksetup{startatroot}

\chapter{The One-Way ANOVA Cell-Means
Model}\label{the-one-way-anova-cell-means-model}

This section introduces the statistical model used for one-way ANOVA.

\section{Model notation}\label{model-notation}

Suppose there are \(k\) groups or treatments. In group \(i\), we observe
\[Y_{i1},Y_{i2},\ldots,Y_{i n_i},
    \qquad i=1,\ldots,k.\] The one-way ANOVA cell-means model is
\[Y_{ij}=\theta_i+\varepsilon_{ij},
    \qquad i=1,\ldots,k,
    \qquad j=1,\ldots,n_i.\] Here:

\begin{itemize}
\item
  \(i\) is the group or treatment index;
\item
  \(j\) is the observation index within group \(i\);
\item
  \(\theta_i\) is the mean of group \(i\);
\item
  \(\varepsilon_{ij}\) is the random error or noise term.
\end{itemize}

\begin{tcolorbox}[enhanced jigsaw, arc=.35mm, bottomrule=.15mm, bottomtitle=1mm, breakable, colback=white, colbacktitle=quarto-callout-note-color!10!white, colframe=quarto-callout-note-color-frame, coltitle=black, left=2mm, leftrule=.75mm, opacityback=0, opacitybacktitle=0.6, rightrule=.15mm, title=\textcolor{quarto-callout-note-color}{\faInfo}\hspace{0.5em}{Definition}, titlerule=0mm, toprule=.15mm, toptitle=1mm]

\textbf{Definition 5} (One-way ANOVA model). The one-way ANOVA model
assumes that \[Y_{ij}\sim \operatorname{Normal}(\theta_i,\sigma^2),\]
independently across all \(i\) and \(j\). Equivalently,
\[\varepsilon_{ij}\sim \operatorname{Normal}(0,\sigma^2)\]
independently.

\end{tcolorbox}

\section{Model assumptions}\label{model-assumptions}

One-way ANOVA depends on three core assumptions.

\begin{tcolorbox}[enhanced jigsaw, arc=.35mm, bottomrule=.15mm, bottomtitle=1mm, breakable, colback=white, colbacktitle=quarto-callout-tip-color!10!white, colframe=quarto-callout-tip-color-frame, coltitle=black, left=2mm, leftrule=.75mm, opacityback=0, opacitybacktitle=0.6, rightrule=.15mm, title=\textcolor{quarto-callout-tip-color}{\faLightbulb}\hspace{0.5em}{Key idea}, titlerule=0mm, toprule=.15mm, toptitle=1mm]

\textbf{One-way ANOVA assumptions.} For the model
\(Y_{ij}=\theta_i+\varepsilon_{ij}\), assume:

\begin{enumerate}
\def\labelenumi{\arabic{enumi}.}
\item
  \(\mathbb{E}(\varepsilon_{ij})=0\) and
  \(\operatorname{Var}(\varepsilon_{ij})=\sigma^2<\infty\);
\item
  the errors are independent and normally distributed;
\item
  the groups have equal variance: \(\sigma_i^2=\sigma^2\) for all
  \(i=1,\ldots,k\).
\end{enumerate}

\end{tcolorbox}

\begin{tcolorbox}[enhanced jigsaw, arc=.35mm, bottomrule=.15mm, bottomtitle=1mm, breakable, colback=white, colbacktitle=quarto-callout-note-color!10!white, colframe=quarto-callout-note-color-frame, coltitle=black, left=2mm, leftrule=.75mm, opacityback=0, opacitybacktitle=0.6, rightrule=.15mm, title=\textcolor{quarto-callout-note-color}{\faInfo}\hspace{0.5em}{Remark}, titlerule=0mm, toprule=.15mm, toptitle=1mm]

\emph{Remark 6}. The equal-variance assumption is the same assumption
used by the pooled two-sample \(t\)-test. ANOVA pools the within-group
variation across all \(k\) groups to estimate the common error variance
\(\sigma^2\).

\end{tcolorbox}

\section{Dot notation}\label{dot-notation}

This subsection introduces the compact dot notation used throughout
ANOVA.

Let \[N=\sum_{i=1}^k n_i\] be the total sample size.

\begin{tcolorbox}[enhanced jigsaw, arc=.35mm, bottomrule=.15mm, bottomtitle=1mm, breakable, colback=white, colbacktitle=quarto-callout-note-color!10!white, colframe=quarto-callout-note-color-frame, coltitle=black, left=2mm, leftrule=.75mm, opacityback=0, opacitybacktitle=0.6, rightrule=.15mm, title=\textcolor{quarto-callout-note-color}{\faInfo}\hspace{0.5em}{Definition}, titlerule=0mm, toprule=.15mm, toptitle=1mm]

\textbf{Definition 7} (Group totals, group means, and grand mean). The
group total for group \(i\) is \[T_{i\cdot}=\sum_{j=1}^{n_i} Y_{ij}.\]
The group mean is \[\bar Y_{i\cdot}=\frac{T_{i\cdot}}{n_i}
    =\frac{1}{n_i}\sum_{j=1}^{n_i}Y_{ij}.\] The overall total is
\[T_{\cdot\cdot}=\sum_{i=1}^k\sum_{j=1}^{n_i}Y_{ij}
    =\sum_{i=1}^k T_{i\cdot}.\] The grand mean is \[\bar Y_{\cdot\cdot}
    =\frac{1}{N}\sum_{i=1}^k\sum_{j=1}^{n_i}Y_{ij}
    =\frac{1}{N}\sum_{i=1}^k T_{i\cdot}
    =\frac{1}{N}\sum_{i=1}^k n_i\bar Y_{i\cdot}.\]

\end{tcolorbox}

\begin{tcolorbox}[enhanced jigsaw, arc=.35mm, bottomrule=.15mm, bottomtitle=1mm, breakable, colback=white, colbacktitle=quarto-callout-note-color!10!white, colframe=quarto-callout-note-color-frame, coltitle=black, left=2mm, leftrule=.75mm, opacityback=0, opacitybacktitle=0.6, rightrule=.15mm, title=\textcolor{quarto-callout-note-color}{\faInfo}\hspace{0.5em}{Remark}, titlerule=0mm, toprule=.15mm, toptitle=1mm]

\emph{Remark 8}. A dot in a subscript means that we sum or average over
that index. Thus \(T_{i\cdot}\) sums over \(j\), while
\(T_{\cdot\cdot}\) sums over both \(i\) and \(j\).

\end{tcolorbox}

\bookmarksetup{startatroot}

\chapter{Example: Liver Deterioration
Data}\label{example-liver-deterioration-data}

This section works through the toxin example from the lecture notes and
computes the descriptive summaries and ANOVA table.

\section{Data table}\label{data-table}

The data record liver deterioration measurements for four groups:
Control, Toxin 1, Toxin 2, and Toxin 3.

{\def\LTcaptype{none} % do not increment counter
\begin{longtable}[]{@{}rrrr@{}}
\toprule\noalign{}
Toxin 1 & Toxin 2 & Toxin 3 & Control \\
\midrule\noalign{}
\endhead
\bottomrule\noalign{}
\endlastfoot
28 & 33 & 18 & 11 \\
23 & 36 & 21 & 14 \\
14 & 34 & 20 & 11 \\
27 & 29 & 22 & 16 \\
& 31 & 24 & \\
& 34 & & \\
\end{longtable}
}

The group descriptive statistics are:

{\def\LTcaptype{none} % do not increment counter
\begin{longtable}[]{@{}lrrrr@{}}
\toprule\noalign{}
Group & \(n_i\) & \(\bar y_{i\cdot}\) & \(s_i^2\) & \(s_i\) \\
\midrule\noalign{}
\endhead
\bottomrule\noalign{}
\endlastfoot
Control & 4 & 13.000 & 6.000 & 2.449 \\
Toxin 1 & 4 & 23.000 & 40.667 & 6.377 \\
Toxin 2 & 6 & 32.833 & 6.167 & 2.483 \\
Toxin 3 & 5 & 21.000 & 5.000 & 2.236 \\
\end{longtable}
}

\begin{tcolorbox}[enhanced jigsaw, arc=.35mm, bottomrule=.15mm, bottomtitle=1mm, breakable, colback=white, colbacktitle=quarto-callout-tip-color!10!white, colframe=quarto-callout-tip-color-frame, coltitle=black, left=2mm, leftrule=.75mm, opacityback=0, opacitybacktitle=0.6, rightrule=.15mm, title=\textcolor{quarto-callout-tip-color}{\faLightbulb}\hspace{0.5em}{Example}, titlerule=0mm, toprule=.15mm, toptitle=1mm]

\textbf{Example 9} (Interpreting the group means). Based on the group
means, which group appears to have the largest liver deterioration?

\end{tcolorbox}

\begin{tcolorbox}[enhanced jigsaw, arc=.35mm, bottomrule=.15mm, bottomtitle=1mm, breakable, colback=white, colbacktitle=quarto-callout-caution-color!10!white, colframe=quarto-callout-caution-color-frame, coltitle=black, left=2mm, leftrule=.75mm, opacityback=0, opacitybacktitle=0.6, rightrule=.15mm, title=\textcolor{quarto-callout-caution-color}{\faFire}\hspace{0.5em}{Solution}, titlerule=0mm, toprule=.15mm, toptitle=1mm]

The group means are \[\bar y_{\text{Control}}=13.000,
    \quad
    \bar y_{\text{Toxin 1}}=23.000,
    \quad
    \bar y_{\text{Toxin 2}}=32.833,
    \quad
    \bar y_{\text{Toxin 3}}=21.000.\] The largest sample mean is for
Toxin 2. This suggests that Toxin 2 may produce the highest average
deterioration, but ANOVA is needed to decide whether the differences are
statistically significant relative to within-group variability.

\end{tcolorbox}

\section{Boxplot interpretation}\label{boxplot-interpretation}

The lecture notes include a boxplot showing that the Control group is
clearly lower than the toxin groups and that Toxin 2 is the highest
group. The boxplot is useful because it displays both location and
within-group spread.

\bookmarksetup{startatroot}

\chapter{Overparameterized Treatment-Effect
Model}\label{overparameterized-treatment-effect-model}

This section explains an alternative but non-identifiable way to write
the same one-way ANOVA model.

\section{Grand mean plus treatment
effects}\label{grand-mean-plus-treatment-effects}

Instead of writing \[Y_{ij}=\theta_i+\varepsilon_{ij},\] one can write
\[Y_{ij}=\mu+\tau_i+\varepsilon_{ij},\] where:

\begin{itemize}
\item
  \(\mu\) is a grand mean or common baseline;
\item
  \(\tau_i\) is the treatment effect for group \(i\);
\item
  \(\varepsilon_{ij}\) is the random error.
\end{itemize}

\section{Non-identifiability}\label{non-identifiability}

The model \[Y_{ij}=\mu+\tau_i+\varepsilon_{ij}\] is overparameterized
unless we impose a constraint. Indeed, for any constant \(c\),
\[\mu'=\mu+c,
    \qquad
    \tau_i'=\tau_i-c\] gives the same fitted mean because
\[\mu'+\tau_i'=(\mu+c)+(\tau_i-c)=\mu+\tau_i.\] Thus, \(\mu\) and the
\(\tau_i\)'s are not uniquely determined.

\begin{tcolorbox}[enhanced jigsaw, arc=.35mm, bottomrule=.15mm, bottomtitle=1mm, breakable, colback=white, colbacktitle=quarto-callout-note-color!10!white, colframe=quarto-callout-note-color-frame, coltitle=black, left=2mm, leftrule=.75mm, opacityback=0, opacitybacktitle=0.6, rightrule=.15mm, title=\textcolor{quarto-callout-note-color}{\faInfo}\hspace{0.5em}{Definition}, titlerule=0mm, toprule=.15mm, toptitle=1mm]

\textbf{Definition 10} (Identifiability). A parameter is identifiable if
different parameter values lead to different probability distributions
for the observed data.

\end{tcolorbox}

\begin{tcolorbox}[enhanced jigsaw, arc=.35mm, bottomrule=.15mm, bottomtitle=1mm, breakable, colback=white, colbacktitle=quarto-callout-tip-color!10!white, colframe=quarto-callout-tip-color-frame, coltitle=black, left=2mm, leftrule=.75mm, opacityback=0, opacitybacktitle=0.6, rightrule=.15mm, title=\textcolor{quarto-callout-tip-color}{\faLightbulb}\hspace{0.5em}{Key idea}, titlerule=0mm, toprule=.15mm, toptitle=1mm]

\textbf{Common identifiability constraint.} A standard way to identify
the treatment-effect model is to impose \[\sum_{i=1}^k \tau_i=0.\] Under
this constraint, \(\mu\) represents the average treatment mean and
\(\tau_i\) represents the deviation of group \(i\) from that average.

\end{tcolorbox}

\bookmarksetup{startatroot}

\chapter{The Classical ANOVA
Hypothesis}\label{the-classical-anova-hypothesis}

This section states the global null hypothesis tested by one-way ANOVA.

\begin{tcolorbox}[enhanced jigsaw, arc=.35mm, bottomrule=.15mm, bottomtitle=1mm, breakable, colback=white, colbacktitle=quarto-callout-note-color!10!white, colframe=quarto-callout-note-color-frame, coltitle=black, left=2mm, leftrule=.75mm, opacityback=0, opacitybacktitle=0.6, rightrule=.15mm, title=\textcolor{quarto-callout-note-color}{\faInfo}\hspace{0.5em}{Definition}, titlerule=0mm, toprule=.15mm, toptitle=1mm]

\textbf{Definition 11} (Classical one-way ANOVA hypothesis). The
classical ANOVA null hypothesis is
\[H_0:\theta_1=\theta_2=\cdots=\theta_k.\] The alternative hypothesis is
\[H_1:\theta_i\neq \theta_j
    \quad \text{for some } i,j.\]

\end{tcolorbox}

\begin{tcolorbox}[enhanced jigsaw, arc=.35mm, bottomrule=.15mm, bottomtitle=1mm, breakable, colback=white, colbacktitle=quarto-callout-note-color!10!white, colframe=quarto-callout-note-color-frame, coltitle=black, left=2mm, leftrule=.75mm, opacityback=0, opacitybacktitle=0.6, rightrule=.15mm, title=\textcolor{quarto-callout-note-color}{\faInfo}\hspace{0.5em}{Remark}, titlerule=0mm, toprule=.15mm, toptitle=1mm]

\emph{Remark 12}. Rejecting \(H_0\) does not mean that all means are
different. It only means that at least two treatment means differ. After
rejecting the global ANOVA null, follow-up comparisons or contrasts are
needed to determine which means differ and how they differ.

\end{tcolorbox}

\begin{tcolorbox}[enhanced jigsaw, arc=.35mm, bottomrule=.15mm, bottomtitle=1mm, breakable, colback=white, colbacktitle=quarto-callout-tip-color!10!white, colframe=quarto-callout-tip-color-frame, coltitle=black, left=2mm, leftrule=.75mm, opacityback=0, opacitybacktitle=0.6, rightrule=.15mm, title=\textcolor{quarto-callout-tip-color}{\faLightbulb}\hspace{0.5em}{Example}, titlerule=0mm, toprule=.15mm, toptitle=1mm]

\textbf{Example 13} (Agricultural ANOVA hypothesis). A study
investigates the effect of fertilizer mixtures on the zinc content of
spinach plants. The treatments are:

{\def\LTcaptype{none} % do not increment counter
\begin{longtable}[]{@{}crrr@{}}
\toprule\noalign{}
Treatment & Magnesium & Potassium & Zinc \\
\midrule\noalign{}
\endhead
\bottomrule\noalign{}
\endlastfoot
1 & 0 & 0 & 0 \\
2 & 0 & 200 & 0 \\
3 & 50 & 200 & 0 \\
4 & 200 & 200 & 0 \\
5 & 0 & 200 & 15 \\
\end{longtable}
}

The global ANOVA hypothesis is
\[H_0:\theta_1=\theta_2=\theta_3=\theta_4=\theta_5.\]

\end{tcolorbox}

\begin{tcolorbox}[enhanced jigsaw, arc=.35mm, bottomrule=.15mm, bottomtitle=1mm, breakable, colback=white, colbacktitle=quarto-callout-caution-color!10!white, colframe=quarto-callout-caution-color-frame, coltitle=black, left=2mm, leftrule=.75mm, opacityback=0, opacitybacktitle=0.6, rightrule=.15mm, title=\textcolor{quarto-callout-caution-color}{\faFire}\hspace{0.5em}{Solution}, titlerule=0mm, toprule=.15mm, toptitle=1mm]

Here \(\theta_i\) is the mean zinc content under fertilizer treatment
\(i\). The null hypothesis says that all five fertilizer mixtures have
the same mean effect. The scientific question is usually not merely
whether the means are equal, but how the treatments compare. For
example, one may compare treatment 1 with treatment 2, or compare
zinc-added treatment 5 with potassium-only treatment 2.

\end{tcolorbox}

\bookmarksetup{startatroot}

\chapter{Linear Combinations and
Contrasts}\label{linear-combinations-and-contrasts}

This section introduces contrasts, which are the most useful way to ask
meaningful comparison questions in ANOVA.

\section{Linear combinations}\label{linear-combinations}

Given numbers or parameters \[t=(t_1,\ldots,t_k)\] and constants
\[a=(a_1,\ldots,a_k),\] a linear combination is
\[\sum_{i=1}^k a_i t_i.\]

\begin{tcolorbox}[enhanced jigsaw, arc=.35mm, bottomrule=.15mm, bottomtitle=1mm, breakable, colback=white, colbacktitle=quarto-callout-note-color!10!white, colframe=quarto-callout-note-color-frame, coltitle=black, left=2mm, leftrule=.75mm, opacityback=0, opacitybacktitle=0.6, rightrule=.15mm, title=\textcolor{quarto-callout-note-color}{\faInfo}\hspace{0.5em}{Definition}, titlerule=0mm, toprule=.15mm, toptitle=1mm]

\textbf{Definition 14} (Contrast). A linear combination
\[\sum_{i=1}^k a_i t_i\] is called a \textbf{contrast} if
\[\sum_{i=1}^k a_i=0.\]

\end{tcolorbox}

\begin{tcolorbox}[enhanced jigsaw, arc=.35mm, bottomrule=.15mm, bottomtitle=1mm, breakable, colback=white, colbacktitle=quarto-callout-tip-color!10!white, colframe=quarto-callout-tip-color-frame, coltitle=black, left=2mm, leftrule=.75mm, opacityback=0, opacitybacktitle=0.6, rightrule=.15mm, title=\textcolor{quarto-callout-tip-color}{\faLightbulb}\hspace{0.5em}{Example}, titlerule=0mm, toprule=.15mm, toptitle=1mm]

\textbf{Example 15} (Basic contrasts). Common contrasts include:

\begin{enumerate}
\def\labelenumi{\arabic{enumi}.}
\item
  Comparing one treatment with a control:
  \[\theta_{\text{Toxin}}-\theta_{\text{Control}}.\] This corresponds to
  coefficients \((1,-1)\) for the treatment and control.
\item
  Comparing the average of two treatments with a control:
  \[\frac{1}{2}(\theta_{T1}+\theta_{T2})-\theta_C.\] The coefficients
  are \(1/2,1/2,-1\), whose sum is zero.
\end{enumerate}

\end{tcolorbox}

\begin{tcolorbox}[enhanced jigsaw, arc=.35mm, bottomrule=.15mm, bottomtitle=1mm, breakable, colback=white, colbacktitle=quarto-callout-caution-color!10!white, colframe=quarto-callout-caution-color-frame, coltitle=black, left=2mm, leftrule=.75mm, opacityback=0, opacitybacktitle=0.6, rightrule=.15mm, title=\textcolor{quarto-callout-caution-color}{\faFire}\hspace{0.5em}{Solution}, titlerule=0mm, toprule=.15mm, toptitle=1mm]

Both examples are contrasts because the coefficients add to zero.
Contrasts are meaningful because they compare means rather than
depending on the arbitrary zero point of the measurement scale.

\end{tcolorbox}

\section{ANOVA null as all contrasts equal
zero}\label{anova-null-as-all-contrasts-equal-zero}

The ANOVA null hypothesis can be expressed using contrasts.

\begin{tcolorbox}[enhanced jigsaw, arc=.35mm, bottomrule=.15mm, bottomtitle=1mm, breakable, colback=white, colbacktitle=quarto-callout-important-color!10!white, colframe=quarto-callout-important-color-frame, coltitle=black, left=2mm, leftrule=.75mm, opacityback=0, opacitybacktitle=0.6, rightrule=.15mm, title=\textcolor{quarto-callout-important-color}{\faExclamation}\hspace{0.5em}{Theorem}, titlerule=0mm, toprule=.15mm, toptitle=1mm]

\textbf{Theorem 16} (ANOVA null and contrasts). \emph{The null
hypothesis \[H_0:\theta_1=\theta_2=\cdots=\theta_k\] is equivalent to
the statement that every contrast is zero: \[\sum_{i=1}^k a_i\theta_i=0
    \qquad \text{for every } a=(a_1,\ldots,a_k) \text{ satisfying } \sum_{i=1}^k a_i=0.\]}

\end{tcolorbox}

\begin{tcolorbox}[enhanced jigsaw, arc=.35mm, bottomrule=.15mm, bottomtitle=1mm, breakable, colback=white, colbacktitle=quarto-callout-note-color!10!white, colframe=quarto-callout-note-color-frame, coltitle=black, left=2mm, leftrule=.75mm, opacityback=0, opacitybacktitle=0.6, rightrule=.15mm, title=\textcolor{quarto-callout-note-color}{\faInfo}\hspace{0.5em}{Proof}, titlerule=0mm, toprule=.15mm, toptitle=1mm]

\emph{Proof.} If \(\theta_1=\cdots=\theta_k=\theta\), then for any
contrast vector \(a\), \[\sum_{i=1}^k a_i\theta_i
    =\theta\sum_{i=1}^k a_i
    =0.\] Conversely, suppose every contrast is zero. Choose the
contrast vector with \(a_i=1\), \(a_j=-1\), and all other entries zero.
Then \[\theta_i-\theta_j=0.\] Since this holds for every pair \(i,j\),
all \(\theta_i\) are equal.~◻

\end{tcolorbox}

\begin{tcolorbox}[enhanced jigsaw, arc=.35mm, bottomrule=.15mm, bottomtitle=1mm, breakable, colback=white, colbacktitle=quarto-callout-tip-color!10!white, colframe=quarto-callout-tip-color-frame, coltitle=black, left=2mm, leftrule=.75mm, opacityback=0, opacitybacktitle=0.6, rightrule=.15mm, title=\textcolor{quarto-callout-tip-color}{\faLightbulb}\hspace{0.5em}{Key idea}, titlerule=0mm, toprule=.15mm, toptitle=1mm]

\textbf{Why contrasts matter.} The global ANOVA test tells us whether
some difference exists. Contrasts tell us which scientific comparison is
being made: treatment versus control, one treatment versus another, or
one group of treatments versus another group.

\end{tcolorbox}

\bookmarksetup{startatroot}

\chapter{Inference for Linear Combinations of
Means}\label{inference-for-linear-combinations-of-means}

This section derives tests and confidence intervals for general linear
combinations and contrasts.

\section{Distribution of group means}\label{distribution-of-group-means}

Under the one-way ANOVA model,
\[Y_{ij}\sim \operatorname{Normal}(\theta_i,\sigma^2),\] so each group
sample mean satisfies \[\bar Y_{i\cdot}
    =\frac{1}{n_i}\sum_{j=1}^{n_i}Y_{ij}
    \sim \operatorname{Normal}\left(\theta_i,\frac{\sigma^2}{n_i}\right).\]
For a linear combination \[\widehat L=\sum_{i=1}^k a_i\bar Y_{i\cdot},\]
we have \[\mathbb{E}(\widehat L)=\sum_{i=1}^k a_i\theta_i\] and, by
independence across groups,
\[\operatorname{Var}(\widehat L)=\sigma^2\sum_{i=1}^k \frac{a_i^2}{n_i}.\]

\begin{tcolorbox}[enhanced jigsaw, arc=.35mm, bottomrule=.15mm, bottomtitle=1mm, breakable, colback=white, colbacktitle=quarto-callout-important-color!10!white, colframe=quarto-callout-important-color-frame, coltitle=black, left=2mm, leftrule=.75mm, opacityback=0, opacitybacktitle=0.6, rightrule=.15mm, title=\textcolor{quarto-callout-important-color}{\faExclamation}\hspace{0.5em}{Theorem}, titlerule=0mm, toprule=.15mm, toptitle=1mm]

\textbf{Theorem 17} (Known-variance standardized statistic). \emph{If
\(\sigma^2\) is known, then \[Z
    =\frac{\sum_{i=1}^k a_i\bar Y_{i\cdot}-\sum_{i=1}^k a_i\theta_i}
    {\sigma\sqrt{\sum_{i=1}^k a_i^2/n_i}}
    \sim \operatorname{Normal}(0,1).\]}

\end{tcolorbox}

\section{\texorpdfstring{Pooled ANOVA estimate of
\(\sigma^2\)}{Pooled ANOVA estimate of \textbackslash sigma\^{}2}}\label{pooled-anova-estimate-of-sigma2}

In practice, \(\sigma^2\) is unknown. ANOVA estimates it using pooled
within-group variation: \[S_p^2
    =\frac{1}{N-k}\sum_{i=1}^k\sum_{j=1}^{n_i}(Y_{ij}-\bar Y_{i\cdot})^2.\]
This is also called the mean squared error: \[\mathrm{MSE}=S_p^2.\]
Under the normal ANOVA model,
\[\frac{(N-k)S_p^2}{\sigma^2}\sim \chi^2_{N-k}.\]

\begin{tcolorbox}[enhanced jigsaw, arc=.35mm, bottomrule=.15mm, bottomtitle=1mm, breakable, colback=white, colbacktitle=quarto-callout-important-color!10!white, colframe=quarto-callout-important-color-frame, coltitle=black, left=2mm, leftrule=.75mm, opacityback=0, opacitybacktitle=0.6, rightrule=.15mm, title=\textcolor{quarto-callout-important-color}{\faExclamation}\hspace{0.5em}{Theorem}, titlerule=0mm, toprule=.15mm, toptitle=1mm]

\textbf{Theorem 18} (Contrast \(t\)-statistic). \emph{For a contrast or
linear combination \[L=\sum_{i=1}^k a_i\theta_i,
    \qquad
    \widehat L=\sum_{i=1}^k a_i\bar Y_{i\cdot},\] the statistic \[T
    =\frac{\widehat L-L}{S_p\sqrt{\sum_{i=1}^k a_i^2/n_i}}\] has a
Student \(t\) distribution with \(N-k\) degrees of freedom:
\[T\sim t_{N-k}.\]}

\end{tcolorbox}

\section{Testing a contrast}\label{testing-a-contrast}

To test \[H_0:\sum_{i=1}^k a_i\theta_i=0
    \qquad \text{versus} \qquad
    H_1:\sum_{i=1}^k a_i\theta_i\neq 0,\] use the statistic \[T
    =\frac{\sum_{i=1}^k a_i\bar Y_{i\cdot}}
    {S_p\sqrt{\sum_{i=1}^k a_i^2/n_i}}.\] At level \(\alpha\), reject
\(H_0\) if \[|T|>t_{N-k,\alpha/2}.\]

\section{Confidence interval for a
contrast}\label{confidence-interval-for-a-contrast}

A \(100(1-\alpha)\%\) confidence interval for the contrast
\[L=\sum_{i=1}^k a_i\theta_i\] is \[\sum_{i=1}^k a_i\bar Y_{i\cdot}
    \pm
    t_{N-k,\alpha/2}
    S_p\sqrt{\sum_{i=1}^k\frac{a_i^2}{n_i}}.\]

\begin{tcolorbox}[enhanced jigsaw, arc=.35mm, bottomrule=.15mm, bottomtitle=1mm, breakable, colback=white, colbacktitle=quarto-callout-tip-color!10!white, colframe=quarto-callout-tip-color-frame, coltitle=black, left=2mm, leftrule=.75mm, opacityback=0, opacitybacktitle=0.6, rightrule=.15mm, title=\textcolor{quarto-callout-tip-color}{\faLightbulb}\hspace{0.5em}{Example}, titlerule=0mm, toprule=.15mm, toptitle=1mm]

\textbf{Example 19} (Comparing two treatments directly). To compare
treatment 1 with treatment 2, use the contrast vector
\[a=(1,-1,0,\ldots,0).\] Then \[L=\theta_1-\theta_2,
    \qquad
    \widehat L=\bar Y_{1\cdot}-\bar Y_{2\cdot}.\] The test statistic is
\[T=\frac{\bar Y_{1\cdot}-\bar Y_{2\cdot}}
    {S_p\sqrt{\frac{1}{n_1}+\frac{1}{n_2}}}.\]

\end{tcolorbox}

\begin{tcolorbox}[enhanced jigsaw, arc=.35mm, bottomrule=.15mm, bottomtitle=1mm, breakable, colback=white, colbacktitle=quarto-callout-caution-color!10!white, colframe=quarto-callout-caution-color-frame, coltitle=black, left=2mm, leftrule=.75mm, opacityback=0, opacitybacktitle=0.6, rightrule=.15mm, title=\textcolor{quarto-callout-caution-color}{\faFire}\hspace{0.5em}{Solution}, titlerule=0mm, toprule=.15mm, toptitle=1mm]

The coefficient vector sums to \(1-1+0+\cdots+0=0\), so it is a
contrast. The statistic has a \(t_{N-k}\) distribution under
\(H_0:\theta_1=\theta_2\). This resembles the two-sample pooled
\(t\)-test, but \(S_p^2\) is estimated using all groups, not only groups
1 and 2.

\end{tcolorbox}

\begin{tcolorbox}[enhanced jigsaw, arc=.35mm, bottomrule=.15mm, bottomtitle=1mm, breakable, colback=white, colbacktitle=quarto-callout-tip-color!10!white, colframe=quarto-callout-tip-color-frame, coltitle=black, left=2mm, leftrule=.75mm, opacityback=0, opacitybacktitle=0.6, rightrule=.15mm, title=\textcolor{quarto-callout-tip-color}{\faLightbulb}\hspace{0.5em}{Example}, titlerule=0mm, toprule=.15mm, toptitle=1mm]

\textbf{Example 20} (Control versus average of two treatments). Suppose
treatment 1 is a control and treatments 2 and 3 are experimental. To
compare the control with the average of treatments 2 and 3, use
\[a=\left(1,-\frac12,-\frac12,0,\ldots,0\right).\] Then
\[L=\theta_1-\frac12(\theta_2+\theta_3),\] and
\[\widehat L=\bar Y_{1\cdot}-\frac12(\bar Y_{2\cdot}+\bar Y_{3\cdot}).\]

\end{tcolorbox}

\begin{tcolorbox}[enhanced jigsaw, arc=.35mm, bottomrule=.15mm, bottomtitle=1mm, breakable, colback=white, colbacktitle=quarto-callout-caution-color!10!white, colframe=quarto-callout-caution-color-frame, coltitle=black, left=2mm, leftrule=.75mm, opacityback=0, opacitybacktitle=0.6, rightrule=.15mm, title=\textcolor{quarto-callout-caution-color}{\faFire}\hspace{0.5em}{Solution}, titlerule=0mm, toprule=.15mm, toptitle=1mm]

The standard error is
\[S_p\sqrt{\frac{1}{n_1}+\frac{1}{4n_2}+\frac{1}{4n_3}}.\] Thus the
statistic for testing \[H_0:\theta_1=\frac12(\theta_2+\theta_3)\] is
\[T=
    \frac{\bar Y_{1\cdot}-\frac12(\bar Y_{2\cdot}+\bar Y_{3\cdot})}
    {S_p\sqrt{\frac{1}{n_1}+\frac{1}{4n_2}+\frac{1}{4n_3}}}.\] Reject
\(H_0\) for large \(|T|\).

\end{tcolorbox}

\begin{tcolorbox}[enhanced jigsaw, arc=.35mm, bottomrule=.15mm, bottomtitle=1mm, breakable, colback=white, colbacktitle=quarto-callout-warning-color!10!white, colframe=quarto-callout-warning-color-frame, coltitle=black, left=2mm, leftrule=.75mm, opacityback=0, opacitybacktitle=0.6, rightrule=.15mm, title=\textcolor{quarto-callout-warning-color}{\faExclamationTriangle}\hspace{0.5em}{Important correction}, titlerule=0mm, toprule=.15mm, toptitle=1mm]

Sign correction For the contrast ``control versus the average of
treatments 2 and 3,'' the natural contrast is
\[\bar Y_{1\cdot}-\frac12(\bar Y_{2\cdot}+\bar Y_{3\cdot}),\] not
\(\bar Y_{1\cdot}-\frac12(\bar Y_{2\cdot}-\bar Y_{3\cdot})\).

\end{tcolorbox}

\bookmarksetup{startatroot}

\chapter{\texorpdfstring{Why Not Use Many Pairwise
\(t\)-Tests?}{Why Not Use Many Pairwise t-Tests?}}\label{why-not-use-many-pairwise-t-tests}

This section explains why ANOVA uses one global \(F\)-test instead of
many separate \(t\)-tests.

The classical ANOVA null hypothesis is
\[H_0:\theta_1=\theta_2=\cdots=\theta_k.\] One could try testing many
pairwise hypotheses: \[H_0:\theta_i=\theta_j
    \qquad \text{versus} \qquad
    H_1:\theta_i\neq \theta_j.\] However, running many separate tests
inflates the probability of at least one false positive.

\begin{tcolorbox}[enhanced jigsaw, arc=.35mm, bottomrule=.15mm, bottomtitle=1mm, breakable, colback=white, colbacktitle=quarto-callout-tip-color!10!white, colframe=quarto-callout-tip-color-frame, coltitle=black, left=2mm, leftrule=.75mm, opacityback=0, opacitybacktitle=0.6, rightrule=.15mm, title=\textcolor{quarto-callout-tip-color}{\faLightbulb}\hspace{0.5em}{Example}, titlerule=0mm, toprule=.15mm, toptitle=1mm]

\textbf{Example 21} (Inflation of Type I error). Suppose each individual
test has Type I error probability \(0.05\) and the tests are
independent. With two tests, the probability of at least one Type I
error is \[1-(1-0.05)^2=1-0.95^2\approx 0.0975.\] With three tests, the
probability is \[1-0.95^3\approx 0.143.\]

\end{tcolorbox}

\begin{tcolorbox}[enhanced jigsaw, arc=.35mm, bottomrule=.15mm, bottomtitle=1mm, breakable, colback=white, colbacktitle=quarto-callout-caution-color!10!white, colframe=quarto-callout-caution-color-frame, coltitle=black, left=2mm, leftrule=.75mm, opacityback=0, opacitybacktitle=0.6, rightrule=.15mm, title=\textcolor{quarto-callout-caution-color}{\faFire}\hspace{0.5em}{Solution}, titlerule=0mm, toprule=.15mm, toptitle=1mm]

Each test has a \(95\%\) chance of making no Type I error. If tests are
independent, the chance that all tests avoid Type I error is \(0.95^r\)
for \(r\) tests. Therefore, the chance of at least one Type I error is
\(1-0.95^r\). This increases as \(r\) increases.

\end{tcolorbox}

\begin{tcolorbox}[enhanced jigsaw, arc=.35mm, bottomrule=.15mm, bottomtitle=1mm, breakable, colback=white, colbacktitle=quarto-callout-tip-color!10!white, colframe=quarto-callout-tip-color-frame, coltitle=black, left=2mm, leftrule=.75mm, opacityback=0, opacitybacktitle=0.6, rightrule=.15mm, title=\textcolor{quarto-callout-tip-color}{\faLightbulb}\hspace{0.5em}{Key idea}, titlerule=0mm, toprule=.15mm, toptitle=1mm]

\textbf{ANOVA advantage.} The ANOVA \(F\)-test provides a single global
test of equal means while controlling the Type I error rate at the
desired level \(\alpha\).

\end{tcolorbox}

\bookmarksetup{startatroot}

\chapter{\texorpdfstring{Union--Intersection View of the ANOVA
\(F\)-Test}{Union--Intersection View of the ANOVA F-Test}}\label{unionintersection-view-of-the-anova-f-test}

This section connects the ANOVA \(F\)-test with the contrast-based view
of the null hypothesis.

Let \[\mathcal A=\left\{a=(a_1,\ldots,a_k):\sum_{i=1}^k a_i=0\right\}\]
be the set of all contrast vectors. The ANOVA null can be written as
\[H_0:\sum_{i=1}^k a_i\theta_i=0
    \quad \text{for all } a\in\mathcal A.\] The alternative is
\[H_1:\sum_{i=1}^k a_i\theta_i\neq 0
    \quad \text{for some } a\in\mathcal A.\] For each contrast \(a\),
define \[\Theta_a=\left\{\theta:\sum_{i=1}^k a_i\theta_i=0\right\}.\]
Then the ANOVA null is the intersection
\[H_0:\theta\in\bigcap_{a\in\mathcal A}\Theta_a.\]

For each contrast, the natural test statistic is \[T_a
    =\frac{\left|\sum_{i=1}^k a_i\bar Y_{i\cdot}\right|}
    {S_p\sqrt{\sum_{i=1}^k a_i^2/n_i}}.\] The union--intersection
principle rejects the global null if some contrast gives strong evidence
against zero: \[\sup_{a\in\mathcal A}T_a>c.\] This maximization leads to
the classical ANOVA \(F\) statistic.

\bookmarksetup{startatroot}

\chapter{Sums of Squares in One-Way
ANOVA}\label{sums-of-squares-in-one-way-anova}

This section defines the three key sums of squares: between-group,
within-group, and total.

\section{Definitions}\label{definitions}

The total variability around the grand mean is decomposed into
between-group and within-group parts.

\begin{tcolorbox}[enhanced jigsaw, arc=.35mm, bottomrule=.15mm, bottomtitle=1mm, breakable, colback=white, colbacktitle=quarto-callout-note-color!10!white, colframe=quarto-callout-note-color-frame, coltitle=black, left=2mm, leftrule=.75mm, opacityback=0, opacitybacktitle=0.6, rightrule=.15mm, title=\textcolor{quarto-callout-note-color}{\faInfo}\hspace{0.5em}{Definition}, titlerule=0mm, toprule=.15mm, toptitle=1mm]

\textbf{Definition 22} (ANOVA sums of squares). The
\textbf{between-treatment sum of squares} is \[\mathrm{SSB}
    =\sum_{i=1}^k n_i(\bar Y_{i\cdot}-\bar Y_{\cdot\cdot})^2.\] The
\textbf{error} or \textbf{within-treatment sum of squares} is
\[\mathrm{SSE}
    =\sum_{i=1}^k\sum_{j=1}^{n_i}(Y_{ij}-\bar Y_{i\cdot})^2.\] The
\textbf{total sum of squares} is \[\mathrm{SST}
    =\sum_{i=1}^k\sum_{j=1}^{n_i}(Y_{ij}-\bar Y_{\cdot\cdot})^2.\]

\end{tcolorbox}

\begin{tcolorbox}[enhanced jigsaw, arc=.35mm, bottomrule=.15mm, bottomtitle=1mm, breakable, colback=white, colbacktitle=quarto-callout-important-color!10!white, colframe=quarto-callout-important-color-frame, coltitle=black, left=2mm, leftrule=.75mm, opacityback=0, opacitybacktitle=0.6, rightrule=.15mm, title=\textcolor{quarto-callout-important-color}{\faExclamation}\hspace{0.5em}{Theorem}, titlerule=0mm, toprule=.15mm, toptitle=1mm]

\textbf{Theorem 23} (Partitioning sums of squares). \emph{The total sum
of squares decomposes as \[\mathrm{SST}=\mathrm{SSB}+\mathrm{SSE}.\]}

\end{tcolorbox}

\begin{tcolorbox}[enhanced jigsaw, arc=.35mm, bottomrule=.15mm, bottomtitle=1mm, breakable, colback=white, colbacktitle=quarto-callout-note-color!10!white, colframe=quarto-callout-note-color-frame, coltitle=black, left=2mm, leftrule=.75mm, opacityback=0, opacitybacktitle=0.6, rightrule=.15mm, title=\textcolor{quarto-callout-note-color}{\faInfo}\hspace{0.5em}{Proof}, titlerule=0mm, toprule=.15mm, toptitle=1mm]

\emph{Proof.} Write \[Y_{ij}-\bar Y_{\cdot\cdot}
    =(Y_{ij}-\bar Y_{i\cdot})+(\bar Y_{i\cdot}-\bar Y_{\cdot\cdot}).\]
Squaring and summing gives \[\begin{aligned}
\mathrm{SST}
&=\sum_{i=1}^k\sum_{j=1}^{n_i}(Y_{ij}-\bar Y_{i\cdot})^2
+\sum_{i=1}^k\sum_{j=1}^{n_i}(\bar Y_{i\cdot}-\bar Y_{\cdot\cdot})^2 \\
&\qquad
+2\sum_{i=1}^k(\bar Y_{i\cdot}-\bar Y_{\cdot\cdot})\sum_{j=1}^{n_i}(Y_{ij}-\bar Y_{i\cdot}).
\end{aligned}\] The cross term is zero because
\[\sum_{j=1}^{n_i}(Y_{ij}-\bar Y_{i\cdot})=0\] for each group.
Therefore, \[\mathrm{SST}=\mathrm{SSE}+\mathrm{SSB}.\]~◻

\end{tcolorbox}

\section{Mean squares}\label{mean-squares}

Define the mean squares \[\mathrm{MSB}=\frac{\mathrm{SSB}}{k-1},
    \qquad
    \mathrm{MSE}=\frac{\mathrm{SSE}}{N-k}.\] The within-group mean
square is also the pooled ANOVA variance estimate:
\[\mathrm{MSE}=S_p^2.\]

\bookmarksetup{startatroot}

\chapter{\texorpdfstring{Chi-Square Distributions and the ANOVA
\(F\)-Statistic}{Chi-Square Distributions and the ANOVA F-Statistic}}\label{chi-square-distributions-and-the-anova-f-statistic}

This section gives the distributional theory behind the ANOVA
\(F\)-test.

\begin{tcolorbox}[enhanced jigsaw, arc=.35mm, bottomrule=.15mm, bottomtitle=1mm, breakable, colback=white, colbacktitle=quarto-callout-important-color!10!white, colframe=quarto-callout-important-color-frame, coltitle=black, left=2mm, leftrule=.75mm, opacityback=0, opacitybacktitle=0.6, rightrule=.15mm, title=\textcolor{quarto-callout-important-color}{\faExclamation}\hspace{0.5em}{Theorem}, titlerule=0mm, toprule=.15mm, toptitle=1mm]

\textbf{Theorem 24} (Chi-square distributions of sums of squares).
\emph{Under the ANOVA assumptions,
\[\frac{\mathrm{SSE}}{\sigma^2}\sim \chi^2_{N-k}.\] If the null
hypothesis \[H_0:\theta_1=\theta_2=\cdots=\theta_k\] is true, then
\[\frac{\mathrm{SSB}}{\sigma^2}\sim \chi^2_{k-1}.\] Also,
\[\frac{\mathrm{SST}}{\sigma^2}\sim \chi^2_{N-1}\] under the null.
Moreover, under the normal ANOVA model, \(\mathrm{SSB}\) and
\(\mathrm{SSE}\) are independent.}

\end{tcolorbox}

\begin{tcolorbox}[enhanced jigsaw, arc=.35mm, bottomrule=.15mm, bottomtitle=1mm, breakable, colback=white, colbacktitle=quarto-callout-important-color!10!white, colframe=quarto-callout-important-color-frame, coltitle=black, left=2mm, leftrule=.75mm, opacityback=0, opacitybacktitle=0.6, rightrule=.15mm, title=\textcolor{quarto-callout-important-color}{\faExclamation}\hspace{0.5em}{Theorem}, titlerule=0mm, toprule=.15mm, toptitle=1mm]

\textbf{Theorem 25} (ANOVA \(F\) statistic). \emph{If
\[H_0:\theta_1=\theta_2=\cdots=\theta_k\] is true, then \[F
    =\frac{\mathrm{SSB}/(k-1)}{\mathrm{SSE}/(N-k)}
    =\frac{\mathrm{MSB}}{\mathrm{MSE}}\] has an \(F\) distribution with
\(k-1\) and \(N-k\) degrees of freedom: \[F\sim F_{k-1,N-k}.\]}

\end{tcolorbox}

\begin{tcolorbox}[enhanced jigsaw, arc=.35mm, bottomrule=.15mm, bottomtitle=1mm, breakable, colback=white, colbacktitle=quarto-callout-note-color!10!white, colframe=quarto-callout-note-color-frame, coltitle=black, left=2mm, leftrule=.75mm, opacityback=0, opacitybacktitle=0.6, rightrule=.15mm, title=\textcolor{quarto-callout-note-color}{\faInfo}\hspace{0.5em}{Remark}, titlerule=0mm, toprule=.15mm, toptitle=1mm]

\emph{Remark 26}. The ANOVA \(F\) statistic compares between-group
variation with within-group variation. Large values of \(F\) indicate
that group means are spread out more than expected from random
within-group error alone.

\end{tcolorbox}

\section{Rejection rule and p-value}\label{rejection-rule-and-p-value}

For an \(\alpha\) level one-way ANOVA test, reject
\[H_0:\theta_1=\cdots=\theta_k\] if
\[F_{\mathrm{obs}}>F_{k-1,N-k,\alpha},\] where \(F_{k-1,N-k,\alpha}\) is
the upper \(\alpha\) critical value of the \(F_{k-1,N-k}\) distribution.

The p-value is
\[p\text{-value}=\mathbb{P}\left(F_{k-1,N-k}\geq F_{\mathrm{obs}}\right).\]

\bookmarksetup{startatroot}

\chapter{One-Way ANOVA Table}\label{one-way-anova-table}

This section summarizes the ANOVA calculations in the standard table
format.

{\def\LTcaptype{none} % do not increment counter
\begin{longtable}[]{@{}lcccc@{}}
\toprule\noalign{}
Source of Variation & Sum of Squares & Degrees of Freedom & Mean Square
& \(F\) Statistic \\
\midrule\noalign{}
\endhead
\bottomrule\noalign{}
\endlastfoot
Between Groups & \(\mathrm{SSB}\) & \(k-1\) &
\(\mathrm{MSB}=\dfrac{\mathrm{SSB}}{k-1}\) &
\(\dfrac{\mathrm{MSB}}{\mathrm{MSE}}\) \\
Within Groups / Error & \(\mathrm{SSE}\) & \(N-k\) &
\(\mathrm{MSE}=\dfrac{\mathrm{SSE}}{N-k}\) & \\
Total & \(\mathrm{SST}\) & \(N-1\) & & \\
\end{longtable}
}

\begin{tcolorbox}[enhanced jigsaw, arc=.35mm, bottomrule=.15mm, bottomtitle=1mm, breakable, colback=white, colbacktitle=quarto-callout-tip-color!10!white, colframe=quarto-callout-tip-color-frame, coltitle=black, left=2mm, leftrule=.75mm, opacityback=0, opacitybacktitle=0.6, rightrule=.15mm, title=\textcolor{quarto-callout-tip-color}{\faLightbulb}\hspace{0.5em}{Key idea}, titlerule=0mm, toprule=.15mm, toptitle=1mm]

\textbf{How to read an ANOVA table.} The ANOVA table separates total
variation into two parts:

\begin{itemize}
\item
  variation explained by group membership, measured by \(\mathrm{SSB}\);
\item
  unexplained within-group variation, measured by \(\mathrm{SSE}\).
\end{itemize}

The ratio \(F=\mathrm{MSB}/\mathrm{MSE}\) is large when group means are
far apart relative to within-group noise.

\end{tcolorbox}

\bookmarksetup{startatroot}

\chapter{Worked ANOVA Example}\label{worked-anova-example}

This section completes the liver deterioration example using the ANOVA
sums of squares and \(F\)-test.

\begin{tcolorbox}[enhanced jigsaw, arc=.35mm, bottomrule=.15mm, bottomtitle=1mm, breakable, colback=white, colbacktitle=quarto-callout-tip-color!10!white, colframe=quarto-callout-tip-color-frame, coltitle=black, left=2mm, leftrule=.75mm, opacityback=0, opacitybacktitle=0.6, rightrule=.15mm, title=\textcolor{quarto-callout-tip-color}{\faLightbulb}\hspace{0.5em}{Example}, titlerule=0mm, toprule=.15mm, toptitle=1mm]

\textbf{Example 27} (ANOVA table for the toxin data). Using the liver
deterioration data, test whether the four group means are equal.

\end{tcolorbox}

\begin{tcolorbox}[enhanced jigsaw, arc=.35mm, bottomrule=.15mm, bottomtitle=1mm, breakable, colback=white, colbacktitle=quarto-callout-caution-color!10!white, colframe=quarto-callout-caution-color-frame, coltitle=black, left=2mm, leftrule=.75mm, opacityback=0, opacitybacktitle=0.6, rightrule=.15mm, title=\textcolor{quarto-callout-caution-color}{\faFire}\hspace{0.5em}{Solution}, titlerule=0mm, toprule=.15mm, toptitle=1mm]

There are \(k=4\) groups and total sample size \[N=4+4+6+5=19.\] The
group means are \[13.000,
    \quad 23.000,
    \quad 32.833,
    \quad 21.000.\] The grand mean is \[\bar y_{\cdot\cdot}=23.474.\]
The sums of squares are \[\mathrm{SSB}=995.904,
    \qquad
    \mathrm{SSE}=190.833,
    \qquad
    \mathrm{SST}=1186.737.\] The degrees of freedom are
\[\mathrm{df}_{\text{between}}=k-1=3,
    \qquad
    \mathrm{df}_{\text{error}}=N-k=15,
    \qquad
    \mathrm{df}_{\text{total}}=N-1=18.\] Thus
\[\mathrm{MSB}=\frac{995.904}{3}=331.968,
    \qquad
    \mathrm{MSE}=\frac{190.833}{15}=12.722.\] The observed \(F\)
statistic is \[F_{\mathrm{obs}}
    =\frac{331.968}{12.722}
    =26.094.\] The ANOVA table is:

{\def\LTcaptype{none} % do not increment counter
\begin{longtable}[]{@{}lrrrr@{}}
\toprule\noalign{}
Source & SS & df & MS & \(F\) \\
\midrule\noalign{}
\endhead
\bottomrule\noalign{}
\endlastfoot
Between groups & 995.904 & 3 & 331.968 & 26.094 \\
Within groups / Error & 190.833 & 15 & 12.722 & \\
Total & 1186.737 & 18 & & \\
\end{longtable}
}

The p-value is approximately \[p\approx 3.35\times 10^{-6}.\] This is
very small, so we reject the null hypothesis that all four group means
are equal. There is strong evidence that at least two treatment means
differ.

\end{tcolorbox}

\begin{tcolorbox}[enhanced jigsaw, arc=.35mm, bottomrule=.15mm, bottomtitle=1mm, breakable, colback=white, colbacktitle=quarto-callout-note-color!10!white, colframe=quarto-callout-note-color-frame, coltitle=black, left=2mm, leftrule=.75mm, opacityback=0, opacitybacktitle=0.6, rightrule=.15mm, title=\textcolor{quarto-callout-note-color}{\faInfo}\hspace{0.5em}{Remark}, titlerule=0mm, toprule=.15mm, toptitle=1mm]

\emph{Remark 28}. The ANOVA test tells us that some difference exists
among the treatment means, but it does not identify which treatment
pairs differ. To answer that question, one should use contrasts or
multiple-comparison procedures.

\end{tcolorbox}

\bookmarksetup{startatroot}

\chapter{Practice Problems}\label{practice-problems-8}

This section gives practice questions that reinforce the main
computations and interpretations.

\begin{tcolorbox}[enhanced jigsaw, arc=.35mm, bottomrule=.15mm, bottomtitle=1mm, breakable, colback=white, colbacktitle=quarto-callout-warning-color!10!white, colframe=quarto-callout-warning-color-frame, coltitle=black, left=2mm, leftrule=.75mm, opacityback=0, opacitybacktitle=0.6, rightrule=.15mm, title=\textcolor{quarto-callout-warning-color}{\faExclamationTriangle}\hspace{0.5em}{Practice Problem}, titlerule=0mm, toprule=.15mm, toptitle=1mm]

\textbf{Practice Problem 29} (Dot notation). Suppose there are three
groups with observations \[(4,6,5),\qquad (8,9),\qquad (3,4,5,6).\]
Compute \(T_{i\cdot}\), \(\bar Y_{i\cdot}\), \(T_{\cdot\cdot}\), and
\(\bar Y_{\cdot\cdot}\).

\end{tcolorbox}

\begin{tcolorbox}[enhanced jigsaw, arc=.35mm, bottomrule=.15mm, bottomtitle=1mm, breakable, colback=white, colbacktitle=quarto-callout-caution-color!10!white, colframe=quarto-callout-caution-color-frame, coltitle=black, left=2mm, leftrule=.75mm, opacityback=0, opacitybacktitle=0.6, rightrule=.15mm, title=\textcolor{quarto-callout-caution-color}{\faFire}\hspace{0.5em}{Solution}, titlerule=0mm, toprule=.15mm, toptitle=1mm]

The group totals are \[T_{1\cdot}=4+6+5=15,
    \qquad
    T_{2\cdot}=8+9=17,
    \qquad
    T_{3\cdot}=3+4+5+6=18.\] The group sample sizes are
\(n_1=3,n_2=2,n_3=4\), so \[\bar Y_{1\cdot}=5,
    \qquad
    \bar Y_{2\cdot}=8.5,
    \qquad
    \bar Y_{3\cdot}=4.5.\] The total is \[T_{\cdot\cdot}=15+17+18=50,\]
and \(N=3+2+4=9\), so \[\bar Y_{\cdot\cdot}=\frac{50}{9}=5.556.\]

\end{tcolorbox}

\begin{tcolorbox}[enhanced jigsaw, arc=.35mm, bottomrule=.15mm, bottomtitle=1mm, breakable, colback=white, colbacktitle=quarto-callout-warning-color!10!white, colframe=quarto-callout-warning-color-frame, coltitle=black, left=2mm, leftrule=.75mm, opacityback=0, opacitybacktitle=0.6, rightrule=.15mm, title=\textcolor{quarto-callout-warning-color}{\faExclamationTriangle}\hspace{0.5em}{Practice Problem}, titlerule=0mm, toprule=.15mm, toptitle=1mm]

\textbf{Practice Problem 30} (Identifying contrasts). For four treatment
means \(\theta_1,\theta_2,\theta_3,\theta_4\), determine which of the
following coefficient vectors define contrasts: \[a=(1,-1,0,0),
    \qquad
    b=\left(1,-\frac12,-\frac12,0\right),
    \qquad
    c=(1,1,0,0).\]

\end{tcolorbox}

\begin{tcolorbox}[enhanced jigsaw, arc=.35mm, bottomrule=.15mm, bottomtitle=1mm, breakable, colback=white, colbacktitle=quarto-callout-caution-color!10!white, colframe=quarto-callout-caution-color-frame, coltitle=black, left=2mm, leftrule=.75mm, opacityback=0, opacitybacktitle=0.6, rightrule=.15mm, title=\textcolor{quarto-callout-caution-color}{\faFire}\hspace{0.5em}{Solution}, titlerule=0mm, toprule=.15mm, toptitle=1mm]

A vector defines a contrast if its coefficients sum to zero.
\[1-1+0+0=0,\] so \(a\) is a contrast. \[1-\frac12-\frac12+0=0,\] so
\(b\) is a contrast. \[1+1+0+0=2\neq 0,\] so \(c\) is not a contrast.

\end{tcolorbox}

\begin{tcolorbox}[enhanced jigsaw, arc=.35mm, bottomrule=.15mm, bottomtitle=1mm, breakable, colback=white, colbacktitle=quarto-callout-warning-color!10!white, colframe=quarto-callout-warning-color-frame, coltitle=black, left=2mm, leftrule=.75mm, opacityback=0, opacitybacktitle=0.6, rightrule=.15mm, title=\textcolor{quarto-callout-warning-color}{\faExclamationTriangle}\hspace{0.5em}{Practice Problem}, titlerule=0mm, toprule=.15mm, toptitle=1mm]

\textbf{Practice Problem 31} (Contrast standard error). Suppose \(k=3\),
\(n_1=5,n_2=5,n_3=10\), and \(S_p^2=4\). Find the standard error of
\[\widehat L=\bar Y_{1\cdot}-\frac12\bar Y_{2\cdot}-\frac12\bar Y_{3\cdot}.\]

\end{tcolorbox}

\begin{tcolorbox}[enhanced jigsaw, arc=.35mm, bottomrule=.15mm, bottomtitle=1mm, breakable, colback=white, colbacktitle=quarto-callout-caution-color!10!white, colframe=quarto-callout-caution-color-frame, coltitle=black, left=2mm, leftrule=.75mm, opacityback=0, opacitybacktitle=0.6, rightrule=.15mm, title=\textcolor{quarto-callout-caution-color}{\faFire}\hspace{0.5em}{Solution}, titlerule=0mm, toprule=.15mm, toptitle=1mm]

The coefficients are \[a_1=1,
    \qquad a_2=-\frac12,
    \qquad a_3=-\frac12.\] The standard error is
\[S_p\sqrt{\sum_{i=1}^3\frac{a_i^2}{n_i}}.\] Since \(S_p=\sqrt{4}=2\),
\[\begin{aligned}
\operatorname{se}(\widehat L)
&=2\sqrt{\frac{1^2}{5}+\frac{(1/2)^2}{5}+\frac{(1/2)^2}{10}} \\
&=2\sqrt{0.2+0.05+0.025} \\
&=2\sqrt{0.275}
\approx 1.049.
\end{aligned}\]

\end{tcolorbox}

\begin{tcolorbox}[enhanced jigsaw, arc=.35mm, bottomrule=.15mm, bottomtitle=1mm, breakable, colback=white, colbacktitle=quarto-callout-warning-color!10!white, colframe=quarto-callout-warning-color-frame, coltitle=black, left=2mm, leftrule=.75mm, opacityback=0, opacitybacktitle=0.6, rightrule=.15mm, title=\textcolor{quarto-callout-warning-color}{\faExclamationTriangle}\hspace{0.5em}{Practice Problem}, titlerule=0mm, toprule=.15mm, toptitle=1mm]

\textbf{Practice Problem 32} (ANOVA sums of squares). For an ANOVA
problem with \(k=4\) groups and \(N=20\) observations, suppose
\[\mathrm{SSB}=120,
    \qquad
    \mathrm{SSE}=180.\] Compute \(\mathrm{SST}\), \(\mathrm{MSB}\),
\(\mathrm{MSE}\), and \(F_{\mathrm{obs}}\).

\end{tcolorbox}

\begin{tcolorbox}[enhanced jigsaw, arc=.35mm, bottomrule=.15mm, bottomtitle=1mm, breakable, colback=white, colbacktitle=quarto-callout-caution-color!10!white, colframe=quarto-callout-caution-color-frame, coltitle=black, left=2mm, leftrule=.75mm, opacityback=0, opacitybacktitle=0.6, rightrule=.15mm, title=\textcolor{quarto-callout-caution-color}{\faFire}\hspace{0.5em}{Solution}, titlerule=0mm, toprule=.15mm, toptitle=1mm]

The total sum of squares is
\[\mathrm{SST}=\mathrm{SSB}+\mathrm{SSE}=120+180=300.\] The degrees of
freedom are \[k-1=3,
    \qquad
    N-k=16.\] Thus \[\mathrm{MSB}=\frac{120}{3}=40,
    \qquad
    \mathrm{MSE}=\frac{180}{16}=11.25.\] The observed \(F\) statistic is
\[F_{\mathrm{obs}}=\frac{40}{11.25}=3.556.\]

\end{tcolorbox}

\begin{tcolorbox}[enhanced jigsaw, arc=.35mm, bottomrule=.15mm, bottomtitle=1mm, breakable, colback=white, colbacktitle=quarto-callout-warning-color!10!white, colframe=quarto-callout-warning-color-frame, coltitle=black, left=2mm, leftrule=.75mm, opacityback=0, opacitybacktitle=0.6, rightrule=.15mm, title=\textcolor{quarto-callout-warning-color}{\faExclamationTriangle}\hspace{0.5em}{Practice Problem}, titlerule=0mm, toprule=.15mm, toptitle=1mm]

\textbf{Practice Problem 33} (Interpreting an ANOVA result). Suppose a
one-way ANOVA test rejects \[H_0:\theta_1=\theta_2=\theta_3=\theta_4.\]
What can and cannot be concluded?

\end{tcolorbox}

\begin{tcolorbox}[enhanced jigsaw, arc=.35mm, bottomrule=.15mm, bottomtitle=1mm, breakable, colback=white, colbacktitle=quarto-callout-caution-color!10!white, colframe=quarto-callout-caution-color-frame, coltitle=black, left=2mm, leftrule=.75mm, opacityback=0, opacitybacktitle=0.6, rightrule=.15mm, title=\textcolor{quarto-callout-caution-color}{\faFire}\hspace{0.5em}{Solution}, titlerule=0mm, toprule=.15mm, toptitle=1mm]

We can conclude that the data provide evidence that at least two
treatment means are different. We cannot conclude that all four means
are different. We also cannot conclude which pairs differ without
follow-up inference, such as planned contrasts or multiple-comparison
procedures.

\end{tcolorbox}

\bookmarksetup{startatroot}

\chapter{Summary}\label{summary-10}

This section introduced the foundations of one-way ANOVA.

\begin{tcolorbox}[enhanced jigsaw, arc=.35mm, bottomrule=.15mm, bottomtitle=1mm, breakable, colback=white, colbacktitle=quarto-callout-tip-color!10!white, colframe=quarto-callout-tip-color-frame, coltitle=black, left=2mm, leftrule=.75mm, opacityback=0, opacitybacktitle=0.6, rightrule=.15mm, title=\textcolor{quarto-callout-tip-color}{\faLightbulb}\hspace{0.5em}{Key idea}, titlerule=0mm, toprule=.15mm, toptitle=1mm]

\textbf{Key takeaways.}

\begin{itemize}
\item
  One-way ANOVA generalizes the two-sample pooled \(t\)-test to \(k\)
  groups.
\item
  The cell-means model is \(Y_{ij}=\theta_i+\varepsilon_{ij}\).
\item
  The global ANOVA null is \(H_0:\theta_1=\cdots=\theta_k\).
\item
  Contrasts are linear combinations \(\sum_i a_i\theta_i\) with
  \(\sum_i a_i=0\).
\item
  The ANOVA decomposition is \(\mathrm{SST}=\mathrm{SSB}+\mathrm{SSE}\).
\item
  The \(F\) statistic is
  \(F=(\mathrm{SSB}/(k-1))/(\mathrm{SSE}/(N-k))\).
\item
  Under the null, \(F\sim F_{k-1,N-k}\).
\end{itemize}

\end{tcolorbox}

\bookmarksetup{startatroot}

\chapter{Chapter 19: Analysis of Variance II --- Applications and
Randomized Block
Designs}\label{chapter-19-analysis-of-variance-ii-applications-and-randomized-block-designs}

This chapter continues ANOVA by emphasizing \textbf{applications,
interpretation, experimental design, and randomized block designs}. We
review the one-way ANOVA table, work through complete numerical
examples, discuss why multiple pairwise tests can inflate Type I error,
and introduce blocking as a way to remove nuisance variation from the
error term.

\begin{tcolorbox}[enhanced jigsaw, arc=.35mm, bottomrule=.15mm, bottomtitle=1mm, breakable, colback=white, colbacktitle=quarto-callout-note-color!10!white, colframe=quarto-callout-note-color-frame, coltitle=black, left=2mm, leftrule=.75mm, opacityback=0, opacitybacktitle=0.6, rightrule=.15mm, title=\textcolor{quarto-callout-note-color}{\faInfo}\hspace{0.5em}{Topics}, titlerule=0mm, toprule=.15mm, toptitle=1mm]

One-way ANOVA review; complete ANOVA computations; ANOVA interpretation;
multiple testing; follow-up comparisons; principles of experimental
design; control, randomization, replication, and blocking; randomized
block design model; RBD sums of squares; RBD ANOVA table; treatment and
block effects.

\end{tcolorbox}

\bookmarksetup{startatroot}

\chapter{Overview}\label{overview-12}

This section shows how ANOVA is used in real data analysis and how
experimental design affects the interpretation of ANOVA results.

The main topics are:

\begin{enumerate}
\def\labelenumi{\arabic{enumi}.}
\item
  review and summary of one-way ANOVA;
\item
  complete numerical ANOVA computations;
\item
  applications and interpretation;
\item
  principles of experimental design;
\item
  randomized block designs.
\end{enumerate}

\bookmarksetup{startatroot}

\chapter{Review and Summary of One-Way
ANOVA}\label{review-and-summary-of-one-way-anova}

This section reviews the one-way ANOVA model and the classical ANOVA
table before moving to applications.

\section{Background}\label{background}

One-way ANOVA is used when we want to compare the means of several
populations or treatment groups.

ANOVA is appropriate in two common settings:

\begin{enumerate}
\def\labelenumi{\arabic{enumi}.}
\item
  \(k\) independent random samples are drawn from \(k\) populations;
\item
  \(k\) different treatments are applied to a homogeneous group of
  experimental units, which is subdivided into \(k\) subgroups.
\end{enumerate}

The scientific question is:
\[\text{Do the treatment groups have the same population mean, or does at least one treatment differ?}\]

\section{One-Way ANOVA Model}\label{one-way-anova-model}

The one-way ANOVA model represents each observation as a group mean plus
random error.

Let \(X_{ij}\) be the \(j\)th observation in group \(i\), where
\[i=1,\ldots,k, \qquad j=1,\ldots,n_i.\] The model is
\[X_{ij}=\mu_i+\epsilon_{ij},\] where \(\mu_i\) is the mean of group
\(i\) and \(\epsilon_{ij}\) is the random error.

\begin{tcolorbox}[enhanced jigsaw, arc=.35mm, bottomrule=.15mm, bottomtitle=1mm, breakable, colback=white, colbacktitle=quarto-callout-note-color!10!white, colframe=quarto-callout-note-color-frame, coltitle=black, left=2mm, leftrule=.75mm, opacityback=0, opacitybacktitle=0.6, rightrule=.15mm, title=\textcolor{quarto-callout-note-color}{\faInfo}\hspace{0.5em}{Definition}, titlerule=0mm, toprule=.15mm, toptitle=1mm]

\textbf{Definition 1} (One-way ANOVA assumptions). The classical one-way
ANOVA model assumes:

\begin{enumerate}
\def\labelenumi{\arabic{enumi}.}
\item
  the samples are independently selected from the \(k\) populations;
\item
  the \(k\) populations are approximately normal;
\item
  all \(k\) population variances are equal;
\item
  equivalently,
  \[\epsilon_{ij}\overset{\text{iid}}{\sim} \operatorname{Normal}(0,\sigma^2).\]
\end{enumerate}

\end{tcolorbox}

\section{Classical ANOVA Hypothesis}\label{classical-anova-hypothesis}

The classical ANOVA test asks whether all treatment means are equal.

The hypotheses are \[H_0: \mu_1=\mu_2=\cdots=\mu_k\] versus
\[H_1: \mu_i\neq \mu_j \quad \text{for some } i,j.\]

Rejecting \(H_0\) means that at least one group mean differs from
another. It does not immediately tell us which means differ. Follow-up
comparisons, contrasts, or multiple-comparison procedures are needed for
that question.

\section{Dot Notation and Means}\label{dot-notation-and-means}

Dot notation provides a compact way to write group totals, group means,
and the grand mean.

Let \[T_{i\cdot}=\sum_{j=1}^{n_i}Y_{ij}\] be the total for group \(i\).
Then the group mean is
\[\overline{Y}_{i\cdot}=\frac{T_{i\cdot}}{n_i}=\frac{1}{n_i}\sum_{j=1}^{n_i}Y_{ij}.\]
The total sample size is \[N=\sum_{i=1}^k n_i.\] The grand total is
\[T_{\cdot\cdot}=\sum_{i=1}^{k}\sum_{j=1}^{n_i}Y_{ij}=\sum_{i=1}^k T_{i\cdot},\]
and the grand mean is
\[\overline{Y}_{\cdot\cdot}=\frac{1}{N}\sum_{i=1}^{k}\sum_{j=1}^{n_i}Y_{ij}
    =\frac{1}{N}\sum_{i=1}^{k}T_{i\cdot}
    =\frac{1}{N}\sum_{i=1}^{k}n_i\overline{Y}_{i\cdot}.\]

\section{Partitioning Sums of
Squares}\label{partitioning-sums-of-squares}

The key idea of ANOVA is to split total variation into between-group
variation and within-group variation.

\begin{tcolorbox}[enhanced jigsaw, arc=.35mm, bottomrule=.15mm, bottomtitle=1mm, breakable, colback=white, colbacktitle=quarto-callout-note-color!10!white, colframe=quarto-callout-note-color-frame, coltitle=black, left=2mm, leftrule=.75mm, opacityback=0, opacitybacktitle=0.6, rightrule=.15mm, title=\textcolor{quarto-callout-note-color}{\faInfo}\hspace{0.5em}{Definition}, titlerule=0mm, toprule=.15mm, toptitle=1mm]

\textbf{Definition 2} (Sums of squares). The total sum of squares is
\[\mathrm{SST}_{\mathrm{total}}=\sum_{i=1}^k\sum_{j=1}^{n_i}(Y_{ij}-\overline{Y}_{\cdot\cdot})^2.\]
The between-group sum of squares is
\[\mathrm{SSB}_{\mathrm{between}}=\sum_{i=1}^k n_i(\overline{Y}_{i\cdot}-\overline{Y}_{\cdot\cdot})^2.\]
The within-group or error sum of squares is
\[\mathrm{SSE}_{\mathrm{error}}=\sum_{i=1}^k\sum_{j=1}^{n_i}(Y_{ij}-\overline{Y}_{i\cdot})^2.\]
The fundamental decomposition is
\[\mathrm{SST}_{\mathrm{total}}=\mathrm{SSB}_{\mathrm{between}}+\mathrm{SSE}_{\mathrm{error}}.\]

\end{tcolorbox}

\begin{tcolorbox}[enhanced jigsaw, arc=.35mm, bottomrule=.15mm, bottomtitle=1mm, breakable, colback=white, colbacktitle=quarto-callout-tip-color!10!white, colframe=quarto-callout-tip-color-frame, coltitle=black, left=2mm, leftrule=.75mm, opacityback=0, opacitybacktitle=0.6, rightrule=.15mm, title=\textcolor{quarto-callout-tip-color}{\faLightbulb}\hspace{0.5em}{Key idea}, titlerule=0mm, toprule=.15mm, toptitle=1mm]

Interpretation of the decomposition The between-group sum of squares
measures how far the group means are from the grand mean. The
within-group sum of squares measures random variation around each group
mean. If the between-group variation is large compared with the
within-group variation, this provides evidence that the treatment means
are not all equal.

\end{tcolorbox}

\section{Mean Squares and the F
Statistic}\label{mean-squares-and-the-f-statistic}

The F statistic compares average between-group variation to average
within-group variation.

Define the mean treatment sum of squares and mean error sum of squares
by
\[\mathrm{MSB}=\frac{\mathrm{SSB}}{k-1}, \qquad \mathrm{MSE}=\frac{\mathrm{SSE}}{N-k}.\]
Under the ANOVA assumptions,
\[\frac{\mathrm{SSE}}{\sigma^2}\sim \chi^2_{N-k}.\] If
\(H_0: \mu_1=\cdots=\mu_k\) is true, then
\[\frac{\mathrm{SSB}}{\sigma^2}\sim \chi^2_{k-1}.\] Therefore, under
\(H_0\), \[F=\frac{\mathrm{MSB}}{\mathrm{MSE}}
    =\frac{\mathrm{SSB}/(k-1)}{\mathrm{SSE}/(N-k)}
    \sim F_{k-1,N-k}.\] Large values of \(F\) provide evidence against
\(H_0\).

\section{The One-Way ANOVA Table}\label{the-one-way-anova-table}

The ANOVA table organizes the decomposition of variation and the F-test.

{\def\LTcaptype{none} % do not increment counter
\begin{longtable}[]{@{}lcccc@{}}
\toprule\noalign{}
Source of Variation & Sum of Squares & df & Mean Square & F Statistic \\
\midrule\noalign{}
\endhead
\bottomrule\noalign{}
\endlastfoot
Between groups & \(\mathrm{SSB}\) & \(k-1\) &
\(\mathrm{MSB}=\mathrm{SSB}/(k-1)\) &
\(F_{\mathrm{obs}}=\mathrm{MSB}/\mathrm{MSE}\) \\
Within groups & \(\mathrm{SSE}\) & \(N-k\) &
\(\mathrm{MSE}=\mathrm{SSE}/(N-k)\) & \\
Total & \(\mathrm{SST}\) & \(N-1\) & & \\
\end{longtable}
}

The p-value is
\[p\text{-value}=\mathbb{P}\left(F_{k-1,N-k}>F_{\mathrm{obs}}\right).\]

\bookmarksetup{startatroot}

\chapter{Example: Comparing Three Golf Ball
Brands}\label{example-comparing-three-golf-ball-brands}

This section gives a complete one-way ANOVA computation for comparing
three treatment means.

\begin{tcolorbox}[enhanced jigsaw, arc=.35mm, bottomrule=.15mm, bottomtitle=1mm, breakable, colback=white, colbacktitle=quarto-callout-tip-color!10!white, colframe=quarto-callout-tip-color-frame, coltitle=black, left=2mm, leftrule=.75mm, opacityback=0, opacitybacktitle=0.6, rightrule=.15mm, title=\textcolor{quarto-callout-tip-color}{\faLightbulb}\hspace{0.5em}{Example}, titlerule=0mm, toprule=.15mm, toptitle=1mm]

\textbf{Example 3} (Golf ball distance by brand). A test was conducted
to compare the mean distance, in yards, traveled by three brands of golf
balls hit by a robotic golfer.

{\def\LTcaptype{none} % do not increment counter
\begin{longtable}[]{@{}ccc@{}}
\toprule\noalign{}
Brand A & Brand B & Brand C \\
\midrule\noalign{}
\endhead
\bottomrule\noalign{}
\endlastfoot
251.2 & 263.2 & 269.7 \\
245.1 & 262.9 & 263.2 \\
248.0 & 265.0 & 277.5 \\
251.1 & 254.5 & 267.4 \\
260.5 & 264.3 & 270.5 \\
\end{longtable}
}

Test whether the mean distances for the three brands are equal.

\end{tcolorbox}

\begin{tcolorbox}[enhanced jigsaw, arc=.35mm, bottomrule=.15mm, bottomtitle=1mm, breakable, colback=white, colbacktitle=quarto-callout-caution-color!10!white, colframe=quarto-callout-caution-color-frame, coltitle=black, left=2mm, leftrule=.75mm, opacityback=0, opacitybacktitle=0.6, rightrule=.15mm, title=\textcolor{quarto-callout-caution-color}{\faFire}\hspace{0.5em}{Solution}, titlerule=0mm, toprule=.15mm, toptitle=1mm]

There are \(k=3\) groups with \(n_1=n_2=n_3=5\), so \(N=15\). The sample
summaries are: \[\begin{array}{c|ccc}
\text{Brand} & \bar{x}_i & s_i^2 & n_i \\
\hline
A & 251.18 & 33.487 & 5 \\
B & 261.985 & 18.197 & 5 \\
C & 269.66 & 27.253 & 5
\end{array}\] The grand mean is \[\bar{x}_{\cdot\cdot}
    =\frac{5(251.18)+5(261.985)+5(269.66)}{15}
    \approx 260.94.\] The between-group sum of squares is
\[\begin{aligned}
\mathrm{SSB}
&=\sum_{i=1}^3 n_i(\bar{x}_i-\bar{x}_{\cdot\cdot})^2 \\
&=5(251.18-260.94)^2+5(261.985-260.94)^2+5(269.66-260.94)^2 \\
&\approx 861.89.
\end{aligned}\] The within-group sum of squares is
\[\mathrm{SSE}=\sum_{i=1}^3(n_i-1)s_i^2
    =4(33.487+18.197+27.253)
    \approx 315.75.\] The total sum of squares is
\[\mathrm{SST}=\mathrm{SSB}+\mathrm{SSE}\approx 861.89+315.75=1177.64.\]
The ANOVA table is: \[\begin{array}{lccccc}
\hline
\text{Source} & \text{SS} & \text{df} & \text{MS} & F & p\text{-value} \\
\hline
\text{Treatment} & 861.89 & 2 & 430.945 & 16.37 & \approx 0.00037 \\
\text{Error} & 315.75 & 12 & 26.312 & & \\
\text{Total} & 1177.64 & 14 & & & \\
\hline
\end{array}\] Thus
\[F_{\mathrm{obs}}=\frac{430.945}{26.312}\approx 16.37.\] Since
\(p\approx 0.00037<0.01\), we reject \(H_0\) at the \(1\%\) significance
level. There is strong evidence that at least one brand has a different
mean travel distance.

\end{tcolorbox}

\begin{tcolorbox}[enhanced jigsaw, arc=.35mm, bottomrule=.15mm, bottomtitle=1mm, breakable, colback=white, colbacktitle=quarto-callout-tip-color!10!white, colframe=quarto-callout-tip-color-frame, coltitle=black, left=2mm, leftrule=.75mm, opacityback=0, opacitybacktitle=0.6, rightrule=.15mm, title=\textcolor{quarto-callout-tip-color}{\faLightbulb}\hspace{0.5em}{Key idea}, titlerule=0mm, toprule=.15mm, toptitle=1mm]

How to read the golf-ball result The ANOVA test does not say that all
three brands differ from one another. It says that the three means are
not all equal. A follow-up comparison, such as pairwise contrasts or a
multiple-comparison method, is needed to identify which brand means
differ.

\end{tcolorbox}

\bookmarksetup{startatroot}

\chapter{Example: Smoking and Heart
Rate}\label{example-smoking-and-heart-rate}

This section gives another full ANOVA computation, emphasizing how sums
of squares are computed from treatment totals.

\begin{tcolorbox}[enhanced jigsaw, arc=.35mm, bottomrule=.15mm, bottomtitle=1mm, breakable, colback=white, colbacktitle=quarto-callout-tip-color!10!white, colframe=quarto-callout-tip-color-frame, coltitle=black, left=2mm, leftrule=.75mm, opacityback=0, opacitybacktitle=0.6, rightrule=.15mm, title=\textcolor{quarto-callout-tip-color}{\faLightbulb}\hspace{0.5em}{Example}, titlerule=0mm, toprule=.15mm, toptitle=1mm]

\textbf{Example 4} (Smoking and heart rate). A study compares heart rate
after exercise for four smoking categories. The data are:

{\def\LTcaptype{none} % do not increment counter
\begin{longtable}[]{@{}ccccc@{}}
\toprule\noalign{}
Subject & Nonsmoker & Light Smoker & Moderate Smoker & Heavy Smoker \\
\midrule\noalign{}
\endhead
\bottomrule\noalign{}
\endlastfoot
1 & 69 & 55 & 66 & 91 \\
2 & 52 & 60 & 81 & 72 \\
3 & 71 & 78 & 70 & 81 \\
4 & 58 & 58 & 77 & 67 \\
5 & 59 & 62 & 57 & 95 \\
6 & 65 & 66 & 79 & 84 \\
\(T_j\) & 374 & 379 & 430 & 490 \\
\(\bar{Y}_j\) & 62.3 & 63.2 & 71.7 & 81.7 \\
\end{longtable}
}

Test \[H_0: \mu_1=\mu_2=\mu_3=\mu_4\] against the alternative that at
least one group mean differs, using \(\alpha=0.05\).

\end{tcolorbox}

\begin{tcolorbox}[enhanced jigsaw, arc=.35mm, bottomrule=.15mm, bottomtitle=1mm, breakable, colback=white, colbacktitle=quarto-callout-caution-color!10!white, colframe=quarto-callout-caution-color-frame, coltitle=black, left=2mm, leftrule=.75mm, opacityback=0, opacitybacktitle=0.6, rightrule=.15mm, title=\textcolor{quarto-callout-caution-color}{\faFire}\hspace{0.5em}{Solution}, titlerule=0mm, toprule=.15mm, toptitle=1mm]

There are \(k=4\) groups and \(n_j=6\) observations in each group, so
\(N=24\). The grand mean is
\[\bar{Y}_{\cdot\cdot}=\frac{374+379+430+490}{24}=69.7.\] The treatment
sum of squares is \[\begin{aligned}
\mathrm{SSTR}
&=\sum_{j=1}^{4}n_j(\bar{Y}_j-\bar{Y}_{\cdot\cdot})^2 \\
&=6\left[(62.3-69.7)^2+(63.2-69.7)^2+(71.7-69.7)^2+(81.7-69.7)^2\right] \\
&\approx 1464.125.
\end{aligned}\] The error sum of squares is computed by summing squared
deviations within each group:
\[\mathrm{SSE}=\sum_{j=1}^{4}\sum_{i=1}^{6}(Y_{ij}-\bar{Y}_j)^2\approx 1594.833.\]
The degrees of freedom are \[\mathrm{df}_{\mathrm{treat}}=k-1=3,
    \qquad
    \mathrm{df}_{\mathrm{error}}=N-k=20.\] The mean squares are
\[\mathrm{MST}=\frac{1464.125}{3}=488.04,
    \qquad
    \mathrm{MSE}=\frac{1594.833}{20}=79.74.\] Thus
\[F_{\mathrm{obs}}=\frac{\mathrm{MST}}{\mathrm{MSE}}
    =\frac{488.04}{79.74}\approx 6.12.\] The critical value for
\(\alpha=0.05\) with \((3,20)\) degrees of freedom is approximately
\[F_{0.05;3,20}\approx 3.10.\] Since \[6.12>3.10,\] we reject \(H_0\).
The p-value is approximately \(0.004\), which gives the same conclusion.

The ANOVA table is: \[\begin{array}{lccccc}
\hline
\text{Source} & \mathrm{df}& \text{SS} & \text{MS} & F & p \\
\hline
\text{Treatment} & 3 & 1464.125 & 488.04 & 6.12 & 0.004 \\
\text{Error} & 20 & 1594.833 & 79.74 & & \\
\text{Total} & 23 & 3058.958 & & & \\
\hline
\end{array}\] There is statistically significant evidence that smoking
level affects mean heart rate after exercise. At least one group's mean
heart rate differs from the others.

\end{tcolorbox}

\bookmarksetup{startatroot}

\chapter{Applications and
Interpretation}\label{applications-and-interpretation}

This section explains what ANOVA conclusions mean in applied work.

\section{What ANOVA tells us}\label{what-anova-tells-us}

ANOVA is a global test for equality of means.

A significant ANOVA F-test tells us that the data provide evidence
against \[H_0: \mu_1=\cdots=\mu_k.\] It does not automatically identify
which group is different. For example, if \(k=4\) and the test is
significant, then many patterns are possible:

\begin{itemize}
\item
  one group differs from the other three;
\item
  two groups differ from the other two;
\item
  the groups increase or decrease in order;
\item
  all groups have different means.
\end{itemize}

\section{Why not just run many
t-tests?}\label{why-not-just-run-many-t-tests}

Multiple pairwise t-tests can inflate the Type I error rate.

If each individual test is run at significance level \(0.05\), then the
chance of at least one false positive increases as the number of tests
increases. For two independent tests,
\[\mathbb{P}(\text{at least one Type I error})=1-(1-0.05)^2
    =1-0.95^2\approx 0.0975.\] For three independent tests,
\[1-0.95^3\approx 0.143.\] The ANOVA F-test provides a single global
test with a controlled Type I error rate.

\section{Follow-up analysis after
ANOVA}\label{follow-up-analysis-after-anova}

After rejecting the global ANOVA null, the next step is usually to
examine scientifically meaningful comparisons.

Examples of follow-up questions include:

\begin{itemize}
\item
  Which treatment differs from the control?
\item
  Which pair of treatments differ?
\item
  Is the average response under several treatments larger than the
  control?
\item
  Is there an ordered pattern across treatment levels?
\end{itemize}

Such questions can be formulated using contrasts, confidence intervals,
or multiple comparison procedures.

\bookmarksetup{startatroot}

\chapter{Design of Experiments}\label{design-of-experiments}

This section connects ANOVA to the design of experiments, because good
design reduces unexplained variation.

Effective experimental design is built on four principles traditionally
associated with Ronald Fisher.

\begin{tcolorbox}[enhanced jigsaw, arc=.35mm, bottomrule=.15mm, bottomtitle=1mm, breakable, colback=white, colbacktitle=quarto-callout-note-color!10!white, colframe=quarto-callout-note-color-frame, coltitle=black, left=2mm, leftrule=.75mm, opacityback=0, opacitybacktitle=0.6, rightrule=.15mm, title=\textcolor{quarto-callout-note-color}{\faInfo}\hspace{0.5em}{Definition}, titlerule=0mm, toprule=.15mm, toptitle=1mm]

\textbf{Definition 5} (Four principles of experimental design). The four
main principles are:

\begin{enumerate}
\def\labelenumi{\arabic{enumi}.}
\item
  \textbf{Control}: Researchers control variables that might influence
  the outcome, often using a control group as a baseline.
\item
  \textbf{Randomization}: Subjects or experimental units are randomly
  assigned to treatment groups.
\item
  \textbf{Replication}: The experiment is repeated across multiple units
  so results are not driven by chance.
\item
  \textbf{Blocking}: Similar experimental units are grouped into blocks
  to reduce nuisance variability within each block.
\end{enumerate}

\end{tcolorbox}

\begin{tcolorbox}[enhanced jigsaw, arc=.35mm, bottomrule=.15mm, bottomtitle=1mm, breakable, colback=white, colbacktitle=quarto-callout-tip-color!10!white, colframe=quarto-callout-tip-color-frame, coltitle=black, left=2mm, leftrule=.75mm, opacityback=0, opacitybacktitle=0.6, rightrule=.15mm, title=\textcolor{quarto-callout-tip-color}{\faLightbulb}\hspace{0.5em}{Key idea}, titlerule=0mm, toprule=.15mm, toptitle=1mm]

Why blocking matters Without blocking, a large nuisance factor may be
hidden inside the error term. With blocking, we explicitly model that
nuisance factor and remove its contribution from the error sum of
squares. This often makes treatment comparisons more sensitive.

\end{tcolorbox}

\bookmarksetup{startatroot}

\chapter{Randomized Block Designs}\label{randomized-block-designs}

This section introduces randomized block designs as a two-way structure
with treatments and blocks.

\section{Motivation}\label{motivation}

A randomized block design is useful when experimental units vary because
of a nuisance factor.

Examples of possible blocks include:

\begin{itemize}
\item
  fields or plots of land;
\item
  patients with similar risk levels;
\item
  machines;
\item
  batches;
\item
  days;
\item
  classrooms or sections.
\end{itemize}

The goal is to compare treatments within relatively homogeneous blocks.

\begin{tcolorbox}[enhanced jigsaw, arc=.35mm, bottomrule=.15mm, bottomtitle=1mm, breakable, colback=white, colbacktitle=quarto-callout-note-color!10!white, colframe=quarto-callout-note-color-frame, coltitle=black, left=2mm, leftrule=.75mm, opacityback=0, opacitybacktitle=0.6, rightrule=.15mm, title=\textcolor{quarto-callout-note-color}{\faInfo}\hspace{0.5em}{Definition}, titlerule=0mm, toprule=.15mm, toptitle=1mm]

\textbf{Definition 6} (Randomized block design). A randomized block
design groups similar experimental units into blocks. Each block
receives all treatments, usually exactly once. Treatments are randomized
within each block.

\end{tcolorbox}

\section{When to Use a Randomized Block
Design}\label{when-to-use-a-randomized-block-design}

A randomized block design is appropriate when blocks capture an
important source of variation.

Use RBD when:

\begin{enumerate}
\def\labelenumi{\arabic{enumi}.}
\item
  experimental units can be partitioned into homogeneous blocks;
\item
  variability among blocks is expected to be large;
\item
  variability within blocks is expected to be small;
\item
  the researcher wants to increase power without increasing sample size.
\end{enumerate}

\begin{tcolorbox}[enhanced jigsaw, arc=.35mm, bottomrule=.15mm, bottomtitle=1mm, breakable, colback=white, colbacktitle=quarto-callout-tip-color!10!white, colframe=quarto-callout-tip-color-frame, coltitle=black, left=2mm, leftrule=.75mm, opacityback=0, opacitybacktitle=0.6, rightrule=.15mm, title=\textcolor{quarto-callout-tip-color}{\faLightbulb}\hspace{0.5em}{Example}, titlerule=0mm, toprule=.15mm, toptitle=1mm]

\textbf{Example 7} (Fertilizers in blocks). Suppose three fertilizers
\(A,B,C\) are applied in four blocks. Each block receives all
treatments:

{\def\LTcaptype{none} % do not increment counter
\begin{longtable}[]{@{}cccc@{}}
\toprule\noalign{}
Block & A & B & C \\
\midrule\noalign{}
\endhead
\bottomrule\noalign{}
\endlastfoot
1 & 18 & 22 & 27 \\
2 & 16 & 25 & 23 \\
3 & 21 & 19 & 30 \\
4 & 20 & 24 & 28 \\
\end{longtable}
}

The blocks may differ in productivity. The randomized block design
removes between-block variation and improves the sensitivity of
treatment comparisons.

\end{tcolorbox}

\begin{tcolorbox}[enhanced jigsaw, arc=.35mm, bottomrule=.15mm, bottomtitle=1mm, breakable, colback=white, colbacktitle=quarto-callout-caution-color!10!white, colframe=quarto-callout-caution-color-frame, coltitle=black, left=2mm, leftrule=.75mm, opacityback=0, opacitybacktitle=0.6, rightrule=.15mm, title=\textcolor{quarto-callout-caution-color}{\faFire}\hspace{0.5em}{Solution}, titlerule=0mm, toprule=.15mm, toptitle=1mm]

The treatment means are
\[\bar{Y}_{A\cdot}=\frac{18+16+21+20}{4}=18.75,\]
\[\bar{Y}_{B\cdot}=\frac{22+25+19+24}{4}=22.50,\]
\[\bar{Y}_{C\cdot}=\frac{27+23+30+28}{4}=27.00.\] These means suggest
that fertilizer \(C\) may produce the largest response. However, ANOVA
for a randomized block design formally compares treatments after
accounting for block-to-block differences.

\end{tcolorbox}

\section{Structure of a Randomized Block
Design}\label{structure-of-a-randomized-block-design}

The randomized block model includes both treatment effects and block
effects.

Let there be \(t\) treatments and \(b\) blocks. The observation for
treatment \(i\) in block \(j\) is modeled as
\[Y_{ij}=\mu+\tau_i+\beta_j+\epsilon_{ij},
    \qquad i=1,\ldots,t,
    \quad j=1,\ldots,b.\] Here:

\begin{itemize}
\item
  \(\mu\) is the overall mean;
\item
  \(\tau_i\) is the treatment effect;
\item
  \(\beta_j\) is the block effect;
\item
  \(\epsilon_{ij}\overset{\mathrm{iid}}{\sim}\operatorname{Normal}(0,\sigma^2)\).
\end{itemize}

Because the model is overparameterized, identifiability constraints are
typically imposed: \[\sum_{i=1}^{t}\tau_i=0,
    \qquad
    \sum_{j=1}^{b}\beta_j=0.\]

\section{Hypotheses for Treatments}\label{hypotheses-for-treatments}

The main goal is usually to test whether treatment effects differ after
accounting for blocks.

The treatment hypothesis is \[H_0: \tau_1=\tau_2=\cdots=\tau_t=0\]
versus \[H_1: \tau_i\neq 0 \quad \text{for at least one } i.\]
Equivalently, we test whether all treatment means are equal after
adjusting for block effects.

\bookmarksetup{startatroot}

\chapter{Computing Sums of Squares in a Randomized Block
Design}\label{computing-sums-of-squares-in-a-randomized-block-design}

This section gives the computational formulas for an RBD ANOVA table.

Let \[T_i=\sum_{j=1}^{b}Y_{ij}\] be the total for treatment \(i\), and
let \[B_j=\sum_{i=1}^{t}Y_{ij}\] be the total for block \(j\). Let
\[G=\sum_{i=1}^{t}\sum_{j=1}^{b}Y_{ij}\] be the grand total. There are
\(tb\) total observations.

\begin{tcolorbox}[enhanced jigsaw, arc=.35mm, bottomrule=.15mm, bottomtitle=1mm, breakable, colback=white, colbacktitle=quarto-callout-note-color!10!white, colframe=quarto-callout-note-color-frame, coltitle=black, left=2mm, leftrule=.75mm, opacityback=0, opacitybacktitle=0.6, rightrule=.15mm, title=\textcolor{quarto-callout-note-color}{\faInfo}\hspace{0.5em}{Definition}, titlerule=0mm, toprule=.15mm, toptitle=1mm]

\textbf{Definition 8} (RBD sums of squares). The total sum of squares is
\[\mathrm{SST}_{\mathrm{total}}
    =\sum_{i=1}^{t}\sum_{j=1}^{b}Y_{ij}^2-\frac{G^2}{tb}.\] The
treatment sum of squares is \[\mathrm{SST}_{\mathrm{trt}}
    =\frac{1}{b}\sum_{i=1}^{t}T_i^2-\frac{G^2}{tb}.\] The block sum of
squares is \[\mathrm{SSB}_{\mathrm{block}}
    =\frac{1}{t}\sum_{j=1}^{b}B_j^2-\frac{G^2}{tb}.\] The error sum of
squares is
\[\mathrm{SSE}=\mathrm{SST}_{\mathrm{total}}-\mathrm{SST}_{\mathrm{trt}}-\mathrm{SSB}_{\mathrm{block}}.\]

\end{tcolorbox}

\section{ANOVA Table for Randomized Block
Design}\label{anova-table-for-randomized-block-design}

The RBD ANOVA table separates total variation into treatment variation,
block variation, and error variation.

{\def\LTcaptype{none} % do not increment counter
\begin{longtable}[]{@{}lcccc@{}}
\toprule\noalign{}
Source & df & SS & MS & F Statistic \\
\midrule\noalign{}
\endhead
\bottomrule\noalign{}
\endlastfoot
Treatments & \(t-1\) & \(\mathrm{SST}_{\mathrm{trt}}\) &
\(\mathrm{MST}=\mathrm{SST}_{\mathrm{trt}}/(t-1)\) &
\(F=\mathrm{MST}/\mathrm{MSE}\) \\
Blocks & \(b-1\) & \(\mathrm{SSB}_{\mathrm{block}}\) &
\(\mathrm{MSB}=\mathrm{SSB}_{\mathrm{block}}/(b-1)\) & \\
Error & \((t-1)(b-1)\) & \(\mathrm{SSE}\) &
\(\mathrm{MSE}=\mathrm{SSE}/[(t-1)(b-1)]\) & \\
Total & \(tb-1\) & \(\mathrm{SST}_{\mathrm{total}}\) & & \\
\end{longtable}
}

The treatment F-test rejects \(H_0\) for large values of
\[F_{\mathrm{trt}}=\frac{\mathrm{MST}}{\mathrm{MSE}}.\] Under the null
hypothesis and model assumptions,
\[F_{\mathrm{trt}}\sim F_{t-1,(t-1)(b-1)}.\]

\bookmarksetup{startatroot}

\chapter{Example: Complete Randomized Block Design
Computation}\label{example-complete-randomized-block-design-computation}

This section completes the fertilizer example using the randomized block
design formulas.

\begin{tcolorbox}[enhanced jigsaw, arc=.35mm, bottomrule=.15mm, bottomtitle=1mm, breakable, colback=white, colbacktitle=quarto-callout-tip-color!10!white, colframe=quarto-callout-tip-color-frame, coltitle=black, left=2mm, leftrule=.75mm, opacityback=0, opacitybacktitle=0.6, rightrule=.15mm, title=\textcolor{quarto-callout-tip-color}{\faLightbulb}\hspace{0.5em}{Example}, titlerule=0mm, toprule=.15mm, toptitle=1mm]

\textbf{Example 9} (RBD computation for fertilizers). Consider the
fertilizer data:

{\def\LTcaptype{none} % do not increment counter
\begin{longtable}[]{@{}cccc@{}}
\toprule\noalign{}
Block & A & B & C \\
\midrule\noalign{}
\endhead
\bottomrule\noalign{}
\endlastfoot
1 & 18 & 22 & 27 \\
2 & 16 & 25 & 23 \\
3 & 21 & 19 & 30 \\
4 & 20 & 24 & 28 \\
\end{longtable}
}

Compute the randomized block ANOVA table and test for treatment
differences.

\end{tcolorbox}

\begin{tcolorbox}[enhanced jigsaw, arc=.35mm, bottomrule=.15mm, bottomtitle=1mm, breakable, colback=white, colbacktitle=quarto-callout-caution-color!10!white, colframe=quarto-callout-caution-color-frame, coltitle=black, left=2mm, leftrule=.75mm, opacityback=0, opacitybacktitle=0.6, rightrule=.15mm, title=\textcolor{quarto-callout-caution-color}{\faFire}\hspace{0.5em}{Solution}, titlerule=0mm, toprule=.15mm, toptitle=1mm]

There are \(t=3\) treatments and \(b=4\) blocks. The treatment totals
are \[T_A=75,\qquad T_B=90,\qquad T_C=108.\] The block totals are
\[B_1=67,\quad B_2=64,\quad B_3=70,\quad B_4=72.\] The grand total is
\[G=273,\] and \(tb=12\).

First compute \[\sum_{i,j}Y_{ij}^2
    =18^2+22^2+27^2+16^2+25^2+23^2+21^2+19^2+30^2+20^2+24^2+28^2
    =6385.\] Thus \[\mathrm{SST}_{\mathrm{total}}
    =6385-\frac{273^2}{12}
    =6385-6210.75
    =174.25.\] The treatment sum of squares is \[\begin{aligned}
\mathrm{SST}_{\mathrm{trt}}
&=\frac{1}{4}(75^2+90^2+108^2)-\frac{273^2}{12} \\
&=\frac{25389}{4}-6210.75 \\
&=136.50.
\end{aligned}\] The block sum of squares is
\[\mathrm{SSB}_{\mathrm{block}}
    =\frac{1}{3}(67^2+64^2+70^2+72^2)-6210.75
    =13.583.\] The error sum of squares is
\[\mathrm{SSE}=174.25-136.50-13.583=24.167.\] The degrees of freedom are
\[\mathrm{df}_{\mathrm{trt}}=t-1=2,
    \quad
    \mathrm{df}_{\mathrm{block}}=b-1=3,
    \quad
    \mathrm{df}_{\mathrm{error}}=(t-1)(b-1)=6,
    \quad
    \mathrm{df}_{\mathrm{total}}=11.\] The mean squares are
\[\mathrm{MST}=\frac{136.50}{2}=68.250,
    \qquad
    \mathrm{MSB}=\frac{13.583}{3}=4.528,
    \qquad
    \mathrm{MSE}=\frac{24.167}{6}=4.028.\] Therefore
\[F_{\mathrm{trt}}=\frac{68.250}{4.028}\approx 16.95.\] The ANOVA table
is: \[\begin{array}{lcccc}
\hline
\text{Source} & \mathrm{df}& \text{SS} & \text{MS} & F \\
\hline
\text{Treatments} & 2 & 136.500 & 68.250 & 16.95 \\
\text{Blocks} & 3 & 13.583 & 4.528 & \\
\text{Error} & 6 & 24.167 & 4.028 & \\
\text{Total} & 11 & 174.250 & & \\
\hline
\end{array}\] The treatment p-value is \[\mathbb{P}(F_{2,6}>16.95),\]
which is small. Thus we reject the null hypothesis of equal treatment
effects. There is strong evidence that the fertilizers do not all have
the same mean response.

\end{tcolorbox}

\bookmarksetup{startatroot}

\chapter{Lab: Randomized Block Designs with
Python}\label{lab-randomized-block-designs-with-python}

This section summarizes the lab-style setup for randomized block
designs.

A typical Python lab can use:

\begin{itemize}
\item
  treatments: \(3\) fertilizers, for example F1, F2, and NF (no
  fertilizer);
\item
  blocks: \(3\) crop types, for example corn, rice, and wheat;
\item
  response: number of months the crops stayed healthy.
\end{itemize}

The model is \[Y_{ij}=\mu+\tau_i+\beta_j+\epsilon_{ij},
    \qquad \epsilon_{ij}\sim \operatorname{Normal}(0,\sigma^2).\]

\begin{tcolorbox}[enhanced jigsaw, arc=.35mm, bottomrule=.15mm, bottomtitle=1mm, breakable, colback=white, colbacktitle=quarto-callout-tip-color!10!white, colframe=quarto-callout-tip-color-frame, coltitle=black, left=2mm, leftrule=.75mm, opacityback=0, opacitybacktitle=0.6, rightrule=.15mm, title=\textcolor{quarto-callout-tip-color}{\faLightbulb}\hspace{0.5em}{Key idea}, titlerule=0mm, toprule=.15mm, toptitle=1mm]

Practical workflow for an RBD analysis

\begin{enumerate}
\def\labelenumi{\arabic{enumi}.}
\item
  Plot the response by treatment and by block.
\item
  Fit a model with treatment and block effects.
\item
  Construct the ANOVA table.
\item
  Test the treatment effect using \(F=\mathrm{MST}/\mathrm{MSE}\).
\item
  Interpret the treatment result after accounting for block variation.
\end{enumerate}

\end{tcolorbox}

\bookmarksetup{startatroot}

\chapter{Practice Problems}\label{practice-problems-9}

This section gives practice problems that reinforce one-way ANOVA and
randomized block designs.

\begin{tcolorbox}[enhanced jigsaw, arc=.35mm, bottomrule=.15mm, bottomtitle=1mm, breakable, colback=white, colbacktitle=quarto-callout-tip-color!10!white, colframe=quarto-callout-tip-color-frame, coltitle=black, left=2mm, leftrule=.75mm, opacityback=0, opacitybacktitle=0.6, rightrule=.15mm, title=\textcolor{quarto-callout-tip-color}{\faLightbulb}\hspace{0.5em}{Practice problem}, titlerule=0mm, toprule=.15mm, toptitle=1mm]

\textbf{Practice Problem 10} (One-way ANOVA from summary statistics).
Suppose three groups have sample sizes \(n_1=n_2=n_3=5\), sample means
\[\bar{x}_1=10,
    \qquad
    \bar{x}_2=14,
    \qquad
    \bar{x}_3=18,\] and sample variances \[s_1^2=4,
    \qquad
    s_2^2=5,
    \qquad
    s_3^2=6.\] Compute the one-way ANOVA table and the observed F
statistic.

\end{tcolorbox}

\begin{tcolorbox}[enhanced jigsaw, arc=.35mm, bottomrule=.15mm, bottomtitle=1mm, breakable, colback=white, colbacktitle=quarto-callout-caution-color!10!white, colframe=quarto-callout-caution-color-frame, coltitle=black, left=2mm, leftrule=.75mm, opacityback=0, opacitybacktitle=0.6, rightrule=.15mm, title=\textcolor{quarto-callout-caution-color}{\faFire}\hspace{0.5em}{Solution}, titlerule=0mm, toprule=.15mm, toptitle=1mm]

The total sample size is \(N=15\) and \(k=3\). The grand mean is
\[\bar{x}_{\cdot\cdot}=\frac{5(10)+5(14)+5(18)}{15}=14.\] The
between-group sum of squares is
\[\mathrm{SSB}=5(10-14)^2+5(14-14)^2+5(18-14)^2=160.\] The within-group
sum of squares is \[\mathrm{SSE}=(5-1)4+(5-1)5+(5-1)6=16+20+24=60.\]
Thus \[\mathrm{MSB}=\frac{160}{2}=80,
    \qquad
    \mathrm{MSE}=\frac{60}{12}=5.\] The observed F statistic is
\[F_{\mathrm{obs}}=\frac{80}{5}=16.\] The ANOVA table is:
\[\begin{array}{lcccc}
\hline
\text{Source} & \text{SS} & \mathrm{df}& \text{MS} & F \\
\hline
\text{Between} & 160 & 2 & 80 & 16 \\
\text{Within} & 60 & 12 & 5 & \\
\text{Total} & 220 & 14 & & \\
\hline
\end{array}\]

\end{tcolorbox}

\begin{tcolorbox}[enhanced jigsaw, arc=.35mm, bottomrule=.15mm, bottomtitle=1mm, breakable, colback=white, colbacktitle=quarto-callout-tip-color!10!white, colframe=quarto-callout-tip-color-frame, coltitle=black, left=2mm, leftrule=.75mm, opacityback=0, opacitybacktitle=0.6, rightrule=.15mm, title=\textcolor{quarto-callout-tip-color}{\faLightbulb}\hspace{0.5em}{Practice problem}, titlerule=0mm, toprule=.15mm, toptitle=1mm]

\textbf{Practice Problem 11} (Why blocking can help). Explain why a
randomized block design can have more power than a completely randomized
one-way ANOVA design.

\end{tcolorbox}

\begin{tcolorbox}[enhanced jigsaw, arc=.35mm, bottomrule=.15mm, bottomtitle=1mm, breakable, colback=white, colbacktitle=quarto-callout-caution-color!10!white, colframe=quarto-callout-caution-color-frame, coltitle=black, left=2mm, leftrule=.75mm, opacityback=0, opacitybacktitle=0.6, rightrule=.15mm, title=\textcolor{quarto-callout-caution-color}{\faFire}\hspace{0.5em}{Solution}, titlerule=0mm, toprule=.15mm, toptitle=1mm]

A randomized block design separates nuisance variation due to blocks
from random error. In a completely randomized design, block-to-block
variation is absorbed into the error term. This can make
\(\mathrm{MSE}\) large and reduce the F statistic for treatment effects.
In an RBD, block variation is removed through the block sum of squares,
often making the residual error smaller. A smaller \(\mathrm{MSE}\)
makes it easier to detect genuine treatment differences.

\end{tcolorbox}

\begin{tcolorbox}[enhanced jigsaw, arc=.35mm, bottomrule=.15mm, bottomtitle=1mm, breakable, colback=white, colbacktitle=quarto-callout-tip-color!10!white, colframe=quarto-callout-tip-color-frame, coltitle=black, left=2mm, leftrule=.75mm, opacityback=0, opacitybacktitle=0.6, rightrule=.15mm, title=\textcolor{quarto-callout-tip-color}{\faLightbulb}\hspace{0.5em}{Practice problem}, titlerule=0mm, toprule=.15mm, toptitle=1mm]

\textbf{Practice Problem 12} (RBD degrees of freedom). A randomized
block design has \(t=5\) treatments and \(b=6\) blocks. Find the degrees
of freedom for treatment, block, error, and total sources of variation.

\end{tcolorbox}

\begin{tcolorbox}[enhanced jigsaw, arc=.35mm, bottomrule=.15mm, bottomtitle=1mm, breakable, colback=white, colbacktitle=quarto-callout-caution-color!10!white, colframe=quarto-callout-caution-color-frame, coltitle=black, left=2mm, leftrule=.75mm, opacityback=0, opacitybacktitle=0.6, rightrule=.15mm, title=\textcolor{quarto-callout-caution-color}{\faFire}\hspace{0.5em}{Solution}, titlerule=0mm, toprule=.15mm, toptitle=1mm]

The degrees of freedom are: \[\mathrm{df}_{\mathrm{treatment}}=t-1=4,\]
\[\mathrm{df}_{\mathrm{block}}=b-1=5,\]
\[\mathrm{df}_{\mathrm{error}}=(t-1)(b-1)=4\cdot 5=20,\] and
\[\mathrm{df}_{\mathrm{total}}=tb-1=30-1=29.\] They add correctly:
\[4+5+20=29.\]

\end{tcolorbox}

\begin{tcolorbox}[enhanced jigsaw, arc=.35mm, bottomrule=.15mm, bottomtitle=1mm, breakable, colback=white, colbacktitle=quarto-callout-tip-color!10!white, colframe=quarto-callout-tip-color-frame, coltitle=black, left=2mm, leftrule=.75mm, opacityback=0, opacitybacktitle=0.6, rightrule=.15mm, title=\textcolor{quarto-callout-tip-color}{\faLightbulb}\hspace{0.5em}{Practice problem}, titlerule=0mm, toprule=.15mm, toptitle=1mm]

\textbf{Practice Problem 13} (Treatment F statistic in an RBD). In a
randomized block design with \(t=4\) treatments and \(b=5\) blocks,
suppose \[\mathrm{SST}_{\mathrm{trt}}=90,
    \qquad
    \mathrm{SSE}=60.\] Compute the treatment F statistic.

\end{tcolorbox}

\begin{tcolorbox}[enhanced jigsaw, arc=.35mm, bottomrule=.15mm, bottomtitle=1mm, breakable, colback=white, colbacktitle=quarto-callout-caution-color!10!white, colframe=quarto-callout-caution-color-frame, coltitle=black, left=2mm, leftrule=.75mm, opacityback=0, opacitybacktitle=0.6, rightrule=.15mm, title=\textcolor{quarto-callout-caution-color}{\faFire}\hspace{0.5em}{Solution}, titlerule=0mm, toprule=.15mm, toptitle=1mm]

The treatment degrees of freedom are \[t-1=3.\] The error degrees of
freedom are \[(t-1)(b-1)=3\cdot 4=12.\] Therefore
\[\mathrm{MST}=\frac{90}{3}=30,
    \qquad
    \mathrm{MSE}=\frac{60}{12}=5.\] The treatment F statistic is
\[F=\frac{\mathrm{MST}}{\mathrm{MSE}}=\frac{30}{5}=6.\]

\end{tcolorbox}

\bookmarksetup{startatroot}

\chapter{Summary}\label{summary-11}

This final section summarizes the key ideas from ANOVA applications and
randomized block designs.

\begin{itemize}
\item
  One-way ANOVA compares \(k\) group means using the ratio
  \(F=\mathrm{MSB}/\mathrm{MSE}\).
\item
  The ANOVA F-test is a global test: it detects whether at least one
  group mean differs.
\item
  The sums of squares decomposition is
  \[\mathrm{SST}=\mathrm{SSB}+\mathrm{SSE}.\]
\item
  Multiple pairwise t-tests can inflate Type I error.
\item
  Good experimental design uses control, randomization, replication, and
  blocking.
\item
  Randomized block designs reduce nuisance variation by comparing
  treatments within homogeneous blocks.
\item
  In an RBD, total variation is decomposed into treatment variation,
  block variation, and error variation.
\end{itemize}

\bookmarksetup{startatroot}

\chapter{Chapter 20: Bayesian
Inference}\label{chapter-20-bayesian-inference}

This chapter collects and organizes the Bayesian ideas used throughout
point estimation, hypothesis testing, interval estimation, and decision
theory. The main message is that once we know the posterior
distribution, we can derive point estimates, tests, credible intervals,
and optimal decisions from one common framework.

\begin{tcolorbox}[enhanced jigsaw, arc=.35mm, bottomrule=.15mm, bottomtitle=1mm, breakable, colback=white, colbacktitle=quarto-callout-note-color!10!white, colframe=quarto-callout-note-color-frame, coltitle=black, left=2mm, leftrule=.75mm, opacityback=0, opacitybacktitle=0.6, rightrule=.15mm, title=\textcolor{quarto-callout-note-color}{\faInfo}\hspace{0.5em}{Topics}, titlerule=0mm, toprule=.15mm, toptitle=1mm]

Bayesian model ingredients; prior, likelihood, posterior, and marginal
likelihood; Bayes estimators; posterior mean, median, mode, and MAP;
conjugate priors; beta-binomial, gamma-Poisson, and normal-normal
models; Bayes risk; Bayesian tests; credible intervals; highest
posterior density regions; loss-function interpretation; practice
problems and solutions.

\end{tcolorbox}

\bookmarksetup{startatroot}

\chapter{Overview}\label{overview-13}

This section collects the Bayesian material that appeared across point
estimation, evaluating estimators, hypothesis testing, evaluating tests,
interval estimation, and evaluating intervals.

Bayesian inference gives a unified method for learning from data.
Instead of treating the parameter as an unknown but fixed constant, the
Bayesian approach represents uncertainty about the parameter using a
probability distribution. Data update this distribution through Bayes'
rule.

\begin{tcolorbox}[enhanced jigsaw, arc=.35mm, bottomrule=.15mm, bottomtitle=1mm, breakable, colback=white, colbacktitle=quarto-callout-tip-color!10!white, colframe=quarto-callout-tip-color-frame, coltitle=black, left=2mm, leftrule=.75mm, opacityback=0, opacitybacktitle=0.6, rightrule=.15mm, title=\textcolor{quarto-callout-tip-color}{\faLightbulb}\hspace{0.5em}{Key idea}, titlerule=0mm, toprule=.15mm, toptitle=1mm]

Main message Bayesian inference is the rule
\[\text{posterior} \propto \text{likelihood} \times \text{prior}.\] Once
the posterior distribution is known, point estimation, hypothesis
testing, and interval estimation can all be derived from it.

\end{tcolorbox}

The core idea is simple but powerful. Before observing data, we describe
prior information by a prior distribution. After observing data, we
update the prior into a posterior distribution. All Bayesian inference
is then based on this posterior distribution.

\bookmarksetup{startatroot}

\chapter{Bayesian Model Ingredients}\label{bayesian-model-ingredients}

This section introduces the basic objects of Bayesian inference:
likelihood, prior, marginal likelihood, and posterior.

Suppose the observed data are
\[X=(X_1,\ldots,X_n), \qquad x=(x_1,\ldots,x_n),\] with sampling density
or mass function \[f(x\mid \theta).\] The parameter \(\theta\) is
unknown. In the Bayesian approach, \(\theta\) is assigned a probability
distribution.

\begin{tcolorbox}[enhanced jigsaw, arc=.35mm, bottomrule=.15mm, bottomtitle=1mm, breakable, colback=white, colbacktitle=quarto-callout-note-color!10!white, colframe=quarto-callout-note-color-frame, coltitle=black, left=2mm, leftrule=.75mm, opacityback=0, opacitybacktitle=0.6, rightrule=.15mm, title=\textcolor{quarto-callout-note-color}{\faInfo}\hspace{0.5em}{Definition}, titlerule=0mm, toprule=.15mm, toptitle=1mm]

\textbf{Definition 1} (Prior distribution). The \textbf{prior
distribution} \(\pi(\theta)\) describes our uncertainty or belief about
\(\theta\) before observing the data.

\end{tcolorbox}

\begin{tcolorbox}[enhanced jigsaw, arc=.35mm, bottomrule=.15mm, bottomtitle=1mm, breakable, colback=white, colbacktitle=quarto-callout-note-color!10!white, colframe=quarto-callout-note-color-frame, coltitle=black, left=2mm, leftrule=.75mm, opacityback=0, opacitybacktitle=0.6, rightrule=.15mm, title=\textcolor{quarto-callout-note-color}{\faInfo}\hspace{0.5em}{Definition}, titlerule=0mm, toprule=.15mm, toptitle=1mm]

\textbf{Definition 2} (Likelihood). After observing \(X=x\), the
\textbf{likelihood function} is \[L(\theta\mid x)=f(x\mid \theta).\] It
measures how compatible each parameter value \(\theta\) is with the
observed data.

\end{tcolorbox}

\begin{tcolorbox}[enhanced jigsaw, arc=.35mm, bottomrule=.15mm, bottomtitle=1mm, breakable, colback=white, colbacktitle=quarto-callout-note-color!10!white, colframe=quarto-callout-note-color-frame, coltitle=black, left=2mm, leftrule=.75mm, opacityback=0, opacitybacktitle=0.6, rightrule=.15mm, title=\textcolor{quarto-callout-note-color}{\faInfo}\hspace{0.5em}{Definition}, titlerule=0mm, toprule=.15mm, toptitle=1mm]

\textbf{Definition 3} (Marginal distribution). The \textbf{marginal
distribution} or \textbf{prior predictive distribution} of the data is
\[m(x)=\int f(x\mid \theta)\pi(\theta)\,d\theta.\] It is the normalizing
constant that makes the posterior integrate to one.

\end{tcolorbox}

\begin{tcolorbox}[enhanced jigsaw, arc=.35mm, bottomrule=.15mm, bottomtitle=1mm, breakable, colback=white, colbacktitle=quarto-callout-note-color!10!white, colframe=quarto-callout-note-color-frame, coltitle=black, left=2mm, leftrule=.75mm, opacityback=0, opacitybacktitle=0.6, rightrule=.15mm, title=\textcolor{quarto-callout-note-color}{\faInfo}\hspace{0.5em}{Definition}, titlerule=0mm, toprule=.15mm, toptitle=1mm]

\textbf{Definition 4} (Posterior distribution). The \textbf{posterior
distribution} of \(\theta\) given \(X=x\) is
\[\pi(\theta\mid x)=\frac{f(x\mid \theta)\pi(\theta)}{m(x)}
=\frac{f(x\mid \theta)\pi(\theta)}{\int f(x\mid \theta')\pi(\theta')\,d\theta'}.\]

\end{tcolorbox}

\begin{tcolorbox}[enhanced jigsaw, arc=.35mm, bottomrule=.15mm, bottomtitle=1mm, breakable, colback=white, colbacktitle=quarto-callout-tip-color!10!white, colframe=quarto-callout-tip-color-frame, coltitle=black, left=2mm, leftrule=.75mm, opacityback=0, opacitybacktitle=0.6, rightrule=.15mm, title=\textcolor{quarto-callout-tip-color}{\faLightbulb}\hspace{0.5em}{Key idea}, titlerule=0mm, toprule=.15mm, toptitle=1mm]

Bayesian updating The denominator \(m(x)\) does not depend on
\(\theta\). Therefore, for many calculations, we write
\[\pi(\theta\mid x)\propto f(x\mid \theta)\pi(\theta).\] This means that
the posterior is proportional to likelihood times prior.

\end{tcolorbox}

\begin{tcolorbox}[enhanced jigsaw, arc=.35mm, bottomrule=.15mm, bottomtitle=1mm, breakable, colback=white, colbacktitle=quarto-callout-tip-color!10!white, colframe=quarto-callout-tip-color-frame, coltitle=black, left=2mm, leftrule=.75mm, opacityback=0, opacitybacktitle=0.6, rightrule=.15mm, title=\textcolor{quarto-callout-tip-color}{\faLightbulb}\hspace{0.5em}{Example}, titlerule=0mm, toprule=.15mm, toptitle=1mm]

\textbf{Example 5} (Bayesian ingredients for coin tossing). Suppose a
coin has unknown probability \(p\) of heads. We toss the coin \(n\)
times and observe \(x\) heads. The likelihood is
\[f(x\mid p)=\binom{n}{x}p^x(1-p)^{n-x}, \qquad 0<p<1.\] If the prior is
\(p\sim\operatorname{Beta}(\alpha,\beta)\), then
\[\pi(p)=\frac{1}{B(\alpha,\beta)}p^{\alpha-1}(1-p)^{\beta-1}.\] Find
the posterior distribution.

\end{tcolorbox}

\begin{tcolorbox}[enhanced jigsaw, arc=.35mm, bottomrule=.15mm, bottomtitle=1mm, breakable, colback=white, colbacktitle=quarto-callout-caution-color!10!white, colframe=quarto-callout-caution-color-frame, coltitle=black, left=2mm, leftrule=.75mm, opacityback=0, opacitybacktitle=0.6, rightrule=.15mm, title=\textcolor{quarto-callout-caution-color}{\faFire}\hspace{0.5em}{Solution}, titlerule=0mm, toprule=.15mm, toptitle=1mm]

Using Bayes' rule,
\[\pi(p\mid x)\propto \binom{n}{x}p^x(1-p)^{n-x}\cdot p^{\alpha-1}(1-p)^{\beta-1}.\]
Ignoring constants that do not depend on \(p\),
\[\pi(p\mid x)\propto p^{\alpha+x-1}(1-p)^{\beta+n-x-1}.\] This is the
kernel of a beta distribution, so
\[p\mid X=x\sim \operatorname{Beta}(\alpha+x,\beta+n-x).\]

\end{tcolorbox}

\bookmarksetup{startatroot}

\chapter{Bayesian Point Estimation}\label{bayesian-point-estimation}

This section explains how to produce a single-number estimate from a
posterior distribution.

In classical point estimation, estimators such as the method of moments
estimator and the maximum likelihood estimator are functions of the
data. In Bayesian inference, once the posterior distribution is
obtained, we can summarize it by its mean, median, mode, or another
loss-optimal point.

\section{Posterior mean, median, and
mode}\label{posterior-mean-median-and-mode}

This subsection introduces three common posterior summaries used as
Bayes point estimates.

\begin{tcolorbox}[enhanced jigsaw, arc=.35mm, bottomrule=.15mm, bottomtitle=1mm, breakable, colback=white, colbacktitle=quarto-callout-note-color!10!white, colframe=quarto-callout-note-color-frame, coltitle=black, left=2mm, leftrule=.75mm, opacityback=0, opacitybacktitle=0.6, rightrule=.15mm, title=\textcolor{quarto-callout-note-color}{\faInfo}\hspace{0.5em}{Definition}, titlerule=0mm, toprule=.15mm, toptitle=1mm]

\textbf{Definition 6} (Posterior mean). The \textbf{posterior mean}
estimator is
\[\widehat\theta_B=\mathbb{E}(\theta\mid x)=\int \theta\pi(\theta\mid x)\,d\theta.\]
It is commonly called the Bayes estimator under squared error loss.

\end{tcolorbox}

\begin{tcolorbox}[enhanced jigsaw, arc=.35mm, bottomrule=.15mm, bottomtitle=1mm, breakable, colback=white, colbacktitle=quarto-callout-note-color!10!white, colframe=quarto-callout-note-color-frame, coltitle=black, left=2mm, leftrule=.75mm, opacityback=0, opacitybacktitle=0.6, rightrule=.15mm, title=\textcolor{quarto-callout-note-color}{\faInfo}\hspace{0.5em}{Definition}, titlerule=0mm, toprule=.15mm, toptitle=1mm]

\textbf{Definition 7} (Posterior median). A \textbf{posterior median} is
any value \(m\) satisfying
\[\mathbb{P}(\theta\le m\mid x)\ge \frac12, \qquad \mathbb{P}(\theta\ge m\mid x)\ge \frac12.\]
It is the Bayes estimator under absolute error loss.

\end{tcolorbox}

\begin{tcolorbox}[enhanced jigsaw, arc=.35mm, bottomrule=.15mm, bottomtitle=1mm, breakable, colback=white, colbacktitle=quarto-callout-note-color!10!white, colframe=quarto-callout-note-color-frame, coltitle=black, left=2mm, leftrule=.75mm, opacityback=0, opacitybacktitle=0.6, rightrule=.15mm, title=\textcolor{quarto-callout-note-color}{\faInfo}\hspace{0.5em}{Definition}, titlerule=0mm, toprule=.15mm, toptitle=1mm]

\textbf{Definition 8} (Posterior mode and MAP estimator). The
\textbf{maximum a posteriori} estimator is
\[\widehat\theta_{\mathrm{MAP}}=\arg\max_{\theta}\pi(\theta\mid x).\]
Since \[\pi(\theta\mid x)\propto f(x\mid \theta)\pi(\theta),\] we can
also write
\[\widehat\theta_{\mathrm{MAP}}=\arg\max_{\theta}\{\log f(x\mid \theta)+\log \pi(\theta)\}.\]

\end{tcolorbox}

\begin{tcolorbox}[enhanced jigsaw, arc=.35mm, bottomrule=.15mm, bottomtitle=1mm, breakable, colback=white, colbacktitle=quarto-callout-tip-color!10!white, colframe=quarto-callout-tip-color-frame, coltitle=black, left=2mm, leftrule=.75mm, opacityback=0, opacitybacktitle=0.6, rightrule=.15mm, title=\textcolor{quarto-callout-tip-color}{\faLightbulb}\hspace{0.5em}{Key idea}, titlerule=0mm, toprule=.15mm, toptitle=1mm]

MAP versus MLE The maximum likelihood estimator maximizes only the
likelihood:
\[\widehat\theta_{\mathrm{MLE}}=\arg\max_\theta \log f(x\mid \theta).\]
The MAP estimator maximizes likelihood plus prior information:
\[\widehat\theta_{\mathrm{MAP}}=\arg\max_\theta \{\log f(x\mid \theta)+\log \pi(\theta)\}.\]
A flat prior makes MAP behave like MLE.

\end{tcolorbox}

\section{Loss functions and Bayes
estimators}\label{loss-functions-and-bayes-estimators}

This subsection connects Bayesian point estimates to decision theory.

A Bayesian point estimator can be chosen by minimizing posterior
expected loss. Let \(a\) be an action, interpreted as an estimate of
\(\theta\). A loss function \(L(\theta,a)\) measures the penalty of
estimating \(\theta\) by \(a\).

\begin{tcolorbox}[enhanced jigsaw, arc=.35mm, bottomrule=.15mm, bottomtitle=1mm, breakable, colback=white, colbacktitle=quarto-callout-note-color!10!white, colframe=quarto-callout-note-color-frame, coltitle=black, left=2mm, leftrule=.75mm, opacityback=0, opacitybacktitle=0.6, rightrule=.15mm, title=\textcolor{quarto-callout-note-color}{\faInfo}\hspace{0.5em}{Definition}, titlerule=0mm, toprule=.15mm, toptitle=1mm]

\textbf{Definition 9} (Posterior expected loss). Given posterior
\(\pi(\theta\mid x)\), the posterior expected loss of action \(a\) is
\[\rho(a\mid x)=\mathbb{E}[L(\theta,a)\mid x]=\int L(\theta,a)\pi(\theta\mid x)\,d\theta.\]
A Bayes action minimizes \(\rho(a\mid x)\).

\end{tcolorbox}

\begin{tcolorbox}[enhanced jigsaw, arc=.35mm, bottomrule=.15mm, bottomtitle=1mm, breakable, colback=white, colbacktitle=quarto-callout-note-color!10!white, colframe=quarto-callout-note-color-frame, coltitle=black, left=2mm, leftrule=.75mm, opacityback=0, opacitybacktitle=0.6, rightrule=.15mm, title=\textcolor{quarto-callout-note-color}{\faInfo}\hspace{0.5em}{Theorem}, titlerule=0mm, toprule=.15mm, toptitle=1mm]

\textbf{Theorem 10} (Common Bayes estimators). \emph{For a real-valued
parameter \(\theta\):}

\begin{enumerate}
\def\labelenumi{\arabic{enumi}.}
\item
  \emph{Under squared error loss \(L(\theta,a)=(a-\theta)^2\), the Bayes
  estimator is the posterior mean.}
\item
  \emph{Under absolute error loss \(L(\theta,a)=|a-\theta|\), the Bayes
  estimator is a posterior median.}
\item
  \emph{Under 0--1 style local loss, the Bayes estimator is a posterior
  mode, or MAP estimator.}
\end{enumerate}

\end{tcolorbox}

\begin{proof}
\emph{Proof.} For squared error loss,
\[\rho(a\mid x)=\mathbb{E}[(a-\theta)^2\mid x]
=a^2-2a\mathbb{E}(\theta\mid x)+\mathbb{E}(\theta^2\mid x).\]
Differentiating with respect to \(a\) gives
\[2a-2\mathbb{E}(\theta\mid x)=0,\] so \(a=\mathbb{E}(\theta\mid x)\).
The absolute loss result follows because a median minimizes expected
absolute deviation. The mode result follows because maximizing posterior
mass in a small neighborhood is equivalent to maximizing the posterior
density.~◻
\end{proof}

\begin{tcolorbox}[enhanced jigsaw, arc=.35mm, bottomrule=.15mm, bottomtitle=1mm, breakable, colback=white, colbacktitle=quarto-callout-tip-color!10!white, colframe=quarto-callout-tip-color-frame, coltitle=black, left=2mm, leftrule=.75mm, opacityback=0, opacitybacktitle=0.6, rightrule=.15mm, title=\textcolor{quarto-callout-tip-color}{\faLightbulb}\hspace{0.5em}{Example}, titlerule=0mm, toprule=.15mm, toptitle=1mm]

\textbf{Example 11} (Beta-binomial posterior mean, variance, and mode).
Suppose \(X\sim\operatorname{Binomial}(n,p)\) and
\(p\sim\operatorname{Beta}(\alpha,\beta)\). If \(X=x\), find the
posterior mean, variance, and mode of \(p\).

\end{tcolorbox}

\begin{tcolorbox}[enhanced jigsaw, arc=.35mm, bottomrule=.15mm, bottomtitle=1mm, breakable, colback=white, colbacktitle=quarto-callout-caution-color!10!white, colframe=quarto-callout-caution-color-frame, coltitle=black, left=2mm, leftrule=.75mm, opacityback=0, opacitybacktitle=0.6, rightrule=.15mm, title=\textcolor{quarto-callout-caution-color}{\faFire}\hspace{0.5em}{Solution}, titlerule=0mm, toprule=.15mm, toptitle=1mm]

The posterior is
\[p\mid x\sim \operatorname{Beta}(\alpha+x,\beta+n-x).\] Let
\[\alpha^*=\alpha+x, \qquad \beta^*=\beta+n-x.\] Then
\[\mathbb{E}(p\mid x)=\frac{\alpha^*}{\alpha^*+\beta^*}
=\frac{\alpha+x}{\alpha+\beta+n}.\] The posterior variance is
\[\operatorname{Var}(p\mid x)=\frac{\alpha^*\beta^*}{(\alpha^*+\beta^*)^2(\alpha^*+\beta^*+1)}
=\frac{(\alpha+x)(\beta+n-x)}{(\alpha+\beta+n)^2(\alpha+\beta+n+1)}.\]
If \(\alpha^*>1\) and \(\beta^*>1\), the posterior mode is
\[\frac{\alpha^*-1}{\alpha^*+\beta^*-2}
=\frac{\alpha+x-1}{\alpha+\beta+n-2}.\]

\end{tcolorbox}

\begin{tcolorbox}[enhanced jigsaw, arc=.35mm, bottomrule=.15mm, bottomtitle=1mm, breakable, colback=white, colbacktitle=quarto-callout-tip-color!10!white, colframe=quarto-callout-tip-color-frame, coltitle=black, left=2mm, leftrule=.75mm, opacityback=0, opacitybacktitle=0.6, rightrule=.15mm, title=\textcolor{quarto-callout-tip-color}{\faLightbulb}\hspace{0.5em}{Example}, titlerule=0mm, toprule=.15mm, toptitle=1mm]

\textbf{Example 12} (Pseudo-count interpretation). Suppose
\(p\sim\operatorname{Beta}(3,5)\) and we observe \(x=5\) heads in
\(n=6\) tosses. Find the posterior distribution and posterior mean.

\end{tcolorbox}

\begin{tcolorbox}[enhanced jigsaw, arc=.35mm, bottomrule=.15mm, bottomtitle=1mm, breakable, colback=white, colbacktitle=quarto-callout-caution-color!10!white, colframe=quarto-callout-caution-color-frame, coltitle=black, left=2mm, leftrule=.75mm, opacityback=0, opacitybacktitle=0.6, rightrule=.15mm, title=\textcolor{quarto-callout-caution-color}{\faFire}\hspace{0.5em}{Solution}, titlerule=0mm, toprule=.15mm, toptitle=1mm]

The posterior is
\[p\mid X=5\sim\operatorname{Beta}(3+5,5+6-5)=\operatorname{Beta}(8,6).\]
The posterior mean is
\[\mathbb{E}(p\mid X=5)=\frac{8}{8+6}=\frac{8}{14}=0.5714.\] The prior
contributes \(3\) prior successes and \(5\) prior failures. The data
contribute \(5\) successes and \(1\) failure. Thus the posterior has
\(8\) pseudo-successes and \(6\) pseudo-failures.

\end{tcolorbox}

\bookmarksetup{startatroot}

\chapter{Conjugate Priors}\label{conjugate-priors-1}

This section studies conjugate priors, which make Bayesian updating
especially simple.

A conjugate prior is useful because the posterior stays in the same
distribution family as the prior. Then Bayesian updating often becomes a
simple rule of adding data summaries to prior hyperparameters.

\begin{tcolorbox}[enhanced jigsaw, arc=.35mm, bottomrule=.15mm, bottomtitle=1mm, breakable, colback=white, colbacktitle=quarto-callout-note-color!10!white, colframe=quarto-callout-note-color-frame, coltitle=black, left=2mm, leftrule=.75mm, opacityback=0, opacitybacktitle=0.6, rightrule=.15mm, title=\textcolor{quarto-callout-note-color}{\faInfo}\hspace{0.5em}{Definition}, titlerule=0mm, toprule=.15mm, toptitle=1mm]

\textbf{Definition 13} (Conjugate prior). For a likelihood function
\(f(x\mid \theta)\), a prior family \(\pi(\theta)\) is called
\textbf{conjugate} if the posterior \(\pi(\theta\mid x)\) belongs to the
same family as the prior.

\end{tcolorbox}

\begin{tcolorbox}[enhanced jigsaw, arc=.35mm, bottomrule=.15mm, bottomtitle=1mm, breakable, colback=white, colbacktitle=quarto-callout-tip-color!10!white, colframe=quarto-callout-tip-color-frame, coltitle=black, left=2mm, leftrule=.75mm, opacityback=0, opacitybacktitle=0.6, rightrule=.15mm, title=\textcolor{quarto-callout-tip-color}{\faLightbulb}\hspace{0.5em}{Key idea}, titlerule=0mm, toprule=.15mm, toptitle=1mm]

Why conjugacy helps If a prior is conjugate, then we can often skip
difficult integration. The posterior distribution is known up to updated
parameters, and the update rule is usually easy to compute.

\end{tcolorbox}

\section{Beta-binomial conjugacy}\label{beta-binomial-conjugacy}

This subsection reviews the most important conjugate pair for
proportions.

\begin{tcolorbox}[enhanced jigsaw, arc=.35mm, bottomrule=.15mm, bottomtitle=1mm, breakable, colback=white, colbacktitle=quarto-callout-note-color!10!white, colframe=quarto-callout-note-color-frame, coltitle=black, left=2mm, leftrule=.75mm, opacityback=0, opacitybacktitle=0.6, rightrule=.15mm, title=\textcolor{quarto-callout-note-color}{\faInfo}\hspace{0.5em}{Theorem}, titlerule=0mm, toprule=.15mm, toptitle=1mm]

\textbf{Theorem 14} (Beta-binomial model). \emph{If
\[X\mid p\sim\operatorname{Binomial}(n,p), \qquad p\sim\operatorname{Beta}(\alpha,\beta),\]
then \[p\mid X=x\sim\operatorname{Beta}(\alpha+x,\beta+n-x).\]}

\end{tcolorbox}

\begin{tcolorbox}[enhanced jigsaw, arc=.35mm, bottomrule=.15mm, bottomtitle=1mm, breakable, colback=white, colbacktitle=quarto-callout-tip-color!10!white, colframe=quarto-callout-tip-color-frame, coltitle=black, left=2mm, leftrule=.75mm, opacityback=0, opacitybacktitle=0.6, rightrule=.15mm, title=\textcolor{quarto-callout-tip-color}{\faLightbulb}\hspace{0.5em}{Example}, titlerule=0mm, toprule=.15mm, toptitle=1mm]

\textbf{Example 15} (Bayesian estimator for a Bernoulli probability).
Suppose
\(X_1,\ldots,X_n\overset{\text{iid}}{\sim}\operatorname{Bernoulli}(p)\)
and \(S=\sum_{i=1}^n X_i\). Let
\(p\sim\operatorname{Beta}(\alpha,\beta)\). Find the Bayes estimator
under squared error loss.

\end{tcolorbox}

\begin{tcolorbox}[enhanced jigsaw, arc=.35mm, bottomrule=.15mm, bottomtitle=1mm, breakable, colback=white, colbacktitle=quarto-callout-caution-color!10!white, colframe=quarto-callout-caution-color-frame, coltitle=black, left=2mm, leftrule=.75mm, opacityback=0, opacitybacktitle=0.6, rightrule=.15mm, title=\textcolor{quarto-callout-caution-color}{\faFire}\hspace{0.5em}{Solution}, titlerule=0mm, toprule=.15mm, toptitle=1mm]

Since \(S\mid p\sim\operatorname{Binomial}(n,p)\), the posterior is
\[p\mid S\sim\operatorname{Beta}(\alpha+S,\beta+n-S).\] Under squared
error loss, the Bayes estimator is the posterior mean:
\[\widehat p_B=\mathbb{E}(p\mid S)=\frac{\alpha+S}{\alpha+\beta+n}.\]
This shrinks the sample proportion \(S/n\) toward the prior mean
\(\alpha/(\alpha+\beta)\).

\end{tcolorbox}

\section{Gamma-Poisson conjugacy}\label{gamma-poisson-conjugacy}

This subsection reviews the conjugate prior for a Poisson rate.

We use the rate parametrization for the gamma distribution:
\[\lambda\sim\operatorname{Gamma}(\alpha,\beta), \qquad
\pi(\lambda)=\frac{\beta^\alpha}{\Gamma(\alpha)}\lambda^{\alpha-1}e^{-\beta\lambda},
\qquad \lambda>0.\]

\begin{tcolorbox}[enhanced jigsaw, arc=.35mm, bottomrule=.15mm, bottomtitle=1mm, breakable, colback=white, colbacktitle=quarto-callout-note-color!10!white, colframe=quarto-callout-note-color-frame, coltitle=black, left=2mm, leftrule=.75mm, opacityback=0, opacitybacktitle=0.6, rightrule=.15mm, title=\textcolor{quarto-callout-note-color}{\faInfo}\hspace{0.5em}{Theorem}, titlerule=0mm, toprule=.15mm, toptitle=1mm]

\textbf{Theorem 16} (Gamma-Poisson model). \emph{If
\[X_1,\ldots,X_n\mid \lambda \overset{\text{iid}}{\sim}\operatorname{Poisson}(\lambda),
\qquad \lambda\sim\operatorname{Gamma}(\alpha,\beta),\] then
\[\lambda\mid x\sim \operatorname{Gamma}\left(\alpha+\sum_{i=1}^n x_i,\,\beta+n\right).\]}

\end{tcolorbox}

\begin{proof}
\emph{Proof.} The likelihood is
\[f(x\mid \lambda)=\prod_{i=1}^n \frac{e^{-\lambda}\lambda^{x_i}}{x_i!}
\propto e^{-n\lambda}\lambda^{\sum x_i}.\] Multiplying by the gamma
prior gives
\[\pi(\lambda\mid x)\propto e^{-n\lambda}\lambda^{\sum x_i}\cdot \lambda^{\alpha-1}e^{-\beta\lambda}
=\lambda^{\alpha+\sum x_i-1}e^{-(\beta+n)\lambda}.\] This is the kernel
of \(\operatorname{Gamma}(\alpha+\sum x_i,\beta+n)\).~◻
\end{proof}

\begin{tcolorbox}[enhanced jigsaw, arc=.35mm, bottomrule=.15mm, bottomtitle=1mm, breakable, colback=white, colbacktitle=quarto-callout-tip-color!10!white, colframe=quarto-callout-tip-color-frame, coltitle=black, left=2mm, leftrule=.75mm, opacityback=0, opacitybacktitle=0.6, rightrule=.15mm, title=\textcolor{quarto-callout-tip-color}{\faLightbulb}\hspace{0.5em}{Example}, titlerule=0mm, toprule=.15mm, toptitle=1mm]

\textbf{Example 17} (Poisson posterior). Suppose \(n=10\) Poisson
observations have total count \(\sum x_i=26\). Let the prior be
\(\lambda\sim\operatorname{Gamma}(2,1)\). Find the posterior mean.

\end{tcolorbox}

\begin{tcolorbox}[enhanced jigsaw, arc=.35mm, bottomrule=.15mm, bottomtitle=1mm, breakable, colback=white, colbacktitle=quarto-callout-caution-color!10!white, colframe=quarto-callout-caution-color-frame, coltitle=black, left=2mm, leftrule=.75mm, opacityback=0, opacitybacktitle=0.6, rightrule=.15mm, title=\textcolor{quarto-callout-caution-color}{\faFire}\hspace{0.5em}{Solution}, titlerule=0mm, toprule=.15mm, toptitle=1mm]

The posterior is
\[\lambda\mid x\sim \operatorname{Gamma}(2+26,1+10)=\operatorname{Gamma}(28,11).\]
For a gamma distribution with rate \(\beta\), the mean is
\(\alpha/\beta\). Therefore,
\[\mathbb{E}(\lambda\mid x)=\frac{28}{11}\approx 2.545.\]

\end{tcolorbox}

\section{Normal-normal conjugacy}\label{normal-normal-conjugacy}

This subsection reviews the conjugate model for a normal mean with known
variance.

Suppose
\[X_1,\ldots,X_n\mid \mu \overset{\text{iid}}{\sim}\operatorname{Normal}(\mu,\sigma^2),\]
where \(\sigma^2\) is known. Let the prior be
\[\mu\sim\operatorname{Normal}(\mu_0,\sigma_0^2).\]

\begin{tcolorbox}[enhanced jigsaw, arc=.35mm, bottomrule=.15mm, bottomtitle=1mm, breakable, colback=white, colbacktitle=quarto-callout-note-color!10!white, colframe=quarto-callout-note-color-frame, coltitle=black, left=2mm, leftrule=.75mm, opacityback=0, opacitybacktitle=0.6, rightrule=.15mm, title=\textcolor{quarto-callout-note-color}{\faInfo}\hspace{0.5em}{Theorem}, titlerule=0mm, toprule=.15mm, toptitle=1mm]

\textbf{Theorem 18} (Normal-normal posterior). \emph{The posterior
distribution of \(\mu\) is normal:
\[\mu\mid x\sim\operatorname{Normal}(m_n,v_n),\] where
\[v_n=\left(\frac{n}{\sigma^2}+\frac{1}{\sigma_0^2}\right)^{-1}
=\frac{\sigma^2\sigma_0^2}{\sigma^2+n\sigma_0^2},\] and
\[m_n=v_n\left(\frac{n\overline x}{\sigma^2}+\frac{\mu_0}{\sigma_0^2}\right)
=\frac{\sigma_0^2\sum_{i=1}^n x_i+\sigma^2\mu_0}{n\sigma_0^2+\sigma^2}.\]}

\end{tcolorbox}

\begin{tcolorbox}[enhanced jigsaw, arc=.35mm, bottomrule=.15mm, bottomtitle=1mm, breakable, colback=white, colbacktitle=quarto-callout-tip-color!10!white, colframe=quarto-callout-tip-color-frame, coltitle=black, left=2mm, leftrule=.75mm, opacityback=0, opacitybacktitle=0.6, rightrule=.15mm, title=\textcolor{quarto-callout-tip-color}{\faLightbulb}\hspace{0.5em}{Key idea}, titlerule=0mm, toprule=.15mm, toptitle=1mm]

Weighted average interpretation The posterior mean is a weighted average
of the sample mean \(\overline x\) and the prior mean \(\mu_0\):
\[m_n=\frac{n\sigma_0^2}{n\sigma_0^2+\sigma^2}\overline x
+\frac{\sigma^2}{n\sigma_0^2+\sigma^2}\mu_0.\] If
\(\sigma_0^2\to\infty\), the prior becomes diffuse and
\(m_n\to\overline x\), the MLE.

\end{tcolorbox}

\begin{tcolorbox}[enhanced jigsaw, arc=.35mm, bottomrule=.15mm, bottomtitle=1mm, breakable, colback=white, colbacktitle=quarto-callout-tip-color!10!white, colframe=quarto-callout-tip-color-frame, coltitle=black, left=2mm, leftrule=.75mm, opacityback=0, opacitybacktitle=0.6, rightrule=.15mm, title=\textcolor{quarto-callout-tip-color}{\faLightbulb}\hspace{0.5em}{Example}, titlerule=0mm, toprule=.15mm, toptitle=1mm]

\textbf{Example 19} (MAP for a normal mean). Suppose
\(X_1,\ldots,X_n\mid \mu\overset{\text{iid}}{\sim}\operatorname{Normal}(\mu,\sigma^2)\)
with known \(\sigma^2\), and suppose
\(\mu\sim\operatorname{Normal}(\mu_0,\sigma_0^2)\). Find the MAP
estimator.

\end{tcolorbox}

\begin{tcolorbox}[enhanced jigsaw, arc=.35mm, bottomrule=.15mm, bottomtitle=1mm, breakable, colback=white, colbacktitle=quarto-callout-caution-color!10!white, colframe=quarto-callout-caution-color-frame, coltitle=black, left=2mm, leftrule=.75mm, opacityback=0, opacitybacktitle=0.6, rightrule=.15mm, title=\textcolor{quarto-callout-caution-color}{\faFire}\hspace{0.5em}{Solution}, titlerule=0mm, toprule=.15mm, toptitle=1mm]

The posterior is normal with mean \(m_n\) and variance \(v_n\). A normal
density is maximized at its mean, so \[\widehat\mu_{\mathrm{MAP}}=m_n
=\frac{\sigma_0^2\sum_{i=1}^n x_i+\sigma^2\mu_0}{n\sigma_0^2+\sigma^2}.\]
Equivalently, \[\widehat\mu_{\mathrm{MAP}}
=\frac{n\sigma_0^2}{n\sigma_0^2+\sigma^2}\overline x
+\frac{\sigma^2}{n\sigma_0^2+\sigma^2}\mu_0.\] Because the posterior is
normal, the posterior mean and posterior mode are the same.

\end{tcolorbox}

\bookmarksetup{startatroot}

\chapter{Bayes Risk and Evaluating
Estimators}\label{bayes-risk-and-evaluating-estimators}

This section explains how Bayesian decision theory evaluates estimators
by averaging risk over a prior distribution.

In frequentist evaluation, the risk \(R(\theta,\delta)\) is viewed as a
function of the unknown parameter \(\theta\). In Bayesian evaluation, we
average this risk over the prior distribution.

\begin{tcolorbox}[enhanced jigsaw, arc=.35mm, bottomrule=.15mm, bottomtitle=1mm, breakable, colback=white, colbacktitle=quarto-callout-note-color!10!white, colframe=quarto-callout-note-color-frame, coltitle=black, left=2mm, leftrule=.75mm, opacityback=0, opacitybacktitle=0.6, rightrule=.15mm, title=\textcolor{quarto-callout-note-color}{\faInfo}\hspace{0.5em}{Definition}, titlerule=0mm, toprule=.15mm, toptitle=1mm]

\textbf{Definition 20} (Risk function). For decision rule \(\delta(X)\)
and loss function \(L(\theta,a)\), the risk function is
\[R(\theta,\delta)=\mathbb{E}_\theta[L(\theta,\delta(X))].\]

\end{tcolorbox}

\begin{tcolorbox}[enhanced jigsaw, arc=.35mm, bottomrule=.15mm, bottomtitle=1mm, breakable, colback=white, colbacktitle=quarto-callout-note-color!10!white, colframe=quarto-callout-note-color-frame, coltitle=black, left=2mm, leftrule=.75mm, opacityback=0, opacitybacktitle=0.6, rightrule=.15mm, title=\textcolor{quarto-callout-note-color}{\faInfo}\hspace{0.5em}{Definition}, titlerule=0mm, toprule=.15mm, toptitle=1mm]

\textbf{Definition 21} (Bayes risk). Given prior distribution
\(\pi(\theta)\), the \textbf{Bayes risk} of \(\delta\) is
\[r(\pi,\delta)=\int R(\theta,\delta)\pi(\theta)\,d\theta.\]
Equivalently,
\[r(\pi,\delta)=\int\int L(\theta,\delta(x))f(x\mid \theta)\pi(\theta)\,dx\,d\theta.\]

\end{tcolorbox}

\begin{tcolorbox}[enhanced jigsaw, arc=.35mm, bottomrule=.15mm, bottomtitle=1mm, breakable, colback=white, colbacktitle=quarto-callout-note-color!10!white, colframe=quarto-callout-note-color-frame, coltitle=black, left=2mm, leftrule=.75mm, opacityback=0, opacitybacktitle=0.6, rightrule=.15mm, title=\textcolor{quarto-callout-note-color}{\faInfo}\hspace{0.5em}{Theorem}, titlerule=0mm, toprule=.15mm, toptitle=1mm]

\textbf{Theorem 22} (Posterior risk minimization). \emph{A Bayes rule
can be found by minimizing posterior expected loss for each observed
\(x\):
\[\delta_B(x)\in \arg\min_a \int L(\theta,a)\pi(\theta\mid x)\,d\theta.\]}

\end{tcolorbox}

\begin{tcolorbox}[enhanced jigsaw, arc=.35mm, bottomrule=.15mm, bottomtitle=1mm, breakable, colback=white, colbacktitle=quarto-callout-tip-color!10!white, colframe=quarto-callout-tip-color-frame, coltitle=black, left=2mm, leftrule=.75mm, opacityback=0, opacitybacktitle=0.6, rightrule=.15mm, title=\textcolor{quarto-callout-tip-color}{\faLightbulb}\hspace{0.5em}{Example}, titlerule=0mm, toprule=.15mm, toptitle=1mm]

\textbf{Example 23} (MSE of a beta-binomial Bayes estimator). Let
\(X_1,\ldots,X_n\overset{\text{iid}}{\sim}\operatorname{Bernoulli}(p)\)
and \(S=\sum X_i\). Consider
\[\widehat p_B=\frac{\alpha+S}{\alpha+\beta+n}.\] Find its frequentist
MSE as a function of \(p\).

\end{tcolorbox}

\begin{tcolorbox}[enhanced jigsaw, arc=.35mm, bottomrule=.15mm, bottomtitle=1mm, breakable, colback=white, colbacktitle=quarto-callout-caution-color!10!white, colframe=quarto-callout-caution-color-frame, coltitle=black, left=2mm, leftrule=.75mm, opacityback=0, opacitybacktitle=0.6, rightrule=.15mm, title=\textcolor{quarto-callout-caution-color}{\faFire}\hspace{0.5em}{Solution}, titlerule=0mm, toprule=.15mm, toptitle=1mm]

Since \(S\sim\operatorname{Binomial}(n,p)\),
\[\mathbb{E}_p(\widehat p_B)=\frac{\alpha+np}{\alpha+\beta+n},
\qquad
\operatorname{Var}_p(\widehat p_B)=\frac{np(1-p)}{(\alpha+\beta+n)^2}.\]
The bias is \[\operatorname{Bias}_p(\widehat p_B)
=\frac{\alpha+np}{\alpha+\beta+n}-p
=\frac{\alpha-(\alpha+\beta)p}{\alpha+\beta+n}.\] Therefore,
\[\operatorname{MSE}_p(\widehat p_B)
=\frac{np(1-p)}{(\alpha+\beta+n)^2}
+\left(\frac{\alpha-(\alpha+\beta)p}{\alpha+\beta+n}\right)^2.\] This
shows the bias-variance tradeoff introduced by the prior.

\end{tcolorbox}

\begin{tcolorbox}[enhanced jigsaw, arc=.35mm, bottomrule=.15mm, bottomtitle=1mm, breakable, colback=white, colbacktitle=quarto-callout-note-color!10!white, colframe=quarto-callout-note-color-frame, coltitle=black, left=2mm, leftrule=.75mm, opacityback=0, opacitybacktitle=0.6, rightrule=.15mm, title=\textcolor{quarto-callout-note-color}{\faInfo}\hspace{0.5em}{Remark}, titlerule=0mm, toprule=.15mm, toptitle=1mm]

\emph{Remark 24} (Why shrinkage can help). The estimator
\(\widehat p_B\) is usually biased as a frequentist estimator, but it
can have smaller MSE than the MLE \(\widehat p=S/n\) for some values of
\(p\), especially when \(n\) is small and the prior information is
reasonable.

\end{tcolorbox}

\bookmarksetup{startatroot}

\chapter{Bayesian Hypothesis Testing}\label{bayesian-hypothesis-testing}

This section explains how Bayesian inference turns hypothesis testing
into posterior probability comparison.

In classical testing, we control Type I error, study power, and often
use p-values or likelihood ratio tests. In Bayesian testing, the
posterior distribution directly assigns probabilities to hypotheses.

\section{Posterior probability tests}\label{posterior-probability-tests}

This subsection defines Bayesian tests through posterior probabilities
of the null and alternative parameter regions.

Let the hypotheses be \[H_0:\theta\in\Theta_0,
\qquad
H_1:\theta\in\Theta_0^c.\] Given posterior \(\pi(\theta\mid x)\),
\[\mathbb{P}(H_0\mid x)=\mathbb{P}(\theta\in\Theta_0\mid x)=\int_{\Theta_0}\pi(\theta\mid x)\,d\theta,\]
and
\[\mathbb{P}(H_1\mid x)=\mathbb{P}(\theta\in\Theta_0^c\mid x)=\int_{\Theta_0^c}\pi(\theta\mid x)\,d\theta.\]

\begin{tcolorbox}[enhanced jigsaw, arc=.35mm, bottomrule=.15mm, bottomtitle=1mm, breakable, colback=white, colbacktitle=quarto-callout-note-color!10!white, colframe=quarto-callout-note-color-frame, coltitle=black, left=2mm, leftrule=.75mm, opacityback=0, opacitybacktitle=0.6, rightrule=.15mm, title=\textcolor{quarto-callout-note-color}{\faInfo}\hspace{0.5em}{Definition}, titlerule=0mm, toprule=.15mm, toptitle=1mm]

\textbf{Definition 25} (Bayesian posterior probability test). A simple
Bayesian decision rule rejects \(H_0\) when
\[\mathbb{P}(H_0\mid x)<\mathbb{P}(H_1\mid x),\] equivalently, when
\[\mathbb{P}(H_0\mid x)<\frac12.\] A more conservative rule rejects
\(H_0\) when \[\mathbb{P}(H_0\mid x)<0.05.\]

\end{tcolorbox}

\begin{tcolorbox}[enhanced jigsaw, arc=.35mm, bottomrule=.15mm, bottomtitle=1mm, breakable, colback=white, colbacktitle=quarto-callout-tip-color!10!white, colframe=quarto-callout-tip-color-frame, coltitle=black, left=2mm, leftrule=.75mm, opacityback=0, opacitybacktitle=0.6, rightrule=.15mm, title=\textcolor{quarto-callout-tip-color}{\faLightbulb}\hspace{0.5em}{Example}, titlerule=0mm, toprule=.15mm, toptitle=1mm]

\textbf{Example 26} (Bayesian one-sided normal mean test). Suppose
\[X_1,\ldots,X_n\mid \mu\overset{\text{iid}}{\sim}\operatorname{Normal}(\mu,\sigma^2),\]
where \(\sigma^2\) is known, and suppose
\[\mu\sim\operatorname{Normal}(\theta,\tau^2).\] Test \[H_0:\mu\le \mu_0
\qquad\text{versus}\qquad
H_1:\mu>\mu_0.\] Find a posterior-probability rejection rule.

\end{tcolorbox}

\begin{tcolorbox}[enhanced jigsaw, arc=.35mm, bottomrule=.15mm, bottomtitle=1mm, breakable, colback=white, colbacktitle=quarto-callout-caution-color!10!white, colframe=quarto-callout-caution-color-frame, coltitle=black, left=2mm, leftrule=.75mm, opacityback=0, opacitybacktitle=0.6, rightrule=.15mm, title=\textcolor{quarto-callout-caution-color}{\faFire}\hspace{0.5em}{Solution}, titlerule=0mm, toprule=.15mm, toptitle=1mm]

The posterior distribution is
\[\mu\mid x\sim\operatorname{Normal}(m_n,v_n),\] where
\[m_n=\frac{\tau^2\sum_{i=1}^n x_i+\sigma^2\theta}{n\tau^2+\sigma^2},
\qquad
v_n=\frac{\sigma^2\tau^2}{\sigma^2+n\tau^2}.\] The posterior probability
of the null is \[\mathbb{P}(H_0\mid x)=\mathbb{P}(\mu\le \mu_0\mid x)
=\Phi\left(\frac{\mu_0-m_n}{\sqrt{v_n}}\right).\] If we reject when
\(\mathbb{P}(H_0\mid x)<1/2\), then because the posterior is normal and
symmetric, this is equivalent to \[m_n>\mu_0.\] That is,
\[\frac{\tau^2\sum_{i=1}^n x_i+\sigma^2\theta}{n\tau^2+\sigma^2}>\mu_0.\]
Equivalently,
\[\overline x>\mu_0+\frac{\sigma^2(\mu_0-\theta)}{n\tau^2}.\]

\end{tcolorbox}

\section{Bayesian tests with loss
functions}\label{bayesian-tests-with-loss-functions}

This subsection connects Bayesian tests to error costs.

Suppose the action space is \(\{a_0,a_1\}\), where \(a_0\) means
``accept or fail to reject \(H_0\)'' and \(a_1\) means ``reject
\(H_0\)''. A generalized 0--1 loss is \[L(\theta,a_0)=
\begin{cases}
0, & \theta\in\Theta_0,\\
c_{II}, & \theta\in\Theta_0^c,
\end{cases}\] and \[L(\theta,a_1)=
\begin{cases}
c_I, & \theta\in\Theta_0,\\
0, & \theta\in\Theta_0^c.
\end{cases}\] Here \(c_I\) is the cost of a Type I error and \(c_{II}\)
is the cost of a Type II error.

\begin{tcolorbox}[enhanced jigsaw, arc=.35mm, bottomrule=.15mm, bottomtitle=1mm, breakable, colback=white, colbacktitle=quarto-callout-note-color!10!white, colframe=quarto-callout-note-color-frame, coltitle=black, left=2mm, leftrule=.75mm, opacityback=0, opacitybacktitle=0.6, rightrule=.15mm, title=\textcolor{quarto-callout-note-color}{\faInfo}\hspace{0.5em}{Theorem}, titlerule=0mm, toprule=.15mm, toptitle=1mm]

\textbf{Theorem 27} (Bayes test under generalized 0--1 loss).
\emph{Under the above loss, reject \(H_0\) if
\[c_I\mathbb{P}(H_0\mid x)<c_{II}\mathbb{P}(H_1\mid x).\] Equivalently,
\[\mathbb{P}(H_0\mid x)<\frac{c_{II}}{c_I+c_{II}}.\]}

\end{tcolorbox}

\begin{proof}
\emph{Proof.} The posterior expected loss of accepting \(H_0\) is
\[\rho(a_0\mid x)=c_{II}\mathbb{P}(H_1\mid x).\] The posterior expected
loss of rejecting \(H_0\) is
\[\rho(a_1\mid x)=c_I\mathbb{P}(H_0\mid x).\] We reject when
\(\rho(a_1\mid x)<\rho(a_0\mid x)\), namely
\[c_I\mathbb{P}(H_0\mid x)<c_{II}\mathbb{P}(H_1\mid x).\] Since
\(\mathbb{P}(H_1\mid x)=1-\mathbb{P}(H_0\mid x)\), this is equivalent to
\[\mathbb{P}(H_0\mid x)<\frac{c_{II}}{c_I+c_{II}}.\]~◻
\end{proof}

\begin{tcolorbox}[enhanced jigsaw, arc=.35mm, bottomrule=.15mm, bottomtitle=1mm, breakable, colback=white, colbacktitle=quarto-callout-note-color!10!white, colframe=quarto-callout-note-color-frame, coltitle=black, left=2mm, leftrule=.75mm, opacityback=0, opacitybacktitle=0.6, rightrule=.15mm, title=\textcolor{quarto-callout-note-color}{\faInfo}\hspace{0.5em}{Remark}, titlerule=0mm, toprule=.15mm, toptitle=1mm]

\emph{Remark 28} (Connection to classical risk). In classical testing,
the risk under generalized 0--1 loss is \[R(\theta,\delta)=
\begin{cases}
c_I\beta(\theta), & \theta\in\Theta_0,\\
c_{II}[1-\beta(\theta)], & \theta\in\Theta_0^c,
\end{cases}\] where
\(\beta(\theta)=\mathbb{P}_\theta(\text{reject }H_0)\) is the power
function. Bayesian testing averages this risk using the posterior or
prior distribution.

\end{tcolorbox}

\bookmarksetup{startatroot}

\chapter{Bayesian Interval
Estimation}\label{bayesian-interval-estimation}

This section introduces Bayesian credible intervals and contrasts them
with frequentist confidence intervals.

In frequentist confidence intervals, the interval is random and the
parameter is fixed. In Bayesian credible intervals, the parameter is
random under the posterior distribution, so probability statements about
\(\theta\) belonging to an observed interval are meaningful within the
model.

\begin{tcolorbox}[enhanced jigsaw, arc=.35mm, bottomrule=.15mm, bottomtitle=1mm, breakable, colback=white, colbacktitle=quarto-callout-note-color!10!white, colframe=quarto-callout-note-color-frame, coltitle=black, left=2mm, leftrule=.75mm, opacityback=0, opacitybacktitle=0.6, rightrule=.15mm, title=\textcolor{quarto-callout-note-color}{\faInfo}\hspace{0.5em}{Definition}, titlerule=0mm, toprule=.15mm, toptitle=1mm]

\textbf{Definition 29} (Credible set). A set \(C(x)\subseteq\Theta\) is
a \(100(1-\alpha)\%\) \textbf{credible set} if
\[\mathbb{P}(\theta\in C(x)\mid x)=\int_{C(x)}\pi(\theta\mid x)\,d\theta=1-\alpha.\]
If \(C(x)=[a,b]\), then \([a,b]\) is called a credible interval.

\end{tcolorbox}

\begin{tcolorbox}[enhanced jigsaw, arc=.35mm, bottomrule=.15mm, bottomtitle=1mm, breakable, colback=white, colbacktitle=quarto-callout-warning-color!10!white, colframe=quarto-callout-warning-color-frame, coltitle=black, left=2mm, leftrule=.75mm, opacityback=0, opacitybacktitle=0.6, rightrule=.15mm, title=\textcolor{quarto-callout-warning-color}{\faExclamationTriangle}\hspace{0.5em}{Warning}, titlerule=0mm, toprule=.15mm, toptitle=1mm]

Interpretation warning A frequentist \(95\%\) confidence interval does
not mean that there is a \(95\%\) probability that the fixed parameter
lies in the observed interval. A Bayesian \(95\%\) credible interval
does allow the statement: \[\mathbb{P}(\theta\in C(x)\mid x)=0.95.\]
This interpretation depends on the prior and the Bayesian model.

\end{tcolorbox}

\section{Equal-tail credible
intervals}\label{equal-tail-credible-intervals}

This subsection introduces credible intervals based on posterior
quantiles.

\begin{tcolorbox}[enhanced jigsaw, arc=.35mm, bottomrule=.15mm, bottomtitle=1mm, breakable, colback=white, colbacktitle=quarto-callout-note-color!10!white, colframe=quarto-callout-note-color-frame, coltitle=black, left=2mm, leftrule=.75mm, opacityback=0, opacitybacktitle=0.6, rightrule=.15mm, title=\textcolor{quarto-callout-note-color}{\faInfo}\hspace{0.5em}{Definition}, titlerule=0mm, toprule=.15mm, toptitle=1mm]

\textbf{Definition 30} (Equal-tail credible interval). A
\(100(1-\alpha)\%\) equal-tail credible interval \([a,b]\) satisfies
\[\mathbb{P}(\theta<a\mid x)=\frac{\alpha}{2},
\qquad
\mathbb{P}(\theta>b\mid x)=\frac{\alpha}{2}.\] Equivalently, \(a\) and
\(b\) are the posterior \(\alpha/2\) and \(1-\alpha/2\) quantiles.

\end{tcolorbox}

\begin{tcolorbox}[enhanced jigsaw, arc=.35mm, bottomrule=.15mm, bottomtitle=1mm, breakable, colback=white, colbacktitle=quarto-callout-tip-color!10!white, colframe=quarto-callout-tip-color-frame, coltitle=black, left=2mm, leftrule=.75mm, opacityback=0, opacitybacktitle=0.6, rightrule=.15mm, title=\textcolor{quarto-callout-tip-color}{\faLightbulb}\hspace{0.5em}{Example}, titlerule=0mm, toprule=.15mm, toptitle=1mm]

\textbf{Example 31} (Beta-binomial credible interval). Suppose
\(X\sim\operatorname{Binomial}(n,p)\),
\(p\sim\operatorname{Beta}(2,2)\), \(n=20\), and \(x=12\). Find the
posterior distribution and describe the \(95\%\) equal-tail credible
interval.

\end{tcolorbox}

\begin{tcolorbox}[enhanced jigsaw, arc=.35mm, bottomrule=.15mm, bottomtitle=1mm, breakable, colback=white, colbacktitle=quarto-callout-caution-color!10!white, colframe=quarto-callout-caution-color-frame, coltitle=black, left=2mm, leftrule=.75mm, opacityback=0, opacitybacktitle=0.6, rightrule=.15mm, title=\textcolor{quarto-callout-caution-color}{\faFire}\hspace{0.5em}{Solution}, titlerule=0mm, toprule=.15mm, toptitle=1mm]

The posterior is
\[p\mid X=12\sim\operatorname{Beta}(2+12,2+20-12)=\operatorname{Beta}(14,10).\]
The \(95\%\) equal-tail credible interval is
\[\left[q_{0.025},q_{0.975}\right],\] where \(q_c\) is the \(c\)th
quantile of the \(\operatorname{Beta}(14,10)\) distribution.
Numerically, this interval is approximately \[[0.385,0.768].\] This
means that, under the beta-binomial Bayesian model,
\[\mathbb{P}(0.385\le p\le 0.768\mid X=12)\approx 0.95.\]

\end{tcolorbox}

\begin{tcolorbox}[enhanced jigsaw, arc=.35mm, bottomrule=.15mm, bottomtitle=1mm, breakable, colback=white, colbacktitle=quarto-callout-tip-color!10!white, colframe=quarto-callout-tip-color-frame, coltitle=black, left=2mm, leftrule=.75mm, opacityback=0, opacitybacktitle=0.6, rightrule=.15mm, title=\textcolor{quarto-callout-tip-color}{\faLightbulb}\hspace{0.5em}{Example}, titlerule=0mm, toprule=.15mm, toptitle=1mm]

\textbf{Example 32} (Normal-normal credible interval). Suppose
\(X_1,\ldots,X_n\mid\mu\overset{\text{iid}}{\sim}\operatorname{Normal}(\mu,\sigma^2)\)
with known \(\sigma^2\), and
\(\mu\sim\operatorname{Normal}(\mu_0,\sigma_0^2)\). Find a
\(100(1-\alpha)\%\) credible interval for \(\mu\).

\end{tcolorbox}

\begin{tcolorbox}[enhanced jigsaw, arc=.35mm, bottomrule=.15mm, bottomtitle=1mm, breakable, colback=white, colbacktitle=quarto-callout-caution-color!10!white, colframe=quarto-callout-caution-color-frame, coltitle=black, left=2mm, leftrule=.75mm, opacityback=0, opacitybacktitle=0.6, rightrule=.15mm, title=\textcolor{quarto-callout-caution-color}{\faFire}\hspace{0.5em}{Solution}, titlerule=0mm, toprule=.15mm, toptitle=1mm]

The posterior is \[\mu\mid x\sim\operatorname{Normal}(m_n,v_n),\] where
\[m_n=\frac{\sigma_0^2\sum_{i=1}^n x_i+\sigma^2\mu_0}{n\sigma_0^2+\sigma^2},
\qquad
v_n=\frac{\sigma^2\sigma_0^2}{\sigma^2+n\sigma_0^2}.\] Therefore a
\(100(1-\alpha)\%\) credible interval is
\[\left[m_n-z_{\alpha/2}\sqrt{v_n},\;m_n+z_{\alpha/2}\sqrt{v_n}\right],\]
where \(z_{\alpha/2}\) satisfies \(\mathbb{P}(Z>z_{\alpha/2})=\alpha/2\)
for \(Z\sim\operatorname{Normal}(0,1)\).

\end{tcolorbox}

\begin{tcolorbox}[enhanced jigsaw, arc=.35mm, bottomrule=.15mm, bottomtitle=1mm, breakable, colback=white, colbacktitle=quarto-callout-tip-color!10!white, colframe=quarto-callout-tip-color-frame, coltitle=black, left=2mm, leftrule=.75mm, opacityback=0, opacitybacktitle=0.6, rightrule=.15mm, title=\textcolor{quarto-callout-tip-color}{\faLightbulb}\hspace{0.5em}{Example}, titlerule=0mm, toprule=.15mm, toptitle=1mm]

\textbf{Example 33} (Gamma-Poisson credible interval). Suppose
\(X_1,\ldots,X_n\mid\lambda\overset{\text{iid}}{\sim}\operatorname{Poisson}(\lambda)\)
and \(\lambda\sim\operatorname{Gamma}(\alpha,\beta)\) in the rate
parametrization. Derive the posterior credible interval for \(\lambda\).

\end{tcolorbox}

\begin{tcolorbox}[enhanced jigsaw, arc=.35mm, bottomrule=.15mm, bottomtitle=1mm, breakable, colback=white, colbacktitle=quarto-callout-caution-color!10!white, colframe=quarto-callout-caution-color-frame, coltitle=black, left=2mm, leftrule=.75mm, opacityback=0, opacitybacktitle=0.6, rightrule=.15mm, title=\textcolor{quarto-callout-caution-color}{\faFire}\hspace{0.5em}{Solution}, titlerule=0mm, toprule=.15mm, toptitle=1mm]

The posterior is
\[\lambda\mid x\sim \operatorname{Gamma}\left(\alpha+\sum_{i=1}^n x_i,\beta+n\right).\]
A \(100(1-\alpha_0)\%\) equal-tail credible interval is
\[\left[q_{\alpha_0/2},q_{1-\alpha_0/2}\right],\] where \(q_c\) is the
\(c\)th quantile of the posterior gamma distribution. For example, if
\(n=10\), \(\sum x_i=26\), and \(\lambda\sim\operatorname{Gamma}(2,1)\),
then \[\lambda\mid x\sim\operatorname{Gamma}(28,11),\] and the interval
is obtained from the quantiles of \(\operatorname{Gamma}(28,11)\).

\end{tcolorbox}

\bookmarksetup{startatroot}

\chapter{Highest Posterior Density
Regions}\label{highest-posterior-density-regions}

This section explains the shortest Bayesian credible regions for
unimodal posterior distributions.

Equal-tail intervals are easy to compute, but they are not always the
shortest credible intervals. For a unimodal posterior density, the
shortest credible region is formed by keeping the parameter values with
highest posterior density.

\begin{tcolorbox}[enhanced jigsaw, arc=.35mm, bottomrule=.15mm, bottomtitle=1mm, breakable, colback=white, colbacktitle=quarto-callout-note-color!10!white, colframe=quarto-callout-note-color-frame, coltitle=black, left=2mm, leftrule=.75mm, opacityback=0, opacitybacktitle=0.6, rightrule=.15mm, title=\textcolor{quarto-callout-note-color}{\faInfo}\hspace{0.5em}{Definition}, titlerule=0mm, toprule=.15mm, toptitle=1mm]

\textbf{Definition 34} (Highest posterior density region). A
\(100(1-\alpha)\%\) \textbf{highest posterior density} (HPD) credible
region is a set
\[C_{\mathrm{HPD}}(x)=\{\theta:\pi(\theta\mid x)\ge k\},\] where \(k\)
is chosen so that
\[\int_{C_{\mathrm{HPD}}(x)}\pi(\theta\mid x)\,d\theta=1-\alpha.\]

\end{tcolorbox}

\begin{tcolorbox}[enhanced jigsaw, arc=.35mm, bottomrule=.15mm, bottomtitle=1mm, breakable, colback=white, colbacktitle=quarto-callout-note-color!10!white, colframe=quarto-callout-note-color-frame, coltitle=black, left=2mm, leftrule=.75mm, opacityback=0, opacitybacktitle=0.6, rightrule=.15mm, title=\textcolor{quarto-callout-note-color}{\faInfo}\hspace{0.5em}{Theorem}, titlerule=0mm, toprule=.15mm, toptitle=1mm]

\textbf{Theorem 35} (Shortest credible interval for unimodal posterior).
\emph{If \(\pi(\theta\mid x)\) is unimodal, then the shortest credible
interval with posterior probability \(1-\alpha\) is the HPD interval.
Its endpoints have equal posterior density, except possibly at a
boundary of the parameter space.}

\end{tcolorbox}

\begin{tcolorbox}[enhanced jigsaw, arc=.35mm, bottomrule=.15mm, bottomtitle=1mm, breakable, colback=white, colbacktitle=quarto-callout-tip-color!10!white, colframe=quarto-callout-tip-color-frame, coltitle=black, left=2mm, leftrule=.75mm, opacityback=0, opacitybacktitle=0.6, rightrule=.15mm, title=\textcolor{quarto-callout-tip-color}{\faLightbulb}\hspace{0.5em}{Example}, titlerule=0mm, toprule=.15mm, toptitle=1mm]

\textbf{Example 36} (Poisson HPD region). Suppose
\[X_1,\ldots,X_n\mid\lambda\overset{\text{iid}}{\sim}\operatorname{Poisson}(\lambda),\]
and use a conjugate gamma prior. If
\[\lambda\mid x\sim\operatorname{Gamma}\left(a+\sum x_i,\, n+b\right)\]
under a rate parametrization, then a \(100(1-\alpha)\%\) HPD credible
region has the form \[\{\lambda:\pi(\lambda\mid x)\ge k\},\] where \(k\)
is chosen to make the posterior probability equal to \(1-\alpha\).

\end{tcolorbox}

\begin{tcolorbox}[enhanced jigsaw, arc=.35mm, bottomrule=.15mm, bottomtitle=1mm, breakable, colback=white, colbacktitle=quarto-callout-caution-color!10!white, colframe=quarto-callout-caution-color-frame, coltitle=black, left=2mm, leftrule=.75mm, opacityback=0, opacitybacktitle=0.6, rightrule=.15mm, title=\textcolor{quarto-callout-caution-color}{\faFire}\hspace{0.5em}{Solution}, titlerule=0mm, toprule=.15mm, toptitle=1mm]

The posterior density is a gamma density. If it is unimodal, the HPD
region contains the highest-density values around the posterior mode.
The cutoff \(k\) is determined by solving
\[\int_{\{\lambda:\pi(\lambda\mid x)\ge k\}}\pi(\lambda\mid x)\,d\lambda=1-\alpha.\]
For example, in the lecture notes case with \(a=b=1\), \(n=10\), and
\(\sum x_i=6\), the posterior is gamma-shaped and the \(90\%\) HPD
credible set is approximately \[[0.253,1.005],\] while a corresponding
equal-tail interval is slightly longer.

\end{tcolorbox}

\begin{tcolorbox}[enhanced jigsaw, arc=.35mm, bottomrule=.15mm, bottomtitle=1mm, breakable, colback=white, colbacktitle=quarto-callout-note-color!10!white, colframe=quarto-callout-note-color-frame, coltitle=black, left=2mm, leftrule=.75mm, opacityback=0, opacitybacktitle=0.6, rightrule=.15mm, title=\textcolor{quarto-callout-note-color}{\faInfo}\hspace{0.5em}{Remark}, titlerule=0mm, toprule=.15mm, toptitle=1mm]

\emph{Remark 37} (Equal-tail versus HPD). Equal-tail intervals split
posterior probability equally between the two tails. HPD intervals
minimize length by taking the most plausible parameter values first. For
symmetric unimodal posteriors, the equal-tail and HPD intervals often
coincide. For skewed posteriors, they are usually different.

\end{tcolorbox}

\bookmarksetup{startatroot}

\chapter{Bayesian Optimality for
Intervals}\label{bayesian-optimality-for-intervals}

This section connects credible intervals to loss-function optimality.

Bayesian interval estimation can be framed as a decision problem: the
action is choosing a set \(C\), and the loss penalizes long intervals
while rewarding coverage of the true parameter.

\begin{tcolorbox}[enhanced jigsaw, arc=.35mm, bottomrule=.15mm, bottomtitle=1mm, breakable, colback=white, colbacktitle=quarto-callout-note-color!10!white, colframe=quarto-callout-note-color-frame, coltitle=black, left=2mm, leftrule=.75mm, opacityback=0, opacitybacktitle=0.6, rightrule=.15mm, title=\textcolor{quarto-callout-note-color}{\faInfo}\hspace{0.5em}{Definition}, titlerule=0mm, toprule=.15mm, toptitle=1mm]

\textbf{Definition 38} (Interval loss function). One simple loss
function for choosing a confidence or credible set \(C\) is
\[L(\theta,C)=b\cdot \operatorname{Length}(C)-\mathbb{1}\{\theta\in C\},\]
where \(b>0\) controls the tradeoff between short length and high
coverage.

\end{tcolorbox}

The corresponding frequentist risk is
\[R(\theta,C)=b\mathbb{E}_\theta[\operatorname{Length}(C(X))]-\mathbb{P}_\theta(\theta\in C(X)).\]
The Bayesian posterior expected loss is
\[\rho(C\mid x)=b\cdot \operatorname{Length}(C)-\mathbb{P}(\theta\in C\mid x).\]

\begin{tcolorbox}[enhanced jigsaw, arc=.35mm, bottomrule=.15mm, bottomtitle=1mm, breakable, colback=white, colbacktitle=quarto-callout-tip-color!10!white, colframe=quarto-callout-tip-color-frame, coltitle=black, left=2mm, leftrule=.75mm, opacityback=0, opacitybacktitle=0.6, rightrule=.15mm, title=\textcolor{quarto-callout-tip-color}{\faLightbulb}\hspace{0.5em}{Key idea}, titlerule=0mm, toprule=.15mm, toptitle=1mm]

Interpretation Large \(b\) prioritizes shorter intervals. Small \(b\)
prioritizes posterior coverage. This makes precise the tradeoff between
precision and uncertainty.

\end{tcolorbox}

\begin{tcolorbox}[enhanced jigsaw, arc=.35mm, bottomrule=.15mm, bottomtitle=1mm, breakable, colback=white, colbacktitle=quarto-callout-tip-color!10!white, colframe=quarto-callout-tip-color-frame, coltitle=black, left=2mm, leftrule=.75mm, opacityback=0, opacitybacktitle=0.6, rightrule=.15mm, title=\textcolor{quarto-callout-tip-color}{\faLightbulb}\hspace{0.5em}{Example}, titlerule=0mm, toprule=.15mm, toptitle=1mm]

\textbf{Example 39} (Normal interval risk). Suppose
\(X\sim\operatorname{Normal}(\mu,\sigma^2)\) with known \(\sigma^2\),
and consider symmetric intervals
\[C(X)=[X-c\sigma,X+c\sigma], \qquad c\ge 0.\] Compute the risk under
\[L(\mu,C)=b\cdot\operatorname{Length}(C)-\mathbb{1}\{\mu\in C\}.\]

\end{tcolorbox}

\begin{tcolorbox}[enhanced jigsaw, arc=.35mm, bottomrule=.15mm, bottomtitle=1mm, breakable, colback=white, colbacktitle=quarto-callout-caution-color!10!white, colframe=quarto-callout-caution-color-frame, coltitle=black, left=2mm, leftrule=.75mm, opacityback=0, opacitybacktitle=0.6, rightrule=.15mm, title=\textcolor{quarto-callout-caution-color}{\faFire}\hspace{0.5em}{Solution}, titlerule=0mm, toprule=.15mm, toptitle=1mm]

The length is \[\operatorname{Length}(C)=2c\sigma.\] The coverage
probability is \[\mathbb{P}_\mu(\mu\in C(X))
=\mathbb{P}_\mu(X-c\sigma\le\mu\le X+c\sigma)
=\mathbb{P}\left(-c\le \frac{X-\mu}{\sigma}\le c\right)
=2\Phi(c)-1.\] Therefore the risk is \[R(c)=b(2c\sigma)-[2\Phi(c)-1].\]
Differentiating, \[R'(c)=2b\sigma-2\phi(c).\] The optimum satisfies
\[\phi(c)=b\sigma.\] Since \(\phi(c)=\frac{1}{\sqrt{2\pi}}e^{-c^2/2}\),
if \(b\sigma\le 1/\sqrt{2\pi}\), then
\[c=\sqrt{-2\log(b\sigma\sqrt{2\pi})}.\] If \(b\sigma>1/\sqrt{2\pi}\),
the minimum occurs at \(c=0\), corresponding to a point estimate.

\end{tcolorbox}

\bookmarksetup{startatroot}

\chapter{Bayesian Inference Workflow}\label{bayesian-inference-workflow}

This section summarizes the steps of Bayesian inference as a reusable
procedure.

\begin{tcolorbox}[enhanced jigsaw, arc=.35mm, bottomrule=.15mm, bottomtitle=1mm, breakable, colback=white, colbacktitle=quarto-callout-tip-color!10!white, colframe=quarto-callout-tip-color-frame, coltitle=black, left=2mm, leftrule=.75mm, opacityback=0, opacitybacktitle=0.6, rightrule=.15mm, title=\textcolor{quarto-callout-tip-color}{\faLightbulb}\hspace{0.5em}{Key idea}, titlerule=0mm, toprule=.15mm, toptitle=1mm]

Bayesian workflow

\begin{enumerate}
\def\labelenumi{\arabic{enumi}.}
\item
  Specify the sampling model \(f(x\mid\theta)\).
\item
  Choose a prior distribution \(\pi(\theta)\).
\item
  Compute the posterior distribution
  \(\pi(\theta\mid x)\propto f(x\mid\theta)\pi(\theta)\).
\item
  For point estimation, report posterior mean, median, MAP, or another
  loss-optimal estimate.
\item
  For testing, compute posterior probabilities of \(H_0\) and \(H_1\),
  then choose an action based on posterior risk.
\item
  For interval estimation, report an equal-tail credible interval or HPD
  credible interval.
\item
  Interpret the answer conditional on the chosen model and prior.
\end{enumerate}

\end{tcolorbox}

\begin{tcolorbox}[enhanced jigsaw, arc=.35mm, bottomrule=.15mm, bottomtitle=1mm, breakable, colback=white, colbacktitle=quarto-callout-warning-color!10!white, colframe=quarto-callout-warning-color-frame, coltitle=black, left=2mm, leftrule=.75mm, opacityback=0, opacitybacktitle=0.6, rightrule=.15mm, title=\textcolor{quarto-callout-warning-color}{\faExclamationTriangle}\hspace{0.5em}{Warning}, titlerule=0mm, toprule=.15mm, toptitle=1mm]

Prior sensitivity Bayesian inference depends on the prior distribution.
With large samples, the likelihood often dominates the prior. With small
samples, the prior can strongly influence posterior estimates, tests,
and intervals.

\end{tcolorbox}

\bookmarksetup{startatroot}

\chapter{Practice Problems}\label{practice-problems-10}

This section provides practice problems that connect the Bayesian
methods from Sections 12--17.

\begin{tcolorbox}[enhanced jigsaw, arc=.35mm, bottomrule=.15mm, bottomtitle=1mm, breakable, colback=white, colbacktitle=quarto-callout-important-color!10!white, colframe=quarto-callout-important-color-frame, coltitle=black, left=2mm, leftrule=.75mm, opacityback=0, opacitybacktitle=0.6, rightrule=.15mm, title=\textcolor{quarto-callout-important-color}{\faExclamation}\hspace{0.5em}{Practice Problem}, titlerule=0mm, toprule=.15mm, toptitle=1mm]

\textbf{Practice Problem 40} (Beta-binomial posterior and estimator).
Suppose \(X\sim\operatorname{Binomial}(30,p)\) and \(x=18\). Let
\(p\sim\operatorname{Beta}(4,6)\).

\begin{enumerate}
\def\labelenumi{\arabic{enumi}.}
\item
  Find the posterior distribution of \(p\).
\item
  Find the posterior mean.
\item
  Find the MAP estimator, assuming the posterior parameters are both
  larger than 1.
\end{enumerate}

\end{tcolorbox}

\begin{tcolorbox}[enhanced jigsaw, arc=.35mm, bottomrule=.15mm, bottomtitle=1mm, breakable, colback=white, colbacktitle=quarto-callout-caution-color!10!white, colframe=quarto-callout-caution-color-frame, coltitle=black, left=2mm, leftrule=.75mm, opacityback=0, opacitybacktitle=0.6, rightrule=.15mm, title=\textcolor{quarto-callout-caution-color}{\faFire}\hspace{0.5em}{Solution}, titlerule=0mm, toprule=.15mm, toptitle=1mm]

\begin{enumerate}
\def\labelenumi{\arabic{enumi}.}
\item
  The posterior is
  \[p\mid x\sim\operatorname{Beta}(4+18,6+30-18)=\operatorname{Beta}(22,18).\]
\item
  The posterior mean is
  \[\mathbb{E}(p\mid x)=\frac{22}{22+18}=\frac{22}{40}=0.55.\]
\item
  The MAP estimator is
  \[\widehat p_{\mathrm{MAP}}=\frac{22-1}{22+18-2}=\frac{21}{38}\approx 0.5526.\]
\end{enumerate}

\end{tcolorbox}

\begin{tcolorbox}[enhanced jigsaw, arc=.35mm, bottomrule=.15mm, bottomtitle=1mm, breakable, colback=white, colbacktitle=quarto-callout-important-color!10!white, colframe=quarto-callout-important-color-frame, coltitle=black, left=2mm, leftrule=.75mm, opacityback=0, opacitybacktitle=0.6, rightrule=.15mm, title=\textcolor{quarto-callout-important-color}{\faExclamation}\hspace{0.5em}{Practice Problem}, titlerule=0mm, toprule=.15mm, toptitle=1mm]

\textbf{Practice Problem 41} (Normal-normal posterior). Suppose
\(X_1,\ldots,X_{16}\mid\mu\overset{\text{iid}}{\sim}\operatorname{Normal}(\mu,4)\)
and \(\overline x=10\). Let \(\mu\sim\operatorname{Normal}(8,9)\).

\begin{enumerate}
\def\labelenumi{\arabic{enumi}.}
\item
  Find the posterior variance \(v_n\).
\item
  Find the posterior mean \(m_n\).
\item
  Give a \(95\%\) credible interval for \(\mu\).
\end{enumerate}

\end{tcolorbox}

\begin{tcolorbox}[enhanced jigsaw, arc=.35mm, bottomrule=.15mm, bottomtitle=1mm, breakable, colback=white, colbacktitle=quarto-callout-caution-color!10!white, colframe=quarto-callout-caution-color-frame, coltitle=black, left=2mm, leftrule=.75mm, opacityback=0, opacitybacktitle=0.6, rightrule=.15mm, title=\textcolor{quarto-callout-caution-color}{\faFire}\hspace{0.5em}{Solution}, titlerule=0mm, toprule=.15mm, toptitle=1mm]

Here \(n=16\), \(\sigma^2=4\), \(\mu_0=8\), and \(\sigma_0^2=9\).

\begin{enumerate}
\def\labelenumi{\arabic{enumi}.}
\item
  \[v_n=\left(\frac{16}{4}+\frac{1}{9}\right)^{-1}
  =\left(4+\frac19\right)^{-1}
  =\frac{9}{37}.\]
\item
  \[m_n=v_n\left(\frac{16\cdot 10}{4}+\frac{8}{9}\right)
  =\frac{9}{37}\left(40+\frac{8}{9}\right)
  =\frac{368}{37}\approx 9.946.\]
\item
  A \(95\%\) credible interval is \[m_n\pm 1.96\sqrt{v_n}
  =9.946\pm 1.96\sqrt{\frac{9}{37}}.\] Since
  \(\sqrt{9/37}\approx 0.493\), the interval is approximately
  \[[8.980,10.912].\]
\end{enumerate}

\end{tcolorbox}

\begin{tcolorbox}[enhanced jigsaw, arc=.35mm, bottomrule=.15mm, bottomtitle=1mm, breakable, colback=white, colbacktitle=quarto-callout-important-color!10!white, colframe=quarto-callout-important-color-frame, coltitle=black, left=2mm, leftrule=.75mm, opacityback=0, opacitybacktitle=0.6, rightrule=.15mm, title=\textcolor{quarto-callout-important-color}{\faExclamation}\hspace{0.5em}{Practice Problem}, titlerule=0mm, toprule=.15mm, toptitle=1mm]

\textbf{Practice Problem 42} (Gamma-Poisson posterior). Suppose
\(X_1,\ldots,X_8\mid\lambda\overset{\text{iid}}{\sim}\operatorname{Poisson}(\lambda)\)
and \(\sum x_i=20\). Let \(\lambda\sim\operatorname{Gamma}(3,2)\) using
the rate parametrization.

\begin{enumerate}
\def\labelenumi{\arabic{enumi}.}
\item
  Find the posterior distribution.
\item
  Find the posterior mean.
\item
  Find the MAP estimator, assuming the posterior shape is larger than 1.
\end{enumerate}

\end{tcolorbox}

\begin{tcolorbox}[enhanced jigsaw, arc=.35mm, bottomrule=.15mm, bottomtitle=1mm, breakable, colback=white, colbacktitle=quarto-callout-caution-color!10!white, colframe=quarto-callout-caution-color-frame, coltitle=black, left=2mm, leftrule=.75mm, opacityback=0, opacitybacktitle=0.6, rightrule=.15mm, title=\textcolor{quarto-callout-caution-color}{\faFire}\hspace{0.5em}{Solution}, titlerule=0mm, toprule=.15mm, toptitle=1mm]

\begin{enumerate}
\def\labelenumi{\arabic{enumi}.}
\item
  The posterior is
  \[\lambda\mid x\sim\operatorname{Gamma}(3+20,2+8)=\operatorname{Gamma}(23,10).\]
\item
  The posterior mean is \[\mathbb{E}(\lambda\mid x)=\frac{23}{10}=2.3.\]
\item
  For a gamma distribution with shape \(a\) and rate \(b\), the mode is
  \((a-1)/b\) when \(a>1\). Thus
  \[\widehat\lambda_{\mathrm{MAP}}=\frac{23-1}{10}=2.2.\]
\end{enumerate}

\end{tcolorbox}

\begin{tcolorbox}[enhanced jigsaw, arc=.35mm, bottomrule=.15mm, bottomtitle=1mm, breakable, colback=white, colbacktitle=quarto-callout-important-color!10!white, colframe=quarto-callout-important-color-frame, coltitle=black, left=2mm, leftrule=.75mm, opacityback=0, opacitybacktitle=0.6, rightrule=.15mm, title=\textcolor{quarto-callout-important-color}{\faExclamation}\hspace{0.5em}{Practice Problem}, titlerule=0mm, toprule=.15mm, toptitle=1mm]

\textbf{Practice Problem 43} (Bayesian test with posterior probability).
Suppose the posterior distribution of a parameter is
\[\theta\mid x\sim\operatorname{Normal}(1.2,0.25).\] Test
\[H_0:\theta\le 0
\qquad\text{versus}\qquad
H_1:\theta>0.\] Compute \(\mathbb{P}(H_0\mid x)\) and decide whether to
reject \(H_0\) using the rule \(\mathbb{P}(H_0\mid x)<0.05\).

\end{tcolorbox}

\begin{tcolorbox}[enhanced jigsaw, arc=.35mm, bottomrule=.15mm, bottomtitle=1mm, breakable, colback=white, colbacktitle=quarto-callout-caution-color!10!white, colframe=quarto-callout-caution-color-frame, coltitle=black, left=2mm, leftrule=.75mm, opacityback=0, opacitybacktitle=0.6, rightrule=.15mm, title=\textcolor{quarto-callout-caution-color}{\faFire}\hspace{0.5em}{Solution}, titlerule=0mm, toprule=.15mm, toptitle=1mm]

The posterior standard deviation is \(\sqrt{0.25}=0.5\). Therefore
\[\mathbb{P}(H_0\mid x)=\mathbb{P}(\theta\le 0\mid x)
=\Phi\left(\frac{0-1.2}{0.5}\right)
=\Phi(-2.4).\] Using the standard normal table,
\[\Phi(-2.4)\approx 0.0082.\] Since \(0.0082<0.05\), we reject \(H_0\)
using this Bayesian posterior-probability rule.

\end{tcolorbox}

\begin{tcolorbox}[enhanced jigsaw, arc=.35mm, bottomrule=.15mm, bottomtitle=1mm, breakable, colback=white, colbacktitle=quarto-callout-important-color!10!white, colframe=quarto-callout-important-color-frame, coltitle=black, left=2mm, leftrule=.75mm, opacityback=0, opacitybacktitle=0.6, rightrule=.15mm, title=\textcolor{quarto-callout-important-color}{\faExclamation}\hspace{0.5em}{Practice Problem}, titlerule=0mm, toprule=.15mm, toptitle=1mm]

\textbf{Practice Problem 44} (Bayesian decision rule with unequal
costs). For testing \(H_0:\theta\in\Theta_0\) versus
\(H_1:\theta\in\Theta_0^c\), suppose \[\mathbb{P}(H_0\mid x)=0.30.\] The
cost of Type I error is \(c_I=5\), and the cost of Type II error is
\(c_{II}=1\). Should we reject \(H_0\)?

\end{tcolorbox}

\begin{tcolorbox}[enhanced jigsaw, arc=.35mm, bottomrule=.15mm, bottomtitle=1mm, breakable, colback=white, colbacktitle=quarto-callout-caution-color!10!white, colframe=quarto-callout-caution-color-frame, coltitle=black, left=2mm, leftrule=.75mm, opacityback=0, opacitybacktitle=0.6, rightrule=.15mm, title=\textcolor{quarto-callout-caution-color}{\faFire}\hspace{0.5em}{Solution}, titlerule=0mm, toprule=.15mm, toptitle=1mm]

Reject \(H_0\) if
\[\mathbb{P}(H_0\mid x)<\frac{c_{II}}{c_I+c_{II}}=\frac{1}{5+1}=\frac16\approx 0.1667.\]
Here \(\mathbb{P}(H_0\mid x)=0.30>0.1667\), so we do not reject \(H_0\).
The high cost of Type I error makes the rejection rule more
conservative.

\end{tcolorbox}

\begin{tcolorbox}[enhanced jigsaw, arc=.35mm, bottomrule=.15mm, bottomtitle=1mm, breakable, colback=white, colbacktitle=quarto-callout-important-color!10!white, colframe=quarto-callout-important-color-frame, coltitle=black, left=2mm, leftrule=.75mm, opacityback=0, opacitybacktitle=0.6, rightrule=.15mm, title=\textcolor{quarto-callout-important-color}{\faExclamation}\hspace{0.5em}{Practice Problem}, titlerule=0mm, toprule=.15mm, toptitle=1mm]

\textbf{Practice Problem 45} (Equal-tail credible interval). Suppose the
posterior distribution is
\[\theta\mid x\sim\operatorname{Normal}(5,4).\] Find a \(90\%\)
equal-tail credible interval.

\end{tcolorbox}

\begin{tcolorbox}[enhanced jigsaw, arc=.35mm, bottomrule=.15mm, bottomtitle=1mm, breakable, colback=white, colbacktitle=quarto-callout-caution-color!10!white, colframe=quarto-callout-caution-color-frame, coltitle=black, left=2mm, leftrule=.75mm, opacityback=0, opacitybacktitle=0.6, rightrule=.15mm, title=\textcolor{quarto-callout-caution-color}{\faFire}\hspace{0.5em}{Solution}, titlerule=0mm, toprule=.15mm, toptitle=1mm]

The posterior standard deviation is \(2\). For a \(90\%\) interval,
\(z_{0.05}\approx 1.645\). Thus \[5\pm 1.645(2)=5\pm 3.29.\] The
credible interval is \[[1.71,8.29].\]

\end{tcolorbox}

\begin{tcolorbox}[enhanced jigsaw, arc=.35mm, bottomrule=.15mm, bottomtitle=1mm, breakable, colback=white, colbacktitle=quarto-callout-important-color!10!white, colframe=quarto-callout-important-color-frame, coltitle=black, left=2mm, leftrule=.75mm, opacityback=0, opacitybacktitle=0.6, rightrule=.15mm, title=\textcolor{quarto-callout-important-color}{\faExclamation}\hspace{0.5em}{Practice Problem}, titlerule=0mm, toprule=.15mm, toptitle=1mm]

\textbf{Practice Problem 46} (HPD interval concept). Suppose a posterior
density is unimodal and skewed to the right. Explain why the \(95\%\)
HPD interval may differ from the \(95\%\) equal-tail interval.

\end{tcolorbox}

\begin{tcolorbox}[enhanced jigsaw, arc=.35mm, bottomrule=.15mm, bottomtitle=1mm, breakable, colback=white, colbacktitle=quarto-callout-caution-color!10!white, colframe=quarto-callout-caution-color-frame, coltitle=black, left=2mm, leftrule=.75mm, opacityback=0, opacitybacktitle=0.6, rightrule=.15mm, title=\textcolor{quarto-callout-caution-color}{\faFire}\hspace{0.5em}{Solution}, titlerule=0mm, toprule=.15mm, toptitle=1mm]

The equal-tail interval places \(2.5\%\) posterior probability in each
tail. The HPD interval instead contains the parameter values with
highest posterior density and chooses the cutoff so that the total
posterior probability is \(95\%\). For a skewed posterior, equal tails
may include some low-density values in the long tail while excluding
higher-density values on the other side. Therefore the HPD interval is
usually shorter and has endpoints with equal posterior density, while
the equal-tail interval is based on posterior quantiles.

\end{tcolorbox}

\bookmarksetup{startatroot}

\chapter{Summary}\label{summary-12}

This section summarizes the role of Bayesian inference across the main
statistical tasks of the course.

\begin{tcolorbox}[enhanced jigsaw, arc=.35mm, bottomrule=.15mm, bottomtitle=1mm, breakable, colback=white, colbacktitle=quarto-callout-tip-color!10!white, colframe=quarto-callout-tip-color-frame, coltitle=black, left=2mm, leftrule=.75mm, opacityback=0, opacitybacktitle=0.6, rightrule=.15mm, title=\textcolor{quarto-callout-tip-color}{\faLightbulb}\hspace{0.5em}{Key idea}, titlerule=0mm, toprule=.15mm, toptitle=1mm]

Section summary

\begin{enumerate}
\def\labelenumi{\arabic{enumi}.}
\item
  Bayesian inference starts with a prior \(\pi(\theta)\) and likelihood
  \(f(x\mid\theta)\).
\item
  Bayes' rule gives the posterior:
  \[\pi(\theta\mid x)\propto f(x\mid\theta)\pi(\theta).\]
\item
  Point estimates can be posterior means, medians, or MAP estimates,
  depending on the loss function.
\item
  Conjugate priors make posterior computation simple.
\item
  Bayesian tests compare posterior probabilities of hypotheses, possibly
  weighted by error costs.
\item
  Bayesian credible intervals give direct posterior probability
  statements.
\item
  HPD regions are shortest credible regions for unimodal posterior
  distributions.
\item
  Bayesian decisions are naturally derived by minimizing posterior
  expected loss.
\end{enumerate}

\end{tcolorbox}

{\def\LTcaptype{none} % do not increment counter
\begin{longtable}[]{@{}lll@{}}
\toprule\noalign{}
Task & Bayesian object & Common answer \\
\midrule\noalign{}
\endhead
\bottomrule\noalign{}
\endlastfoot
Point estimation & Posterior distribution & Posterior mean, median,
MAP \\
Testing & Posterior hypothesis probability & Reject if posterior risk is
smaller \\
Interval estimation & Posterior credible probability & Equal-tail or HPD
credible interval \\
Model updating & Prior and likelihood & Posterior distribution \\
Decision theory & Posterior expected loss & Bayes action \\
\end{longtable}
}

\bookmarksetup{startatroot}

\chapter{References and Further
Reading}\label{references-and-further-reading}

This page collects the main references, open resources, software
documentation, and chapter-by-chapter reading suggestions for
\textbf{Foundations of Probability and Statistical Theory}.

The book is designed to be self-contained for MATH 5010, but students
are encouraged to consult the references below for additional examples,
exercises, and alternative explanations.

\section{Main References}\label{main-references}

\subsection{Primary statistical theory
references}\label{primary-statistical-theory-references}

\begin{enumerate}
\def\labelenumi{\arabic{enumi}.}
\item
  Casella, G., \& Berger, R. L. (2002). \emph{Statistical Inference}
  (2nd ed.). Duxbury.
\item
  DeGroot, M. H., \& Schervish, M. J. (2012). \emph{Probability and
  Statistics} (4th ed.). Pearson.
\item
  Hogg, R. V., McKean, J. W., \& Craig, A. T. (2019). \emph{Introduction
  to Mathematical Statistics} (8th ed.). Pearson.
\end{enumerate}

\subsection{Probability references}\label{probability-references}

\begin{enumerate}
\def\labelenumi{\arabic{enumi}.}
\setcounter{enumi}{3}
\item
  Grinstead, C. M., \& Snell, J. L. (2012). \emph{Introduction to
  Probability}. American Mathematical Society.
\item
  Ross, S. M. (2019). \emph{A First Course in Probability} (10th ed.).
  Pearson.
\item
  Durrett, R. (2019). \emph{Probability: Theory and Examples} (5th ed.).
  Cambridge University Press.
\end{enumerate}

\subsection{Open educational
resources}\label{open-educational-resources}

\begin{enumerate}
\def\labelenumi{\arabic{enumi}.}
\setcounter{enumi}{6}
\item
  LibreTexts Statistics. \emph{Statistics Library}.\\
  \url{https://stats.libretexts.org/Bookshelves}
\item
  Blitzstein, J. K., \& Hwang, J. (2019). \emph{Introduction to
  Probability} (2nd ed.). CRC Press.
\item
  Wasserman, L. (2004). \emph{All of Statistics: A Concise Course in
  Statistical Inference}. Springer.
\end{enumerate}

\section{Computational References}\label{computational-references}

\subsection{Python and scientific
computing}\label{python-and-scientific-computing}

\begin{enumerate}
\def\labelenumi{\arabic{enumi}.}
\item
  NumPy Documentation.\\
  \url{https://numpy.org/doc/}
\item
  SciPy Documentation.\\
  \url{https://docs.scipy.org/doc/scipy/}
\item
  pandas Documentation.\\
  \url{https://pandas.pydata.org/docs/}
\item
  Matplotlib Documentation.\\
  \url{https://matplotlib.org/stable/contents.html}
\item
  statsmodels Documentation.\\
  \url{https://www.statsmodels.org/stable/index.html}
\item
  scikit-learn Documentation.\\
  \url{https://scikit-learn.org/stable/}
\end{enumerate}

\subsection{Bayesian computation}\label{bayesian-computation}

\begin{enumerate}
\def\labelenumi{\arabic{enumi}.}
\setcounter{enumi}{6}
\item
  PyMC Documentation.\\
  \url{https://www.pymc.io/projects/docs/en/stable/}
\item
  Stan Documentation.\\
  \url{https://mc-stan.org/users/documentation/}
\end{enumerate}

These Bayesian computation tools are not required for the core theory,
but they are useful for students who want to extend Chapter 20 into
modern Bayesian modeling and posterior simulation.

\section{Quarto, Jupyter, and Google
Colab}\label{quarto-jupyter-and-google-colab}

\begin{enumerate}
\def\labelenumi{\arabic{enumi}.}
\item
  Quarto Documentation.\\
  \url{https://quarto.org/docs/books/}
\item
  Jupyter Notebook Documentation.\\
  \url{https://docs.jupyter.org/}
\item
  Google Colab.\\
  \url{https://colab.research.google.com/}
\end{enumerate}

The computer labs in this book are written as Jupyter notebooks. They
can be run locally or opened in Google Colab through the links on the
\hyperref[computer-labs-and-interactive-pages]{Labs and Interactive HTML
Pages} page.

\section{Additional Advanced
References}\label{additional-advanced-references}

\begin{enumerate}
\def\labelenumi{\arabic{enumi}.}
\item
  Lehmann, E. L., \& Casella, G. (1998). \emph{Theory of Point
  Estimation} (2nd ed.). Springer.
\item
  Lehmann, E. L., \& Romano, J. P. (2005). \emph{Testing Statistical
  Hypotheses} (3rd ed.). Springer.
\item
  Robert, C. P., \& Casella, G. (2004). \emph{Monte Carlo Statistical
  Methods} (2nd ed.). Springer.
\item
  Liu, J. S. (2001). \emph{Monte Carlo Strategies in Scientific
  Computing}. Springer.
\item
  Gelman, A., Carlin, J. B., Stern, H. S., Dunson, D. B., Vehtari, A.,
  \& Rubin, D. B. (2013). \emph{Bayesian Data Analysis} (3rd ed.). CRC
  Press.
\item
  Bernardo, J. M., \& Smith, A. F. M. (1994). \emph{Bayesian Theory}.
  Wiley.
\item
  Montgomery, D. C. (2017). \emph{Design and Analysis of Experiments}
  (9th ed.). Wiley.
\end{enumerate}

\section{Notes on Attribution}\label{notes-on-attribution}

These lecture notes, labs, and interactive pages are prepared for MATH
5010. The structure follows the course sequence from probability
foundations to modern statistical inference, computation, and Bayesian
methods. External references are listed to support further study and are
not required unless assigned by the instructor.




\end{document}
